\documentclass[11pt,letterpaper]{article}
\usepackage[T1]{fontenc}
\usepackage{lmodern,amsmath,amssymb,amsthm,microtype,xcolor,xurl,booktabs,enumitem}
\usepackage[letterpaper,left=1.02in,right=1.02in,top=.92in,bottom=.95in]{geometry}
\usepackage{multicol}
\usepackage{longtable}
\usepackage[breakable,skins]{tcolorbox}
\newtcolorbox{altbox}{breakable,enhanced jigsaw,colback=white,colframe=black!75,boxrule=0.7pt,arc=0pt,
  left=5pt,right=5pt,top=3pt,bottom=3pt,colbacktitle=black!10,coltitle=black,fonttitle=\bfseries\small,
  title={Alternative to the Lean proof}}
\newtcolorbox{illbox}{breakable,enhanced jigsaw,colback=white,colframe=black!45,boxrule=0.4pt,arc=0pt,
  left=5pt,right=5pt,top=3pt,bottom=3pt,colbacktitle=black!4,coltitle=black,fonttitle=\bfseries\small,
  title={Illustration --- not used in the proof}}
\usepackage{hyperref}
\definecolor{source}{HTML}{395E6F}
\hypersetup{hypertexnames=false,colorlinks=true,linkcolor=black,urlcolor=source,pdftitle={The Alpoge-Claude proof of the Riemann hypothesis explained for the Riemann zeta function}}
\newtheorem*{lemma}{Lemma}
\newtheorem{theorem}{Theorem}[section]
\newtheorem*{prop}{Proposition}
\newtheorem*{thmA}{Estimate A}
\newtheorem*{thmB}{Estimate B}
\newtheorem*{thmC}{Estimate C}
\newcommand{\Z}{\mathbb Z}\newcommand{\Q}{\mathbb Q}\newcommand{\C}{\mathbb C}
\newcommand{\re}{\operatorname{Re}}
\newcommand{\quoted}{\textsf{[quoted]}}
\setlength{\parskip}{4pt plus 1pt}
\setlist{itemsep=2pt,topsep=3pt}
\renewcommand{\thesection}{I.\arabic{section}}
\renewcommand{\theequation}{I.\arabic{equation}}
\title{\vspace{-11.5mm}The Alp\"oge--Claude proof of the Riemann hypothesis\\[2pt] explained for the Riemann zeta function}
\author{Claude Fable 5.2, Codex Astra 6.2, and Gemini Carbon\\[2pt] at the behest of Michael O. Rubinstein}
\date{Working draft, 22 November 2026}
\newcounter{partno}\setcounter{partno}{1}
\renewcommand{\theHequation}{\arabic{partno}.\arabic{equation}}
\renewcommand{\theHsection}{\arabic{partno}.\arabic{section}}
\renewcommand{\theHtheorem}{\arabic{partno}.\arabic{section}.\arabic{theorem}}
\begin{document}

\maketitle\thispagestyle{empty}
\vspace{-14mm}


\vspace{3mm}
\begin{center}\begin{minipage}{.86\textwidth}\small
\textbf{Abstract.} We explain the argument of Alp\"oge and Claude for the Riemann hypothesis in the case of
the Riemann zeta function. The analytic step shows that upper bounds $G(x)\ll_\varepsilon x^{1/2+\varepsilon}$, for every $\varepsilon>0$,
for a smoothed sum $G$ over prime powers, valid throughout intervals $[Y,Y^{3/2}]$ with $Y$ arbitrarily large,
exclude zeros to the right of the critical line. We derive these bounds from the main inequality (M), an upper
bound for the mean of a weighted sum over prime powers after subtraction of its main term, and two estimates
for Dirichlet characters. We prove the two estimates, for suitably chosen auxiliary primes, from theorems in the literature, and explain how the
divisibility and the size of theta lifts enter (M). We quote (M) from the Lean development and prove a conditional reduction of it.
This is a work in progress. Part III reduces (M), and with it the Riemann hypothesis, to a list of numbered
hypotheses, each stated in classical terms (\S\ref{as:sec}); none of them is proved outright here, and the last
includes that the primes with which the construction is made satisfy the first estimate. It proves
parts of the argument behind several of them, works the simplest case, the lemniscate, in full
(\S\ref{el:sec}), and lists the 195 statements that the Lean development takes under the heading of
published results (\S\ref{cs:sec}). (M) itself is not proved here.
\end{minipage}\end{center}

\vspace{1mm}
\begin{center}\fbox{\parbox{.92\textwidth}{\small\textbf{Note.} This is a working guide to Alp\"oge's
Lean-verified proof of the Riemann hypothesis, specialised to the Riemann zeta function. Its goal is a
human-readable exposition of the proof in the style of classical analytic number theory. It is a work in
progress, written as an aid to studying the proof.

\smallskip\emph{What is in these pages.} Parts I and II deduce the Riemann hypothesis from one inequality, called (M)
throughout, and two estimates for Dirichlet characters. Part III proves the two estimates for suitably chosen auxiliary primes (\S III.1--III.2). It
reduces (M), and with it the Riemann hypothesis, to thirteen numbered hypotheses, each stated in classical terms
(\S\ref{as:sec}); the thirteenth includes that the primes with which the construction is made satisfy the first
estimate. It defines the objects of the construction (\S\ref{ob:sec}), constructs the two sets of columns outright
(\S III.3), and proves parts of the argument behind several of the hypotheses, in an abstract setting or from
named statements of the Lean development (\S III.4--III.9). The lemniscate illustrates the size and divisibility
calculations; the theorems stated for it are proved from classical results (\S\ref{el:sec}). And it lists the 195 statements that
the Lean development takes under the heading of published results (\S\ref{cs:sec}). Labelled boxes mark selected passages
that give a proof different from the Lean development's, or that are illustrations not used in the proof;
smaller departures and models are noted where they occur.

\smallskip\emph{What is not.} None of the thirteen hypotheses is proved outright here. Behind them stand, taken from
the Lean development: the integrality of the theta kernels; the inner-product formula for the theta lifts, the
covolumes of the lattices and the local indices; the local data at the level primes, the count of forms and the
covolume of a lattice of forms; the trace-formula estimates and the size and subspace conditions of
Corollary~T2; the bound for Galois cohomology and the comparison inequality that links it to the
denominators; the existence of the objects of the construction; and the supply of all these data at every
scale of an interval. The bibliographic attributions proposed in the census (\S\ref{cs:sec}) are from memory.

\smallskip The tables at the start of Part III and in \S\ref{as:sec} say, statement by statement, what is proved in
these pages, what is proved from named statements, and what is taken without proof from the Lean development.
\S I.12 lists what remains to be written up.}}\end{center}


{\small\setcounter{tocdepth}{1}\makeatletter\renewcommand{\l@part}[2]{\par\addvspace{4pt}\penalty-200\noindent\textbf{#1}\par\nobreak}\renewcommand{\@pnumwidth}{1.6em}\renewcommand{\l@section}{\@dottedtocline{1}{0em}{2.7em}}\makeatother
\setlength{\columnsep}{22pt}\begin{multicols}{2}[\noindent\textbf{Contents} {\footnotesize (entries are links)}]\makeatletter\@starttoc{toc}\makeatother\end{multicols}}

\noindent\textbf{What this is.} In November 2026 Levent Alp\"oge announced a proof, found with the AI system
Claude, of the Riemann hypothesis. It exists as a formal proof in the Lean proof assistant, about 17 million
lines long. This paper explains how the argument works for the Riemann zeta function, for a reader who knows
a first course in analytic number theory: the prime number theorem, Dirichlet characters, and Mellin
transforms.

\noindent\textbf{How it is organised.} Part I is an exposition: the strategy first, then each step with its
motivation and with examples computed on real primes. Part II gives the details and proofs to which Part I
refers. A reader can stop after Part I and know how the proof goes; Part II is there to check each step.
Part III, which is being written, is to prove the statements that Parts I and II quote. It proves Estimates B and C
and constructs the two auxiliary sets of residue classes (the ``sets of columns'' of \S I.5). It reduces the main inequality (M),
and with it the Riemann hypothesis, to a list of numbered hypotheses, each stated in classical terms
(\S\ref{as:sec}); none of them is proved outright here. It proves parts of the argument behind several of them,
some conditionally on statements of the Lean development which it names, and takes the
rest from that development; the table at the start of Part III says which. It works the simplest case, the
lemniscate, in full (\S\ref{el:sec}). A census lists the 195 statements that the Lean development takes under
the heading of published results (\S\ref{cs:sec}).

\noindent\textbf{What is proved here and what is not.} The argument rests on three statements. One
of them, the main inequality (M), is the arithmetic heart of the matter. Two further estimates (B and C) concern
auxiliary Dirichlet characters. We state B and C exactly, and (M) in the form in which it is used, and explain why one should expect them. That the
three together imply the Riemann hypothesis is proved in Part II. Estimates B and C are proved in Part III,
from theorems on Dirichlet $L$-functions that are cited there. But (M) is not yet proved here: we are working through its proof in the Lean development and adding the details to Part III section by section.
Of what it needs, Part III so far constructs the two sets of columns (\S III.3) and proves: the analytic part of the
bound for the size of the theta kernels, with four conditional passages (\S III.4); a local divisibility estimate for power series at
inert primes, under the integrality of the theta kernels in the coordinates of the height-two logarithm
(hypothesis (B)), for orbits of length two and with a further duality input for longer orbits (\S III.5); the lemmas on
lattices that carry this estimate to lattices and part of the sum over primes, the
main theorem there being conditional on two statements of the Lean development (\S III.6); that the size inequality
follows from eleven named inputs (\S III.7); that the dimension bound follows from three statements of the Lean
development on traces, ranks and pairings and the size and subspace conditions of Corollary~T2 (\S III.8); and
the step from the bound for Galois cohomology and a quoted comparison inequality to a bound for
denominators (\S III.9); and that (M), at one sufficiently large scale,
follows from twelve numbered hypotheses, and the Riemann hypothesis from these and a thirteenth, which concerns
the scales and the auxiliary primes (\S\ref{as:sec}). None of the thirteen hypotheses is proved outright here.
Part II proves the elementary
lemmas. The rest is for now
\emph{quoted} or assumed: hypothesis (B) and the statement on valuations built on it, the arithmetic inputs of the size estimate,
the trace-formula estimates themselves, the bound for Galois cohomology itself and the comparison inequality, the construction of the maps and their
weights, and those of the hypotheses of \S\ref{as:sec} that no section treats. The simplest case, the lemniscate, is
worked in full in \S\ref{el:sec}; \S\ref{cs:sec} is a census of the 195 statements that the Lean development takes
under the heading of published results. Statements marked \quoted\ are taken
from the Lean development without proof. The full proof of (M) is where almost all of the formal
development goes, and writing it up is the work in progress; \S I.12 lists what remains. Every numerical example, except where stated, can be reproduced with the scripts listed
under Computations, after Part III, or named where the numbers occur.

\bigskip
\begin{center}\Large\bfseries Part I. How the proof works\end{center}
\phantomsection\addcontentsline{toc}{part}{Part I. How the proof works}

\section*{I.0\quad The strategy in eight steps}\addcontentsline{toc}{section}{I.0\quad The strategy in eight steps}

\textbf{Chebyshev's idea about counting primes.} Chebyshev \cite{Che52} bounded the number of primes using
the prime factorization of factorials. In the simplest form of the
argument, due to Erd\H{o}s \cite{Erd32}, one considers the middle binomial coefficient $\binom{2n}{n}$. It is
divisible by every prime between $n$ and $2n$, and it is at most $4^n$. So the product of
those primes is at most $4^n$. That is an upper bound for a count of primes, obtained from two facts about
one integer: what divides it, and how big it is.

\textbf{One piece of notation.} For a real number $v$ we write $v_+$ for its \emph{positive part}:
$v_+=v$ if $v>0$, and $v_+=0$ otherwise. For example $5_+=5$ and $(-3)_+=0$. A sum of terms of the form
$(\cdots)_+$ is automatically cut off where the terms would become negative.

\textbf{The plan.} The proof of the Riemann hypothesis described here is Chebyshev's argument, run with a far
better ``integer''. There are eight steps. After each we name the main ingredients, and say where it is
treated and whether it is proved in this paper.

\begin{enumerate}[leftmargin=2em]
\item \textbf{Zeros to primes.} The Riemann hypothesis follows from one-sided bounds
$G(x)\le x^{1/2+\varepsilon}$, for every $\varepsilon>0$, for a certain smoothed count of prime powers $G(x)$, and it is enough to have
it on a sparse sequence of long intervals.
\emph{Ingredients:} a variant of Landau's theorem on Mellin transforms of functions of constant sign; the
Mellin transform \eqref{F} of $G$ and Chebyshev's bound. \emph{Where:} \S I.1; proved in Part II, \S1.
\item \textbf{Chebyshev's principle, for lattices.} If $N\ne0$ and $p^{e_p}\mid N$ for each of several distinct primes $p$, then
$\prod_pp^{e_p}\le|N|$. To bound primes from above one compares the divisibility of an object with its size.
In the proof the object is not one integer but a linear map between lattices of large rank.
\emph{Ingredients:} Chebyshev's method; the slope inequality for a map of lattices (Bost), and the
filtration of a lattice by slopes (Stuhler, Grayson). \emph{Where:} \S I.2, \S I.7; the inequality is proved
in the form needed in Part III, \S\ref{as:sec}.
\item \textbf{Divisibility with no floor functions.} For $n!$ the exponent of $p$ is
$\sum_h\lfloor n/p^h\rfloor$, and the rounding in the floor functions is why the classical method gives
constants and not asymptotics. Derivatives of order $a$ of suitable functions attached to complex
multiplication have a smooth lower bound for their divisibility, up to one rounding, at the primes that stay
prime (the \emph{inert} primes): for the lemniscate sine and every inert prime $p$,
$v_p(A_a)\ge\lfloor ap/(p^2-1)\rfloor>a(1/p+1/p^3+\cdots)-1$ (\S\ref{el:sec}). At the other primes there is no
such systematic gain. For theta series it is not automatic: the simplest one fails (\S I.4.1). It holds for the
particular theta series of the proof, in particular coordinates, and that is the statement taken from the
Lean development. Subtracting a weight $w$ turns each term $a/p^h$ into
$(a/p^h-w)_+=w\,(x/p^h-1)_+$ with $x=a/w$: the kernel that $G$ uses, magnified by $w$.
\emph{Ingredients:} formal groups of height two (Lubin and Tate, Honda); the reduction of curves with
complex multiplication (Deuring). \emph{Where:} \S I.3, \S I.4; Part III, \S III.5 (conditionally on the integrality hypothesis (B) and, for longer orbits, a
duality statement, both taken from the Lean development) and \S III.6 (in part; its main theorem is
conditional on two further statements, and the bounds for the losses are quoted). The simplest case, the
lemniscate, is proved in full in \S\ref{el:sec}.
\item \textbf{Size, through the linear term.} For the comparison to be useful the size must be known not
just to leading order $a\log a$ but through the term linear in $a$. The map is a chain of $m$ theta lifts
between unitary groups; the kernels are written in classical terms here for the definite ones
(\S\ref{ob:sec}). A bound for its norm, together with the covolume estimates, identifies
the main term through the term linear in $a$, and the linear terms on the two sides agree.
\emph{Ingredients:} Gaussian integrals ($\|z^a\|^2=a!$); the period formula (Kronecker's limit formula;
Lerch, Chowla and Selberg); the mass formula (Siegel, Shimura) and the Siegel--Weil formula; the inner
product formula for theta lifts (Rallis); Weyl's dimension formula. \emph{Where:} \S I.4.2, \S I.7;
Part III, \S III.4 and \S III.7, from inputs that are quoted.
\item \textbf{Denominators.} A chain of lifts need not map integral lattices to integral lattices. The
lattices are assembled prime by prime, using the freedom to twist the intermediate lifts by characters
without changing the final period. Comparing these choices costs an error at each lift. A quoted
comparison bounds the total error at each lift in terms of the lengths of certain Galois cohomology groups
of the field, generalisations of class groups, and Theorem~G of \S III.9 bounds their weighted sum.
\emph{Ingredients:} the seesaw identity (Kudla) and Bessel periods; Galois representations attached to
automorphic forms; Iwasawa theory and $p$-adic Hodge theory. This is the part of the proof that uses the
most mathematics from outside analytic number theory. \emph{Where:} \S I.7; Part III, \S III.9, where the
bound itself is quoted.
\item \textbf{Rank.} Each source lattice has smaller rank than the target, and one map controls only the
largest slope of the target. A family of maps is used, and their images are shown to be spread evenly, so
that together they fill almost the whole target and the \emph{mean} slope is controlled.
\emph{Ingredients:} a relative trace formula (Jacquet), for a second moment of periods; a lemma of linear
algebra. \emph{Where:} \S I.7; Part III, \S III.8, from three quoted statements. Steps 2 to 6 together give
the main inequality (M) (\S I.5); it is deduced from numbered hypotheses in Part III, \S\ref{as:sec}.
\item \textbf{Balance.} With $m$ lifts there is an error $x/m$, made once at the two ends of the chain, and
an error $m$, accumulated over the steps. They balance at $m\approx\sqrt x$. This is where the exponent
$\frac12$ comes from. \emph{Where:} \S I.8.
\item \textbf{From inert primes to all primes.} The sum in (M) runs over prime powers in certain residue
classes, the cells of a table. Averaging over translates removes the column condition on the
distinguished cells. Their contribution gives $G(x)$ plus three sums twisted by quadratic characters of two
auxiliary primes $q$ and $d$, and for well-chosen $q$ and $d$ the three sums are small on long intervals;
Estimate~C controls the remaining cells.
\emph{Ingredients:} zero-free regions and the exceptional zero (Landau, Page); a zero-density estimate
(Montgomery); primes in progressions to large moduli (Thorner and Zaman); the prime number theorem in
progressions; a set of columns of the table chosen by a construction of Behrend. \emph{Where:} \S I.5,
\S I.6; Part II, \S3--5; proved in Part III, \S III.1--III.3, from cited theorems.
\end{enumerate}

\noindent In short: steps 1 and 8 are proved in this paper, from cited theorems of analytic number theory,
and step 7 is a choice of parameters. Steps 2 to 6 make up the main inequality (M), which is reduced to
numbered hypotheses and is not proved here.

\noindent\textbf{Logical skeleton.}
\begin{equation}\label{D:1}
\underbrace{\text{(M)}}_{\text{\quoted}},\ \underbrace{\text{B},\ \text{C}}_{\text{Part III}}
\ \xrightarrow{\ \S\mathrm{I}.6\ }\ \text{Estimate A}
\ \xrightarrow{\ \S\mathrm{I}.6\ }\ G(x)\le x^{1/2+\varepsilon}\ \text{on long intervals}
\ \xrightarrow{\ \S\mathrm{I}.1\ }\ \text{RH}.
\end{equation}
Three statements enter: the main inequality (M), which is quoted (Part III, \S\ref{as:sec} reduces it,
and with it the Riemann hypothesis, to thirteen numbered hypotheses; none is proved outright), and two estimates for auxiliary
Dirichlet characters, one at the scale $x$ (Estimate B) and one at a much larger scale (Estimate C), which
are proved in Part III from cited theorems. The order of
choices is this. For each small $\varepsilon$ there are intervals $[Z,Z^{3/2}]$ with $Z$ arbitrarily large
such that every $x$ in the interval has a pair of auxiliary primes $q,d$, chosen at a scale $Y$ with
$x^{1/J}\le Y\le x$ ($J\ge1$ depends on $\varepsilon$ only), for which Estimates B and C hold at $x$ and for
which the construction is made at $x$. The argument is for each $x$ separately. (Theorem~\ref{thm:pair}
gives one pair for all $x$ in $[Y,Y^{J}]$; that the construction can be made at each $x$ with such a pair,
with constants independent of $x$, is Hypothesis (H12) of Part III, \S\ref{as:sec}.)

\section{From zeros to a sum over primes}

\textbf{The statistic.} Let $\Lambda(n)=\log p$ if $n=p^h$ is a prime power and $0$ otherwise. Define
\begin{equation}\label{D:2}
 S(x)=\sum_{n\le x}\Lambda(n)\Big(\frac xn-1\Big),\qquad
 G(x)=S(x)-x\log x+(1+\gamma)x ,
\end{equation}
where $\gamma$ is Euler's constant. Equivalently $S(x)=\sum_n\Lambda(n)(x/n-1)_+$, where $v_+=\max(v,0)$: with the
positive part, the condition $n\le x$ under the sum can be omitted. Compared with Chebyshev's $\psi(x)=\sum_{n\le x}\Lambda(n)$, the sum
$S$ weights $n$ by the \emph{ramp} $(x/n-1)_+$: as a function of $n$ it is large for small $n$, decreases,
reaches $0$ at $n=x$, and is $0$ for $n>x$. Also
$S(x)=x\int_1^x\psi(t)\,dt/t^2$. The two subtracted terms are the main term of $S$, so $G$ is a remainder.

\textbf{Why this statistic.} The Mellin transform of the ramp is simple:
$\int_1^\infty (x/n-1)_+\,x^{-s-1}dx=n^{-s}/(s(s-1))$ for $\re s>1$. Summing over $n$ with weights
$\Lambda(n)$,
\begin{equation}\label{F}
 F(s)=\int_1^\infty G(x)x^{-s-1}dx
 =-\frac{\zeta'(s)}{\zeta(s)}\cdot\frac1{s(s-1)}-\frac1{(s-1)^2}+\frac{1+\gamma}{s-1}.
\end{equation}
The last two terms cancel the pole at $s=1$ (both terms of its principal part). What remains are simple poles of $F$ at the zeros $\rho$
of $\zeta$, with residue $-1/(\rho(\rho-1))$ at a simple zero, and a pole at $s=0$. (The trivial zeros
$-2,-4,\ldots$ also give poles; they produce the term $O(x^{-2})$ below.) Shifting the contour in
the inverse transform gives the explicit formula (compare \cite[Ch.~17]{Dav00})
\begin{equation}\label{D:4}
 G(x)=\log 2\pi-\sum_\rho\frac{x^\rho}{\rho(\rho-1)}+O(x^{-2}),
\end{equation}
the sum being over the nontrivial zeros. It converges absolutely, because of the factor $\rho(\rho-1)$; that
is the benefit of smoothing. One has $\sum_\rho\frac1{\rho(1-\rho)}=2+\gamma-\log4\pi=0.0461\ldots$\,. If every zero has real part
$\frac12$ then $|\rho(\rho-1)|=\rho(1-\rho)$, and therefore, for large $x$,
$|G(x)-\log2\pi|\le 0.0462\sqrt x$. A zero with
real part $\beta>\frac12$ would contribute a term of size $x^\beta$.

\emph{A sanity check on real primes.} Summing over the prime powers up to $2\cdot10^6$ gives
$G(10^4)=1.749$, $G(10^5)=3.671$, $G(10^6)=11.040$, $G(2\cdot10^6)=-12.268$, to be compared with
$\log2\pi=1.838$ and $0.0462\sqrt{10^6}=46.2$.

\textbf{One side is enough.} The terms of a zero off the line and of its complex conjugate together form a real function of $x$ that
oscillates and takes both signs. So one expects a zero with real part $\beta$ to make $G(x)$
exceed $+cx^{\beta}$ for some arbitrarily large $x$. Thus an \emph{upper} bound for $G$ should already
exclude it. That an upper bound suffices follows from a theorem of Landau on Mellin transforms of functions of one sign \cite{Lan05}; see \cite[\S15.1]{MV07}.

\begin{lemma}[one-sided criterion] Let $0<\theta<1$. If $G(x)\le x^\theta$ for all large $x$, then
$\zeta(s)\ne0$ for $\re s>\theta$.
\end{lemma}
\begin{proof} The integral $H(s)=\int_1^\infty\big(x^\theta-G(x)\big)x^{-s-1}dx$ equals $1/(s-\theta)-F(s)$,
and its integrand is nonnegative for large $x$. Landau's theorem says that a Mellin transform of an
eventually nonnegative function has a singularity at the real point of its abscissa of convergence. But by
\eqref{F} the function $H$ is regular at every real $s>\theta$: the only real singularities of $F$ in
$s>0$ would come from real zeros of $\zeta$ in $(0,1)$, and there are none. Hence the abscissa is at most
$\theta$, the integral converges for $\re s>\theta$, and $F$ is holomorphic there. A zero of $\zeta$ in that
half-plane would be a pole of $F$.
\end{proof}

The arithmetic argument below produces the upper bound only for $x$ in certain long intervals. The lemma
survives this weakening.

\begin{lemma}[the same, on intervals] Let $0<\theta<1$. If there are arbitrarily large $Y$ such that
$G(x)\le x^\theta$ for all $x\in[Y,Y^{3/2}]$, then $\zeta(s)\ne0$ for $\re s>\theta$.
\end{lemma}

\noindent The proof replaces Landau's theorem by an estimate of high derivatives of $F$, which see only one
interval at a time; it is one page and is given in full in Part II, \S1.

\begin{altbox}
\emph{A second proof, from a published theorem on changes of sign.} Kaczorowski and Pintz
\cite[Theorem~1 and Corollary~1]{KP86} proved the following quantitative form of Landau's theorem for meromorphic transforms (a sharpening of a theorem of P\'olya). Let $f$ be real
on $(0,\infty)$, and let $\Phi(s)=\int_0^\infty f(x)x^{-s-1}dx$ converge absolutely for $\re s$ large.
Suppose that $\Phi$ is regular for $\re s>\Theta$ but in no half-plane $\re s>\Theta-\eta$ with $\eta>0$;
that $\Phi$ is meromorphic in $\re s\ge\Theta-c_0$ for some $c_0>0$; and that $\Phi$ is regular at the real
point $s=\Theta$. Then for every $\alpha>0$ and every $T\ge T_0(f,\alpha)$ the interval $[T^{1-\alpha},T]$ contains a
change of sign of $f$; in the sense of their definition this means two points of it at which $f$ is
nonzero and of opposite signs. (The third hypothesis is not displayed in the statement of
\cite{KP86}; it is the case in which their least boundary ordinate $\gamma$ is positive, with $\gamma=+\infty$
allowed, the only one their proof treats, and without it the
conclusion fails for $f(x)=x^\Theta$ ($x\ge1$), $f(x)=0$ ($x<1$), with $\Phi(s)=1/(s-\Theta)$. We read the
statement, the beginning of the proof and its last step (\cite[(3.29)--(3.30), p.~182]{KP86}), not the whole
proof.) To deduce the lemma, suppose that $\zeta$ has a zero with real part greater than $\theta$, and let
$\Theta\le1$ be the supremum of the real parts of the zeros. Put $f(x)=G(x)-2x^\theta$ for $x\ge1$ and
$f(x)=0$ for $x<1$. Then $\Phi(s)=F(s)-2/(s-\theta)$ for $\re s>1$, since $G(x)=O(x\log x)$. By \eqref{F},
$\Phi$ is meromorphic in the plane; its poles with $\re s>0$ are the zeros of $\zeta$ in the critical strip
and the point $\theta$, which lies to the left of $\Theta$. So $\Phi$ is regular for $\re s>\Theta$ and in
no larger half-plane. It is regular at the real point $\Theta$: there is no real zero of $\zeta$ in $(0,1)$,
and $F$ is regular at $s=1$. Hence $f$ takes a positive value in $[T^{1-\alpha},T]$ for every large $T$.
Take $\alpha=\frac14$ and $T=Y^{3/2}$: the interval $[Y^{9/8},Y^{3/2}]$ lies in $[Y,Y^{3/2}]$, where
$f(x)\le-x^\theta<0$. This is a contradiction when $Y$ is large. The argument does not need the supremum
$\Theta$ to be attained. It explains the shape of the intervals: any intervals $[Y,Y^{1+c}]$ with a fixed
$c>0$ would do.
The Lean development proves the lemma by high derivatives of $F$; Part~II, \S1 gives a version of that
argument.
\end{altbox}

\textbf{Conclusion of step 1.} To prove the Riemann hypothesis it suffices to show: for every
$\varepsilon>0$ there are arbitrarily large $Y$ with $G(x)\le x^{1/2+\varepsilon}$ throughout
$[Y,Y^{3/2}]$. (Zeros come in pairs $\rho$, $1-\rho$, so no zeros to the right of the line means none to the
left.) \emph{From now on the goal is an upper bound for a sum over primes.}

\section{Chebyshev's method, and where it stalls}

\textbf{The method.} Let $N=\binom{2n}{n}$. Every prime $p$ with $n<p\le2n$ divides $(2n)!$ and not
$(n!)^2$, so it divides $N$. Also $N\le(1+1)^{2n}=4^n$. A positive integer divisible by distinct primes is at
least their product, so
\begin{equation}\label{D:5}
 \sum_{n<p\le2n}\log p\ \le\ \log N\ \le\ 2n\log 2 .
\end{equation}
In general: \emph{if $\Delta$ is a nonzero integer and $p^{e_p}$ divides $\Delta$ for each prime $p$ in some
set, then}
\begin{equation}\label{det}
 \sum_p e_p\log p\ \le\ \log|\Delta| .
\end{equation}
Divisibility is a statement about each prime separately; size is a statement about the real numbers. The
inequality \eqref{det} is the bridge, and it always points the same way: it bounds primes from \emph{above}.
By step 1 that is the direction we need. This is the arithmetic analogue of Stepanov's count of zeros with
multiplicity; see Part II, \S2.

\textbf{The best classical integer.} For $\Delta=n!$ both sides are known exactly. Legendre's formula \cite{Leg08} says
the exponent of $p$ in $n!$ is $\sum_{h\ge1}\lfloor n/p^h\rfloor$. Hence
\begin{equation}\label{leg}
 \log n!=\sum_{r\le n}\Lambda(r)\Big\lfloor\frac nr\Big\rfloor .
\end{equation}
The left side is $n\log n-n+O(\log n)$ by Stirling. If the floor were not there, the right side would be
$n\sum_{r\le n}\Lambda(r)/r$, which is $S(n)+\psi(n)$: essentially our statistic.

\textbf{Where it stalls.} The floor \emph{is} there. Writing $\lfloor n/r\rfloor=n/r-\{n/r\}$, the
fractional parts contribute $\sum_{r\le n}\Lambda(r)\{n/r\}$, which is of size $n$: as large as the term
$-n$ on the left. All information about $\psi(n)-n$ is hidden in the fractional parts, and \eqref{leg}
cannot separate it. This is why Chebyshev's method yields $c_1n\le\psi(n)\le c_2n$ and no more.

\textbf{What would be needed.} A nonzero ``integer'' $\Delta$ such that:
\begin{enumerate}[label=(W\arabic*),leftmargin=3em]
\item the exponent of $p^h$ in $\Delta$ is a \emph{smooth} function of $p^h$, like the ramp $(x/p^h-1)_+$,
with no floor function, or with rounding that is negligible;
\item $\log|\Delta|$ is known with an error much smaller than its linear term, in fact smaller by a factor
of about $x^{-1/2}$, where $x$ is the scale at which the primes are tested (\S I.1);
\item the primes that occur are all primes, or a set of primes we can relate to all primes.
\end{enumerate}
The rest of the proof constructs such an object. No single integer does all three; the final object is a
linear map between lattices, and \eqref{det} becomes an inequality between the index of a lattice and the
norm of a map.

\section{A divisible number with smooth exponents}

\textbf{An example to compute with.} The lemniscate sine $\mathrm{sl}(z)$ is the inverse function of
$t\mapsto\int_0^t du/\sqrt{1-u^4}$. It satisfies $\mathrm{sl}''=-2\,\mathrm{sl}^3$, and from this its Taylor
coefficients can be computed exactly. Let $A_n=\mathrm{sl}^{(n)}(0)$, the $n$th derivative at the origin.
These are integers, studied by Hurwitz \cite{Hur99}; $A_n=0$ unless $n\equiv1\bmod4$ (both statements are
proved in \S\ref{el:sec}, Lemma~Le4):
\begin{equation}\label{D:8}
 A_1=1,\quad A_5=-12,\quad A_9=3024,\quad A_{13}=-4390848,\quad
 A_{17}=21224560896=2^8\cdot3^6\cdot7^2\cdot11\cdot211 .
\end{equation}
The curve $y^2=1-x^4$ has an extra symmetry $x\mapsto ix$. One says it has \emph{complex multiplication} by
$\Z[i]$. A prime $p\equiv3\bmod4$ remains prime in $\Z[i]$ (it is \emph{inert}); a prime $p\equiv1\bmod4$
factors (it \emph{splits}).

\begin{center}\small
\begin{tabular}{r|rr|rr|rr||cc}
\toprule
 & \multicolumn{6}{c||}{exponent of $p$ in $A_n$, and $n\,p/(p^2-1)$, for inert $p$}&\multicolumn{2}{c}{split $p$}\\
$n$ & $p=3$ & $3n/8$ & $p=7$ & $7n/48$ & $p=11$ & $11n/120$ & $p=5$ & $p=13$\\
\midrule
17 & 6 & 6.4 & 2 & 2.5 & 1 & 1.6 & 0 & 0\\
49 & 19 & 18.4 & 7 & 7.1 & 4 & 4.5 & 0 & 0\\
101 & 37 & 37.9 & 14 & 14.7 & 9 & 9.3 & 0 & 0\\
201 & 75 & 75.4 & 29 & 29.3 & 18 & 18.4 & 0 & 0\\
\bottomrule
\end{tabular}
\end{center}

\noindent For every inert prime $p$ and every $n$ the exponent of $p$ in $A_n$ is at least
$\lfloor n\,p/(p^2-1)\rfloor$, where
\begin{equation}\label{D:9}
 n\cdot\frac{p}{p^2-1}=n\Big(\frac1p+\frac1{p^3}+\frac1{p^5}+\cdots\Big)
\end{equation}
is the prediction. That is, the exponent is at least the prediction minus $1$. This is proved in
\S\ref{el:sec} (Theorem~Le18). In the range computed ($n\le2001$, $n\equiv1\bmod4$) equality with the
rounded prediction holds for most $n$. For the inert primes $3,7,11,19,23,31$ and $n\le201$ it exceeds $n\,p/(p^2-1)$ by at most $1.65$; this ceiling does not persist for larger $n$
(for $p=3$ the excess is $4.625$ at $n=1945$ and $5.625$ at $n=5833$). The split primes $5,13,17,29$ divide none of the $A_n$ with $n\equiv1\bmod4$ (proved for all such $n$ in Proposition~Le20a of \S\ref{el:sec}). (Split primes are not excluded
altogether: $157$ divides $A_{101}$. The point is that nothing forces them to divide. What is proved is that a split prime $p$ never divides $A_p$: Proposition~Le20 of \S\ref{el:sec}.)

\textbf{Why.} Near the origin the addition law of an elliptic curve is a \emph{formal group}: a power series
$X\oplus Y$ in two variables. Its logarithm $l(X)$ is the power series with $l(X\oplus Y)=l(X)+l(Y)$; here
$l(t)=\int_0^tdu/\sqrt{1-u^4}$ and $\mathrm{sl}=l^{-1}$. For a curve with complex multiplication and an
inert prime $p$, multiplication by $p$ on the formal group looks like $X\mapsto X^{p^2}$ modulo $p$, and the
logarithm behaves like the model series
\begin{equation}\label{D:10}
 l(X)=X+\frac{X^{p^2}}{p}+\frac{X^{p^4}}{p^2}+\cdots .
\end{equation}
(For the lemniscate both statements are proved in \S\ref{el:sec}, Theorem~Le13 and Corollary~Le16. There the
signs alternate: the series is $X-X^{p^2}/p+X^{p^4}/p^2-\cdots$, and multiplication by $p$ is
$X\mapsto-X^{p^2}$ modulo $p$, with no further terms (Corollary~Le16\,(ii)). This changes no valuation.)
(From here on we write $a$ for the order of the derivative, as in the rest of the paper.) For this model one proves by induction that the coefficient of $z^a$ in the inverse series has $p$-adic
valuation at least $-(a-1)/(p^2-1)$ (Part II, \S2; Proposition~Le15 of
\S\ref{el:sec} is the refinement that includes the factorial: $a!$ times this coefficient is divisible by
$p^{\,p(a-1)/(p^2-1)}$). An $a$th derivative
multiplies that coefficient by $a!$, whose valuation is $a/(p-1)-O(\log a)$. The difference is
\begin{equation}\label{D:11}
 a\Big(\frac1{p-1}-\frac1{p^2-1}\Big)=a\cdot\frac p{p^2-1}=a\Big(\frac1p+\frac1{p^3}+\cdots\Big).
\end{equation}
\emph{Only odd powers of $p$ appear.} That is the signature of an inert prime. At a split prime the formal group is of a different type (height one): its logarithm has the denominator
$p^j$ already at degree $p^j$ (for the lemniscate: Proposition~Le20 of \S\ref{el:sec}), the inverse series loses about $a/(p-1)$, and this cancels the whole
valuation of $a!$. There is no gain. That an elliptic curve with complex multiplication has the first type
at inert primes and the second at split primes is a theorem of Deuring \cite{Deu41}; for the shape of the
logarithm see Honda \cite{Hon70} and Lubin--Tate \cite{LT65}.

\textbf{What this achieves, and what it does not.} We have integers whose exponents are smooth in $p$:
wish (W1), up to a small amount of rounding. Three defects remain, and each drives a part of the real proof.
\begin{enumerate}[label=(D\arabic*),leftmargin=3em]
\item \emph{Rounding.} An error of size about $1$ in each exponent costs $\sum_{p\le n}\log p\approx n$ in \eqref{det}:
again as large as the linear term. This is the floor-function problem of \S I.2 in another guise.
\item \emph{Only half the primes.} Only inert primes appear, weighted by $1-\chi(p)$, that is, $2$ for inert $p$ and $0$ for split $p$, where $\chi$ is the
quadratic character of the field ($\chi(p)=-1$ for inert $p$).
\item \emph{Size.} $|A_n|=4\,n!\,R^{-n-1}(1+o(1))$ for $n\equiv1\bmod4$, with $R=1.854\ldots$ (Theorem~Le9 of \S\ref{el:sec}, where $R=\varpi/\sqrt2$ with the constant $\varpi$ of \S I.4.1), so $\log|A_n|\approx n\log n$. But the
inert primes account only for $\sum_{p\equiv3(4),\,p\le n}\frac np\log p=\frac12n\log n+O(n)$. The inequality
\eqref{det} is true and useless: it is off by a factor $2$ in the \emph{leading} term (Proposition~Le19 of \S\ref{el:sec}).
(At $n=101$: inert primes up to $n$ contribute $167.7$ of $\log|A_n|=306.8$.)
\end{enumerate}
For the objects of the real proof, the divisibility at inert primes is written up in \S III.5, conditionally on the integrality of the theta kernels in the coordinates of the height-two logarithm (hypothesis (B)), which is taken from the Lean development. For one explicit example, built from the lemniscate, hypothesis (B) is proved in \S\ref{el:sec}.

\section{Repairing the three defects}

The lemniscate numbers $A_n$ were derivatives of a function attached to an elliptic curve. The real proof
keeps the idea \emph{differentiate a function of arithmetic origin and look at what divides the result},
but changes the function. This section says which functions are differentiated, what
``differentiating'' means for them, and what each choice achieves. It has four parts: the functions
(\S I.4.1), then the repair of (D3), (D1) and (D2) in that order (\S I.4.2--I.4.4), and finally a list of what
the reader can and cannot check directly (\S I.4.5).

\subsection*{I.4.1\quad What is differentiated: theta series with a variable}

\textbf{The simplest case.} Let $\Lambda=\Z[i]\subset\C$, the lattice of Gaussian integers, and put
\begin{equation}\label{D:12}
 \Theta(z)=\sum_{\lambda\in\Lambda}\exp\big(-\pi|\lambda|^2+\lambda z\big)\qquad(z\in\C).
\end{equation}
The Gaussian factor makes the sum converge for every $z$, so $\Theta$ is an entire function. Its $n$th
derivative at the origin is
\begin{equation}\label{D:13}
 T_n=\Theta^{(n)}(0)=\sum_{\lambda\in\Lambda}\lambda^n e^{-\pi|\lambda|^2}.
\end{equation}
So \emph{differentiating $n$ times at the origin means inserting the weight $\lambda^n$ into the sum over the
lattice}. More generally, for a polynomial $P$ one has
$P(\tfrac{d}{dz})\Theta(0)=\sum_\lambda P(\lambda)e^{-\pi|\lambda|^2}$: a weighted count of lattice points.
Because $\Lambda$ is preserved by multiplication by $i$, the sum $T_n$ vanishes unless $4\mid n$.

\textbf{Three things visible in this example.} They can all be checked by direct computation (script
\texttt{e6\_hecke\_theta.py}).
\begin{enumerate}[leftmargin=2em]
\item \emph{A period.} $T_0=1.18034\ldots=\sqrt2\,\varpi/\pi$ (a classical evaluation; proved in \S\ref{el:sec}, Proposition~Le7), where
$\varpi=2\int_0^1dt/\sqrt{1-t^4}=2.62205\ldots$ is the lemniscate constant: the \emph{period} of the curve
$y^2=1-x^4$. And $T_4/T_0=0.12131\ldots=\frac14(\varpi/\pi)^4$ (Proposition~Le7).
\item \emph{Rationality.} After division by $T_0$ and by $(\varpi/\pi)^n$, the values are rational numbers:
\begin{equation}\label{D:14}
 r_k=\frac{T_{4k}}{T_0\,(\varpi/\pi)^{4k}}:\qquad
 r_1=\tfrac14,\quad r_2=\tfrac{17}{16},\quad r_3=-\tfrac{111}{64},\quad r_4=\tfrac{12513}{256},
 \quad r_5=-\tfrac{146079}{1024},\ \ldots
\end{equation}
The $r_k$ are rational with denominator dividing $4^k$ (proved in \S\ref{el:sec}, Corollary~Le8); the values shown are exact.
The reason is the identity $\Theta(z)/T_0=\sigma(u)^2\wp(u)$
with $u=\varpi z/2\pi$, where $\sigma$ and $\wp$ are the Weierstrass functions of the lattice with $g_2=4$,
$g_3=0$, whose Taylor coefficients are rational; this identity is proved in \S\ref{el:sec} (Theorem~Le6).
It gives $r_k$ with denominator dividing $4^k$. That lattice sums of
this kind are algebraic multiples of powers of a period is a classical theorem; compare Damerell
\cite{Dam70}. So derivatives of $\Theta$ are \emph{arithmetic} objects, once the period is removed. Each derivative costs one factor $\varpi/\pi$; this is why a term linear in the order
of differentiation, $-a\log\Omega$, appears in the size estimate of \S I.7.
\item \emph{Gaussian size.} By the Cauchy--Schwarz inequality,
$|T_n|^2\le\big(\sum_\lambda|\lambda|^{2n}e^{-\pi|\lambda|^2}\big)\,T_0$, and the first factor is comparable
to $\int_{\C}|z|^{2n}e^{-\pi|z|^2}dz=n!/\pi^n$. So $|T_n|$ is at most about $\sqrt{n!/\pi^n}$: the square root
of a factorial, not a factorial. Numerically, $\log|T_{120}|=119.0$, while
$\tfrac12\log(120!/\pi^{120})=160.2$ and $\log(120!)=457.8$.
\end{enumerate}

\textbf{The general case.} In the proof, $\Lambda$ is replaced by a lattice of the following kind. Let $E$ be
a number field with complex multiplication (it is specified in \S I.4.4 and in Part II, \S6.1), and let $L$,
$L'$ be lattices over its ring of integers, of ranks $r$ and $r'$, each with a positive definite Hermitian
form. The role of
$\Lambda$ is played by the lattice of linear maps from $L$ to $L'$. The variable $z$ becomes a complex
$r\times r'$ matrix, one for each embedding of $E$ in $\C$ from a fixed set of $2n$ embeddings (a \emph{CM
type}; $2n$ is half the degree of $E$, and $n$ is the number of columns of the table of \S I.5). The theta series $\Theta_{L,L'}(z)$ is a holomorphic function of
all these matrix entries, and one applies to it a differential operator $P(\partial)$, where $P$ is a
polynomial in the entries:
\begin{equation}\label{entry}
 P(\partial)\,\Theta_{L,L'}(0)=\sum_\nu P_\nu\,\nu!\,c_\nu,\qquad
 \Theta_{L,L'}(z)=\sum_\nu c_\nu z^\nu,\quad P(z)=\sum_\nu P_\nu z^\nu .
\end{equation}
Here $\nu$ runs over exponent vectors and $\nu!$ is the product of the factorials of its entries. As in
item 2 above, these numbers become algebraic only after the period has been removed; in the proof this is
done by writing the theta series in algebraic coordinates from the start.
\emph{The numbers \eqref{entry}, so normalised, are the divisible objects of the proof} \quoted.

\textbf{Why differentiate at all?} There are three separate reasons.
\begin{itemize}
\item \emph{It is forced by the construction.} The numbers \eqref{entry}, for suitable $P$, are exactly the
entries of the matrix of a theta lift (\S I.7) between spaces of forms of prescribed \emph{weights}. The
degree of $P$ at a given embedding is one of the parameters (weights) of these spaces. We call it the
\emph{order} $a$, and reserve the word weight for the parameters $w$ of \S I.4.3. It can be chosen
freely, and it will be chosen huge.
\item \emph{Differentiation creates divisibility.} This is the mechanism of \S I.3: in a suitable algebraic
coordinate the function has $p$-integral Taylor coefficients, an $a$th derivative multiplies a coefficient
by $a!$, and at an inert prime the net gain is $a\,(1/p+1/p^3+\cdots)$.
\item \emph{Differentiation costs little in size.} By the Gaussian computation in item 3 above, an $a$th
derivative has squared size about $a!$.
\end{itemize}

\textbf{A warning, with numbers.} The simplest theta function above does \emph{not} have the full
divisibility. At $n=120$ the exponent of $3$ in $r_{30}$ (known exactly, by Corollary~Le8 of \S\ref{el:sec}) is $14$, whereas $n\cdot\frac{3}{8}=45$. (For the
lemniscate sine of \S I.3 it was the full amount; for that function the size was too large.) Having both
properties at once, full divisibility at inert primes \emph{and} Gaussian size, is a theorem about the
particular theta series used in the proof, written in particular coordinates. It is the deepest local
statement in the argument, and it is quoted, not reproved, here \quoted. (See \S III.5 for what is proved:
the divisibility at inert primes is written up in \S III.5, conditionally on the integrality of the theta kernels in the coordinates of the height-two logarithm (hypothesis (B)), which is taken from the Lean development.)

\subsection*{I.4.2\quad Repair of (D3): squared Gaussian norms}

Recall the defect: for the lemniscate numbers, divisibility at inert primes accounted for only half of
$\log|A_n|$. The repair is point (i): the size becomes about $\sqrt{a!}$ in place of $a!$. Point (ii) fixes
the units in which the two sides are then compared.

\textbf{(i) Size is measured by a squared norm, and it is Gaussian.} Give the polynomials in a complex
variable $z$ the inner product with weight $\pi^{-1}e^{-|z|^2}$. Then the monomials are orthogonal and
\begin{equation}\label{D:16}
 \|z^a\|^2=\int_{\C}|z|^{2a}\,\pi^{-1}e^{-|z|^2}\,dz=\int_0^\infty t^ae^{-t}\,dt=a! .
\end{equation}
For several variables the norms multiply: $\|z^\nu\|^2=\nu!$. Put
$\|P\|^2=\sum_\nu|P_\nu|^2\nu!$. Applying the Cauchy--Schwarz inequality to \eqref{entry}, written as
$\sum_\nu(P_\nu\sqrt{\nu!})(c_\nu\sqrt{\nu!})$, gives
\begin{equation}\label{CS}
 \big|P(\partial)\Theta(0)\big|^2\ \le\ \|P\|^2\cdot\sum_{\nu}\nu!\,|c_\nu|^2 ,
\end{equation}
the second sum being over the monomials that occur in $P$. For $P=z^a$ the first
factor is $a!$. The second factor is a Gaussian-weighted norm of the theta function; bounding it is the
analytic part of the size estimate. So the \emph{squared} size of an entry is about $a!$ times lower-order
factors: its logarithm is $a\log a+O(a)$, the term of order $a$ being computed exactly in \S I.7.

\textbf{(ii) Both sides are measured for the squared norm.} Here and below $d$ is an auxiliary prime, $d\equiv3\bmod4$, and the inert primes are those inert in $\Q(\sqrt{-d})$; see \S I.4.4. After the period normalisation the
entries \eqref{entry} are algebraic numbers, in a field containing $\Q(\sqrt{-d})$. If such a number $\beta$
is a nonzero algebraic integer divisible by $p^e$, for any prime $p$, then the \emph{average} of
$\log|\sigma(\beta)|^2$ over the complex embeddings $\sigma$ of the field is at least $2e\log p$. This is the
product formula. (At a single embedding it can fail, because of units; only the average is controlled.) In
an imaginary quadratic field there is just one pair of embeddings and the statement is elementary: if
$\beta=p^e\beta_0$ with $\beta_0$ a nonzero integer of that field, then $|\beta|^2=p^{2e}|\beta_0|^2\ge p^{2e}$. So
divisibility by $p^e$ contributes $2e\log p$ to the logarithm of the squared size.

\textbf{Result.} With $e\approx a/p$ at each inert prime $p\le a$, the divisibility side is
\begin{equation}\label{D:18}
 \sum_{p\ \mathrm{inert}}2\cdot\frac ap\log p\ =\ a\log a+O(a),
\end{equation}
since the inert primes are half of all primes and $\sum_{p\le a}\frac{\log p}p=\log a+O(1)$ (Mertens'
theorem for the residue classes modulo $d$; the $O(a)$ depends on $d$). The size side is
$a\log a+O(a)$, the term of order $a$ being computed exactly in \S I.7. \emph{The leading terms now agree}, and the comparison has become a
statement about the terms of order $a$. Compare (D3): there the size was $\log|A_n|^2\approx2n\log n$
against the same $n\log n$ from divisibility.

\subsection*{I.4.3\quad Repair of (D1): a huge order, and a weight}

\textbf{What the exponent really is.} For the entries \eqref{entry} the exponent of an inert prime is not
$a\,(1/p+1/p^3+\cdots)$ but a modification of it. Fix one of the $2n$ embeddings; call it the
\emph{place}. Let $a$ be the order (the degree of $P$) at that place. Then, schematically,
\begin{equation}\label{expo}
 \text{exponent of an inert prime } p\ \ \ge\ \sum_{h\ \mathrm{odd}}\Big(\frac a{p^h}-w_h\Big)_+
 \ -\ (\text{small losses}),
\end{equation}
where $w_h>0$ is a \emph{weight} attached to another place: the one reached from the given place by the
Frobenius automorphism of $p^h$ followed by complex conjugation. The display keeps only the main term. The
statement proved in the source is an inequality for sums of such exponents over the $m$ lifts and over the
places, with several loss terms: from the digits of the integers involved, from repeated visits of a
Frobenius orbit to the same place, and from indices of lattices (Part II, \S6.4) \quoted\ (see \S III.5 for what is proved).

\textbf{The Frobenius and the places.} The field $E$ lies inside a field generated by roots of unity. The
Frobenius of $p$ is the automorphism of $E$ induced by raising those roots of unity to the $p$th power. It
permutes the $4n$ embeddings of $E$, and it depends only on residue classes of $p$. The $2n$ places are
chosen embeddings, one from each pair of complex conjugates. At a prime that is inert in $\Q(\sqrt{-d})$ an
odd power of the Frobenius carries each chosen embedding to one that is not chosen; composing with complex
conjugation brings it back. This is the reason for the conjugation above. Since the Frobenius
depends only on a residue class of $p^h$, residue classes of prime powers will appear in \S I.5.

\textbf{Where the subtraction of $w$ comes from.} The polynomial $P$ is a product, over the $2n$ places, of
minors (determinants of sub-matrices) of the matrix variable at each place. Its shape at each place is
described by a diagram of boxes. At the chosen place the diagram has a first row of length about $a$. At
another place $\tau'$ the polynomial contains about $w_{\tau'}$ minors, where $w_{\tau'}$ is the weight
parameter at $\tau'$. At the scale $p^h$ the Frobenius links the chosen place to the place $\tau_h$, and each
minor there can absorb only one derivative in a linked direction: in $\det(X+sb,\ Y+tb)$ there is no term in
$st$, $s^2$ or $t^2$. Each derivative absorbed at $\tau_h$ costs one unit of the gain at the scale $p^h$. There are
about $w_h$ minors at $\tau_h$, each absorbing at most one, and at most $a/p^h$ units are available; so the
loss at that scale is at most $\min(a/p^h,w_h)$, and what survives is $(a/p^h-w_h)_+$ (Part II, \eqref{II:t6.1}).
In the simplest case (a single index; Part II, \eqref{II:t6.1a}) the loss comes from the digits of the integers involved, by Legendre's
formula, and is of size $O((\log a)^2)$ for each prime (Part II, \S6.4). The losses in the statement the
source proves are larger and of several kinds, as said above.

\textbf{Example.} Take $a=81$, $p=3$, and $w_h=2$ for all $h$.
\begin{center}\small
\begin{tabular}{c|rrrr}
\toprule
$h$ & $1$ & $3$ & $5$ & $7$\\
\midrule
$p^h$ & $3$ & $27$ & $243$ & $2187$\\
$a/p^h$ & $27$ & $3$ & $0.33$ & $0.04$\\
$a/p^h-w$ & $25$ & $1$ & $-1.67$ & $-1.96$\\
$(a/p^h-w)_+$ & $25$ & $1$ & $0$ & $0$\\
\bottomrule
\end{tabular}
\end{center}
\noindent The exponent is at least $25+1=26$, up to the losses. Every term with $p^h>a/w=40.5$ is zero.

\textbf{The choice of $a$ and $w$, and what it achieves.} Let $x$ be the scale at which we want information
about primes. Put
\begin{equation}\label{D:20}
 b=\exp(x^\varepsilon),\qquad w\approx b,\qquad a\approx xb .
\end{equation}
Then $a/w\approx x$, and the elementary identity
\begin{equation}\label{ramp}
 \Big(\frac a{p^h}-w\Big)_+=w\,\Big(\frac{x}{p^h}-1\Big)_+\qquad(x=a/w)
\end{equation}
shows what has been gained. (In the example, $x=40.5$ and for $p^h=3$ the right side is
$2\cdot(40.5/3-1)=25$.)
\begin{itemize}
\item \emph{Only prime powers up to $x$ are involved,} at the places that carry the weight $b$. The right
side of \eqref{ramp} vanishes for $p^h\ge x$. Although we differentiate about $xb$ times, the primes that
matter at those places are those below $x$. (At the places with small weights the cut-off is much larger;
see below.)
\item \emph{The kernel is the ramp of \S I.1.} The factor $(x/p^h-1)_+$ is exactly the weight with which
prime powers enter $G(x)$. This is where the smoothing in the definition of $G$ comes from.
\item \emph{Rounding becomes invisible.} Each prime power below $x$ now contributes about $b$ times its
ramp value. The bound \eqref{expo} is used only where its main term is positive; elsewhere one uses only
that the exponent is not negative, which is the integrality of the entries. The loss terms are of lower
order, and after dividing the whole inequality by $b$ they are absorbed in the error of (M). The
precise accounting is part of the quoted statement. This removes defect (D1), and with it the
floor-function problem of \S I.2.
\end{itemize}

\textbf{Why $b$ must be exponentially large.} To defeat rounding, a power of $x$ would do.
One other step needs more. There are $2n$ places and each needs a weight, but only some of them can
be given the weight $b$ (the columns that carry it form the set $T_2$ constructed in \S III.3). The others are given very small weights $w'$, so that they test primes at a scale
$a/w'\ge\frac13\exp(x^{\varepsilon/3})$; only at such a scale is the distribution of primes in the relevant
progressions known unconditionally with a good error term (\S I.6(b)). The weights in the construction are positive integers, so
$w'\ge1$; this forces $a\ge\frac13\exp(x^{\varepsilon/3})$, and hence $b$ exponential in a power of $x$. The size of $b$ also enters the control of denominators (\S I.7).
What is bounded there is a weighted sum of logarithms of orders, $\sum_\ell\log\ell\sum_ah_{\ell,a}$ in the
notation of Part II, \S6.8, and $b$ is chosen so that the available bound is at most $\sqrt b$. After division
by $b$ this is negligible.

\subsection*{I.4.4\quad Repair of (D2): characters, and a table}

\textbf{Only inert primes: use a character.} Let $d$ be a prime with $d\equiv3\bmod4$ and let
$\chi_d(r)=(r/d)$ be the Legendre symbol. The primes inert in $\Q(\sqrt{-d})$ are those with
$\chi_d(p)=-1$. A sum over odd powers of inert primes is a sum over all prime powers with the coefficient
$1-\chi_d(r)$:
\begin{equation}\label{D:22}
 1-\chi_d(p^h)=\begin{cases}2,& p \text{ inert},\ h\text{ odd},\\ 0,&\text{all other prime powers with } p\ne d.\end{cases}
\end{equation}
(With this coefficient $2$ the sum matches the factor $2$ in $2e\log p$ of \S I.4.2(ii).) So the divisibility side is
$\sum_r\Lambda(r)\,(1-\chi_d(r))\,(\cdots)$: the sum for $\zeta$ \emph{minus} the sum for $L(s,\chi_d)$. To
isolate $\zeta$ one must know that the $\chi_d$-sum is small. This is arranged by choosing $d$ well
(\S I.6(e)), and it is not circular.

\textbf{Many places: a table.} By \eqref{expo} the weight met by a prime power $r=p^h$ depends on the place
to which the Frobenius of $r$ leads. The field used in the proof is
\begin{equation}\label{D:23}
 E=\Q\big(\zeta_N+\zeta_N^{-1},\ \sqrt{-q},\ \sqrt{-d}\big),\qquad \zeta_N=e^{2\pi i/N},\quad N=4n,
\end{equation}
with $n$ a large power of $2$ and $q\equiv3\bmod4$ a second auxiliary prime. Starting from the fixed place, the place reached by a prime power $r$ is determined by
two things: the residue class of $r$ modulo $N$ up to sign (the Frobenius sends $\zeta_N+\zeta_N^{-1}$ to
$\zeta_N^r+\zeta_N^{-r}$), and the quadratic symbol $\chi_q(r)$ (the Frobenius multiplies $\sqrt{-q}$ by $\chi_q(r)$; the complex conjugation, which enters for the odd powers of inert primes that contribute here, then changes its sign, which only interchanges the names of the two rows). For general $r$ prime to $2q$ this is a labelling of residue classes; the prime $d$ itself is treated separately, with coefficient $\log d$ (Part II, \S4.1). So the places can be arranged in a \emph{table}, with two rows (the value of $\chi_q$) and
$n$ columns (the residue class), and each prime power is assigned to a cell of the table. The inequality
that comes out of the construction, stated in \S I.5, is a bound for a sum over prime powers in which each
prime power carries the weight of its cell. \S I.6 then removes the dependence on the cell by averaging.

\subsection*{I.4.5\quad What the reader can check, and what is quoted}

\begin{center}\small
\begin{tabular}{p{.62\textwidth}p{.3\textwidth}}
\toprule
Statement in this section & Status\\
\midrule
$T_n$, the period $\varpi$, the rational numbers $r_k$, the size of $T_n$ & computed; script \texttt{e6\_hecke\_theta.py}\\
$\|z^a\|^2=a!$ and inequality \eqref{CS} & proved above\\
Identity \eqref{ramp} and the example $a=81$ & elementary\\
Leading terms agree: $\sum_{p\text{ inert}}\frac{2a}p\log p=a\log a+O(a)$ & Mertens' theorem in progressions\\
A determinant absorbs one derivative per direction; the factorial count & proved in Part II, \S6.4\\
The entries of the lift are the numbers \eqref{entry} & quoted\\
The exponent bound \eqref{expo} for the theta series of the proof & written up in \S III.5, conditionally on the integrality of the theta kernels in the coordinates of the height-two logarithm (hypothesis (B)), which is taken from the Lean development\\
\bottomrule
\end{tabular}
\end{center}

\section{The inequality that comes out}

Here is the statement. It has the form: there exist data (weights, two sets of columns, integers $m$, $k$)
with the properties listed, such that the inequality (M) holds. The properties are given here in outline and
exactly in Part II, \S4.2--4.3. The statement looks elaborate; the example after it shows that it is
elementary to state.

\textbf{Data.} Fix a small $\varepsilon>0$ and a scale $x$. Let $q\equiv3\bmod8$ and $d\equiv7\bmod16$ be
primes with $q\mid d+1$ (how they are chosen: \S I.6). They are very large: $\log q$ is about $Y^{\varepsilon/5}$ and $\log d$ about
$Y^{\varepsilon/5}(\log Y)^2$, where $Y$ is the scale at which the pair is chosen; $x^{1/J}\le Y\le x$ with
$J\ge1$ depending on $\varepsilon$ only (\S I.6(e); Part III, \S\ref{as:sec}, \S7).
Let $n$ be a power of $2$ with $x^{2+8\varepsilon}\le n\le2x^{2+8\varepsilon}$, and $N=4n$. Let $H$ be the group of odd residue classes modulo $N$ with $r$ and $-r$ identified; it
has $n$ elements, written $[r]$. Let $b=\exp(x^\varepsilon)$.

\textbf{The table.} Picture a table with two rows $v\in\{0,1\}$ and $n$ columns indexed by $H$. An odd
prime power $r=p^h$, $p\ne q$, is assigned to the cell
\begin{equation}\label{D:24}
 \text{row } v(r)=\tfrac12(1+\chi_q(r)),\qquad\text{column } [r] .
\end{equation}
The $2n$ cells are the $2n$ places of \S I.4.3. Each cell $i$ carries a positive weight $w_i$. A cell of row $1$ is \emph{distinguished} when its column lies in a set of columns of density $\delta$; its weight
is then (very nearly) $b$. All other weights are tiny by comparison: so small that $a/w_i\ge\frac13\exp(x^{\varepsilon/3})$.

\textbf{The prime sum and its main term.} For a derivative order $a$ with $|a/b-x|\le1$ (it depends on the label $t$ introduced below, and is written
$a_t$ in Part II), put
\begin{align}\label{D:25}
 P&=\sum_{p\ne2,q}\ \sum_{h\ge1}\big(1-\chi_d(p^h)\big)\log p\,
     \Big(\frac a{p^h}-w_{\mathrm{cell}(p^h)}\Big)_+ ,\\
 M&=a\Big(\log a-1-\gamma+\frac{L'}{L}(1,\chi_d)-\langle\log w\rangle\Big),
\end{align}
where $\langle\log w\rangle$ is the average of $\log w_i$ over the $2n$ cells. $P$ is the divisibility side:
by \eqref{ramp} it is a sum of ramps, one for each cell, at the scale $a/w_i$. $M$ is its expected value
if primes were perfectly distributed among the cells (\S I.6(a); proved in Part II, \S5.1). So $P-M$ measures the
irregularity of primes in the cells of the table, at the scales $a/w_i$.

The primes $2$ and $q$ can be omitted without changing anything: $d\equiv7\bmod8$ gives $\chi_d(2)=1$, and
$q\mid d+1$ gives $\chi_d(q)=1$ by quadratic reciprocity, so the coefficient $1-\chi_d$ vanishes on their
powers. This is what the congruence conditions on $q$ and $d$ are for (Part II, \S4.1). (The powers of $d$
do occur, with coefficient $\log d$, but only at the large scale, since $d>x$.)

\textbf{The inequality.} The construction depends on three further choices: a translate $u\in H$ (which
moves the distinguished cells to other columns), a label $t$ in a set $T\subset H$ (which shifts the column
of every prime power), and an index $j=1,\dots,m$ numbering the $m$ repetitions. Write $P_{u,t,j}$,
$M_{u,t,j}$ accordingly.
\begin{equation}\label{D:26}\tag{M}
 \boxed{\ \frac1n\sum_{u\in H}\ \sum_{t\in T}\ \sum_{j=1}^m\big(P_{u,t,j}-M_{u,t,j}\big)\ \le\ \mathcal E,
 \qquad \frac{\mathcal E}{B}\le x^{C\varepsilon}\Big(\frac xm+\frac{xm}{k}+m+\frac k{n^2}\Big)\ }
\end{equation}
Here $B=|T|\,m\,b\,\delta/2$ is the natural normalisation ($\delta$ is the proportion of distinguished
columns), $k$ is a further parameter of the construction (a rank), and $C$ is an absolute constant. The
inequality is asserted for $x\ge x_0(\varepsilon)$ \quoted.
In words: \emph{on average over translates, the primes are not more numerous than expected in the
table, up to the four errors shown.} It is an upper bound, as every instance of \eqref{det} is.
The set $T$ and the set of distinguished columns (called $T_2$ in Parts II and III) are constructed in
Part III, \S III.3; how the construction uses them is part of the quoted statement.

More precisely: there are intervals $[Z,Z^{3/2}]$ with $Z$ arbitrarily large such that for every $x$ in
them the construction provides primes $q,d$ for which Estimates B and C hold at $x$ (\S I.6(e)), chosen at
a scale $Y$ with $x^{1/J}\le Y\le x$, where $J\ge1$ depends on $\varepsilon$ only, and
$m=x^{1/2+O(\varepsilon)}$, $k=x^{9/2+O(\varepsilon)}$,
sets $T$, $T_2$ with $|T_2|/n=\delta\ge x^{-\varepsilon}$ and $|T|/n\ge\frac12x^{-\varepsilon/8}$, orders $a_t$
and weights as above, for which (M) holds. The constants $C$ and $x_0(\varepsilon)$ do not depend on $x$,
$Y$, $q$, $d$. This is what Part III, \S\ref{as:sec} states exactly, as Hypothesis (H12); it is assumed
there.

\textbf{A table small enough to write out.} Take $q=11$, $d=1319$, $N=16$, so $n=4$ and
$H=\{[1],[3],[5],[7]\}$. These primes satisfy the congruences: $q\equiv3\bmod8$, $d\equiv7\bmod16$, and
$q\mid d+1=1320$. The primes $3$, $5$, $7$ and $11$ are squares modulo $1319$, so they split in
$\Q(\sqrt{-1319})$ and have coefficient $0$. The prime $23$ is a square modulo $11$ and $23\equiv7\bmod16$: it
goes in row $1$, column $[7]$, with coefficient $2\log23$ because $23$ is a non-square modulo $1319$. The
primes up to $61$ with non-zero coefficient fall as follows.
\begin{equation}\label{D:27}
 \begin{array}{c|cccc}
 v\backslash\text{column} &[1]&[3]&[5]&[7]\\ \hline
 0&17&13,\,29,\,61& - &41\\
 1&47& - &37,\,59&23
 \end{array}
\end{equation}
With $x=1000$, $b=10$, $a=10^4$, weight $10$ in cell $(1,[1])$ and $0.05$ elsewhere, summing over the prime powers
gives $P=84701.032$ and $M=84463.394$: the excess $P-M$ is $237.64$, against a main term of nearly $10^5$.
Moving the distinguished cell to the other three columns gives excesses $-300.26$, $-176.73$, $229.13$.
The mean of the four excesses is $-2.55$. (The example only
illustrates the identities, which hold for any positive weights. Its parameters do not satisfy all the
conditions under which (M) is proved. There $x$ is large, $q$ and $d$ exceed a bound that depends on the
ranks of the construction and is at least $541$ (here $q=11$ is too small), the odd primes up to a bound
that is at least $11$ split in $\Q(\sqrt{-d})$ (here they do up to $11$, but $13$ is inert), the construction
even asks that they divide $d+1$ (here $7\nmid1320$), and the weights are integers. See \S III.10. The first version of this paper used $q=3$, $d=23$, which violates
more of these conditions.)

\section{From the table to $G(x)$}

This step is classical analytic number theory, and it is proved in full, from (M), B and C, in Part II, \S4--5.
Here are the ideas.

\textbf{(a) Each cell is an arithmetic progression.} The prime powers in a given cell are those in one pair
of residue classes $\pm c$ modulo $N$ with a prescribed value of $\chi_q$; the coefficient $1-\chi_d$ then
selects the odd powers of primes inert in $\Q(\sqrt{-d})$. For fixed moduli and scale tending to infinity,
the prime number theorem in progressions gives each cell the fraction $1/(2n)$ of the primes, half of them
inert, and $M$ is defined as the resulting main term, with its constants. The identity
$P-M=\sum_iw_iR_i(a/w_i)$, where $R_i(y)$ is the remainder of cell $i$ at scale $y$ (Part II, \eqref{II:t5.2}), is then
exact (Part II, \eqref{II:t5.3}). No uniformity is claimed at this point. At the scale $x$ there are far more cells
than primes, and a single $R_i(x)$ need not be small; the argument uses only the average over translates
in (c). At the large scale the uniformity is Estimate C.

\textbf{(b) Two scales.} A cell with weight $w_i$ tests primes up to $a/w_i$. For a distinguished cell this
is $a/b\approx x$: the scale we care about. For every other cell it is at least $\frac13\exp(x^{\varepsilon/3})$,
an enormous scale. \emph{Why arrange this?} The construction needs weights on all $2n$ cells, but we want
information at scale $x$ only. At the enormous scale the distribution of primes in these progressions can be controlled with a saving
better than any power of $x$. This is Estimate C below. It rests on the classical zero-free region for
Dirichlet $L$-functions; it is available because $\log q$ and $\log d$ are much smaller than
$x^{\varepsilon/3}$, and because $q$ and $d$ are chosen so that no exceptional real zero interferes. So those
cells contribute a negligible amount. In the example they account for $0.15$ of the mean excess $-2.55$.

\textbf{(c) Averaging over translates.} As $u$ runs over $H$, the distinguished cells visit every column
equally often. So the average over $u$ no longer sees the residue class modulo $N$. It sees only the two
quadratic conditions: row $1$ means $\chi_q(r)=1$, and the coefficient is $1-\chi_d(r)$. The average of the
distinguished part is therefore
\begin{equation}\label{D:28}
 \frac{\delta b}{2}\sum_{r\le x}\Lambda(r)\big(1+\chi_q(r)\big)\big(1-\chi_d(r)\big)\Big(\frac xr-1\Big)
 \quad\text{minus its main term.}
\end{equation}
In the example this accounts for the other $-2.70$ of the $-2.55$: it is $\frac{10}8R_{11,1319}(1000)$, with $R_{11,1319}(1000)=-2.163$ (Part~II, \S5.3).

\textbf{(d) Expanding the product.} Since $(1+\chi_q)(1-\chi_d)=1+\chi_q-\chi_d-\chi_q\chi_d$, the last
expression is $\frac{\delta b}2R_{q,d}(x)$ with
\begin{equation}\label{D:29}
 R_{q,d}(x)=G(x)+G_{\chi_q}(x)-G_{\chi_d}(x)-G_{\chi_q\chi_d}(x),
\end{equation}
where, for a primitive character $\chi\ne1$, $G_\chi(x)=\sum_{r\le x}\Lambda(r)\chi(r)(x/r-1)+x\,\frac{L'}L(1,\chi)$ is the analogue of $G$ for
$L(s,\chi)$. \emph{The function $G$ appears with coefficient $+1$.} (The construction gives this at the scale $a/b$,
with $|a/b-x|\le1$; the passage to $x$ costs $O(1+\log x)$, see Part II, \S5.4.) Dividing (M) by $B$, and using Estimates B
and C below to bound the error terms (Part II, \S5.4), gives:

\begin{thmA} Let $x$ be a point of one of the intervals of \S I.5, with its primes $q,d$. For the parameters
of the construction at $x$, with $m=x^{1/2+O(\varepsilon)}$, $k=x^{9/2+O(\varepsilon)}$, $n=x^{2+O(\varepsilon)}$ (exactly: $m\asymp x^{1/2-\varepsilon/16}$, $k\asymp x^{9/2+10\varepsilon}$, $n\asymp x^{2+8\varepsilon}$),
\begin{equation}\label{D:30}
 R_{q,d}(x)\ \le\ x^{C\varepsilon}\Big(\frac xm+\frac{xm}k+m+\frac k{n^2}\Big)\ \le\ x^{1/2+C'\varepsilon}
 \qquad(x\ge x_0(\varepsilon)).
\end{equation}
\end{thmA}

\textbf{(e) The three twisted sums.} To pass from $R_{q,d}$ to $G$ we need:

\begin{thmB} Let $0<\varepsilon<\frac1{52}$. For all large $Y$ there are primes $q\equiv3\bmod8$, $d\equiv7\bmod16$ with $q\mid d+1$, $\log q\asymp Y^{\varepsilon/5}$, $\log d\asymp Y^{\varepsilon/5}(\log Y)^2$, such that $|G_\chi(x)|\le x^{1/2+\varepsilon}$ for
$\chi=\chi_q,\ \chi_d,\ \chi_q\chi_d$ and all $x$ in the interval $[Y,Y^{3/2}]$. (Proved in Part III as Theorem~\ref{thm:pair}(d), in a
stronger form: exponent $\frac12+\frac\varepsilon2$, and any interval $[Y,Y^A]$. With $A=J$ this covers every
$x$ with $x^{1/J}\le Y\le x$, as in \S I.5.)
\end{thmB}

\begin{thmC} For the same primes $q$, $d$ and every $x\in[Y,Y^{3/2}]$, $Y$ large, every one of the $4n$ primitive Dirichlet characters $\chi$ that occur when
the cells of the table are expanded in characters (their conductors divide $4nqd$) satisfies
$|G_\chi(y)|\le y\exp(-\kappa x^{\varepsilon/8})$ for all $y\ge\frac13\exp(x^{\varepsilon/3})$, with a constant
$\kappa>0$; here $G_\chi=G$ for the trivial character. (Proved in Part III as Theorem~\ref{thm:C}, where $\kappa$ is called $c'$ and may be any positive number
once $x$ is large. The proof uses a further property of $q$ and $d$, part (e) of Theorem~\ref{thm:pair}.
The hypotheses of that theorem hold for every $x\ge Y$ once $Y$ is large, so the estimate holds as well for every $x$ with $x^{1/J}\le Y\le x$, as in \S I.5.)
\end{thmC}

\emph{Is this circular?} It looks like assuming the Riemann hypothesis for three $L$-functions. It is not,
for a reason of scale. The moduli $q$ and $d$ are astronomically larger than $x$: their logarithms are
powers of $x$. The quantity $G_\chi(x)$ has two parts. One is the constant $\frac{L'}L(1,\chi)$, an infinite
sum over prime powers. The other is the finite sum over $r\le x$. The tail of the constant beyond $Z=Y^{600}$ is shown to be small using a zero-density estimate for
zeros close to the line $\re s=1$; this discards a small proportion of the moduli, and uses nothing about
zeros near $\re s=\frac12$. What remains is a sum over $r\le Z$, which is still very short compared
with $q$ and $d$. Mean-square estimates over the modulus show that for \emph{most} pairs $(q,d)$ such a sum
is as small as a random sum, about $\sqrt x$. One only needs one good pair. For C see (b) above. Estimates B and C are proved in Part III,
from standard theorems on Dirichlet $L$-functions that are cited there.

\textbf{Conclusion.} From A and B (A itself uses (M), B and C): $G(x)\le R_{q,d}(x)+3x^{1/2+\varepsilon}\le x^{1/2+C''\varepsilon}$ on
the interval. By \S I.1 this proves the Riemann hypothesis.

\section[Where (M) comes from]{Where the main inequality (M) comes from}

This is the part quoted from the Lean development; Part III has begun its proof: \S III.4 (the analytic part
of item 2 below) and \S III.5 (item 1 below, conditionally on one hypothesis from the Lean development).
Sections III.6--III.9 prove further steps, each from named statements of the Lean development, which are
quoted: the passage from derivatives to lattices, in part (\S III.6); the deduction of the size inequality
from eleven inputs (\S III.7); the deduction of the dimension bound of item 3 from three statements
(\S III.8); and the step from the bound for Galois cohomology to the bound for denominators (\S III.9). The
table at the start of Part III says what each section proves, proves conditionally, or quotes. The
construction of the maps and weights is not proved: Part III, \S\ref{ob:sec} defines the objects and takes
their existence from the Lean development, as seven numbered hypotheses. The assembly is written as a reduction: Part III, \S\ref{as:sec} deduces (M), at one scale, from twelve numbered hypotheses, and the Riemann hypothesis from these and a thirteenth. None of the hypotheses is proved outright, and (M) is still not proved. We describe the objects and the three things that have
to be proved about them.

\textbf{Theta series.} A theta series counts representations by a quadratic form. For the form $a^2+b^2$,
\begin{equation}\label{D:31}
 \Theta(z)=\sum_{a,b\in\Z}s^{a^2+b^2}=1+4s+4s^2+4s^4+8s^5+\cdots,\qquad s=e^{2\pi iz}.
\end{equation}
Theta series have two properties at once: they are sums over a
lattice, so their coefficients are integers with arithmetic meaning; and they are Gaussian in nature
(the terms have absolute value $e^{-2\pi y(a^2+b^2)}$, where $y$ is the imaginary part of $z$), so their sizes are governed by Gaussian integrals. These are the two
properties needed in \S I.4: divisibility of derivatives (\S I.4.1) and Gaussian size (the repair of (D3), \S I.4.2).

\textbf{A theta lift} is a map defined by a theta series in two sets of variables, used as the kernel
of an integral: $(\mathcal Tf)(h)=\int\Theta(g,h)\overline{f(g)}\,dg$. (It is conjugate-linear in $f$.) We describe the
case in which both groups are \emph{definite unitary groups}, for which everything is finite. (In the
construction some of the groups are not compact; Part III, \S\ref{ob:sec} gives the classical form for the
definite ones only.) In the definite case a ``function on the group'' is a function on
a finite set of lattices $L_1,\dots,L_c$ (the classes in a genus), the integral is a finite sum weighted by
$1/|\mathrm{Aut}\,L_i|$, and $\mathcal T$ is a matrix. Its entries are derivatives, at the origin, of the
theta series attached to pairs of lattices. So \emph{a theta lift is a matrix of derivatives of theta
functions}, between two spaces that carry natural integer lattices (functions with integer values) and
natural norms.

\textbf{The role of \eqref{det}.} Let $\mathcal T:L\to M$ be an injective linear map between lattices in
Euclidean spaces. If $\mathcal T(L)\subset p^eM$ for some prime power, then, comparing volumes,
\begin{equation}\label{slope}
 e\log p\ \le\ \log\|\mathcal T\|+\mu_{\max}(M)-\mu(L),
\end{equation}
where $\|\mathcal T\|$ is the operator norm and $\mu=-\frac1{\mathrm{rank}}\log(\text{covolume})$ is the
\emph{slope} of a lattice ($\mu_{\max}$: the largest slope of a sublattice). For $L=M=\Z$ this is
\eqref{det}. In the construction the lattices are modules over the integers of a field with complex
multiplication, and one works with squared norms: \eqref{slope} is doubled, so that $2e\log p$ is compared
with $\log\|\mathcal T\|^2$, as in \S I.4. Three things must be established.

\textbf{1. Divisibility.} \emph{The matrix of the lift is divisible by $p^{e}$ with $e$ of the form
$\sum_{h\text{ odd}}(a/p^h-w)_+$ at inert primes.} The mechanism is that of \S I.3, applied to theta functions of lattices over a field with complex
multiplication: derivatives of order $a$ gain $a(1/p+1/p^3+\cdots)$. The subtraction of $w$ has an elementary source, explained in \S I.4.3:
each minor can absorb only one derivative in a given direction. Composing $m$ lifts multiplies the matrices, so
the exponents add: this is the analogue of a function vanishing to order $m$ in Stepanov's proof of the
Riemann hypothesis for curves over finite fields. The divisibility at inert primes is written up in \S III.5, conditionally on the integrality of the theta kernels in the coordinates of the height-two logarithm (hypothesis (B)), which is taken from the Lean development. The passage from this estimate to lattices is the subject of \S III.6. Proved there: the lemmas on slopes of lattices, the form of the prime sum and its splitting by classes of primes, and bounds for several of its terms. The theorem on the glued lattice is conditional on two statements of the Lean development. The bounds for the eleven loss sums and for their total are quoted, and so is the statement on valuations itself.

\textbf{2. Size.} \emph{The squared norm of the lift has logarithm at most
$a(\log a-1-\log\Omega_E)+\cdots$, with the linear term known exactly.} The leading term is the Gaussian
integral $\|z^a\|^2=a!$. The constants $\Omega$ and $\Omega_E$ are \emph{periods}: $\Omega$ belongs to an elliptic curve, and
$\Omega_E$, defined below, to the field $E$. A derivative with respect to the algebraic invariant differential
and one with respect to the analytic variable differ by a constant factor, the period; an $a$th derivative
acquires its $a$th power, and $\Omega$ is the square of that factor. Its value is
\begin{equation}\label{D:33}
 \log\Omega=\gamma+\log\frac{2\pi}{\sqrt d}-\frac{L'}{L}(1,\chi_d).
\end{equation}
In this form the identity is Kronecker's limit formula; $\Omega$ is then a geometric mean, over the ideal
classes of $\Q(\sqrt{-d})$, of squared periods of the corresponding elliptic curves. Expressing it through
values of the Gamma function is the formula of Lerch and of Chowla and Selberg \cite{CS67}, whose name the
source uses. For the field $\Q(i)$, with $d$ replaced by $4$ and $\chi_d$ by the character modulo $4$, exponentiating this formula gives $\Omega=2\pi(\varpi/\pi)^2=\Gamma(\tfrac14)^4/(4\pi^2)$ exactly:
that is $2\pi$ times the square of the factor $\varpi/\pi$ met in \S I.4.1.
The constant $\Omega$ above is the period for the elliptic curve. The constant that enters the size bound
for the field $E$ is $\Omega_E$, with
\begin{equation}\label{D:34}
 \log\Omega_E=\log\Omega-\log\frac{\mathrm{rd}_E}{\sqrt d}
 =\gamma+\log\frac{2\pi}{\mathrm{rd}_E}-\frac{L'}{L}(1,\chi_d),
\end{equation}
where $\mathrm{rd}_E$ is the root discriminant of $E$. Part II, \S6.6 gives both normalisations; the
cancellation of the linear terms is the same.
This is how $\gamma$ and $L'/L(1,\chi_d)$ enter the main term $M$. They are the same constants that
appear in the prime sum, by \S I.6(a); that the two sides produce the same linear term is the reason
the method can work at all. The remaining linear terms come from the slopes in \eqref{slope}: the number of
lattice classes counted with weights $1/|\mathrm{Aut}|$ (the Siegel mass formula) and the dimensions of the spaces (the Weyl dimension formula).
After an exact cancellation of factorials, the total is $M$.

\textbf{3. Denominators and rank.} Two things can spoil \eqref{slope}.
\begin{itemize}
\item \emph{Denominators.} The lift need not map integer-valued functions to integer-valued functions; one
must find the right lattices. They are assembled prime by prime, using the freedom to twist the lift by
characters. Comparing these choices costs an error at each lift. A quoted comparison bounds the total error
at each lift in terms of the lengths of certain Galois cohomology groups of the field (generalisations of
class groups); Proposition~G2 of \S III.9 combines it with the bound for those lengths. A bound, by $\sqrt b$, for the logarithm of the product of the orders of these groups
(over all primes $\ell$ and all twists up to $x^{400}$) is proved
separately in the Lean development (it is quoted here); it uses the theory of automorphic forms and Galois representations, and is a large part of the
formal development. Part III, \S III.9 proves the step from this bound and the quoted comparison to the bound for denominators, and
the vanishing of the groups in some cases over $\Q$ from cited theorems; the bound itself is quoted, with
an account of its proof in the Lean development.
\item \emph{Rank.} A single map of small rank controls only the \emph{maximal} slope of its target, which
can be far from the average and would lose the main term. One therefore uses many maps. Their images cannot
all be perpendicular to one another: the dimensions of the images add up to more than the dimension of the
target. What is
shown is a bound for their total overlap (a second moment; a trace formula gives it). It says that the images
are spread evenly: the weighted sum, over the maps, of the adjoint of the map followed by the map is close
in mean square to a multiple of the identity of the target. A lemma of linear
algebra then says the images fill almost the whole target, so that the \emph{mean} slope is controlled.
Part III, \S III.8 proves that the dimension bound of this lemma follows from three statements of the Lean
development (on traces, ranks and pairings) and stated size conditions; the three statements, and the
theorems behind them, are quoted.
\end{itemize}

\textbf{The elementary parts.} Several steps in this section are elementary and are proved, with worked
examples, in Part II, \S6: the persistence of integrality through a chain of maps;
the minor congruence; the factorial count (from Legendre's formula) giving the gain
$a\,p/(p^2-1)$; the Gaussian norm; the cancellation that leaves $M$; the elementary slope
inequality; and the linear-algebra lemma on images that are spread evenly in this sense.
The analytic part of the size estimate is proved in part in Part III, \S III.4: the Gaussian bound for one
block, for one map and for a chain, with four passages conditional on inputs from the Lean development.
Part III, \S III.7 proves that the size inequality \eqref{III:sizefinal} follows from eleven named inputs,
and that an index-corrected mean of periods equals $\log\Omega_E$ under hypotheses stated there. The eleven
inputs are quoted, among them the inner-product formula, the covolumes and the local index bounds. So the
size bound itself (Part II, \eqref{II:t6.5}, first line) remains quoted.

\section{Why the exponent is one half}

After division by $B$, the errors in (M) are, up to factors $x^{C\varepsilon}$,
\begin{equation}\label{D:35}
 \frac xm\qquad\frac{xm}{k}\qquad m\qquad\frac k{n^2}.
\end{equation}
\begin{itemize}
\item $x/m$: errors made \emph{once}, in the sizes of the lattices at the two ends of the chain of maps. The gain from divisibility
is collected $m$ times, so these are divided by $m$. This is why one composes many lifts.
\item $m$: errors accumulated over the $m$ steps. The budget of each step (ranks, weights, denominators) itself grows in proportion to $m$, so the total is of order $m^2$ before normalisation; the factor $m$ contained in $B$ leaves $m$.
\item $xm/k$ and $k/n^2$: costs of the auxiliary parameters $k$ (the rank) and $n$ (the size of the table).
With $k\approx x^{9/2}$ and $n\approx x^2$ the first is negligible ($x^{-3}$) and the second is again of
order $\sqrt x$.
\end{itemize}
The terms $x/m$ and $m$ are balanced when $m\approx\sqrt x$, and both equal $\sqrt x$. \emph{That is where the
exponent $\frac12$ comes from.} If the second cost were $m^2$ instead of $m$ after normalisation, one would have $x/m+m^2$, balanced at
$m=x^{1/3}$, and the same analytic argument would give only that $\zeta(s)\ne0$ for $\re s>\frac23$, provided the corresponding inequality could be proved with the other errors below $x^{2/3+\varepsilon}$. (Whether such an inequality can be proved without the estimate for Galois cohomology is a separate question, which we do not address.) So the precision of
the comparison of lattices is decisive.

\section{Dictionary with Alp\"oge's announcement}

The announcement\footnote{\url{https://alpo.ge/remark.pdf}, pages 3--4.} lists twenty steps, in the order in
which they were found between 6 October and 16 November 2026. It is a record of discovery, so it contains
earlier formulations and motivating ideas as well as steps of the final argument. This section says where
each step appears in this paper. The correspondence is our reading. For steps 1, 2, 5 and 7--20, and for the Poincar\'e bundle in step 3,
it rests on statements of the Lean development that were read in the course of this work. Steps 4 and 6,
and the rest of step 3, were not found in the Lean development and are matched from the description in the
announcement.

\textbf{Steps that are in this paper.}

\begin{center}\small
\begin{tabular}{rp{.43\textwidth}p{.43\textwidth}}
\toprule
 & Step in the announcement & In this paper\\
\midrule
1 & Stepanov determinant & \S I.2: a nonzero integer divisible by $\prod p^{e_p}$ is at least that large, \eqref{det}; the model is Stepanov's count of points on curves\\
5 & Switch to $\sum_{n\le X}\Lambda(n)(X/n-1)-X\log X+(1+\gamma)X$ & \S I.1: this is $G(X)$\\
8 & Iterated theta lifts between totally definite unitary groups; a chain of length about $\sqrt X$ & \S I.7 (theta lift; composition); \S I.8 ($m\approx\sqrt x$); in the construction some of the groups are indefinite (\S\ref{ob:sec})\\
9 & Norms of the lifts; cancellation using Chowla--Selberg, Weyl, Siegel mass & \S I.7, ``Size''; Part II, \S6.6--6.7; the norms, in part: Part III, \S III.4; the size inequality from eleven quoted inputs: \S III.7\\
10 & Lower bound for the divisibility of the $m$-fold lift & \S I.4.3, \eqref{expo}; \S I.7, ``Divisibility''; written up in Part III, \S III.5, conditionally on the integrality of the theta kernels in the coordinates of the height-two logarithm (hypothesis (B)), which is taken from the Lean development; the passage to lattices, in part: \S III.6\\
12 & Seesaw: the period of the iterated lift is a Bessel period independent of the Hecke characters & \S I.7, ``Denominators'' (the freedom to twist by characters); Part II, \S6.8\\
13 & Bost slope inequality & \S I.7, the inequality for a map of lattices; the elementary case is proved in Part II, \S6.8; a form sufficient for the deduction of (M) is proved in Part III, \S\ref{as:sec}, Lemma~M3\\
14 & Cost of each lift; ``one would stop at $2/3$''; need for ``division'' & \S I.8: $x/m+m^2$ gives $\frac23$, $x/m+m$ gives $\frac12$\\
15 & Glue an integral structure across Hecke characters (a ``gcd'') & \S I.7, ``Denominators'': lattices assembled prime by prime; Part II, \S6.8\\
16 & Relative trace formula; approximate orthogonality & \S I.7, ``Rank''; Part II, \S6.9--6.10\\
17 & Classes in $H^1$ not killed by twisting; a weakening of Vandiver's conjecture & \S I.7, ``Denominators'': the bound by $\sqrt b$ for Galois cohomology; Part II, \S6.8\\
20 & GRH & here only $\zeta$, the case of the trivial character\\
\bottomrule
\end{tabular}
\end{center}

\textbf{Steps that appear here in another form, or only in part.}

\begin{itemize}
\item \emph{Step 2} (the invariant derivation $D$ on a curve with complex multiplication satisfies
$D^p\equiv0$ at an inert prime; divisibility of $D^a\,\Omega\cdot\theta(0)$). This is the idea of \S I.3 and
\S I.4.1: derivatives on such a curve, the period ($\varpi$ in \S I.4.1, $\Omega$ in \S I.7), and exponents involving odd powers of $p$ only.
For the lemniscate sine the stated exponent is a theorem: at an inert prime $p$ the $n$th derivative at zero,
$A_n=\mathrm{sl}^{(n)}(0)$, is divisible by $p^{\lfloor np/(p^2-1)\rfloor}$ (\S\ref{el:sec}). For the Weierstrass function
$\wp$ of the same curve a divisibility with this exponent, for the derivatives at zero of $\wp(z)-z^{-2}$ times a
simple factor, is in the literature: for an odd inert prime $p$, $n\ge2$ and $b\in\Z[i]$ prime to $p$, the
derivative of order $n$ times $(1-b^{n+2})(1-p^n)$ is divisible by $p^{\lfloor np/(p^2-1)\rfloor}$
(Katz \cite{Kat81}; Bannai and Kobayashi \cite[Thm.~5.1]{BK16}; we have not consulted
these papers, and the statement was compared with \cite{BK16} by a reviewer). For the theta function $\Theta$ of $\Z[i]$ in \S I.4.1 the exponent does not hold: the exponent we
computed is smaller, and for a closely related two-variable function Bannai and Kobayashi \cite{BK10} found
coefficients that are not integral at $p=7$ (also not consulted). So the object of the announcement is not that
theta function. For the lemniscate curve the first statement of the step, that the $p$-th power of the invariant derivation vanishes modulo an inert prime $p$, is Proposition~Le20a\,(ii) of \S\ref{el:sec}. For what is proved in this paper about the divisibility at an inert prime, see Part III,
\S III.5.
\item \emph{Steps 4 and 6} (positivity $W(f)\ge0$ in Weil's explicit formula, obtained from a determinant; a
zero-free region from such a positivity). These are an earlier form of our first step. In the final
argument the statement ``an upper bound excludes zeros'' is the Landau-type lemma of \S I.1.
\item \emph{Step 7} (determinants containing Gan--Gross--Prasad periods; interlacing weights). In the Lean
development the period pairing at the end of the chain of lifts is identified with a Bessel period (\S I.7;
Part II, \S6.2 and \S6.8; source file [19]), and a condition of interlacing on the weights governs the
passage between consecutive unitary groups (source file [23]). We do not use the Gan--Gross--Prasad
terminology in the text.
\item \emph{Steps 11 and 19} (an Iwasawa tower; a construction of Behrend). The tower is the field $E$ with
$n$ a large power of $2$, which gives the columns of the table (\S I.4.4, \S I.5). Behrend's construction
is used to choose the set $T_2$ of distinguished columns (Part II, \S4.2). The set is constructed in
Part III, \S III.3, following source file [27], which uses the sphere construction of Behrend as formalised
in Mathlib. (The definitions in [17] are an earlier formulation; the final theorem does not use them.) The announcement itself expects Behrend's construction to be unnecessary.
\item \emph{Step 18} (covolumes; the index of theta series; congruences between lifts and forms). These are
the slopes of the lattices in \S I.7. The elementary inequality for slopes is proved in Part II, \S6.8, and the forms used for (M) in Part III, \S\ref{as:sec} (Lemma~M3; Proposition~M4, from quoted statements and a cited fact on the filtration by slopes); the covolumes and the indices are quoted or assumed.
\end{itemize}

\textbf{A step largely absent from the final proof.} \emph{Step 3} (an identity $qN=M^k+R$; a quadric;
automorphy through a theta function on a sixfold cover of $\mathrm{GL}_5$; the Poincar\'e bundle). We found
no trace of the identity or of a sixfold cover in the Lean development (quadrics occur in some file names, in what appear to be unrelated lemmas on orbit measures); step 16 of the
announcement refers to them as ``a different half-plane argument''. A Poincar\'e bundle does occur in the archive, on an auxiliary abelian scheme, in a construction of the line
bundle whose section gives the theta kernel in algebraic coordinates (\S I.4.1; Part II, \S6.5). The
declaration we tested from that file is not among those the final theorem depends on, so we cannot say
that the final proof uses it.

\textbf{Steps of the proof that the announcement does not list.} The choice of the auxiliary primes $q$ and
$d$, with the three character sums and the zero-density and mean-square estimates that control them
(\S I.6(e)); the two scales, and the size of $b$ (\S I.4.3, \S I.6(b)); and the form of Landau's lemma on
intervals (\S I.1; Part II, \S1). These are in the Lean development; file names are given in Part II. The
first is proved in Part III, \S III.1, the estimate at the large scale in \S III.2, and the lemma in Part II, \S1.

\section{The same proof for Dirichlet $L$-functions}

The Lean proof treats all Dirichlet $L$-functions at once, and $\zeta$ is the case of the trivial character. This section says what changes for another character. It is described from the statements of the Lean development and is not worked out in this paper.

Let $\psi$ be a primitive Dirichlet character of conductor $M$. Its powers $1,\psi,\psi^2,\dots$ form a finite cyclic group; let $o$ be the number of distinct powers, that is, the order of $\psi$. For $1\le j<o$ let $\psi_j$ be the primitive character that induces $\psi^j$ (the power $\psi^j$ itself need not be primitive), and let $G_{\psi_j}$ be the function of \S I.6 attached to it:
\begin{equation}\label{D:36}
 G_{\psi_j}(y)=\sum_{n\le y}\Lambda(n)\psi_j(n)\Big(\frac yn-1\Big)+y\,\frac{L'}{L}(1,\psi_j).
\end{equation}
The statistic $G$ is replaced by the combination
\begin{equation}\label{D:37}
 G^{\psi}(y)=G(y)+\frac1{3o+1}\sum_{j=1}^{o-1}\re\,G_{\psi_j}(y).
\end{equation}
For $n$ prime to $M$, the term $\Lambda(n)(y/n-1)$ occurs in it with the coefficient $1+\frac1{3o+1}\sum_{j=1}^{o-1}\re\,\psi^j(n)$, which equals $\frac{4o}{3o+1}$ if $\psi(n)=1$ and $\frac{3o}{3o+1}$ otherwise. For $n$ not prime to $M$ the coefficient is $1+\frac1{3o+1}\sum_{j=1}^{o-1}\re\,\psi_j(n)$, which can differ from this formula, since $\psi_j(n)$ may be nonzero where $\psi^j(n)=0$; it is still at least $\frac{3o}{3o+1}$. So every coefficient is positive. Multiplied by $3o+1$, the combination is the statistic attached to the product of $\zeta(s)^{3o+1}$ and the functions $L(s,\psi_j)$, $1\le j<o$; at a zero of this product the multiplicities add, so every zero of $\zeta$ or of an $L(s,\psi_j)$ gives a pole of the transform. The lemma of \S I.1 then applies with one addition. The proof for $\zeta$ uses that there is no pole on the real segment $0<s<1$; for $L$-functions a real zero $\beta$ cannot be excluded in advance. At such a zero the residue of the transform is a positive multiple of $1/(\beta(1-\beta))$. If it is the only pole nearest to the real point $\sigma$, the derivatives at $\sigma$ have size about $(\sigma-\beta)^{-j-1}$, and the fixed sign of the residue makes them positive; this contradicts the bound of Part II, \S1 when $\beta>\theta$; the Lean proof of the lemma has this case as a separate branch. The conclusion is: if for every $\varepsilon>0$ the bound $G^\psi(y)\le y^{1/2+\varepsilon}$ holds on intervals $[Y,Y^{3/2}]$ with $Y$ arbitrarily large, then $\zeta$ and every $L(s,\psi_j)$ have no zero with real part greater than $\frac12$. (For $o=1$ this is $G$ and $\zeta$.)

The rest changes as follows. The field $E$ is enlarged so that $\psi$ becomes one of its characters. The columns of the table are then indexed by a larger group: it consists of $o$ cosets of the subgroup on which the label coming from $\psi$ is trivial, and the common size of these cosets depends on how much $\psi$ shares with the residue classes modulo $4n$ (for the character $(\frac2\cdot)$ modulo $8$, for instance, the number of columns does not grow). The set $T$ is chosen with slightly different densities on different groups of columns; this is where the weight $\frac1{3o+1}$ comes from. Estimate B is needed for the twists $\xi\chi_q^i\chi_d^j$ of all characters $\xi$ modulo $M$, and Estimate C for the characters of the larger field; the theorems of \S III.1 and \S III.2 are stated in that generality. Everything else is the same: the main inequality (M), with its theta lifts, divisibility, size, trace formula and Galois cohomology, is the same construction for the larger field, and the passage from the table to one prime sum is the same computation with more characters.

\section{Notes on the literature}

This paper explains an argument due to Alp\"oge and Claude \cite{Alp26}. Apart from that argument it uses, or
quotes, results of many mathematicians. This section says which, and where. Two kinds of use should be
distinguished.

\textbf{(a) Published results used in the proofs given in this paper.}
\begin{itemize}
\item \emph{Chebyshev's method} (\S I.0, \S I.2): Chebyshev \cite{Che52}; the form with the middle binomial
coefficient is due to Erd\H{o}s \cite{Erd32}. The exponent of a prime in $n!$ is
due to Legendre \cite{Leg08}. Textbook accounts: \cite[Ch.~4]{Apo76}, \cite[Ch.~7]{Dav00}.
\item \emph{Landau's theorem} on Mellin transforms (or Dirichlet series) of functions of constant sign
(\S I.1): Landau \cite{Lan05}; see \cite[\S15.1]{MV07}. A second proof of the version on intervals, from a theorem of Kaczorowski and Pintz \cite{KP86} on changes of sign, is in \S I.1; we read the statement of that theorem, the beginning of its proof and its last step, and a reviewer compared the statement with the article. The version on intervals (Part II, \S1) is proved in
full in this paper, following the Lean development.
\item \emph{The explicit formula} and the value $\sum_\rho\frac1{\rho(1-\rho)}=2+\gamma-\log4\pi$ (\S I.1):
von Mangoldt; see \cite[Ch.~12 and 17]{Dav00}. The use of positivity in an explicit formula goes back to
Weil \cite{Wei52}.
\item \emph{The prime number theorem in arithmetic progressions}, for a fixed modulus, and
\emph{Mertens' theorem} in progressions (\S I.4.2, \S I.6; Part II, \S5): de la Vall\'ee Poussin; see
\cite[Ch.~20 and 22]{Dav00}, \cite[Ch.~11]{MV07}.
\item \emph{Theorems on Dirichlet $L$-functions used in Part III.} For Estimate B (\S III.1): the zero-free
region with at most one exceptional zero, Landau and Page \cite{Pag35}, \cite[Ch.~14]{Dav00}; primes in
progressions to large moduli, Thorner and Zaman \cite{TZ21}, after Linnik and Gallagher \cite{Gal70},
\cite{Bom87}; a zero-density estimate, Montgomery \cite{Mon71}; Landau's bound for $L'/L$ in a zero-free
region, \cite{Tit86}, \cite{MV07}. For Estimate C (\S III.2): the zero-free region, the count of zeros in a
window and the truncated explicit formula, \cite[Ch.~14, 16, 19]{Dav00}, \cite[Ch.~11, 12]{MV07}. These are
cited, not proved. The mean value over the auxiliary primes is obtained from primes in progressions and not
from a large sieve; compare \cite{Bom65}, \cite{DH66}, \cite{HB95}.
\item \emph{Sets without three-term progressions} (\S I.9, steps 11 and 19; Part III, \S III.3): Behrend
\cite{Beh46}. The construction is carried out, with proof, in \S III.3.
\item \emph{The filtration of a lattice by slopes} (Part III, \S\ref{as:sec}, (C6); used in the proof of
Proposition~M4): Stuhler \cite{Stu76}, Grayson \cite{Gra84}; cited from memory.
\item \emph{Stepanov's method} (\S I.2, \S I.7; Part II, \S2), as the model for counting with multiplicity:
Stepanov \cite{Ste69}; see Bombieri \cite{Bom74}.
\item \emph{The lemniscate functions} and their arithmetic (\S I.3): Hurwitz \cite{Hur99}. The main statements of
\S I.3 and \S I.4.1 about them are proved in \S\ref{el:sec}, which lists what is proved and what is only
computed; no novelty is claimed for them. \emph{Formal groups
with complex multiplication}, which explain the exponent $a\,p/(p^2-1)$ at inert primes: Lubin and Tate
\cite{LT65}, Honda \cite{Hon70}. That the reduction is of this type exactly at the inert primes: Deuring
\cite{Deu41}. Algebraicity of the lattice sums of \S I.4.1: compare Damerell \cite{Dam70}.
\item \emph{The Gaussian inner product on polynomials}, $\|z^a\|^2=a!$ (\S I.4.2; Part II, \S6.5): the
Fock, or Bargmann, space \cite{Bar61}.
\item Hadamard's inequality, the Chinese remainder theorem and the Cauchy--Schwarz inequality are used in
Part II, \S6 without further reference.
\end{itemize}

\textbf{(b) Results behind the quoted statements.} A census of the 195 statements that the Lean development
takes under the heading of published results is in \S\ref{cs:sec}. The statements marked \quoted\ are proved in the Lean
development, which contains its own proofs of the theorems it needs. The underlying mathematics has the
following sources. We cite them for the classical statements; we have not compared the formalised
statements with the published ones in detail.
\begin{itemize}
\item \emph{Theta series as kernels} and the representation that governs them (\S I.4.1, \S I.7): Weil
\cite{Wei64}; the correspondence between forms on a pair of groups: Howe \cite{How79}. \emph{Seesaw
identities}: Kudla \cite{Kud84}. \emph{The inner product formula} for theta lifts: Rallis \cite{Ral87}.
\emph{Bessel periods} and the restriction problems in which they occur: Gan, Gross and Prasad \cite{GGP12}.
\item \emph{The period formula} (\S I.7; Part II, \S6.6). In the form used, it is Kronecker's limit formula
(see \cite{Wei76}); the evaluation through Gamma values is due to Lerch, and to Chowla and Selberg
\cite{CS67}.
\item \emph{The mass formula} for definite lattices (\S I.7; Part II, \S6.6): Siegel \cite{Sie35}; for
Hermitian lattices, Shimura, see Gan, Hanke and Yu \cite{GHY01}. \emph{The Siegel--Weil formula}
(\S III.7, \S III.8): Weil \cite{Wei65}; the precise hypotheses and normalisations used are those stated in
those sections. \emph{The dimension formula}: Weyl
\cite{Wey25}.
\item \emph{Slopes of Euclidean lattices} and the maximal slope (\S I.7; Part II, \S6.8, \S6.10): Stuhler
\cite{Stu76}, Grayson \cite{Gra84}. \emph{The slope inequality} for maps of lattices: Bost \cite{Bos96}.
\item \emph{The relative trace formula} (\S I.7, ``Rank''): Jacquet \cite{Jac86}.
\item \emph{The bound for Galois cohomology} (\S I.7, ``Denominators''; Part II, \S6.8). The announcement
describes it as a weakening of Vandiver's conjecture. According to our reading of the Lean development its
proof uses the vanishing of the Iwasawa $\mu$-invariant (Ferrero and Washington \cite{FW79}): a statement
under that name is proved in the development and the final theorem depends on it. The development also
contains statements named after Kato's Euler systems \cite{Kat04}; some of these are among the declarations the
final theorem depends on and one that we tested is not. The proof
also uses modularity lifting theorems. We did not find the main conjecture of Iwasawa theory
(Mazur and Wiles \cite{MW84}) stated under that name. An upper bound for a cohomology length in terms of the
valuation of a Bernoulli number with explicit integral factors is present; it corresponds to a consequence of
the main conjecture, and we have not compared the two. A bound for the whole group $H^2$ at every prime is also
present. Finiteness of these groups is due to Soul\'e (\cite{Sou79}, from memory). Soul\'e
\cite[Thms.~2 and~3]{Sou03} gives a quantitative bound away from small primes: for the primes $p\ge4n'-1$,
$p\ge[L:\Q]+2$, where $L$ is the field and $n'$ is the twist (Soul\'e writes $n$; it is not the $n$ of this
paper). This was compared with the article by a reviewer, not by us. We have not
identified a published estimate in the full small-prime range. Part III, \S\ref{gc:sec} describes the statement and the structure of its proof in the development, and says
which published theorems the cases correspond to; we have not checked how closely the formal proofs follow the
published ones.
\end{itemize}

\textbf{What is new.} The combination described in this paper, namely iterated theta lifts as a source of
multiplicity, the comparison of divisibility with size through slopes of lattices carried to the term
linear in the order, and the resulting one-sided bound for $G(x)$, is the contribution of \cite{Alp26}.
The exposition, the worked examples and the computations in this paper are ours.

\section{Summary: what is proved where}

\begin{center}\small
\begin{tabular}{p{.47\textwidth}p{.45\textwidth}}
\toprule
Statement & Status\\
\midrule
One-sided bound for $G$ on intervals $\Rightarrow$ RH (\S I.1) & proved (Part II, \S1)\\
Model divisibility $a\,p/(p^2-1)$; lemniscate numbers (\S I.3) & proved for the model series (Part II, \S2) and for the lemniscate numbers (\S\ref{el:sec}, Theorem~Le18)\\
(M), B and C $\Rightarrow$ Estimate A $\Rightarrow$ the bound for $G$ (\S I.6) & proved (Part II, \S3--5); the identities used are checked numerically on an example\\
Estimate B: the auxiliary primes (\S I.6) & proved (Part III, \S III.1), from four cited theorems\\
Estimate C: the characters at the large scale (\S I.6) & proved (Part III, \S III.2), from three cited theorems and Theorem~\ref{thm:pair}(e)\\
The sets of columns $T_2$ and $T$ (\S I.5) & constructed (Part III, \S III.3); their use in (M) is through two hypotheses of Part III, \S\ref{as:sec}\\
(M), size (\S I.7, item 2) & the analytic estimates for one map and for a chain are proved in an abstract setting (Part III, \S III.4); that the maps of the construction are instances of it is quoted; four passages of \S III.4 are conditional; the arithmetic inputs are quoted\\
(M), divisibility at inert primes (\S I.7, item 1) & written up in \S III.5, conditionally on the integrality of the theta kernels in the coordinates of the height-two logarithm (hypothesis (B)), which is taken from the Lean development; elementary lemmas proved (Part II, \S6)\\
(M), from derivatives to lattices (\S I.7, item 1) & proved (Part III, \S III.6): lemmas on slopes of lattices, the splitting of the prime sum, bounds for several of its terms. Conditional on two statements of the Lean development: the theorem on the glued lattice. Quoted: the bounds for the eleven loss sums and their total; the statement on valuations\\
(M), arithmetic inputs of the size bound (\S I.7, item 2) & proved (Part III, \S III.7): the size inequality follows from eleven named inputs; a mean of periods equals $\log\Omega_E$ under stated hypotheses. The eleven inputs are quoted\\
(M), rank (\S I.7, item 3) & proved (Part III, \S III.8): the dimension bound follows from three statements of the Lean development and stated size conditions. The three statements are quoted; elementary lemmas proved (Part II, \S6)\\
(M), denominators (\S I.7, item 3) & proved (Part III, \S III.9): the step from the bound for Galois cohomology to the bound for denominators. The bound itself is quoted; elementary lemmas proved (Part II, \S6)\\
(M), the assembly & reduced to twelve numbered hypotheses (Part III, \S\ref{as:sec}); the hypotheses are derived from quoted statements, quoted, or assumed\\
(M) at every scale of the intervals, with the primes of Estimates B and C & a thirteenth hypothesis, assumed; with it the Riemann hypothesis follows from the hypotheses (Part III, \S\ref{as:sec}, Corollary~M11)\\
(M), the construction of the maps and weights & assumed (Part III, \S III.10)\\
Error budget and exponent $\frac12$ (\S I.8) & proved from the stated form of (M)\\
\bottomrule
\end{tabular}
\end{center}

\clearpage
\noindent\textbf{What remains to be written up.} For these pages to be a complete human-readable proof, the
following would still have to be written out. They are grouped by the part of the proof they serve; the
table of Hypotheses in Part~III, \S\ref{as:sec}, and the table at the start of Part~III give the details.
\begin{enumerate}[leftmargin=2em,itemsep=1pt]
\item \emph{The construction.} The field $E$ is explicit once $n$, $q$, $d$ are fixed (\S\ref{ob:sec}). Over
it one must construct the Hermitian spaces and lattices, the chain of theta lifts with its ranks, levels and
weights, and the translates and labels, with the properties of \S\ref{ob:sec}; verify the identifications
left open in the dictionary of \S\ref{ob:sec}, \S7; and supply these data at every scale of the intervals,
with constants independent of the scale and with auxiliary primes for which Estimate~B holds (Hypothesis
(H12) of \S\ref{as:sec}; Estimate~C then follows, \S III.2). \S\ref{ob:sec} gives no classical description
of the forms on the groups that are not compact.
\item \emph{Divisibility.} Hypothesis (B), the integrality of the theta kernels at inert primes in the
coordinates of the height-two logarithm (proved here only for one example, \S\ref{el:sec}); in \S III.5 the
duality input for longer orbits, the operator on the ordinary rows, the passage to all frames and labels, and
the remaining classes of primes; in \S III.6 the two statements (Val) and (Lat) behind the glued lattice, the
local comparisons (a)--(d) needed for Hypothesis (H6)(ii), and the maximal-slope comparison used in
Lemma~M10 for (H6)(iii); and the eleven losses with the required bound for their total (Hypothesis (H7)).
\item \emph{Size.} The eleven arithmetic inputs of the size inequality (\S III.7), among them the
inner-product formula for the lifts, the covolumes of the lattices and the local indices; that the maps of
the construction are instances of the setting of \S III.4; and the conditional passages there.
\item \emph{Levels, counts, covolumes and slopes.} The local level correction (H2); the comparison of the
heights of two words (H3); in (H4), the lower bound for the level with its allowance (a count of forms) and
the comparison of the levels of two words; the covolume of the first lattice of the chain, a lattice of forms (H5); and the
range of the slopes (H9). The comparison of heights (H3) is elementary in kind; the count of forms and the
covolume in (H5) are deep. (The upper bound in (H4) is a condition of the construction.)
\item \emph{Rank.} For the family of the construction: the moment estimates (S), which come from a relative
trace formula; the bound (I) for the projection loss, which the source relates to a count of good weight
tuples and to residual irreducibility, and which also gives (H10); the pairing statement (P), from
rationality and the identification with a Bessel pairing; and the rank, size and subspace conditions of
Corollary~T2, which also give (H11) and the first half of (H0)(g) (\S III.8). For the mean-slope
inequality (H8), the pairings of \S III.6 and \S III.8 and the map of \S III.7 must be shown to be one map,
with the same norms, so that the integrality, height and rank bounds apply to it together.
\item \emph{Denominators.} The bound for Galois cohomology, Theorem~G, and the comparison inequality that
bounds the total twist error at each lift in terms of cohomology lengths (\S III.9; Proposition~G2 deduces the
bound for denominators from the two). For several steps of the proof of Theorem~G in the Lean development
we found no published counterpart (\S III.9, \S6, list~(C)).
\item \emph{The published results.} For each substantive outside result used in the prose proof, a
checked reference with matching hypotheses, or a proof. The census of \S\ref{cs:sec} records the 195 formal
entries that the Lean development gathers under this heading, including repetitions and one unused slot;
the bibliographic matches proposed there are from memory and have not been checked systematically, and
several of the entries are weak forms of the published theorems (for instance the local Langlands
correspondence, \S III.9).
\end{enumerate}
\noindent What is complete: the deduction of the Riemann hypothesis from (M) and Estimates~B and~C (Parts~I
and~II); Estimates~B and~C, from cited theorems on Dirichlet $L$-functions (\S III.1--III.2); the sets of
columns (\S III.3); and the deduction of (M), and of the Riemann hypothesis, from the thirteen Hypotheses
(\S\ref{as:sec}): (H0)--(H11) at one sufficiently large scale give (M) there, and (H12) supplies them on the
required intervals. So the census of \S\ref{cs:sec} is not the only gap. Items~1--6 carry the substance of
(M), and they are where most of the Lean development goes.

\medskip
\noindent The formal proof in Lean establishes the Riemann hypothesis for all Dirichlet $L$-functions at
once, and with it the statements quoted here, in its own formulation. The case of $\zeta$ described here is
the case of the trivial character; it uses the same construction.

\medskip
\noindent Sources (the Lean files) are at the end of Part II; the list of scripts follows Part III.

\newpage
\setcounter{equation}{0}\renewcommand{\theequation}{II.\arabic{equation}}
\stepcounter{partno}\begin{center}\Large\bfseries Part II. Details and proofs\end{center}
\phantomsection\addcontentsline{toc}{part}{Part II. Details and proofs}

\noindent Part II can be read after \S I.1 and \S I.5: it uses $G(x)$ from \S I.1 and otherwise repeats its
definitions. It is numbered separately. Its sections are 1--6. Its displays are numbered (II.$n$), except
two that carry names, (A) and (M).\footnote{In the version of 21 November 2026 eighteen displays of Part II carried section numbers instead; they are now numbered as follows: (1.1)$\to$(II.3), (4.1)$\to$(II.14), (4.2)$\to$(II.16), (5.1)$\to$(II.21), (5.2)$\to$(II.22), (5.3)$\to$(II.24), (5.4)$\to$(II.29), (5.5)$\to$(II.30), (5.6)$\to$(II.33), (5.7)$\to$(II.36), (6.1)$\to$(II.47), (6.1a)$\to$(II.48), (6.2)$\to$(II.50), (6.3)$\to$(II.54), (6.4)$\to$(II.57), (6.5)$\to$(II.58), (6.6)$\to$(II.59), (6.7)$\to$(II.69).} The main inequality is called (M) in all three parts, and Estimate A of
Part I is formula (A) of Part II. Section 3 of Part II states the
reduction to one prime sum again, with all definitions, so that sections 4--6 can be read on their own.
References in square brackets with numbers are to the list of source files at the end of Part II. Some letters are
used with more than one meaning in Part II (for instance $N$, $L$, $T$, $h$ and $a$); each is defined where it
is used. Formula \eqref{II:t5.6} of Part II is Estimate C of Part I.
Where sections 3--5 attribute a statement to ``the source'', a proof is now in Part III for the following:
the choice of the primes $q$, $d$ and the bound for the three twists (section 3; Part III, \S III.1);
formula \eqref{II:t5.6} (Part III, \S III.2); the existence of the sets $T$, $T_2$ with the stated densities
(section 4; Part III, \S III.3). The divisibility at inert primes of sections 6.3--6.4 is written up in Part III, \S III.5, conditionally on the integrality of the theta kernels in the coordinates of the height-two logarithm (hypothesis (B)), which is taken from the Lean development. Part III, \S III.6--III.9 prove further steps towards (M), each from named statements of the Lean development: the passage to lattices in part (\S III.6), the size inequality from eleven quoted inputs (\S III.7), the dimension bound of section 6.9 from three quoted statements (\S III.8), and the step from the bound for Galois cohomology to the bound for denominators (\S III.9); Part III, \S\ref{as:sec} deduces (M), at one scale, from twelve numbered hypotheses, and the Riemann hypothesis from these and a thirteenth; the table at the start of Part III gives the status of each. The main inequality (M) is still not proved: it is quoted in Part II.

{\setlength{\parindent}{0pt}\setlength{\parskip}{.32\baselineskip plus 1pt}
\section*{1. What a zero forces}\addcontentsline{toc}{section}{1. What a zero forces}

For $\re s>1$, termwise integration gives
\begin{equation}\label{D:38}
 F(s)=\int_1^\infty G(x)x^{-s-1}\,dx
 =-\frac{\zeta'(s)}{s(s-1)\zeta(s)}
   -\frac1{(s-1)^2}+\frac{1+\gamma}{s-1}.
\end{equation}
Indeed, the transform of $(x/n-1)_+$ is $n^{-s}/[s(s-1)]$, where $v_+=\max(v,0)$. The Laurent expansion of $\zeta$ at $1$ cancels both terms of the principal part there. A zero $\rho$ of multiplicity $\nu$ in the critical strip gives $F$ a simple pole of residue $-\nu/[\rho(\rho-1)]$. The poles at $s=0$ and at the trivial zeros $-2,-4,\ldots$ lie outside the half-plane $\re s>\theta$ used below. The poles in the critical strip are nonreal, since the alternating zeta series gives $\zeta(\sigma)<0$ for $0<\sigma<1$.

A bound valid for every large $x$ would allow the usual oscillation argument of Landau. We know the bound only on certain intervals. To use one such interval, we differentiate the transform many times. Its new weight is concentrated on that interval, and a nearby pole forces large derivatives.

\textbf{Lemma (the one-sided criterion on intervals; \S I.1).} \textit{Let $0<\theta<1$. If $G(x)\le x^\theta$ throughout $[Y,Y^{3/2}]$ for arbitrarily large $Y$, then $\zeta$ has no zero in $\re s>\theta$.}

Write $u=\log x$, $g(u)=G(e^u)$ and, for each integer $j\ge0$, let $A_j(s)=(-1)^jF^{(j)}(s)/j!$ be the signed $j$th Taylor coefficient. Let $\eta=\max(0.8e^{0.2},1.2e^{-0.2})=0.9825\ldots$ and $c=0.008$, which is less than $\log((1+\eta)/(2\eta))=0.00887\ldots$; these are absolute constants. Fix $\sigma>2$ so large that $(\sigma-\theta)/(\sigma-2)\le e^{c}$. With the constants above this holds for $\sigma\ge252$ and every $\theta$ in $(0,1)$. Differentiating the Mellin integral gives
\begin{equation}\label{D:39}
 A_j(\sigma+it)=\int_0^\infty g(u)e^{-itu}K_j(u)\,du,
 \qquad K_j(u)=u^je^{-\sigma u}/j!.
\end{equation}
On each interval satisfying the lemma's bound, put $U=\log Y$ and choose $j$ nearest to $\tfrac54(\sigma-2)U$. Then
\begin{equation}
 |A_j(\sigma+it)|+A_j(\sigma)
 \le 2\int_0^\infty \max(g(u),0)K_j(u)\,du
 \ll (\sigma-\theta)^{-j},                                      \label{II:t1.1}
\end{equation}
uniformly for real $t$. The first inequality follows from the triangle inequality and $|g|+g=2\max(g,0)$. On $[U,3U/2]$ integrate $e^{\theta u}K_j(u)$.

Outside the window, use $|g(u)|\ll e^{2u}$, since partial summation of Chebyshev's bound $\sum_{n\le x}\Lambda(n)\ll x$ gives $\sum_{n\le x}\Lambda(n)/n\ll\log x$ and hence $|G(x)|\ll x\log x$ for $x\ge2$. The density proportional to $u^je^{-(\sigma-2)u}$ is concentrated about $j/(\sigma-2)\sim5U/4$.

\textbf{Claim.} \textit{For $U$ large and $j$ nearest to $\tfrac54(\sigma-2)U$, the part of $\int_0^\infty u^je^{-(\sigma-2)u}\,du$ outside $[U,3U/2]$ is at most $C_0e^{-cj}$ times the whole, with $C_0$ independent of $U$.} Indeed, on putting $v=(\sigma-2)u/j$ the density is proportional to $(ve^{1-v})^j$, and the window becomes an interval in $v$ whose ends tend to $4/5$ and $6/5$ as $U\to\infty$. The base $ve^{1-v}$ has its unique maximum $1$ at $v=1$. Off the window it is bounded by a number $\eta'<1$ that tends to $\eta$ as $U\to\infty$, so the tail integral is at most $\eta'^{\,j-1}\int_0^\infty ve^{1-v}\,dv$. A small fixed interval about $1$ has base at least $(1+\eta')/2$. So the fraction outside is $O\big((2\eta'/(1+\eta'))^j\big)$, and this is $O(e^{-cj})$ for $U$ large, because $c<\log((1+\eta)/(2\eta))$. The tail is therefore $O(e^{-cj}(\sigma-2)^{-j-1})=O((\sigma-\theta)^{-j})$, by the choice of $\sigma$. This proves \eqref{II:t1.1} along a sequence $j\to\infty$.

Suppose now there is a pole to the right of $\theta$. Let $R$ be the smaller of $\sigma-\theta$ and the distance from $\sigma$ to the nearest pole in $\re s>0$. Some pole has horizontal distance from the line $\re s=\sigma$ strictly less than $R$. If $R=\sigma-\theta$, this follows from the supposition. Otherwise take a nearest pole to $\sigma$ and use its nonzero imaginary part. Take $t_0=\operatorname{Im}\rho$ for that pole $\rho$; then the pole nearest to $\sigma+it_0$ is at distance less than $R$. Only finitely many poles lie within $R+1$ of it, and each pair is equidistant from at most one point of the line $\re s=\sigma$ (both lie to its left). So $t$ near $t_0$ can be chosen so that $z=\sigma+it$ has a unique nearest pole, at distance $d<R$.

Choose $d<r<R$ so that the disk of radius $r$ about $z$ contains just that pole. Subtracting its principal part and applying Cauchy's estimate gives $|A_j(z)|\ge c_1 d^{-j}-O(r^{-j})$ with $c_1>0$ (the principal part $c/(s-\rho)$ contributes exactly $c\,(z-\rho)^{-j-1}$ to $A_j(z)$). At the real centre $\sigma$ the same estimate gives $|A_j(\sigma)|\ll r^{-j}$. This contradicts \eqref{II:t1.1} for large $j$, since $d<r<\sigma-\theta$. This proof uses only the poles in two bounded disks. The Lean proof [1] reaches the same statement by a slightly different route.

\newpage
\section*{2. How divisibility counts primes}\addcontentsline{toc}{section}{2. How divisibility counts primes}

Start with a polynomial $P(X)$ of degree $D$. If $P$ has distinct roots $\alpha_1,\ldots,\alpha_s$, each of multiplicity at least $m\ge1$, then $\prod_{j=1}^s(X-\alpha_j)^m$ divides $P$. Comparing degrees gives $ms\le D$.

Stepanov uses the same count for a \textbf{nonzero rational function on a smooth projective curve over a finite field}. Count points over the algebraic closure of that field (the curve is assumed geometrically irreducible), and include the points at infinity. For a nonconstant function the zeros and poles are the inverse images of $0$ and of infinity under its map to the projective line. Each inverse image, counted with multiplicity, has the degree of that map. Thus the total multiplicity of the zeros equals the total order $D$ of the poles. Hence at most $\lfloor D/m\rfloor$ points can be zeros of order at least $m$. On the projective line the function is a quotient of polynomials. A polynomial of positive degree $D$ has just one pole, of order $D$ at infinity.

Over the integers, write $v_p(c)$ for the exponent of a prime $p$ in a nonzero rational number $c$ (numerator exponent minus denominator exponent), and put $v_p(0)=+\infty$. The corresponding count is
\begin{equation}\label{D:41}
 \sum_p e_p\log p\le\log|\Delta|,\qquad
 \Delta\in\mathbb Z\setminus\{0\},\quad 0\le e_p\le v_p(\Delta).
\end{equation}
To bound the left side one needs an independent upper bound for $|\Delta|$. For a square integer matrix of nonzero determinant $\Delta$, for example, divisibility gives a lower bound, while the lengths of its columns give an upper bound by Hadamard's inequality. Section 6 explains the corresponding comparison for lattices. (The file [2] states such an inequality for an earlier form of the argument. That form is stated in the archive but, according to our dependency check, is not used by the final theorem.)

\textbf{Why odd powers of a prime occur.} Here is the local model in rational arithmetic. For a prime $p$ form the series
\begin{equation}\label{D:42}
 l(X)=X+\sum_{j\ge1}X^{p^{2j}}/p^j.
\end{equation}
It has a unique inverse under substitution, $U(z)=z+\cdots$, given recursively by $U=z-\sum_{j\ge1}U^{p^{2j}}/p^j$. Put $L=p^2-1$. We claim that the coefficient of $z^a$ in $U$ has valuation at least $-(a-1)/L$ for each integer $a\ge1$. This holds for $a=1$. In a product of $r=p^{2j}$ copies of $U$, the degrees $i_1,\ldots,i_r$ contributing to $z^a$ are positive and sum to $a$. Each is less than $a$. By induction their coefficients have product of valuation at least $-(a-r)/L$. Division by $p^j$ costs $j\le(r-1)/L$, proving the claim.

For a series $f$ with integer coefficients, the same product calculation gives the bound $-(a-1)/L$ for the coefficient of $z^a$ in $f(U(z))$. Its $a$th derivative at zero is $a!$ times this coefficient. Counting multiples of $p,p^2,\ldots$ in $a!$ gives $v_p(a!)=\sum_{h\ge1}\lfloor a/p^h\rfloor=a/(p-1)+O(1+\log_p a)$, where $\lfloor y\rfloor$ is the greatest integer at most $y$ and $\log_p a=\log a/\log p$. The derivative therefore has valuation at least
\begin{equation}\label{D:43}
 a\left(\frac1{p-1}-\frac1{p^2-1}\right)-O(1+\log_p a)
 =a\left(\frac1p+\frac1{p^3}+\frac1{p^5}+\cdots\right)-O(1+\log_p a).
\end{equation}
The series models differentiation on an elliptic curve with complex multiplication (extra endomorphisms) at primes remaining prime in its imaginary quadratic field. The source proves a weighted version for its arithmetic maps. At odd powers $p^h$ it yields the term $(a/p^h-w)_+$ for a positive weight $w$ (quoted; see section 6.4). When $w=b$ and $a=bx$, this is precisely $b(x/p^h-1)_+$. [3]

\newpage
\section*{3. The reduction to one prime sum}\addcontentsline{toc}{section}{3. The reduction to one prime sum}

For an odd prime $D$ let $\chi_D(r)$ equal $0$ if $D\mid r$, $1$ if $r$ is a nonzero square modulo $D$, and $-1$ otherwise. For a nonprincipal Dirichlet character $\chi$ put
\begin{equation}\label{D:44}
 c_\chi=\frac{L'}{L}(1,\chi),\qquad
 G_\chi(x)=\sum_{r\le x}\Lambda(r)\chi(r)(x/r-1)+c_\chi x .
\end{equation}
A Dirichlet character is a homomorphism from the invertible classes modulo an integer to the nonzero complex numbers, extended by zero to other classes. It is nonprincipal if it takes a value different from $1$ on the invertible classes. The function $L(s,\chi)$ is the continuation of $\sum_{r\ge1}\chi(r)r^{-s}$ from $\re s>1$, and the prime denotes differentiation in $s$.

For distinct odd primes $q,d$ write $\eta=\chi_q\chi_d$ for the pointwise product. Consider the single prime sum
\begin{equation}\label{D:45}
 S_{q,d}(x)=\sum_{r\le x}\Lambda(r)(1+\chi_q(r))(1-\chi_d(r))(x/r-1).
\end{equation}
Away from $q,d$, the coefficient of $\Lambda(r)$ before smoothing is $4$ if $r$ is a square modulo $q$ and a nonsquare modulo $d$, and zero otherwise. The exact identity obtained by expanding the two factors is
\begin{equation}\label{D:46}
 \begin{aligned}
 R_{q,d}(x)&=S_{q,d}(x)-x\log x+(1+\gamma+c_{\chi_q}-c_{\chi_d}-c_\eta)x\\
           &=G(x)+G_{\chi_q}(x)-G_{\chi_d}(x)-G_\eta(x).
 \end{aligned}
\end{equation}
Thus a one-sided bound for $R_{q,d}$ and absolute bounds for three quadratic character sums suffice.

\textbf{The reduction.} Fix $0<\varepsilon<1/52$. By Theorems~\ref{thm:pair} and~\ref{thm:C} of Part III, for every large $Y$ there are primes $q\equiv3\bmod8$, $d\equiv7\bmod16$ with $q\mid d+1$ for which statement 2 below (Estimate B) and formula \eqref{II:t5.6} (Estimate C) hold for $Y\le x\le Y^{3/2}$. Statement 1 (Estimate A) is then deduced in sections 4 and 5 from the main inequality (M), which is quoted, together with Estimates B and C. The argument of sections 4 and 5 is for one $x$ at a time, and uses (M), Estimate B and Estimate C at that $x$ for the same $q$, $d$, $n$. So nothing in it changes if the pair depends on $x$. In the source, and in Hypothesis (H12) of Part III, \S\ref{as:sec}, the intervals are $[Z,Z^{3/2}]$ with $Z$ arbitrarily large, and the pair used at $x$ is chosen at a scale $Y$ with $x^{1/J}\le Y\le x$, where $J\ge1$ depends on $\varepsilon$ only; $Y$ need not be the left end $Z$, and the pair may depend on $x$. For a pair given by Theorem~\ref{thm:pair} with $A=J$ at the scale $Y$, Estimates B and C hold at every $x$ with $Y\le x\le Y^J$; that the construction is made at $x$ with such a pair is part of (H12). Part III, \S\ref{as:sec}, \S7 states exactly what is needed, as (H12), and repeats the deduction with it (Corollary~M11).

\par\noindent\textbf{1.}\enspace \textbf{The residue sum.} There are positive integers $m,k,n$ depending on $x$ such that
\begin{equation}\label{D:47}
    R_{q,d}(x)\le C_\varepsilon x^{C\varepsilon}
            \left(\frac{x}{m}+\frac{xm}{k}+m+\frac{k}{n^2}\right). \tag{A}
\end{equation}
   Here $C$ is absolute and $C_\varepsilon$ depends only on $\varepsilon$. The orders are $m=x^{1/2+O(\varepsilon)}$, $k=x^{9/2+O(\varepsilon)}$, $n=x^{2+O(\varepsilon)}$. Sections 4 and 5 derive this from (M).
\par\noindent\textbf{2.}\enspace \textbf{The three twists.} For $\chi=\chi_q,\chi_d,\eta$ the selection gives $|G_\chi(x)|\le x^{1/2+\varepsilon}$. To see what is estimated, put $\Psi_\chi(t)=\sum_{r\le t}\Lambda(r)\chi(r)$. Partial summation and $\sum_{r\ge1}\Lambda(r)\chi(r)/r=-c_\chi$ give
\begin{equation}\label{D:48}
     G_\chi(x)=-x\int_x^\infty \Psi_\chi(t)\,\frac{dt}{t^2}.
\end{equation}
   The convergence here follows from the prime number theorem for a fixed nonprincipal character. The proof bounds the proportion of bad choices by mean-square estimates, first in $q$, then in $d$ for each good $q$. This gives a common pair on the whole interval. It is given in Part III, \S III.1, with the exponent $\frac12+\frac\varepsilon2$; the Lean proof is in [9, 10, 29].

\textbf{Conclusion.} The terms $x/m$ and $m$ balance at $\sqrt{x}$. With the stated orders, $k/n^2=x^{1/2+O(\varepsilon)}$ and $xm/k$ is smaller. By the identity above, $G(x)\le C'_\varepsilon x^{1/2+C'\varepsilon}$ on these intervals, with an absolute $C'$. Given $1/2<\theta<1$, choose $\varepsilon$ so that $1/2+C'\varepsilon<\theta$. On all sufficiently large selected intervals $G(x)\le x^\theta$. Section 1 excludes every zero to the right of $\theta$. Letting $\theta$ decrease to $1/2$ and using the symmetry $\rho\mapsto1-\rho$ from the functional equation proves RH from (M) and Estimates B and C. [1, 6, 12]

\newpage
\section*{4. The main inequality (M)}\addcontentsline{toc}{section}{4. The main inequality (M)}

\subsection*{4.1 Which prime powers are counted?}

Fix $x$ and primes $q,d$ for which Estimates B and C hold at $x$ (section 3). For an odd prime power $r=p^h$ with $p\ne q$, attach two labels. The first is the row $v(r)=(1+\chi_q(r))/2$, equal to $1$ if $r$ is a nonzero square modulo $q$, and $0$ otherwise. For the second, choose $n=2^e$ at the scale $x^{2+8\varepsilon}\le n\le 2x^{2+8\varepsilon}$ and put $N=4n$. Write $[r]=\{r,-r\}\bmod N$. The $n$ pairs of odd residues form a group $H$ under multiplication. Its elements will label columns. Choosing a fixed $t\in H$ puts $r$ in column $t[r]$.

The coefficient at $r=p^h$, before inserting a weight, is
\begin{equation}\label{D:49}
 c(r)=(1-\chi_d(r))\log p.
\end{equation}
If $p\ne d$ it is $2\log p$ for an odd power of a nonsquare modulo $d$ and zero otherwise. At $p=d$ it is $\log d$ for every power. The two omitted primes $2,q$ have $\chi_d(2)=\chi_d(q)=1$. Indeed $d\equiv7\bmod 8$ gives the first identity. For the second, $q,d\equiv3\bmod 4$ and $d\equiv-1\bmod q$, so quadratic reciprocity gives $(q/d)=-(d/q)=-(-1/q)=1$. So $c(r)$ vanishes at every power of $2$ and of $q$; omitting these primes changes nothing.

\textbf{A table small enough to write out.} Take $q=11$, $d=1319$, $N=16$ and $t=[1]$. Then $q\equiv3\bmod8$, $d\equiv7\bmod16$ and $q\mid d+1=1320$. The four columns are $[1],[3],[5],[7]$. For example $[3]^2=[9]=[7]$ and $[3]^3=[5]$. The nonzero squares modulo $11$ are $1,3,4,5,9$. In each cell below are the primes at most $61$ with nonzero coefficient; the primes $3,5,7,19,31,43,53$ are squares modulo $1319$ and have coefficient $0$. A dash denotes an empty cell in this range.
\begin{equation}\label{D:50}
 \begin{array}{c|cccc}
 v\backslash s &[1]&[3]&[5]&[7]\\ \hline
 0&17&13,29,61& - &41\\
 1&47& - &37,59&23
 \end{array}
\end{equation}
For each cell, the primes in that cell with nonzero coefficient have density $1/16$ among all primes, as shown in section 5. To place $13$, observe that $13\equiv2\bmod11$ and $[13]=[3]$. To place a higher power one repeats both tests on that power. The prime $23$ goes in $(1,[7])$ with coefficient $2\log 23$. Its square disappears because $\chi_{1319}(23^2)=1$. Its cube has the same row and column as $23$. At the exceptional prime $1319$, its first power lands in $(0,[7])$, its square in $(1,[1])$, and each has coefficient $\log 1319$. The example only illustrates the identities of section 5, which hold for any positive weights; its parameters do not satisfy all the conditions of (M) (Part~III, \S\ref{ob:sec}, \S3). The first version of this paper used $q=3$, $d=23$.

\newpage
\subsection*{4.2 How a weight determines a cutoff}

Let $b=\exp(x^\varepsilon)$. The construction supplies nonempty subsets $T,T_2$ of $H$, an integer $a_t>0$ for each $t\in T$ with $|a_t/b-x|\le1$, and $m$ maps of initial rank $k$ (the chain of theta lifts of section 6.2; $m$ and $k$ are the parameters in (A)). Write $|S|$ for the cardinality of a finite set $S$ and $\delta=|T_2|/n$. For each translate $u\in H$ and map $1\le j\le m$ there are $2n$ positive weights $w_{u,j}(v,s)$. In the construction they are integers. The distinguished cells are $(1,s)$ with $s\in uT_2$. Their weights satisfy $b\le w_{u,j}(1,s)\le b+O(b/(nx^2))$. At every other cell the ratio $a_t/w_{u,j}(v,s)$ is at least $\exp(x^{\varepsilon/3})/3$. Thus the same table tests primes at two markedly different scales.

Fix $u,t,j$, write $a=a_t$ and $w=w_{u,j}$. Each power in cell $i=(v,s)$ contributes
\begin{equation}\label{D:51}
 c(r)(a/r-w_i)_+
 =c(r)w_i\big(a/(rw_i)-1\big)_+,\qquad z_+=\max(z,0).
\end{equation}
The contribution ends at $r=a/w_i$. Adding all cells gives the finite sum
\begin{equation}
 P_{u,t,j}=\sum_{\substack{p\ {\rm odd\ prime}\\p\ne q}}\ \sum_{h\ge1}
 (1-\chi_d(p^h))\log p\,
 \big(a_t/p^h-w_{u,j}(v(p^h),t[p^h])\big)_+ .                 \label{II:t4.1}
\end{equation}
In the running example take $T=T_2=\{[1]\}$, one table, $x=1000$, $b=10$, $a=10000$ and $u=[1]$. Here $\delta=1/4$. The weights and cutoffs, in the same column order as above, are
\begin{equation}\label{D:53}
 \begin{array}{c|rrrr}
 w_i &[1]&[3]&[5]&[7]\\ \hline
 0&0.05&0.05&0.05&0.05\\
 1&10&0.05&0.05&0.05
 \end{array}\qquad
 \begin{array}{c|rrrr}
 a/w_i &[1]&[3]&[5]&[7]\\ \hline
 0&200000&200000&200000&200000\\
 1&1000&200000&200000&200000
 \end{array}
\end{equation}
For example $23$ contributes $2\log 23\,(10000/23-0.05)$, whereas $47$ contributes $2\log 47\,(10000/47-10)$, and the prime $1319$ contributes $\log1319\,(10000/1319-0.05)$. Translating by $u=[7]$ moves the weight $10$ to column $[7]$ in row $1$. Now the subtracted weight at $23$ is $10$, and at $47$ it is $0.05$.

For a reproducible finite calculation run through the primes $p\le200000$ and the successive powers $p,p^2,\ldots\le200000$. Drop $2,11$, compute the coefficient and the cell as above, and add its positive part. This gives $P_{[1],t,1}=84701.0317$ to four decimals. The same number will be obtained by grouping the terms in a different order in section 5.

\newpage
\subsection*{4.3 What the arithmetic construction must prove}

The quantity to subtract from \eqref{II:t4.1} is
\begin{equation}
 \langle\log w\rangle=\frac{1}{2n}\sum_i\log w_i,\qquad
 M_{u,t,j}=a_t\big(\log a_t-1-\gamma+c_{\chi_d}-\langle\log w\rangle\big).
                                                                 \label{II:t4.2}
\end{equation}
Section 5 derives each term by ordinary prime summation. Section 6 explains why the arithmetic estimate on the maps yields the same expression. The match through the term linear in $a_t$ is essential.

Every translate has seven weights equal to $0.05$ and one equal to $10$. Substitution in \eqref{II:t4.2} gives
\begin{equation}\label{D:55}
 \begin{aligned}
 M&=10000\big(\log 10000-1-\gamma-1.520227887516\ldots\\[-2pt]
 &\hspace{39mm}-(\log 10+7\log 0.05)/8\big)=84463.3942\ldots .
 \end{aligned}
\end{equation}
For $u=[1]$ the excess is $237.6375$ to four decimals. The other three excesses are $-300.2570$, $-176.7314$, and $229.1329$, giving a mean of $-2.5545$. (M) bounds their average over $u$, summed over the labels $t$ and the maps $j$.

\textbf{The main inequality (M).} At every $x$ of intervals $[Z,Z^{3/2}]$ with $Z$ arbitrarily large the source constructs the above data, for a pair $q,d$ attached to $x$ (Part III, \S\ref{as:sec}, Hypothesis (H12)), with $\delta\ge x^{-\varepsilon}$, $|T|/n\ge x^{-\varepsilon/8}/2$, and an upper error $\mathcal E\ge0$. (Sets $T_2$, $T$ with these densities are constructed in Part III, \S III.3. The rest of the data is quoted.) With $B=|T|mb\delta/2$ it proves
\begin{equation}\label{D:56}
 \frac1n\sum_{u\in H}\sum_{t\in T}\sum_{j=1}^m(P_{u,t,j}-M_{u,t,j})\le\mathcal E,\qquad
 \frac{\mathcal E}{B}\le x^{C\varepsilon}
 \left(\frac{x}{m}+\frac{xm}{k}+m+\frac{k}{n^2}\right).          \tag{M}
\end{equation}
Here $C$ is absolute, and the inequality is asserted for $x\ge x_0(\varepsilon)$. The first inequality compares integral divisibility with complex size. To check the powers in the second, put $h=|T|$, $h_2=|T_2|$ and $a=xb$. The source proves more than the mean: for its data at an admissible scale $x$ and every translate $u\in H$,
\begin{equation}\label{D:57}
 \sum_{t\in T}\sum_{j=1}^m(P_{u,t,j}-M_{u,t,j})\le x^{(C_{\rm abs}+e_L+3)\varepsilon}\left[
 h a(1+m^2/k)+m\left(2nmbx^\varepsilon+
                 \frac{h_2kb x^\varepsilon}{n(n x^{-\varepsilon})}\right)\right]
\end{equation}
for $x\ge x_0(\varepsilon)$, where $C_{\rm abs},e_L\ge0$ are constants of the source that do not depend on $x$ or $\varepsilon$ ($e_L=0$ on the route of the final theorem; the value of $C_{\rm abs}$ was not determined). This is the statement the final theorem uses. Part III, \S\ref{as:sec}, Theorem~M deduces this inequality from the Hypotheses (H0)--(H11) of that section, which are not proved in this paper, as soon as $x^\varepsilon\ge79(1+K)^2(1+C_H)(1+e_K)$ for the constants $K,C_H,e_K$ of the Hypotheses. The right side does not depend on $u$, so the mean over $u$ satisfies the same bound; call it $\mathcal E$. Division by $B=hmb\delta/2$ bounds the three bracketed contributions respectively by
\begin{equation}\label{D:58}
 \frac2\delta\left(\frac{x}{m}+\frac{xm}{k}\right),\quad
 \frac{4n}{h\delta}mx^\varepsilon,\quad
 2\frac{n}{h}\frac{k}{n^2}x^{2\varepsilon}.
\end{equation}
Here $h_2=n\delta$. The bounds $\delta\ge x^{-\varepsilon}$ and $h/n\ge x^{-\varepsilon/8}/2$ give (M) with $C=C_{\rm abs}+e_L+6$. This computation explains why multiplicity matters. A cost incurred once yields $x/m$ after division, while the lattice comparison incurs a cost of order $m$ after division. Their balance occurs at $\sqrt{x}$. [3, 5, 6]

\newpage
\section*{5. From the table to a single prime sum}\addcontentsline{toc}{section}{5. From the table to a single prime sum}

\subsection*{5.1 Subtracting the main term cell by cell}

Keep $t$ fixed; the quantities $S_i$, $\Psi_i$, $C_i$, $R_i$ below depend on $t$, through the column $t[r]$. Let $\sum_{r\in i}$ denote a sum over prime powers assigned to cell $i$ by section 4. The smoothed sum in that cell and its unsmoothed counting function are
\begin{equation}\label{D:59}
 S_i(y)=\sum_{r\in i}c(r)(y/r-1)_+,\qquad
 \Psi_i(z)=\sum_{\substack{r\in i\\r\le z}}c(r).
\end{equation}
For $y\ge1$, integrating each step of $\Psi_i$ proves the first identity below. The second follows term by term from \eqref{II:t4.1}.
\begin{equation}
 S_i(y)/y=\int_1^y\Psi_i(z)\frac{dz}{z^2},\qquad
 P_{u,t,j}=\sum_i w_iS_i(a/w_i).                              \label{II:t5.1}
\end{equation}
For fixed moduli, among primes not dividing $2qd$ the fraction contributing to a given cell is $1/(4n)$. One half are nonsquares modulo $d$, one half have the specified sign modulo $q$, and $1/n$ have the specified pair modulo $N$. The Chinese remainder theorem makes these conditions independent among the invertible residue classes. The prime number theorem in progressions gives their density among primes. The coefficient $2\log p$ doubles this density in $\Psi_i$. Thus $\Psi_i(z)$ has main term $\mu z$, where $\mu=1/(2n)$.

The fixed-modulus prime number theorem gives an error $O_A(z/(\log z)^A)$ as $z\to\infty$ for each fixed $A>0$. Higher prime powers and the one exceptional prime fit within this bound. The constants may depend on the moduli.

Set $C_i=\int_1^\infty(\Psi_i(z)-\mu z)\,dz/z^2$, which converges by that estimate. Subtract the integral from $y$ to infinity to obtain the exact formula
\begin{equation}
 S_i(y)=\mu y\log y+C_i y+R_i(y),\qquad
 R_i(y)=-y\int_y^\infty(\Psi_i(z)-\mu z)\frac{dz}{z^2}.        \label{II:t5.2}
\end{equation}
In particular $R_i(y)=o(y)$ for fixed moduli, meaning $R_i(y)/y\to0$. The uniform estimate needed when the moduli vary with $x$ comes later.

What is $\sum_i C_i$? Write $S_1(y)$ for the ordinary smoothed von Mangoldt sum and $S_\chi(y)$ for the same sum with coefficient $\chi(r)$. Then
\begin{equation}\label{D:62}
 \sum_i S_i(y)= S_1(y)-S_{\chi_d}(y).
\end{equation}
This identity includes powers of $d$, and both sides vanish termwise at $p=2,q$. For a nonprincipal character, partial summation and the convergent series for $-L'(1,\chi)/L(1,\chi)$ give $S_\chi(y)/y\to-c_\chi$. For the ordinary sum, the constant in $\sum_{r\le y}\Lambda(r)/r-\log y$ is $-\gamma$, by the Laurent expansion of $-\zeta'/\zeta$ at $1$ and partial summation. Subtracting $y^{-1}\sum_{r\le y}\Lambda(r)\to1$ gives the linear coefficient $-1-\gamma$. Comparison with \eqref{II:t5.2} yields
\begin{equation}
 \sum_i C_i=-1-\gamma+c_{\chi_d},\qquad
 P_{u,t,j}-M_{u,t,j}=\sum_i w_iR_i(a/w_i).                   \label{II:t5.3}
\end{equation}
For the last equality expand $\mu a\sum_i\log(a/w_i)=a\log a-a\langle\log w\rangle$ and insert the sum of the constants.

\newpage
\subsection*{5.2 Arithmetic check of the centering}

For $q=11,d=1319,n=4,t=[1]$, every cell has $\mu=1/8$. The constants and sums below allow \eqref{II:t5.1}--\eqref{II:t5.3} to be checked on the individual cells. Here $y_0=1000$ and $y_1=200000$.
\begin{equation}\label{D:64}
 \begin{array}{cc|r|rr|rr}
 v&s&C_i&S_i(y_0)&S_i(y_1)&R_i(y_0)&R_i(y_1)\\ \hline
 0&[1]&-.260399&622.935&253248.083&19.864&176.026\\
 0&[3]& .185272&1005.562&341839.067&-43.179&-367.240\\
 0&[5]&-.769229&158.362&150230.420&64.122&-1075.524\\
 0&[7]&-.356666&476.905&234628.730&-29.898&810.184\\
 1&[1]&-.538998&352.750&198411.282&28.278&1059.057\\
 1&[3]&-.574173&256.299&189878.990&-32.998&-438.252\\
 1&[5]&-.443019&404.588&217066.374&-15.862&518.334\\
 1&[7]&-.340232&542.738&236579.027&19.500&-526.431
 \end{array}
\end{equation}
Summing the eight constants before rounding gives $-3.097443552\ldots=-1-\gamma+c_{\chi_{1319}}$. For $u=[1]$ multiply the $(1,[1])$ entry at $y_0$ by $10$, and each of the other seven entries at $y_1$ by $0.05$. Doing this in the $S$ columns, with the entries before rounding, gives $84701.0317$. Doing it in the $R$ columns gives $237.6375$. The same operations at the four translates give
\begin{equation}\label{D:65}
 \begin{array}{c|rrr}
 u&P_{u,t,1}&M_{u,t,1}&P_{u,t,1}-M_{u,t,1}\\ \hline
 {[1]}&84701.0317&84463.3942&237.6375\\
 {[3]}&84163.1372&84463.3942&-300.2570\\
 {[5]}&84286.6628&84463.3942&-176.7314\\
 {[7]}&84692.5272&84463.3942&229.1329
 \end{array}
\end{equation}
\textbf{How to compute $C_i$ without fitting data.} For this table let $\xi$ be the character of the pairs modulo $16$ with $\xi([3])=\sqrt{-1}$. For $i=(v,s)$ put $\beta=2v-1$. As a function of a prime power, the coefficient of cell $i$ divided by $\Lambda$ has the following finite character expansion:
\begin{equation}\label{D:66}
 \frac18\sum_{j,l=0}^1\sum_{k=0}^3
 \beta^j(-1)^l\overline{\xi(s)^k}\,(\chi_{11}^j\chi_{1319}^l\xi^k)_{\rm prim}.
\end{equation}
The subscript uses the character of smallest modulus inducing the product on units, so a factor with exponent zero is the constant one even at its own prime. For units, the two sign sums are each $2$ when the signs agree and zero otherwise. The four powers of $\xi$ sum to $4$ precisely in column $s$. This gives coefficient $2$. At $1319$, terms with $l=1$ vanish, giving coefficient $1$ in the prescribed cell. At $2$ and $11$ the $\chi_{1319}$ factor sums to zero. Substituting the linear coefficient $-c_\chi$ for each nontrivial character and $-1-\gamma$ for the trivial one gives the displayed $C_i$. For general $t,n$ use the column $t^{-1}s$, the $n$ characters of $H$, and coefficient $1/(2n)$. [3]

\newpage
\subsection*{5.3 What translation preserves}

Fix $t,j$. First replace each distinguished weight by $b$. For a given column $s$, the equation $s=uz$, $z\in T_2$, gives exactly $|T_2|$ choices $u=sz^{-1}$. Thus for each admissible prime power $r$
\begin{equation}\label{D:67}
 \frac1n\sum_{u\in H}\mathbf1_{\{v(r)=1,\ t[r]\in uT_2\}}
   =\delta\,\mathbf1_{\{\chi_q(r)=1\}}
   =\frac{\delta}{2}(1+\chi_q(r)).
\end{equation}
Here $\mathbf1$ equals one when its condition holds and zero otherwise. Multiplication by $1-\chi_d(r)$ gives the right contribution at the three special primes as well, with zero understood at the two omitted primes. Hence, with $a=a_t$,
\begin{equation}
 \frac1n\sum_u\sum_{i\ {\rm distinguished\ at}\ u} b S_i(a/b)
           =\frac{\delta b}{2}S_{q,d}(a/b).                       \label{II:t5.4}
\end{equation}
This is an identity for every common argument $y=a/b$. The map sending a sum with expansion $\alpha y\log y+\beta y+o(y)$ to its remainder after subtracting those two terms is linear. The coefficients are unique by taking limits. Applying it to \eqref{II:t5.4} gives
\begin{equation}
 \frac1n\sum_u\sum_{i\ {\rm distinguished\ at}\ u} b R_i(a/b)
            =\frac{\delta b}{2}R_{q,d}(a/b).                      \label{II:t5.5}
\end{equation}
In particular, with $\eta=\chi_q\chi_d$, the right side involves $G+G_{\chi_q}-G_{\chi_d}-G_\eta$. The plus sign on $G$ comes from expanding $(1+\chi_q)(1-\chi_d)$.

In the example each of the four row-$1$ cells is distinguished once. From the $S_i(1000)$ column we get
\begin{equation}\label{D:70}
 \tfrac{10}{4}(352.749714+256.298804+404.588288+542.737992)
       \simeq3890.936997 .
\end{equation}
Direct summation with the two quadratic signs gives the same value $(10/8)S_{11,1319}(1000)$. For \eqref{II:t5.5}, the four remainders give $-2.703354\ldots$ to six decimals, or $(10/8)R_{11,1319}(1000)$. There is a separate check using the three $L$-values:
\begin{equation}\label{D:71}
 \begin{array}{c|rrrr|r}
 y&G&G_{\chi_{11}}&G_{\chi_{1319}}&G_{\chi_{11}\chi_{1319}}&R_{11,1319}\\ \hline
 1000&1.531&-2.299&-8.297&9.692&-2.163\\
 1000000&11.040&-266.971&912.895&-618.737&-550.088
 \end{array}
\end{equation}
To six decimals the three constants are $c_{\chi_{11}}=-0.084218$, $c_{\chi_{1319}}=-1.520228$, and $c_\eta=-0.779618$. They can be computed by grouping the Dirichlet series into residue classes and differentiating its Laurent expansion at $1$. The mean excess in section 5.2 is $-2.554507$ to six decimals. Its distinguished part accounts for $-2.703354$ and the other cells for $0.148847$, rounded separately. The latter need an estimate in the asymptotic proof. [3, 6]

\newpage
\subsection*{5.4 Discarding the other scale with a uniform estimate}

Divisibility uses all the cells, but only the distinguished cells should measure primes near $x$. Small weights at the other cells push their cutoffs far enough out for the prime number theorem to control them. The fixed-modulus limit $R_i(y)/y\to0$ gives no rate uniform in $q,d,N$. For the moduli of section 3 one has, for each primitive character in the cell expansion and for $y\ge \exp(x^{\varepsilon/3})/3$,
\begin{equation}
 |G_\chi(y)|\le y\,\vartheta,\qquad
 \vartheta=\exp(-\kappa x^{\varepsilon/8}),\qquad \kappa>0,       \label{II:t5.6}
\end{equation}
with $G$ for the trivial character. The constant $\kappa$ does not depend on $x$, $q$ or $d$, and the bound is for $x\ge x_0(\varepsilon,\kappa)$. This is Estimate C. It is proved in Part III, \S III.2 (Theorem~\ref{thm:C}), where $\kappa$ is called $c'$. At any nondistinguished cell $i$ use \eqref{II:t5.6} at $y=a/w_i$. Each remainder is a combination $R_i=\sum_\chi f_\chi G_\chi$ of the functions $G_\chi$ for the characters of the table, with $\sum_\chi|f_\chi|=2$: the sum of the absolute values of the $4n$ Fourier coefficients is $4n/(2n)=2$. Therefore
\begin{equation}\label{D:73}
 |w_i R_i(a/w_i)|\le2a\vartheta,\qquad
 \frac1{B n}\sum_{u,t,j}\sum_{i\ {\rm nondistinguished}}
        |w_i R_i(a_t/w_i)|\ll \frac{nx}{\delta}\,\vartheta=o(1).
\end{equation}
Indeed there are at most $2n$ cells per table, $a_t/b\le x+1$, and $nx/\delta$ grows polynomially in $x$. Exponential decay in \eqref{II:t5.6} absorbs this factor.

For a distinguished cell let $w_i=b+\Delta_i$ with $0\le\Delta_i\ll b/(nx^2)$. The function $w\mapsto(a/r-w)_+$ changes by at most $\Delta_i$ and only $r\le a/b$ can contribute. Since $0\le c(r)\le2\Lambda(r)$, Chebyshev's estimate bounds the total change in the prime sum by $O(x\Delta_i)$. In its main term $aC_i$ is unchanged, and $\mu a\log(w_i/b)\le\mu a\Delta_i/b$. Consequently
\begin{equation}\label{D:74}
 |w_i R_i(a/w_i)-b R_i(a/b)|\ll x\Delta_i,\qquad
 \frac{1}{\delta b/2}\sum_{i\ {\rm distinguished}} x\Delta_i\ll 1/x.
\end{equation}
For the second bound, there are $\delta n=|T_2|$ distinguished cells in each table. Summing over $u,t,j$ and dividing by $Bn$ gives $O(1/x)$.

Finally $|a_t/b-x|\le1$. When $y$ moves by at most $1$ from $x$, $S_{q,d}(y)$ and $y\log y$ change by $O(1+\log x)$, since $\sum_{r\le x+1}\Lambda(r)/r\ll\log x$. For $\chi=\chi_q,\chi_d,\eta$, Theorem~\ref{thm:pair}(d) of Part III gives $|c_\chi|\le4\log x$ (it also follows from Estimate B and the definition of $G_\chi$). The same estimate for the variation thus holds for the centering terms. Combining \eqref{II:t5.3}, \eqref{II:t5.5}, and the two error bounds gives
\begin{equation}
 \frac1n\sum_{u,t,j}(P_{u,t,j}-M_{u,t,j})
       =B\,R_{q,d}(x)+O_\varepsilon\big(B(1+\log x)\big).        \label{II:t5.7}
\end{equation}
With (M), this gives $R_{q,d}(x)\le\mathcal E/B+O_\varepsilon(1+\log x)$, proving (A). Only an upper bound follows. In the finite example the changes of distinguished weight and argument are both zero, so before estimation \eqref{II:t5.7} is the numerical split of the mean excess into $-2.703354$ and $0.148847$. [6, 11]

\newpage
\section*{6. Where the main inequality (M) comes from}\addcontentsline{toc}{section}{6. Where the main inequality (M) comes from}

Sections 6.1--6.2 construct the maps from theta series. Sections 6.3--6.4 explain why their derivatives are divisible by powers of inert primes. Sections 6.5--6.7 bound the size of the same derivatives and identify the main term. Sections 6.8--6.10 control denominators and show that enough independent images survive. We prove the elementary calculations in each step and say which arithmetic statements are quoted.

\subsection*{6.1 A theta lift in classical terms}

A theta series records representations of integers by a quadratic form. For instance put $L=\mathbb Z[i]$, $Q(a+bi)=a^2+b^2$ and $z$ in the complex upper half-plane. With $s=\exp(2\pi i z)$, its theta series is
\begin{equation}\label{D:76}
 \Theta_L(z)=\sum_{\lambda\in L}s^{Q(\lambda)}
     =1+4s+4s^2+4s^4+8s^5+\cdots .
\end{equation}
The coefficient of $s^5$ is $8$ because the solutions are $(\pm1,\pm2)$ and $(\pm2,\pm1)$. Absolute convergence follows from $\sum_{a,b}\exp(-2\pi\,\Im z\,(a^2+b^2))<\infty$. Fourier orthogonality recovers a coefficient by pairing: for each $y>0$,
\begin{equation}\label{D:77}
 e^{10\pi y}\int_0^1 \Theta_L(u+iy)e^{-2\pi i5u}\,du=8.
\end{equation}
Indeed $\int_0^1 e^{2\pi i(k-5)u}\,du$ is one for $k=5$ and zero for every other integer $k$.

To make this a lift, allow the lattice to vary within a finite collection of isometry classes $L_i$. Attach to each class a value $f(i)$. Weight it by $h_i=1/|\operatorname{Aut}L_i|$, the reciprocal of the number of its symmetries in the chosen group. The theta series are now a kernel depending on both $i$ and $z$. The lifted function is
\begin{equation}\label{D:78}
 \Theta(f)(z)=\sum_i h_i\overline{f(i)}\,\Theta_{L_i}(z).
\end{equation}
In the rank-one Hermitian example, unitary symmetries of $\mathbb Z[i]$ are multiplication by $1,-1,i,-i$. On its single class with $f=1$ this normalization gives $\Theta_L/4$, so the same Fourier pairing gives $8/4=2$. This example shows what integration against a theta kernel does. The high-rank construction has different groups and test functions.

For the source let $E=\mathbb Q(\zeta_N+\zeta_N^{-1},\sqrt{-q},\sqrt{-d})$, with $\zeta_N=\exp(2\pi i/N)$, and let $F$ denote its maximal real subfield. Write $\overline{\phantom{a}}$ for complex conjugation on $E$. A Hermitian matrix $J$ satisfies $\overline J^{\,t}=J$. Its unitary group consists of matrices $g$ with $\overline g^{\,t}Jg=J$. It is definite if its form at every real embedding of $F$ is definite. At a fixed integral level and scalar weight, forms on a definite group can be viewed as functions on finitely many lattice classes. At higher weight the values are vectors and must transform under the finite stabilizers.

For a pair of unitary groups, the general lift has the same two operations:
\begin{equation}\label{D:79}
 \Theta_\varphi(g,h)=\sum_{\lambda\in V(F)}(\omega(g,h)\varphi)(\lambda),
 \qquad
 \mathcal T_\varphi f(h)=\int_{[U(J)]}\Theta_\varphi(g,h)\overline{f(g)}\,dg.
\end{equation}
Here $V$ is a finite-dimensional coordinate space for the theta sum; $V(F)$ denotes its rational points over $F$. The notation $[U(J)]$ denotes integration over the group at all places modulo its rational points, with its invariant measure. The test function $\varphi$ decays rapidly at infinity and is locally constant of compact support at finite primes. The Weil representation $\omega$ is the action of the two groups on these test functions. Lattice indicators at the finite primes restrict the sum to lattices. In the displayed convention the lift is linear in $\varphi$ and conjugate-linear in $f$. [13] (The linearity statement of [13] is a statement in the archive, not used by the final theorem.)

\newpage
\subsection*{6.2 Composing maps and clearing their denominators}

For a finite form of the calculation let $I_j$ be the finite set of classes at stage $j$, $h_j(i)$ its weights, and $K_j(i,i')$ the matrix of the $j$th map. In the linear convention write
\begin{equation}\label{D:80}
 (T_j f)(i')=\sum_{i\in I_j} h_j(i)K_j(i,i')f(i).
\end{equation}
Complex conjugates from a pairing are absorbed by taking the conjugate source space. Expanding two such sums gives the composite kernel $\sum_{i'}K_j(i,i')h_{j+1}(i')K_{j+1}(i',i'')$. Thus successive maps really do use the same intermediate index.

In the source the definite ranks are $N_j=k+j(\tau_0+1)$ for an odd integer $\tau_0$. The rank conditions [26] include $m$ even and $m\ge2$, $k$ even and $k\ge m+3$, $\tau_0\ge k+m+4$, $\tau_0>3m$, and $m\le2n\le k$. Since $n\ge4$, the smallest case is $m=2$, $k=8$, $\tau_0=15$, with ranks $8,24,40$. The full chain also needs field, level, and lattice data. Since $m$ is even and $m\ge2$, the chain has at least two steps.

At the end of the chain the complex pairing of the entire chain can be expressed as a single period integral over a subgroup. The source proves this identity with prescribed normalization factors and coordinate changes. Its use is that the final period can be compared at different choices of the intermediate character twists. [7, 14]

\textbf{The integrality assertion for a finite chain.} Let $S$ be a subring of a commutative ring $R$. Suppose all $K_j(i,i')$ and starting values belong to $S$. If scalars $D_j\in R$ satisfy $D_j h_j(i)\in S$ for every $i$, then every final value multiplied by $\prod_{1\le j\le m}D_j$ belongs to $S$. For one map this follows by moving $D_1$ into its finite sum. Applying the same argument to the already cleared values proves the statement by induction. For vector values, replace $S$ by a lattice at each stage and require each $K_j(i,i')$ to carry the preceding lattice into the next. This is the assertion stated as $\mathrm{ChainInt}$, with its vector version, in [14] (a statement in the archive, not used by the final theorem).

For example, on one class let the successive weights be $1/2,1/3$, the kernel entries $5,7$, and the starting value $1$. The output is $35/6$, and clearing the two weights gives $35$. At the prime $5$, two entries divisible respectively by $5^2$ and $5^3$ would give a cleared output divisible by $5^5$. In a sum the exponent can increase through cancellation. It is always at least the common lower bound on the summands. This elementary induction explains how divisibility persists. It requires separate local estimates to put significant prime powers into each kernel entry.

\textbf{Where a period enters.} A kernel evaluated in complex coordinates need not be integral. An algebraic local coordinate and a complex analytic coordinate on the same elliptic curve differ at the tangent space by a nonzero scale. An $a$th derivative acquires its $a$th power. Taking logarithms gives a term linear in $a$. In the global comparison the local gains, the period, and the lattice indices must all pertain to the same normalized maps. We next give a local example in which one can track every derivative, and then the estimates that take its place in the chain.

\newpage
\subsection*{6.3 A genuine local example and an elementary minor calculation}

Let $s(z)$ be the inverse near $0$ of $t\mapsto\int_0^t(1-u^4)^{-1/2}\,du$, and let $A_j=s^{(j)}(0)$. The lemniscate curve $y^2=1-x^4$ has the symmetry $(x,y)\mapsto(ix,y)$ fixing the origin $(0,1)$. This yields complex multiplication by $\mathbb Z[i]$. Differentiating the inverse function gives
\begin{equation}\label{D:81}
 s(0)=0,\quad s'(0)=1,\quad (s')^2=1-s^4,\quad s''=-2s^3.
\end{equation}
The product rule gives an integer recursion, with $A_0=0,A_1=1$,
\begin{equation}\label{D:82}
 A_{j+2}=-2\!\sum_{a+b+c=j}\frac{j!}{a!\,b!\,c!}\,A_a A_b A_c.
\end{equation}
For $j=3$, only $(1,1,1)$ contributes, giving $A_5=-12$. For $j=7$ there are three permutations of $(5,1,1)$, giving $A_9=-2\cdot3\cdot42\cdot(-12)=3024$. Uniqueness of the series solution gives $s(iz)=is(z)$, so only indices $j\equiv1\bmod4$ can contribute. The exact computation gives
\begin{equation}\label{D:83}
 A_{17}=21224560896=2^8\,3^6\,7^2\,11\,211 .
\end{equation}
The primes $3,7,11$ remain prime in $\mathbb Z[i]$. Their exponents are $6,2,1$; the model in section 2 predicts the respective leading terms $17\cdot3/8$, $17\cdot7/48$, $17\cdot11/120$. At the split primes $5,13$ the exponents here are zero. At $j=201$ the computed valuations at $3,7$ are $75,29$. These valuations are computed, not proved in general; section 2 proves the bound only for its model series.

There is a quantitative gap between this example and \eqref{II:t4.1}. At $j=101$ the contribution to $\log|A_j|$ from inert primes at most $j$ is $167.67$ out of $306.77$. At $j=121$ it is $203.58$ out of $388.67$. The density calculation explains the scale: $\sum_{p\le j,\ p\equiv3\bmod4}(\log p)/p\sim\tfrac12\log j$, so the model exponents $j/p$ at these primes give about $\tfrac12 j\log j$. The main inequality (M) requires agreement even to the term proportional to $j$. So the numbers $A_j$ cannot serve in (M): they are too large by a factor $2$ in the leading term (defect (D3) of \S I.3).

\textbf{A local calculation from the source.} Let $A$ be a commutative ring with an element $\pi_0$ satisfying $\pi_0^2=p$, where $p$ is prime. For an integer $e\ge0$ put $P=p^e$. Let $M,N'$ be square matrices with $N'_{bc}\equiv M_{bc}^P\bmod p$. Suppose all but the zeroth coordinates of row vectors $u,u'$ are divisible by $\pi_0$. Write $X=uM$, $Y=u'N'$. The minor used in [15] satisfies
\begin{equation}\label{D:84}
 \pi_0\mid Y_c X_0^P-Y_0 X_c^P .
\end{equation}
Indeed, modulo $\pi_0$, $X_c=u_0M_{0c}$ and $Y_c=u'_0 M_{0c}^P$. Substitution makes the two products equal. For $p=3$, $\pi_0=\sqrt3$, take $P=3$, $M=N'=\left(\begin{smallmatrix}1&1\\0&1\end{smallmatrix}\right)$, $u=(1,\pi_0)$, $u'=(1,0)$. The minor for $c=1$ is $1-(1+\pi_0)^3=-3\pi_0-3\pi_0^2-\pi_0^3$, visibly divisible by $\pi_0$. The estimate on the full lift requires many such congruences and their behavior under differentiation.

\newpage
\subsection*{6.4 How the weight cuts down the local exponent}

Here is a calculation which accounts for the positive part in \eqref{II:t4.1}. A directional Hasse derivative of a polynomial $Q$ is the coefficient of $Z^\beta$ in $Q(X+\sum_j b_j Z_j)$, where $b_j$ are vectors, $Z_j$ indeterminates, and $\beta$ a list of nonnegative exponents. It has no factorial denominator. For a product $\prod_c Q_c$ the coefficient is the sum of products over decompositions $\beta=\sum_c\beta_c$. If for a given set $S$ of directions the coefficient for $Q_c$ vanishes whenever $\sum_{j\in S}(\beta_c)_j>m_c$, then the coefficient for the product vanishes whenever $\sum_{j\in S}\beta_j>\sum_c m_c$: each term exceeds the allowance of some factor. This is the product-vanishing lemma in [21].

For a determinant the allowance has a simple origin. In $\det(X+s b,Y+t b)$, where $X,Y,b$ are columns of length two, the coefficient of $st$ is $\det(b,b)=0$, and there is no term in $s^2$ or $t^2$. Each minor can take at most one such hit. Products of $C$ minors can take at most $C$. With several directions confined to specified columns, the number of available columns provides another upper bound. These are constraints on nonvanishing terms before any prime estimate is made.

For an inert prime $p$, write $\gamma_p=p/(p^2-1)=\sum_{h\ge1,\ h\ {\rm odd}}p^{-h}$. Suppose a local expansion of derivative order $a$ supplies a valuation lower bound $\gamma_p a-\sum_h i_h-D_0$, where $D_0$ is a loss, and the nonvanishing terms have
\begin{equation}\label{D:85}
 0\le i_h\le \min(a/p^h,w_h)\qquad (h\ {\rm odd}).
\end{equation}
Then the gain after charging these terms is at least
\begin{equation}
 \sum_{h\ {\rm odd}}\left(a/p^h-\min(a/p^h,w_h)\right)-D_0
     =\sum_{h\ {\rm odd}}(a/p^h-w_h)_+-D_0.                      \label{II:t6.1}
\end{equation}
The first bound $p^hi_h\le a$ is a degree bound. The second is supplied by column and minor capacities in the translated weight data. For example $p=3$, $a=81$, $w_h=2$ gives $25$ at $h=1$, $1$ at $h=3$, and zero thereafter. Obtaining all these bounds together is the arithmetic content. In the source the factorial valuations in the polynomial substitution give the initial lower bound, other terms have to be paid for, and a local lattice index must also be subtracted. [19, 21]

\textbf{The factorial count, with proof.} The initial lower bound in the last paragraph comes from valuations of factorials. Here is the source's statement in full, with an elementary proof. Let $p$ be prime and put $\rho_p=1/(p^2-1)$, so that $\gamma_p+\rho_p=1/(p-1)$. For an integer $n\ge0$ let $\lambda(n)$ be the number of digits of $n$ in base $p$, with $\lambda(0)=0$. Let $S$ be a finite set and $S_0\subset S$. Suppose given nonnegative integers as follows:
\begin{enumerate}[label=(\roman*),leftmargin=2.6em,itemsep=0pt,topsep=2pt]
\item $\alpha_s$ and $\ell_s$ for $s\in S$, with $\alpha_s=0$ outside $S_0$;
\item lists $e_{s,1},\ldots,e_{s,r_1}$ with sum $\alpha_s$;
\item lists $a_{s,1},\ldots,a_{s,r_0r_1}$ with sum at most $\ell_s$;
\item $i_{s,h}$ for $h$ in a finite set of nonnegative integers, zero outside $S_0$, with $p^{2h+1}i_{s,h}\le\alpha_s$;
\item with $A=\sum_s\alpha_s$ and $\Lambda=\sum_s\ell_s$, one has $\sum_{s}\sum_j a_{s,j}+\sum_{s,h}i_{s,h}=\Lambda$.
\end{enumerate}
Let $V$ be the sum of $v_p(n!)$ over all the $e_{s,j}$, $a_{s,j}$, $i_{s,h}$. Then for every integer $i_2$ with $pi_2\le\Lambda$,
\begin{equation}
 V\ \ge\ \gamma_pA+\rho_p(A+\Lambda)+i_2-\Delta,\qquad
 \Delta=\sum_{s\in S}\big(r_1\lambda(\alpha_s)+r_0r_1\lambda(\ell_s)\big)
        +\sum_{s\in S_0}\lambda(\alpha_s)^2 .                    \label{II:t6.1a}
\end{equation}
\emph{Proof.} Legendre's formula is $v_p(n!)=(n-\sigma_p(n))/(p-1)$, where $\sigma_p(n)$ is the sum of the digits of $n$ in base $p$. Each digit is at most $p-1$, so $v_p(n!)\ge n/(p-1)-\lambda(n)$. The integers in the three families add up to $A+\Lambda$, by the two hypotheses on sums. Hence $V\ge(A+\Lambda)/(p-1)-\sum\lambda(n)$. Now count digits. Each $e_{s,j}$ is at most $\alpha_s$, and there are $r_1$ of them: at most $r_1\lambda(\alpha_s)$. Each $a_{s,j}$ is at most $\ell_s$, and there are $r_0r_1$ of them: at most $r_0r_1\lambda(\ell_s)$. If $i_{s,h}\ge1$ then $p^{2h+1}\le\alpha_s$, so $h<\lambda(\alpha_s)$ and at most $\lambda(\alpha_s)$ values of $h$ occur; each such $i_{s,h}$ is at most $\alpha_s$. These contribute at most $\lambda(\alpha_s)^2$, and only for $s\in S_0$. Thus $\sum\lambda(n)\le\Delta$. Finally $(A+\Lambda)/(p-1)=\gamma_pA+\gamma_p\Lambda+\rho_p(A+\Lambda)$, and $\gamma_p\Lambda\ge\Lambda/p\ge i_2$. $\square$

In \eqref{II:t6.1a} the term $\gamma_pA$ is the model gain $A(1/p+1/p^3+\cdots)$ of section 2. The term $\rho_p(A+\Lambda)$ pays exactly for the denominators $1/(p^2-1)$ per unit of degree met there in the inverse series, and $i_2$ pays for one further division. The loss $\Delta$ is small: per index $s$ it is at most $(r_1+r_0r_1+\lambda)\lambda$, where $\lambda$ is the larger of the digit counts $\lambda(\alpha_s)$, $\lambda(\ell_s)$. For example, with $p=3$, one index $s$, $r_0=r_1=1$, $\alpha=81$, $\ell=0$ and no $i_{s,h}$: $V=v_3(81!)=40$, while $\gamma_3A+\rho_3A=81/2=40.5$ and $\Delta=\lambda(81)+\lambda(81)^2=30$. The inequality $40\ge40.5-30$ holds with room to spare; the main terms agree to within one. This lemma supplies valuations of individual factorials. It does not identify which factorials occur in the theta coefficients, which is part of the quoted local theorem. [21]

\textbf{The coefficient $2\log p$.} For the coefficient field $L_0$ and a prime ideal $\mathfrak p$ over $p$, write $e_{\mathfrak p},f_{\mathfrak p}$ for the ramification index and residue degree. The normalized local weight is $\lambda_{\mathfrak p}=\log \operatorname{Norm}(\mathfrak p)/[L_0:\mathbb Q]$. The source measures a local gain $y$ in units $2\lambda_{\mathfrak p}e_{\mathfrak p}y$. Since $\operatorname{Norm}(\mathfrak p)=p^{f_{\mathfrak p}}$ and $\sum_{\mathfrak p\mid p}e_{\mathfrak p}f_{\mathfrak p}=[L_0:\mathbb Q]$, one unit at every place over $p$ sums to $2\log p$. For example in $\mathbb Q(i)$ above $3$, $e=1$, $f=2$ and $\lambda=(\log9)/2=\log3$. The factor $2$ is the squared-norm normalisation of \S I.4.2(ii). At the ramified prime $d$ the splitting condition has half weight and gives $\log d$. These are the coefficients $c(r)$ of section 4.1; section 5.2 checked that the character expansion reproduces them. [3, 22]

The concrete minor, the Hasse product rule and \eqref{II:t6.1} are three elementary ingredients of the local argument. Applying them to all the normalized theta kernels, choosing the usable places, and bounding the losses after summation requires the source's local divisibility theorem. In its statement the left side is the sum of the odd-prime-power expressions \textbf{minus} the orbit and rounding losses. The right side is the sum of actual local gains over the $m$ maps \textbf{minus} the local lattice correction. [19] The divisibility at inert primes is written up in Part III, \S III.5, conditionally on the integrality of the theta kernels in the coordinates of the height-two logarithm (hypothesis (B)), which is taken from the Lean development; under that hypothesis \S III.5 proves the lower bound for the valuation of a product of minors applied to one such function (for Frobenius orbits of length two; for longer orbits it also uses a duality statement cited from the Lean), with the positive parts of \eqref{II:t6.1}, and it proves without hypothesis the estimate of section 2 for the inverse of the logarithm in several variables; it does not prove the passage to lattices, and the congruence of section 6.3 is not used there. The passage to lattices is treated in Part III, \S III.6: the lemmas on lattices and the splitting of the prime sum are proved there, the theorem on the glued lattice is conditional on two statements of the source, and the bounds for the loss sums are quoted.

\newpage
\subsection*{6.5 Which theta function is differentiated?}

At step $j$ let $N_j=k+j(\tau_0+1)$ be the dimension of the positive definite Hermitian space at that stage (the space $\mathsf X_j$ of Part~III, \S\ref{ob:sec}), where $\tau_0$ is the odd rank parameter. The coordinates of the holomorphic theta function, for a pair of lattice classes $L,L'$ at consecutive stages and a complex embedding of the coefficient field, are the entries of a rectangular matrix at each weight place:
\begin{equation}\label{D:88}
 z=(z_{\tau,\alpha,\beta}),\qquad
 1\le\alpha\le N_j,\quad 1\le\beta\le N_{j+1}.
\end{equation}
Write its Taylor series at zero as $\Theta_{L,L'}(z)=\sum_{\nu}c_\nu z^\nu$, using a multi-index $\nu=(\nu_g)$ with $g=(\tau,\alpha,\beta)$ and $\nu!=\prod_g\nu_g!$. For vectors in the consecutive spaces of forms there is a polynomial $P(z)=\sum_\nu P_\nu z^\nu$ determined by their weights and pairing. In the formulation of [23] (a statement in the archive; see the start of Part III) the local expression in a matrix entry is
\begin{equation}
 P(\partial)\Theta_{L,L'}(0)=\sum_{\nu} P_\nu\,\nu!\,c_\nu .       \label{II:t6.2}
\end{equation}
This explains why the polynomial minors and the factorials occur together. At each place the monomials in $P$ have prescribed total degree. The global matrix coefficient sums the displayed expressions over preceding classes with reciprocal stabilizer weights. The source also describes the theta function with algebraic coordinates at finite places so that the same normalized pairing has the local divisibility. That identification of rational structures is a theorem of the construction. That the entries used by the final theorem are of the form \eqref{II:t6.2} is quoted. [23]

On the complex side the polynomial has the squared Gaussian norm
\begin{equation}\label{D:90}
 \|P\|_{\rm F}^2=\sum_\nu |P_\nu|^2\nu! .
\end{equation}
Here is a proof of its simplest case. In the complex plane give area the density $\pi^{-1}e^{-|z|^2}$. Integration in angle makes different powers orthogonal, and the squared norm of $z^a$ is
\begin{equation}\label{D:91}
 2\int_0^\infty r^{2a+1}e^{-r^2}\,dr
     =\int_0^\infty t^a e^{-t}\,dt=a! .
\end{equation}
The last integral follows by $a$ integrations by parts. For several independent variables the integrals multiply. For a positive width $\Omega$, the density $(\pi\Omega)^{-1}\exp(-|z|^2/\Omega)$ gives the squared monomial norm $a!\Omega^a$. In the dual coefficient estimate the scale appears with the inverse power. Weighted Cauchy--Schwarz gives the inequality stated in the archive [23]
\begin{equation}\label{D:92}
 |P(\partial)\Theta(0)|^2\le
    \|P\|_{\rm F}^2\sum_{\nu\in\operatorname{supp}P}\nu!|c_\nu|^2 ,
\end{equation}
where $\operatorname{supp}P$ consists of its monomials with nonzero coefficients. Apply Cauchy--Schwarz to the two lists $P_\nu\sqrt{\nu!}$ and $c_\nu\sqrt{\nu!}$.

The analytic estimates bound the coefficient factor using the Gaussian-weighted norm of the theta function on a fundamental domain. The polynomial factor is bounded by a product of hook factors of the weights times the squared norms of the two input vectors. In one row of length $a$ the hook factor is exactly $a!$. The full operator bound also contains the period, level, multiplicities and theta mass. [23, 24] The estimates for the two factors are proved in Part III, \S III.4, apart from the four conditional passages listed at the start of that section.

\textbf{Comparison with section 6.3.} In the lattice comparison the archimedean map height is the average over infinite places $v$ of $\log\|T_v\|_{\rm op}^2$, where the operator norm uses the two chosen Hermitian metrics. Divisibility by $p^e$ thus contributes $2e\log p$ when measured at all places above $p$, and a degree-$a$ monomial contributes $\log(a!)=a\log a-a+O(\log(a+1))$ in the squared Gaussian norm. The difference from section 6.3 is visible in $A_a=a![z^a]s(z)$, where $[z^a]$ means the coefficient of $z^a$. The nearest pole of $s$ is at finite positive distance $R$; the radius formula gives $\limsup_a|[z^a]s|^{1/a}=1/R$. Along a subsequence, therefore, $\log|A_a|^2=2\log(a!)+O(a)$, while the squared Gaussian norm contributes just one factorial. [24, 25]

\newpage
\textbf{A natural question.} Why not differentiate a classical entire theta function, which already has Gaussian size? For the lemniscatic lattice the Weierstrass function $\sigma$ of $y^2=4x^3-4x$ has integer derivatives $\sigma^{(n)}(0)$; they vanish unless $n\equiv1\bmod4$, and the nonzero ones satisfy $\log|\sigma^{(n)}(0)|\le\tfrac12n\log n+O(n)$, by the Gaussian growth of $\sigma$ and Cauchy's formula. Exact computation for $n\le241$ shows that their divisibility at inert primes falls well short of the model: the exponent of $3$ is $13$ at $n=101$ and $28$ at $n=241$, against $3n/8=37.9$ and $90.4$. So Gaussian size does not by itself bring the divisibility with it. The source's local theorem establishes both properties together for its theta kernels, in coordinates chosen for that purpose, and it is quoted here. The computation is in the file \texttt{e4\_sigma.py} listed under Computations. [21, 23]

\subsection*{6.6 Size of a derivative, its period, and the count of forms}

The model for the size of a derivative is this. Let $f$ be an entire function with $|f(z)|\le S\,e^{|z|^2/2}$ for all $z$, where $S>0$. Cauchy's formula on the circle $|z|=\sqrt a$ gives $|f^{(a)}(0)|\le S\,a!\,(e/a)^{a/2}$ for $a\ge1$. By Stirling's formula, if $f^{(a)}(0)\ne0$, then $\log|f^{(a)}(0)|^2\le a\log a-a+2\log S+O(\log(a+1))$: one factorial for the square. Rescaling the variable by the period adds $-a\log\Omega_E$, with $\Omega_E$ as defined below. This is the model at the start of Part III, \S III.4. It identifies the factorial and scale terms of \eqref{II:t6.5}. The bound for the norm of an actual composed map needs additional estimates.

Here is the exact period normalization appearing in the comparison. It is the normalization for the field $E$, which is the one of the source. Let $\mathrm{rd}_E=|\mathrm{disc}\,E|^{1/[E:\mathbb Q]}$ be the root discriminant of $E$, and put $v=\mathrm{rd}_E/2$. The normalized period $\Omega_E$ is the positive number with
\begin{equation}
 \log\Omega_E=\gamma+\log\pi-\log v-c_{\chi_d}
            =\gamma+\log(2\pi/\mathrm{rd}_E)-c_{\chi_d}.                 \label{II:t6.3}
\end{equation}
This is the Chowla--Selberg identity in the form used. The period $\Omega_E$ measures squared size. In [16] the period is not constructed from a specific curve: the statement is that a curve and a normalisation exist for which \eqref{II:t6.3} holds. It is important here which period normalization is used. Rescaling the differential changes the linear term. [16]

\textbf{Remark. For the elliptic curve itself}, with complex multiplication by $\mathbb Q(\sqrt{-d})$, the same identity holds with $v=\sqrt d/2$, the imaginary part of the quadratic lattice parameter $(1+i\sqrt d)/2$, and with the period $\Omega$ of the curve: $\log\Omega=\gamma+\log(2\pi/\sqrt d)-c_{\chi_d}=\gamma+\log\pi-\log v-c_{\chi_d}$. In this form the identity is Kronecker's limit formula, and $\Omega$ is a geometric mean over the ideal classes of $\mathbb Q(\sqrt{-d})$. For $\mathbb Q(i)$, with $d$ replaced by $4$, exponentiating the same expression gives $\Omega=2\pi(\varpi/\pi)^2$, with $\varpi$ the lemniscate constant. For $d=1319$ the formula gives $\log\Omega=\gamma+\log(2\pi/\sqrt{1319})+1.520227887516\ldots$. The two normalizations differ by $\log\Omega-\log\Omega_E=\log(\mathrm{rd}_E/\sqrt d)$, and the sum of the logarithm of the period and $\log v$ is the same in both: $\log\Omega_E+\log(\mathrm{rd}_E/2)=\log\Omega+\log(\sqrt d/2)$. Only this sum enters \eqref{II:t6.6}. So \eqref{II:t6.6} is the same in either normalization, and the numerical check in section 6.7 may use either. The two lines of \eqref{II:t6.5}, taken separately, hold only with $\Omega_E$ and $v=\mathrm{rd}_E/2$.

The other source of linear terms is a count of forms. The mass of a definite genus is the sum of reciprocal stabilizer orders over its lattice classes, as in section 6.1. The dimension of a representation with strictly decreasing shifted integral weights $z_1,\ldots,z_r$ is
\begin{equation}\label{D:94}
 D_r(z)=\frac{V_r(z)}{\prod_{j=0}^{r-1}j!},\qquad
 V_r(z)=\prod_{1\le i<j\le r}(z_i-z_j).
\end{equation}
The scale factor in the normalized mass changes from rank $r-1$ to $r$ according to the mass formula. Write $f=2n$ for the number of weight places. In the algebraic identity stated in [17] (the statement of [17] is not used by the final theorem; a statement of the same name elsewhere in the archive is), $M_r$ denotes the mass and $\kappa_r>0$ collects the extra local factors:
\begin{equation}\label{D:95}
 \frac{M_r}{M_{r-1}}=
 \big((r-1)!(v/\pi)^r\big)^f(2v)^{-f/2}\kappa_r .
\end{equation}
Here $v=\mathrm{rd}_E/2$, as in \eqref{II:t6.3}, and $M_r$ is unrelated to the subtracted main term $M_{u,t,j}$. For two ranks $r_0<r_1$ (for instance two consecutive ranks $N_j$, $N_{j+1}$ of the chain of section 6.2), let $u_\rho,w_\rho$ be their decreasing weight lists at each of the $f$ places $\rho$. Multiplying the mass ratios and the $f$ dimension ratios gives precisely
\begin{equation}
 \frac{M_{r_1}\prod_\rho D_{r_1}(w_\rho)}{M_{r_0}\prod_\rho D_{r_0}(u_\rho)}
 =\prod_\rho\frac{V_{r_1}(w_\rho)}{V_{r_0}(u_\rho)}
   (v/\pi)^{f\sum_{r=r_0+1}^{r_1} r}(2v)^{-(r_1-r_0)f/2}
   \prod_{r=r_0+1}^{r_1}\kappa_r .                                   \label{II:t6.4}
\end{equation}
Indeed the extra $(r-1)!$ in the mass cancels the extra $(r-1)!$ in each dimension denominator. Taking logarithms is the count identity on ordered weights. For a numerical algebra check take one place, $r_0=2$, $r_1=3$, and $u=(3,1),w=(5,3,0)$. The two dimensions are $2$ and $15$. The ratio of dimensions is $15/2$. The factorial from the mass is $2$, leaving $15=V_3(w)/V_2(u)=30/2$ before the scale and local factors. Equation \eqref{II:t6.4} follows directly from the displayed mass recurrence. The automorphic mass formula supplying that recurrence is quoted. The passage from the count identity \eqref{II:t6.4} to the second line of \eqref{II:t6.5} below (the definition of the level $\mathcal T$, the factors $\kappa_r$, the error $E_{\rm cnt}$) is also quoted.

\newpage
\subsection*{6.7 Why the linear term has precisely the required sign}

Write $\mathcal T$ for the positive auxiliary level attached to one choice of lattice data. It is different from the set $T$ of derivative labels. Let $\ell_g=\log(g_0g_1)$ for the two auxiliary level primes $g_0,g_1$. The size and count bounds [7, 19] have the following form at a fixed translate, with $\Omega_E$ and $v=\mathrm{rd}_E/2$ as in \eqref{II:t6.3}. Here $h_B$ stands for the combination $\mu_{\rm CC}+h-F_0$ of three quantities defined after the display, and $E_{\rm size},E_{\rm cnt}$ are the respective error allowances:
\begin{equation}
 \begin{aligned}
 h_B &\le \sum_{t\in T}\sum_{j=1}^m a_t(\log a_t-1-\log\Omega_E)
               -\sum_{t\in T}a_t\log\mathcal T+E_{\rm size},\\
 \log\mathcal T &\ge
  \sum_{j=1}^m\big(\langle\log w_{u,j}\rangle+\log v-\log\pi\big)
             +m\ell_g-E_{\rm cnt}.
 \end{aligned}                                                   \label{II:t6.5}
\end{equation}
The three quantities in $h_B=\mu_{\rm CC}+h-F_0$ are these. First, $h$ is an upper bound for the height of the pairing map, the average over the archimedean places of the logarithm of the squared operator norm (section 6.5). Second, $\mu_{\rm CC}$ is the mean slope of the common lattice in which all the pairing maps take their values. Third, $F_0$ is the part of the local divisibility that comes from the places other than those of the top weights.

Neither offset survives. The term $F_0$ also appears, with the opposite sign, in the slope of the source lattice after gluing (section 6.8), and the two cancel. The term $\mu_{\rm CC}$ is the slope of the target in the inequality of section 6.10, where it appears on both sides. So one may read \eqref{II:t6.5} as a bound for the height $h$ alone, as in Part I. An earlier formulation [18] writes the same pair of inequalities for a single quantity.

The local divisibility comparison bounds the prime sum by $h_B+m\ell_g\sum_t a_t$ plus losses; the second term is a level correction. Insert the \textit{lower} bound for $\log\mathcal T$ in the first line of \eqref{II:t6.5}. Since each $a_t$ is positive, this gives an upper bound. The terms $m\ell_g\sum_t a_t$ cancel, and the coefficient of $a_t$ at step $j$ is
\begin{equation}
 \begin{aligned}
 &\log a_t-1-\log\Omega_E-\langle\log w_{u,j}\rangle-\log v+\log\pi\\
 &\hspace{10mm}=\log a_t-1-\gamma+c_{\chi_d}-\langle\log w_{u,j}\rangle .
 \end{aligned}                                                    \label{II:t6.6}
\end{equation}
The equality is \eqref{II:t6.3}. Thus the sum of the surviving terms is exactly $\sum_{t,j}M_{u,t,j}$. The two errors in this substitution are $E_{\rm size}+E_{\rm cnt}\sum_t a_t$. Other local and lattice losses must still be added to them. The main term $M$ needs all three: Stirling's formula ($\log a-1$), the period \eqref{II:t6.3} ($-\gamma+c_{\chi_d}$), and the count ($-\langle\log w\rangle$).

In the finite table one can check the cancellation even though no level or lattice has been constructed for that table. For $a=10000$ take the single-step expression on the first line of \eqref{II:t6.6}, multiply by $a$, and use the normalization of the elliptic curve (the remark in section 6.6): put $v=\sqrt{1319}/2$. The period combination $-\log\Omega-\log v+\log\pi$, which is the same as $-\log\Omega_E-\log(\mathrm{rd}_E/2)+\log\pi$, equals $-\gamma-1.520227887516\ldots$. The remaining two entries are $\log10000-1$ and $-(\log10+7\log0.05)/8$. This is the value $84463.3942$ of $M$ in \S4.3, which the prime sums of \S5.2 reproduce as $P-(P-M)$. An uncancelled constant per derivative in \eqref{II:t6.6} would survive multiplication by $a_t\simeq bx$; division by $B$ would leave an error at least of order $x$.

\textbf{What is quoted in this comparison.} The period formula \eqref{II:t6.3}, the mass formula, the construction and count of the auxiliary levels, and the local level correction are the arithmetic inputs; they are quoted. The analytic norm bound is proved in part in Part III, \S III.4. The factorial cancellation \eqref{II:t6.4} and all signs in \eqref{II:t6.6} are algebraic consequences proved above. The next comparison explains why the divisibility can be applied to enough independent vectors at once.

\newpage
\subsection*{6.8 From one determinant to lattices}

A lattice $L$ of rank $r$ in a Euclidean space is the integer span of $r$ independent vectors. Its covolume, measured in its span, is the volume of the parallelepiped on a basis. Put $\mu(L)=-\log(\operatorname{covol}L)/r$. Let $T$ be an injective linear map with $T(L)\subset M$, and let $\|T\|$ be its Euclidean operator norm. (This section uses the plain normalisation, $\mu$ and $\log\|T\|$. In \eqref{II:t6.5}--\eqref{II:t6.6} the height $h$ is the logarithm of the \emph{squared} norm, and $\mu_{\rm CC}$ is normalised to match, so it is twice a slope in the present sense, up to an additive constant depending only on the field. The two normalisations must not be mixed.) The volume of $T(L)$ is at most $\|T\|^r$ times that of $L$, as one sees using an orthonormal basis of the span and Hadamard's inequality. Taking logarithms,
\begin{equation}\label{D:99}
 \mu(L)\le \mu(T(L))+\log\|T\| .
\end{equation}
Suppose first that $M$ has the same rank $r$. Then $T(L)$ has finite index in $M$ and $\operatorname{covol}T(L)=[M:T(L)]\operatorname{covol}M$, so
\begin{equation}\label{D:100}
 \mu(L)+\frac1r\log[M:T(L)]\le \mu(M)+\log\|T\| .
\end{equation}
This is the exact form of the principle that divisibility is bounded by size: the index is on the left and the norm on the right. For $T=\operatorname{diag}(6,10)$ and $L=M=\mathbb Z^2$ the index is $60$, both slopes are zero and $\|T\|=10$. The inequality reads $\tfrac12\log60\le\log10$, that is, $60\le100$. The factors $2,3,5$ of the index and the two column lengths refer to the same map.

If $L$ has smaller rank than $M$, the image $T(L)$ is only a sublattice of smaller rank, and the first inequality gives
\begin{equation}\label{D:101}
 \mu(L)\le \mu_{\max}(M)+\log\|T\|,
\end{equation}
where $\mu_{\max}$ is the supremum of slopes of nonzero sublattices. A single map of small rank therefore controls only the maximal slope of its target. Sections 6.9 and 6.10 use many maps together in order to reach the mean slope $\mu(M)$.

A rational matrix presents a further problem: local divisibility may apply at different primes in different bases. The source uses equality of a period pairing under changes of intermediate twists to place the local lattices in a common rational vector space. One can then choose a favorable lattice at each prime and put these choices together.

\begin{illbox}
Over $\mathbb Q$ the last operation has a short proof. Let $L_p$ be a full lattice over $\mathbb Z_p$ inside $\mathbb Q_p^r$, equal to $\mathbb Z_p^r$ outside a finite set. Here $\mathbb Q_p$ is the completion of $\mathbb Q$ for the $p$-adic absolute value, and $\mathbb Z_p$ its ring of integers. Choose a positive integer $D$ with $D\mathbb Z_p^r\subseteq L_p\subseteq D^{-1}\mathbb Z_p^r$ at all primes. Let $L$ be the set of $x$ in $D^{-1}\mathbb Z^r$ whose image in $D^{-1}\mathbb Z_p^r/D\mathbb Z_p^r$ lies in $L_p/D\mathbb Z_p^r$ for every $p\mid D$. The Chinese remainder theorem applied after multiplying by $D$ identifies the ambient finite quotient with $(\mathbb Z/D^2\mathbb Z)^r$, the product of its prime parts. The resulting lattice has the prescribed localization at every prime. For example allow a denominator $2$ in the first coordinate only at $2$, and a denominator $3$ in the second only at $3$. The result is $\tfrac12\mathbb Z\oplus\tfrac13\mathbb Z$. Its covolume is $1/6$. This proves the gluing step over $\mathbb Q$. The source uses the corresponding construction over a number field, with its locally chosen character twists.
\end{illbox}

\textbf{The denominators that remain.} For a prime $\ell$ and positive twist index $a$, the relevant finite group consists of first Galois cohomology classes over $E$ with values in $\mathbb Q_\ell/\mathbb Z_\ell(-a)$, unramified away from $\ell$ and zero locally above $\ell$. The twist means that Galois acts by the inverse $a$th power of its action on roots of unity of $\ell$-power order. A finite group of order $\ell^h$ costs $h\log\ell$ in a logarithmic index. Write $h_{\ell,a}$ for the length of this group as a $\mathbb Z_\ell$-module. On the selected window the supplied bound, after passage to the working scale, is $h_{\ell,a}<\infty$ for $1\le a<\lfloor x^{400}\rfloor$ (the number of twists used is at most $x^{400}$; see Part III, \S\ref{gc:sec}) and
\begin{equation}\label{II:shabound}
 \sum_{\ell\in S}(\log\ell)\sum_{1\le a<\lfloor x^{400}\rfloor}h_{\ell,a}\le\sqrt b
\end{equation}
for every finite set $S$ of primes. This is the uniform assertion needed here. (In [8] the bound is stated for $a<x^{401}$ at a slightly different scale; the form above is what is used.) Even for $\zeta$ the bound is needed over the field $E$, not over $\mathbb Q$. The statement on these extension groups, the comparison inequality relating those lengths to the numerical denominator cost, and the independence of the normalized global period are quoted. [8, 19] The slope inequalities proved in this section are the elementary case, for lattices over $\mathbb Z$. A slope inequality over the integers of a number field, with the height of a map on the right, is proved in Part III, \S\ref{as:sec}, Lemma~M3, for an injective map and for the maps induced on quotients; this is the form that the deduction of (M) there uses. The form with minimal slopes, which this section marked as quoted and which the source states, follows from it (apply Lemma~M3 to the injective map induced on the quotient by the kernel: $\mu_{\min}(\Lambda)\le\mu(\Lambda/(\Lambda\cap\ker\phi))\le\mu_{\max}(\Lambda')+\mathrm{ht}(\phi)$), and is not needed there (section 6.10).

\newpage
\subsection*{6.9 Why many independent images improve the estimate}

A large norm for one map could be concentrated on one direction. A second moment distinguishes this possibility. For a positive Hermitian matrix $Q$ of size $D$, Cauchy--Schwarz on its nonzero eigenvalues gives
\begin{equation}\label{D:103}
 \operatorname{rank}Q\ge (\operatorname{tr}Q)^2/\operatorname{tr}(Q^2)
 \quad (Q\ne0).
\end{equation}
Two maps with images the coordinate axes in $\mathbb C^2$ give $Q=I$ when one adds their orthogonal projections. The bound gives rank two. Two maps onto the same axis give $Q=\operatorname{diag}(2,0)$, with the same trace $2$ but second moment $4$, and the bound gives rank one.

Here is the version needed when some source directions have already been discarded. Let $B_i:V_i\to H$ be maps of finite-dimensional complex inner product spaces, with $D=\dim H>0$, $S=\sum_i\dim V_i>0$ and positive weights $c_i$. A superscript $\dagger$ denotes the adjoint. Put $K_i=c_iB_i^\dagger B_i$ and $Q=\sum_i c_i B_iB_i^\dagger$. Let $U_i$ be subspaces of $V_i$ with $\sum_i(\dim V_i-\dim U_i)\le\alpha S$, and suppose
\begin{equation}\label{D:104}
 \sum_i\|K_i-I\|_{\rm HS}^2\le e S,\quad 0\le e\le1/4,\qquad
 \operatorname{tr}(Q^2)\le(1+\eta)(\operatorname{tr}Q)^2/D.
\end{equation}
Here $\|A\|_{\rm HS}^2=\operatorname{tr}(A^\dagger A)$. \textbf{Lemma.} \textit{Under these hypotheses the codimension of $\sum_i B_i U_i$ in $H$ satisfies}
\begin{equation}\label{D:105}
 \operatorname{codim}\Big(\sum_i B_iU_i\Big)\le\big(2\eta+4\alpha+4\sqrt e\big)D .
\end{equation}
(The proof below gives the sharper bound with $\eta$ in place of $2\eta$. The form with $\eta$ is the one that is used: in section 6.10 for the chain of the source, and in Part III, \S\ref{tr:sec}, Lemma~T1(d).) The display, with $2\eta$, is the finite-dimensional assertion $\mathrm{FrameLA}$ stated in [20] (a statement in the archive, not used by the final theorem; the chain uses a variant, $\mathrm{FrameLA1}$, stated for subspaces and with $\eta$ in place of $2\eta$; see Part III, \S\ref{tr:sec}, \S2).

\textbf{Proof.} We may take $0\le\alpha\le1$, since losing the entire space makes a larger bound automatic. Let $P$ project onto the missing target space, of dimension $zD$. Then $PB_i$ vanishes on $U_i$. Hence $\operatorname{tr}(PQ)$ is at most the sum of the traces of $K_i$ over the omitted source directions, since $c_iB_i^\dagger PB_i\le K_i$ and it vanishes on $U_i$. Write $K_i=I+(K_i-I)$ there and apply Cauchy--Schwarz. Also apply it to the total trace:
\begin{equation}\label{D:106}
 \operatorname{tr}(PQ)\le(\alpha+\sqrt{\alpha e})S,\qquad
 \operatorname{tr}Q\ge (1-\sqrt e)S\ge S/2 .
\end{equation}
Writing $\bar q=\operatorname{tr}Q/D$, the second-moment hypothesis gives $\|Q-\bar q I\|_{\rm HS}^2\le\eta(\operatorname{tr}Q)^2/D$. Pairing with $P$ yields
\begin{equation}\label{D:107}
 z\le\operatorname{tr}(PQ)/\operatorname{tr}Q+\sqrt{\eta z}
   \le 2\alpha+2\sqrt e+(\eta+z)/2.
\end{equation}
This proves even $z\le\eta+4\alpha+4\sqrt e$. The specialization to the actual maps requires a lower first moment and an upper second moment of their period pairings; the latter comes from a relative trace formula. A trace formula equates a sum over a basis of functions with a sum of geometric integrals. Here it controls correlations between the images of different maps. These automorphic estimates are quoted. [7, 19, 20] In the source the index $i$ is a word: a string of $k-2$ binary letters which chooses, at every weight place, one of the weights that interlace a fixed weight of the group in $k-1$ variables; each word also has its own level prime. The space $V_i$ is attached to the word, and $H$ is common to all words. The character twists of section 6.8 are a different index (one twist for each place) and do not enter this family. There is one such family for each translate $u$ of section 4; the weights $w_{u,j}$ depend on the translate and not on the word. See Part III, \S\ref{tr:sec}.

\newpage
\subsection*{6.10 From the rank bound to (M)}

It remains to see why a count of directions can control a mean slope. In this section $H$ is a bound for heights and $u$ is a spread of slopes; they are not the target space $H$ of section 6.9 and the translate $u$ of section 4. Let $d_1,\ldots,d_D$ be real slopes, $a$ a reference value, and assume $d_j\ge a-H$ with $H\ge0$. Put $N(t)=\#\{j:d_j<a-t\}$. Integrating the indicator of each deficit gives
\begin{equation}
 \frac1D\sum_j d_j\ge a-\frac1D\sum_j(a-d_j)_+
     =a-\int_0^H \frac{N(t)}D\,dt .                              \label{II:t6.7}
\end{equation}
For example, slopes $2,3,5$ and $a=4$ give $N(t)=2$ on $0<t<1$, $1$ on $1<t<2$, and $0$ thereafter. The mean deficit is $(2+1)/3=1$. The lower bound $3$ is below the true mean $10/3$ because the largest slope exceeded $4$. This is the layer-cake formula; the form used by the final theorem is in [19].

Here is how an integrable count is obtained in [20]. If $R$ slopes $y_j$ have mean at least $s$ and all are at most $s+u$, $u>0$, the fraction below $s-t$ is at most $u/(u+t)$. Indeed, with fraction $z$ below, their mean is at most $z(s-t)+(1-z)(s+u)$. Apply the rank bound of section 6.9 after discarding those directions; so $\alpha=u/(u+t)$ in the lemma of section 6.9, which gives the term $4u/(u+t)$ below.

(Passing from the bound on the missing directions to a count of the slopes of the target below a given level uses a slope inequality for the maps induced on quotients, and properties of the filtration of a lattice by slopes. This step is no longer quoted as a whole: it is derived in Part III, \S\ref{as:sec}, Proposition~M4. There the slope inequality is Lemma~M3, which is proved, and the properties of the filtration are a cited fact, (C6) of that section (Stuhler, Grayson; cited from memory), which is not proved in this paper. Proposition~M4 takes a rank bound of the form of section 6.9 (the conclusion of Corollary~T2 of \S III.8) as a hypothesis. The derivation there is not the one of the source, which uses the statements of [25].)

In the notation of [20], for numbers $s,h,L,u,H,\eta,e$ with $u>0$ and $L,H,\eta,e\ge0$, the hypotheses take the form
\begin{equation}\label{D:109}
 d_j\ge s-h-H,\quad
 \frac{\#\{j:d_j<s-h-L-t\}}{D}\le 2\eta+\frac{4u}{u+t}+4\sqrt e
 \quad(0\le t\le H).
\end{equation}
Putting $a=s-h-L$ in \eqref{II:t6.7} and integrating gives an error at most $(2\eta+4\sqrt e)H+4u\log(1+H/u)$. In particular it proves the looser bound stated in [20] (the chain of the source uses a variant of it, which is a hypothesis in [7] and is applied for each translate in [28]; see below)
\begin{equation}\label{D:110}
 \frac1D\sum_j d_j\ge s-h-L-u\big(1+8\log(1+H/u)\big)
                      -(4\eta+8\sqrt e)H .
\end{equation}
The trace estimates make $\sqrt e$ exponentially small on the chosen intervals: $e=2e_0$ with $\sqrt{e_0}\le\exp(-x^{\varepsilon/8})$, where $e_0$ is the constant in the bound for the diagonal Gram blocks (Part III, \S\ref{tr:sec}). For $\eta$ the source has $\eta=5x^{-C_H}+3\sqrt{\alpha_{\rm ram}}$. Here $\alpha_{\rm ram}$ is a bound for the proportion of source directions lost at the level primes (it is not the $\alpha=u/(u+t)$ above), the source takes $\sqrt{\alpha_{\rm ram}}=\exp(-x^{\varepsilon/8})$, and $C_H$ is the exponent in the bound $H\le x^{C_H}b$. For all sufficiently large $x$, $(3+4\sqrt2)\exp(-x^{\varepsilon/8})\,x^{C_H}\le1$. Since $\eta=5x^{-C_H}+3\sqrt{\alpha_{\rm ram}}$ and $e=2e_0$, this gives $(\eta+4\sqrt e)x^{C_H}\le6$, and so $(\eta+4\sqrt e)H\le6b$. So $\eta$ is bounded only by a negative power of $x$, chosen to match $H$. The power of $x$ comes from the number of maps, $\lceil3\mathsf G^2x^{C_H}\rceil$, where $\mathsf G=\exp(C_{\mathsf G}\,m^5f\,x^{\varepsilon})$ is the constant in the quoted bounds for the Gram blocks, with a constant $C_{\mathsf G}>0$ (Part III, \S\ref{tr:sec}). In the variant the chain uses, the error bracket is $4u\log(1+H/u)+(\eta+4\sqrt e)H$, with $\eta$ and not $2\eta$. It is multiplied by $1/(1-\delta_{\rm p})$, where $\delta_{\rm p}$ is the proportion of the target not spanned by its lattice (exponentially small, quoted). The lower bound also has $-\log|\mathsf W|$ as a separate term, which this factor does not multiply. Here $\mathsf W$ is the set of words and $\log|\mathsf W|\le(k-3)\log2$, which is negligible; the sizes of $H$ and of $E_{\rm cnt}$ are part of the quoted statements and are not derived here. The bound for the spread $u$ and the arithmetic control of $L$ are also needed. This argument uses a mean slope of the target, where the elementary inequality of section 6.8 by itself would only control a maximal slope.

\textbf{Assembly.} The local derivative bounds place the prime sum in the integral lattices. The period identity makes it possible to compare the different local choices in one space. The height and counting comparisons \eqref{II:t6.5}--\eqref{II:t6.6} give the prime sum's main term. Finally the mean-slope bound accounts for losses from indices and directions which must be discarded. The exact output of the source's comparison is this: for each translate, the required divisibility sum is at most its main term plus an aggregated error [28]. The hypotheses of that statement are defined in [7]. The archive also has the form averaged over the translates, in [7]; the final theorem does not use it. The size estimates bound the error by the unnormalized expression of section 4.3. These are the inputs [5, 7, 19, 28] used at (M).

Once the quoted estimates are granted, sections 4 and 5 give (A), and section 3 concludes by balancing $x/m$ and $m$. A cost of $m^2$ would instead give $m=x^{1/3}$ and exponent $2/3$, if the corresponding inequality were available with the other errors under control.

\par\medskip\noindent\fbox{\parbox{\dimexpr\linewidth-2\fboxsep-2\fboxrule\relax}{\textbf{Proved in Part II, section 6.} Finite-chain integrality, the minor congruence, the Hasse product rule, the reduction \eqref{II:t6.1}, the factorial count \eqref{II:t6.1a}, the Gaussian norm and the Cauchy--Schwarz bound of section 6.5, the count identity \eqref{II:t6.4} and the cancellation \eqref{II:t6.6}, the elementary slope inequality, rational lattice gluing (an illustration over $\mathbb Q$), the rank bound and its integration. (Sections 1--5 are proved in full, given (M) and Estimates B, C.)\par\smallskip \textbf{Quoted from the Lean source.} The normalized theta maps and period identity, their simultaneous local divisibilities and losses, the period and mass formulas, the level bounds, and the norm bounds (apart from the analytic estimates of Part III, \S III.4); the denominator estimates; and the two moments for this family. Part III, \S III.6--III.9 deduce parts of these from named statements of the source, which remain quoted. Part III, \S\ref{as:sec} deduces (M) at one scale from twelve numbered hypotheses, and the Riemann hypothesis from these and a thirteenth; none of the hypotheses is proved outright. See the table at the start of Part III.}}

\newpage
\section*{Sources}\addcontentsline{toc}{section}{Sources}
\small\setlength{\parskip}{2pt}

Alpöge published the Lean source as one archive, \path{lean.tar.gz}, split into fourteen download parts at \href{https://alpo.ge/}{alpo.ge} (listed in footnote 1 of \href{https://alpo.ge/remark.pdf}{his announcement}). Unpacked, it contains 175{,}187 Lean files. The SHA-256 digest of the joined archive is \path{9631b4f020199e580ba86178cb56c6e9b09f690c2369a80dc4334aa93f49c7fc}. The numbered references below are files of that archive; all except [4] are cited in this paper. Each is given by its path from the root of the archive.

\par\hangindent=2.2em\noindent\makebox[2.2em][l]{[1]}{\small \path{lean/src/Development/AutomorphicArithmetic/DedekindLFunctions/LandauTheorems/}\linebreak\path{landauWindow.lean}}

\par\hangindent=2.2em\noindent\makebox[2.2em][l]{[2]}{\small \path{lean/src/Formalism/ProductFormula/StatementOCCM.lean}}

\par\hangindent=2.2em\noindent\makebox[2.2em][l]{[3]}{\small \path{lean/src/Formalism/SizeWindowBrackets/FirstPrimeNeSecondPrime.lean}}

\par\hangindent=2.2em\noindent\makebox[2.2em][l]{[4]}{\small \path{lean/src/Formalism/StatementCM.lean}} {\small (definitions for the earlier form of the argument; see the remark on [2] in section 2)}

\par\hangindent=2.2em\noindent\makebox[2.2em][l]{[5]}{\small \path{lean/src/Development/AutomorphicArithmetic/DedekindLFunctions/LandauTheorems/}\linebreak\path{ChTop.lean}}

\par\hangindent=2.2em\noindent\makebox[2.2em][l]{[6]}{\small \path{lean/src/Formalism/SizeWindowBrackets/GrowthBoundGivesRiemannHypothesis.lean}}

\par\hangindent=2.2em\noindent\makebox[2.2em][l]{[7]}{\small \path{lean/src/Formalism/RankGapBounds/MeanBoundSecondForm.lean}}

\par\hangindent=2.2em\noindent\makebox[2.2em][l]{[8]}{\small \path{lean/src/Development/AutomorphicArithmetic/CharacterBoxes/ExistsPageN.lean}}

\par\hangindent=2.2em\noindent\makebox[2.2em][l]{[9]}{\small \path{lean/src/Development/AutomorphicArithmetic/MeanSquareBounds/pairEx.lean}}

\par\hangindent=2.2em\noindent\makebox[2.2em][l]{[10]}{\small \path{lean/src/Development/AutomorphicArithmetic/CharacterBoxes/WinPairs.lean}}

\par\hangindent=2.2em\noindent\makebox[2.2em][l]{[11]}{\small \path{lean/src/Development/AutomorphicArithmetic/PrimitiveCosets/CoreGm.lean}}

\par\hangindent=2.2em\noindent\makebox[2.2em][l]{[12]}{\small \path{lean/src/Supplements/Kernels/MainTheorems.lean}}

\par\hangindent=2.2em\noindent\makebox[2.2em][l]{[13]}{\small \path{lean/src/Development/AutomorphicArithmetic/AdelicQuotients/ThetaLiftLinear.lean}}

\par\hangindent=2.2em\noindent\makebox[2.2em][l]{[14]}{\small \path{lean/src/Formalism/ThetaLiftIndependence/Statement.lean}}

\par\hangindent=2.2em\noindent\makebox[2.2em][l]{[15]}{\small \path{lean/src/Formalism/SharpRestriction/GrowthBounds/CoarseMultiplicity/KeyMinorDvd.lean}}

\par\hangindent=2.2em\noindent\makebox[2.2em][l]{[16]}{\small \path{lean/src/Development/AutomorphicSpectrum/Proofs/GeometryProofs/}\newline\path{ChowlaSelbergPeriodStatement.lean}}

\par\hangindent=2.2em\noindent\makebox[2.2em][l]{[17]}{\small \path{lean/src/Formalism/BehrendSets/StatementE.lean}}

\par\hangindent=2.2em\noindent\makebox[2.2em][l]{[18]}{\small \path{lean/src/Formalism/ThickAverageLoss/PartnersToFactorsAt.lean}}

\par\hangindent=2.2em\noindent\makebox[2.2em][l]{[19]}{\small \path{lean/src/Formalism/SharpRestriction/GlobalTheoremInputs.lean}}

\par\hangindent=2.2em\noindent\makebox[2.2em][l]{[20]}{\small \path{lean/src/Formalism/CorankShare/StatementFR.lean}}

\par\hangindent=2.2em\noindent\makebox[2.2em][l]{[21]}{\small \path{lean/src/Formalism/ThetaLogCoordinates/UniformThetaSubstitution/StatementAB.lean}}

\par\hangindent=2.2em\noindent\makebox[2.2em][l]{[22]}{\small \path{lean/src/Formalism/LocalRingCardinality/DataParams}\textsubscript{2}\path{.lean}} {\small (the file name contains the subscript character U+2082)}

\par\hangindent=2.2em\noindent\makebox[2.2em][l]{[23]}{\small \path{lean/src/Formalism/SharpRestriction/GlobalSchurBounds/FlipTwoThree.lean}}

\par\hangindent=2.2em\noindent\makebox[2.2em][l]{[24]}{\small \path{lean/src/Formalism/SharpRestriction/GlobalSchurBounds/}\newline\path{IntegrableOfContinuousCircleMeasure.lean}}

\par\hangindent=2.2em\noindent\makebox[2.2em][l]{[25]}{\small \path{lean/src/Formalism/HarderNarasimhan/StatementPl.lean}}

\par\hangindent=2.2em\noindent\makebox[2.2em][l]{[26]}{\small \path{lean/src/Formalism/LocalRingCardinality/Params.lean}}

\par\hangindent=2.2em\noindent\makebox[2.2em][l]{[27]}{\small \path{lean/src/Formalism/FrobeniusOrbits/FrobeniusOrbitsProofs/ThickExists.lean}}

\par\hangindent=2.2em\noindent\makebox[2.2em][l]{[28]}{\small \path{lean/src/Formalism/SizeWindowBrackets/ChainSlopeWindows/LevelOfCongruence.lean}}

\par\hangindent=2.2em\noindent\makebox[2.2em][l]{[29]}{\small \path{lean/src/Development/AutomorphicArithmetic/CharacterBoxes/PairGood.lean}}
}
\newpage
\sloppy
\setcounter{section}{0}\setcounter{equation}{0}
\renewcommand{\thesection}{III.\arabic{section}}
\renewcommand{\theequation}{III.\arabic{equation}}
\stepcounter{partno}\begin{center}\Large\bfseries Part III. Proofs of the quoted statements\end{center}
\phantomsection\addcontentsline{toc}{part}{Part III. Proofs of the quoted statements}

\noindent Parts I and II deduce the Riemann hypothesis from statements marked \quoted. Part III is meant to
prove those statements, in the same classical language, so that the three parts together form a complete
proof. It is being written in stages, and the table below records what has been done.

{\small\setlength{\LTleft}{\fill}\setlength{\LTright}{\fill}
\begin{longtable}{p{.36\textwidth}p{.56\textwidth}}
\toprule
Statement of Parts I and II & Status in Part III\\
\midrule\endhead
\bottomrule\endfoot
Estimate B: the auxiliary primes $q$, $d$ (\S III.1) & proved, from four cited theorems (Q1)--(Q4)\\
Estimate C: the characters of $E$ at the large scale (\S III.2) & proved, from three cited theorems (Z1)--(Z3) and Theorem~\ref{thm:pair}(e)\\
The sets $T_2$ and $T$ of columns (\S III.3) & constructed, with proof (elementary). In (M) they enter through the densities (H0)(h) and the loss bound (H7) of \S\ref{as:sec}; (H7) is quoted\\
The inequality (M), size of the theta kernels (\S I.7, item 2; \S III.4) & the analytic estimates for one block, one map and a chain are proved in an abstract setting; that the maps of the construction are instances of it is quoted (item (vi) of the list in \S III.4); four passages are conditional (the reference mean square, the eigenvalues of a lift, the row-by-row comparison, the pairing at the end of the chain); the arithmetic inputs listed at the end of the section are not proved\\
The inequality (M), divisibility at inert primes (\S I.7, item 1; \S III.5) & written up in \S III.5, conditionally on the integrality of the theta kernels in the coordinates of the height-two logarithm (hypothesis (B)), which is taken from the Lean development. Under (B) the local estimate for power series is proved for orbits of length two; for longer orbits it also uses a duality statement cited from the Lean. That the theta kernels satisfy (B), the frames and labels, and the remaining primes are not proved\\
The inequality (M), from derivatives to lattices (\S III.6) & proved: the lemmas on slopes of lattices, the form of the prime sum and its splitting by classes of primes, the digit losses, the inequality at one class of primes, and bounds for several terms of the sum over primes. Quoted: the bounds for the eleven loss sums and for their total. Conditional on two statements of the source (on valuations and on lattices): the theorem on the glued lattice. A model case only: the comparison of the two sides. Quoted: the statement on valuations itself\\
The inequality (M), arithmetic inputs of the size bound (\S III.7) & proved: the size inequality \eqref{III:sizefinal} follows from eleven named inputs; an index-corrected mean of periods equals $\log\Omega_E$, under the hypotheses on the periods and the field stated there (its identification with the periods of the input (S4) rests on two further statements of the source and one that is proved there). The eleven inputs are quoted, among them the inner-product formula, the covolumes and the local index bounds\\
The inequality (M), the trace-formula estimates (\S III.8) & proved: the dimension bound of Part II, \S6.9 follows from three statements of the source (on traces, ranks and pairings) and the size and subspace conditions of Corollary~T2; a finite model of the trace estimates. The three statements, and the theorems behind them, are quoted\\
The inequality (M), the bound for Galois cohomology (\S III.9) & proved: the step from the bound and a quoted comparison inequality to the bound for denominators, and vanishing in some cases over $\Q$ from cited theorems. The bound itself and the comparison inequality are quoted, with an account of the proof of the bound in the source\\
The objects of the construction (\S III.10) & definitions, in classical terms. Proved: the structure of the field $E$ and (from cited theorems) how primes split in it; the rule for the cells of the table; slopes of lattices and of their duals; convergence, size and Taylor coefficients of the theta series of a pair of lattices. Assumed from the source, as seven numbered hypotheses: that the spaces, lattices, parameters and weights exist with the stated properties, that the maps of the chain are given by these theta series, and that the kernels are algebraic, integral and of the stated size. A dictionary relates the settings of \S III.4--III.9 to these objects; some of its rows are not verified\\
The inequality (M), the assembly (\S\ref{as:sec}) & proved: (M) for one scale $x$ follows from twelve numbered Hypotheses (H0)--(H11) (Theorem~M, Corollary~M9). None of the twelve is proved outright. Six are derived, there or in \S III.6--III.8, from named statements of the source, which are quoted or assumed: (H0), (H1), (H6), (H8), (H10), (H11). One, (H7), is quoted. Four, (H2), (H3), (H5), (H9), are assumed, and (H4) is assumed in part\\
The Riemann hypothesis from numbered Hypotheses (\S\ref{as:sec}, \S7) & proved: a thirteenth Hypothesis, (H12), implies the Riemann hypothesis (Corollary~M11), by Theorem~M, Parts~I and~II and \S III.2. (H12) says that data satisfying (H0)--(H11) exist at every $x$ of intervals $[Z,Z^{3/2}]$ with $Z$ arbitrarily large, with constants that do not depend on $x$ and with primes $q$, $d$ for which Estimates~B and~C hold at $x$. (H12) is assumed. \S III.1 proves that primes as in its clause (b) exist. Estimate~C for such primes is proved in \S III.2; it is not part of (H12)\\
The hypotheses of \S III.10 & assumed; not proved in these pages\\
\end{longtable}}

\noindent So at present Parts I--III prove the Riemann hypothesis from the inequality (M), on the intervals of
Part~II, \S3, and from the cited theorems; (M) is not proved. \S\ref{as:sec} reduces (M) at one scale to the
Hypotheses (H0)--(H11), and the Riemann hypothesis to the Hypotheses (H0)--(H12) and the cited theorems
(Corollary~M11 there). None of the thirteen Hypotheses is proved outright in this paper;
\S III.4--III.10 treat them to the extent stated in the table, conditionally as stated there. That the
Hypotheses hold at every scale of such intervals, with constants that do not depend on the scale and for
primes for which Estimates~B and~C hold, is Hypothesis (H12). It is asserted by the source and is assumed
here.

\noindent The statements below are translated from the Lean development of the announcement. In Part III
``the source'' always means this Lean development. For \S III.1--III.3 each is a statement that the
final theorem actually depends on (for \S III.4--III.10 see the remarks in those sections): we computed the set of all declarations on which the Lean theorem
\texttt{riemannHypothesis} depends, and checked that these belong to it.

\noindent A caution about this check. A declaration belongs to that set as soon as the proof term of the final
theorem refers to it, directly or through other declarations. That is necessary for a statement to be used in the
proof, but it is not sufficient: a statement that is carried along inside a record of statements is counted even if
the argument never consults it. An example: among the declarations on which \texttt{riemannHypothesis} depends in
this sense is one whose statement is the proposition \texttt{True}. So in these
pages ``the final theorem depends on it'' means membership in that set, and nothing more. Where we say that a
statement is not among the declarations the final theorem depends on, that is conclusive: such a statement is not
used. The census of \S\ref{cs:sec} gives a second example: ten of the statements listed there are never read
in the theorems they are passed through, as far as this was followed (eight are entries of a record that are
never taken out of it, and two occur only in the proofs of these), and the final theorem depends on each of
them.

\noindent For \S III.1--III.3 the check was also made for the lemmas of the Lean proofs that the sections
follow (for instance \texttt{QH.thinH}, \texttt{twistK}, \texttt{meanSqK}, \texttt{bad\_share},
\texttt{tailK}, \texttt{count\_bad}, \texttt{hostThickExists}, \texttt{topsChoiceH}): all are among the
declarations the final theorem depends on.

\emph{A correction to the source references of Part II.} The same computation shows that several
statements cited in Part II are present in the archive but are \emph{not} among the $1{,}478{,}198$
declarations on which \texttt{riemannHypothesis} depends: the Fock norm and jet definitions of [23]; the
statements \texttt{ChainInt} of [14], \texttt{FrameLA} and \texttt{FrameLast} of [20], and
\texttt{CountIdentity} of [17] (a statement of the same name elsewhere is used); the definitions concerning
Behrend sets in [17]; and the inequality of [6] in the form quoted. These are
parallel or earlier formulations. The same holds for the linearity statement of [13]. The mathematics of Part II does not depend on them, since what Part II
proves it proves directly; but for these items the phrase ``the source's statement'' should be read as ``a
statement in the archive''. The statements that are used include: the Landau lemma on intervals [1], the
auxiliary primes [9, 29], the weighted sums \texttt{ASK} and \texttt{MAIN} [3], the height of a map [25], the
Chowla--Selberg period [16], the minors congruence [15], the digit inequality \texttt{StarArith} [21], the
layer-cake formula, the Bessel identity and the local comparisons of [19], the bound for each translate
\texttt{each\_unit\_bound\_statement} of [28] with the hypotheses defined in [7] (the averaged statement
\texttt{MeanBoundSecondForm} of [7] is not used), and the bound for Galois cohomology [8].

\section{The auxiliary primes: proof of Estimate B}

For a primitive Dirichlet character $\chi\ne1$ put, as in Part I,
\begin{equation}\label{D:111}
 G_\chi(y)=\sum_{n\le y}\Lambda(n)\chi(n)\Big(\frac yn-1\Big)+y\,\frac{L'}{L}(1,\chi).
\end{equation}
For an odd prime $p$ write $\chi_p$ for the Legendre symbol modulo $p$.

\begin{theorem}[auxiliary primes; \texttt{QH.pairEx} in {[9]}, statement in {[29]}]\label{thm:pair}
Let $M\ge1$ be an integer, $0<\varepsilon<\frac1{52}$ and $A\ge1$. There is $Y_1$ such that for every
$Y\ge Y_1$ there are primes $q$ and $d$ with the following properties.
\begin{enumerate}[label=\textup{(\alph*)},leftmargin=2.2em]
\item $q\equiv3\bmod8$ and $\frac12Y^{\varepsilon/5}<\log q\le Y^{\varepsilon/5}$.
\item $d\equiv7\bmod16$ and $\frac12Y^{\varepsilon/5}(\log Y)^2<\log d\le Y^{\varepsilon/5}(\log Y)^2$.
\item $q$ divides $d+1$, and so does every odd prime $p\le Y^{\varepsilon/5}$. Neither $q$ nor $d$ divides $M$.
\item Let $\xi$ be a Dirichlet character modulo $M$ and let $(i,j)\in\{0,1\}^2$, $(i,j)\ne(0,0)$. Let $\chi$ be
the primitive character that induces $\xi\chi_q^i\chi_d^j$. Then for all $y$ with $Y\le y\le Y^A$
\begin{equation}\label{D:112}
 |G_\chi(y)|\le y^{1/2+\varepsilon/2}\qquad\text{and}\qquad\Big|\frac{L'}{L}(1,\chi)\Big|\le4\log y .
\end{equation}
\item Let $\rho$ be a real primitive character whose conductor divides $8M$, let $(i,j)\ne(0,0)$, and let
$\chi$ be the primitive character that induces $\rho\chi_q^i\chi_d^j$. Then $L(s,\chi)\ne0$ in the region
\begin{equation}\label{D:113}
 \re s>1-\tfrac1{400},\qquad |\mathrm{Im}\,s|\le 8M\exp\!\big(Y^{\varepsilon/5}\big)\exp\!\big(j\,Y^{\varepsilon/5}(\log Y)^2\big).
\end{equation}
\end{enumerate}
\end{theorem}

\noindent\emph{Remarks.} (1) For $\zeta$ one takes $M=1$ and $A\ge\frac32$. Then (d) gives Estimate B of Part I, with $x=y$, the interval $[Y,Y^{3/2}]$ contained in $[Y,Y^A]$, and the
exponent $\frac12+\frac\varepsilon2$ in place of $\frac12+\varepsilon$.

(2) The conductors are far larger than $y$: $\log q$ is a power of $Y$. So (d) cannot be a statement about
a fixed character as $y\to\infty$. It is a statement about well-chosen $q$ and $d$.

(3) Condition (c) has a consequence that is not visible in Part I. If $p$ is an odd prime with
$p\le Y^{\varepsilon/5}$, then $d\equiv-1\bmod p$, and since $d\equiv3\bmod4$ quadratic reciprocity gives
$\chi_d(p)=\big(\tfrac{-1}p\big)(-1)^{(p-1)/2}=1$. Also $\chi_d(2)=1$. So \emph{every prime up to
$Y^{\varepsilon/5}$ splits in $\Q(\sqrt{-d})$}; the inert primes, which carry the argument of Part I, all
exceed $Y^{\varepsilon/5}$.

(4) Part (e) is used for Estimate C: it excludes exceptional real zeros for the real characters whose
conductor is divisible by $q$ or $d$.


\subsection*{Proof of Theorem \ref{thm:pair}}

The proof chooses $q$ at random from a set of primes, then $d$ at random from a second set, and shows that
most choices work. Two things must hold for the chosen primes: a sum over $n\le Z$ must be small (this comes
from a mean square over the primes), and a tail $n>Z$ must be small (this comes from a zero-free box, which
holds for all but few primes by a zero-density theorem). Throughout, $\tau=\varepsilon/2$,
\begin{equation}\label{D:114}
 Z=Y^{400A},\qquad z=Y^{\varepsilon/5},\qquad z'=Y^{\varepsilon/5}(\log Y)^2 ,
\end{equation}
and $Y$ is large in terms of $M$, $\varepsilon$, $A$. Constants $C$, $c$ are absolute, or depend only on $A$, unless indicated.

\textbf{Cited theorems.}
\begin{enumerate}[label=\textup{(Q\arabic*)},leftmargin=2.6em]
\item \emph{Zero-free region, with uniqueness} (Landau, Page; \cite[Ch.~14]{Dav00}, \cite{Pag35}; in this
explicit form \cite[\S1]{TZ21}). Let $c_1=\frac1{13}$. For every modulus $m\ge2$, the functions $L(s,\chi)$,
$\chi$ modulo $m$, have no zero with
$\beta\ge1-c_1/\log(m(|\gamma|+3))$, except possibly one simple real zero $\beta_1$ of one real nonprincipal
character $\chi_1$. Moreover $\beta_1\le1-c/m$; this follows from the bound $\beta_1\le1-c'm^{-1/2}$ recorded in \cite[\S1]{TZ21}. This is the exceptional zero meant in (Q2).
\item \emph{Primes in progressions to large moduli} (Thorner and Zaman \cite[Corollary 1.4]{TZ21}; results
of this type go back to Linnik and to Gallagher \cite{Gal70}, see also \cite{Bom87}). There is an absolute
$c_3>0$ such that for coprime $m\ge2$ and $b$, and $t\ge m^{12}$,
\begin{equation}\label{D:115}
 \sum_{p\le t,\ p\equiv b\,(m)}\log p=\Lambda_b(t)\big(1+O(\kappa(t))\big),\qquad
 \Lambda_b(t)=\frac1{\varphi(m)}\Big(t-\chi_1(b)\frac{t^{\beta_1}}{\beta_1}\Big),
\end{equation}
with $\kappa(t)=\exp(-c_3\log t/\log m)+\exp(-c_3(\log t)^{3/5}(\log\log t)^{-1/5})$. The term with $\chi_1$
is present only if the exceptional zero of (Q1) exists for the modulus $m$.
\item \emph{Zero density} (Montgomery \cite[Theorem 12.2]{Mon71}). The number of primitive characters $\chi$
of conductor at most $Q$ for which $L(s,\chi)$ has a zero with $\beta>1-\frac1{400}$, $|\gamma|\le Q$, is at
most $CQ^{1/4}$.
\item \emph{The logarithmic derivative in a zero-free box} (a standard consequence of Landau's lemma
\cite[\S3.9]{Tit86}, \cite[Lemma 6.3]{MV07}, and the convexity bound for $L(s,\chi)$). Let $\chi\ne1$ be primitive of conductor at most $Q$, and suppose $L(s,\chi)\ne0$ for
$\re s>1-\frac1{400}$, $|\mathrm{Im}\,s|\le Q$. Then for each $\delta>0$ there is $K_\delta$ with
$|\frac{L'}L(s,\chi)|\le K_\delta\log Q$ for $\re s\ge1-\frac1{400}+\delta$, $|\mathrm{Im}\,s|\le Q-1$.
\end{enumerate}
\emph{The form of (Q2) that is used.} Let $U\ge m^{24}$, and let
$\pi^*(U;m,b)$ be the number of primes $p\equiv b\bmod m$ with $U^{1/2}<p\le U$. Then
\begin{equation}\label{III:Q2}
 \big|\pi^*(U;m,b)-\mathcal M_b\big|\le k\,\mathcal M_b,\qquad
 \mathcal M_b=\frac1{\varphi(m)}\int_{U^{1/2}}^U\big(1-\chi_1(b)t^{\beta_1-1}\big)\frac{dt}{\log t},
\end{equation}
with $k=A_0\big(e^{-c_4\log U/\log m}+e^{-\sqrt{c_4\log U}}\big)$. \emph{Proof.} On $[U^{1/2},U]$ we have $t\ge m^{12}$, and $(\log t)^{3/5}(\log\log t)^{-1/5}\ge\frac12\sqrt{\log U}$, so that
the relative error in (Q2), with its implied constant, is at most $k/20$ for suitable $A_0$, $c_4$. Write $\theta(t)$ for the sum in (Q2). By partial summation
$\pi^*=\theta(U)/\log U-\theta(U^{1/2})/\log U^{1/2}+\int_{U^{1/2}}^U\theta(t)\,dt/(t\log^2t)$. Replacing $\theta$ by
$\Lambda_b$ gives exactly $\mathcal M_b$, because $\Lambda_b'(t)=(1-\chi_1(b)t^{\beta_1-1})/\varphi(m)$. The function
$\Lambda_b$ is positive and increasing, so the error is at most
$\frac k{20}\Lambda_b(U)\big(\frac1{\log U}+\frac2{\log U}+\frac2{\log U}\big)$. Finally
$\Lambda_b(U)\le4\int_{U/2}^U\Lambda_b'\le4\log U\cdot\mathcal M_b$: if $\chi_1(b)=1$ the derivative $\Lambda_b'$ is
increasing, and otherwise it lies between $1/\varphi(m)$ and $2/\varphi(m)$. \qed

The Lean development proves all four. For (Q2) and (Q3) it has its own proofs: (Q2) from a zero-density
estimate and the Deuring--Heilbronn phenomenon, (Q3) by a mollifier and the Hal\'asz--Montgomery inequality.
In (Q3) the published theorem is much stronger
than what is used.

\textbf{The sets of primes.} Let $m_0$ be a positive integer divisible by $8$, $a_0$ an integer prime to $m_0$, and $U>1$. Put
\begin{equation}\label{D:117}
 P=P(m_0,a_0,U)=\{q\text{ prime}:\ U^{1/2}<q\le U,\ q\equiv a_0\bmod m_0\}.
\end{equation}
Call $(m_0,a_0,U)$ \emph{admissible} if
\begin{equation}\label{III:adm}
 \tfrac14(\log Y)^2\log m_0\le\log U,\qquad Y^{\varepsilon/6}\le\log U\le Y^{\varepsilon/4},
\end{equation}
and every prime factor of $m_0$ is either at most $Y^{1/4}$ or greater than $Z^2$.

\textbf{Step 1: characters over the primes of $P$.} \emph{Let $(m_0,a_0,U)$ be admissible, let $B>0$, and let $Y$ be large in terms of $B$ as well. Then
$|P|\ge U^{9/10}$. Moreover, among the odd squarefree $s$ with $1<s\le Z^2$ and $(s,m_0)=1$, at most one fails to satisfy the
following bound; if there is one, call it $n^*$:}
\begin{equation}\label{III:twist}
 \Big|\sum_{q\in P}\Big(\frac qs\Big)\Big|\le Y^{-B}|P| .
\end{equation}
\emph{Proof.} Apply \eqref{III:Q2} with modulus $m=m_0s$. By \eqref{III:adm}, $\log m\le\log m_0+800A\log Y$ and
$\log U/\log m\ge c(\log Y)^2$ (if $\log m_0\le\log Y$ use $\log U\ge Y^{\varepsilon/6}$). So $U\ge m^{24}$ and $k\le Y^{-B-1}$. Every $q\in P$ exceeds $Z^2\ge s$, since $\log q>\frac12Y^{\varepsilon/6}$, so $q\nmid s$. Split $P$ into residue classes $b$ modulo $m_0s$ with $b\equiv a_0\bmod m_0$,
$(b,s)=1$; the symbol $(\frac bs)$ is constant on each, and the class contains $\pi^*(U;m,b)$ elements of
$P$. So
\begin{equation}\label{D:120}
 \sum_{q\in P}\Big(\frac qs\Big)=\sum_b\Big(\frac bs\Big)\mathcal M_b+O\Big(k\sum_b\mathcal M_b\Big).
\end{equation}
\emph{The error.} Since $k\le\frac12$, \eqref{III:Q2} gives $\mathcal M_b\le2\pi^*(U;m,b)$, so
$\sum_b\mathcal M_b\le2|P|$. (This is why a relative error is needed: $P$ may be small if the modulus $m_0$
itself has an exceptional zero.)

\emph{The main terms.} The part of $\mathcal M_b$ without $\chi_1$ is the same for all $b$, and
$b\mapsto(\frac bs)$ is a nontrivial character on the classes $b\equiv a_0\bmod m_0$ ($s$ is squarefree, odd,
prime to $m_0$ and greater than $1$). So these parts cancel. The parts with $\chi_1$ cancel in the same way
unless $\chi_1(b)(\frac bs)$ is constant on those classes. In that case call $s$ \emph{bad} if in addition
$\beta_1>1-\frac{c_1}{3\log(m_0Z^2)}$.

\emph{At most one $s$ is bad.} If $s_1$, $s_2$ are both bad, apply (Q1) to the modulus $m_0s_1s_2$; since
$\log(3m_0s_1s_2)\le3\log(m_0Z^2)$, both exceptional zeros lie in its region, so the two exceptional
characters are induced by the same character. Hence $(\frac u{s_1})=(\frac u{s_2})$ for $u\equiv1\bmod m_0$
prime to $s_1s_2$, and $s_1=s_2$. So at most one $s=n^*$ is bad.

\emph{The other $s$.} If $s$ is not bad and the parts with $\chi_1$ do not cancel, then
$t^{\beta_1-1}\le U^{-(1-\beta_1)/2}\le U^{-c_1/(6\log(m_0Z^2))}\le Y^{-B-1}$ on the range of integration, so the
part with $\chi_1$ is at most $Y^{-B-1}$ times the part without it, hence at most $2Y^{-B-1}\mathcal M_b$.
Altogether $|\sum_{q\in P}(\frac qs)|\le(2k+4Y^{-B-1})|P|\le Y^{-B}|P|$.

\emph{The size of $P$.} Apply \eqref{III:Q2} with $m=m_0$, $b=a_0$: $|P|\ge\frac12\mathcal M_{a_0}$. By (Q1),
$1-\beta_1\ge c/m_0$, so $1-t^{\beta_1-1}\ge c'/m_0$ for $t\ge U/2$, and
$\mathcal M_{a_0}\ge c'U/(2m_0^2\log U)$. Since $\log m_0\le4\log U/(\log Y)^2$, this exceeds $2U^{9/10}$. \qed

\textbf{Step 2: the mean square.} For $y\ge1$, an odd prime $q$ and a function $\kappa$ with $|\kappa(n)|\le1$ put
\begin{equation}\label{D:121}
 S(y,q,\kappa)=\sum_{n\le Z}\Lambda(n)\min\Big(1,\frac yn\Big)\kappa(n)\Big(\frac nq\Big).
\end{equation}
\emph{Let $(m_0,a_0,U)$ be admissible. Then for $Y\le y\le Y^A$ and every such $\kappa$,}
\begin{equation}\label{III:ms}
 \sum_{q\in P}|S(y,q,\kappa)|^2\le C\,y\log y\,|P| .
\end{equation}
\emph{Proof.} Write $c_n=\Lambda(n)\min(1,y/n)$. Expanding the square,
\begin{equation}\label{D:123}
 \sum_{q\in P}|S|^2\le\sum_{n,n'\le Z}c_nc_{n'}\Big|\sum_{q\in P}\Big(\frac{nn'}q\Big)\Big| .
\end{equation}
Every $q\in P$ exceeds $Z^2\ge nn'$, since $\log q>\frac12Y^{\varepsilon/6}$. Let $s$ be the product of the
primes that do not divide $m_0$ and occur in $nn'$ to an odd power. All $q\in P$ lie in one residue class
modulo $m_0$, and $8\mid m_0$; so by quadratic reciprocity and its supplements $(\frac{nn'}q)=\pm(\frac qs)$
with a sign that does not depend on $q$. By Step~1 the inner sum is at most $Y^{-A-1}|P|$ unless $s=1$ or
$s=n^*$, when we bound it by $|P|$.

The terms $n,n'$ are prime powers. Put $\sigma(p)=\sum_{p^j\le Z}c_{p^j}$. One checks $\sigma(p)\le3\log y$,
and $\sigma(p)\le2y\log p/p$ for $p>y$, whence $\sum_p\sigma(p)^2\le Cy\log y$ by Chebyshev's bound. Call $n$
\emph{neutral} if it is an even power of a prime or a power of a prime dividing $m_0$; let
$R$ be the sum of $c_n$ over such $n\le Z$. Even powers contribute at most $C\sqrt y$. The primes dividing
$m_0$ that are at most $Z$ are at most $Y^{1/4}$, so there are at most $Y^{1/4}$ of them, each contributing
at most $3\log y$. So $R\le C\sqrt y$.

\emph{Pairs with $s=1$.} Either $n$, $n'$ are powers of the same prime, or both are neutral.
Their contribution is at most $\big(\sum_p\sigma(p)^2+R^2\big)|P|\le Cy\log y\,|P|$.

\emph{Pairs with $s=n^*$.} Since $n$, $n'$ are prime powers, such pairs exist only if $n^*$ is a prime or a
product of two primes; then one of $n,n'$ is a power of
a fixed prime while the other is a power of a second fixed prime or is neutral. The
contribution is at most $(2\cdot9\log^2y+6\log y\cdot R)|P|\le Cy|P|$.

\emph{All other pairs.} Since $\sum_{n\le Z}c_n\le y\sum_{n\le Z}\Lambda(n)/n\le Cy\log Z$, they contribute at
most $Y^{-A-1}(Cy\cdot400A\log Y)^2|P|\le Cy|P|$, because $y\le Y^A$. \qed

\textbf{Step 3: all $y$ at once.} \emph{Let $(m_0,a_0,U)$ be admissible and $\kappa$ as before. The number of
$q\in P$ for which $|S(y,q,\kappa)|>\frac12y^{1/2+\tau}$ for some $y\in[Y,Y^A]$ is at most $C_\tau Y^{-\tau}|P|$.}

\noindent\emph{Proof.} As a function of $y$, $S$ is continuous and piecewise linear, with derivative
$D(y)=\sum_{y<n\le Z}\Lambda(n)\kappa(n)(\frac nq)/n$. The point is that $y\,D(y)=S(y,q,\kappa_y)$, where
$\kappa_y$ is $\kappa$ restricted to $n>y$. So \eqref{III:ms}, which is uniform in $\kappa$, gives also
$\sum_q|D(y)|^2\le C\log y\,|P|/y$. On a block $I=[W,\min(2W,Y^A)]$ with $Y\le W$, by the Cauchy--Schwarz inequality,
$|S(y)|^2\le2|S(W)|^2+2W\int_I|D|^2$, so
$\sum_q\max_{y\in I}|S(y)|^2\le CW\log W\,|P|$. By Chebyshev's inequality the number of $q$ with
$\max|S|>\frac12W^{1/2+\tau}$ on the block is at most $CW^{-2\tau}\log W\,|P|\le(C/\tau)W^{-\tau}|P|$. Sum
over the blocks $W=2^iY$. \qed

\textbf{Step 4: the tail.} For a primitive $\chi\ne1$ and $w\ge1$ put
\begin{equation}\label{D:124}
 t_\chi(w)=-\frac{L'}L(1,\chi)-\sum_{n\le w}\frac{\Lambda(n)\chi(n)}n ,
\end{equation}
the tail of the series for $-\frac{L'}L(1,\chi)$. \emph{Let $\chi$ have conductor at most $Q$ and suppose
$L(s,\chi)\ne0$ for $\re s>1-\frac1{400}$, $|\mathrm{Im}\,s|\le Q$. Then for $Q\ge2$, $w\ge2$, $0<\eta\le1$ and
$0<\delta\le\frac1{800}$,}
\begin{equation}\label{III:tail}
 |t_\chi(w)|\le K'_\delta\Big((\log Q)^2w^{-1/400+\delta}+\frac{w\log Q}{\eta\,Q}\Big)+\frac{2(w\eta+1)\log(2w)}{w}.
\end{equation}
\emph{Proof.} Let $G_\chi$ be as above. For $h=w\eta$ one has directly from the definitions
\begin{equation}\label{D:126}
 G_\chi(w+h)-G_\chi(w)=-h\,t_\chi(w)+\sum_{w<n\le w+h}\Lambda(n)\chi(n)\Big(\frac{w+h}n-1\Big),
\end{equation}
and the last sum is at most $h(h+1)\log(2w)/w$ in absolute value. On the other hand, by Perron's formula
with the kernel $y^s/(s(s-1))$,
\begin{equation}\label{D:127}
 G_\chi(y)-y\frac{L'}L(1,\chi)=\frac1{2\pi i}\int_{\re s=2}-\frac{L'}L(s,\chi)\frac{y^s}{s(s-1)}ds .
\end{equation}
Move the part $|\mathrm{Im}\,s|\le Q/2$ of the contour to the line $\re s=\sigma_1=1-\frac1{400}+\delta$. No
zero is crossed. The pole at $s=1$ has residue $-y\frac{L'}L(1,\chi)$, which cancels the term on the left.
So $G_\chi(y)$ is the integral over the new contour. Take the difference at $y=w+h$ and $y=w$. On the line
$\re s=\sigma_1$ we have $|(w+h)^s-w^s|\le4\eta|s|w^{\sigma_1}$, which cancels one factor $|s|$, and (Q4)
bounds the logarithmic derivative; this part is at most $K_\delta\log Q\cdot\eta\,w^{\sigma_1}\log Q$. The
horizontal segments and the remaining parts of the line $\re s=2$ contribute at most $K_\delta\log Q\cdot w^2/Q$.
Divide by $h=w\eta$. \qed

\emph{Two consequences.} Let $\chi$ be as in Step 4, with $Y^{\varepsilon/6}\le\log Q\le Y^{\varepsilon/4}$.

(a) \emph{$Y^A|t_\chi(Z)|\le Y^\varepsilon$.} Take $w=Z=Y^{400A}$, $\eta=Y^{-2A}$, $\delta=\varepsilon/(1600A)$ in
\eqref{III:tail}. The first term gives $Y^A\cdot Y^{\varepsilon/2}\cdot Y^{-A+\varepsilon/4}$; the second is
at most $Y^{404A}e^{-Y^{\varepsilon/6}}$; the third is at most $801A\,Y^{-A}\log Y$. This is where the numbers $400$
and $\frac1{400}$ are matched: $Z^{-1/400}=Y^{-A}$ pays for the factor $y\le Y^A$.

(b) \emph{If $\log Q\le3z'$ and $\varepsilon<\frac1{52}$, then $|\frac{L'}L(1,\chi)|\le4\log Y$.} Take
$w=(\log Q)^{1000}$, $\eta=(\log Q)^{-5}$ and $\delta=\frac1{4000}$, so that
$(\log Q)^2w^{-1/400+\delta}=(\log Q)^{-1/4}$; then \eqref{III:tail} gives $|t_\chi(w)|\le C$, and
$|\sum_{n\le w}\Lambda(n)\chi(n)/n|\le\log w+C=1000\log\log Q+C$. Now
$\log\log Q\le\frac\varepsilon5\log Y+2\log\log Y+\log3$. So
$|\frac{L'}L(1,\chi)|\le200\varepsilon\log Y+2000\log\log Y+C$, and $200\varepsilon<\frac{50}{13}$. This step
needs $\varepsilon<\frac1{50}$; it is the one that comes closest to using the hypothesis $\varepsilon<\frac1{52}$.

\textbf{Step 5: from $S$ and the tail to $G_\chi$.} For $y\le Z$, and $\chi(n)=\kappa(n)(\frac nq)$ for $n\le Z$,
\begin{equation}\label{III:split}
 G_\chi(y)=-S(y,q,\kappa)-y\,t_\chi(Z).
\end{equation}
Indeed
\begin{equation}\label{D:129}
 -S=-\sum_{n\le y}\Lambda(n)\chi(n)-y\sum_{y<n\le Z}\frac{\Lambda(n)\chi(n)}n,\qquad
 -y\,t_\chi(Z)=y\frac{L'}L(1,\chi)+y\sum_{n\le Z}\frac{\Lambda(n)\chi(n)}n .
\end{equation}

\textbf{Step 6: most primes of $P$ are good.} \emph{Let $(m_0,a_0,U)$ be admissible, let $1\le F\le U^{1/2}$
with $\log(FU)\le Y^{\varepsilon/4}$, and let $\Xi$ be a set of at most $R$ primitive characters of conductor
at most $F$. Then $P$ is not empty, and for all $q\in P$ outside a set of at most $\frac12|P|$ elements (for
$Y$ large in terms of $R$), every $\xi\in\Xi$ has the following properties. The character $\chi=\xi\chi_q$
is primitive and nontrivial of conductor at most $FU$; $L(s,\chi)\ne0$ for $\re s>1-\frac1{400}$,
$|\mathrm{Im}\,s|\le FU$; and $|G_\chi(y)|\le y^{1/2+\tau}$ for $Y\le y\le Y^A$.}

\noindent\emph{Proof.} Since $q>U^{1/2}\ge F$, the prime $q$ does not divide the conductor of $\xi$, so
$\xi\chi_q$ is primitive, and $q$ is determined by it. Hence by (Q3) the number of $q\in P$ for which some
primitive character of the form $\xi'\chi_q$, with $\xi'$ of conductor at most $F$, has a zero in the box of
height $FU$ is at most $C(FU)^{1/4}\le CU^{3/8}$, which is at most $\frac14U^{9/10}\le\frac14|P|$. By Step~3
with $\kappa=\xi$, used for each $\xi\in\Xi$, the number of $q$ with $|S(y,q,\xi)|>\frac12y^{1/2+\tau}$ for
some $y$ and some $\xi$ is at most $RC_\tau Y^{-\tau}|P|\le\frac14|P|$. For the remaining $q$,
\eqref{III:split} and consequence (a) of Step~4 with $Q=FU$ give
$|G_\chi(y)|\le\frac12y^{1/2+\tau}+Y^\varepsilon\le y^{1/2+\tau}$. \qed

\textbf{Step 7: the choice of $q$ and $d$.}

\emph{The prime $q$.} Take $m_0=8$, $a_0=3$, $U=e^z$. This is admissible. Take $F=8M$ and for $\Xi$ all
primitive characters of conductor dividing $8M$, the trivial one included. Step~6 gives a prime
$q\equiv3\bmod8$ with $\frac z2<\log q\le z$ such that every $\xi\chi_q$ has the three properties, with box
height $8Me^z$.

\emph{The prime $d$.} Now take $m_0=16\,q\prod_{3\le p\le z}p$ and let $a_0$ be the class with
$a_0\equiv7\bmod16$ and $a_0\equiv-1$ modulo $q$ and modulo every odd prime $p\le z$. Take $U=e^{z'}$. Then
$\log m_0\le4z=4z'/(\log Y)^2$, and the prime factors of $m_0$ are at most $z\le Y^{1/4}$, or equal to $q$,
which exceeds $Z^2$. So $(m_0,a_0,U)$ is admissible. Take $F=8Me^z\le U^{1/2}$ and
$\Xi=\{\xi\chi_q^a\}$ with $\xi$ as before and $a\in\{0,1\}$. Step~6 gives a prime $d\equiv7\bmod16$ with
$\frac{z'}2<\log d\le z'$, with $d\equiv-1$ modulo $q$ and modulo every odd prime up to $z$, such that every
$\xi\chi_q^a\chi_d$ has the three properties, with box height $8Me^{z+z'}$.

\emph{Verification.} Parts (a)--(c) hold by construction; $q,d>8M$. For (d): a character $\xi$ modulo $M$ is
induced by a primitive character of conductor dividing $M$, so the primitive character inducing
$\xi\chi_q^i\chi_d^j$ is one of those treated, and $|G_\chi(y)|\le y^{1/2+\varepsilon/2}$ on $[Y,Y^A]$. The
bound $|\frac{L'}L(1,\chi)|\le4\log Y\le4\log y$ is consequence (b) of Step~4, with
$\log Q\le\log(8M)+z+z'\le3z'$. Part (e) is the zero-free box. \qed

\emph{Remarks.} (1) \emph{There is no circularity.} Nothing of the strength of the Riemann hypothesis is
assumed for the characters $\chi$. The bound $y^{1/2+\tau}$ is the mean-square size of a sum of fixed length
$Z=Y^{400A}$ over a family of primes $q$ with $\log q\ge\frac12Y^{\varepsilon/6}$, far beyond that length. As
$q$ varies, the symbols $(\frac nq)$, $n\le Z$, behave like independent signs; reciprocity turns the
statement into one about primes $q$ in progressions to moduli $m_0s$ much smaller than $q$, where (Q2)
applies. The zero-free box is a weak property of the same kind, and it is proved for all but few $q$, by
(Q3).

(2) Part (e) is proved for every primitive $\rho$ of conductor dividing $8M$, real or not.

(3) The mean value \eqref{III:ms} is obtained without a large sieve inequality. A reader who knows
Heath-Brown's large sieve for real characters \cite{HB95} will recognise the statement; the route above
uses primes in progressions.

\section{The large scale: proof of Estimate C}

This section proves Estimate C from three standard theorems on Dirichlet $L$-functions, which we cite, and
from part (e) of Theorem~\ref{thm:pair}. No zero-density estimate and no mean value over characters is
used here. In the Lean development this is the theorem \texttt{QH.thinH} with the lemmas it calls.

\textbf{Notation.} For a primitive character $\chi$ modulo $N$ (the case $N=1$, $\chi=1$ included) put
\begin{equation}\label{D:130}
 \psi_0(u,\chi)=\sum_{n\le u}\Lambda(n)\chi(n)-\begin{cases}u&\chi=1\\0&\chi\ne1.\end{cases}
\end{equation}
Let $G_\chi$ be as in \S III.1 for $\chi\ne1$, and $G_1=G$. Zeros $\rho=\beta+i\gamma$ are counted with
multiplicity.

\begin{theorem}[Estimate C; \texttt{QH.thinH}]\label{thm:C}
Let $B\ge1$, $0<\varepsilon<\frac1{52}$ and $c'>0$. There is $x_0=x_0(B,\varepsilon,c')$ such that the following holds for
$x\ge x_0$. Let $q\ne d$ be odd primes and $Q$ a positive integer prime to $qd$, divisible by $8$, with odd
part at most $B$, and with $\log(Qqd)\le2x^{\varepsilon/5}(\log x)^2$. Suppose that for every real primitive
character $\chi\ne1$ whose conductor divides $Qqd$ and is divisible by $q$ or by $d$, $L(\beta,\chi)\ne0$ for
real $\beta>1-\frac1{400}$. Then for every primitive character $\chi$ whose conductor divides $Qqd$,
\begin{equation}\label{D:131}
 |G_\chi(y)|\le y\exp\!\big(-c'x^{\varepsilon/8}\big)\qquad\text{for all }y\ge\tfrac13\exp\!\big(x^{\varepsilon/3}\big).
\end{equation}
\end{theorem}

\noindent In the application $Q$ is the power of $2$ equal to $4n$ (for $\zeta$), the characters are those of
the field $E$, and $q$, $d$ are the primes of Theorem~\ref{thm:pair} with $Y\le x\le Y^{3/2}$; then
$\log(Qqd)\le2x^{\varepsilon/5}(\log x)^2$ by parts (a), (b) of that theorem, and the hypothesis on real zeros
is its part (e) with $M=1$ (see Step~6; for a general $Q$ one takes $M$ divisible by the odd part of $Q$).
In Part I and in formula \eqref{II:t5.6} of Part II the constant $c'$ is called $\kappa$.

\textbf{The three theorems cited.} Proofs are in Davenport \cite[Ch.~14, 16, 19]{Dav00} and Montgomery and
Vaughan \cite{MV07}: Theorem 11.3 for (Z1), Theorem 10.17 for (Z2), and Theorems 12.5 and 12.10 with the
discussion after Corollary 12.11 for (Z3), which is stated there in a slightly different form. The Lean development proves all three.
\begin{enumerate}[label=\textup{(Z\arabic*)},leftmargin=2.6em]
\item \emph{Zero-free region} (de la Vall\'ee Poussin, Gronwall, Landau, Titchmarsh). There is an absolute
constant $c_1>0$ such that, for every $N\ge1$ and every character $\chi$ modulo $N$, if $L(\rho,\chi)=0$ and
$\beta\ge1-c_1/\log(N(|\gamma|+2))$, then $\chi$ is real and nonprincipal and $\rho$ is real. (For $N=1$ this
says that $\zeta$ has no zero in the region.)
\item \emph{Zeros in a window.} There is an absolute $C_1$ such that for every primitive $\chi$ modulo $N$ and
every real $t$, the number of zeros with $0<\beta<1$ and $|\gamma-t|\le1$ is at most $C_1\log(N(|t|+2))$.
\item \emph{Truncated explicit formula.} There is an absolute $C_2$ such that for every primitive $\chi$
modulo $N$ and $2\le T\le u$,
\begin{equation}\label{D:132}
 \Big|\psi_0(u,\chi)+\sum_{|\gamma|<T}\frac{u^\rho}{\rho}\Big|
 \le C_2\Big(\frac{u\log^2(NuT)}{T}+\log(Nu)+\sum_{|\gamma|<1}\frac1{|\rho|}\Big),
\end{equation}
the sums being over zeros with $0<\beta<1$.
\end{enumerate}
We also use the functional equation, in the form: if $\chi$ is primitive and $0<\beta<1$, then $1-\rho$ is a
zero of $L(s,\overline\chi)$; and Dirichlet's theorem $L(1,\chi)\ne0$.

\textbf{Step 1: an integral for $G_\chi$.} \emph{For primitive $\chi$ and $y\ge1$,}
\begin{equation}\label{III:int}
 G_\chi(y)=-y\int_y^\infty\psi_0(u,\chi)\,\frac{du}{u^2}.
\end{equation}
\emph{Proof.} Let $\delta=1$ if $\chi=1$ and $\delta=0$ otherwise, so that
$\psi_0(u,\chi)=\psi(u,\chi)-\delta u$ with $\psi(u,\chi)=\sum_{n\le u}\Lambda(n)\chi(n)$. For real $s>1$,
integrating the Dirichlet series term by term,
\begin{equation}\label{D:134}
 H(s)=\int_1^\infty\psi_0(u,\chi)\,u^{-s-1}du=-\frac1s\frac{L'}{L}(s,\chi)-\frac{\delta}{s-1}.
\end{equation}
By Step~4 below (which does not use the present step), $|\psi_0(u,\chi)|u^{-2}$ is integrable; for this one
needs, for a real $\chi$, only some interval $(1-c_\chi,1)$ free of zeros, which exists because
$L(1,\chi)\ne0$. So we may let $s\to1^+$ under the integral. This gives $H(1)=-\frac{L'}L(1,\chi)$ for
$\chi\ne1$, and $H(1)=-1-\gamma$ for $\chi=1$, from the expansion
$-\frac{\zeta'}\zeta(s)=\frac1{s-1}-\gamma+O(s-1)$. On the other hand, by partial summation,
\begin{equation}\label{D:135}
 y\int_1^y\psi(u,\chi)\frac{du}{u^2}=\sum_{n\le y}\Lambda(n)\chi(n)\Big(\frac yn-1\Big).
\end{equation}
Hence $y\int_1^y\psi_0u^{-2}du$ equals this sum minus $\delta y\log y$. Subtracting $yH(1)$ gives
$-y\int_y^\infty\psi_0u^{-2}du=\sum_{n\le y}\Lambda(n)\chi(n)(y/n-1)-\delta y\log y-yH(1)$, which is
$G_\chi(y)$ in both cases. This is also where the constant $1+\gamma$ in the definition of $G$ comes
from. \qed

\textbf{Step 2: the zero-free region, as used.} Let $c_0>0$ and $c=\min(c_1,c_0)$. \emph{Let $\chi$ be a
character modulo $N$ and $\lambda\ge\max(1,\log N)$. If $\chi$ is real and nonprincipal, suppose
$L(\beta,\chi)\ne0$ for $1-c_0<\beta<1$. Then every zero with $\beta<1$ satisfies}
\begin{equation}\label{III:zfr}
 \beta\le1-\frac{c}{\lambda+\log(|\gamma|+3)} .
\end{equation}
\emph{Proof.} If not, $\beta>1-c_1/\log(N(|\gamma|+2))$, so by (Z1) $\chi$ is real and nonprincipal and
$\rho=\beta$ is real. But then $\beta>1-c/(\lambda+\log3)\ge1-c_0$, since $c\le c_0$ and $\lambda\ge1$; this is against the assumption. \qed

\emph{Zeros are not close to $0$.} If moreover $\chi$ is primitive, a zero with $0<\beta<1$ and $|\gamma|\le2$
has $\beta\ge c/(3\lambda)$. Indeed $1-\rho$ is a zero of $L(s,\overline\chi)$, and \eqref{III:zfr} applies to
$\overline\chi$ (if $\overline\chi$ is real it equals $\chi$); so $1-\beta\le1-c/(\lambda+\log5)$, and
$\log5\le2\lambda$.

\textbf{Step 3: the sum of $1/|\rho|$.} \emph{Under the assumptions of Step 2 with $\chi$ primitive, for
$T\ge1$,}
\begin{equation}\label{D:137}
 \sum_{|\gamma|<T}\frac1{|\rho|}\le C_3(\lambda+\log T)^2,\qquad C_3=6C_1\Big(\frac9c+4\Big).
\end{equation}
\emph{Proof.} Group the zeros by $i=\lfloor|\gamma|\rfloor$, $0\le i\le T$. By (Z2), applied to the two windows centred at $\pm i$, each group has at most
$2C_1\log(N(i+2))\le6C_1(\lambda+\log T)$ zeros. For $i\le1$ (so $|\gamma|<2$) use $1/|\rho|\le1/\beta\le3\lambda/c$; for
$i\ge2$ use $1/|\rho|\le1/i$. Since $\sum_{i\le T}1/i\le1+\log T$, the sum is at most
$6C_1(\lambda+\log T)\big(6\lambda/c+1+\log T\big)$. \qed

\textbf{Step 4: the bound for $\psi_0$.} \emph{Under the same assumptions there are $C_4$ and $c_2>0$,
depending only on $c_0$, such that for $u\ge e$, with $v=\log u$,}
\begin{equation}\label{III:psi}
 |\psi_0(u,\chi)|\le C_4\,u\,(v+\lambda)^2\exp\Big(-\frac{c_2\,v}{\lambda+\sqrt v}\Big).
\end{equation}
\emph{Proof.} Take $T=\exp(\sqrt v)$ in (Z3); then $2\le T\le u$. Put $c_2=\min(c/3,1)$ and
$\mathcal E=\exp(-c_2v/(\lambda+\sqrt v))$. For $|\gamma|<T$ we have $\log(|\gamma|+3)\le\sqrt v+2$, so by
\eqref{III:zfr} $\beta\le1-c/(\lambda+\sqrt v+2)\le1-c_2/(\lambda+\sqrt v)$, and $|u^\rho|\le u\mathcal E$. By
Step~3, $\sum_{|\gamma|<T}|u^\rho/\rho|\le C_3(\lambda+\sqrt v)^2u\mathcal E$. The error terms of (Z3):
$1/T=e^{-\sqrt v}\le\mathcal E$ because $c_2v/(\lambda+\sqrt v)\le\sqrt v$, and
$\log(NuT)\le2(v+\lambda)$, so the first is at most $4(v+\lambda)^2u\mathcal E$. Since $u\mathcal E\ge1$, the
second is at most $(v+\lambda)^2u\mathcal E$, and the third at most $C_3\lambda^2u\mathcal E$. \qed

\textbf{Step 5: integration.} \emph{Let $A>0$ and $c'>0$. For large $x$, if
$y\ge\frac13\exp(x^{\varepsilon/3})$, $1\le\lambda\le Ax^{\varepsilon/5}(\log x)^2$, and \eqref{III:psi} holds
for $u\ge y$, then $y\int_y^\infty|\psi_0(u,\chi)|u^{-2}du\le y\exp(-c'x^{\varepsilon/8})$.}

\noindent\emph{Proof.} Substitute $u=e^v$. The integral is at most
$C_4\int_V^\infty(v+\lambda)^2\exp(-c_2v/(\lambda+\sqrt v))\,dv$ with $V=\log y\ge\frac12x^{\varepsilon/3}$. For
$v\ge\frac12x^{\varepsilon/3}$ and large $x$ we have $\lambda\le v^{61/100}$, because
$v^{61/100}\ge2^{-1}x^{\varepsilon/5}x^{\varepsilon/300}$ and $2A(\log x)^2\le x^{\varepsilon/300}$. Also
$\sqrt v\le v^{61/100}$ and $\lambda\le v$. Hence the exponent is at least $\frac{c_2}2v^{39/100}$, and
$(v+\lambda)^2\le4v^2$. Now $v^{39/100}\ge\frac12x^{13\varepsilon/100}$, so
\begin{equation}\label{D:139}
 \exp\big(-\tfrac{c_2}2v^{39/100}\big)\le\exp\big(-\tfrac{c_2}8x^{13\varepsilon/100}\big)\exp\big(-\tfrac{c_2}4v^{39/100}\big),
\end{equation}
and $\int_0^\infty v^2\exp(-\frac{c_2}4v^{39/100})dv$ is a finite constant. So the integral is at most a
constant times $\exp(-\frac{c_2}8x^{13\varepsilon/100})$. Since $\frac{13}{100}>\frac18$, this is at most
$\exp(-c'x^{\varepsilon/8})$ for large $x$. \qed

\emph{Where the exponents come from.} The classical saving is
$\exp(-c\log y/(\log N+\sqrt{\log y}))$. Here $\log N$ is about $x^{\varepsilon/5}$ and $\log y$ about
$x^{\varepsilon/3}$, so the saving is $\exp(-x^{\varepsilon/3-\varepsilon/5-o(1)})$, and
$\frac13-\frac15=\frac2{15}>\frac18$.

\textbf{Step 6: exceptional zeros.} It remains to produce $c_0>0$, depending only on $B$, such that every
real primitive $\chi\ne1$ with conductor $N\mid Qqd$ has no zero in $(1-c_0,1)$.

\emph{(i) $N$ divisible by $q$ or $d$.} This is the hypothesis of the theorem, with $\frac1{400}$.

\emph{(ii) $N$ prime to $qd$.} Then $N\mid Q$. The conductor of a real primitive character is not divisible
by $16$ or by the square of an odd prime. (If $p^2\mid N$ with $p$ odd, each $u\equiv1\bmod N/p$ is a square
modulo $N$: $1+\frac Npt\equiv\big(1+\frac Npt\frac{p+1}2\big)^2$. So $\chi(u)=1$ and $\chi$ is defined modulo
$N/p$. Similarly $1+\frac N2t\equiv(1+\frac N4t)^2$ when $16\mid N$.) So $N$ divides $8$ times the odd part
of $Q$, and $N\le8B$. There are finitely many characters of conductor at most $8B$. For each, $L(1,\chi)\ne0$
and $L(s,\chi)$ is continuous, so $L(\beta,\chi)\ne0$ on some interval $(1-c_\chi,1]$. Let $c_0$ be the
smallest of the $c_\chi$ and of $\frac1{400}$. (For $\zeta$ only the real characters of conductor $4$ and $8$ occur. The constant can be made explicit
in $B$ by (Q1) of \S III.1: a real zero $\beta$ of a character of conductor $N\le8B$ has
$1-\beta\ge\min(c_1/\log(3N),\,c/N)$.)

\emph{How part (e) of Theorem~\ref{thm:pair} gives (i).} A real primitive character $\chi$ of conductor
$N\mid Qqd$ divisible by $q$ or $d$ factors as $\rho\,\chi_q^i\chi_d^j$ with $(i,j)\ne(0,0)$ and $\rho$ real
primitive of conductor dividing $8$ times the odd part of $Q$. Part (e), on the real axis, says
$L(\beta,\chi)\ne0$ for $\beta>1-\frac1{400}$.

\textbf{Conclusion.} Let $\chi$ be primitive of conductor $N\mid Qqd$, and
$\lambda=\max(1,\log N)\le3x^{\varepsilon/5}(\log x)^2$. By Step~6 the assumption of Step~2 holds with $c_0$.
So \eqref{III:psi} holds for all $u\ge e$, and Steps~1 and~5 give the theorem. \qed

\emph{Remark.} The Lean proof obtains the limits in Step~1, and the convergence of the integral, from
the Siegel--Walfisz theorem. The argument above uses only \eqref{III:psi} for the one character in
question.

\section{The sets of columns $T_2$ and $T$}

The table of \S I.5 has $n$ columns, indexed by the group $H$ of \S I.5. It is cyclic of order $n$; fix an isomorphism with $A=\Z/n\Z$ and write it
additively, so that the translates $uT_2$ of Part II become $u_0+T_2$. The construction singles out
two sets of columns. Both are needed only through the combinatorial properties proved here; how (M)
uses them is in \S\ref{as:sec} ((H0)(h) and (H7) there). Throughout, $x$ is large, $0<\varepsilon<\frac1{52}$, and
$x^{2+8\varepsilon}\le n\le2x^{2+8\varepsilon}$.

\subsection*{The set $T_2$: a construction of Behrend}

\begin{theorem}[\texttt{hostThickExists} in {[27]}; \texttt{hostThickOrbit}, \texttt{thickFrame2}]\label{thm:T2}
Let $L\ge2$ be an integer with $L\le C_L\log x$. For $x$ large in terms of $\varepsilon$ and $C_L$ there is
a set $T_2$ of positive integers with the following properties.
\begin{enumerate}[label=\textup{(\alph*)},leftmargin=2.2em]
\item $4Ls<n$ for every $s\in T_2$, and $|T_2|\ge x^{-\varepsilon}n$.
\item \emph{No relation.} If $s_1,s_2,s_3\in T_2$ are distinct and $i,j$ are integers with $0<|i|,|j|\le L$,
$i\ne j$, then $j(s_2-s_1)\ne i(s_3-s_1)$.
\item \emph{Arithmetic progressions meet translates of $T_2$ rarely.} Let $u_0,t,h\in\Z/n\Z$. If the additive
order of $h$ exceeds $4L$, then among any $L$ consecutive terms $t+ih$ at most two lie in
$u_0+T_2$. If the order of $h$ is at most $4L$, at most one element of $\Z/n\Z$ on the progression lies in $u_0+T_2$.
\end{enumerate}
\end{theorem}

\noindent For $i=1$, $j=2$ property (b) says that $T_2$ contains no three-term arithmetic progression. The
proof is Behrend's \cite{Beh46}: points on a sphere, written in a large base.

\emph{Proof.} \emph{A sphere.} Let $S,D\ge2$ be integers. The $D^S$ vectors $v\in\{0,\dots,D-1\}^S$ have
$|v|^2<SD^2$, so some sphere $|v|^2=k$ contains at least $D^{S-2}/S$ of them. Call this set $\Sigma$. Three
points of a sphere are never collinear. More precisely: \emph{if $v_1,v_2,v_3\in\Sigma$ and
$j(v_2-v_1)=i(v_3-v_1)$ with $i,j\ne0$, $i\ne j$, then $v_2=v_1$.} Indeed put $b=v_2-v_1$, $c=v_3-v_1$. Since
$|v_1+b|^2=|v_1|^2$ we have $2v_1\cdot b+|b|^2=0$, and likewise for $c$. Substitute $b=\frac ijc$ in the first relation and multiply by $j^2/i$: this gives
$2j\,v_1\cdot c+i|c|^2=0$. Subtract $j$ times the second relation: $(i-j)|c|^2=0$. So $c=0$ and then $b=0$.

\emph{Base $B=4LD$.} Map $v\mapsto s(v)=\sum_mv_mB^m$, and let $T$ be the image of $\Sigma$. Suppose
$j(s(v_2)-s(v_1))=i(s(v_3)-s(v_1))$ with $i,j$ as in (b). The digits
$c_m=j(v_{2,m}-v_{1,m})-i(v_{3,m}-v_{1,m})$ satisfy $|c_m|\le2L(D-1)<B$ and $\sum_mc_mB^m=0$; hence every
$c_m=0$ (look at the lowest nonzero digit modulo $B$). By the statement about the sphere, $v_2=v_1$. This
is (b) for $T$.

\emph{Parameters.} Put $u=\sqrt{\log n\,\log L}$, $S=\lfloor u/\log L\rfloor$, $a=(\log n-\log L-3)/S$ and
$D=\lfloor e^a/(4L)\rfloor$. For large $x$ one has $S\ge32$, $a\ge7+\log L$, hence $D\ge e^a/(8L)\ge2$, and $a\le2u$.
Then $B^S\le e^{Sa}<n/(4L)$, so $4Ls<n$ on $T$, and
\begin{equation}\label{D:140}
 |T|\ge\frac{D^{S-2}}S\ge\frac nL\exp\big(-32\sqrt{\log n\,\log L}\big).
\end{equation}
Since $\log L\le\log\log x+O(1)$, the right side exceeds $4x^{-\varepsilon}n$ for large $x$. Remove $0$ from
$T$ to get $T_2$. This proves (a) and (b).

\emph{Property (c).} Let the order of $h$ exceed $4L$, and suppose $t+ah$, $t+bh$, $t+ch$ lie in $u_0+T_2$
with $a<b<c<a+L$: they equal $u_0+s_a$, $u_0+s_b$, $u_0+s_c$. The $s$'s are distinct because $(b-a)h$,
$(c-a)h$, $(c-b)h\ne0$. Eliminating $h$,
\begin{equation}\label{D:141}
 X=(c-a)(s_b-s_a)-(b-a)(s_c-s_a)\equiv0\bmod n .
\end{equation}
But $|X|<n$, because $4Ls<n$ for each $s$. So $X=0$, against (b). If the order $m$ of $h$ is at most $4L$
and two points of the progression are $u_0+s_1$, $u_0+s_2$, then $n\mid m(s_1-s_2)$ and $|m(s_1-s_2)|<n$, so
$s_1=s_2$. \qed

\begin{altbox}
\noindent\emph{Remark (a fixed dimension suffices).} The choice of $S$ above, which grows with $x$, gives
the lower bound $L^{-1}\exp(-32\sqrt{\log n\,\log L})$ for the density before $0$ is removed, far more than the theorem states. For the density
$x^{-\varepsilon}$ of (a) one may fix the dimension once: let $S$ be an integer with $S>12/\varepsilon$, and
put $D=\lfloor(n/(4L))^{1/S}/(4L)\rfloor$. Then $B^S=(4LD)^S\le n/(4L)$, so $4Ls<n$ on $T$ as before. For
large $x$ we have $D\ge(n/(4L))^{1/S}/(8L)\ge2$, and
\begin{equation}\label{D:142}
 |T|\ \ge\ \frac{D^{S-2}}S\ \ge\ n\cdot n^{-2/S}\cdot\frac{(8L)^{-S}}S .
\end{equation}
Here $n\le2x^{2+8\varepsilon}\le x^3$, so $n^{-2/S}\ge x^{-6/S}\ge x^{-\varepsilon/2}$; and
$(8L)^{-S}/S\ge(8C_L\log x)^{-S}/S\ge4x^{-\varepsilon/2}$ for large $x$, because $S$ does not depend on $x$.
So $|T|\ge4x^{-\varepsilon}n$, and the rest of the proof is unchanged. The proof above follows the source,
which optimises $S$. (This remark is due to a reviewer.)
\end{altbox}

\subsection*{The set $T$: a random choice}

Let $\delta_1=x^{-\varepsilon/8}$. The set $T\subset A$ is chosen so that $|T|$ is about $\delta_1n$ and so that,
for most primes $p$ and most $t\in T$, the points of $T$ on the orbit of $t$ under the Frobenius of $p$ are
evenly spread. The precise meaning is this.

\emph{Good orbits.} Let $n'\ge1$ and $W\subset\{0,\dots,n'-1\}$, read cyclically. Fix integers $L_2,K,r\ge1$,
and put $u=\lceil1/\delta_1\rceil$, $\ell=(6r+1)u$,
$M_0=\lceil(5+8\varepsilon)\,10^{6r+1}\log x\rceil$, $G'=(M_0+5)\ell$ and
$N_1=\lceil\max(6eG'\delta_1,(5+8\varepsilon)\log_2x)\rceil+1$. Cut $\{0,\dots,n'-1\}$ into $\lfloor n'/\ell\rfloor$ consecutive
\emph{tiles} of length $\ell$ (a remainder shorter than $\ell$ is ignored, and tiles are numbered cyclically), and each tile into $6r+1$ \emph{cells} of length $u$, numbered $1,\dots,6r+1$. A tile is \emph{marked}
if its odd-numbered cells contain no point of $W$ and each even-numbered cell contains exactly one point of
$W$, at an even position. Call $W$ \emph{good} if
\begin{enumerate}[label=\textup{(G\arabic*)},leftmargin=2.8em]
\item every cyclic window of length $L_2$ contains at most $K$ points of $W$;
\item among any $M_0+1$ consecutive tiles (cyclically) at least one is marked;
\item every cyclic window of length $G'$ contains at most $N_1$ points of $W$.
\end{enumerate}

\begin{theorem}[\texttt{topsGroup}, \texttt{topsChoiceH}]\label{thm:T}
Let $L_b>0$, $K=\lceil8(4+8\varepsilon)/\varepsilon\rceil$ and $r=K+1$.
There are $C_b$ and $x_0$ such that for $x\ge x_0$ the following holds. Let $L_2\le L_b$. Let $\mathcal P$ be
a finite set of primes, and for $p\in\mathcal P$ let $h_p\in A$ have order $n_p\ge x^{\varepsilon/4}$. For each
$p$ and each coset of $\langle h_p\rangle$ fix a base point $b$ on it; for $t\in A$ let
$W_p(t,T)=\{0\le i<n_p:\ b+ih_p\in T\}$, $b$ being the base point of the coset of $t$. Then there is
$T\subset A$ with
\begin{equation}\label{D:143}
 \tfrac12\delta_1n\le|T|\le2\delta_1n\qquad\text{and}\qquad
 \sum_{p\in\mathcal P}\frac{\log p}p\,\#\{t\in T:\ W_p(t,T)\text{ is not good}\}
 \le4C_bx^{-2}\,\delta_1n\sum_{p\in\mathcal P}\frac{\log p}p .
\end{equation}
\end{theorem}

\noindent In the application $\mathcal P$ consists of the inert primes below $\exp(x^{\varepsilon/3})$ whose
Frobenius has order at least $x^{\varepsilon/4}$ in $A$, and $L_2$ is bounded by a constant depending only on $\varepsilon$ (the source gives $2(620\,C_o/\varepsilon+1)$ for primes $p>x^{\varepsilon/(5C_o)}$, with a constant $C_o\ge2$ that is defined in the source). That $L_2$ is
bounded uses Remark~(3) of \S III.1: every inert prime exceeds $Y^{\varepsilon/5}$.

\emph{Proof.} Choose $T$ at random: each element of $A$ belongs to $T$ independently with probability
$\delta_1$. (For Dirichlet characters other than the trivial one the source uses two different
probabilities, $\delta_1$ and $\frac34\delta_1$, on different cosets. For $\zeta$ there is one coset.)

\emph{Size.} $|T|$ has mean $\delta_1n$ and variance at most $\delta_1n$.

\emph{One orbit.} Fix $p$ and $t$, write $n'=n_p$, and condition on $t\in T$. The other points of the orbit
remain independent. We bound the probability that $W=W_p(t,T)$ is not good.

(G1) fails only if some window of length $L_2$ contains $K$ points other than $t$. For one window this has
probability at most $\binom{L_2}K\delta_1^K\le L_b^Kx^{-K\varepsilon/8}\le L_b^Kx^{-(4+8\varepsilon)}$. There
are $n'$ windows.

(G2). A given tile not containing $t$ is marked with probability at least
\begin{equation}\label{D:144}
 \big((1-\delta_1)^u\big)^{3r+1}\Big(\big\lfloor\tfrac u2\big\rfloor\delta_1(1-\delta_1)^{u-1}\Big)^{3r}\ge10^{-(6r+1)},
\end{equation}
since $u\delta_1$ is between $1$ and $1.01$. Among $M_0+1$ consecutive tiles at most one contains $t$;
the other $M_0$ are disjoint, hence independent. The probability that none is marked is at most
$(1-10^{-(6r+1)})^{M_0}\le x^{-(5+8\varepsilon)}$. There are at most $n'$ starting tiles. (Here and for (G1), (G3)
one needs $L_2\le n'$ and $G'=(M_0+5)\ell\le n'$, so that windows and runs of tiles do not wrap around; this holds for large $x$ because
$n'\ge x^{\varepsilon/4}$, $L_2\le L_b$ and $\ell\le2(6r+1)x^{\varepsilon/8}$.)

(G3) fails only if some window of length $G'$ contains $N_1$ points other than $t$; then it contains
$j=N_1-1$ such points. Since
$j\ge6eG'\delta_1$, the probability of the latter is at most $\binom{G'}{j}\delta_1^j\le(eG'\delta_1/j)^j\le6^{-j}$, which is at
most $x^{-(5+8\varepsilon)}$. There are $n'$ windows.

Since $n'\le n\le2x^{2+8\varepsilon}$, the three bounds add up to at most $C_bx^{-2}$. So the expected number
of $t\in T$ with $W_p(t,T)$ not good is at most $C_bx^{-2}\delta_1n$, for every $p$.

\emph{Choice.} We may assume $\mathcal P$ is not empty. Let $V=(|T|-\delta_1n)^2/(2\delta_1n)$ and let $\mathcal B$ be the left side of the inequality
of the theorem divided by four times its expected upper bound $C_bx^{-2}\delta_1n\sum_p\frac{\log p}p$. Then
$\mathbb E(V)\le\frac12$ and $\mathbb E(\mathcal B)\le\frac14$. So some $T$ has $V+\mathcal B\le1$. For it
$\mathcal B\le1$, which is the second assertion, and $\big||T|-\delta_1n\big|\le\sqrt{2\delta_1n}\le\frac12\delta_1n$,
which gives the first. \qed

\section{The size of the theta kernels}

This section concerns the upper bound for the size of the maps: inequality \eqref{II:t6.5} of Part II, first line.
The following are proved below: the one-variable model; the coefficient estimate in
several variables and in several blocks, with the period; the bound for the polynomial on the unit sphere;
the bound for one map and for a chain of maps; the bound for a theta function by its mean square; and the
hook product. Four passages are conditional: the reference mean square; the eigenvalues of a lift; the
comparison, row by row, of the kernel degrees with the hook product; and the norm of the pairing at the end
of the chain. The last two are deductions from inputs of the source that are named where they are used.
The passage from these estimates to \eqref{II:t6.5} is in \S\ref{sa:sec}: Theorem~\ref{sa:thm} deduces the first line of
\eqref{II:t6.5} from eleven named inputs. The inputs are the arithmetic statements listed at the end of this section;
none of them is proved in these pages. So the first line of \eqref{II:t6.5} is reduced to quoted statements, not
proved.
We begin with the one-variable model. It shows where the two features of the bound come from: the size
$\sqrt{n!}$ and the sign of the term with the period.

\textbf{The model.} \emph{Let $F$ be an entire function with $|F(z)|\le S\,e^{|z|^2/2}$ for all $z$, where $S>0$. Then for
$n\ge1$}
\begin{equation}\label{D:145}
 |F^{(n)}(0)|\le S\,n!\,\Big(\frac en\Big)^{n/2},
\end{equation}
so that, if $F^{(n)}(0)\ne0$,
\begin{equation}\label{D:146}
 \log|F^{(n)}(0)|^2\le2\log S+n\log n-n+O(\log(n+1)).
\end{equation}
\emph{Proof.} By Cauchy's formula on the circle $|z|=r$, $|F^{(n)}(0)|\le n!\,S\,e^{r^2/2}r^{-n}$. Take
$r=\sqrt n$. The second form follows from Stirling's formula, $\log n!=n\log n-n+O(\log(n+1))$. \qed

\noindent So a Gaussian bound for the function gives, for the \emph{square} of its $n$th derivative, the size
$n\log n-n$: this is $\log n!$. Compare a function with a pole at finite distance, such as $1/(1-z)$: its $n$th
derivative at $0$ is $n!$, and the logarithm of the square is $2\log n!$. That was the situation of the
lemniscate sine in \S I.3. Theta functions
satisfy such a Gaussian bound. This is point (i) of \S I.4.2.

\emph{The period.} If the coordinate is changed by a factor, $z=\lambda w$ with $|\lambda|^2=1/\Omega$, then
the $n$th derivative with respect to $w$ is $\lambda^n$ times that with respect to $z$, and the bound for the
logarithm of its square becomes
\begin{equation}\label{D:147}
 n\log n-n-n\log\Omega+2\log S+O(\log(n+1)).
\end{equation}
This is the shape $a(\log a-1-\log\Omega)$ of \eqref{II:t6.5}, with its sign. The constant relevant there is
$\Omega_E$, not the elliptic $\Omega$ of \S I.7; the two differ by the factor $\mathrm{rd}_E/\sqrt d$.

The paragraphs that follow carry the model over to the setting of the proof: several complex variables, a
polynomial with values in linear maps in place of a single derivative, several blocks of variables (one for
each place), and sums over classes of lattices. Those up to and including the bound for a theta function by its mean square, and the hook product, are complete as stated; the four passages named above are conditional on inputs from the source. The first of them assumes a Gaussian bound $S$ for the function and a bound for
the polynomial on the unit sphere; both assumptions are then proved in an abstract setting; for the first, the mean square of one
reference theta function is bounded under the inputs (I1)--(I5) listed at the head of that paragraph. That the actual theta kernels
and polynomials of the construction are instances of this setting is taken from the source: the hypotheses
of the majorant lemma with $v=\mathrm{rd}_E$, the form of the constant $S$ with its factors $\beta$ and $c$,
and the identification of the polynomial $P$ of the construction with a substitution of the kind treated
below.

\textbf{Several variables: a coefficient estimate on a sphere.}
Let $F(z)=\sum_\alpha c_\alpha z^\alpha$ be entire on $\mathbb C^d$, $d\ge1$,
with $|F(z)|\le S\exp(\|z\|^2/2)$, where $\|z\|^2=\sum_j|z_j|^2$.
Let $P(\bar z)=\sum_{|\alpha|=n}A_\alpha\bar z^\alpha$ be homogeneous,
with coefficients in a finite-dimensional normed space, for example a space
of linear maps. Here $\alpha$ is a $d$-tuple of nonnegative integers,
$|\alpha|=\sum_j\alpha_j$, $z^\alpha=\prod_j z_j^{\alpha_j}$ and
$\alpha!=\prod_j\alpha_j!$. Suppose, for unit vectors $u$, that the following holds, factors with $\nu_i=0$ being omitted:
\begin{equation}\label{D:148}
 \|P(\bar u)\|\le\prod_i s_i(u)^{\nu_i},\qquad
 s_i(u)\ge0,\quad \sum_i s_i(u)^2=1,\quad
 \nu_i\ge0,\quad\sum_i\nu_i=n.
\end{equation}
Put $h(v)=v\log v-v$ for $v>0$ and $h(0)=0$.
Then
\begin{equation}\label{D:149}
 \left\|\sum_{|\alpha|=n}c_\alpha\alpha! A_\alpha\right\|
 \le S\, e^{1/2}(2n+2)^d
       \exp\left(\tfrac12\sum_i h(\nu_i)\right).
\end{equation}

\emph{Proof.} Let $du$ be probability measure on the unit sphere in
$\mathbb C^d$. Independent rotations of the coordinates show that the mean
of $u^\beta\bar u^\alpha$ vanishes unless $\alpha=\beta$. Its value for
$\alpha=\beta$ is $\alpha!(d-1)!/(n+d-1)!$. Indeed, integrate
$|z^\alpha|^2$ against $\pi^{-d}e^{-\|z\|^2}$ first in the separate complex
coordinates and then in polar coordinates. The two radial moments needed
are $\int_0^\infty t^j e^{-t}\,dt=j!$ for $j=d-1,n+d-1$.
Thus for $r>0$ the sum to be bounded equals
\begin{equation}\label{D:150}
 \frac{(n+d-1)!}{(d-1)!\,r^n}\int F(ru) P(\bar u)\,du.
\end{equation}
Termwise integration is valid by uniform convergence on the sphere.
For $n>0$ the weighted arithmetic--geometric mean inequality, with weights
$\nu_i/n$, gives $\prod_i s_i(u)^{\nu_i}\le
\prod_{\nu_i>0}(\nu_i/n)^{\nu_i/2}$. Take $r=\sqrt n$.
After removing $S\exp(\frac12\sum_i h(\nu_i))$, the remaining factor is
\begin{equation}\label{D:151}
 n!(e/n)^n {n+d-1\choose d-1}
 \le en(n+1)^{d-1}\le e^{1/2}(2n+2)^d .
\end{equation}
The first inequality follows by comparing $\sum_{j<n}\log j$ with
$\int_1^n\log t\,dt$ and counting the weak compositions of $n$ into $d$
parts by their first $d-1$ entries. When $n=0$ the estimate follows at once
from $|F(0)|\le S$ and $\|A_0\|\le1$. \qed

\textbf{Several blocks, and the period.} For several blocks write $z=(z_\tau)_\tau$. Assume $|F(z)|\le S e^{\|z\|^2/2}$,
where $\|z\|^2=\sum_\tau\|z_\tau\|^2$. Take $P$ homogeneous of degree
$n_\tau$ in each block and assume the product of the preceding bounds for
$P$ on the product of unit spheres. Integrate over that product, using radius
$\sqrt{n_\tau}$ in each block of positive degree and setting the variables
of degree-zero blocks of $F$ to zero. This gives the product of the
displayed factors, with $S$ occurring once. Multiplying all coordinates
in block $\tau$ of $F$ by a scalar $\lambda_\tau$ multiplies the extracted
sum by $\prod_\tau\lambda_\tau^{n_\tau}$. Thus its squared norm acquires
the factor $\prod_\tau\Omega_\tau^{-n_\tau}$ if
$|\lambda_\tau|^2=\Omega_\tau^{-1}$.

The second assumption of the estimate, the bound $\|P(\bar u)\|\le\prod_is_i(u)^{\nu_i}$ on the unit sphere,
is proved in the next two paragraphs. The variables of a block form a matrix; the $s_i(u)$ are its singular
values; and $P(\bar u)$ acts on a space of polynomials by substituting a multiple of the matrix of
variables. (Lean: \texttt{fock\_subst\_le} and \texttt{blockBound}.)

\textbf{Where the bound for the polynomial comes from.}
The numbers $s_i$ are the singular values of the matrix of variables.
The integers $\nu_i$ are its prescribed column degrees. We prove the
estimate directly from the two properties defining the polynomials.
Let $\nu_1\ge\cdots\ge\nu_r\ge0$ be integers. Write $W_\nu$ for the
polynomials $f(X)$ on complex $r$ by $r$ matrices which have degree
$\nu_i$ in column $i$ (every monomial has that degree) and are unchanged
when a multiple of an earlier column is added to a later column.
Use the squared Gaussian norm
\begin{equation}\label{D:152}
 \|f\|_G^2=\pi^{-r^2}\int |f(X)|^2 e^{-\sum_{i,j}|X_{ij}|^2}\,dX .
\end{equation}
Here and below integration on a space of complex matrices uses ordinary
Lebesgue measure on the real and imaginary parts of its entries.
The same definition, with one factor $\pi^{-1}$ per entry, applies to
rectangular matrices. If $A$ is $r$ by $b$, with $b\ge r$, let
$s_1\ge\cdots\ge s_r\ge0$ be its singular values. For a $b$ by $r$
variable matrix $Y$ we claim that
\begin{equation}\label{D:153}
       \|f(AY)\|_G\le\Big(\prod_{i=1}^r s_i^{\nu_i}\Big)\|f\|_G,
       \qquad \sum_i s_i^2=\sum_{i,j}|A_{ij}|^2 .
\end{equation}
Factors of exponent zero are taken as one.

First consider a monomial occurring in $f$, and let $d_i$ be its
degree in row $i$. Then
\begin{equation}\label{D:154}
 \sum_{i\le k}d_i\le\sum_{i\le k}\nu_i\quad(1\le k<r),\qquad
 \sum_i d_i=\sum_i\nu_i .
\end{equation}
To see this, multiply the first $k$ rows of $X$ by an indeterminate $t$.
On the dense set where the upper left $k$ by $k$ minor is invertible,
subtract combinations of the first $k$ columns from each later column
to make its first $k$ entries zero. These operations commute with the row
scaling and preserve $f$. Now $t$ occurs only in the first $k$ columns,
so the resulting polynomial in $t$ has degree at most $\sum_{i\le k}\nu_i$.
This holds for all $X$ by polynomial identity. Equality for $k=r$ is
total homogeneity.

Row substitutions by unitary matrices preserve both the Gaussian norm
and the two defining properties of $W_\nu$. Hence the singular value
decomposition of $A$ reduces the claim to diagonal row scaling by the
$s_i$, with the unused rows of $Y$ integrated out. For positive $s_i$,
partial summation of the degree inequalities against the decreasing
numbers $\log s_i$ gives $\sum_i d_i\log s_i\le\sum_i\nu_i\log s_i$.
Thus the multiplier of every monomial is at most $\prod_i s_i^{\nu_i}$.
Distinct monomials are orthogonal in the Gaussian norm by coordinate
rotations. This proves the claim. Zero singular values follow by
continuity, and the sum of their squares is the trace of $AA^*$.

\textbf{Why the norm factors cancel.}
In the application the input degree vector is $\mu$, and the polynomial
is multiplied by $\det(X)^e$ before substitution, for an integer $e\ge0$.
This gives precisely $\nu_i=\mu_i+e$. For a degree vector $\mu$ let
$p_\mu$ be the product of the leading $i$ by $i$ minors raised to
$\mu_i-\mu_{i+1}$, where $\mu_{r+1}=0$. The norm on $W_\mu$ is
$\|f\|_G/\|p_\mu\|_G$. In this norm multiplication by $\det^e$ is
an isometry from $W_\mu$ to its image in $W_\nu$.

Indeed, write a nonsingular Gaussian matrix as $X=QR$ by
Gram--Schmidt on its columns, with $Q$ unitary and $R$ upper triangular
with positive diagonal. The defining column operations give
$f(QR)=f(Q)\prod_i R_{ii}^{\mu_i}$. In computing $\|f\|_G^2$, we may
first average $|f(UX)|^2$ over unitary $U$ with its invariant probability
measure, since the Gaussian measure of $X$ is unchanged. The average
of $|f(UQ)|^2$ is the average of $|f(U)|^2$. Hence the squared Gaussian
norm is the unitary average of $|f(U)|^2$ times a factor depending only
on $\mu$. Singular matrices have measure zero. Dividing by the same
formula for $p_\mu$ cancels the factor. Since $p_\nu=\det^e p_\mu$
and $|\det U|=1$, the normalised norm is unchanged by the shift.
Padding $\nu$ with zeros to length $b$ preserves the norm of $p_\nu$,
because it then depends only on the upper left $r$ by $r$ variables.
Thus the normalised substitution has norm at most $\prod_i s_i^{\nu_i}$.
For independent matrix blocks the maps form a tensor product; its norm
is the product of the individual norms, as follows by singular value
decomposition. On each unit sphere $\sum_i s_i^2=1$. This supplies the
polynomial bound used in the coefficient estimate.

\textbf{Sums over classes, and composing the maps.}
Fix the number field $K$ over which the matrices of the maps are defined.
Let $v$ run over its archimedean places, with one embedding chosen for
each conjugate pair, and let $s$ be the number of these places.
For $0\le j\le m$ let $V_j$ be a finite-dimensional vector space over $K$.
At $v$ its complexification has a norm $\|\cdot\|_{j,v}$.
Consider $K$-linear maps $K_j:V_j\to V_{j+1}$ for $0\le j<m$.
We write $\|K_j\|_v$ for the operator norm in these two norms.

Here is the precise information about one map needed for the chain.
The classes at step $j$ are partitioned into $a_j$ groups, each
containing at most $H_j$ classes. Write $I_j$ for their union and give
each class $p$ a weight $0<w_j(p)\le1$. At each place a vector $x$ is
represented by a tuple $f=(f_p)_{p\in I_j}$ in a complex normed space,
with $\|x\|_{j,v}^2=\sum_p w_j(p)\|f_p\|^2$. Suppose that $K_jx$ is
represented in the corresponding tuple space at step $j+1$ by
\begin{equation}\label{D:155}
    g_q=\sum_{p\in I_j}w_j(p)\epsilon_{pq}T_{pq}f_p
          \quad(q\in I_{j+1}),\qquad
    |\epsilon_{pq}|=1,\qquad \|T_{pq}\|^2\le b_j(v).
\end{equation}
The phases and the block operators may depend on $j$ and $v$.
Here $b_j(v)\ge0$ is a squared bound, as supplied by the preceding
coefficient estimate. Put $D_j(v)=a_j a_{j+1}H_jH_{j+1}b_j(v)$.
Then $\|K_j\|_v^2\le D_j(v)$. Indeed, Cauchy--Schwarz gives
\begin{equation}\label{D:156}
   \|g_q\|^2\le b_j(v)\Big(\sum_p w_j(p)\Big)
                  \sum_p w_j(p)\|f_p\|^2 .
\end{equation}
Multiply by $w_{j+1}(q)$ and sum over $q$. The two sums of weights are
at most $|I_j|\le a_jH_j$ and $|I_{j+1}|\le a_{j+1}H_{j+1}$.
This proves the bound also when the tuples representing vectors form
only a subspace of the space of all tuples.

Let $P=K_{m-1}\cdots K_0$. If $P\ne0$, define its archimedean
height by $h(P)=s^{-1}\sum_v\log\|P\|_v^2$. At each $v$ a nonzero
entry of the matrix of $P$ stays nonzero under the embedding of $K$.
Consequently every $K_j$ has positive norm there, and every factor of
$D_j(v)$ is positive. Successive application of the operator norm gives
\begin{equation}\label{D:157}
   \log\|P\|_v^2\le\sum_{j=0}^{m-1}\log\|K_j\|_v^2
                  \le\sum_{j=0}^{m-1}\log D_j(v).
\end{equation}
Averaging over $v$ proves
\begin{equation}\label{D:158}
 h(P)\le\sum_{j=0}^{m-1}\left\{
  \frac1s\sum_v\log b_j(v)
    +\log H_j+\log H_{j+1}+\log a_j+\log a_{j+1}\right\}.
\end{equation}
Thus the contribution of each complex place can be estimated before
the logarithms are averaged. Applying this result to the arithmetic
maps requires the stated tuple representation with the same norm at
the output of one map and the input of the next. The final period
pairing, when appended to the chain, requires its own norm bound.

\emph{Combining the two.} Take for $b_j(v)$ the square of the block estimate, with the period:
\begin{equation}\label{D:159}
 \log b_j(v)\le2\log S_{j,v}+\sum_\tau\Big(\sum_ih(\nu_{\tau,i})-n_\tau\log\Omega_\tau+2d_\tau\log(2n_\tau+2)+1\Big),
\end{equation}
where $d_\tau$ is the number of variables of block $\tau$.
Then the height of the composite of $m$ maps is at most the sum over $j$ of the averages over $v$ of these
quantities, plus the logarithms of the numbers of classes. Since $\sum_ih(\nu_i)\le h(n)=n\log n-n$ when
$\sum_i\nu_i=n$, this has the shape of the first line of \eqref{II:t6.5}, provided the constants $S_{j,v}$ and the
numbers of classes are under control, and provided the periods $\Omega_\tau$ average to $\Omega_E$.

\textbf{Where the Gaussian bound comes from.} The first assumption above, $|F(z)|\le S\,e^{\|z\|^2/2}$, holds
for theta functions because they are quasi-periodic, and the constant $S$ is controlled by a mean square.
\emph{Let $A$ be a lattice in $\C^g$ with a fundamental domain $D$ of volume $1$, and let $v\ge1$ be such that
$\|x\|^2\ge1/v$ for every nonzero $x\in A$. Let $G$ be entire on $\C^g$, and suppose
$N(z)=|G(z)|^2e^{-\pi\|z\|^2}$ is invariant under translation by $A$. Then for every $z$}
\begin{equation}\label{D:160}
 |G(z)|^2e^{-\pi\|z\|^2}\le3\,(2gv+1)^g\int_DN(w)\,dw .
\end{equation}
\emph{Proof.} Let $r^2=1/(4v)$ and let $B$ be the ball of radius $r$ about $z$. Two translates of $B$ by
different elements of $A$ are disjoint, since their centres are at distance at least $2r$. As $N$ is
periodic, $\int_BN\le\int_DN$. Put $\varphi(u)=G(u)\exp(-\pi\langle u-z,z\rangle)$, where
$\langle u,z\rangle=\sum_ju_j\bar z_j$; it is entire in $u$ and $\varphi(z)=G(z)$. Now
$|\varphi(z)|^2$ is at most the mean of $|\varphi|^2$ over $B$: expand $\varphi$ in its Taylor series about
$z$ and integrate $|\varphi|^2$ over each sphere centred at $z$; the mixed terms integrate to zero and the
constant term contributes $|\varphi(z)|^2$. The volume of $B$ is $\pi^gr^{2g}/g!$. From
$\|u\|^2=\|z\|^2+2\re\langle u-z,z\rangle+\|u-z\|^2$ one gets
$|\varphi(u)|^2e^{-\pi\|z\|^2}=e^{\pi\|u-z\|^2}N(u)\le3N(u)$ on $B$, because $\pi r^2\le\pi/4<1$. Hence
\begin{equation}\label{D:161}
 |G(z)|^2e^{-\pi\|z\|^2}\le3\,g!\Big(\frac{4v}{\pi}\Big)^g\int_BN\le3\,(2gv+1)^g\int_DN,
\end{equation}
using $g!\le g^g$ and $4/\pi<2$. \qed

\noindent After the change of variable $z\mapsto z/\sqrt\pi$ this is the assumption of the estimate, with
$S^2=3(2gv+1)^g\int_DN$. In the source (\texttt{sup\_le\_mean\_square}, \texttt{sup\_translate\_le}) it is
applied to the theta functions attached to a pair of lattices of ranks $N$, $N'$, with $g=2nNN'$ and with
$v$ the root discriminant of $E$; so for one such theta function $\log S^2\le g\log(2g\,\mathrm{rd}_E+1)+\log3+\log\int_DN$.

\emph{The constant in the application is larger.} The kernel whose coefficients are algebraic integers is
not one such theta function but a finite combination of translates of a reference theta function. In the
source the constant that enters the block estimate is
\begin{equation}\label{III:Sform}
 S=\beta\,c^{\,2[E:\Q]NN'}\sqrt{3(2g\,\mathrm{rd}_E+1)^g}\;\Theta_0^{1/2},
\end{equation}
where $c$ is a modulus, $\beta$ bounds the coefficients of the combination, and $\Theta_0$ is the mean
square of the reference theta function at the given embedding. The factor in front of the square root
bounds the sum of the absolute values of the coefficients. The bound for $\Theta_0$, with its arithmetic normalisation, is the subject
of the next paragraph. (All the Lean declarations named in this section are among the declarations the final theorem depends
on, by the dependency computation: \texttt{sup\_le\_mean\_square}, \texttt{sup\_translate\_le},
\texttt{fock\_subst\_le}, \texttt{blockBound}, \texttt{blockBound\_of\_fockJet},
\texttt{gaussianCoeffIntegral}, \texttt{chain\_coupled\_bound}, \texttt{divClausesR\_of}; in the passage
``Row by row'', \texttt{weights\_rows}, \texttt{nu\_eq\_originRow}, \texttt{aPlus\_eq\_originBig},
\texttt{origin\_rewrite}, \texttt{origin\_interlace}, \texttt{originRow}, \texttt{originBig},
\texttt{RowsOriginOff}, \texttt{RowsOriginTop}, \texttt{Sw}; and in the passage ``The pairing at the end of
the chain'', \texttt{brStd}, \texttt{brConst}, \texttt{brAt}, \texttt{BranchMap}, \texttt{sum\_pair\_le}.)

\textbf{The reference mean square $\Theta_0$.} \emph{This paragraph is conditional on the following
inputs, which are not proved here.} The notation is explained below the list.
\begin{enumerate}[label=\textup{(I\arabic*)},leftmargin=2.6em,itemsep=1pt,topsep=3pt]
\item The singular modulus $j$ is an algebraic integer (classical).
\item The conjugates of a curve with complex multiplication by $\mathcal O_K$ run equally often through the
ideal classes of $K$.
\item The class mean period formula: $\langle\log A_\rho\rangle=\log\pi+\gamma+\log(2\pi/\sqrt d)-c_d$. (It follows from Kronecker's limit formula, \cite[Ch.~VIII, \S8]{Wei76}, applied to the decomposition of the zeta function of $\Q(\sqrt{-d})$ over ideal classes.)
\item The tensor-power description of the normalised reference section of ranks $N$, $N'$: its squared
norm is $\Theta_{0,\rho}=H_\rho^{fNN'}$.
\item The reference section at $\rho$ is $\theta/\theta(0)$, for the characteristic whose zero is the
marked half-period, with the metric $e^{-2\pi yv^2}$ in the coordinates $z=u+v\tau$.
\end{enumerate}
\emph{Conclusion.} Under (I1)--(I5),
\begin{equation}\label{D:163}
 \langle\log\Theta_{0,\rho}\rangle\le\frac{fNN'}2\big(\log\sqrt d-\gamma+c_d\big).
\end{equation}
Let $K=\mathbb Q(\sqrt{-d})$ have discriminant $-d$, with $d$ prime and
$d\equiv7\pmod8$ (these are the hypotheses of the application; the calculation below does not use them). Write $\chi_d$ for its quadratic character and
$c_d=L'(1,\chi_d)/L(1,\chi_d)$. Let $M$ be a number field of definition
for the elliptic curve with complex multiplication by $\mathcal O_K$
and its marked point of order two, enlarged if necessary. Average over
all $[M:\mathbb Q]$ embeddings $\rho:M\hookrightarrow\mathbb C$, denoting
this average by brackets. (The chain of maps is defined over a field containing $M$, and
$H_{\bar\rho}=H_\rho$; so an average over the archimedean places of that field, as in the bound for the
composite, is the same as this average.) The periods at $\rho$ form a lattice
$\Lambda=\omega(\mathbb Z+\tau\mathbb Z)$, $y=\Im\tau>0$. The differential
is scaled so that the Weierstrass invariants satisfy
$g_2^3-27g_3^2=1$. Thus $A_\rho=|\omega|^2y$ is the area of this
particular lattice. The marked point and its conjugates determine which
of the three even theta functions to use.

In coordinates $z=u+v\tau$ on the torus, put
\begin{equation}\label{D:164}
 \theta(z)=\sum_{m\in\mathbb Z+\alpha}
   e^{\pi i m^2\tau+2\pi i m(z+\beta)},\qquad
 (\alpha,\beta)\in\{(0,0),(0,\tfrac12),(\tfrac12,0)\}.
\end{equation}
Choose the characteristic whose zero is the marked half-period in these
coordinates. Write $T=|\theta(0)|>0$. The normalisation of the section
at the origin divides it by $\theta(0)$; this is input (I5). In Jacobi coordinates its
squared norm, averaged over the torus with total mass one, is
\begin{equation}\label{D:165}
 H_\rho=\int_0^1\!\!\int_0^1
   \frac{|\theta(u+v\tau)|^2}{T^2}e^{-2\pi yv^2}\,du\,dv
          =\frac{1}{\sqrt{2y}\,T^2}.
\end{equation}
Indeed, integration in $u$ kills the mixed Fourier terms. The remaining
integrals of $e^{-2\pi y(m+v)^2}$ tile the real line. Absolute convergence
on the closed parallelogram justifies termwise integration.

\emph{The arithmetic normalisation.} On the curve
$Y^2=4x^3-g_2x-g_3$, let $e$ be the $x$-coordinate of the marked
point, and $e',e''$ the other two roots of the cubic. Put
$a=e-e'$, $b=e-e''$, $P=ab=3e^2-g_2/4$, and $X=P^{-3}\in M$.
All three roots are distinct. The discriminant identity gives
$16P^2(a-b)^2=1$, while the sum of the roots gives
$g_2=\tfrac43((a-b)^2+P)$. Substituting these identities into
$j=1728g_2^3$ yields
\begin{equation}\label{D:166}
                    (X+16)^3-jX=0.
\end{equation}
The singular modulus $j$ is an algebraic integer by complex
multiplication. The displayed monic equation therefore makes $X$ a
nonzero algebraic integer. This is the only integrality calculation
needed here.

For clarity, the classical Jacobi identities used to recover the
absolute value of $X$ are as follows. The three half-period
differences, up to sign, are $\pi^2\omega^{-2}$ times the fourth
powers of the three even theta constants. The two differences through
$e$ use the other two constants. The product of all three constants
is $2\eta(\tau)^3$, where
$\eta(\tau)=e^{\pi i\tau/12}\prod_{n\ge1}(1-e^{2\pi i n\tau})$.
Finally, the discriminant identity
$1=(2\pi/\omega)^{12}\eta(\tau)^{24}$ gives
$|\omega|=2\pi|\eta(\tau)|^2$. Hence at each embedding
\begin{equation}\label{D:167}
 |\rho X|=(T/|\eta(\tau)|)^{12},\qquad
 A_\rho=4\pi^2 y|\eta(\tau)|^4,\qquad
 \log H_\rho=\tfrac12\log(2\pi^2)-\tfrac12\log A_\rho
                 -\tfrac16\log|\rho X|.
\end{equation}
Here the same algebraic number $X$ works at every embedding because
its formula uses just the marked coordinate $e$ and $g_2$.

The conjugates of a curve with complex multiplication by $\mathcal O_K$
run equally often through the ideal classes of $K$ (input (I2)). By the class mean
period formula of \S I.7 in the discriminant-one scale (input (I3)),
$\langle\log A_\rho\rangle=\log\pi+\gamma+
 \log(2\pi/\sqrt d)-c_d$, where $\gamma$ is Euler's constant.
Since $|N_{M/\mathbb Q}X|$ is a positive integer, averaging gives
\begin{equation}\label{D:168}
 \langle\log H_\rho\rangle
 =\tfrac12(\log\sqrt d-\gamma+c_d)
 -\frac{\log |N_{M/\mathbb Q}X|}{6[M:\mathbb Q]}
 \le \tfrac12(\log\sqrt d-\gamma+c_d).
\end{equation}
The reference section of ranks $N,N'$ has squared norm
$\Theta_{0,\rho}=H_\rho^{fNN'}$, where $f=2n$ is the number of weight places ($n$ as in Part II).
This last identity uses the tensor-power description of the normalised
reference section; the integral of a product on independent elliptic
coordinates factors. Thus its average logarithm is at most $fNN'$
times the displayed bound. This is the conclusion stated at the head of the paragraph.

\emph{The step that uses $L_{\rm sc}$.} This step is separate from the conclusion above. In the application
$|c_d|\le4\log x$, by Theorem~\ref{thm:pair}(d) with $M=1$ and $\chi=\chi_d$. Put
$v_{\rm host}=\mathrm{rd}_E/2$, and let $L_{\rm sc}$ be the logarithmic scale of the Lean development; it
is defined in the source. In the parameters of the chain one has
$\log\sqrt d-\gamma\le\log v_{\rm host}$ (the inequality $\sqrt d\le v_{\rm host}$ itself holds for
$N\ge8$ but fails for $N=4$, $q=3$) and
$\log v_{\rm host}+|c_d|\le2L_{\rm sc}$; both inequalities are taken from the source. So the contribution
is at most $fNN' L_{\rm sc}$.

\textbf{The hook product.}
The hook product enters as an exact Gaussian integral. Given integers
$\lambda_1\ge\cdots\ge\lambda_s\ge0$, form a diagram with $\lambda_i$
boxes in row $i$, left justified. The hook length of a box is one plus
the numbers of boxes to its right in the same row and below it in the
same column. Let $H(\lambda)$ be the product of all hook lengths,
with empty product one. For a square complex matrix $X$ of size $s$ put
\begin{equation}\label{D:169}
 p_\lambda(X)=\prod_{i=1}^s
       \det(X_{1:i,1:i})^{\lambda_i-\lambda_{i+1}},
       \qquad \lambda_{s+1}=0 .
\end{equation}
With the Gaussian probability norm used above,
$\|p_\lambda\|_G^2=H(\lambda)$.

Here is a proof by removing full columns. Write $X=QR$ as in the proof
of the polynomial estimate. Conditional on the preceding columns,
$R_{ii}^2$ is the squared length of a standard complex Gaussian vector
in dimension $s-i+1$. Its density is $t^{s-i}e^{-t}/(s-i)!$ on $t>0$,
by polar integration. This conditional law is constant, so the
$R_{ii}^2$ are independent. Averaging first over unitary matrices on the
left as above gives
\begin{equation}\label{D:170}
 \|p_\lambda\|_G^2=
 \left(\int_{U(s)}|p_\lambda(U)|^2\,dU\right)
 \prod_{i=1}^s\frac{(\lambda_i+s-i)!}{(s-i)!},
\end{equation}
where $dU$ has total mass one. Put $a=\lambda_s$ and
$\mu_i=\lambda_i-a$. Then $p_\lambda=(\det X)^a p_\mu$, so the
unitary integrals for $\lambda$ and $\mu$ agree. The ratio of their
squared Gaussian norms is consequently
\begin{equation}\label{D:171}
       \prod_{i=1}^s\frac{(\lambda_i+s-i)!}{(\lambda_i-a+s-i)!}.
\end{equation}
These are exactly the hook lengths in the first $a$ columns, since the
hook at $(i,j)$ there has length $\lambda_i-j+1+s-i$. Deleting those
columns leaves all other hook lengths unchanged. The last row is now
empty, and $p_\mu$ depends only on the upper left $(s-1)$ by $(s-1)$
matrix. Induction on $s$ proves the formula, including $a=0$.

\textbf{The eigenvalues of a lift: what the inner-product formula contributes.} \emph{This paragraph is
conditional on the following hypotheses, which are taken from the source and not proved here.} The
notation is explained below the list.
\begin{enumerate}[label=\textup{(E\arabic*)},leftmargin=2.6em,itemsep=1pt,topsep=3pt]
\item At each weight place the polynomial in the theta lift is a product of two polynomials of the form
$p_\lambda$ on disjoint sets of variables.
\item On each non-empty simultaneous Hecke component of the target, the Rallis inner-product formula,
evaluated with the chosen local test functions, gives $TT^*=\lambda I$ with the scalar $\lambda$ of the
form displayed below; in it $D_0\ge1$, and the level index times the squared norm of the finite test
function equals $1/D_S$.
\item The parameter conditions: $1\le f\le\varrho$, $1\le r\le\varrho/2$, $t\ge r+5$, $\varrho\ge8$,
$L\ge5.5$, $\log\varrho\le L$, and $\log(d_\rho+2)\le2L$.
\item The bound for the mass $M_0$, the bounds for the normalising factors $D_S$, $D_{S^c}$, and the
formal-degree estimate for $J_\rho$, as displayed below.
\item The local parameter bound $|x|\le(\mathrm N\mathfrak p)^{(r-t-2)/2}$ for the parameters of the
unramified Euler factors.
\item The lift used at the distinguished embedding is $S=PT$ with $P$ an orthogonal projection on the
target, and the Hecke components span the target.
\item At the other embeddings the lower bound of the conclusion holds as well.
\end{enumerate}
\emph{Conclusion.} Under (E1)--(E7) every positive eigenvalue of $SS^*$ is at least
\begin{equation}\label{D:172}
 \exp\Big(\sum_\rho\log H(A_\rho^+)-8f\varrho^2L\Big).
\end{equation}

\noindent The lower bound for these eigenvalues is one of the inputs listed at the end of this section.
Fix one step of the chain. At each of the $f=2n$ weight places $\rho$,
the polynomial in the theta lift is a product of two polynomials of the
form $p_\lambda$ on disjoint sets of variables, by (E1). Write $A_\rho^+$ and
$A_\rho^-$ for their diagrams; their row lengths are the column degrees
of the two products of minors, arranged in decreasing order. Reversing
rows or columns in a corner minor leaves its absolute value unchanged
after the corresponding change of variables. The Gaussian integrals
over disjoint coordinates factor, giving the squared norm
$\prod_\rho H(A_\rho^+)H(A_\rho^-)$.

Let $T$ be the full theta lift between its finite-dimensional complex
inner-product spaces. On each non-empty simultaneous Hecke component of the target
the Rallis inner-product formula, evaluated with the chosen local test
functions, gives $TT^*=\lambda I$, by (E2). In the normalisation of the source
its scalar has the form
\begin{equation}\label{D:173}
 \lambda= D_0\,M_0\,
       \frac{E_0}{D_S D_{S^c}}\,
       \prod_\rho H(A_\rho^+)H(A_\rho^-)J_\rho .
\end{equation}
Here $D_0\ge1$ is a product of representation dimensions, $M_0$ is
the mass (the sum of reciprocal stabilizer orders), $E_0$ is the
absolute value of the unramified Euler product, and $D_S,D_{S^c}$
are the positive normalising factors at the exceptional finite places
and their complement. The level index times the squared norm of the
finite test function equals $1/D_S$. The factor $J_\rho$ is the integral
of the squared absolute value of a unit vector's matrix coefficient
on the real unitary group at $\rho$, for the measure of the formula; it is defined in the source.

Here are the bounds on these factors that suffice. Here $k$, $m$ and $t$ are parameters of the construction
(in Part II, \S6.2: the initial rank $k$, the number $m$ of lifts, and the odd parameter written $\tau_0$
there), and $f=2n$ is the number of weight places. Write $r=k+m-1$,
$\varrho=r+t$ for the rank parameter, $g_0,g_1$ for the two level primes,
$b$ for the weight bound, and
$L=\log(b\,\mathrm{rd}_E\,g_0g_1\varrho)$. (In Part II, \S6.6 the letter $h$ denotes $k+m-1$, which is $r$ here.)
By (E3) the parameter conditions give
$1\le f\le\varrho$, $1\le r\le\varrho/2$, $t\ge r+5$, $\varrho\ge8$, $L\ge5.5$ and $\log\varrho\le L$.
The mass bound, the finite Euler factor estimates and the
formal-degree estimate used in the source give
\begin{equation}\label{III:massbd}
 M_0\ge\frac1{(8 f^2)^\varrho\,\varrho!},\qquad E_0\ge e^{-2rf},
 \qquad D_S,D_{S^c}\le e^{4rf},\qquad
 -\log J_\rho\le2+r^2\log\bigl(2\sqrt2\,r^2\varrho(d_\rho+2)\bigr),
\end{equation}
where $d_\rho=|A_\rho^+|+|A_\rho^-|$ counts their boxes and $\log(d_\rho+2)\le2L$. The bounds for $M_0$,
$D_S$, $D_{S^c}$ and $J_\rho$ are (E4). The bound for $E_0$ follows from (E5):
the factors of $E_0$ are $|1-x|^{-1}$ with $r$ parameters
at each unramified prime ideal $\mathfrak p$ of $E$. Their absolute
values satisfy $|x|\le(\mathrm N\mathfrak p)^{(r-t-2)/2}
\le(\mathrm N\mathfrak p)^{-2}$, by the local parameter bound and
$t\ge r+5$. At most $[E:\mathbb Q]=2f$ prime ideals lie above any
rational prime $p$, and each has norm at least $p$, so $\sum_{\mathfrak p}(\mathrm N\mathfrak p)^{-2}
\le2f\sum_{p}p^{-2}\le2f$, the last sum being over the rational primes. Taking logarithms and using
$\log|1-x|\le|x|$ proves the stated bound for $E_0$.
The logarithmic losses are at most $4\varrho L$ from the mass, $10rf$ from
the finite factors, and $f(2+6r^2L)$ from the real integrals.
Their sum is at most $8f\varrho^2L$ by the preceding inequalities.
Since every hook length is at least one, this gives
\begin{equation}\label{D:175}
             \log\lambda\ge
             \sum_\rho\log H(A_\rho^+)-8f\varrho^2L .
\end{equation}
This deduction uses the stated arithmetic and analytic estimates for
the test functions.

Finally, by (E6), the lift used at the distinguished embedding is $S=PT$ with
$P$ an orthogonal projection on the target, and the Hecke components span
the target. Since $TT^*$ is self-adjoint, the scalar bounds on those
components give $TT^*\ge cI$, where $c$ is the exponential of the
last lower bound. On the range of $P$ one has
$\langle SS^*x,x\rangle=\|T^*x\|^2\ge c\|x\|^2$. This bounds every
positive eigenvalue of $SS^*$. At the other embeddings the same bound is taken from the Lean development;
this is (E7).

\textbf{Row by row: why only the first row survives.} \emph{This passage is conditional: it is a deduction
from inputs taken from the source, listed below.}
It compares the terms $\sum_ih(\nu_i)$ of the block estimate with the hook product of the preceding
paragraph, one row at a time. Fix a step $1\le j\le m$ of the chain (steps are numbered from 1 here, as in the source) and one weight place $\rho$.
Let $k$ and $t$ be the parameters of the preceding paragraph, with $k\ge2$. The terms used are these.
\begin{itemize}[leftmargin=2.4em,itemsep=1pt,topsep=3pt]
\item The \emph{kernel degrees} are the column degrees of the polynomial of the kernel of this step in the
block at this place. We assume, as input (o) from the source, that these are the degrees $\nu_i$ to which
the block estimate is applied (one block for each weight place); in the Lean development this is the case: the degrees fed to the block
estimate are identified with these numbers (\texttt{weights\_rows}, \texttt{nu\_eq\_originRow}). We call them also the lengths of the kernel rows. There are at least $k$ of them.
\item $A_\rho^+$ is the positive-signature diagram of the preceding paragraph, and $H(A_\rho^+)$ its hook
product.
\item The weight shift $w$ at this step, the weight parameters $u_1,\ldots,u_k$ of the domain at this place, and
the parameters $v_1,\ldots,v_{k-1}$ of the positive-signature rows are as in the source (Lean:
\texttt{origin\_rewrite}, \texttt{origin\_interlace}, with the formulas \texttt{originRow},
\texttt{originBig}). We do not define them here.
\item A weight place is \emph{ordinary} or \emph{distinguished}, as in the source (Lean: the two cases
\texttt{RowsOriginOff} and \texttt{RowsOriginTop}). At a distinguished place, which we denote by
$\mathfrak t$, there are two integers $a_{\mathfrak t}\ge0$ and $b_{\mathfrak t}\ge0$; $a_{\mathfrak t}$ is
the number $a_t$ of the inequality at the end of this section. The index $\mathfrak t$ is not the
parameter $t$.
\item $\Delta$ is the width of the weight interval (in the source, \texttt{Sw}).
\item $B$ is a bound for the sum of the kernel degrees other than the first $k$.
\item $L_a=\log(a\,\mathrm{rd}_E)$, with $a$ as in the inequality at the end of this section.
\end{itemize}
Put
\begin{equation}\label{D:176}
 \varsigma=\frac{k+(2j-1)(t+1)}2,\quad W=w-\varsigma,\quad \delta=(j-1)t,
\end{equation}
\begin{equation}\label{D:177}
 x_i=u_i-\frac{k+1-2i}{2},\quad y_i=v_i-\frac{k-2i}{2}.
\end{equation}
Assumed here, from the source:
\begin{enumerate}[label=\textup{(\alph*)},leftmargin=2.4em,itemsep=1pt,topsep=3pt]
\item \emph{The degree formulas.} $W\ge0$, the first $k$ kernel degrees are given by the formulas for $\nu_i$ below, and
$k-1$ of the rows of $A_\rho^+$ have the lengths $R_i$ below:
\begin{equation}\label{D:178}
  \nu_i=W+x_i\quad(1\le i\le k),\qquad
  R_i=W+y_i+\delta\quad(1\le i<k).
\end{equation}
There may be further rows on either side. (Lean: \texttt{nu\_eq\_originRow},
\texttt{aPlus\_eq\_originBig}.)
\item \emph{Strict interlacing.} $u_i>v_i>u_{i+1}$ for $i<k$.
\item \emph{Integrality and parity.} The degrees $\nu_i$ and $R_i$ are nonnegative integers; $k\ge2$ is even,
$t\ge1$ is odd, and the weight $w$ is an integer. Then $\varsigma$, the $x_i$ and the $y_i$ are integers.
(The argument below uses only that $x_i-y_i=\nu_i-R_i+\delta$ and $y_i-x_{i+1}=R_i-\delta-\nu_{i+1}$ are
integers, and this follows from the integrality of $\nu_i$, $R_i$ and $\delta$.)
\item \emph{The endpoint bounds.} At an ordinary place, $x_1\le\Delta$ and $x_k\ge0$. At a distinguished
place, $x_1=a_{\mathfrak t}$, $x_k=-b_{\mathfrak t}$, $y_{k-1}\ge0$ and $y_1\le\Delta$.
\item \emph{The extra degrees and the sizes.} The sum of the kernel degrees other than the first $k$ is at
most $B$, and $B\le a\,\mathrm{rd}_E$; $w+\Delta\le a\,\mathrm{rd}_E$ at an ordinary place, and
$w+a_{\mathfrak t}\le a\,\mathrm{rd}_E$ at a distinguished place. The size hypotheses of the source permit
$B=m(t+1)\eta$, with $\eta$ the bound used there for the additional rows.
\end{enumerate}
None of (a)--(e) is proved in these pages.

\emph{The comparison of rows.} Strict interlacing and integrality give
$x_i\ge y_i\ge x_{i+1}$, since both differences are integers greater
than $-1/2$. Consequently
\begin{equation}\label{D:179}
                         R_i\ge\nu_{i+1}\quad(1\le i<k).
\end{equation}
In particular every one of the first $k$ kernel rows except the first
has its own row of at least the same length in $A_\rho^+$.

Recall that $h(v)=v\log v-v$ for $v>0$ and $h(0)=0$. For integers $d\ge0$
one has $h(d)\le\log(d!)$ by integration of $\log x$.
The hooks along a row of length $R$ are at least $R,R-1,\ldots,1$.
Hence $\log H(A_\rho^+)\ge\sum_{i<k}\log(R_i!)\ge\sum_{i=2}^k h(\nu_i)$.
Each remaining kernel degree $d$ satisfies $h(d)\le d\,L_a$ (by (e), since $d\le B$), so
\begin{equation}\label{D:180}
           \sum_{\text{all kernel rows}}h(\nu_i)-\log H(A_\rho^+)
                 \le h(\nu_1)+B\,L_a .
\end{equation}
At an ordinary place $x_1\le\Delta$ and $w+\Delta\le a\,\mathrm{rd}_E$. Thus the right side
is at most $(w+\Delta+B)L_a$. At a distinguished place $x_1=a_{\mathfrak t}$ and
$w+a_{\mathfrak t}\le a\,\mathrm{rd}_E$. Integrating $h'(v)=\log v$ gives
$h(a_{\mathfrak t}+W)\le h(a_{\mathfrak t})+W L_a$, also when $a_{\mathfrak t}=0$. The right side is
therefore at most $h(a_{\mathfrak t})+(w+B)L_a$. This accounts for the one
large factorial term left at each distinguished place.

\emph{The degree that multiplies the period.} The same comparison explains the period term. At an
ordinary place $0\le x_k\le x_1\le\Delta$. Here the lower bound $x_k\ge0$ is the hypothesis in (d); it
does not follow from $\nu_k\ge0$, and the bound below fails without it. At a distinguished place the
inner weights satisfy $0\le y_{k-1}\le y_1\le\Delta$, hence $0\le x_i\le\Delta$ for
$1<i<k$; the last endpoint is $x_k=-b_{\mathfrak t}$ with
$0\le b_{\mathfrak t}\le W$, since $\nu_k\ge0$. Write $a_*=a_{\mathfrak t}$ at a
distinguished place and $a_*=0$ otherwise. Summing the row lengths gives
\begin{equation}\label{D:181}
  \left|\sum_{\text{all kernel rows}}\nu_i-kw-a_*\right|
      \le k(\varsigma+\Delta)+w+B .
\end{equation}
Thus the degree that multiplies $-\log\Omega_E$ is $k\,w+a_{\mathfrak t}$ at a distinguished place, and
$k\,w$ at an ordinary place, up to the displayed error. The part $k\,w$ cancels only through the comparison
of periods, item (iii) of the list at the end of this section. Nothing here identifies the period of the
block with the period in the covolume.

\emph{What this proves, and what it assumes:} from (a)--(e) it proves that each of the first $k$ kernel
rows after the first is matched by a row of $A_\rho^+$ of at least the same length, that the terms
$\sum_ih(\nu_i)$ exceed $\log H(A_\rho^+)$ by at most $(w+\Delta+B)L_a$ at an ordinary place and
$h(a_{\mathfrak t})+(w+B)L_a$ at a distinguished place, and that the total degree is $k\,w+a_*$ up to
$k(\varsigma+\Delta)+w+B$; it assumes (o) and (a)--(e), which are statements about the weights of the construction
taken from the source and not proved here.

\textbf{The pairing at the end of the chain.} \emph{This passage is conditional: the form of the last map,
its level factor and the form of the final functional are inputs from the source, marked below.}
Here the rank decreases from $\Psi=\psi+1$ to $\psi$, with $\psi\ge1$. The last map of the chain goes from
polynomials of weight $\lambda$ on $\Psi$ by $\Psi$ matrices (the space $W_{\lambda'}$, in the notation
$W_\nu$ above) to polynomials of weight $\sigma$ on $\psi$ by $\psi$ matrices (the space $W_{\sigma'}$);
$\mathcal T$ is the level of Part II, \S6.7. We give the calculation
at one place. Let $\lambda_1\ge\cdots\ge\lambda_\Psi$ and
$\sigma_1\ge\cdots\ge\sigma_\psi$ be the integer weights, with
$\lambda_i\ge\sigma_i\ge\lambda_{i+1}$ and $\lambda_\Psi=\sigma_\psi$.
Subtract their common last entry to get polynomial degrees
$\lambda'_i,\sigma'_i$. In the polynomial spaces used above, divide each Gaussian norm by the Gaussian
norm of the corresponding highest polynomial, $p_{\lambda'}$ or $p_{\sigma'}$. Set
\begin{equation}\label{D:182}
 a_i=\sigma_i-\lambda_{i+1},\quad b_i=\lambda_i-\sigma_i,\quad
 d_0=\sum_{i=1}^\psi a_i,\quad
 \ell=\sum_{i=1}^\Psi\lambda_i-\sum_{i=1}^\psi\sigma_i=\lambda_1-d_0 .
\end{equation}
The exponent $\ell$ uses the original integer weights.

Write the matrix of variables as its first $\psi$ rows and a last row
$(z_1,\ldots,z_\psi)$, omitting its last column, whose degree is zero.
Let $C$ extract the coefficient of $z_1^{b_1}\cdots z_\psi^{b_\psi}$ and
project it orthogonally on the target weight space. Let $D_i$ be the
upper left $i$ by $i$ minor and let $D_i'$ be that minor with row $i$
replaced by the last row. Put
\begin{equation}\label{D:183}
 \xi=\prod_{i=1}^\psi D_i^{a_i}(D_i')^{b_i},\qquad
 \kappa=\frac{\|\xi\|_G^2}{\|p_{\lambda'}\|_G^2}.
\end{equation}
Then $C\xi=p_{\sigma'}$. Indeed the exponent $b_\psi$ forces every factor
$D'_\psi$ to use $z_\psi$. Descending induction forces each $D'_i$ to
contribute $z_iD_{i-1}$, with $D_0=1$. The remaining product is
$p_{\sigma'}$.

The next identity is needed for the norm of $C$: $C^*p_{\sigma'}=\xi/\kappa$.
The operator $C$ commutes with unitary substitutions on the first $\psi$
rows, as does its adjoint. Such an intertwiner also commutes with
complex invertible row substitutions: differentiate on the unitary
group and take complex linear combinations, since skew-Hermitian
matrices span all matrices over $\mathbb C$. Thus $C^*p_{\sigma'}$
has row degrees $\sigma'_i$ in the first $\psi$ rows and is invariant
under adding an earlier row to a later one among them. (Rows are meant here, where the definition of
$W_\nu$ used columns: this step uses that $p_{\sigma'}$ has row degrees $\sigma'_i$ and is unchanged when
an earlier row is added to a later one, as well as the corresponding properties for columns.) There is at most
one such polynomial up to scalar in the weight space of the domain. Indeed,
where the leading minors are nonzero, these row operations and the
allowed column operations reduce the upper $\psi$ by $\psi$ matrix to a
diagonal matrix. On this slice row and column degrees force a single
monomial, $\prod_i X_{ii}^{\sigma'_i}z_i^{b_i}$. Polynomial identity
extends uniqueness to all matrices. The polynomial $\xi$ has these
properties, and pairing with $\xi$ gives the asserted scalar.
The unitary row translates of $p_{\sigma'}$ span the target. To see
this, any nonzero invariant orthogonal complement would contain a
vector of lexicographically largest row degree among its diagonal
torus weights. Adding an earlier row to a later one fixes this vector,
since every nonconstant term of the substitution raises that weight.
Reduction to a diagonal matrix as above now identifies this vector
with a multiple of $p_{\sigma'}$. It could not lie in the complement.
It follows by equivariance that $CC^*=I/\kappa$ on the target.

\emph{The integral normalisation is an input from the source.} There the map is $\vartheta\kappa C$
(Lean: \texttt{brStd}, with the constant \texttt{brConst}),
where $\vartheta$ is the following nonzero integer. For each $i\le\psi$ take all
integers $v$ from $\lambda_{i+1}$ to
$\min(\lambda_i,\lambda_{i+1}+d_0)$, omitting $\sigma_i$.
Multiply $v-i-(\sigma_j-j)$ over these $i,v$ and over $1\le j\le\psi$.
Interlacing gives $v-i>\sigma_j-j$ for $i<j$ and the reverse inequality
for $i>j$. The omission excludes zero when $i=j$. Put
$\Gamma=\lambda_1-\lambda_\Psi+\Psi$ and
$d_1=\sum_{i=1}^\psi\min(\lambda_i-\lambda_{i+1},d_0)$.
Each factor has absolute value at most $\Gamma$ and there are $\psi d_1$ of them.
What the source shows to be integral is the transposed map $\vartheta\kappa C^*$ (Lean:
\texttt{BranchMap}); the map $\vartheta\kappa C$ itself can have denominators. A map and its adjoint have
the same norm, so the statements about norms here are not affected. By $CC^*=I/\kappa$, the squared norm
of either map is $\vartheta^2\kappa$.

One can also bound $\kappa$ by the Gaussian calculation above. The two
polynomials in its ratio, $\xi$ and $p_{\lambda'}$, have identical column degrees, and both are unchanged
when a multiple of an earlier column is added to a later column. So both lie in $W_{\lambda'}$, the
argument by $X=QR$ applies to both, and their radial factors cancel. Equal column degrees alone would not
be enough. Every minor of a unitary matrix has absolute value
at most one. In addition, the Gaussian hook formula and the radial
moments give
\begin{equation}\label{D:184}
 \kappa\le\left(\int_{U(\Psi)} |p_{\lambda'}(U)|^2\,dU\right)^{-1}
 =\frac{1}{H(\lambda')}\prod_{i=1}^\Psi
       \frac{(\lambda'_i+\Psi-i)!}{(\Psi-i)!}
 \le\prod_{i=1}^\Psi {\lambda'_i+\Psi-i\choose \Psi-i}
 \le \Gamma^{\Psi(\Psi-1)/2}.
\end{equation}
The middle bound uses the row factorial lower bound for $H(\lambda')$.

\emph{The level factor is taken from the source, not derived here} (Lean: \texttt{brAt}).
At level $\mathcal T$ the map has the additional scalar
$\varpi^{-\ell}$ with $|\varpi|^2=\mathcal T$ in orthonormal coordinates. Its
squared norm is consequently $\mathcal T^{-\ell}\vartheta^2\kappa$.
Factors at independent places multiply. Taking logarithms, the full
bound is
\begin{equation}\label{D:185}
 -\sum_\rho \ell_\rho\log\mathcal T+
 \sum_\rho\bigl(2\psi d_{1,\rho}+\Psi(\Psi-1)/2\bigr)\log \Gamma_\rho .
\end{equation}
The final functional pairs these images with the target by a weighted
sum over lattice classes, with reciprocal stabilizer orders as weights. That it has this form is an input
from the source. Cauchy--Schwarz for this finite sum gives precisely the same bound for
the squared pairing on unit vectors. Here unit vectors are tuples indexed by the classes, of norm one in
the weighted class norms $\sum_pw(p)\|f_p\|^2$ of the paragraph on sums over classes, with $w(p)$ the
reciprocal stabilizer order; and the squared pairing is the squared absolute value of the sum over the
classes $p$ of $w(p)$ times the inner product of the image of $f_p$ with the $p$th entry of the second
tuple (Lean: \texttt{sum\_pair\_le}).

\emph{What this proves, and what it assumes:} it proves that $CC^*=I/\kappa$ in the normalised norms, so
that $\vartheta\kappa C$ and its adjoint have squared norm $\vartheta^2\kappa$, that
$|\vartheta|\le\Gamma^{\psi d_1}$ and $\kappa\le\Gamma^{\Psi(\Psi-1)/2}$, and the displayed bound for the
squared pairing; it assumes interlacing of the two weights with a common last entry, that the last map of
the construction is $\vartheta\kappa C$ up to the level factor, that this factor is $\varpi^{-\ell}$, and
that the final functional is the weighted sum over classes, and it does not substitute the actual weights
into $\ell$ and $d_1$.

\textbf{Relation to the Lean proof.} The estimate for one block, with the constant
$e^{1/2}(2n+2)^d\exp(\frac12\sum_ih(\nu_i))$ and the factor $\Omega_\tau^{-n_\tau}$, is the statement of the
Lean development that the final theorem depends on (\texttt{blockBound\_of\_fockJet}, with
\texttt{gaussianCoeffIntegral}); the Lean statement is for the square of the norm, which is why the period enters as $\Omega_\tau^{-n_\tau}$;
with the factor $|I|\,|J|$ from the sum over classes it becomes the bound for one lift. The source proves it through an integral over $\C^d$ against the Gaussian measure; the proof
above uses the sphere instead. For a chain of $m$ lifts the source bounds the norm of each map at each
archimedean place in this way, multiplies the bounds (the norm of a composition is at most the product of
the norms), and averages the logarithms over the places (\texttt{chain\_coupled\_bound}). The hook products and the
lower bound for the eigenvalues of a lift enter the assembly as a term subtracted for each map: in the
source the bound for the height has the form $h_A+\sum_j(h_{B,j}-h_{X,j})$, where $h_{B,j}$ comes from the
block estimate and $h_{X,j}$ from the eigenvalue bound.

\textbf{What remains to be written.} In the Lean development the passage from this estimate to the first
line of \eqref{II:t6.5} is an assembly (\texttt{divClausesR\_of}) of the following further statements, each proved
elsewhere in the development. The assembly itself is carried out in \S\ref{sa:sec}, which proves that the first line of \eqref{II:t6.5} follows from a list of such statements, and proves one statement on the periods. The statements themselves are not proved in these pages. Until they are, the first line of \eqref{II:t6.5} remains a quoted statement.
\begin{enumerate}[label=\textup{(\roman*)},leftmargin=2.4em]
\item \emph{The inputs \textup{(I1)--(I5)} of the reference mean square above}: the integrality of the singular modulus
(classical), that the conjugates of the curve run equally often through the ideal classes, the formula for
the mean of $\log A_\rho$ over the classes (the class mean period formula), the description of the
reference section as a tensor power, and the identification of the reference section at $\rho$ with
$\theta/\theta(0)$ and its metric. Also: that the chain of maps is defined over a field containing $M$, with
$H_{\bar\rho}=H_\rho$; the two inequalities for $v_{\rm host}$ and $L_{\rm sc}$ used in the last step
there; and the bounds for $\beta$ and $c$ in the
constant $S$.
\item \emph{The pairing at the end of the chain, and the eigenvalues of a lift.} The norm of the pairing at
the end of the chain and the row-by-row comparison are proved above, conditionally on the inputs named
there. What remains: the integral realisation of the last map, with its level factor; the substitution of
the actual weights of the construction (into $\ell$ and $d_1$ for the pairing, and into the inputs (o) and (a)--(e)
of the comparison); and, for the lower bound for the eigenvalues of a lift, the hypotheses (E1)--(E7) of the conditional
paragraph above: the form of the polynomial at each weight place (E1); the Rallis inner-product formula
with the chosen test functions and the form of its scalar (E2); the parameter conditions (E3); the bound
for the mass, the bounds for the normalising factors and the formal-degree estimate (E4); the local
parameter bound (E5); the form $S=PT$ of the lift and the spanning by Hecke components (E6); and the bound
at the other embeddings (E7).
\item \emph{The periods}: that the average over the archimedean places of $\log\Omega_\tau$ is the constant
$\log\Omega_E$ of \S I.7 (the constant relevant here is $\Omega_E$, not the elliptic $\Omega$; the two differ
by the factor $\mathrm{rd}_E/\sqrt d$), and that the same constant appears in the covolume of the lattices. This is where
the Chowla--Selberg formula is used, and where the two occurrences of $\log\Omega$ cancel up to a small error.
\item \emph{Local index bounds} at the finite places, including the term $F_0$ of Part II, \S6.7, which is
supplied by the divisibility statement (item 1 of \S I.7).
\item \emph{The passage between levels}, and bounds for class numbers and for the constants of the kernels.
\item \emph{That the kernels and polynomials of the construction are instances of the abstract setting}:
the lattice hypothesis with $v=\mathrm{rd}_E$, the form \eqref{III:Sform} of $S$, the identification of $P$ with
a substitution, and the tuple representation of each map with matching norms.
\end{enumerate}
The inequality that results is, in the notation of Part II, \S6.7,
\begin{equation}\label{III:sizefinal}
 \mu_{\rm CC}+h-F_0\le\sum_{t\in T}\Big(m\big(a_t\log a_t-a_t-a_t\log\Omega_E\big)-a_t\log\mathcal T\Big)
 +C\,|T|\,a\,L_a+C\,m\,(B_0+N_c)\,x^{2\varepsilon},
\end{equation}
where $a=x\,b$ (the orders $a_t$ are not assumed to be at most $a$; Part II allows $a_t$ up to $(x+1)\,b$), $L_a=\log(a\,\mathrm{rd}_E)$, $B_0=2n\,m\,b\,x^{\varepsilon}$ and $N_c=|T_2|\,k\,b\,x^{2\varepsilon}/n^2$.
(In the Lean record \texttt{RecNums.rhs} the last factor is $x^{\varepsilon}$ times a constant, and the form above is weaker; the final Lean statement has the factor $x^{C_{\rm abs}\varepsilon}$ with a free constant $C_{\rm abs}$, and the form above is its case $C_{\rm abs}=2$. The
value of $N_c$ is the one fixed in the Lean file \texttt{DataArith.lean}.)
Here $h$ is an upper bound for the height of the pairing map, $\mu_{\rm CC}$ is the mean slope of the lattice in which the pairing maps take their values (in the source it is the dual of a fixed lattice), in the normalisation for squared norms, as in Part II, \S6.7, and $F_0$ is as in Part II, \S6.7. After division by $|T|\,m\,b\,\delta/2$, where $\delta=|T_2|/n$, the term $C\,|T|\,a\,L_a$ gives the
error $x/m$ of (M), the summand with $B_0$ gives $m$, and the summand with $N_c$ gives $k/n^2$ (each up to a
factor $x^{O(\varepsilon)}$); the fourth error, $xm/k$, does not come from this inequality.

\section{Divisibility at inert primes}

This section concerns the lower bound for the exponent of an inert prime: inequality \eqref{expo} of
\S I.4.3. It proves the local estimate behind that bound, for power series, under one hypothesis of
integrality, called (B) below. In this section the prime is called $\ell$, as in the Lean files; it is the
inert prime $p$ of Part~I. Several letters used here have other meanings elsewhere in the paper, among them
$\beta$, $\lambda$, $\varepsilon$, $a$, $p$, $d$, $x$, $N$, $D$, $q$, $k$, $m$, $n$, $f$, $c$, $\sigma$, $e$ and $\psi$; in this section each has the meaning
given where it is introduced.

\emph{Proved outright.} The valuation of the coefficients of the inverse of the height-two logarithm, in
several variables (Lemma~1); the integrality lemma (Lemma~2); the optimisation for the top row (Lemma~6),
which produces the sum $\sum_{h\ \mathrm{odd}}(A/\ell^h-w_h)_+$; weak duality for the linear program of \S5;
and the closed form of that program for orbits of length two.

\emph{Proved conditionally on hypothesis \textup{(B)}.} Lemmas~3 and~4: the passage from (B) to a bound for
a series $\beta$ in the exponent, with the gain from $1/(\ell^2-1)$ to $1/(\ell-1)$. Proposition~5: the bound
for one Taylor coefficient of a function $U\cdot\exp(\pm\beta)$, with weights on the variables. Theorem~J,
for orbits of length two: the bound for a product of minors applied to such a function.

\emph{Conditional on \textup{(B)} and on strong duality, cited from the source.} Theorem~J for orbits of
length four or more. The cited statement is the equality of a linear program with its dual (Lean:
\texttt{strongDuality}). This case was not tested end to end.

\emph{Taken from the source without proof.} The transcription of the top-level inequality
\texttt{LinkG\_at} (\S1 below); that the theta kernels of the construction have the form
$U\cdot\exp(\pm\beta)$ and satisfy (B); the dictionary between sites and places; the frames and labels that
Theorem~J requires; the passage from Taylor coefficients (jets) to lattices and to the inequality that the
final theorem uses; and the remaining primes. These are listed at the end.

\subsection*{1.\quad The statement in the source}

\textbf{The top-level inequality.} The final theorem depends on the Lean statement \texttt{LinkG}, of
the file \texttt{GlobalTheoremInputs.lean}. It does so through the theorems \texttt{closed\_linkG\_of\_duty}
and \texttt{linkG\_concrete5}.
Its third clause is the inequality of Part~II, \S6.4, last paragraph. We state it. The scale is
$x>1$, $0<\varepsilon<1/52$, $b=\exp(x^\varepsilon)$, $a=xb$, $X'=\exp(x^{\varepsilon/3})$, $b'=a/X'$,
$w''=b'/(xn)$, $X''=a/w''$, $p_0=k^2x^{26}$, $\delta_1=x^{-\varepsilon/8}$. Here $m$ (the number of lifts),
$k$ and $t_0$ are integer parameters of the construction (the fields \texttt{m}, \texttt{k}, \texttt{t} of
\texttt{D.R}); $n$ is a field of the parent structure (\texttt{ord}); and $f=2n$ and $N_m=k+m(t_0+1)$ are
defined quantities, not fields. $\mathcal T$ is a finite index set (\texttt{D.host.TT}). The
following are \emph{given data} (fields of the structure \texttt{SizeData}$_2$): integers $a_t$ for
$t\in\mathcal T$ with $a\le a_t\le a+a/(kx^{13})$; real
numbers $\mathrm{gain}(j,w)$ and $\tau(w)$ for $1\le j\le m$ and $w$ a place of the coefficient field; real
numbers $\mathrm{wPart}(j,t,p,i)\ge1$; subsets $\mathcal T_w\subset\mathcal T$; and real numbers
$\mathrm{Sc}$, $\eta$, $\Lambda_G$. Put
\begin{equation}\label{dv:V}
 V(j,t,p)=\sum_{\substack{i\ \mathrm{odd}\\ i<\lfloor\log_pa_t\rfloor+2}}
   \Big(\frac{a_t}{p^i}-\mathrm{wPart}(j,t,p,i)\Big)_+ .
\end{equation}
This is the Lean development's \texttt{vonMangoldtClassSum}. For constants $C$ and $b_o$ put
\begin{equation}\label{dv:fl}
 \mathrm{fl}_G(p)=C\,f\,N_m^2\,\log_pa+C\,f\,(\log_pa)^2,\qquad
 \mathrm{lo}_G(p)=C\,(\mathrm{Sc}+k+\eta)\log_p(2a)+\mathrm{pr}(p)+\mathrm{or}(p),
\end{equation}
where $\mathrm{pr}(p)$ and $\mathrm{or}(p)$ vanish unless $p$ is of the class \texttt{inertClass}
(defined below), and for such $p$, with $r_G=\lceil8(4+8\varepsilon)/\varepsilon\rceil+1$ and
$e_p=\lceil\log_p(k^2x^{78})\rceil+1$,
\begin{equation}\label{dv:lo}
\begin{aligned}
 \mathrm{pr}(p)&=\begin{cases}C\,r_G\,b'\log x,& p<p_0,\\
   C\big(b'/x+\eta+\mathrm{Sc}+(w''+t_0)\,\delta_1^{-1}\log x\big)+b_o,& p_0\le p<X',\\ 0,&p\ge X',\end{cases}\\
 \mathrm{or}(p)&=C\log x\Big(\frac{\Lambda_G\,a}{p\,k}+\begin{cases}a\,p^{-e_p},&p<p_0\\ a/p^2,&p\ge p_0\end{cases}\Big)
 \quad(p<X'),\qquad \mathrm{or}(p)=0\quad(p\ge X').
\end{aligned}
\end{equation}
In the definition of $\mathrm{pr}(p)$ the three cases are taken in this order, and the first that applies is
used, as in the Lean (an \texttt{if \ldots\ else if \ldots\ else}). This matters only when $X'\le p_0$.

The third clause of \texttt{LinkG\_at} says: for every place $w$ whose residue characteristic $p$ satisfies
$p\le X''$ and is of class \texttt{signal} or \texttt{inertClass}, one has
$0\le\mathrm{gain}(j,w)\le4\,|\mathcal T|\,a/p$ for $1\le j\le m$, and
\begin{equation}\label{dv:link}
 \sum_{j=1}^m\ \sum_{t\in\mathcal T_w}\Big(V(j,t,p)-\mathrm{lo}_G(p)\Big)\ -\ m\,\mathrm{fl}_G(p)
 \ \le\ \sum_{j=1}^m\mathrm{gain}(j,w)\ -\ \tau(w).
\end{equation}
The statement \texttt{LinkG} asserts this, with four further clauses, for all large $x$, with a constant $C$
and with $b_o=b^{\delta(x)}$ for some function $\delta(x)\to0$. (The classes: a prime not dividing the
conductor, inert in the field $\Q(\sqrt{-d})$ of Part~I (this $d$ is the prime of Part~I; everywhere else in this
section $d$ is the length of an orbit), is of class \texttt{signal} or \texttt{inertClass} according to the
value $+1$ or $-1$ of one further sign; \texttt{DataHost.lean}.)

\textbf{What this statement is.} The inequality \eqref{dv:link} is an inequality between the numbers
$\mathrm{gain}$, $\tau$, $\mathrm{wPart}$. In the Lean it is not a statement about valuations. A separate
field of the family (\texttt{fills}, of type \texttt{ParametersFill}) identifies these numbers with explicit
formulas (\texttt{gainP}, \texttt{tauP}), and the proof of \eqref{dv:link} (\texttt{linkG\_at\_core}) rewrites by \texttt{ParametersFill} and
concludes by linear arithmetic from a hypothesis \texttt{AuxBlock8.Conclusion1}, which already is the
inequality for the formulas, together with \texttt{Conclusion0}, \texttt{lo\_le\_loG} and
\texttt{fl\_le\_flG}. The calculation with the formulas is in \texttt{stmtF}. The statements which identify $\mathrm{gain}$ with a valuation of theta
kernels are other hypotheses and other theorems of the development. We do not claim more than this about
\eqref{dv:link}. The Lean theorem in which a valuation of Taylor coefficients is bounded below by a sum of
the shape \eqref{dv:V} is \texttt{ThmAOp} (\texttt{StatementI.lean}, proved in \texttt{ThmAOp.lean}); the
final theorem depends on its proof \texttt{thmAOp}: this is confirmed by the dependency computation
(see \S6\,(iii)). The rest of this section is about that theorem.

\textbf{One function, several places.} Let $K$ be a field of characteristic zero with a non-archimedean
valuation $v$, normalised by $v(\ell)=1$. Let $d\ge2$ be even. The elements of $\Z/d$ are called
\emph{sites}; they stand for the places in one orbit of the Frobenius of $\ell$, the site $s+1$ being
reached from $s$ by the Frobenius followed by complex conjugation (this identification is made in the
source and is assumed here). Let $N_0,N_1\ge1$. The variables are $Z_{s,q,c}$ and $X_{s,i,c}$, for $s$ a
site, $q,i<N_0$, $c<N_1$. A \emph{frame} is a family $\xi(s,i,q)\in K$ with $v\ge0$; it gives the substitution
$Z_{s,q,c}=\sum_i\xi(s,i,q)X_{s,i,c}$. The function is
\begin{equation}\label{dv:theta}
 \Theta=U\big(g(Z)\big)\cdot\exp\Big(\pm\sum_{q,c}\beta\big((Z_{s,q,c})_{s}\big)\Big),
\end{equation}
where $U$ is a power series in the variables $Z$ with coefficients of valuation $\ge0$, where $g$ is the
inverse of the logarithm \eqref{dv:log} below, taken for the following data (Lean: \texttt{entLog}): the
index set $\mathcal Y$ of \eqref{dv:log} is the set of the variables $Z_{s,q,c}$, and
$\pi(s,q,c)=(s-2,q,c)$ (for $d=2$ this $\pi$ is the identity, and the logarithm is the one-variable
logarithm in each variable); and where $\beta$ is a power series in $d$ variables, one for
each site, subject to the hypothesis (B) of \S3. A \emph{shape} is, for each site $s$, a list of column
lengths $\mathrm{len}(s,1),\ldots,\mathrm{len}(s,\nu_1(s))$ in $[1,N_0]$. Put $\nu_2(s)=\#\{c:\mathrm{len}(s,c)\ge2\}$ and
$\mathrm{size}(s)=\sum_c\mathrm{len}(s,c)$. Let $\mathit{TT}$ be a set of sites (the \emph{top} sites; the Lean
uses the letters \texttt{TT} both for this set and for the index set $\mathcal T$ above). The polynomial
is $P=\prod_s\prod_c\det\big(X_{s,i,I_{s,c}(k)}\big)_{i,k<\mathrm{len}(s,c)}$ for chosen columns $I_{s,c}(k)<N_1$.
The row $i=0$ at a top site is a \emph{top row}; its degree in every monomial of $P$ is $\nu_1(s)$. Put
\begin{equation}\label{dv:shape}
 \mathrm{cap}(s)=\begin{cases}\nu_2(s),&s\in\mathit{TT}\\ \nu_1(s),&s\notin\mathit{TT},\end{cases}\qquad
 D(s)=\begin{cases}\mathrm{size}(s)-\nu_1(s),&s\in\mathit{TT}\\ \mathrm{size}(s),&s\notin\mathit{TT}.\end{cases}
\end{equation}
Let $d_{\rm top}(s)$ be the least odd $e\ge1$ with $s+e\in\mathit{TT}$, or $\infty$. Let $W(s)$ be the
largest $\mathrm{cap}(s-i)$ over odd $i\ge1$ with $\ell^i\le\mathrm{size}(s-i)$ (zero if there is none);
here $i$ runs over all odd integers $\ge1$, not only over $i<d$ (Lean: \texttt{Wcap}). Put
$\gamma=\ell/(\ell^2-1)$, $\rho=1/(\ell^2-1)$, so that $\gamma=\sum_{h\ \mathrm{odd}}\ell^{-h}$ and
$\gamma+\rho=1/(\ell-1)$. The exponent of the Lean theorem is
\begin{equation}\label{dv:B}
\begin{aligned}
 B&=\sum_{s\in\mathit{TT}}\big(\mathrm{top}(s)-\gamma W(s)\big)+F(D)-dN_0N_1(\log_\ell\bar a+1),\\
 \mathrm{top}(s)&=\sum_{\substack{h\ \mathrm{odd}\\ h<d_{\rm top}(s)}}\Big(\frac{\nu_1(s)}{\ell^h}-\mathrm{cap}(s+h)\Big)_+
   +\sum_{\substack{h\ \mathrm{odd}\\ h>d_{\rm top}(s)}}\frac{\nu_1(s)}{\ell^h},
\end{aligned}
\end{equation}
where $\bar a\ge1$ is an integer with $\mathrm{size}(s)\le\bar a$ for all $s$, and $F(D)\ge0$ is the value of a linear
program defined in \S5. \texttt{ThmAOp} asserts $v\big(P(\partial)\Theta(0)\big)\ge B$ under hypotheses on the
frame stated in Theorem~J; here $P(\partial)\Theta(0)=\sum_\kappa P_\kappa\,\kappa!\,\Theta_\kappa$ as in
\eqref{II:t6.2} of Part~II, with $\Theta$ expanded in the variables $X$.

\emph{One top site.} If $\mathit{TT}=\{s_0\}$ then $d_{\rm top}(s_0)=\infty$, because $d$ is even. With
$A=\nu_1(s_0)$ and $w_h=\nu_1(s_0+h)$ the bound reads
\begin{equation}\label{dv:onerow}
 v\big(P(\partial)\Theta(0)\big)\ \ge\ \sum_{h\ \mathrm{odd}}\Big(\frac A{\ell^h}-w_h\Big)_+
   -\gamma W(s_0)+F(D)-dN_0N_1(\log_\ell\bar a+1).
\end{equation}
This is \eqref{expo}: $A$ is the length of the first row at the chosen place, and $w_h$ is the number of
minors at the place reached by the $h$th power of the Frobenius followed by conjugation.

\subsection*{2.\quad Weights, and the inverse of the logarithm}

\textbf{Weights.} Let $f=\sum_nf_nX^n$ be a power series in finitely many variables $X_i$, with
coefficients in $K$. Let $r=(r_i)$ be real numbers $\ge0$ and $w$ a real number. Write $n\cdot r=\sum_in_ir_i$
and $|n|=\sum_in_i$. We say that $f$ has \emph{weight $w$ at radii $r$}, and write $f\in\mathcal W(r,w)$, if
$v(f_n)\ge w-n\cdot r$ for all $n$. This says that the Gauss norm of $f$ on the polydisc $v(X_i)\ge r_i$ is
at most $\ell^{-w}$. (Lean: \texttt{WtGe}.) The following rules are immediate from $v(x+y)\ge\min(v(x),v(y))$.
\begin{enumerate}[label=\textup{(W\arabic*)},leftmargin=3.2em]
\item If $f\in\mathcal W(r,w)$ and $h\in\mathcal W(r,w')$ then $fh\in\mathcal W(r,w+w')$. If $r\le r'$ then
$\mathcal W(r,w)\subset\mathcal W(r',w)$.
\item If $f\in\mathcal W(r,w)$ and the series $a_i$ have no constant term and $a_i\in\mathcal W(r',r_i)$,
then $f(a)\in\mathcal W(r',w)$. (A monomial $f_na^n$ lies in $\mathcal W(r',\,w-n\cdot r+n\cdot r)$.)
\item If $Y$ has no constant term and $Y\in\mathcal W(r,\frac1{\ell-1})$, then
$\exp(Y)-1$ and $\log(1+Y)$ lie in $\mathcal W(r,\frac1{\ell-1})$.
\item The coefficient of $X^\kappa$ in $fh$, in $\exp(Y)$ and in $\log(1+Y)$ depends only on the
coefficients of $X^n$ with $n\le\kappa$; so (W1) and (W3) hold with ``for all $n$'' replaced by ``for all
$n\le\kappa$''.
\end{enumerate}
For (W3): $Y^k\in\mathcal W(r,\frac k{\ell-1})$, and $v(k)\le v(k!)=(k-\sigma_\ell(k))/(\ell-1)\le(k-1)/(\ell-1)$,
where $\sigma_\ell(k)\ge1$ is the sum of the digits of $k$ in base $\ell$. So $Y^k/k!$ and $Y^k/k$ lie in
$\mathcal W(r,\frac1{\ell-1})$. (Lean: \texttt{MulWtLe}, \texttt{RadMono}, \texttt{SubstWt},
\texttt{LogStep}. The Lean statement \texttt{ExpWtLe} is not (W3): it gives weight $0$ for the exponential
of a series, truncated at $\kappa$, which is what Proposition~5 uses.)

\textbf{Lemma 1 (inverse of the logarithm).} \emph{Let $\mathcal Y$ be a finite set and
$\pi:\mathcal Y\to\mathcal Y$ a map. For $y\in\mathcal Y$ put}
\begin{equation}\label{dv:log}
 L_y(X)=\sum_{e\ge0}\ell^{-e}\,X_{\pi^e(y)}^{\ell^{2e}} .
\end{equation}
\emph{There is a unique family $g=(g_y)$ of power series without constant term with $L_y(g)=X_y$ for all
$y$. It satisfies $g_y\in\mathcal W(\rho,\rho)$, where the radius is $\rho$ for every variable:}
\begin{equation}\label{dv:ginv}
 v\big(\text{coefficient of }X^n\text{ in }g_y\big)\ \ge\ -\big(|n|-1\big)\rho .
\end{equation}
\emph{Proof.} The equation is $g_y=X_y-\sum_{e\ge1}\ell^{-e}g_{\pi^e(y)}^{\,\ell^{2e}}$. It determines the
coefficients by induction on $|n|$, since the terms with $e\ge1$ involve only coefficients of lower degree.
The bound holds for $|n|=1$. Let $|n|\ge2$ and $r=\ell^{2e}$, $e\ge1$. The coefficient of $X^n$ in an $r$th
power is a sum of products of $r$ coefficients, of degrees $|n_1|,\ldots,|n_r|\ge1$ with sum $|n|$. By
induction such a product has valuation at least $-(|n|-r)\rho$. Division by $\ell^e$ costs $e$, and
$e\le1+\ell^2+\cdots+\ell^{2(e-1)}=(r-1)\rho$. \qed

\noindent This is Part~II, \S2, in several variables. (Lean: \texttt{GenLogWeight},
\texttt{generalized\_log\_weight}; the series \eqref{dv:log} is \texttt{genLog}. Also $g(L(X))=X$:
\texttt{GenLogInverse}.) In one variable ($\pi$ the identity) a computation confirms
\eqref{dv:ginv} for $\ell=3,5,7$ and all degrees up to $801$, $1201$, $1201$: the minimum of
$v(g_n)+(n-1)\rho$ is $0$, attained exactly at $n=\ell^{2e}$. The bound cannot be improved. (See Proposition~Le15 of \S\ref{el:sec} for the refinement that includes the factorial.)

\subsection*{3.\quad From integrality to a bound for $\beta$}

In \eqref{dv:theta} the factor $U(g(Z))$ lies in $\mathcal W(\rho,0)$, by Lemma~1 and (W2). After
multiplication by $\kappa!$ this gives the exponent $|\kappa|\,(\frac1{\ell-1}-\rho)=|\kappa|\gamma$ of
\S I.3. For the exponential factor to do as well, by (W3) the series $\beta$ should have weight
$\frac1{\ell-1}$ at radii $\rho$. Integrality gives directly only the weight $\rho$. The gain from $\rho$ to
$\frac1{\ell-1}$ is the content of this paragraph. It uses that the degrees of the monomials are divisible by $\ell+1$. The proof given below also uses that the
coefficients lie in a field with value group $\Z$ (an unramified field); the source carries this hypothesis
(\texttt{Unram}, in \texttt{BHyp}), and it is kept as (B1). The statements do not need it: see the Remarks
after Theorem~J.

\textbf{Hypothesis (B).} Let $K_0\subset K$ be a subfield with $v(K_0^\times)=\Z$. Let $\zeta\in K$ be a
primitive root of unity of order $\ell^d-1$. For a site $s$, represented by $0\le s<d$, put
$\mathrm{lw}(s)=(-1)^s\ell^s$ (the letters are those of the Lean). The series $\beta$, in variables $X_s$ ($s\in\Z/d$), satisfies:
(B1) its coefficients lie in $K_0$; (B2) $\beta\big((\zeta^{\mathrm{lw}(s)}X_s)_s\big)=\beta(X)$; (B3) every
monomial of $\beta$ has degree $\ge1$ in the variables at even sites, degree $\ge1$ in the variables at odd
sites, and total degree $\ge3$; (B4) for the two series
\begin{equation}\label{dv:de}
\begin{aligned}
 \Delta_1\beta&=\beta(x_{\rm ev}+y_{\rm ev},x_{\rm od})-\beta(x_{\rm ev},x_{\rm od})-\beta(y_{\rm ev},x_{\rm od}),\\
 \Delta_2\beta&=\beta(x_{\rm ev},x_{\rm od}+y_{\rm od})-\beta(x_{\rm ev},x_{\rm od})-\beta(x_{\rm ev},y_{\rm od}),
\end{aligned}
\end{equation}
in which the variables at the even (resp.\ odd) sites are doubled, the series
$\exp(\Delta_i\beta)\circ L$ has all coefficients of valuation $\ge0$. Here $L$ is \eqref{dv:log} for the
doubled set of variables, with $\pi$ the shift $s\mapsto s-2$ of the site. (Lean: \texttt{BHyp};
$\Delta_1,\Delta_2$ are \texttt{de1}, \texttt{de2}; $L$ is \texttt{letterLog}.) Hypothesis (B4) says that
$\exp(\Delta_i\beta)$ becomes integral when written in the coordinates $T=g(X)$ of the formal group. It is
assumed here, from the source.

\textbf{Lemma 2 (quantization).} \emph{Let $c\in K_0$ and let $N$ be an integer divisible by $\ell+1$. If
$v(c)\ge\rho-N\rho$, then $v(c)\ge\frac1{\ell-1}-N\rho$.}

\emph{Proof.} Let $c\ne0$, $k=v(c)\in\Z$ and $N=(\ell+1)m$. Multiply $k\ge\rho-(\ell+1)m\rho$ by $\ell^2-1$:
\begin{equation}\label{dv:quant}
 (\ell+1)\big(k(\ell-1)+m\big)\ \ge\ 1 .
\end{equation}
The number $z=k(\ell-1)+m$ is an integer and $(\ell+1)z\ge1$; so $z\ge1$. Divide $z\ge1$ by $\ell-1$:
$k\ge\frac1{\ell-1}-\frac m{\ell-1}=\frac1{\ell-1}-N\rho$. \qed

\noindent Here $k$ is the valuation of the coefficient and $m$ is its degree divided by $\ell+1$. The
hypothesis loses the factor $\ell+1$ because $N\rho=m/(\ell-1)$ has denominator $\ell-1$ only. (Lean:
\texttt{Quant}, \texttt{weightQuantization}; the step $(\ell+1)z\ge1\Rightarrow z\ge1$ is at lines 83--91 of
\texttt{OneLetterRadiusIdentity.lean}. The Lean statement \texttt{Quant} has as hypothesis the congruence
of the torus, that $\ell^d-1$ divides $\sum_xn_x\,\mathrm{lw}(s_x)$ in the notation of the proof of Lemma~3;
it does not have the hypothesis that $\ell+1$ divides $N$.) Without the divisibility the conclusion fails: for $N=1$ the value
$k=0$ is allowed and $\frac1{\ell-1}-\rho>0$.

\textbf{Lemma 3.} \emph{Assume \textup{(B)}. Then $\Delta_1\beta$ and $\Delta_2\beta$ lie in
$\mathcal W(\rho,\frac1{\ell-1})$.}

\emph{Proof.} Let $\Delta$ be one of them and $E=\exp(\Delta)\circ L$. By (B3) $\Delta$ has no constant
term, so $E$ has constant term $1$. Since $L\circ g$ is the identity, $\exp(\Delta)-1=(E-1)\circ g$. By
(B4), Lemma~1 and (W2), a monomial of degree $j\ge1$ of $E-1$ contributes a series in
$\mathcal W(\rho,j\rho)\subset\mathcal W(\rho,\rho)$. So $\exp(\Delta)-1\in\mathcal W(\rho,\rho)$. Its
coefficients lie in $K_0$, by (B1). By (B2), $\Delta$ and $\exp(\Delta)$ are unchanged when each variable at
the site $s$ is multiplied by $\zeta^{\mathrm{lw}(s)}$. So a monomial $X^n$ with nonzero coefficient has
$\sum n_x\mathrm{lw}(s_x)\equiv0$ modulo $\ell^d-1$, where $s_x$ is the site of the variable $x$. Now $\ell+1$
divides $\ell^d-1$, because $d$ is even, and $\mathrm{lw}(s)\equiv(-1)^s(-1)^s=1$ modulo $\ell+1$. Hence $\ell+1$
divides $|n|$. Lemma~2 gives $\exp(\Delta)-1\in\mathcal W(\rho,\frac1{\ell-1})$, and (W3) gives the same
for $\Delta=\log(1+(\exp(\Delta)-1))$. \qed

\noindent (Lean: \texttt{asmS\_diff} in \texttt{ConstantCoeffEvenDerivativeEqZero.lean}, with
\texttt{ExpDeInv} for the divisibility.) A computation tests the chain of Lemmas 1--3 in one
variable: for $Y=\varphi(g(X))$ with $\varphi$ an integral series in $T^{\ell+1}$, and $\Delta=\log(1+Y)$,
the minimum over $n\le240$ of $v(\Delta_n)-(\frac1{\ell-1}-n\rho)$ is $0$ for $\ell=3,5,7$, attained at
$n=\ell+1$. For $\varphi(T)=T$ the bound fails (the minimum is $-11/8$, $-29/24$, $-55/48$).

\textbf{Lemma 4 (from the differences to $\beta$).} \emph{Assume \textup{(B)}. Call a monomial
\emph{exceptional} if it is $X_s^{\ell^\alpha}X_{s'}^{\ell^{\alpha'}}$ with $s$ even, $s'$ odd,
$(\alpha,\alpha')\ne(0,0)$ and $s+\alpha=s'+\alpha'$ in $\Z/d$. Then every coefficient $c_n$ of $\beta$
satisfies}
\begin{equation}\label{dv:beta}
 v(c_n)\ \ge\ \frac1{\ell-1}-|n|\rho-\mathrm{ex}(n),\qquad \mathrm{ex}(n)=1\ \text{if $X^n$ is exceptional},\quad
 \mathrm{ex}(n)=0\ \text{otherwise.}
\end{equation}
\emph{Proof.} Write $X^n=X_{\rm ev}^{\,u}X_{\rm od}^{\,u'}$ with multi-indices $u,u'$ on the even and the odd
sites. By \eqref{dv:de}, the coefficient of $x_{\rm ev}^{\,j}\,y_{\rm ev}^{\,u-j}\,x_{\rm od}^{\,u'}$ in
$\Delta_1\beta$ is $c_n\prod_s\binom{u_s}{j_s}$, for every $j$ with $0\le j\le u$, $j\ne0$, $j\ne u$. By
Lemma~3, $v(c_n)\ge\frac1{\ell-1}-|n|\rho-\mu(u)$, where $\mu(u)$ is the least valuation of such a product
of binomial coefficients; this holds if $|u|\ge2$. If $u$ involves two sites, take $j$ equal to $u$ at one
site and $0$ elsewhere: $\mu(u)=0$. Let $u$ be supported at one site, with $u_s$ not a power of $\ell$. If $u_s$ has at least two nonzero
digits in base $\ell$, let $j$ be its leading digit times the corresponding power of $\ell$; if
$u_s$ is $\ell^k$ times one digit between $2$ and $\ell-1$, let $j=\ell^k$. By Lucas's congruence $\binom{u_s}{j}$ is prime to
$\ell$, so $\mu(u)=0$. If $u_s=\ell^\alpha$ with $\alpha\ge1$ then $\mu(u)=1$, attained at $j=\ell^{\alpha-1}$. The same
holds for $u'$ and $\Delta_2\beta$. By (B3), $|u|\ge2$ or $|u'|\ge2$. So $\mathrm{ex}(n)=0$ unless $u$ and $u'$
are each a power of $\ell$ at one site, not both equal to $1$. In that case
$X^n=X_s^{\ell^\alpha}X_{s'}^{\ell^{\alpha'}}$, and if $c_n\ne0$ then (B2) gives
$\ell^{s+\alpha}\equiv\ell^{s'+\alpha'}$ modulo $\ell^d-1$, that is $s+\alpha=s'+\alpha'$ in $\Z/d$. \qed

\noindent (Lean: the statement is \texttt{LemmaS}, proved as \texttt{lemmaS}; the step of Lemma~4 is
\texttt{BinomExtract}. The proof of Lemma~4 above is ours; the Lean proof of \texttt{BinomExtract} was not
read. The facts on binomial coefficients were checked for $u_s\le1500$.)
So $\beta$ has the weight $\frac1{\ell-1}$ that is needed, except at the exceptional monomials, where one
unit is missing. The exceptional monomials pair a site $s$ with the site $s'=s+(\alpha-\alpha')$, at an odd
distance. They are the reason for the weights of the next paragraph.

\subsection*{4.\quad One Taylor coefficient, with weights on the variables}

\textbf{Proposition 5.} \emph{Assume \textup{(B)}. Let $X_i$ $(i\in\iota)$ be finitely many variables. Let
$Z_y=\sum_iM_{y,i}X_i$ be linear forms with $v(M_{y,i})\ge0$, one for each variable $Z_y$ of
\eqref{dv:theta}, and let $\Theta$ be \eqref{dv:theta} written in the variables $X$. Let
$\mathcal B=\sum_{q,c}\beta\big((Z_{s,q,c})_s\big)=\sum_n\mathcal B_nX^n$. Let $\kappa$ be a multi-index and
let $r_i\ge0$ be real numbers such that}
\begin{equation}\label{dv:cond}
 \ell^\alpha r_i+\ell^{\alpha'} r_j\ \ge\ 1\qquad\text{whenever }X_i^{\ell^\alpha}X_j^{\ell^{\alpha'}}\text{ divides }
 X^\kappa\text{ and }v\big(\mathcal B_{\ell^\alpha e_i+\ell^{\alpha'} e_j}\big)<\tfrac1{\ell-1}-(\ell^\alpha+\ell^{\alpha'})\rho .
\end{equation}
\emph{The case $i=j$ is included; the monomial is then $X_i^{\ell^\alpha+\ell^{\alpha'}}$. Then, with
$\lambda(k)=\lfloor\log_\ell k\rfloor+1$ for $k\ge1$ and $\lambda(0)=0$,}
\begin{equation}\label{dv:jet}
 v\big(\kappa!\,\Theta_\kappa\big)\ \ge\ \sum_i\kappa_i\,(\gamma-r_i)\ -\ \sum_i\lambda(\kappa_i).
\end{equation}
\emph{Proof.} First, $\mathcal B_n$ has valuation $\ge\frac1{\ell-1}-|n|\rho$ unless $n=\ell^\alpha e_i+\ell^{\alpha'} e_j$
for some $i,j,\alpha,\alpha'$ ($i=j$ allowed); and $\ge\frac1{\ell-1}-1-|n|\rho$ always. Indeed, by \eqref{dv:beta} and (W2)
only the exceptional monomials $c\,X_s^{\ell^\alpha}X_{s'}^{\ell^{\alpha'}}$ of $\beta$ matter. For a linear form
$Z=\sum_iM_iX_i$ one has $Z^{\ell^\alpha}=\sum_iM_i^{\ell^\alpha}X_i^{\ell^\alpha}+\ell R$, with $R$ a polynomial with
coefficients of valuation $\ge0$. So all monomials of $Z^{\ell^\alpha}Z'^{\ell^{\alpha'}}$ other than the
$X_i^{\ell^\alpha}X_j^{\ell^{\alpha'}}$ carry a factor $\ell$. Now put $r'_i=\rho+r_i$. For $n\le\kappa$ we get
$v(\mathcal B_n)\ge\frac1{\ell-1}-n\cdot r'$: this is clear if $v(\mathcal B_n)\ge\frac1{\ell-1}-|n|\rho$, and
otherwise $n=\ell^\alpha e_i+\ell^{\alpha'} e_j$ and $n\cdot r'=|n|\rho+\ell^\alpha r_i+\ell^{\alpha'} r_j\ge|n|\rho+1$ by
\eqref{dv:cond}. By (W3) and (W4) the coefficients of $X^n$ in $\exp(\pm\mathcal B)$, $n\le\kappa$, have
valuation $\ge-n\cdot r'$. The factor $U(g(Z))$ lies in $\mathcal W(\rho,0)\subset\mathcal W(r',0)$ by
Lemma~1 and (W2), (W1). By (W1), $v(\Theta_\kappa)\ge-\kappa\cdot r'$. Finally
$v(k!)=(k-\sigma_\ell(k))/(\ell-1)\ge\frac k{\ell-1}-\lambda(k)$, and $\frac1{\ell-1}-\rho=\gamma$. \qed

\noindent (Lean: \texttt{LemmaR}, with \texttt{PureCross} for the first step and \texttt{FactVal} for the
last; the Lean loss is $|\iota|\,(\log_\ell|\kappa|+1)$. \texttt{LemmaR} has one more hypothesis than
Proposition~5: each variable $X_i$ belongs to one site. Proposition~5 does not need it. The statement \texttt{LemmaR} is in the archive,
proved as \texttt{lemmaR\_of}, but the dependency computation shows that \texttt{LemmaR} is not among the
declarations the final theorem depends on. Proposition~5 is proved here independently. By occurrences of
names the chain uses \texttt{LemmaRA}.) Each variable pays its radius $r_i$ out of the gain
$\gamma$. The radii are needed only where an exceptional monomial is present with a coefficient that is not
divisible by one more power of $\ell$.

\textbf{One variable.} Take one variable $X$ and $\Theta=U(g(X))\exp(\mathcal B(X))$, where
$v(\mathcal B_n)\ge\frac1{\ell-1}-n\rho$ for $n\notin S$ and $v(\mathcal B_n)\ge\frac1{\ell-1}-1-n\rho$ for
$n\in S=\{\ell^\alpha+\ell^{\alpha'}:\alpha+\alpha'\text{ odd}\}$. The second half of the proof of Proposition~5 applies to such a series: what is needed is
$nr\ge1$ for $n\in S$, $n\le\kappa$. The least element of $S$ is $\ell+1$, so $r=\frac1{\ell+1}$ will do,
and $\gamma-\frac1{\ell+1}=\rho$. Thus
\begin{equation}\label{dv:onevar}
 v\big(\kappa!\,\Theta_\kappa\big)\ \ge\ \kappa\gamma-\lambda(\kappa)\ \ \text{if }\mathcal B\text{ has no terms in }S,
 \qquad v\big(\kappa!\,\Theta_\kappa\big)\ \ge\ \kappa\rho-\lambda(\kappa)\ \ \text{in general.}
\end{equation}
In the first case the factorial contributes $\kappa/(\ell-1)=\kappa(\gamma+\rho)$ and the function costs
$\kappa\rho$. In the second the exceptional terms can cost a further $\kappa/(\ell+1)$. (Lean, for the second:
\texttt{OneLetterRadius}. The Lean statement \texttt{OneLetterR} is for several variables through an
integral matrix, with exponent $(\mathrm{sgn}/2)\sum\beta(Z,\mathrm{ep}\cdot Z)$.) A computation tests \eqref{dv:onevar} with exact
rational arithmetic for $\ell=3,5,7$ and all $\kappa\le240$, with $U$ a random series with integer
coefficients that are units at $\ell$, and $\mathcal B_n$ of the least valuation allowed. The minimum over $\kappa$ of the left side
minus the right side is $\frac14,\frac16,\frac18$ in the first case and $\frac34,\frac56,\frac78$ in the
second. In the second case the first bound fails: at $\kappa=240$ the left side is below
$\kappa\gamma-\lambda(\kappa)$ by $57$, $37$, $28$. So the cost $\kappa/(\ell+1)$ is real in this model
($240/(\ell+1)=60,40,30$).

\subsection*{5.\quad The choice of the weights}

In Proposition~5 the bound is $\sum_i\kappa_i(\gamma-r_i)$, and the radii are free subject to
\eqref{dv:cond}. The exponent \eqref{dv:B} comes from a choice of radii. Two elementary optimisations are
involved.

\textbf{Lemma 6 (the top radius).} \emph{Let $A\ge0$, let $c_j\ge0$ for odd $j\ge1$, and let $J\ge1$ be odd.
There is $\omega\ge\ell^{-J}$ with}
\begin{equation}\label{dv:top}
 A\omega+\sum_{\substack{j\ \mathrm{odd}\\ j<J}}c_j\,(1-\ell^j\omega)_+\ \le\
 \sum_{\substack{j\ \mathrm{odd}\\ j<J}}\min\Big(c_j,\frac A{\ell^j}\Big)+\frac A{\ell^J}.
\end{equation}
\emph{Consequently}
\begin{equation}\label{dv:topgain}
 A(\gamma-\omega)-\sum_{\substack{j\ \mathrm{odd}\\ j<J}}c_j\,(1-\ell^j\omega)_+\ \ge\
 \sum_{\substack{j\ \mathrm{odd}\\ j<J}}\Big(\frac A{\ell^j}-c_j\Big)_++\sum_{\substack{j\ \mathrm{odd}\\ j>J}}\frac A{\ell^j}.
\end{equation}
\emph{Proof.} Let $j_0$ be the least odd $j<J$ with $c_j\ge A/\ell^j$, and $j_0=J$ if there is none. Take
$\omega=\ell^{-j_0}$. On the left of \eqref{dv:top} the terms with $j\ge j_0$ vanish, and each term
with $j<j_0$ is at most $c_j=\min(c_j,A/\ell^j)$. The term $A\omega=A/\ell^{j_0}$ equals
$\min(c_{j_0},A/\ell^{j_0})$ if $j_0<J$, and is the last term on the right if $j_0=J$. For
\eqref{dv:topgain} use $\gamma A=\sum_{j\ \mathrm{odd}}A/\ell^j$ and $x-\min(c,x)=(x-c)_+$. \qed

\noindent (The archive contains this statement as \texttt{TopRadius}, proved as
\texttt{top\_radius\_holds} with the same choice of $\omega$. The dependency computation shows that neither
is among the declarations the final theorem depends on. Lemma~6 is proved here independently. The declaration
the final theorem uses for the top-row term has been identified by reading the proof terms: it is \texttt{asmA\_top}, which the proof of \texttt{thmAOp} reaches through \texttt{thmAJetA} and \texttt{rowRadiusAssembly}; the bound for one Taylor coefficient on that route is \texttt{lemma\_R\_log\_loss}. The lemma was checked on
$20000$ random cases.) Here $\omega$ is the radius of the variables of a top row of length
$A$, and $c_j$ bounds the degree in those variables at the site $s+j$ which are linked to the top row by an
exceptional monomial. Such a variable needs the radius $(1-\ell^j\omega)_+$, by \eqref{dv:cond}. The
right side of \eqref{dv:topgain}, with both of its sums, is $\mathrm{top}(s)$ of \eqref{dv:B} for $A=\nu_1(s)$,
$c_j=\mathrm{cap}(s+j)$ and $J=d_{\rm top}(s)$: the second sum of \eqref{dv:topgain}, over odd $j>J$, is the
second sum in the definition of $\mathrm{top}(s)$. The lemma asserts that a suitable $\omega$ exists, not that
\eqref{dv:topgain} holds for every $\omega\ge\ell^{-J}$. If
$d_{\rm top}(s)=\infty$ and only the $j$ with $\ell^j\le A$ occur on the left, the same choice (with
$\omega=0$ if there is no $j_0$) gives $\sum_{j\ \mathrm{odd}}(A/\ell^j-c_j)_+$ on the right.

\textbf{The linear program.} Let $D:\Z/d\to\mathbb R_{\ge0}$. A \emph{flow} is a family $\psi(s,i)\ge0$, for
sites $s$ and odd $i<d$, with
\begin{equation}\label{dv:flow}
 \sum_{i}\ell^i\,\psi(s,i)+\sum_{i}\psi(s-i,i)\ \le\ D(s)\qquad\text{for all }s .
\end{equation}
Let $\mathrm{maxFlow}(D)$ be the supremum of $\sum_{s,i}\psi(s,i)$ over flows, and
$F(D)=\gamma\sum_sD(s)-\mathrm{maxFlow}(D)$. A \emph{dual} is a family $0\le\sigma(s)\le1$ with
$\ell^i\sigma(s)+\sigma(s+i)\ge1$ for all $s$ and all odd $i<d$. (Lean: \texttt{IsFlow}, \texttt{Fnum},
\texttt{IsDual}; and \texttt{lpF\_eq\_Fnum} for the form used elsewhere.)

\emph{Weak duality, proved here.} For a flow $\psi$ and a dual $\sigma$,
\begin{equation}\label{dv:weak}
 \sum_{s,i}\psi(s,i)\ \le\ \sum_{s,i}\psi(s,i)\big(\ell^i\sigma(s)+\sigma(s+i)\big)
 =\sum_s\sigma(s)\Big(\sum_i\ell^i\psi(s,i)+\sum_i\psi(s-i,i)\Big)\ \le\ \sum_s\sigma(s)D(s).
\end{equation}
So $\sum_sD(s)(\gamma-\sigma(s))\le F(D)$ for every dual $\sigma$.

\emph{Strong duality, cited.} For $\ell\ge3$ and $d$ even there is a dual $\sigma$ with
$\sum_sD(s)(\gamma-\sigma(s))=F(D)$, and with $\ell^i\sigma(s)+\sigma(s+i)\ge1$ for \emph{all} odd $i\ge1$
(Lean: \texttt{strongDuality}, of type \texttt{LemmaFa}). For a dual with the last property, if
$s+\alpha=s'+\alpha'$ and $\alpha+\alpha'$ is odd, then
\begin{equation}\label{dv:rad}
 \ell^\alpha\sigma(s)+\ell^{\alpha'}\sigma(s')\ \ge\ 1:
\end{equation}
if $\alpha>\alpha'$ this is $\ell^{\alpha'}$ times $\ell^{\alpha-\alpha'}\sigma(s)+\sigma(s+(\alpha-\alpha'))\ge1$;
if $\alpha'>\alpha$ exchange the roles of $(s,\alpha)$ and $(s',\alpha')$ (and $\alpha\ne\alpha'$, since
$\alpha+\alpha'$ is odd).
So such a $\sigma$, used as radius for all variables at the site $s$, satisfies \eqref{dv:cond}.

\emph{Orbits of length two, proved here.} Let $d=2$. Then
\begin{equation}\label{dv:d2}
 \mathrm{maxFlow}(D)=\min\Big(D_0,\,D_1,\,\frac{D_0+D_1}{\ell+1}\Big),\quad
 F(D)=\begin{cases}\rho\,(D_0+D_1),&D_1/\ell\le D_0\le\ell D_1,\\
 \gamma D_0-(1-\gamma)D_1,&\ell D_1\le D_0,\\ \gamma D_1-(1-\gamma)D_0,&\ell D_0\le D_1.\end{cases}
\end{equation}
Indeed, in the first case $\psi=\big(\frac{\ell D_0-D_1}{\ell^2-1},\frac{\ell D_1-D_0}{\ell^2-1}\big)$ is a flow
and $\sigma=(\frac1{\ell+1},\frac1{\ell+1})$ a dual, both of value $\frac{D_0+D_1}{\ell+1}$; in the second
$\psi=(D_1,0)$ and $\sigma=(0,1)$, of value $D_1$; the third is symmetric. By \eqref{dv:weak} these are optimal.
These duals satisfy the inequalities for all odd $i$. (Checked exactly on $20000$ cases, and strong duality
numerically for $d=2,4,6$.) Note that $F(D)$ is \emph{not} $\sum_{h\ \mathrm{odd}}(D_0/\ell^h-D_1)_+$:
for $\ell=5$ and $D=(1000,10)$ the two are $200.42$ and $190$; for $D=(1000,250)$ they are $52.08$ and $0$.
One has $F(D)\ge\sum_{h\ \mathrm{odd}}(D_0/\ell^h-D_1)_+$, with equality only for $D_1=0$. In the source the
sum of the shape \eqref{dv:V} is the term $\mathrm{top}(s)$. Here it comes from Lemma~6; on the route of the final theorem it is the statement \texttt{asmA\_top}. The linear program
governs the rows other than the top rows.

\textbf{Theorem J.} \emph{Assume \textup{(B)}, with $\ell$ an odd prime and $d\ge2$ even. Let a shape, a set
$\mathit{TT}$, the polynomial $P$ and the integer $\bar a$ be as in \S1. Call a row $(s,i)$ \emph{ordinary} if it
is not a top row. Call a top site $s_0$, an odd integer $j$ and an ordinary row $(s,i)$ \emph{linked} if $s=s_0+j$, and
$\ell^j\le\nu_1(s_0)$ if $j>0$, and $\ell^{-j}\le\mathrm{size}(s)$ if $j<0$. Let $\mathrm{spec}$ be an
injective map from a set of ordinary rows to pairs $(s_0,j)$, such that $s_0$, $j$, $(s,i)$ are linked
whenever $\mathrm{spec}(s,i)=(s_0,j)$; the pair is the \emph{label} of the row. Here $j$ is an integer, not a class modulo $d$: $(s_0,j)$ and
$(s_0,j+d)$ are different labels for the same site. Suppose that whenever $s_0$,
$j$, $(s,i)$ are linked and $(s_0,j)$ is not the label of $(s,i)$, the frame satisfies}
\begin{equation}\label{dv:frame}
 v\Big(\sum_{q}\xi(s_0,0,q)^{\ell^{j}}\,\xi(s,i,q)\Big)\ge1\quad(j>0),\qquad
 v\Big(\sum_{q}\xi(s_0,0,q)\,\xi(s,i,q)^{\ell^{-j}}\Big)\ge1\quad(j<0).
\end{equation}
\emph{Then $v\big(P(\partial)\Theta(0)\big)\ge B'$, where
$B'=\sum_{s\in\mathit{TT}}\big(\mathrm{top}(s)-\gamma W(s)\big)+F(D)-\max_\kappa\sum_y\lambda(\kappa_y)$, the
maximum being over the monomials $X^\kappa$ of $P$ and the sum over the variables $y$. This is the loss that
the proof gives. Since $\sum_y\lambda(\kappa_y)\le dN_0N_1(\log_\ell\bar a+1)$, one has $B'\ge B$ with $B$ as in
\eqref{dv:B}, which is the exponent of the Lean statement. For $d=2$ this is proved
here; for $d\ge4$ the proof uses the cited strong duality.}

\emph{Proof.} Since $P$ has integer coefficients, it suffices to prove $v(\kappa!\Theta_\kappa)\ge B'$ for
each monomial $X^\kappa$ of $P$. A determinant is linear in each row. So the degree of $X^\kappa$ in the
variables of the row $(s,i)$ is the number of $c$ with $\mathrm{len}(s,c)>i$. It is $\nu_1(s)$ for $i=0$, and at most
$\nu_2(s)$ for $i\ge1$. Hence the degree is $\nu_1(s)$ in a top row, at most $\mathrm{cap}(s)$ in an
ordinary row, and $D(s)$ in all ordinary rows at $s$ together; and each $\kappa_y\le\bar a$.

\emph{Radii.} Let $\sigma$ be an optimal dual for $D$, as above. For $s\in\mathit{TT}$ let $\omega(s)$
be given by Lemma~6 with $A=\nu_1(s)$, $c_j=\mathrm{cap}(s+j)$ and $J=d_{\rm top}(s)$ (and by the remark
after it if $d_{\rm top}(s)=\infty$). Give the variables of a top row the radius $\omega(s)$. Give the
variables of an ordinary row $(s,i)$ the radius $\sigma(s)$ if the row has no label;
$\max(\sigma(s),1-\ell^j\omega(s_0))$ if its label is $(s_0,j)$ with $j>0$; and
$\max(\sigma(s),\ell^{-|j|})$ if $j<0$. (Lean: \texttt{rowRad}.)

\emph{The condition \eqref{dv:cond}.} With $Z_{s,q,c}=\sum_i\xi(s,i,q)X_{s,i,c}$, a monomial
$X_y^{\ell^\alpha}X_{y'}^{\ell^{\alpha'}}$, $y=(s,i,c)$, $y'=(s',i',c')$, has in $\mathcal B$ the coefficient
$c_n\sum_q\xi(s,i,q)^{\ell^\alpha}\xi(s',i',q)^{\ell^{\alpha'}}$ if $c=c'$, and $0$ otherwise, where $c_n$ is the
coefficient of $X_s^{\ell^\alpha}X_{s'}^{\ell^{\alpha'}}$ in $\beta$ (it is $0$ by (B3) if $s$ and $s'$ have the
same parity). By \eqref{dv:beta} the condition concerns
only the case $s+\alpha=s'+\alpha'$, $\alpha+\alpha'$ odd, and only if the sum over $q$ has valuation $<1$.
If both rows are ordinary, their radii are at least $\sigma(s)$, $\sigma(s')$, and \eqref{dv:rad} applies.
If both are top rows, say with $\alpha>\alpha'$, then $d_{\rm top}(s)\le\alpha-\alpha'$, so
$\ell^\alpha\omega(s)\ge\ell^\alpha\ell^{-(\alpha-\alpha')}\ge1$. There remains the case of a top row and an ordinary row; by symmetry let $(s,0)$ be the top row and
$(s',i')$ the ordinary row, and $j=\alpha-\alpha'$, so $s'=s+j$. Let $j>0$. Since the monomial divides $X^\kappa$,
$\ell^j\le\ell^\alpha\le\nu_1(s)$. If the label of $(s',i')$ is $(s,j)$, then
$\ell^\alpha\omega(s)+\ell^{\alpha'}(1-\ell^j\omega(s))=\ell^{\alpha'}\ge1$. If not, \eqref{dv:frame}
holds, and $\sum_q\big(\xi(s,0,q)^{\ell^j}\xi(s',i',q)\big)^{\ell^{\alpha'}}\equiv
\big(\sum_q\xi(s,0,q)^{\ell^j}\xi(s',i',q)\big)^{\ell^{\alpha'}}$ modulo $\ell$ has valuation $\ge1$: the condition
is void. Let $j<0$. Then $\ell^{|j|}\le\ell^{\alpha'}\le\mathrm{size}(s')$. If the label is $(s,j)$ then
$\ell^{\alpha'}\ell^{-|j|}=\ell^\alpha\ge1$. If not, \eqref{dv:frame} applies in the same way.

\emph{The sum.} By Proposition~5 it remains to show $\sum_y\kappa_y(\gamma-r_y)\ge
\sum_{s\in\mathit{TT}}(\mathrm{top}(s)-\gamma W(s))+F(D)$. Use $\max(\sigma,u)\le\sigma+u_+$. The terms with
$\sigma$ give $\sum_sD(s)(\gamma-\sigma(s))=F(D)$. The labels $(s_0,j)$ with $j>0$ cost at most
$\mathrm{cap}(s_0+j)(1-\ell^j\omega(s_0))_+$ each, because $\mathrm{spec}$ is injective and one row has
degree at most $\mathrm{cap}$; with the top row this gives $\mathrm{top}(s_0)$ by \eqref{dv:topgain}. The
labels $(s_0,j)$ with $j<0$ cost at most $\mathrm{cap}(s_0-|j|)\,\ell^{-|j|}\le W(s_0)\,\ell^{-|j|}$, and
$\sum_{j\ \mathrm{odd}}\ell^{-j}=\gamma$. The loss in \eqref{dv:jet} is $\sum_y\lambda(\kappa_y)$, as in $B'$. It is at most
$dN_0N_1(\log_\ell\bar a+1)$, because there are $dN_0N_1$ variables and each $\kappa_y\le\bar a$. \qed

\noindent (Lean: \texttt{ThmAJetA} for one $\kappa$, \texttt{ThmAAssembly} for the sum, \texttt{ThmARadii}
for the condition, \texttt{ThmAOp} for the result. The Lean statement \texttt{ThmAOp} is more general: it
allows an integral linear operator on the ordinary rows before the minors are formed. That generality is
not proved here. The inequality for the sum was tested on $3000$ random shapes.)

\textbf{Remarks.} The following are observations of the test of this section. They are recorded
without proof, and nothing above depends on them.

\emph{(1) The field $K_0$.} The field $K_0$, hypothesis (B1) and Lemma~2 are not used in Lemmas 3--4,
Proposition~5 and Theorem~J, except in the proof of Lemma~3 as written. The reason is this. In that proof
the series $E=\exp(\Delta)\circ L$ is invariant under the torus, because $L$ commutes with it. Hence every
monomial of $E-1$ has degree $j$ divisible by $\ell+1$, so $j\ge\ell+1$, and it contributes a series in
$\mathcal W(\rho,j\rho)\subset\mathcal W(\rho,(\ell+1)\rho)=\mathcal W(\rho,\frac1{\ell-1})$. This is the
conclusion of Lemma~3. The Lean statement \texttt{BinomExtract} has no $K_0$ either.

\emph{(2) Strong duality.} The strong duality cited from the source is the duality theorem of linear
programming together with two elementary reductions. First, replacing $\sigma$ by $\min(\sigma,1)$ keeps the
constraints and does not increase $\sum_s\sigma(s)D(s)$. Second, for odd $i\ge d$ write $i=i'+kd$ with $i'$
odd and $1\le i'<d$; then $\ell^i\sigma(s)+\sigma(s+i)\ge\ell^{i'}\sigma(s)+\sigma(s+i')\ge1$, so every dual
satisfies the inequalities for all odd $i\ge1$. The equality was confirmed in every instance computed
($d=2,\ldots,10$).

\emph{(3) The prime $2$.} That $\ell$ is odd is not used in the proofs of this section, and Lemma~1 holds
for $\ell=2$. No series with (B) was constructed for $\ell=2$, so Lemmas 3--4 and Theorem~J were not tested
for $\ell=2$.

\begin{illbox}
The example below, and its completion in \S\ref{el:sec}, is ours and not the source's; it is not used in the
proof.
\end{illbox}
\textbf{An example.} The test of this section constructed an explicit series $\beta$ for orbits of length
two ($d=2$). Let $\ell\equiv3\bmod4$. Let $E_c$ be the curve $Y^2=X^3-cX$, with $c$ in $\Z[i]$ and
$c^{(\ell^2-1)/4}\equiv-1$ modulo $\ell$ ($c=-i$ for $\ell=3$, $11$ and $19$, $c=-1-i$ for $\ell=7$). It is a twist
of the lemniscate curve, the curve of the lemniscate sine of \S I.3. Let $\sigma_c$ be the Weierstrass
$\sigma$-function of $E_c$ (with $g_2=4c$, $g_3=0$), let $t=-X/Y$, and put
$m(z)=\log\big(\sigma_c(z)/t(z)\big)$ and $\tilde\beta(z,w)=m(z+w)-m(z)-m(w)$. The series $\beta(X_0,X_1)$ is
the part of $\tilde\beta$ that is invariant under $(X_0,X_1)\mapsto(\zeta X_0,\zeta^{-\ell}X_1)$; the field
is $K_0=\Q_\ell(i)$. Hypotheses (B1)--(B4) hold for this $\beta$ in all degrees: this is proved in \S\ref{el:sec}
(Proposition~Le22). A computation in exact arithmetic agrees with it: the series $\exp(\Delta_1\beta)\circ L$ has
coefficients integral at $\ell$ up to total degree $84$ for $\ell=3$, $72$ for $\ell=7$ and $144$ for
$\ell=11$. On this $\beta$, Lemmas~3 and~4 were checked up to degree $800$, and
equality holds in \eqref{dv:beta} at exceptional monomials: there the valuation of the coefficient is
exactly $\frac1{\ell-1}-|n|\rho-1$. So the tested coefficients attain the loss of one unit. Since (B) holds for this $\beta$ in all degrees (Proposition~Le22 of \S\ref{el:sec}), the term $\mathrm{ex}(n)$ cannot be omitted: equality holds at every exceptional monomial of degree $\le800$. That equality holds at all exceptional monomials is not proved. Theorem~J was checked on
this $\beta$ in exact arithmetic, for minors of size one, with $\ell=3,7,11,19$: the bound holds, both with
$B$ and with $B'$, in each of $1\,041\,934$ comparisons, of which $147\,296$ have $B>0$ and $266\,970$ have
$B'>0$. It also holds in every case computed with minors of size two at one site.

\textbf{Comparison with \S I.4.3.} The description there is \eqref{dv:onerow}. The weight $w_h$ is
$\mathrm{cap}(s_0+h)$: the number of minors at the place reached by the $h$th power of the Frobenius
composed with conjugation. It enters in two steps. First, an exceptional monomial of $\beta$ links the top
row at $s_0$ to a row at $s_0+h$, by \eqref{dv:beta}; only such a link forces a radius, by \eqref{dv:cond}.
Second, a monomial of a product of minors has degree at most $\mathrm{cap}$ in one row; this is the
statement of \S I.4.3 that a minor absorbs one derivative in a linked direction. Lemma~6 then balances the
radius of the top row against the radii of the linked rows, and the result is $(A/\ell^h-w_h)_+$ at each
odd $h$. The labels with \eqref{dv:frame} express that at each $h$ only one row at the place $s_0+h$ is
linked. The congruence of Part~II, \S6.3 is a statement of the same kind as \eqref{dv:frame}. The source does not
use it by name at this point: see \S6\,(ii).

\subsection*{6.\quad Assumed here, from the source}

\begin{enumerate}[label=\textup{(\roman*)},leftmargin=2.4em]
\item \emph{The form \eqref{dv:theta} of the theta kernels, and hypothesis \textup{(B)}.} In the Lean this is
the statement \texttt{Loc.AdmB}, of the file \texttt{ProductZerothFirstParts.lean}. It says that there are a unit $e$ and a series $\beta$ such
that the kernel, in the coordinates $e\cdot L$, equals $U\cdot\exp(\pm\sum\beta)\circ L$ with $U$ integral
(clause \texttt{B0}; also \texttt{OrbitThetaExpansion}), with coefficients in the fraction field of the
Witt ring (\texttt{B1}), torus invariance (\texttt{B2}), the condition on degrees (\texttt{B3}), and the
integrality (B4) (\texttt{B4}). The hypothesis \texttt{hServ.pdv} (\texttt{SelectionServed}) of the final theorem does not contain
\texttt{Loc.AdmB} literally: it gives \texttt{BeOK} (B1--B4 for each orbit representative) and
\texttt{ThetaExpansionAllPlaces}. How these give \texttt{Loc.AdmB} was not read. Nothing of it is proved
here. The passage from
\texttt{Loc.AdmB} to \texttt{BHyp} (\texttt{tieB}) was not read.
\item \emph{The frames.} Theorem~J needs a labelling and \eqref{dv:frame}. In the Lean the theorem for all
frames at once is \texttt{ThmAGauss}: every coefficient of the polynomial \texttt{frameJet} in the frame
variables has valuation $\ge B$, if $2\log_\ell\bar a+2\le N_0-1$. Its proof (\texttt{thmAGauss\_of}) names
\texttt{residuallyGenericExists}: there are frames whose entries are residually independent units
(\texttt{IsGeneric}). Three further inputs are supplied to it as hypotheses, in \texttt{thmAGauss}: a transfer
of the hypotheses between fields (\texttt{admissibleTransfer}, of type \texttt{TransferB}), the evaluation
of the frame polynomial (\texttt{FrameJetEval}, \texttt{GaussPow}), and the existence of labels
(\texttt{labelsExist\_statement}, of type \texttt{LabelsExist}). The statement \texttt{TransferB} contains no
frame: it transports \texttt{Unram}, \texttt{BHyp}, the inverse $g$ of the logarithm and the integrality of
$U$ along a norm-preserving homomorphism of fields. The proofs of these inputs were not read. The congruence
\texttt{key\_minor\_dvd} of Part~II, \S6.3 is not used by name in this step: no file under
\texttt{Formalism/GaussJetBounds} mentions it. Its only use by name is in its own file; that file is
imported by \texttt{Loc4ResOfLine} and \texttt{LiftSubstStart}, in the directory \texttt{CoarseMultiplicity},
which were not read.
\item \emph{From Taylor coefficients to the statement used.} The comparison of $F(D)$ and $\mathrm{top}(s)$
with the weights of the construction (\texttt{normCor7}; \texttt{split\_bound\_eleven\_one}, conditional on
\texttt{Loc.AdmB}); the factorial count \texttt{StarArith} and the inequality \texttt{BCore}, which is of the form \eqref{II:t6.1}
of Part~II; the identification of sites with places and of $\nu_1$ with the weights. Then the loss terms
$\mathrm{lo}_G$, $\mathrm{fl}_G$ of \eqref{dv:fl}; the lattice clauses of \texttt{LinkG\_at};
\texttt{stLoss7}; \texttt{placesF\_min'}. For \texttt{stLoss7} and \texttt{placesF\_min'} only the statements
were read. On the route from \texttt{LinkG} to the final theorem that was read, the clauses about
valuations of kernels are proved for the zero chain (\texttt{linkG\_without\_kernels}); how \texttt{ThmAOp}
enters the final theorem was not read. The dependency computation shows that \texttt{thmAOp} is among the
declarations the final theorem depends on. By occurrences of names it reaches the final theorem through \texttt{thmAGauss}, \texttt{thmA\_of},
\texttt{split\_bound\_eleven\_one} and \texttt{sum\_over\_primes\_service}, that is, through
\texttt{SelectionServed}.
\item \emph{The other primes.} The primes of class \texttt{inertClass} are at least $13$ for large $x$
(\texttt{thirteen\_le}, which has a hypothesis that $x$ is large and uses the assumed splitting of the
small primes, \texttt{split\_small}). The primes of the
other classes, and the primes above $X''$, are treated by other clauses, not discussed here.
\end{enumerate}

\textbf{What this section proves, and what it assumes.} \emph{Proved outright:} Lemma~1 (the inverse
of the logarithm), Lemma~2 (integrality), Lemma~6 (the optimisation for the top row), weak duality
\eqref{dv:weak}, and the closed form \eqref{dv:d2} for orbits of length two. \emph{Proved conditionally on
hypothesis \textup{(B)}:} Lemmas~3 and~4, Proposition~5, and Theorem~J for $d=2$. \emph{Conditional on
\textup{(B)} and on strong duality, which is cited from the source} (\texttt{strongDuality}): Theorem~J for
$d\ge4$; this case was not tested end to end. \emph{Taken from the source without proof:} the transcription
of the top-level inequality \texttt{LinkG\_at} in \S1; that the theta kernels satisfy (B) (item (i)); the
dictionary between sites and places; the frames and labels (item (ii)); the passage from jets to lattices
(item (iii)); the remaining primes (item (iv)). Theorem~J is proved without the operator on the ordinary
rows that the Lean statement \texttt{ThmAOp} allows.

\emph{Relation to the declarations the final theorem depends on.} By the dependency computation, the
following declarations named or used in this section are among them: \texttt{strongDuality},
\texttt{weightQuantization}, \texttt{lemmaS\_of}, \texttt{thmAOp},
\texttt{binomial\_coefficient\_extraction}, \texttt{lpF\_eq\_Fnum}, \texttt{residuallyGenericExists},
\texttt{theorem\_A\_generic} and \texttt{gauss\_sum\_powers}. The four statements \texttt{TopRadius},
\texttt{top\_radius\_holds}, \texttt{LemmaR} and \texttt{ThmABound} are in the archive but are not used by
the final theorem. So Lemma~6 and the weighted bound for one Taylor coefficient (Proposition~5) are proved
here independently of the route of the final theorem, and the declarations the final theorem uses instead are \texttt{asmA\_top} and \texttt{lemma\_R\_log\_loss}. Their statements are consistent with Lemma~6 and Proposition~5: for a finite cut-off, and for the set of all odd $j<J$, \texttt{asmA\_top} is \eqref{dv:topgain} with $\mathrm{top}(s)$ moved to the other side, and its proof takes the same radius (for an infinite cut-off it is the remark after Lemma~6; see \S\ref{lt:sec}, \S4), and Proposition~5 is the stronger of the two bounds for a coefficient.

\textbf{What remains to be written.} The passage from the estimate of this section to lattices is taken up in \S\ref{lt:sec}; some of the items below are treated there, as that section states.
\begin{enumerate}[label=\textup{(\roman*)},leftmargin=2.4em]
\item A proof of strong duality for $d\ge4$ (in the Lean: compactness and a separation theorem); see
Remark~(2) after Theorem~J.
\item Theorem~J with the operator on the ordinary rows, and the deduction of the statement for all frames
(\texttt{ThmAGauss}); the role of the congruence of Part~II, \S6.3.
\item Hypothesis (B) for the theta kernels of the construction: the height-two formal group at an inert
prime, the series $\beta$, and the integrality (B4). This is the arithmetic input of the whole section.
\item The passage from Theorem~J to \eqref{dv:link}: items (iii) and (iv) above, with the loss terms.
\item An example for $d\ge4$ with an exceptional monomial; none is known to us. (That the series $\beta$ of
the example after Theorem~J satisfies (B1)--(B4) in all degrees is no longer open: it is proved in
\S\ref{el:sec}, Proposition~Le22.) For $d\ge4$ Lemmas 3--4 and
Theorem~J were tested only on a series without exceptional monomials and on the combinatorial steps.
\item To follow, in the source, how the statement for one Taylor coefficient on the route of the final
theorem, \texttt{lemma\_R\_log\_loss}, is used in place of Proposition~5. (The statements
\texttt{TopRadius}, \texttt{top\_radius\_holds}, \texttt{LemmaR} and \texttt{ThmABound} of the archive are
not used by it.)
\end{enumerate}


\section{Divisibility: from derivatives to lattices}\label{lt:sec}

This section concerns the passage from the local estimate of \S III.5 to the prime sum $P_{u,t,j}$ of
\eqref{II:t4.1} of Part~II, the left side of the main inequality (M). In \S III.5 one function at one prime is
studied. What (M) needs is a statement about lattices over a number field, at all primes at once, with the
losses added up. The section states that passage, proves its elementary parts, and names the statements of
the source that carry the arithmetic. For the prime sum it does this much: it splits the sum and bounds
several of its terms. The estimates for the eleven loss sums and their total bound \eqref{lt:loss} are
quoted. Theorem~\ref{thm:T} alone does not give that bound. The section does not prove the divisibility
side of (M). The divisibility of
the period pairing is assumed, as statement \textup{(Val)} of \S2. Also assumed are a statement on
lattices, the inequality \eqref{dv:link} at the primes of one of the two classes, the estimates for the
eleven loss sums with their total bound, and the slope inequality for many maps (derived in
\S\ref{as:sec}, Proposition~M4, from the rank bound of \S III.8 and a cited fact on the filtration by
slopes).

\emph{Notation.} The notation of \S III.5, \S1 is kept: the scale $x,\varepsilon,b,a,b',X',X'',p_0,\delta_1$;
the parameters $m,k,t_0,n,f,N_m$; the set $\mathcal T$; the numbers $a_t$, $\mathrm{gain}$, $\tau$,
$\mathrm{wPart}$, $V(j,t,p)$, $\mathrm{lo}_G$, $\mathrm{fl}_G$, $\mathrm{Sc}$, $\eta$, and the constant $C$
of \eqref{dv:fl}. So $a$ is the number $xb$, as in \S III.4 and \S III.5, with $a\le a_t$. And $\tau(w)$ is the number in \eqref{dv:link}; it is not the modulus $\tau$ of \S III.4. The
letter $u$ is the translate of Part~II, \S4.2, as in $P_{u,t,j}$. The letters introduced here are listed at
the end of \S1.

\emph{The letters are local.} The letters of this section have the meanings given here in this section
only. Other sections of Part~III use some of them differently, and some objects have another name there.
The main cases are these. $\tau$ is a weight place in \S III.7. $\mathcal T$ is the level in \S III.4 and
\S III.7, and $\mathcal T(\mathsf w)$ in \S III.8 is the level prime of a word, which is $\mathrm{Lv}_\alpha$
here. The members $\alpha$ are the words $\mathsf w$ of \S III.8, where $\alpha$ is a constant. $U$ is a
power series in \S III.5 and a space on the target side in \S III.7. $\gamma$ is $\ell/(\ell^2-1)$ in
\S III.5 and in the paragraph on \texttt{asmA\_top} of \S4, and Euler's constant in \S6. $\mathrm{ka}_j$ is
the number $\varsigma$ of \S III.4, and $c_{\chi_d}$ is the $c_d$ of \S III.4 and \S III.7. $\lambda_w$ is
written $\mathrm{wt}(y)$ in \S III.7, and the $\lambda_{\mathsf w}$ of \S III.8 is another number.
$\mu_{0,\alpha}$ is not the $\mu_0$ of \S III.7 or of \S III.9. $\Xi$ is an operator in \S III.7 and
\S III.8. $L_0$ is the field called $K$ in \S III.9, where $L_0$ is another field.

\emph{Proved outright.} Three lemmas on lattices over the ring of integers of a number field: the slope of
a sublattice (Lemma~L1), the slope of a lattice multiplied by a fractional ideal, with the cost of rounding
(Lemma~L2), and the integrality of a pairing after such a multiplication (Lemma~L3). The bound for the
slope when the rank drops (Lemma~L5). The identity between the prime sum \eqref{II:t4.1} and sums over single primes
(Lemma~L6). Three elementary comparisons of losses (Lemmas~L7 to~L9). An inequality of the shape of
\eqref{dv:link} at the primes of class \texttt{signal}, for the explicit formulas of the source and with
explicit loss terms (Proposition~L10); it gives \eqref{dv:link} itself when the constants satisfy the
conditions \eqref{lt:consts}. The splitting of the prime sum by classes of primes (Proposition~L11) and three of
the estimates that bound its terms (Lemmas~L12 to~L14; the last uses Theorem~\ref{thm:T} and one cited
estimate).

\emph{Proved from two statements of the source.} Theorem~L4. Given the valuation statement \textup{(Val)}
of \S2, the glued lattice pairs integrally with the target. Given the lattice statement \textup{(Lat)} of
\S2, its slope exceeds that of the source lattice by the total gain, minus explicit losses.

\emph{A model, not a statement about the source.} Proposition~L15 shows how the terms cancel when there
is one injective map between spaces of equal dimension. That hypothesis does not hold in the source.

\emph{Taken from the source without proof.} The statement \textup{(Val)}, that is, the divisibility of the
period pairing of the chain of maps. In the Lean it rests on three assumptions (\texttt{adm2}, the
hypotheses \texttt{Tied} and \texttt{CutUp} of \texttt{val2AllAtRF}, and the field \texttt{l1b}; see \S2). It is where Theorem~J of \S III.5
is used, and it is conditional on hypothesis \textup{(B)} there. The statement \textup{(Lat)}. The
inequality \eqref{dv:link} at the primes of class \texttt{inertClass}. The estimates for the eleven sums over primes
of the losses, and their total bound \eqref{lt:loss}. The slope inequality for many maps (Part~II, \S6.9--6.10; derived in \S\ref{as:sec}, Proposition~M4, from
the rank bound of \S III.8 and a cited fact on the filtration by slopes). These are listed in \S7.

\subsection*{1.\quad Lattices over a number field: three lemmas}

\textbf{Places and weights.} Let $L_0$ be a number field of degree $n_0$, totally complex, with ring of
integers $\mathfrak o$. (In the source $L_0$ is the coefficient field, \texttt{LL}.) A \emph{place} $w$ is a
nonzero prime ideal $\mathfrak p_w$ of $\mathfrak o$; $p_w$ is its residue characteristic, $e_w$ its
ramification index over $p_w$, $v_w$ the valuation with $v_w(\mathfrak p_w)=1$, and $\mathfrak o_w$ the
localisation of $\mathfrak o$ at $\mathfrak p_w$, a discrete valuation ring. For an $\mathfrak o_w$-submodule
$\Lambda_w$ of an $L_0$-vector space and an integer $z$ (which may be negative) we write $\mathfrak p_w^{z}\Lambda_w$ for $(\mathfrak p_w\mathfrak o_w)^z\Lambda_w$.
As in Part~II, \S6.4, put
\begin{equation}\label{lt:lam}
 \lambda_w=\frac{\log\operatorname{Norm}(\mathfrak p_w)}{n_0},\qquad
 \sum_{w\mid p}\lambda_we_w=\log p,\qquad \mathrm{sz}_w(\sigma)=2\lambda_we_w\,\sigma\quad(\sigma\in\mathbb R).
\end{equation}
The middle identity holds because the sum over $w\mid p$ of $e_w$ times the residue degree of $w$ is $n_0$.
So a gain of $\sigma$ at every place above $p$, measured in units $\mathrm{sz}_w$, has total
$2\sigma\log p$. (Lean: \texttt{lw}, \texttt{SizeData$_2$.sz}, \texttt{FactEtaMean}. The Lean states
\texttt{FactEtaMean} for a field that is Galois over $\Q$. The identity holds for every number field, by
the argument just given.)

\textbf{Lattices and slopes.} Let $U$ be an $L_0$-vector space of finite dimension $\dim U\ge1$, and fix a Haar measure
on the real vector space $U\otimes_{\Q}\mathbb R$. A \emph{lattice} in $U$ is a finitely generated
$\mathfrak o$-submodule $\Lambda$ that spans $U$ over $L_0$. Its rank is $\mathrm{rk}\,\Lambda=\dim U$; it
is a discrete subgroup of $U\otimes\mathbb R$ of rank $n_0\dim U$, with a covolume. The slope divides by
the rank. So wherever a slope is taken, in Lemmas L1, L2 and L5 and in Theorem~L4, the space has positive
dimension and the lattice has rank $\ge1$. Put
\begin{equation}\label{lt:mu}
 \mu(\Lambda)=-\frac{2}{n_0\,\mathrm{rk}\,\Lambda}\,\log\operatorname{covol}(\Lambda)+\mathrm{cst}_{L_0},\qquad
 \mathrm{cst}_{L_0}=-\log2+\frac{\log|\mathrm{disc}\,L_0|}{n_0}.
\end{equation}
This is the slope in the normalisation for squared norms: twice the slope of Part~II, \S6.8, up to the
constant. (Lean: \texttt{muDef}, with the measure given by Hermitian forms at the complex places. In
Lemmas L1--L5 only differences $\mu(\Lambda)-\mu(\Lambda')$ for lattices in the same space occur; they do
not depend on the measure or on $\mathrm{cst}_{L_0}$.) For a finite $\mathfrak o$-module write $\lg_w$ of
it for the length of its localisation at $w$: the number of factors $\mathfrak o/\mathfrak p_w$ in a
composition series. For a lattice $\Lambda$ write $\Lambda_w=\mathfrak o_w\Lambda\subset U$. Here $\Lambda$
always denotes a lattice; the number $\Lambda_G$ of \eqref{dv:lo} is not used in this section.

\textbf{Lemma L1 (a sublattice).} \emph{Let $\Lambda'\subset\Lambda$ be lattices in $U$, with
$\dim U\ge1$. Then}
\begin{equation}\label{lt:index}
 [\Lambda:\Lambda']=\prod_w\operatorname{Norm}(\mathfrak p_w)^{\lg_w(\Lambda/\Lambda')},\qquad
 \mu(\Lambda')=\mu(\Lambda)-\frac{2}{\mathrm{rk}\,\Lambda}\sum_w\lambda_w\,\lg_w(\Lambda/\Lambda').
\end{equation}
\emph{Proof.} The quotient $\Lambda/\Lambda'$ is a finite $\mathfrak o$-module. It is the direct sum of its
localisations, and each composition factor $\mathfrak o/\mathfrak p_w$ has
$\operatorname{Norm}(\mathfrak p_w)$ elements. This gives the index. Then
$\operatorname{covol}(\Lambda')=[\Lambda:\Lambda']\operatorname{covol}(\Lambda)$ and
$\log\operatorname{Norm}(\mathfrak p_w)=n_0\lambda_w$. \qed

\noindent (Lean: \texttt{muDef\_index}, \texttt{mu\_index\_sum}.)

\textbf{Lemma L2 (multiplication by an ideal).} \emph{Let $\Lambda$ be a lattice in $U$, with
$\dim U\ge1$, let $z_w$ be
integers, zero for almost all $w$, and put $I=\prod_w\mathfrak p_w^{-z_w}$, a fractional ideal.}
\begin{enumerate}[label=\textup{(\roman*)},leftmargin=2.4em,itemsep=1pt,topsep=2pt]
\item \emph{$I\Lambda$ is a lattice with $(I\Lambda)_w=\mathfrak p_w^{-z_w}\Lambda_w$ for every $w$, and it
is the only lattice with these localisations.}
\item $\mu(I\Lambda)=\mu(\Lambda)+2\sum_w\lambda_wz_w$.
\item \emph{Let $\sigma_w$ be real numbers, zero for almost all $w$, let $z_w=\lfloor\sigma_w\rfloor$, and
let $\mathrm{Rd}$ be a finite set of primes containing $p_w$ for every $w$ with $\sigma_w\notin\Z$. Then}
\begin{equation}\label{lt:round}
 2\sum_w\lambda_w\sigma_w-2\sum_{p\in\mathrm{Rd}}\log p\ \le\ \mu(I\Lambda)-\mu(\Lambda)\ \le\ 2\sum_w\lambda_w\sigma_w .
\end{equation}
\end{enumerate}
\emph{Proof.} (i) Localisation commutes with products, and $I\mathfrak o_w=\mathfrak p_w^{-z_w}\mathfrak o_w$.
A lattice is determined by its localisations: if $y\in U$ lies in $\Lambda_w$ for all $w$, the ideal
$\{\beta\in\mathfrak o:\beta y\in\Lambda\}$ is contained in no $\mathfrak p_w$, so it is $\mathfrak o$ and
$y\in\Lambda$. (ii) Suppose first that all $z_w\le0$, so that $I\subset\mathfrak o$ and
$I\Lambda\subset\Lambda$. The module $\Lambda_w$ is finitely generated and torsion-free over the
discrete valuation ring $\mathfrak o_w$, hence free of rank $\mathrm{rk}\,\Lambda$. So the localisation
of $\Lambda/I\Lambda$ at $w$ is a sum of $\mathrm{rk}\,\Lambda$ copies of
$\mathfrak o_w/\mathfrak p_w^{-z_w}\mathfrak o_w$, of length $-z_w\,\mathrm{rk}\,\Lambda$, and Lemma~L1
gives (ii). In general write $I=I_1I_2^{-1}$ with ideals $I_1,I_2\subset\mathfrak o$. Then
$I_2\cdot(I\Lambda)=I_1\Lambda$. The case just proved, applied to $(I_2,I\Lambda)$ and to $(I_1,\Lambda)$,
gives $\mu(I\Lambda)-2\sum_w\lambda_wv_w(I_2)=\mu(\Lambda)-2\sum_w\lambda_wv_w(I_1)$, and
$v_w(I_1)-v_w(I_2)=-z_w$. (iii) By (ii) the middle term is $2\sum_w\lambda_wz_w$. Now
$0\le\sigma_w-z_w<1$, and $\sigma_w-z_w=0$ unless $p_w\in\mathrm{Rd}$. Finally
$\sum_{w\mid p}\lambda_w\le\sum_{w\mid p}\lambda_we_w=\log p$. \qed

\noindent (Lean: \texttt{ShiftEx}, \texttt{ShiftMu} with its proof \texttt{shift\_slope\_mu},
\texttt{RoundingCost}; the ideal is \texttt{shiftIdeal}, the lattice $I\Lambda$ is \texttt{shiftLat}. The
factor $2$ in (ii) is the squared normalisation.) The lemma and its proof hold for any number field. For
example over $\Q$, where $\lambda_p=\log p$: for $\Lambda=\Z^2$ and $I=\frac12\Z$ the covolume of
$I\Lambda$ is $\frac14$, and the slope \eqref{lt:mu} rises by $2\log2$.

\textbf{Lemma L3 (integrality after the multiplication).} \emph{Let $U,U'$ be $L_0$-vector spaces and
$\langle\ ,\ \rangle:U\times U'\to L_0$ a bilinear map. Let $\Lambda\subset U$ and
$\mathcal R\subset U'$ be lattices. For each place $w$ let $\mathrm{Gen}_w\subset U'$ be a set that
generates $\mathcal R_w$ over $\mathfrak o_w$, and let $z_w$ be integers, zero for almost all $w$, such
that}
\begin{equation}\label{lt:val}
 \langle y,y'\rangle=0\quad\text{or}\quad v_w\big(\langle y,y'\rangle\big)\ge z_w
 \qquad\text{for all }y\in\Lambda,\ y'\in\mathrm{Gen}_w .
\end{equation}
\emph{Let $I=\prod_w\mathfrak p_w^{-z_w}$. Then $\langle\tilde y,\tilde y'\rangle\in\mathfrak o$ for all
$\tilde y\in I\Lambda$ and $\tilde y'\in\mathcal R$. If moreover $v_{w_0}(\langle y,y'\rangle)=z_{w_0}$ for
some place $w_0$, some $y\in\Lambda$ and some $y'\in\mathrm{Gen}_{w_0}$, then
$\mathfrak p_{w_0}^{-1}I\Lambda$ does not have this property.}

\emph{Proof.} An element of $L_0$ lies in $\mathfrak o$ if its valuation is $\ge0$ at every place. Fix
$w$. The element $\tilde y$ is a finite sum of products $\beta y$ with $\beta\in I$ and $y\in\Lambda$. The
element $\tilde y'$ is a finite sum of products $\beta'y'$ with $\beta'\in\mathfrak o_w$ and
$y'\in\mathrm{Gen}_w$. So $\langle\tilde y,\tilde y'\rangle$ is a sum of terms
$\beta\beta'\langle y,y'\rangle$, and each term is zero or has valuation at least $-z_w+0+z_w=0$ at $w$.
For the second assertion take $\beta\in\mathfrak p_{w_0}^{-1}I$ with $v_{w_0}(\beta)=-z_{w_0}-1$, and
$\beta'\in\mathfrak o\smallsetminus\mathfrak p_{w_0}$ with $\beta'y'\in\mathcal R$. Then
$v_{w_0}(\langle\beta y,\beta'y'\rangle)=-1$. \qed

\noindent No slope occurs in Lemma~L3, and it holds also for spaces of dimension $0$.
(Lean: \texttt{ShiftPairing}, \texttt{DivIntoDualShift}. There the hypothesis is stated with real
numbers and integers at most their ceilings, which is the same since valuations are integers; and the
local forms differ from $\langle\ ,\ \rangle$ by units at $w$, which changes nothing.)

\textbf{Remarks.} (1) \emph{Places are independent.} By Lemma~L2\,(i) the localisation of $I\Lambda$ at $w$
depends only on $z_w$, and by (ii) the gains add over the places. This holds also for two places above the
same prime. It is the form over a number field of the gluing over $\Q$ of Part~II, \S6.8, where the Chinese
remainder theorem was used. (2) \emph{Maps compose.} Suppose that there are $m$ maps and $m+1$ lattices,
and that for $1\le j\le m$ the $j$th map sends the localisation at $w$ of the lattice before it into
$\mathfrak p_w^{z_{w,j}}$ times the localisation of the lattice after it. Then the composite sends the
first localisation into $\mathfrak p_w^{z_{w,1}+\cdots+z_{w,m}}$ times the last. So the exponents of the
$m$ maps of the chain add. This is the integrality assertion of Part~II, \S6.2. (3) \emph{A sublattice
costs.} If $\Lambda'_w\subset\mathfrak p_w^{z}\Lambda_w$ with an integer $z\ge0$, then
$\lg_w(\Lambda/\Lambda')\ge z\,\mathrm{rk}\,\Lambda$, and by Lemma~L1 the slope drops by at least
$2\lambda_wz$ at $w$.

\emph{Checks.} The three lemmas were tested in exact arithmetic over $\Q(\sqrt{-5})$, whose class number
is $2$, with the ramified prime above $2$, the two primes above $3$ and above $7$, and the inert prime
$11$; lattices of rank $1$ to $3$, not all free. In $60$ cases the two sides of Lemma~L2\,(ii) agree to
$2\cdot10^{-15}$; with $\sum\lambda_wz_w$ in place of $2\sum\lambda_wz_w$ they do not. In $25$ cases the
index formula of Lemma~L1 holds exactly, the lengths being computed place by place. In $12$ cases the
pairing of Lemma~L3 is integral on $I\Lambda$ and is not on $\mathfrak p_{w_0}^{-1}I\Lambda$.

\emph{Letters introduced in this section.} The field $L_0$, $\mathfrak o$, $n_0$; a place $w$ with
$\mathfrak p_w,p_w,e_w,v_w,\mathfrak o_w$; $\lambda_w$, $\mathrm{sz}_w$; spaces $U$, $U'$; a lattice
$\Lambda$, $\mathrm{rk}$, $\mu$, $\mathrm{cst}_{L_0}$, $\lg_w$; integers $z_w$, real numbers $\sigma_w$, an
ideal $I$, a set of primes $\mathrm{Rd}$; elements $y,\tilde y$ of $U$ and $y',\tilde y'$ of $U'$, and
scalars $\beta,\beta'$ in $L_0$; a member $\alpha$, with $U_\alpha$, $\Lambda^{\rm src}_\alpha$,
$\Lambda^{\rm div}_\alpha$, $\Lambda^{\rm gl}_\alpha$, $I_\alpha$, $\mathrm{Lv}_\alpha$,
$\mathrm{rk}^*_\alpha$, $\mathrm{Rd}_\alpha$; the lattice $\mathcal R$, the sets $\mathrm{Gen}_w$;
$\mathrm{sg}_\alpha(w)$, $\phi(j,w)$, $\mathrm{err}(\alpha,j,w)$, $\mathrm{io}(w)$; $G_{\rm top}$, $F_0$,
$E_{\rm err}(\alpha)$, $\mathrm{Lo}(\alpha)$; the group $H$, $h_p$, $n_p$, an exponent $i$ of a prime,
$\mathrm{cell}$, $\mathrm{wt}_j$, $V_\infty$, an orbit $\mathcal Q$, $\mathcal T(\mathcal Q)$,
$n(\mathcal Q)$, $\kappa(\mathcal Q)$, $\mathrm{Srv}(p)$, $\mathrm{ka}_j$, $\mathrm{gain}^{\rm plain}$,
$\mathrm{gain}^{\rm srv}$, $\mathrm{Pr}$, $\mathrm{Ex}$, $\mathrm{dl}$, $C_1$, $C_2$, $\mathrm{dig}_p$,
the constants $C_{\rm pl}$, $C_{\rm lv}$, $C_{\rm den}$ of \S2 and $C_{\rm r}$ of \S5,
$\mathrm{odd}_p$, real numbers $A$, $c$, $c_i$, $\psi_i$, counts $n_{\rm var}$, $n_{\rm st}$, $n_{\rm top}$; sets of primes $\Xi$, $\Xi_{\rm s}$, $\Xi_1$, $\Xi_{\rm c}$; $\mathrm{def}(p)$,
$\mathcal T_p$, $\mathrm{Sig}(p)$, $\mathrm{bw}$, $\theta$, $a'$, $Y$; and in \S6 $\mathrm{Per}_\alpha$,
$\mathrm{ht}_\alpha$ and the error terms $E_{\rm size},E_{\rm lev},E_{\rm cnt},E_{\rm src}$. The number of
digits of an integer in base $p$, written $\lambda(k)$ in \S III.5, is written $\mathrm{dig}_p(k)$ here,
because $\lambda_w$ is the weight of a place. The level of a member, written $\mathcal T$ in Part~II,
\S6.7, is written $\mathrm{Lv}_\alpha$ here, because $\mathcal T$ is the set of labels, as in \S III.5.

\subsection*{2.\quad The glued lattice: statement and proof}

\textbf{The data.} Fix $x$ and one choice of the data of the construction (a ``record'' of the source).
The record contains one translate $u\in H$ (Lean: the field \texttt{Host.u0}). It is the translate $u$ of
Part~II, \S4.2, and the weights of the record are the weights for this $u$ (\S4 below). The maps of the
record are indexed by a finite set. Its elements $\alpha$ are the \emph{members} (Lean: \texttt{Word}). A
member is not a translate. All members of a record have the same translate, and the prime sum $P_{u,t,j}$
does not depend on the member (the Lean statement \texttt{Loss\_at} contains no member). A member enters
through its lattices, its level and its error terms. For each member $\alpha$ the source gives:
\begin{itemize}[leftmargin=1.6em,itemsep=1pt,topsep=2pt]
\item an $L_0$-vector space $U_\alpha$ of dimension $\ge1$ with a lattice $\Lambda^{\rm src}_\alpha$, the
\emph{source lattice} (Lean: \texttt{MramL}; it is the preimage of the first lattice of the chain of maps;
the positive dimension is the field \texttt{a\_pos} of the inputs of \texttt{divClausesR\_of}, \S3);
\item a lattice $\Lambda^{\rm div}_\alpha\subset\Lambda^{\rm src}_\alpha$ (Lean: \texttt{NuLat});
\item a space $U'$ and a lattice $\mathcal R\subset U'$, the same for all $\alpha$ (Lean: \texttt{RR}), and
a bilinear pairing $\langle\ ,\ \rangle_\alpha:U_\alpha\times U'\to L_0$, the period pairing of the whole
chain (Lean: \texttt{perCol}; the bilinear form is \texttt{col.B});
\item a prime $\mathrm{Lv}_\alpha$, the level of the member (the level $\mathcal T$ of Part~II, \S6.7), and
an integer $\mathrm{rk}^*_\alpha\ge\mathrm{rk}\,\Lambda^{\rm src}_\alpha$ (Lean: \texttt{duS}).
\end{itemize}
Independently of $\alpha$ except where shown, there are real numbers, zero outside a finite set of places:
$\mathrm{gain}(j,w)\ge0$, $\phi(j,w)$, $\mathrm{err}(\alpha,j,w)$, $\tau(w)\ge0$, $\mathrm{io}(w)\ge0$, for
$1\le j\le m$ (Lean: the fields \texttt{gain}, \texttt{F}, \texttt{E}, \texttt{tau}, \texttt{io} of
\texttt{SizeData}$_2$). Here $\mathrm{gain}$ and $\tau$ are the numbers of \eqref{dv:link}. Put
\begin{equation}\label{lt:sg}
 \mathrm{sg}_\alpha(w)=\sum_{j=1}^m\big(\mathrm{gain}(j,w)+\phi(j,w)-\mathrm{err}(\alpha,j,w)\big),
\end{equation}
\begin{equation}\label{lt:totals}
 G_{\rm top}=\sum_w\mathrm{sz}_w\Big(\sum_j\mathrm{gain}(j,w)-\tau(w)\Big),\quad
 F_0=\sum_{j,w}\mathrm{sz}_w\big(\phi(j,w)\big),\quad
 E_{\rm err}(\alpha)=\sum_{j,w}\mathrm{sz}_w\big(\mathrm{err}(\alpha,j,w)\big).
\end{equation}
(Lean: \texttt{sg}, \texttt{Ga}, \texttt{Fo}, \texttt{Ee}. In the Lean the gain is first set to zero at
the places above a set of ``dropped'' primes, see \S5; we assume this done.) So
\begin{equation}\label{lt:sgsum}
 \sum_w\mathrm{sz}_w\big(\mathrm{sg}_\alpha(w)-\tau(w)\big)=G_{\rm top}+F_0-E_{\rm err}(\alpha).
\end{equation}
The number $F_0$ is the $F_0$ of Part~II, \S6.7. The \emph{glued lattice} is
\begin{equation}\label{lt:N}
 \Lambda^{\rm gl}_\alpha=I_\alpha\,\Lambda^{\rm div}_\alpha,\qquad
 I_\alpha=\prod_w\mathfrak p_w^{-\lfloor e_w\,\mathrm{sg}_\alpha(w)\rfloor}.
\end{equation}
(Lean: \texttt{Nu}, \texttt{Iu}, and the statement \texttt{NprimeDef}, which is \eqref{lt:N}.) The factor
$e_w$ converts an exponent of $p_w$ into an exponent of $\mathfrak p_w$.

\textbf{The two statements taken from the source.}
\begin{enumerate}[leftmargin=3.4em,itemsep=2pt,topsep=2pt]
\item[\textup{(Val)}] \emph{Valuation.} For every member $\alpha$ and every place $w$ there is a set
$\mathrm{Gen}_w\subset U'$ generating $\mathcal R_w$ over $\mathfrak o_w$ with
\begin{equation}\label{lt:V}
 \langle y,y'\rangle_\alpha=0\quad\text{or}\quad
 v_w\big(\langle y,y'\rangle_\alpha\big)\ \ge\ \big\lfloor e_w\,\mathrm{sg}_\alpha(w)\big\rfloor
 \qquad(y\in\Lambda^{\rm div}_\alpha,\ y'\in\mathrm{Gen}_w).
\end{equation}
\item[\textup{(Lat)}] \emph{Lattices.} For every member $\alpha$ and every place $w$ with
$p_w\ne\mathrm{Lv}_\alpha$: $\tau(w)$ is an integer,
$(\Lambda^{\rm div}_\alpha)_w\subset\mathfrak p_w^{e_w\tau(w)}(\Lambda^{\rm src}_\alpha)_w$, and
\begin{equation}\label{lt:L}
 \lg_w\big(\Lambda^{\rm src}_\alpha/\Lambda^{\rm div}_\alpha\big)\ \le\
 \mathrm{rk}^*_\alpha\,e_w\,\big(\tau(w)+\mathrm{io}(w)\big).
\end{equation}
\end{enumerate}
Note what \textup{(Lat)} does not give. It has $\mathrm{rk}^*_\alpha$, not
$\mathrm{rk}\,\Lambda^{\rm src}_\alpha$, on the right of \eqref{lt:L}. And it says nothing at the places
above $\mathrm{Lv}_\alpha$.

\emph{The statement \textup{(Val)} in the Lean.} It is not one Lean statement. It is assumed here; in
the Lean it is obtained from three assumptions.
\begin{enumerate}[label=\textup{(\arabic*)},leftmargin=2.4em,itemsep=1pt,topsep=2pt]
\item \texttt{LocalQuadraticParts.adm2}. Let $\mu_w$ be the twist chosen at $w$ (\texttt{muAt}). The
letters of this item are those of the Lean: $\mu_w$, with a place as index, is a twist and not a slope; the
element $F$ below has nothing to do with the number $F_0$; and $\widetilde M$, $V_0$ have nothing to do
with $M_{u,t,j}$ and $V$. The statement is: for every member $\alpha$, every place $w$, every $y\in\Lambda^{\rm div}_\alpha$, and every
$F$ that lies both in the lattice $\widetilde M(\mu_w)$ (\texttt{Mtilde}) and in the subspace
$V_0(\mu_w)$ (\texttt{V0}) of the ambient space of the twist, there is an element of $L_0$ whose image
under the fixed embedding of $L_0$ in $\mathbb C$ (\texttt{embLL}) is the period of $y$ against $F$; and
this element is zero or has valuation at least $\lfloor e_w\mathrm{sg}_\alpha(w)\rfloor$ at $w$. Both
restrictions on $F$ are part of the statement. It says nothing for an element of the lattice outside
$V_0(\mu_w)$, and nothing for an element of $V_0(\mu_w)$ outside the lattice. That the period is the
embedded value of an element of $L_0$ is also part of the statement; \textup{(Val)} presupposes it.
\item The hypotheses of the theorem \texttt{val2AllAtRF}. This theorem turns \textup{(1)} into the same
bound for the matrix form of the pairing, on a lattice \texttt{Nu} given in coordinates. Besides
\textup{(1)} it has three hypotheses. Two are carried by structures. \texttt{ColBrF.Tied} says that the
matrix form is tied to the period: the twist at $w$ is the one of the matrix form, each element of the
lattice of the matrix form maps into $V_0$ and is an integral combination of the generators of
$\widetilde M$, and the matrix pairing has the period as its embedded value. \texttt{CutUp}, for every twist
that is used, says that an element of $V_0$ that is an integral combination of the generators of
$\widetilde M$ is an integral combination of elements of $\widetilde M\cap V_0$. This is what the two
restrictions in \textup{(1)} require. The third hypothesis is the compatibility with the source of
\textup{(1)}: every element of \texttt{Nu} is mapped by the coordinate map \texttt{eNu} into
\texttt{NuLat}, that is, into $\Lambda^{\rm div}_\alpha$. It is not an assumption. Where
\texttt{val2AllAtRF} is applied (in the files \texttt{RecBridge.lean} and \texttt{RecBridgeTied.lean}) it
is supplied by the theorem \texttt{nu\_into}, from the definition \eqref{lt:N} (\texttt{NprimeDef}) and two
identities of lattices: \texttt{shift}, a field of the package, and \texttt{lat}, a field of the structure
\texttt{HLatTie}, which the package carries for each member as its field \texttt{hNu}.
\item The structure field \texttt{l1b} (of type \texttt{PairingFactorsThroughPeriods}). It says that the
local form at $w$ equals a unit at $w$ times $\langle y,y'\rangle_\alpha$, where $y'$ is the image of the
element of the local lattice under a fixed linear map.
\end{enumerate}
That the sets $\mathrm{Gen}_w$ generate $\mathcal R_w$ is not an assumption. The set $\mathrm{Gen}_w$ is the
image of the local lattice of \textup{(1)}. The lattice $\mathcal R$ is \emph{defined} as the lattice
glued from these images (\texttt{RR}, by \texttt{glueLat}), and the theorem \texttt{latticeGlue} proves
that this is a lattice with these localisations (Lemma~L2\,(i) is the case of an ideal). The passage to one
bilinear pairing is faithful: the Lean theorem \texttt{intoDual} applies \texttt{DivIntoDualShift}
(Lemma~L3) with one $L_0$-bilinear form. The quantifiers of \textup{(Val)} are those of the Lean: all
members, all places, all $y\in\Lambda^{\rm div}_\alpha$. The final theorem depends on \texttt{adm2} (the
dependency computation confirms this). It also depends on \texttt{val2AllAtRF} and on
\texttt{latticeGlue}: in the seventh run of the dependency computation both are marked as used. Whether the
hypotheses \texttt{Tied} and \texttt{CutUp} of \textup{(2)} and the field of \textup{(3)} are discharged
for the construction on the route of the final theorem was not determined. The dependency computation does
not say how hypotheses are discharged, and the proof terms were not followed that far.

\emph{The statement \textup{(Lat)} in the Lean.} It is the part of the first clause of \texttt{LinkG\_at}
that is used here, read with \texttt{NuDef}. That clause is for the places $w$ with
$p_w\ne\mathrm{Lv}_\alpha$ and has three parts. The first is the inclusion, with the exponent
$\lfloor e_w\tau(w)\rfloor$. The second is
$\mathtt{lenMN}(\alpha,w)\le\mathrm{rk}^*_\alpha\,(\tau(w)+\mathrm{io}(w))$. The third says: if the set
\texttt{VV} of the Lean at the prime $p_w$ is empty, then $\tau(w)=0$, $\mathrm{io}(w)=0$ and
$\mathtt{lenMN}(\alpha,w)=0$. The third part is not used in Theorem~L4. The number \texttt{lenMN} is not
the local length. The last clause of \texttt{NuDef} says that
$\lg_w(\Lambda^{\rm src}_\alpha/\Lambda^{\rm div}_\alpha)=\mathtt{lenMN}(\alpha,w)\cdot e_w$. With the
second part this gives \eqref{lt:L}, and it is where the factor $e_w$ there comes from. For the
integrality of $\tau(w)$ two statements are needed. The theorem \texttt{exists\_tau\_eq\_natCast} says that
the number given by the formula \texttt{tauP} is a natural number. It does not mention the $\tau$ of the
record. That $\tau(w)$ equals \texttt{tauP} at $p_w$ is part of the hypothesis \texttt{fills} of \S3.

\textbf{Theorem L4 (the glued lattice).} \emph{Let $\alpha$ be a member. Write
$\mathrm{rk}=\mathrm{rk}\,\Lambda^{\rm src}_\alpha$ and $\mathrm{rk}^*=\mathrm{rk}^*_\alpha$; so
$\mathrm{rk}^*\ge\mathrm{rk}\ge1$.}
\begin{enumerate}[label=\textup{(\roman*)},leftmargin=2.4em,itemsep=1pt,topsep=2pt]
\item \emph{Assume \textup{(Val)}. Then $\langle y,y'\rangle_\alpha\in\mathfrak o$ for all
$y\in\Lambda^{\rm gl}_\alpha$ and $y'\in\mathcal R$. That is, the pairing maps $\Lambda^{\rm gl}_\alpha$
into the dual lattice $\mathcal R^\vee=\{\xi\in\mathrm{Hom}_{L_0}(U',L_0):
\xi(\mathcal R)\subset\mathfrak o\}$.}
\item \emph{Without any assumption,}
\begin{equation}\label{lt:mueq}
 \mu(\Lambda^{\rm gl}_\alpha)=\mu(\Lambda^{\rm src}_\alpha)
   +\sum_w2\lambda_w\big\lfloor e_w\,\mathrm{sg}_\alpha(w)\big\rfloor
   -\frac2{\mathrm{rk}}\sum_w\lambda_w\,\lg_w\big(\Lambda^{\rm src}_\alpha/\Lambda^{\rm div}_\alpha\big).
\end{equation}
\item \emph{Assume \textup{(Lat)}. Let $\mathrm{Rd}_\alpha$ be the set of primes $p_w$ with
$e_w\,\mathrm{sg}_\alpha(w)\notin\Z$. Then}
\begin{equation}\label{lt:slope}
 \mu(\Lambda^{\rm gl}_\alpha)\ \ge\ \mu(\Lambda^{\rm src}_\alpha)+G_{\rm top}+F_0-E_{\rm err}(\alpha)-\mathrm{Lo}(\alpha),
\end{equation}
\emph{where}
\begin{equation}\label{lt:Lodef}
\begin{aligned}
 \mathrm{Lo}(\alpha)={}&\sum_{w:\ p_w\ne\mathrm{Lv}_\alpha}
   \mathrm{sz}_w\Big(\frac{\mathrm{rk}^*\,\mathrm{io}(w)+(\mathrm{rk}^*-\mathrm{rk})\,\tau(w)}{\mathrm{rk}}\Big)\\
 &+\sum_{w:\ p_w=\mathrm{Lv}_\alpha}\Big(\frac2{\mathrm{rk}}\,\lambda_w\,
   \lg_w\big(\Lambda^{\rm src}_\alpha/\Lambda^{\rm div}_\alpha\big)-\mathrm{sz}_w\big(\tau(w)\big)\Big)
   +2\sum_{p\in\mathrm{Rd}_\alpha}\log p .
\end{aligned}
\end{equation}
\end{enumerate}
The three terms of $\mathrm{Lo}(\alpha)$ are: the loss of the lattice $\Lambda^{\rm div}_\alpha$ beyond
$\tau$, with the cost of the rank $\mathrm{rk}^*\ge\mathrm{rk}$; the places above the level, where
\textup{(Lat)} gives no bound and the length is kept as it is; and the rounding. If
$\mathrm{rk}^*=\mathrm{rk}$ the first term is the sum of $\mathrm{sz}_w(\mathrm{io}(w))$ over the places $w$ with $p_w\ne\mathrm{Lv}_\alpha$.

\emph{Proof.} (i) This is Lemma~L3 with $\Lambda=\Lambda^{\rm div}_\alpha$ and
$z_w=\lfloor e_w\mathrm{sg}_\alpha(w)\rfloor$; the hypothesis \eqref{lt:val} is \eqref{lt:V}.

(ii) Lemma~L2\,(ii) for the pair $(I_\alpha,\Lambda^{\rm div}_\alpha)$ gives
$\mu(\Lambda^{\rm gl}_\alpha)=\mu(\Lambda^{\rm div}_\alpha)+\sum_w2\lambda_w\lfloor e_w\mathrm{sg}_\alpha(w)\rfloor$.
Lemma~L1 for $\Lambda^{\rm div}_\alpha\subset\Lambda^{\rm src}_\alpha$ gives
$\mu(\Lambda^{\rm div}_\alpha)=\mu(\Lambda^{\rm src}_\alpha)-\frac2{\mathrm{rk}}\sum_w\lambda_w\lg_w(\Lambda^{\rm src}_\alpha/\Lambda^{\rm div}_\alpha)$.

(iii) We bound the two sums in \eqref{lt:mueq}. \emph{The first sum.} Apply Lemma~L2\,(iii) with
$\sigma_w=e_w\,\mathrm{sg}_\alpha(w)$ and $\mathrm{Rd}=\mathrm{Rd}_\alpha$. These numbers vanish outside a
finite set of places, and $2\lambda_w\sigma_w=\mathrm{sz}_w(\mathrm{sg}_\alpha(w))$. So
\begin{equation}\label{lt:pf1}
 \sum_w2\lambda_w\big\lfloor e_w\,\mathrm{sg}_\alpha(w)\big\rfloor\ \ge\
 \sum_w\mathrm{sz}_w\big(\mathrm{sg}_\alpha(w)\big)-2\sum_{p\in\mathrm{Rd}_\alpha}\log p .
\end{equation}
\emph{The second sum, at a place $w$ with $p_w\ne\mathrm{Lv}_\alpha$.} By \eqref{lt:L} and Lemma~L5 below
(whose proof is one line and uses nothing of this theorem),
\begin{equation}\label{lt:pf2}
 \frac2{\mathrm{rk}}\,\lambda_w\,\lg_w\big(\Lambda^{\rm src}_\alpha/\Lambda^{\rm div}_\alpha\big)\ \le\
 \mathrm{sz}_w\big(\tau(w)\big)
 +\mathrm{sz}_w\Big(\frac{\mathrm{rk}^*\,\mathrm{io}(w)+(\mathrm{rk}^*-\mathrm{rk})\,\tau(w)}{\mathrm{rk}}\Big).
\end{equation}
\emph{The second sum, at a place $w$ with $p_w=\mathrm{Lv}_\alpha$.} Here nothing is assumed. We write the
term as $\mathrm{sz}_w(\tau(w))$ plus the difference
$\frac2{\mathrm{rk}}\lambda_w\lg_w(\Lambda^{\rm src}_\alpha/\Lambda^{\rm div}_\alpha)-\mathrm{sz}_w(\tau(w))$.
\emph{Conclusion.} Insert \eqref{lt:pf1}, \eqref{lt:pf2} and the last identity into \eqref{lt:mueq}. All
sums are finite. The result is
\begin{equation}\label{D:226}
 \mu(\Lambda^{\rm gl}_\alpha)\ \ge\ \mu(\Lambda^{\rm src}_\alpha)
 +\sum_w\mathrm{sz}_w\big(\mathrm{sg}_\alpha(w)-\tau(w)\big)-\mathrm{Lo}(\alpha),
\end{equation}
with $\mathrm{Lo}(\alpha)$ as in \eqref{lt:Lodef}. By \eqref{lt:sgsum} this is \eqref{lt:slope}. \qed

\noindent Of \textup{(Lat)} the proof of (iii) uses only \eqref{lt:L}. The inclusion
$(\Lambda^{\rm div}_\alpha)_w\subset\mathfrak p_w^{e_w\tau(w)}(\Lambda^{\rm src}_\alpha)_w$ explains the
signs. The gain $\mathrm{gain}(j,w)$ enters $\Lambda^{\rm gl}_\alpha$ with a plus sign and $\tau(w)$ with a
minus sign for the following reason. The divisibility \eqref{lt:V} holds on $\Lambda^{\rm div}_\alpha$, not
on $\Lambda^{\rm src}_\alpha$. By the inclusion and Remark~(3), $\Lambda^{\rm div}_\alpha$ is smaller than
$\Lambda^{\rm src}_\alpha$ by at least $\tau(w)$ at $w$, in the units $\mathrm{sz}_w$. The net exponent at
$w$ is $\sum_j\mathrm{gain}(j,w)-\tau(w)$: the right side of \eqref{dv:link}.

\textbf{Lemma L5 (when the rank drops).} \emph{Let $\Lambda'\subset\Lambda$ be lattices, $w$ a place,
$\mathrm{rk}=\mathrm{rk}\,\Lambda\ge1$, and let $\mathrm{rk}^*\ge\mathrm{rk}$, $\tau(w)\ge0$ and
$\mathrm{io}(w)\ge0$ be real numbers with
$\lg_w(\Lambda/\Lambda')\le\mathrm{rk}^*e_w(\tau(w)+\mathrm{io}(w))$. Then}
\begin{equation}\label{lt:low}
 \frac2{\mathrm{rk}}\,\lambda_w\,\lg_w(\Lambda/\Lambda')-\mathrm{sz}_w\big(\tau(w)\big)\ \le\
 \mathrm{sz}_w\Big(\frac{\mathrm{rk}^*\,\mathrm{io}(w)+(\mathrm{rk}^*-\mathrm{rk})\,\tau(w)}{\mathrm{rk}}\Big).
\end{equation}
\emph{Proof.} The left side is at most
$2\lambda_we_w\big(\frac{\mathrm{rk}^*}{\mathrm{rk}}(\tau(w)+\mathrm{io}(w))-\tau(w)\big)$, and
$\frac{\mathrm{rk}^*}{\mathrm{rk}}(\tau(w)+\mathrm{io}(w))-\tau(w)=
\frac{\mathrm{rk}^*\mathrm{io}(w)+(\mathrm{rk}^*-\mathrm{rk})\tau(w)}{\mathrm{rk}}$. \qed

\noindent (Lean: \texttt{ClaimLow}. There the rank deficit is a number \texttt{de} with
$\mathrm{rk}=(1-\mathtt{de})\,\mathrm{rk}^*$, and the right side of \eqref{lt:low} reads
$\mathrm{sz}_w\big((\mathrm{io}(w)+\mathtt{de}\,\tau(w))/(1-\mathtt{de})\big)$, which is the same number.)

\emph{Checks.} Theorem~L4 was tested in exact arithmetic over $\Q(\sqrt{-5})$ on $200$ lattices of rank
$2$, with $\mathrm{rk}^*$ from $2$ to $5$, with the level prime $2$, $3$, $7$ or $11$, with a sublattice
that is smaller at a place above the level, and with every hypothesis of \textup{(Lat)} verified on the
data. The identity \eqref{lt:mueq} holds exactly in all cases, and \eqref{lt:slope} holds in all cases. In
$60$ cases the data were chosen with equality in \eqref{lt:L} and with $e_w\mathrm{sg}_\alpha(w)$ integral;
then \eqref{lt:slope} is an equality, exactly, in all $60$. On these $60$ cases three weaker forms of
$\mathrm{Lo}(\alpha)$ were tried. With $\sum_w\mathrm{sz}_w(\mathrm{io}(w))$ in place of the first two
terms (the loss one would get from \eqref{lt:L} with $\mathrm{rk}$ in place of $\mathrm{rk}^*$ at all
places) the inequality fails in $48$. Without the term at the level it fails in $15$, which are all the
cases where that term is positive. Without the factor for the rank it fails in $46$, which are all the
cases where $\mathrm{rk}^*>\mathrm{rk}$ and $\tau+\mathrm{io}$ is not zero. On the other $140$ cases,
without the rounding term it fails in $15$. Lemma~L5 was also tested on $20000$ random data in the form of
the Lean statement \texttt{ClaimLow}; without the factor for the rank the inequality fails in a third of
the cases.

\textbf{The Lean statement, and what it compares with.} The Lean statement of the slope bound
(\texttt{SlopesN\_at}, proved from \texttt{Places\_at} by \texttt{lower\_bound\_simplified}) is not
\eqref{lt:slope} word for word. It reads
\begin{equation}\label{lt:Lo}
 \mu(\Lambda^{\rm gl}_\alpha)\ \ge\ \mu_{0,\alpha}+G_{\rm top}+F_0-E_{\rm err}(\alpha)
 -\Big(C_{\rm pl}\sum_w\lambda_we_w\mathrm{io}(w)+C_{\rm pl}\,|\mathcal T|\,m\Big(\frac ak+b'\Big)\log\mathrm{Lv}_\alpha
   +C_{\rm pl}\,X''+1\Big),
\end{equation}
where $C_{\rm pl}$ is a constant of the source (it is not the constant $C$ of \eqref{dv:fl}), and
$\mu_{0,\alpha}$ is the slope of the \emph{first lattice of the chain}, not of
$\Lambda^{\rm src}_\alpha$, which is its preimage. It comes with an upper bound of the same shape for the
largest slope of a sublattice of $\Lambda^{\rm gl}_\alpha$. The passage from \eqref{lt:slope} to
\eqref{lt:Lo} uses four statements of the source, which are assumed here (Lean: \texttt{adm26},
\texttt{HypLow2}, \texttt{places\_at\_of}).
\begin{enumerate}[label=\textup{(\alph*)},leftmargin=2.4em,itemsep=1pt,topsep=2pt]
\item \emph{The rank.} $\mathrm{rk}=(1-1/(\mathrm{Lv}_\alpha-1))^f\,\mathrm{rk}^*$ (\texttt{adm26}). Hence
$(\mathrm{rk}^*-\mathrm{rk})/\mathrm{rk}^*\le2f/\mathrm{Lv}_\alpha$, which the source bounds by $\frac12$.
So $\mathrm{rk}^*/\mathrm{rk}\le2$, and the part of the first term of \eqref{lt:Lodef} with
$\mathrm{io}$ is at most $4\sum_w\lambda_we_w\mathrm{io}(w)$. This factor is absorbed in the constant $C_{\rm pl}$
of \eqref{lt:Lo}; it is not part of the ``$+1$''.
\item \emph{The places above the level.} There $\tau(w)=0$; this is the hypothesis \texttt{htauT} of
\texttt{places\_at\_of}, not a clause of \texttt{adm26}. The second term of \eqref{lt:Lodef} is then at
most $C_{\rm pl}|\mathcal T|m(a/k+b')\log\mathrm{Lv}_\alpha$ (the last clause of \texttt{adm26}). This is
the second term of \eqref{lt:Lo}.
\item \emph{The rounding.} $2\sum_{p\in\mathrm{Rd}_\alpha}\log p\le C_{\rm pl}X''$. The theorem
\texttt{places\_at\_of} does not prove this: it has the bound as a hypothesis (\texttt{NprimeMuR}, the
statement that $\mu(\Lambda^{\rm gl}_\alpha)-\mu(\Lambda^{\rm div}_\alpha)$ differs from
$\sum_w\mathrm{sz}_w(\mathrm{sg}_\alpha(w))$ by at most $C_{\rm pl}X''$). The source proves that hypothesis
elsewhere (\texttt{round\_le}, \texttt{nprimeMu\_at}, \texttt{nprimeMuR\_at}, \texttt{linkMu}). There the
rounding is counted at every place with $\mathrm{sg}_\alpha(w)\ne0$; the primes below these places include
$\mathrm{Rd}_\alpha$. The set $\mathrm{Rd}_\alpha$ is not contained in $[1,X'']$. By two statements of the
source (\texttt{ZeroAbove} for $\mathrm{gain}$ and $\mathrm{err}$, the field \texttt{F\_above} for $\phi$;
together \texttt{sg\_eq\_zero}), $\mathrm{sg}_\alpha(w)=0$ if
$p_w>X''$, unless $p_w$ is $\mathrm{Lv}_\alpha$, $g_0$, $g_1$, $d$ or a prime factor of the denominator
$D_\alpha$ of \S3. So
\begin{equation}\label{D:229}
 \sum_{p\in\mathrm{Rd}_\alpha}\log p\ \le\ X''\log4+\log\mathrm{Lv}_\alpha+\log g_0+\log g_1+\log d+\log D_\alpha ,
\end{equation}
by Chebyshev's bound for the primes up to $X''$ (\texttt{sum\_log\_primes\_le}), and because the sum of
$\log p$ over the prime factors of $D_\alpha$ is at most $\log D_\alpha$
(\texttt{sum\_log\_primeFactors\_le}). This display is the statement \texttt{nprimeMu\_at}. The source then
uses that each of $\log g_0,\log g_1,\log d$ is at most $X''$ (from the size conditions \texttt{SizeG},
\texttt{SizeQD}), and two further hypotheses on the family:
$\log\mathrm{Lv}_\alpha\le C_{\rm lv}X''$ and $\log D_\alpha\le C_{\rm den}X''$ (\texttt{MuTu},
\texttt{MuDGam}). This gives the constant $2(\log4+3+C_{\rm lv}+C_{\rm den})$ in \texttt{NprimeMuR}
(\texttt{nprimeMuR\_at}, \texttt{linkMu}). It also assumes that the glued lattices are finitely generated
and have positive rank (\texttt{MuFG}, \texttt{MuRank}).
These hypotheses are assumed here. The names of this item were not put to the dependency computation.
\item \emph{The first lattice of the chain, and the ``$+1$''.} The source has
$\mu(\Lambda^{\rm src}_\alpha)=\mu_{0,\alpha}-f\log(\mathrm{Lv}_\alpha-1)/(\mathrm{Lv}_\alpha-2)$ exactly
(\texttt{adm26}), and $\sum_w\mathrm{sz}_w(\tau(w))\le b^2$ (a hypothesis of \texttt{places\_at\_of}; the source proves it as \texttt{ta\_le}). With (a), the part of the first term of
\eqref{lt:Lodef} with $\tau$ is at most $4fb^2/\mathrm{Lv}_\alpha$. The source requires
$4fb^2/\mathrm{Lv}_\alpha+f\log(\mathrm{Lv}_\alpha-1)/(\mathrm{Lv}_\alpha-2)\le1$ (\texttt{HypLow2}). So the
cost of the rank on $\tau$ and the change from $\mu(\Lambda^{\rm src}_\alpha)$ to $\mu_{0,\alpha}$ are
together at most $1$. This is the single ``$+1$'' of \eqref{lt:Lo}. There is one such term, not two.
\end{enumerate}

\subsection*{3.\quad What \texttt{LinkG} is, and where the valuation statement lives}

The statement \texttt{LinkG} of the source has five clauses (\texttt{LinkG\_at}). They are:
(1) the lattice clause, of which \textup{(Lat)} is the part used in \S2; (2) a bound for the auxiliary numbers \texttt{Sv}; (3) the
inequality \eqref{dv:link} of \S III.5, with $0\le\mathrm{gain}(j,w)\le4|\mathcal T|a/p_w$, at the places of
class \texttt{signal} or \texttt{inertClass} with $p_w\le X''$; (4) a clause that holds by definition; (5) a
bound for the mean of certain numbers $\kappa(\mathcal Q)$, used in \S5. None of them mentions a theta
function, a Taylor coefficient, or a pairing. As \S III.5 already said, \eqref{dv:link} is an inequality
between the numbers $\mathrm{gain},\tau,\mathrm{wPart}$. A hypothesis on the family (the field
\texttt{fills}, of type \texttt{ParametersFill}) says that these numbers are given by explicit formulas
(\texttt{gainP}, \texttt{tauP}, \texttt{fill}, \texttt{partner}, \texttt{TTP}); they are recalled in \S4.
With that hypothesis, clause (3) is arithmetic about the formulas. Proposition~L10 below carries this out
at the primes of class \texttt{signal}, with its own loss terms.

\emph{The zero chain.} The proof of \texttt{LinkG} in the source (\texttt{linkG\_without\_kernels}) goes
through an intermediate statement \texttt{admG1}, one part of which does speak of the valuation of a chain
of maps on a local lattice (\texttt{vKge}). On the route to the final theorem that part is applied to the
chain in which every map is zero (\texttt{zeroCh}). The valuation of $0$ is infinite, so the part holds
for a trivial reason (\texttt{vKge\_zero}). So \texttt{LinkG} is proved without any information on the
theta kernels. This is consistent: \texttt{LinkG} is the bookkeeping, and the lattice clause (1).

\emph{Where the divisibility is asserted.} It is the statement \textup{(Val)}. The final theorem depends
on \texttt{LocalQuadraticParts.adm2}. The route, read downward from the place where it is consumed, is
this.
\begin{enumerate}[label=\textup{(\alph*)},leftmargin=2.4em,itemsep=1pt,topsep=2pt]
\item \texttt{divClausesR\_of} takes a structure of inputs, \texttt{DivInputsR}, and produces four
clauses. One field of \texttt{DivInputsR} is \texttt{val2}, that is \eqref{lt:V}. The other fields are
independent of it. From \texttt{val2} comes the first clause only: the integrality of Theorem~L4\,(i), by
\texttt{intoDual} and \texttt{divIntoDualShift} (Lemma~L3), which use the field \texttt{l1b}. The second
clause bounds the squares of the operator norms of the pairing at the complex places by numbers whose mean
logarithm is the height of the record minus a constant. It comes
from the height inputs (the fields \texttt{ch}, \texttt{hn}, \texttt{nP}, \texttt{nKcb}, \texttt{eig6}).
The third clause is the inequality of sizes, first line of \eqref{II:t6.5} of Part~II, with its term $-F_0$. It
comes from the size inputs (among them \texttt{L3main}, \texttt{Qb}, \texttt{deQ}, \texttt{errB}, and the
inequality \texttt{constant\_ge\_twenty\_eight}). The fourth clause is the slope identity of
Lemma~L2\,(ii), by \texttt{shiftMu}; it uses that \texttt{Nu} is a lattice of positive rank
(\texttt{Nu\_lat}, \texttt{a\_pos}). So the valuation statement does not imply the height bound or the
inequality of sizes. The inequality of sizes is the subject of \S III.4.
\item \texttt{val2AllAtRF} deduces \texttt{val2} from \texttt{adm2}, by linearity, under the hypotheses
described in \S2. In the seventh run of the dependency computation \texttt{val2AllAtRF} is marked as used
by the final theorem. The archive has a second formulation, \texttt{DivisionPackage.val2}. A later run of the dependency computation
shows that the final theorem does not depend on it; of the five declarations named \texttt{val2}, it depends on
\texttt{DivInputsR.val2} only.
\item \texttt{admissible\_level\_two} proves \texttt{adm2}, place by place. At the places above
$\mathrm{Lv}_\alpha$ and above the two auxiliary level primes $g_0,g_1$ of Part~II, \S6.7 it uses separate
inputs. At all other places it uses the statement \texttt{Conclusion1} of the file
\texttt{WeilJacquetShalikaLaw.lean}: \emph{for every $y$ in the local lattice at $w$ that defines
$\Lambda^{\rm div}_\alpha$, the image of $y$ under the composite of the $m$ maps of the chain is divisible,
in the last lattice of the chain, by $\mathfrak p_w$ to the power
$e_w\big(\sum_j(\mathrm{gain}(j,w)+\phi(j,w)-\mathrm{err}(\alpha,j,w))+\mathtt{sp}\cdot v_{p_w}(D_\alpha)\big)$.}
Here $D_\alpha$ is an integer denominator of the member, $v_{p_w}$ is the $p_w$-adic valuation, and
\texttt{sp} is a real parameter of the Lean statement. On the route \texttt{sp} is $1$: the statement
\texttt{Serv} contains \texttt{Conclusion1} with this value. The term with $D_\alpha$ does not reach
\eqref{lt:V}. The theorem \texttt{period\_geTU}, which passes from the chain to the period, lowers the
exponent by exactly $v_{p_w}(D_\alpha)$ (\texttt{vDG}), and the two cancel. This is the statement about
kernels.
\item \texttt{Conclusion1} is part of \texttt{ServF}, which \texttt{sum\_over\_primes\_service} proves
from the structure \texttt{SumOverPrimesHypotheses}, which has the fields \texttt{twU} and \texttt{pdv},
a bound for the denominators $D_\alpha$
(\texttt{DGamSize}), and the theorem \texttt{split\_bound\_eleven\_one}. The latter is the local
estimate: it rests on \texttt{thmAOp}, the Lean form of Theorem~J of \S III.5. The hypothesis \texttt{pdv}
contains hypothesis (B) of \S III.5 for the kernels of the construction.
\end{enumerate}
So there are two branches. One (\texttt{LinkG}) says what the numbers $\mathrm{gain}$ and $\tau$ are and
that they dominate the sums $V$; it needs no kernel. The other (\texttt{ServF}, then \textup{(Val)}) says
that the kernels are divisible by these numbers; it needs Theorem~J and hypothesis (B), and it is assumed
here. They meet in Theorem~L4 and in \S5.

\subsection*{4.\quad From Theorem~J to a gain at a place}

\textbf{Cells, and sums over one prime.} Let $H$ be the group of Part~II, \S4.1, of order $n$. It is the
group of \S III.3, and it is written additively here, as there and as in the Lean; $\mathcal T\subset H$.
A \emph{cell} is a pair $(0,t)$ or $(1,t)$ with $t\in H$; the first entry is the \emph{row}. There are
$f=2n$ cells. For an odd prime $p$ let $h_p\in H$ be its class, of order $n_p$. The classes of primes are
those of the Lean (\texttt{primeClass}): a prime $p\nmid2qd$ is \texttt{splitK} if $\chi_d(p)=1$; if
$\chi_d(p)=-1$ it is \texttt{signal} when $\chi_q(p)=1$ and \texttt{inertClass} when $\chi_q(p)=-1$. For an
odd prime $p\ne q$, $t\in H$ and an integer $i\ge1$ put
\begin{equation}\label{lt:cell}
 \mathrm{cell}(p,t,i)=\Big(\frac{1+\chi_q(p)^i}2,\ t+i\,h_p\Big).
\end{equation}
This is the cell of the prime power $p^i$ in the table of Part~II, \S4.1, for the column $t$. For odd $i$
the row is $1$ if $p$ is \texttt{signal} and $0$ if $p$ is \texttt{inertClass}. (Lean, for odd $i$:
\texttt{partner}.) For each $j$ there are weights $\mathrm{wt}_j>0$ on the cells (Lean: \texttt{fill});
these are the weights of Part~II, \S4.2, for the translate $u$ of the record and the map $j$. The
hypothesis \texttt{fills} says $\mathrm{wPart}(j,t,p,i)=\mathrm{wt}_j(\mathrm{cell}(p,t,i))$. Put
\begin{equation}\label{lt:Vdef}
 V_\infty(j,t,p)=\sum_{i\ \mathrm{odd}}\Big(\frac{a_t}{p^i}-\mathrm{wt}_j\big(\mathrm{cell}(p,t,i)\big)\Big)_+ ,
\end{equation}
the sum being over all odd $i\ge1$. The sum $V(j,t,p)$ of \eqref{dv:V} has the same terms, but only
those with $p^{i-1}\le a_t$. \emph{If all the weights are $\ge1$, then $V_\infty=V$}: the terms with
$p^i>a_t$ vanish, and \eqref{dv:V} contains every $i$ with $p^i\le a_t$. In the source the weights that
occur are $\ge1$ (\S III.5, \S1). Without this hypothesis the two sums can differ; see the example below.

\textbf{Lemma L6 (the prime sum).} \emph{Let the weights be any positive numbers. The sum \textup{\eqref{II:t4.1}}
of Part~II, for these weights, is}
\begin{equation}\label{lt:P}
 P_{u,t,j}=\sum_{\substack{p\ \mathtt{signal}\\\text{or }\mathtt{inertClass}}}2\log p\ V_\infty(j,t,p)\ +\
 \log d\sum_{i\ge1}\Big(\frac{a_t}{d^i}-\mathrm{wt}_j\big(\mathrm{cell}(d,t,i)\big)\Big)_+ .
\end{equation}
\emph{Proof.} In \eqref{II:t4.1} the coefficient of $p^i$ is $(1-\chi_d(p^i))\log p$. For $p\ne d$ it is $2\log p$
if $\chi_d(p)=-1$ and $i$ is odd, and $0$ otherwise. The row and the column of $p^i$ are those of
\eqref{lt:cell}. The prime $d$ gives the last term, with coefficient $\log d$ for every $i$. The primes $2$
and $q$ are absent, and split primes have coefficient $0$. \qed

\noindent In the running example of Part~II, \S4.2 ($q=11$, $d=1319$, $n=4$, $a_t=10000$) the weights are
$0.05$ and $10$. Computing the right side of \eqref{lt:P} gives $84701.0317$ for the translate $[1]$, the
value of Part~II, and the four values for the four translates agree with \eqref{II:t4.1} to $10^{-9}$. The example
uses $V_\infty$, not $V$, because its weights $0.05$ are below $1$. Here the two happen to agree: no odd power
$p^i$ of a prime of class \texttt{signal} or \texttt{inertClass} has $p^{i-1}>10^4$ and $10^4/p^i>0.05$.
(For the example of the first version, $q=3$, $d=23$, the right side with $V$ is smaller by $0.309064$ for
the translate $[1]$: the powers $5^7$ and $11^5$ are lost.) With the weights $10$ and $1$ the difference is
$0$ in any case. The primes of class \texttt{inertClass}
contribute $48943.2040$ for every translate: their odd powers lie in row $0$, which contains no
distinguished cell. The primes of class \texttt{signal} contribute $35703.7167$ for the translate $[1]$.
The distinguished weights carry the information on the zeros. Among the unramified primes, only
\texttt{signal} primes sample a distinguished weight; this is the reason for the name. The powers of $d$
contribute the separate last term of Lemma~L6. In general that term depends on the translate, because the
even powers of $d$ lie in row $1$. In this example only $d$ itself is below $a_t/0.05$; it lies in the cell
$(0,[7])$, since $\chi_{11}(1319)=-1$ and $1319\equiv7\bmod16$, and the last term is
$\log1319\,(10^4/1319-0.05)=54.1110$ for every translate.

\textbf{The formulas for the gain.} From here on the weights are $\ge1$, as in the source, so that
$V_\infty=V$. Let $p$ be \texttt{signal} or \texttt{inertClass}. The group generated by $h_p$ (by $2h_p$
if $p$ is \texttt{signal}) acts on $H$ by translation. For an orbit $\mathcal Q$ write
$\mathcal T(\mathcal Q)=\mathcal T\cap\mathcal Q$. (The Lean fixes a base point on each orbit,
\texttt{reps}.) Put $\mathrm{ka}_j=\frac12\big(k+(2j-1)(t_0+1)\big)$, with the integer parameter $t_0$ of
\S III.5. The source has $0\le\mathrm{ka}_j\le\mathrm{Sc}$ for $1\le j\le m$ (\texttt{kaJ\_nonneg} for the lower
bound, \texttt{kaJ\_le} for the upper bound). For an
orbit $\mathcal Q$ put
\begin{equation}\label{lt:g0}
 \mathrm{gain}^{\rm plain}(j,\mathcal Q)=\sum_{t\in\mathcal T(\mathcal Q)}\
 \sum_{\substack{i\ \mathrm{odd}\\ i\notin\mathrm{Ex}(p,t)}}
 \Big(\frac{a_t-8\,\mathrm{Sc}}{p^i}-\mathrm{wt}_j\big(\mathrm{cell}(p,t,i)\big)+\mathrm{ka}_j-4\,\mathrm{Sc}\Big)_+
 -\mathrm{dl}(\mathcal Q),
\end{equation}
\begin{equation}\label{lt:dl}
 \mathrm{dl}(\mathcal Q)=C_1\,n(\mathcal Q)\,N_m^2\log_p(2a)+C_2\,|\mathcal T(\mathcal Q)|\,\big(\log_p(2a)\big)^2 .
\end{equation}
Here $n(\mathcal Q)$ is the number of \emph{sites} of the orbit. For an \texttt{inertClass} prime the
sites are the points of the orbit, and $n(\mathcal Q)=n_p$. For a \texttt{signal} prime a site is a pair
consisting of the parity of an exponent $i\ge0$ and the column $t+i\,h_p$, with $t\in\mathcal Q$; there are
two parities and, for each, as many columns as $\mathcal Q$ has points, so $n(\mathcal Q)$ is twice the
order of $2h_p$. (Lean: the sites are the values of \texttt{partner}, and $n(\mathcal Q)$ is \texttt{nP}.
The sites are not the cells \eqref{lt:cell} of the powers of $p$: the two agree for odd $i$ only.) The
numbers $C_1,C_2\ge0$ are constants, and $\mathrm{Ex}(p,t)$ is the set
of odd $i$ for which $\mathrm{cell}(p,t,i)=(0,t')$ with $t'\in\mathcal T$, $t'\ne t$, and $i$ is the least
positive exponent with this property for $t'$. (Lean: \texttt{g0}, \texttt{dl}, \texttt{Excl}, \texttt{al},
\texttt{capU}, \texttt{kaJ}; the number $4\,\mathrm{Sc}$ is \texttt{cG}.) The gain is
\begin{equation}\label{lt:gainP}
\begin{aligned}
 \mathrm{gain}(j,p)={}&\sum_{\mathcal Q\in\mathrm{Srv}(p)}
   \max\big(\mathrm{gain}^{\rm srv}(j,\mathcal Q)-\mathrm{Pr}(j,\mathcal Q),\,\mathrm{gain}^{\rm plain}(j,\mathcal Q),\,0\big)\\
   &+\sum_{\mathcal Q\notin\mathrm{Srv}(p)}\max\big(\mathrm{gain}^{\rm plain}(j,\mathcal Q),0\big)
\end{aligned}
\end{equation}
for $p\le X''$, and $\mathrm{gain}(j,p)=0$ for $p>X''$ and for \texttt{inertClass} primes with $h_p=0$.
Here $\mathrm{Srv}(p)$ is the set of \emph{served} orbits, defined in \S5; it is empty unless $p$ is
\texttt{inertClass} and $p<X'$. The numbers $\mathrm{gain}^{\rm srv}$ and $\mathrm{Pr}$ are further
explicit expressions (Lean: \texttt{gU}, \texttt{Pr}): $\mathrm{gain}^{\rm srv}$ is \eqref{lt:g0} without
the exclusion of $\mathrm{Ex}$ and with the full $(a_t-8\,\mathrm{Sc})\,p/(p^2-1)$ for the points $t$ at
even positions on the orbit, by the identity \eqref{II:t6.1} of Part~II. By \texttt{fills},
$\mathrm{gain}(j,w)=\mathrm{gain}(j,p_w)$ and $\tau(w)$ depends only on $p_w$; we write $\tau(p)$ for it.

\textbf{Comparison with Theorem~J.} Theorem~J of \S III.5 is applied to one orbit: the sites are the
sites of the orbit just defined, the top sites are the sites whose column lies in $\mathcal T(\mathcal Q)$
(for a \texttt{signal} prime: those of even parity among them), the number
$a_t-8\,\mathrm{Sc}$ plays the role of the length of the first row at the top site $t$, and
$\mathrm{wt}_j-\mathrm{ka}_j+4\,\mathrm{Sc}$ that of the capacity at a site. The term $\mathrm{top}$ of
\eqref{dv:B} then has the shape of the double sum in \eqref{lt:g0}, and the last term of \eqref{dv:B} has
the shape of the first term of \eqref{lt:dl} (Lemma~L8). This dictionary, the treatment of the excluded
exponents, the other terms of \eqref{dv:B}, and the passage from one function to the lattices of the
chain are in the source (\texttt{normCor7}, \texttt{split\_bound\_eleven\_one},
\texttt{orbit\_entry\_bounds\_hold}, \texttt{lift\_bounds\_of\_essential}, then
\texttt{WeilJacquetShalikaLaw.Proof.stmt\_proved}); they are assumed here.

\emph{The statement \texttt{asmA\_top}.} The final theorem uses this statement for the top row (it is
marked as used in the seventh run of the dependency computation). It was read and compared with Lemma~6 of \S III.5. In this paragraph $s$ is a site, $\ell$ is the prime, and
$\nu_1$, $\mathrm{cap}$, $\mathrm{top}$, $d_{\rm top}$ are as in \S III.5, \S1; $i$ is an odd integer. The
Lean statement is: for every site $s$ and every finite set of odd integers $i\ge1$ there is a real number
$\omega\ge0$, with $\omega\ge\ell^{-d_{\rm top}(s)}$ if $d_{\rm top}(s)$ is finite, such that
\begin{equation}\label{lt:asmA}
 \mathrm{top}(s)+\nu_1(s)\,\omega+\sum_{i\ \text{in the set}}\mathrm{cap}(s+i)\,\big(1-\ell^i\omega\big)_+
 \ \le\ \nu_1(s)\,\frac{\ell}{\ell^2-1}.
\end{equation}
It relates to Lemma~6 of \S III.5 as follows.
\begin{enumerate}[label=\textup{(\alph*)},leftmargin=2.4em,itemsep=1pt,topsep=2pt]
\item \emph{If $d_{\rm top}(s)$ is finite}, take the set of all odd $i<d_{\rm top}(s)$. Then \eqref{lt:asmA}
is \eqref{dv:topgain} of that lemma, with $A=\nu_1(s)$, $c_i=\mathrm{cap}(s+i)$ and with $d_{\rm top}(s)$ as the
odd bound $J$ of the lemma: the same statement. Here $\mathrm{top}(s)$ is the whole of the second line of
\eqref{dv:B}: the sum over odd $h<d_{\rm top}(s)$ and the sum of $\nu_1(s)/\ell^h$ over odd
$h>d_{\rm top}(s)$ (the Lean development's \texttt{topTerm} has both). It is the right side of \eqref{dv:topgain}, and
$\nu_1(s)\,\ell/(\ell^2-1)=\gamma A$, with $\gamma=\ell/(\ell^2-1)$ as in \S III.5. So \eqref{lt:asmA} is \eqref{dv:topgain} with $\mathrm{top}(s)$ moved
to the left. Both assert only that some $\omega$ exists; neither holds for every $\omega\ge\ell^{-d_{\rm top}(s)}$.
For example, for $\nu_1(s)=9$, $\ell=3$, $d_{\rm top}(s)=1$ one has $\mathrm{top}(s)=\frac38$ and
$\gamma A=\frac{27}8$; both hold with equality at $\omega=\frac13$ and both fail at $\omega=\frac38$. A larger set gives nothing more, because the terms with $i\ge d_{\rm top}(s)$ vanish when
$\omega\ge\ell^{-d_{\rm top}(s)}$. The Lean proof takes the same radius as the proof of the lemma:
$\omega=\ell^{-i_0}$, where $i_0$ is the least odd $i<d_{\rm top}(s)$ with
$\mathrm{cap}(s+i)\ge\nu_1(s)/\ell^i$, and $i_0=d_{\rm top}(s)$ if there is none.
\item \emph{If $d_{\rm top}(s)=\infty$}, the Lean statement is not Lemma~6, which has a finite bound. It is
the remark after Lemma~6. The Lean allows any finite set of odd $i$; the remark has the restriction
$\ell^i\le\nu_1(s)$, and the Lean does not.
\item The Lean has $\ell$ any integer $\ge2$, and $\nu_1$, $\mathrm{cap}$ natural numbers.
\end{enumerate}
(Tested: $25600$ exact cases for (a), $30000$ for the radius; and $40000$ further exact cases in which the two inequalities were evaluated from the definitions at the radius of the lemma and at an arbitrary $\omega$, and agreed.)

\emph{The statement \texttt{lemma\_R\_log\_loss}.} The final theorem also uses this statement (seventh
run). It is the bound for one Taylor coefficient on the route of the final theorem. Its comparison with
Proposition~5 is in \S III.5, in the paragraph on the declarations the final theorem depends on:
Proposition~5 is the stronger of the two bounds for a coefficient. The formal statement
(\texttt{LemmaRA}, in the file \texttt{GaussJetBounds/StatementG.lean}) is for a multi-index whose entries
are all at most an integer $a\ge1$, and it assumes a test on the pure monomials below the multi-index. Its
exponent has the explicit loss $|\iota|\,(\log_\ell a+1)$, where $\iota$ is the finite index set of the
variables and $\ell$ is the prime. (In this sentence $a$ is the bound of the Lean statement.)

We now prove what is elementary: three comparisons of losses, and then an inequality of the shape of
\eqref{dv:link} at the \texttt{signal} primes.

\textbf{Lemma L7 (shifted positive parts).} \emph{Let $p\ge2$, $A\ge0$, $c\ge0$, and for each odd $i\ge1$
let $\psi_i\ge1$ and $c_i\ge0$ be real numbers. Let $\mathrm{odd}_p(A)$ be the number of odd $i\ge1$ with
$p^i\le A$. Then}
\begin{equation}\label{lt:shiftpos}
 \sum_{i\ \mathrm{odd}}\Big(\frac{A-2c}{p^i}-\psi_i+c_i-c\Big)_+\ \ge\
 \sum_{i\ \mathrm{odd}}\Big(\frac A{p^i}-\psi_i\Big)_+-c\Big(1+\frac2p\Big)\mathrm{odd}_p(A) .
\end{equation}
\emph{The left side may be $+\infty$. If $A\ge1$, then $\mathrm{odd}_p(A)\le\frac12(\log_pA+1)$.}

\emph{Proof.} For each $i$, $\frac{A-2c}{p^i}-\psi_i+c_i-c\ge\frac A{p^i}-\psi_i-c(1+\frac2{p^i})$. The positive
part is increasing, and subtracting a number $\ge0$ lowers a positive part by at most that number. If
$p^i>A$ then $A/p^i-\psi_i<0$, so the term on the right is $0$ and nothing is lost. The odd $i$ with
$p^i\le A$ are the odd $i$ with $1\le i\le\log_pA$; for $A\ge1$ their number is at most
$\frac12(\log_pA+1)$. \qed

\noindent The bound for $\mathrm{odd}_p(A)$ needs $A\ge1$: for $0<A<1/p$ the count is $0$ and
$\frac12(\log_pA+1)$ is negative. In the use below $A=a_t\ge1$. The left side of \eqref{lt:shiftpos} is
finite if $c_i\le c+\psi_i$ for all large $i$; this holds in the use below.

\textbf{Lemma L8 (three digit losses, one bound).} \emph{Let $p$ be a prime and $\bar a\ge p$ an integer.
Then $\mathrm{dig}_p(\bar a)\le\log_p\bar a+1\le2\log_p\bar a$, and:}
\begin{enumerate}[label=\textup{(\roman*)},leftmargin=2.4em,itemsep=1pt,topsep=2pt]
\item \emph{(the loss of \S III.5) the sum of $\mathrm{dig}_p$ over any $n_{\rm var}$ integers in
$[0,\bar a]$ is at most $n_{\rm var}\,\mathrm{dig}_p(\bar a)$; in the last term of \eqref{dv:B} the number
$n_{\rm var}$ of variables is the number $n_{\rm st}$ of sites times the two numbers $N_0,N_1$ of
\S III.5;}
\item \emph{(the loss $\Delta$ of \textup{\eqref{II:t6.1a}} of Part~II) if the index set there has $n_{\rm st}$
elements, its subset has $n_{\rm top}$ elements, the multiplicities there are $r_0,r_1$, and all the
integers whose digits are counted are $\le\bar a$, then
$\Delta\le n_{\rm st}\,(r_1+r_0r_1)\,\mathrm{dig}_p(\bar a)+n_{\rm top}\,\mathrm{dig}_p(\bar a)^2$;}
\item \emph{(the loss of the Lean) if $N_0,N_1,r_0,r_1\le N_m$, both bounds are at most
$2\,n_{\rm st}\,N_m^2\,\mathrm{dig}_p(\bar a)+n_{\rm top}\,\mathrm{dig}_p(\bar a)^2$. With
$n_{\rm st}=n(\mathcal Q)$, $n_{\rm top}=|\mathcal T(\mathcal Q)|$ and $\bar a\le2a$ this is at most
\eqref{lt:dl} with $C_1=4$, $C_2=4$.}
\end{enumerate}
\emph{Proof.} $\mathrm{dig}_p$ is increasing and $\mathrm{dig}_p(\bar a)=\lfloor\log_p\bar a\rfloor+1$.
(i) and (ii) are term-by-term bounds; for (ii) recall that $\Delta$ is a sum over the index set of $r_1$
digit counts and $r_0r_1$ digit counts, plus a sum over the subset of squares of digit counts.
(iii) $N_0N_1\le N_m^2$, $r_1+r_0r_1\le2N_m^2$ and $\mathrm{dig}_p(\bar a)\le2\log_p(2a)$. \qed

\noindent The paper has three digit losses: one digit count per variable (\S III.5); one per factorial,
with a square for the top sites (Part~II, \eqref{II:t6.1a}; Lean: \texttt{StarArith}, and the loss in
\texttt{gainA}); and the closed form \eqref{lt:dl} per orbit. They are three different quantities. The
lemma shows only that \eqref{lt:dl} is an upper bound for the first two. For example, with $p=3$,
$\bar a=80$, two sites, $N_0=N_1=2$ and $r_0=r_1=1$, the largest value of the first loss is $32$ and that
of the second, when both sites are in the subset, is $48$.

\textbf{Lemma L9 (the sum over the orbits).} \emph{For every prime $p$ of class \texttt{signal} or
\texttt{inertClass}, $\sum_{\mathcal Q}n(\mathcal Q)\le f$ and
$\sum_{\mathcal Q}|\mathcal T(\mathcal Q)|=|\mathcal T|$, the sums being over the orbits. Hence}
\begin{equation}\label{lt:dlsum}
 \sum_{\mathcal Q}\mathrm{dl}(\mathcal Q)\ \le\ C_1\,f\,N_m^2\log_p(2a)+C_2\,|\mathcal T|\big(\log_p(2a)\big)^2 .
\end{equation}
\emph{Proof.} The orbits partition $H$. For \texttt{inertClass} the orbit has $n_p$ points and
$n(\mathcal Q)=n_p$, so the sum is $n=f/2$. For \texttt{signal} the orbits are the cosets of the group
generated by $2h_p$, and $n(\mathcal Q)$ is twice their length, so the sum is $2n=f$. \qed

\noindent The right side of \eqref{lt:dlsum} is the term $\mathrm{fl}_G(p)$ of \eqref{dv:fl}, up to the
constant and to $\log_p(2a)\le2\log_pa$ (which holds when $a\ge2$), since $|\mathcal T|\le n=f/2$. This is
why $\mathrm{fl}_G$ is subtracted once for each of the $m$ maps in \eqref{dv:link}, and not once for each
$t$.

\textbf{Proposition L10 (the shape of \eqref{dv:link} at a \texttt{signal} prime).} \emph{Let $p\le X''$
be of class \texttt{signal}, let the weights be $\ge1$, and let $\mathrm{gain}(j,p)$ be given by
\eqref{lt:gainP} with $0\le\mathrm{ka}_j\le\mathrm{Sc}$. Then $\mathrm{Srv}(p)=\emptyset$, and for
$1\le j\le m$}
\begin{equation}\label{lt:signal}
 \mathrm{gain}(j,p)\ \ge\ \sum_{t\in\mathcal T}\Big(V(j,t,p)-4\,\mathrm{Sc}\Big(1+\frac2p\Big)\mathrm{odd}_p(a_t)\Big)
   -C_1\,f\,N_m^2\log_p(2a)-C_2\,|\mathcal T|\big(\log_p(2a)\big)^2 ,
\end{equation}
\emph{and $0\le\mathrm{gain}(j,p)\le\frac43\sum_{t\in\mathcal T}a_t/p\le4|\mathcal T|a/p$. Summing
\eqref{lt:signal} over $j$ gives, at the places $w$ above $p$,}
\begin{equation}\label{lt:signalsum}
 \sum_{j=1}^m\sum_{t\in\mathcal T}\Big(V(j,t,p)-\mathrm{lo}'(p)\Big)-m\,\mathrm{fl}'(p)\ \le\
 \sum_{j=1}^m\mathrm{gain}(j,w)-\tau(w),
\end{equation}
\emph{with $\tau(w)=0$,
$\mathrm{lo}'(p)=2\,\mathrm{Sc}\,(1+\frac2p)(\log_p(2a)+1)$ and
$\mathrm{fl}'(p)=C_1fN_m^2\log_p(2a)+C_2|\mathcal T|(\log_p(2a))^2$. This is \eqref{dv:link} with
$\mathcal T_w=\mathcal T$ and with $\mathrm{lo}'$, $\mathrm{fl}'$ in place of $\mathrm{lo}_G$,
$\mathrm{fl}_G$. If moreover}
\begin{equation}\label{lt:consts}
 C\ge\tfrac{20}3,\qquad C\ge2C_1,\qquad C\ge2C_2,\qquad a\ge2,\qquad \eta\ge0,\qquad p\le2a,
\end{equation}
\emph{then $\mathrm{lo}'(p)\le\mathrm{lo}_G(p)$ and $\mathrm{fl}'(p)\le\mathrm{fl}_G(p)$, and
\eqref{dv:link} itself holds at the places above $p$.}

\emph{Proof.} Served orbits exist only for \texttt{inertClass} primes, so $\mathrm{Srv}(p)=\emptyset$ and
$\mathrm{gain}(j,p)=\sum_{\mathcal Q}\max(\mathrm{gain}^{\rm plain}(j,\mathcal Q),0)\ge\sum_{\mathcal Q}\mathrm{gain}^{\rm plain}(j,\mathcal Q)$.
For a \texttt{signal} prime and odd $i$, $\mathrm{cell}(p,t,i)$ lies in row $1$, so $\mathrm{Ex}(p,t)$ is
empty. Apply Lemma~L7 with $A=a_t$, $c=4\,\mathrm{Sc}$, $\psi_i=\mathrm{wt}_j(\mathrm{cell}(p,t,i))$ and
$c_i=\mathrm{ka}_j$, sum over $t\in\mathcal T(\mathcal Q)$ and over the orbits $\mathcal Q$, and use
Lemma~L9. Since the weights are $\ge1$, the sum on the right of \eqref{lt:shiftpos} is $V(j,t,p)$. This
gives \eqref{lt:signal}. For the upper bound, $\mathrm{wt}_j-\mathrm{ka}_j+4\,\mathrm{Sc}\ge0$, so each inner
sum in \eqref{lt:g0} is at most $\sum_{i\ \mathrm{odd}}a_t/p^i=a_tp/(p^2-1)\le\frac43a_t/p$, and
$a_t\le2a$, $\frac83\le4$. For \eqref{lt:signalsum}: $\tau(w)=0$ when $\mathrm{Srv}(p)=\emptyset$ (the
definition of \texttt{tauP}), and $\mathrm{odd}_p(a_t)\le\frac12(\log_p(2a)+1)$ because $1\le a_t\le2a$.
For the last assertion, a \texttt{signal} prime is odd, so $1+\frac2p\le\frac53$; and $p\le2a$, so
$\log_p(2a)+1\le2\log_p(2a)$. Hence
$\mathrm{lo}'(p)\le\frac{20}3\mathrm{Sc}\log_p(2a)\le C(\mathrm{Sc}+k+\eta)\log_p(2a)=\mathrm{lo}_G(p)$,
since $\mathrm{Sc},k,\eta\ge0$ and the terms $\mathrm{pr}$, $\mathrm{or}$ of \eqref{dv:lo} vanish at a
\texttt{signal} prime. For $a\ge2$, $\log_p(2a)\le2\log_pa$. So
$C_1fN_m^2\log_p(2a)\le2C_1fN_m^2\log_pa$, and
$C_2|\mathcal T|(\log_p(2a))^2\le4C_2|\mathcal T|(\log_pa)^2\le2C_2f(\log_pa)^2$ because
$|\mathcal T|\le n=f/2$. \qed

\noindent The proposition does not show that the constants of the source satisfy \eqref{lt:consts}. In the
Lean the comparison of the losses is \texttt{lo\_le\_loG} and \texttt{fl\_le\_flG}. They say: for every real
number $C_*$, the Lean development's own losses \texttt{lo} and \texttt{fl} with the constant $C_*$ are at most
\texttt{loG} and \texttt{flG} with the constant $4|C_*|+1$. This is the source's version of the conditions
\eqref{lt:consts}.
(The sums in \eqref{lt:g0} are finite because the capacities are $\ge0$. Proposition~L10 was tested on
$3000$ random toy data, with $H$ cyclic of order $8$, $16$, $32$, and again on $30170$ exact cases: no
violation, and the lower bound is attained; without the loss terms the inequality fails in nearly all
cases.) For the primes of class \texttt{inertClass} the same computation applies to the orbits that are
not served, but there $\mathrm{Ex}(p,t)$ is not empty and the labels on these orbits do not belong to
$\mathcal T_w$. On the served orbits the source uses $\mathrm{gain}^{\rm srv}-\mathrm{Pr}$ and pays
$\tau(w)$; these orbits satisfy the conditions (G1)--(G3) of \S III.3. That case is assumed
(\texttt{AuxBlock8.Pf.stmtF}, \texttt{linkG\_at\_core}). The terms $\mathrm{pr}(p)$ and $\mathrm{or}(p)$ of
\eqref{dv:lo} occur only there.

\subsection*{5.\quad The sum over the primes}

In this subsection the weights are $\ge1$, the record and its translate $u$ are fixed, and
$\mathrm{gain}(j,w)$ and $\tau(w)$ depend only on $p_w$. The set $\mathcal T_w$ of \eqref{dv:link} also
depends only on $p_w$. From here on it is written $\mathcal T_p$, with $p=p_w$.

\textbf{Proposition L11 (splitting).} \emph{Let $\Xi$ be a finite set of primes of class \texttt{signal}
or \texttt{inertClass}. Suppose that $V(j,t,p)=0$ for all $j$, $t$ and all primes $p\notin\Xi$ of these two
classes, and that}
\begin{equation}\label{lt:notinert}
 \mathrm{gain}(j,w)=0\ \ (1\le j\le m)\quad\text{and}\quad\tau(w)=0\qquad
 \text{at every place $w$ with }p_w\notin\Xi .
\end{equation}
\emph{Put $\mathrm{def}(p)=\sum_{j}\sum_{t\in\mathcal T}V(j,t,p)-\big(\sum_j\mathrm{gain}(j,p)-\tau(p)\big)$.
Then}
\begin{equation}\label{lt:spliteq}
 \sum_{t\in\mathcal T}\sum_{j=1}^mP_{u,t,j}-G_{\rm top}=\sum_{p\in\Xi}2\log p\ \mathrm{def}(p)
 +\sum_{t\in\mathcal T}\sum_{j=1}^m\log d\sum_{i\ge1}
   \Big(\frac{a_t}{d^i}-\mathrm{wt}_j\big(\mathrm{cell}(d,t,i)\big)\Big)_+ .
\end{equation}
\emph{Let $\Xi_{\rm s}\subset\Xi$ be a set of primes at which \eqref{dv:link} holds with a set
$\mathcal T_p\subset\mathcal T$, and suppose $\mathrm{gain}(j,p)=\tau(p)=0$ for
$p\in\Xi\smallsetminus\Xi_{\rm s}$. Then}
\begin{equation}\label{lt:split}
\begin{aligned}
 \mathrm{def}(p)&\le\sum_j\sum_{t\in\mathcal T\smallsetminus\mathcal T_p}V(j,t,p)+m\,|\mathcal T_p|\,\mathrm{lo}_G(p)+m\,\mathrm{fl}_G(p)
 &&(p\in\Xi_{\rm s}),\\
 \mathrm{def}(p)&=\sum_j\sum_{t\in\mathcal T}V(j,t,p)&&(p\in\Xi\smallsetminus\Xi_{\rm s}).
\end{aligned}
\end{equation}
\emph{Proof.} By \eqref{lt:notinert} only the places above the primes of $\Xi$ contribute to
$G_{\rm top}$, and by \eqref{lt:lam}, $G_{\rm top}=\sum_{p\in\Xi}2\log p\,(\sum_j\mathrm{gain}(j,p)-\tau(p))$.
With Lemma~L6 and $V_\infty=V$ this is \eqref{lt:spliteq}. For $p\in\Xi_{\rm s}$ split
$\sum_{t\in\mathcal T}$ into $\mathcal T_p$ and its complement and insert \eqref{dv:link}. \qed

\noindent The hypothesis \eqref{lt:notinert} is needed: $G_{\rm top}$ is a sum over all places, and $\Xi$
contains only primes of the two classes. In the source it holds at the places above the split primes,
above $2$, $q$, $d$, $g_0$, $g_1$ and above the levels. The Lean field \texttt{not\_inert} gives the
vanishing at every place whose prime is of neither of the two classes. That the level primes, among them
$g_0$ and $g_1$, are of class \texttt{splitK} is another field (\texttt{level\_split}). At $g_0$, $g_1$ the
gain vanishes also by the field \texttt{gain\_g}. Without it
\eqref{lt:spliteq} fails: a gain of $1$ at one place above a split prime changes the left side and not
the right. (Lean: \texttt{SplitEq}, \texttt{SplitClasses}. Tested on $15000$ random data with
\eqref{lt:notinert}: no failure. Without it \eqref{lt:spliteq} fails in all $15000$. If the labels outside
$\mathcal T_p$ are forgotten, \eqref{lt:split} fails.) So the prime sum exceeds $G_{\rm top}$ by four kinds
of terms: the losses $\mathrm{lo}_G,\mathrm{fl}_G$ at the primes where \eqref{dv:link} is used; the labels
not served at those primes; the whole of $V$ at the other primes; and the prime $d$.

\textbf{Lemma L12 (one prime).} \emph{For $p\ge2$, $V(j,t,p)\le a_t\,p/(p^2-1)\le\frac43\,a_t/p$.}
\emph{Proof.} $\sum_{i\ \mathrm{odd}}p^{-i}=p/(p^2-1)$, and $p^2/(p^2-1)\le\frac43$. \qed

\textbf{The classes of primes in the sum.} The source sorts the primes of class \texttt{signal} or
\texttt{inertClass} as follows (\texttt{clAt}, \texttt{NoGain}). Say that a prime $p$ is \emph{in a band}
if $p\ge p_0$ and some power $p^i$ lies within a factor $1\pm\mathrm{bw}$ of $X'$, for a small number
$\mathrm{bw}>0$ of the source. Below, $C_{\rm r}$ is a further constant of the source (Lean: \texttt{Cr}; it
is not the constant $C$ of \eqref{dv:fl}). A \emph{candidate} is an \texttt{inertClass} prime $p<X'$ whose class $h_p$
has order $n_p\ge x^{\varepsilon/4}$ and which is not in a band. For a candidate put
$\mathrm{Sig}(p)=m\sum\kappa(\mathcal Q)$, the sum over the \emph{good} orbits $\mathcal Q$, those on
which $\mathcal T$ satisfies (G1)--(G3) of \S III.3 with the parameters there; $\kappa(\mathcal Q)\ge0$ is a
number attached to the orbit (\texttt{kapv}). The set of served orbits $\mathrm{Srv}(p)$ of a candidate is
the set of good orbits if $\mathrm{Sig}(p)\le x^{-\varepsilon}$, and is empty otherwise; for all other
primes $\mathrm{Srv}(p)=\emptyset$. The set of served labels is
\begin{equation}\label{lt:Tp}
 \mathcal T_p=\mathcal T\ \ \text{if $p$ is \texttt{signal} or }p\ge X',\qquad
 \mathcal T_p=\{t\in\mathcal T:\ \text{the orbit of $t$ is in }\mathrm{Srv}(p)\}\ \ \text{otherwise.}
\end{equation}
There are four classes. A prime belongs to the first class in the following list whose condition it
satisfies; read in this order the list is a partition. (As conditions, ``high'' and ``first'' overlap.)
\begin{enumerate}[label=\textup{\arabic*.},leftmargin=2.0em,itemsep=1pt,topsep=2pt]
\item \emph{high}: $p>X''$. The gain is $0$. All weights are at least $a/X''$ (Lean: the field
\texttt{off\_thin} of the structure \texttt{Inst}, with $a/X''=w''$; the other cells have larger weights,
by \texttt{off\_top} and \texttt{off\_thick}). Since moreover
$a_t\le a(1+1/(kx^{13}))$, $V$ vanishes unless $p\le X''(1+1/(kx^{13}))$: a short interval.
\item \emph{second}: candidates with $\mathrm{Sig}(p)>x^{-\varepsilon}$. The gain is set to $0$.
\item \emph{first}: \texttt{inertClass}, $p<X'(1+C_{\rm r}\,\mathrm{bw})$, with $n_p<x^{\varepsilon/4}$ or $p$ in a
band; and \texttt{inertClass} with $n_p$ odd and $n_p>1$. The gain is set to $0$.
\item \emph{served}: all others. There \eqref{dv:link} holds with $\mathcal T_p$ as in \eqref{lt:Tp}.
\end{enumerate}
The class ``served'' contains the \texttt{inertClass} primes with $h_p=0$ and $p\ge X'(1+C_{\rm r}\,\mathrm{bw})$.
Their gain is $0$ by \eqref{lt:gainP}, and $V$ vanishes there, because $a_t/p$ is below the weights of the
cells $(0,t)$, $t\in\mathcal T$. These weights exceed $b'=a/X'$ (Lean: the field \texttt{bT\_ge} of
\texttt{SizeData}$_2$ for the label, and \texttt{off\_top} of \texttt{Inst} for the rest of the weight,
which is at least $b'/x$). The primes of the classes high, second and first are the ``dropped'' primes of \S2. The
other primes need no gain: split primes have coefficient $0$ in \eqref{II:t4.1}, and so do $2$ and $q$. Every odd
prime up to $x^{\varepsilon/(5C_o)}$ splits, where $C_o$ is the constant of the source that appears
after Theorem~\ref{thm:T} (Lean: \texttt{split\_small}, an assumed property of $d$, with
\texttt{three\_le\_split}; the Lean writes the exponent as \texttt{ez}$/5$, and \texttt{ez} is
$\varepsilon$ divided by a constant \texttt{J} of the scale, so $C_o$ is \texttt{J}). The paragraph after
Theorem~\ref{thm:T} writes $C_o\ge2$. The Lean structure of the scale has only $1\le\mathtt J$
(\texttt{one\_le\_J}). Whether $\mathtt J\ge2$ on the route was not determined; the statement that these
primes split does not depend on it. So the sum starts above that bound. At the
split primes, at $g_0$, $g_1$ and at the levels the numbers $\mathrm{gain}$, $\tau$, $\mathrm{io}$ are
$0$, and only $\phi$ and $\mathrm{err}$ occur. The cut-offs have this meaning: $X'=a/b'$ is where $a/p$
falls below the weights of the cells $(0,t)$, $t\in\mathcal T$; $X''$ is where it falls below the smallest
weights; below $p_0$ the source uses a different local argument.

\textbf{Lemma L13 (a label that is not served lies on a bad orbit).} \emph{Let $p$ be a served prime and
$t\in\mathcal T\smallsetminus\mathcal T_p$. Then $p$ is of class \texttt{inertClass}, $p<X'$,
$n_p\ge x^{\varepsilon/4}$, and the orbit of $t$ under $h_p$ is not good.}

\emph{Proof.} By \eqref{lt:Tp}, $p$ is \texttt{inertClass} and $p<X'$. Since $p$ is neither first nor
second, $n_p\ge x^{\varepsilon/4}$, $p$ is not in a band, and $\mathrm{Sig}(p)\le x^{-\varepsilon}$. So
$p$ is a candidate and $\mathrm{Srv}(p)$ is the set of good orbits. \qed

\noindent (Lean: \texttt{unserved\_bad}, which the final theorem uses by the seventh run of the
dependency computation.) This is where the conditions (G1)--(G3) are used in the sum over
primes. Since served orbits are good by definition, the conditions are also available in the proof of
\eqref{dv:link} on the served orbits (\S4); that proof was not read.

\textbf{Lemma L14 (the three small sums).} \emph{Let $\Xi_1$ be the set of \texttt{inertClass}
primes $p<X'$ with $n_p\ge x^{\varepsilon/4}$.}
\begin{enumerate}[label=\textup{(\roman*)},leftmargin=2.4em,itemsep=1pt,topsep=2pt]
\item \emph{Let $L_b$, $C_b$, $x_0$ and $L_2$ be as in Theorem~\ref{thm:T}, with $L_2\le L_b$, and let
$x\ge x_0$. Let $\mathcal T$ be the set $T$ of that theorem, for the set of primes $\Xi_1$.
If $x^{2\varepsilon/3}\ge16C_b$ and $\sum_{p<X'}\frac{\log p}p\le2x^{\varepsilon/3}$, then}
\begin{equation}\label{lt:bad}
 \sum_{p\in\Xi_1}\frac{\log p}p\,\#\{t\in\mathcal T:\text{ the orbit of $t$ under $h_p$ is not good}\}
 \ \le\ |\mathcal T|\,x^{-2+\varepsilon}.
\end{equation}
\item \emph{Hence the labels not served cost at most}
$\sum_{p\ \mathrm{served}}2\log p\sum_j\sum_{t\notin\mathcal T_p}V(j,t,p)\le\frac83\,m\,a'\,|\mathcal T|\,x^{-2+\varepsilon}$,
\emph{where $a'=\max_ta_t$.}
\item \emph{Let $\theta>0$ and $\mathrm{Sig}(p)\ge0$ on a finite set $\Xi_{\rm c}$ of primes. Then}
\begin{equation}\label{lt:markov}
 2\!\!\sum_{p\in\Xi_{\rm c},\ \mathrm{Sig}(p)>\theta}\frac{\log p}p\ \le\
 \sum_{p\in\Xi_{\rm c}}\frac{\log p}p\Big(2\cdot[\mathrm{Sig}(p)>\theta]+4\,\mathrm{Sig}(p)\Big)
 \ \le\ \Big(\frac2\theta+4\Big)\sum_{p\in\Xi_{\rm c}}\frac{\log p}p\,\mathrm{Sig}(p).
\end{equation}
\emph{Here $[\,\cdot\,]$ is $1$ if the condition holds and $0$ otherwise. Hence, if the middle sum is at
most $x^{-9}$, the primes of the second class cost at most} $\frac83\,m\,a'\,|\mathcal T|\,x^{-9}$.
\end{enumerate}
\emph{Proof.} (i) Theorem~\ref{thm:T} bounds the left side by
$4C_bx^{-2}\delta_1n\cdot2x^{\varepsilon/3}$, and $\delta_1n\le2|\mathcal T|$. So it is at most
$16C_bx^{-2+\varepsilon/3}|\mathcal T|\le x^{-2+\varepsilon}|\mathcal T|$. (ii) By Lemma~L12 each $V$ is at
most $\frac43a'/p$; by Lemma~L13 the labels counted are those of \eqref{lt:bad}. (iii) The first inequality
is clear. For the second, $[\mathrm{Sig}(p)>\theta]\le\mathrm{Sig}(p)/\theta$. The cost of the second class
is, by Lemma~L12, at most $\frac83ma'|\mathcal T|\sum_{\mathrm{second}}\frac{\log p}p$. \qed

\noindent (Lean: \texttt{ChBad\_at} is the inequality of Theorem~\ref{thm:T}; \texttt{ChBadp\_at} is
\eqref{lt:bad}; \texttt{OutD} is (ii); \texttt{lhsW}, \texttt{ChW\_at}, \texttt{term\_le},
\texttt{SecondD} are (iii). The bound $\sum_{p<Y}\frac{\log p}p\le2\log Y$ for $Y\ge1$ is the Lean development's
\texttt{mertens}; it is a form of Mertens' first theorem and is cited. By the seventh run of the
dependency computation the final theorem uses \texttt{ChBadp\_at}.) Since $a'\le2xb$, the two costs
are at most $6\,m\,|\mathcal T|\,b\,x^{-1+\varepsilon}$ and $6\,m\,|\mathcal T|\,b\,x^{-8}$. Let
$B=|\mathcal T|\,m\,b\,\delta/2$ be the unit of (M), with $\delta=|T_2|/n$ as in Part~II, \S4.2; this $\delta$ is not the $\delta_1$ of the scale. Since
$\delta\ge x^{-\varepsilon}$, the two costs divided by $B$ are at most $12\,x^{-1+2\varepsilon}$ and
$12\,x^{-8+\varepsilon}$. With $\varepsilon<\frac1{52}$ these are small for large $x$ only: the first
bound is $1.31$ at $x=10$.

\emph{Where the hypothesis of \textup{(iii)} comes from.} The numbers $a_t$ and the weights are still
free in short intervals. Clause (5) of \texttt{LinkG\_at} bounds the mean of
$\kappa(\mathcal Q)$ over these choices by $|\mathcal T(\mathcal Q)|\,x^{-12}$. So the mean of
$\mathrm{Sig}(p)$ is at most $m|\mathcal T|x^{-12}$ for each candidate $p$, and the mean of the right side
of \eqref{lt:markov} with $\theta=x^{-\varepsilon}$ is at most
$(2x^{\varepsilon}+4)\,m|\mathcal T|x^{-12}\cdot2x^{\varepsilon/3}$. This is at most $x^{-9}$ if
\begin{equation}\label{lt:sizes}
 m\le x^{1/2+\varepsilon/100},\qquad|\mathcal T|\le2x^{2+9\varepsilon},\qquad x\ge10^5 .
\end{equation}
Indeed the product is then at most $24\,x^{-9.5+10.35\,\varepsilon}\le24\,x^{-9.3}$. The first two bounds of
\eqref{lt:sizes} follow from three statements of the source: $m\le x^{\varepsilon/100}\sqrt{x\delta_1}$
(\texttt{SizeM}), $|\mathcal T|\le2n\delta_1$ (\texttt{SizeTT}) and $n\le x^{2+8\varepsilon+\varepsilon/100}$
(\texttt{SizeN}), with $\delta_1=x^{-\varepsilon/8}$. They are needed, and
without a bound for $m|\mathcal T|$ the inequality is false. Some choice of the weights does at least as
well as the mean. (Lean: \texttt{chosenW\_of}, \texttt{size\_alt\_alt}. The latter has the three statements
as hypotheses, and it assumes $24\le x^{1/4}$ where \eqref{lt:sizes} has $x\ge10^5$; the condition of
\eqref{lt:sizes} suffices for the inequality as stated here. The bound for the mean of $\kappa$
is part of \texttt{LinkG} and is assumed.)

\textbf{The other terms, and the total.} There remain: the losses $m|\mathcal T_p|\mathrm{lo}_G(p)$ and
$m\,\mathrm{fl}_G(p)$ at the served primes; the first class; the high class; and the prime $d$. For the
terms of $\mathrm{lo}_G,\mathrm{fl}_G$ in $\log_p$ one has, for $Y\ge2$,
$\sum_{p\le Y}\log p\cdot\log_p(2a)=\#\{p\le Y\}\cdot\log(2a)$ and
$\sum_{p\le Y}\log p\,(\log_pa)^2\le\#\{p\le Y\}\cdot(\log a)^2/\log2$. With $Y=X''$, and with the bounds
$\mathrm{Sc},\eta\le\exp(x^{\varepsilon/4})$ of the source, these are negligible against
$b=\exp(x^{\varepsilon})$ for large $x$. The term $b'/x$ of $\mathrm{pr}(p)$, summed with weight $2\log p$
over $p<X'$, gives $2\,(b'/x)\sum_{p<X'}\log p\le2\log4\cdot(b'/x)X'=2\log4\cdot b$ for each pair $(j,t)$,
by Chebyshev's bound $\sum_{p\le Y}\log p\le Y\log4$ (cited): a constant times the unit
$|\mathcal T|mb$. So this term alone is of the size of the right side of \eqref{lt:loss} below. The
complete statement is taken from the source: \emph{with the constant $C\ge0$ of \eqref{dv:fl} and with
$C'=1$, for large $x$}
\begin{equation}\label{lt:loss}
 \sum_{t\in\mathcal T}\sum_{j=1}^mP_{u,t,j}-G_{\rm top}\ \le\ (3C+13)\,|\mathcal T|\,m\,b,\qquad
 \sum_w\lambda_we_w\,\mathrm{io}(w)\ \le\ C'\,|\mathcal T|\,m\,b .
\end{equation}
(Lean: the constants $3C+13$ and $C'=1$ are those of \texttt{joint\_discriminant\_eventually}
(\texttt{ExcessAndPrimeLogSumUnitBounds}), which contains the bounds for the eleven sums; its proof was not
read. There $C$ is $\max(C_{\rm G},0)$ for the constant $C_{\rm G}$ of \texttt{LinkG}. The statement on the
route is \texttt{stLoss7}. It proves \texttt{Loss\_at} with exponent $0$ from \texttt{LinkG}, the size of
$\mathcal T$, \eqref{lt:bad} and the hypothesis of Lemma~L14\,(iii). It has one constant for both
inequalities of \eqref{lt:loss}, namely $3(1+|C_{\rm G}|+C_{\rm ord})+13$, where $C_{\rm ord}\ge0$ is a
constant of \texttt{joint\_discriminant\_eventually}. This is at least $3C+13$ and at least $1$, so it is the weaker statement. In the Lean the left side is written
with a sum \texttt{ASK}; that it equals $\sum P_{u,t,j}$ is Lemma~L6 together with the statement
\texttt{askIsF}, not read. By the seventh run of the dependency computation the final theorem uses
\texttt{stLoss7}.) The inequality \eqref{lt:loss} is not proved here: the estimates for the eleven sums and
their total bound are quoted. Proposition~L11 and Lemmas L12 to~L14 split the sum and bound the terms named
in them, and no more. In particular Theorem~\ref{thm:T} gives \eqref{lt:bad}, one of the inputs of
\texttt{stLoss7}; alone it does not give \eqref{lt:loss}.

\subsection*{6.\quad A model computation: how the terms would cancel}

This subsection does not prove the prime side of (M), for the source or for one member. It proves an
inequality under a hypothesis, called (a) below, that does not hold in the source. Its purpose is to show
which term of the earlier subsections meets which term of Part~II, \S6.7. The actual argument, for all
members of a record together and without hypothesis (a), is in \S\ref{as:sec}. There Proposition~M7
deduces the inequality for one translate from numbered Hypotheses, among them \eqref{lt:Lo}, which
follows from Theorem~L4(iii) and the statements (a)--(d) of \S2, and \eqref{lt:loss}. Theorem~L4(i)
enters through the mean-slope inequality (Hypothesis (H8) there), which Proposition~M4 derives from the
rank bound of \S III.8 and which takes the place of hypothesis (a). The Hypotheses of \S\ref{as:sec} are
not proved in this paper.

\begin{illbox}
Proposition~L15 is a model: its hypothesis (a) does not hold in the source, and it is not used in the proof.
\end{illbox}
\textbf{Proposition L15 (a model; hypothesis \textup{(a)} does not hold in the source).} \emph{Fix a
member $\alpha$ of a record with translate $u$. Assume the conclusions \textup{(i)} and \textup{(iii)} of
Theorem~L4, the first inequality of \eqref{lt:loss}, and:}
\begin{enumerate}[label=\textup{(\alph*)},leftmargin=2.4em,itemsep=1pt,topsep=2pt]
\item \emph{(model) the pairing gives an injective map
$\mathrm{Per}_\alpha:U_\alpha\to\mathrm{Hom}_{L_0}(U',L_0)$, and $\dim U_\alpha=\dim U'$; the measures on
$U_\alpha\otimes\mathbb R$ and on $\mathrm{Hom}_{L_0}(U',L_0)\otimes\mathbb R$ come from Euclidean norms
that are orthogonal sums over the $n_0/2$ complex places of $L_0$; and $\mathrm{ht}_\alpha$ is the mean
over these places of the logarithm of the square of the operator norm of $\mathrm{Per}_\alpha$ at the
place. \textbf{This hypothesis does not hold in the source.} There the rank of
$\Lambda^{\rm src}_\alpha$ is smaller than the rank of $\mathcal R$, there are many members, and no
inequality for one member is available: the source has only an inequality for a mean over the members of
the record (Lean: \texttt{Frame\_at}), and that inequality has further error terms (the logarithm of the
number of members, a term for the spread of the slopes, two error terms from second moments multiplied by
a range of slopes, and a factor close to $1$);}
\item \emph{(size) $\mu(\mathcal R^\vee)+\mathrm{ht}_\alpha-F_0\le\sum_{t\in\mathcal T}\big(m(a_t\log a_t-a_t-a_t\log\Omega_E)-a_t\log\mathrm{Lv}_\alpha\big)+E_{\rm size}$;}
\item \emph{(level) $|E_{\rm err}(\alpha)-m\log(g_0g_1)\sum_ta_t|\le E_{\rm lev}$;}
\item \emph{(count) $\log\mathrm{Lv}_\alpha\ge\sum_{j}\big(\langle\log\mathrm{wt}_j\rangle+\log(\mathrm{rd}_E/2)-\log\pi\big)+m\log(g_0g_1)-E_{\rm cnt}$;}
\item \emph{(source lattice) $\mu(\Lambda^{\rm src}_\alpha)\ge-E_{\rm src}$.}
\end{enumerate}
\emph{Here $\Omega_E$, $\mathrm{rd}_E$ are as in \textup{\eqref{II:t6.3}} of Part~II, $g_0,g_1$ are the level primes
of Part~II, \S6.7, $\pi=3.14159\ldots$, $\gamma$ is Euler's constant, and $\langle\log\mathrm{wt}_j\rangle$ is the mean of
$\log\mathrm{wt}_j$ over the $2n$ cells. Then, with
$M_{u,t,j}=a_t(\log a_t-1-\gamma+c_{\chi_d}-\langle\log\mathrm{wt}_j\rangle)$ as in \textup{\eqref{II:t4.2}} of
Part~II,}
\begin{equation}\label{lt:M}
 \sum_{t\in\mathcal T}\sum_{j=1}^m\big(P_{u,t,j}-M_{u,t,j}\big)\ \le\ (3C+13)|\mathcal T|mb+\mathrm{Lo}(\alpha)
 +E_{\rm size}+E_{\rm lev}+E_{\rm cnt}\sum_ta_t+E_{\rm src} .
\end{equation}
\emph{This is not \textup{(M)}. The inequality \textup{(M)} is a mean over the translates $u$; a record has
one translate, and the members $\alpha$ are not translates.}

\emph{Proof.} By Theorem~L4\,(i), $\mathrm{Per}_\alpha(\Lambda^{\rm gl}_\alpha)\subset\mathcal R^\vee$, and
by (a) both are lattices of the same rank, say $\mathrm{rk}$, in $\mathrm{Hom}_{L_0}(U',L_0)$. At each
complex place the space has real dimension $2\,\mathrm{rk}$, so the map multiplies volumes there by at
most the operator norm to the power $2\,\mathrm{rk}$. Taking the product over the $n_0/2$ places, the
covolume of $\mathrm{Per}_\alpha(\Lambda^{\rm gl}_\alpha)$ is at most
$\exp(\frac12n_0\,\mathrm{rk}\,\mathrm{ht}_\alpha)$ times that of $\Lambda^{\rm gl}_\alpha$. With
\eqref{lt:mu} this is
$\mu(\Lambda^{\rm gl}_\alpha)\le\mu(\mathrm{Per}_\alpha(\Lambda^{\rm gl}_\alpha))+\mathrm{ht}_\alpha\le\mu(\mathcal R^\vee)+\mathrm{ht}_\alpha$:
the first inequality of Part~II, \S6.8, in the squared normalisation, followed by Lemma~L1. Now
\begin{equation}\label{D:251}
 G_{\rm top}\le\mu(\Lambda^{\rm gl}_\alpha)-\mu(\Lambda^{\rm src}_\alpha)-F_0+E_{\rm err}(\alpha)+\mathrm{Lo}(\alpha)
 \le\big(\mu(\mathcal R^\vee)+\mathrm{ht}_\alpha-F_0\big)+E_{\rm src}+E_{\rm err}(\alpha)+\mathrm{Lo}(\alpha)
\end{equation}
by \eqref{lt:slope} and (e). The term $F_0$ has cancelled against the $F_0$ of (b). Insert (b), then (c)
and (d): the terms $m\log(g_0g_1)\sum_ta_t$ cancel, and the coefficient of $a_t$ for the map $j$ is
$\log a_t-1-\log\Omega_E-\langle\log\mathrm{wt}_j\rangle-\log(\mathrm{rd}_E/2)+\log\pi$, which is
$\log a_t-1-\gamma+c_{\chi_d}-\langle\log\mathrm{wt}_j\rangle$ by \eqref{II:t6.3} of Part~II. So
$G_{\rm top}\le\sum_{t,j}M_{u,t,j}+E_{\rm size}+E_{\rm lev}+E_{\rm cnt}\sum_ta_t+E_{\rm src}+\mathrm{Lo}(\alpha)$.
Add the first inequality of \eqref{lt:loss}. \qed

\noindent This is the computation of Part~II, \S6.7, with the lattices made explicit. It differs from the
source in four ways. (1) Hypothesis (a): in the source one map controls only the largest slope of a
sublattice of the target (Part~II, \S6.8), and the inequality
$\mu(\Lambda^{\rm gl}_\alpha)\le\mu(\mathcal R^\vee)+\mathrm{ht}_\alpha$ is replaced by the mean-slope
argument of Part~II, \S6.9--6.10, applied to all members at once. (2) The errors and the inputs of that
argument (second moments, the filtration by slopes) have no counterpart in \eqref{lt:M}. (3)
\eqref{lt:M} is for one member and one translate, and (M) is a mean over translates. (4) The inputs
(b)--(e) are modelled on Lean statements but are not those statements: (b) and (c) on clauses 7 and 8 of
\texttt{LinkDIV\_at}, (d) on \texttt{Count\_at}. For (e), the Lean statement \texttt{SlopesM\_at} bounds
the slope $\mu_{0,\alpha}$ of the first lattice of the chain, not $\mu(\Lambda^{\rm src}_\alpha)$; the two
differ by the term $f\log(\mathrm{Lv}_\alpha-1)/(\mathrm{Lv}_\alpha-2)$ of \S2\,(d). The proposition shows
what each lattice statement of this section is for: Theorem~L4\,(i) lets the size of the map bound the
slope of $\Lambda^{\rm gl}_\alpha$; \eqref{lt:slope} turns that into a bound for $G_{\rm top}$;
\eqref{lt:loss} turns $G_{\rm top}$ into the prime sum. (The arithmetic was tested on $25000$ random data
satisfying the hypotheses: no violation, equality when all hypotheses are equalities, and failure with a
wrong sign in three places.)

\subsection*{7.\quad Assumed here, from the source}

\begin{enumerate}[label=\textup{(\roman*)},leftmargin=2.4em]
\item \emph{The valuation statement \textup{(Val)}.} In the Lean it rests on three assumptions:
\texttt{LocalQuadraticParts.adm2}, for the elements that lie both in the lattice $\widetilde M(\mu_w)$
and in $V_0(\mu_w)$ (the final theorem depends on it); the hypotheses \texttt{ColBrF.Tied}
and \texttt{CutUp} of \texttt{val2AllAtRF}; and the structure field \texttt{l1b}. The final theorem depends
on \texttt{val2AllAtRF} (seventh run of the dependency computation). The third hypothesis of that theorem,
the inclusion into $\Lambda^{\rm div}_\alpha$, is supplied by \texttt{nu\_into}. It was not determined
whether \texttt{Tied}, \texttt{CutUp} and \texttt{l1b} are discharged on the route of the final theorem.
That the local images generate
$\mathcal R_w$ is a definition (\texttt{RR}) with the theorem \texttt{latticeGlue}, not an assumption. The
statement \texttt{adm2} is proved in the source by \texttt{admissible\_level\_two} from: the statement
\texttt{Conclusion1} on the composite of the chain (\S3\,(c)), which is part of \texttt{ServF}; the
integrality of the chain and of the period at the places where no gain is claimed
(\texttt{MemberChainIntegral}, \texttt{TwistedPointIntegral}, \texttt{PeriodClassSumWeyl},
\texttt{TwistedPeriodAdmissible}); separate statements at the places above $\mathrm{Lv}_\alpha$, $g_0$,
$g_1$; and the hypothesis \texttt{fills}. None of it is proved here.
\item \emph{\texttt{ServF} from Theorem~J.} \texttt{sum\_over\_primes\_service} deduces \texttt{ServF}
from \texttt{split\_bound\_eleven\_one} (which rests on \texttt{thmAOp}, \S III.5, and is conditional on
\texttt{Loc.AdmB}) and from the fields \texttt{twU}, \texttt{pdv} of its hypothesis
\texttt{SumOverPrimesHypotheses}; the
field \texttt{pdv} contains hypothesis (B) for the kernels (\texttt{EvenConditionsInert},
\texttt{JetIntegralityBis}, \texttt{ThetaKernelShapeAtD}, \texttt{Shape}). The steps between
(\texttt{orbit\_entry\_bounds\_hold}, \texttt{lift\_bounds\_of\_essential},
\texttt{four\_orbit\_conditions\_of\_proof}, \texttt{stmt\_proved}) were not read. In particular the
dictionary of \S4 between Theorem~J and \eqref{lt:g0} is not proved here: the frames, the labels, the
identification of the sites of an orbit with places, and the length of the first row.
\item \emph{The lattice statement \textup{(Lat)}}, with the integrality of $\tau$ (the first two parts of
the first clause of \texttt{LinkG\_at}; the last clause of \texttt{NuDef}, for the factor $e_w$;
\texttt{exists\_tau\_eq\_natCast} together with \texttt{fills}), and the definition \eqref{lt:N}
(\texttt{NprimeDef}). For the passage from \eqref{lt:slope} to the Lean statement \eqref{lt:Lo}: the rank
formula, the bound at the places above the level, the comparison with the first lattice of the chain
(\texttt{adm26}), the inequality of \texttt{HypLow2}, the bound $\sum_w\mathrm{sz}_w(\tau(w))\le b^2$
(\texttt{ta\_le}), and the bound for the rounding term. The last is the hypothesis \texttt{NprimeMuR} of
\texttt{places\_at\_of}; the source proves it by \texttt{linkMu} from \texttt{ZeroAbove}, from Chebyshev's
bound up to $X''$, and from bounds for $\log\mathrm{Lv}_\alpha$, $\log g_0$, $\log g_1$, $\log d$ and
$\log D_\alpha$ (\S2\,(c)), with the finite generation and the positivity of the ranks of the glued
lattices (\texttt{MuFG}, \texttt{MuRank}). That $\tau(w)=0$ above the level is the hypothesis \texttt{htauT}
of \texttt{places\_at\_of}. The archive also has a
statement \texttt{RoundingTheta}; it was not put to the dependency computation.
\item \emph{The inequality \eqref{dv:link} at the \texttt{inertClass} primes}, and the bound for the mean
of $\kappa(\mathcal Q)$ (clauses 3 and 5 of \texttt{LinkG\_at}; \texttt{linkG\_at\_core},
\texttt{AuxBlock8.Pf.stmtF}). The hypothesis \texttt{fills}, which identifies the data with the formulas of
\S4, and the bounds \texttt{kaJ\_nonneg}, \texttt{kaJ\_le}. At the \texttt{signal} primes: that the constants of the source satisfy
\eqref{lt:consts}, or the Lean development's own comparison \texttt{lo\_le\_loG}, \texttt{fl\_le\_flG}.
\item \emph{The sums over primes}: the estimates for the eleven loss sums and their total bound
\eqref{lt:loss} (\texttt{stLoss7}, \texttt{joint\_discriminant\_eventually}); the sizes \eqref{lt:sizes} (\texttt{SizeM}, \texttt{SizeTT}, \texttt{SizeN});
the vanishing \eqref{lt:notinert} (\texttt{not\_inert}, with \texttt{level\_split} and \texttt{gain\_g}); the two classical estimates cited in \S5.
\item \emph{The slope inequality for many maps} (Part~II, \S6.9--6.10; \texttt{Frame\_at}), which the
source has in the place of hypothesis (a) of Proposition~L15, and the Lean statements on which the inputs
(b)--(e) of that proposition are modelled. The inequality (b) is the subject of \S III.4.
\item \emph{The properties of $d$}: every odd prime up to $x^{\varepsilon/(5C_o)}$ splits
(\texttt{split\_small}), with $C_o$ the Lean development's \texttt{J}, for which the Lean has $1\le\mathtt J$.
\item \emph{The translates.} How the records for the different translates $u$ enter the mean over $u$ in
(M) was not examined.
\end{enumerate}

\textbf{What this section proves, and what it assumes.} \emph{Proved outright:} Lemmas L1, L2, L3 and L5
(lattices); Lemma~L6 (the prime sum as a sum over single primes); Lemmas L7, L8, L9 (losses);
Proposition~L10 (an inequality of the shape of \eqref{dv:link} at the \texttt{signal} primes, for the
formulas of the source, with explicit loss terms; it is \eqref{dv:link} under \eqref{lt:consts});
Proposition~L11, Lemmas L12 and L13, and Lemma~L14 given Theorem~\ref{thm:T} and one cited estimate.
For the prime sum this amounts to the following: the section splits the sum and bounds several of its
terms. The estimates for the eleven loss sums and their total bound \eqref{lt:loss} are quoted;
Theorem~\ref{thm:T} alone does not give that bound.
\emph{Proved from statements of the source:} Theorem~L4: part (i) from \textup{(Val)}, part (iii) from
\textup{(Lat)}; part (ii) needs neither. \emph{A model:} Proposition~L15, whose hypothesis (a) does not
hold in the source. \emph{Assumed:} the items of \S7. The section does not prove that any theta kernel is
divisible by anything: that is \textup{(Val)}, which rests on Theorem~J and on hypothesis (B). It does not
prove the prime side of (M).

\textbf{What remains.}
\begin{enumerate}[label=\textup{(\roman*)},leftmargin=2.4em]
\item The passage from Theorem~J to \texttt{Conclusion1}: frames and labels, the dictionary of \S4, the
chain of $m$ maps, and the local lattice that defines $\Lambda^{\rm div}_\alpha$. Then the passage from
\texttt{adm2} to \textup{(Val)}: the hypotheses \texttt{Tied} and \texttt{CutUp} of \texttt{val2AllAtRF}
and the field \texttt{l1b}, that is, the comparison of the local pairings at different twists.
\item The inequality \eqref{dv:link} on the served orbits of the \texttt{inertClass} primes, which are
good orbits in the sense of (G1)--(G3), with the terms $\mathrm{pr}(p)$, $\mathrm{or}(p)$ and $\tau$; the
bound for the mean of $\kappa(\mathcal Q)$; and the comparison of constants \eqref{lt:consts} at the
\texttt{signal} primes.
\item The lattice statement \textup{(Lat)}, and the statements of \S2\,(a)--(d) that lead from
\eqref{lt:slope} to \eqref{lt:Lo}; these are four items, which contain at least eight statements of the
source.
\item The eleven sums over primes behind \eqref{lt:loss}, and their total bound.
\item The mean-slope argument for all members of a record, which the source has in the place of
hypothesis (a) of Proposition~L15, and the mean over the translates in (M): these are now written, in
\S\ref{as:sec} (Proposition~M4, Proposition~M7, Corollary~M9), as a deduction of (M) from the Hypotheses
(H0)--(H11) of that section. The Hypotheses are not proved; items (i)--(iv) above are among them or behind
them. The upper bound for the largest slope of $\Lambda^{\rm gl}_\alpha$ that accompanies \eqref{lt:Lo}
is proved in \S\ref{as:sec}, Lemma~M10, from the inclusion in \textup{(Lat)}, from \S2\,(b), and from one
further statement of the source, all three assumed.
\end{enumerate}


\section{The size bound: the arithmetic inputs}\label{sa:sec}

This section concerns the passage from the analytic estimates of \S III.4 to the inequality
(\ref{III:sizefinal}) at the end of \S III.4, which is the first line of \eqref{II:t6.5} of Part II. In the source this
passage is one theorem (\texttt{divClausesR\_of}) with a long list of hypotheses.

\emph{Proved here.} First: a finite list of numerical inputs, named (S1)--(S8), (W1), (W2) and (G) below,
implies the inequality, with explicit constants (Theorem~\ref{sa:thm}). The proof includes a lemma on the
norm of a map that factors through a quotient, the averaging of the bounds over the archimedean places,
and the cancellation of the period. Second: a formula for a mean of logarithms of periods, defined by the
integral lines of the field $E$ (the Proposition of subsection 4). It is deduced from the class mean
formula for the elliptic curve (inputs (I2), (I3) of \S III.4), from four statements about ideals, three
of which are proved and one cited, and from one fact about the field $E$. Third: (W1) and (W2) follow from
the conclusions of the paragraph ``Row by row'' of \S III.4 and three facts taken from the source (the
remark after the list of inputs).

\emph{Assumed here, from the source.} All the inputs (S1)--(S8), (W1), (W2), (G). For each we give the
statement and say what kind of theorem it is. In particular the period input (S4) is assumed. The
Proposition of subsection 4 does not prove it. (S4) follows from that Proposition only together with
three further statements, (P1)--(P3), listed at the end of subsection 4. (P1) and (P3) are taken from
the source. (P2) is proved in subsection 4. Nothing in this section proves that the lattices and maps of the construction exist or have
the stated properties.

\subsection*{1.\quad Notation}

\emph{Letters introduced in this section:} $k_d$, $\mathcal P$, $\mathcal J$, $U$, $U_0$, $\bar U$, $Y$,
$\iota$, $\Pi$, $\Phi$, $\varkappa_v$, $\lambda_{j,p}(v)$, $\mathcal L$, $\mathcal R$, $\mathcal R_0$,
$\mu(\cdot)$, $\mu_{\max}(\cdot)$, $\mu_0$, $\mathcal H$, $\nu_{j\tau i}$, $H_{j\tau}$, $\Omega_\tau(v)$, $h_\Pi$, $\Sigma_j(v)$, $r_\nu$, $n_{j\tau}$, $p_j$, $w_{j\tau}$,
$\mathfrak e_{j\tau}$, $\mathfrak d_{j\tau}$, $S_{\rm c}$, $\lambda_0$, $C_0$, $c_e$, $C_h$,
$\mathfrak F$, $E_{\rm loc}$, $E_3$, $l_{3,j}$, $l_{6,j}$, $\mathrm{de}$, $\mathrm{la}$, $H_{\max}$, $H_*$,
$Q_0$, $\mathcal U_1$, $\mathcal U_2$, $\mathcal V$, $\mathcal E_1,\dots,\mathcal E_8$, $E_b$,
$L_{\rm pass}$, $\mathcal I$, $\mathfrak g_j$, $v_{\rm h}$, $\beta^{*}_j(v)$, $\Theta_1(v)$,
$\mathrm{cl}_j$, $\mathrm{cl}_j'$, $L_g$, $\lambda_{1,\tau}$, $b_\tau$, $\mathfrak m_\tau$,
$\mathfrak m_\tau'$, $e_\Pi$, the constants $C_2,C_{\rm th},C_3,C_4,c_5,c_6,C_T,C_8,C_e,C_\Pi,C_H$, and,
in subsection 4, $\mathfrak D$, $\mathfrak b$, $\mathfrak a_0$, $\mathfrak l_\tau$, $\mathcal M$,
$\mathcal M_{\mathfrak c}$, $\mathcal M^{*}$, $\mathfrak c^{*}$. Bold letters
($\mathbf u$, $\mathbf x$, $\mathbf y$, $\mathbf e_\tau$, $\mathbf a_i$, $\mathbf m$) denote vectors.
The lists contain the letters that are used in more than one subsection. A letter that is used in one
proof, in one remark, or in one paragraph of subsection 5 only is introduced where it is used and is not
listed. Examples: $\Xi$, $\mathfrak s_p$ in Step 1; $\Delta$, $B$, $\varsigma$, $C_1$ in remark (1);
$\mathfrak Q$, $z$, $\mathfrak A$ in subsection 4; $w_E$, $D_0$, $\mathrm{wt}(y)$, $\mathrm{ram}(y)$ in
subsection 5.

\emph{Letters that keep the meaning they have in \S III.4:} $x$, $\varepsilon$, $b$, $a$, $m$, $k$,
$f=2n$, $\eta$, $T$, $a_t$, $\mathcal T$, $\Omega_E$, $\mathrm{rd}_E$, $L_a$, $B_0$, $N_c$, $L_{\rm sc}$,
$h(\cdot)$, $H(\cdot)$, $h_A$, $h_{B,j}$, $h_{X,j}$, $g_0$, $g_1$, $r=k+m-1$, $\varrho$, $L$, $M$,
$A_\rho$, $\Lambda$, $\gamma$, $c_d$, the terms $\mu_{\rm CC}$, $F_0$ and the constant $C$ of
(\ref{III:sizefinal}), and, in (S2), $\ell$, $d_0$, $d_1$, $\lambda_1$ (each with the
index $\tau$ of the place). \emph{Letters taken from \S III.5:} $N_m$, $t_0$, $w''$, $b'$.

\emph{Letters that need a remark.}
\begin{itemize}[leftmargin=1.6em,itemsep=1pt,topsep=3pt]
\item $K$ is the number field over which the matrices of the maps are defined, as in the paragraph
``Sums over classes'' of \S III.4. In the paragraph ``The reference mean square'' \S III.4 writes $K$
also for $\Q(\sqrt{-d})$; here that field is $k_d$.
\item $K_j$ $(1\le j\le m)$ are the maps of the chain. Steps are numbered from $1$, as in ``Row by row''
and in the source; ``Sums over classes'' numbers them from $0$.
\item $T$ is the set of distinguished places and $\mathcal T$ is the level, as in (\ref{III:sizefinal}). The
other uses of $T$ in \S III.4 (for $|\theta(0)|$ and for the theta lift) do not occur here. In \S III.5
the letter $\mathcal T$ is an index set.
\item $\tau$ always denotes a weight place, that is, an element of $\mathcal P$; $t$ denotes an element
of $T$; and $\rho$ denotes an embedding of the field $M$ into $\C$, as in ``The reference mean square''.
The letter $\rho$ occurs only in subsection 4 and in quotations from \S III.4. \S III.4 writes $\tau$ for the weight place in ``Several
blocks'' and $\rho$ in ``The eigenvalues of a lift'', ``Row by row'' and ``The pairing at the end of the
chain''; what is $A_\rho^{\pm}$, $J_\rho$, $d_\rho$ there is $A_\tau^{\pm}$, $J_\tau$, $d_\tau$ here.
\item $h$ is the function $h(u)=u\log u-u$, and, in the last display of Theorem~\ref{sa:thm} only, the
term $h$ of (\ref{III:sizefinal}); \S III.4 does the same. Heights of maps are written $\mathrm{ht}$. The numbers
$h_{B,j}$ are not the $h_B$ of Part II, \S6.7.
\item $N_m$ is the rank of the last space of the chain; in the source $N_m=k+m(t_0+1)$. The letters $N$,
$N'$ are left to \S III.4, where they are the ranks of a pair. $t_0$ is the parameter called $t$ in
\S III.4. $\varrho$ and $L$ are those of ``The eigenvalues of a lift'' (Lean: \texttt{nV}, \texttt{Lp}).
\item $\eta$ is the bound for the additional rows, as in input (e) of ``Row by row''; it is not
Dedekind's function. $\pi$ is the number. $\varkappa_v$ is not the ratio $\kappa$ of \S III.4.
$\Theta_1(v)$ is not the $\Theta_0$ of \S III.4; see (S1) and subsection 5.
\item The labels (W1), (W2) are not the rules (W1)--(W4) of \S III.5, and (G) is a label, not a letter.
Parts (a), (b) of Theorem~\ref{sa:thm} are not the inputs (a)--(e) of ``Row by row''.
\item \S III.5 uses $x$, $a$, $\varepsilon$, $\lambda$, $p$, $d$, $\beta$ and $N$ with meanings of its
own, as it says. Here they follow \S III.4.
\item Some letters have several roles inside this section. The index tells them apart.
\begin{itemize}[leftmargin=1.4em,itemsep=0pt,topsep=1pt]
\item $\lambda$: $\lambda_0$ is the bound for $\mathrm{la}$ in (S7). $\lambda_{1,\tau}$ is the first weight
at $\tau$ in (S2), the $\lambda_1$ of \S III.4; it is not related to $\lambda_0$. $\lambda_{j,p}(v)$ are the
eigenvalue factors of (S3). $\lambda$ without index is the scalar of the source's evaluation, in
subsection 5 only.
\item $p$: $p_j$ is a number. $p$ without index is the index of the subspaces $\bar U_p$, in (S3) and in
Step 1 only.
\item $l$: $\ell_\tau$ is the exponent of the level in (S2). $l_{3,j}$ and $l_{6,j}$ are the real numbers
of (S5), (S6) and of (S3). $\mathfrak l_\tau$ is the integral line of subsection 4.
\item $H$: $H(\cdot)$ is the hook product of a diagram, and $\log H_{j\tau}$ is a number of the record.
$H_{\max}$ and $H_*$ are bounds for slopes in (S7); they are not hook products. $H_\rho$ occurs once, in
a quotation from \S III.4 in subsection 5.
\item $E$: $E$ is the field. $E_3$, $E_b$, $E_K$, $E_X$, $E_{\rm loc}$ and
$\mathcal E_1,\dots,\mathcal E_8$ are error terms. $E_0$, in subsection 5 only, is the Euler product of
\S III.4.
\item $d$: $d$ is the prime. $d_{0,\tau}$, $d_{1,\tau}$ are the degrees in (S2). $d_\tau$, in subsection 5
only, is the $d_\rho$ of \S III.4. $\mathfrak d_{j\tau}$ is the error term of (W2). $\mathrm{de}$ is one
symbol.
\item $w$: $w_{j\tau}$ is the weight shift, written $w$ in remark (1) for fixed $j,\tau$. $w''$ is a
parameter. $w_E$, in subsection 5 only, is the number of roots of unity of $E$.
\item $S_{\rm c}$ is a parameter. In $D_{S^c}$ (subsection 5) the index $S^c$ is the complement of a
finite set $S$ of primes.
\item $\varphi$ is an element of $V_0$; $\phi$ is the finite test function of subsection 5. $\rho$ is an
embedding; $\varrho$ is a rank. $r=k+m-1$; $r_\nu$ is the number of kernel degrees. $B_0$ is the scale
above; $B$ occurs in remark (1) only. $\Sigma_j(v)$ is a number, not a sum sign.
\end{itemize}
\end{itemize}

\emph{Letters and the other sections.} The letters introduced here are local to this section. The
sections of Part III do not share one notation, and we did not rename across sections. Each section
declares its own letters. The clashes most likely to confuse are these. $C_0$ and $C_H$ are other
constants in \S III.8. $\mathcal R$ is another object in \S III.6 and in \S III.9. $\mathcal M$,
$\mathcal H$, $\mathcal I$, $\mathfrak m$, $\varkappa$ and $\Theta_1$ have other meanings in \S III.8.
$U$ is a unitary matrix in \S III.4 and another space in \S III.6. $Y$ has other meanings in
\S III.4--III.6, III.8 and III.9. $\Phi$ is a test function in \S III.8. $\iota$ is an index set in
\S III.5 and an isomorphism in \S III.9. $\Pi$ and $\mathfrak d$ have other meanings in \S III.9. $C_2$
is another constant in \S III.6, and $\mu_0$ another number in \S III.9. No statement of this section
uses a letter in the sense of another section, except the letters listed above as kept from \S III.4 and
\S III.5.

\emph{Scales.} $x>1$ and $\varepsilon>0$ are real numbers; $L_{\rm sc}=x^{\varepsilon}$,
$b=e^{L_{\rm sc}}$, $a=x\,b$, $L_a=\log(a\,\mathrm{rd}_E)$, $B_0=f\,m\,b\,x^{\varepsilon}$, and $N_c\ge0$.
(In the application $f=2n$ and $N_c=|T_2|\,k\,b\,x^{2\varepsilon}/n^2$; the proofs below use only
$N_c\ge0$.) As in \S III.4, $h(u)=u\log u-u$ for $u>0$ and $h(0)=0$. So
$a_t\log a_t-a_t$ means $0$ when $a_t=0$ (the convention $0\cdot\log0=0$).

\emph{Index sets.} $\mathcal J=\{1,\dots,m\}$ indexes the lifts. $\mathcal P$ is the set of weight
places; it has at most $f$ elements. $T\subseteq\mathcal P$ is the set of distinguished places, with a
real number $a_t\ge0$ for each $t\in T$. When an element of $T$ is called $\tau$ this number is written
$a_\tau$. For a statement $\mathsf S$ the symbol $[\mathsf S]$ is $1$ if $\mathsf S$ holds and $0$
otherwise.

\emph{The field and its places.} $K$ is the number field of the first remark above. In the source it is
totally complex. $v$ runs over its archimedean places, $s$ is their number, and
$\langle u\rangle=s^{-1}\sum_vu(v)$ for a real function $u$ of $v$.

\emph{The chain and its end.} A Hermitian norm at $v$ on a $K$-vector space $W$ means a positive definite
Hermitian norm on the complex space $W\otimes_{K,v}\C$. $V_0,\dots,V_m$ are finite-dimensional $K$-vector spaces with a Hermitian
norm at each $v$, and $K_j:V_{j-1}\to V_j$ $(j\in\mathcal J)$ are $K$-linear maps. $U$ is a
finite-dimensional $K$-vector space with a positive definite Hermitian norm at each $v$, and
$U_0\subseteq U$ is a $K$-subspace. $\bar U=U/U_0$, and $\bar{\mathbf u}\in\bar U$ is the class of
$\mathbf u\in U$. $\bar U$ carries at each $v$ the quotient norm,
$\|\mathbf x\|=\min\{\|\mathbf u\|:\bar{\mathbf u}=\mathbf x\}$, where $\mathbf x$ is in $\bar U\otimes_{K,v}\C$ and $\mathbf u$ runs over $U\otimes_{K,v}\C$ (the minimum is attained there). $Y$ is a $K$-vector space with a
Hermitian norm at each $v$, and $\iota:\bar U\to Y$ is a $K$-linear bijection. $\Pi:V_m\to U^{*}$ is a
$K$-linear map into the dual space (the pairing at the end of the chain). Dual spaces carry the dual
norms. The map of the inequality is the $K$-linear map $\Phi:V_0\to Y^{*}$ determined by
\begin{equation}\label{sa:fact}
 \Phi(\varphi)\big(\iota\,\bar{\mathbf u}\big)=\Pi\big(K_m\cdots K_1\varphi\big)(\mathbf u)
 \qquad(\varphi\in V_0,\ \mathbf u\in U).
\end{equation}
It exists as soon as the right side vanishes for $\mathbf u\in U_0$; this is part of input (S3). If
$\Phi\ne0$, its height is $\mathrm{ht}(\Phi)=\langle\log\|\Phi\|_v^2\rangle$; \S III.4 writes $h(P)$ for
the height of a map $P$. In the source the chain is a chain of matrix steps, and one step may contain
several lifts (Lean: a map \texttt{lf} from the lifts to the steps); the bound (S1) is then a bound for
the matrix of a step. We state everything for one lift in each step. The argument is the same.

\emph{Lattices and slopes.} A lattice $\mathcal L$ is a finitely generated $\mathcal O_K$-module that
spans a $K$-vector space with norms at every $v$. $\mu(\mathcal L)$ denotes the normalised slope of the
source (\texttt{muDef}): twice the slope of $\mathcal L$ for these norms, plus a constant depending only
on $K$. It is the normalisation for squared norms. $\mathcal L^\vee$ is the dual lattice in the dual
space, with the dual norms, and $\mu_{\max}(\mathcal L)$ is the largest normalised slope of a sublattice.
$\mathcal R\subset Y$ is the lattice whose dual receives the values of $\Phi$, and
$\mu_{\rm CC}=\mu(\mathcal R^\vee)$. There is a second lattice $\mathcal R_0$, in a space $Y_0$ of the
same kind, built from the same data at a lower level; $\mu_0=\mu(\mathcal R_0^\vee)$.

\emph{The record.} The following real numbers are given: $\log\mathcal T\ge0$; $\log\Omega_E$; for
$j\in\mathcal J$ and $\tau\in\mathcal P$ nonnegative kernel degrees $\nu_{j\tau i}$ $(1\le i\le r_\nu)$
with sum $n_{j\tau}=\sum_i\nu_{j\tau i}$, a number $\log H_{j\tau}$ (in the application the logarithm of
the hook product $H(A^+_\tau)$ of step $j$), a number $w_{j\tau}$ (the weight shift), and two error terms
$\mathfrak e_{j\tau}$, $\mathfrak d_{j\tau}$; for $j\in\mathcal J$ a number $p_j$; and parameters
$k,N_m,S_{\rm c},t_0,\eta,w'',b',\lambda_0,\varrho,L$ and constants $C_0,c_e,C_h$. (In the source the
kernel degrees are padded by zeros to $r_\nu=N_m$.) Put
\begin{align}\label{D:253}
 h_A&=-\sum_{t\in T}a_t\log\mathcal T+C_h\big(|T|\,b'N_m+f(m\eta+S_{\rm c})N_m+fN_m^2\big)L_a,\\
 h_{B,j}&=\sum_{\tau\in\mathcal P}\Big(\sum_ih(\nu_{j\tau i})-n_{j\tau}\log\Omega_E\Big)
           +C_hfN_m^2\big(L_a+m^3L_{\rm sc}\big),\\
 h_{X,j}&=\sum_{\tau\in\mathcal P}\log H_{j\tau}-C_hf\,\varrho^2L,\\
 \mathcal H&=h_A+\sum_{j\in\mathcal J}\big(h_{B,j}-h_{X,j}\big).
\end{align}
These are the quantities $h_A$, $h_{B,j}$, $h_{X,j}$ of the paragraph ``Relation to the Lean proof'' of
\S III.4 (Lean: \texttt{RecNums.hA}, \texttt{hB}, \texttt{hX}, \texttt{hu}).

\subsection*{2.\quad The inputs}

\begin{enumerate}[label=\textup{(S\arabic*)},leftmargin=2.8em,itemsep=3pt,topsep=3pt]
\item \emph{The chain bound.} For $j\in\mathcal J$ and every $v$, $\log\|K_j\|_v^2\le\Sigma_j(v)$, where
\begin{align}\label{D:254}
 \Sigma_j(v)={}&\sum_{\tau\in\mathcal P}\Big(\sum_ih(\nu_{j\tau i})-n_{j\tau}\log\Omega_\tau(v)\Big)
   +|\mathcal P|\,\mathfrak g_j\log\Theta_1(v)\\
  &+|\mathcal P|\,\mathfrak g_j\log\big(8+32\,|\mathcal P|\,\mathfrak g_j\,v_{\rm h}\big)
   +2\log\beta^{*}_j(v)\\
  &+\sum_{\tau\in\mathcal P}\big(1+2\mathfrak g_j\log(2n_{j\tau}+2)\big)+\log\mathrm{cl}_j+\log\mathrm{cl}_j' .
\end{align}
Here $\Omega_\tau(v)>0$ (the local squared period), $\Theta_1(v)>0$ (the mean square of the reference
theta function, per variable), $\beta^{*}_j(v)>0$ (the constant of the kernel),
$\mathrm{cl}_j,\mathrm{cl}_j'>0$ (numbers of classes), $\mathfrak g_j$ and $v_{\rm h}$ are real numbers, $v_{\rm h}\ge0$,
with, for constants $C_2,C_{\rm th},C_4,c_5,c_6,C_T,C_8,C_e\ge0$ and a real number $L_g$ (in the
application $\log(g_0g_1)$):
\begin{gather}\label{D:255}
 \log\beta^{*}_j(v)\le C_2fm^2N_m^2(L_g+\log\mathcal T)+C_efN_m^2L_g,\qquad
 \langle\log\Theta_1\rangle\le C_{\rm th}L_{\rm sc},\\
 \log\mathrm{cl}_j,\ \log\mathrm{cl}_j'\le C_4fN_m^2(L_{\rm sc}+\log\mathcal T+L_g),\qquad
 0\le\mathfrak g_j\le N_m^2,\\
 \log\big(8+32|\mathcal P|\mathfrak g_jv_{\rm h}\big)\le c_5L_a,\qquad\log(2n_{j\tau}+2)\le c_6L_a,\\
 L_g\le C_8L_{\rm sc},\qquad\log\mathcal T\le C_TmL_{\rm sc},
\end{gather}
and $C_{\rm th}+c_5+2c_6+1+2C_4+2C_eC_8\le C_h$, $2(C_2+C_4)(C_T+C_8)\le C_h$. Since
$\mathfrak g_j\ge0$ and $v_{\rm h}\ge0$, the number $8+32|\mathcal P|\mathfrak g_jv_{\rm h}$ is at least
$8$; so its logarithm is defined. The condition $v_{\rm h}\ge0$ is assumed here for this reason only. The
record of the source does not contain it; in the source's construction $v_{\rm h}=\mathrm{rd}_E/2>0$
(see subsection 5).
\item \emph{The pairing at the end.} There is a real number $h_\Pi$ with $\|\Pi\|_v^2\le e^{h_\Pi}$ for
every $v$, and there are real numbers $\ell_\tau,\lambda_{1,\tau},d_{0,\tau},d_{1,\tau},b_\tau$,
$\mathfrak m_\tau,\mathfrak m_\tau'$ $(\tau\in\mathcal P)$ and $e_\Pi$, $C_\Pi\ge0$ with
\begin{gather}\label{D:256}
 h_\Pi\le-\sum_{\tau\in\mathcal P}\ell_\tau\log\mathcal T
   +\sum_{\tau\in\mathcal P}\big(2\,\mathfrak m_\tau+\mathfrak m_\tau'\big)+e_\Pi,\qquad
 \ell_\tau=\lambda_{1,\tau}-d_{0,\tau},\\
 \lambda_{1,\tau}=a_\tau-\tfrac{N_m-k}2\ (\tau\in T),\qquad
 \lambda_{1,\tau}\ge S_{\rm c}-\tfrac{N_m-1}2\ (\tau\notin T),\\
 d_{0,\tau}\le S_{\rm c}+m\eta+[\tau\in T]\,b_\tau,\quad b_\tau\le2b'\ (\tau\in T),\quad
 d_{1,\tau}\le3d_{0,\tau}+S_{\rm c}+\eta,\\
 \mathfrak m_\tau\le(N_m-1)\,d_{1,\tau}\log(2a),\quad\mathfrak m_\tau'\le\tfrac{N_m(N_m-1)}2\log(2a),\quad
 e_\Pi\le C_\Pi fN_m^2L_a,
\end{gather}
and $2(16+C_T+C_\Pi)\le C_h$. The number $e_\Pi$ is one real number; it does not depend on $v$. We read
the first line as the bound of the paragraph ``The pairing at the end of the chain'' of \S III.4 at the
place $\tau$, with $\psi\le N_m-1$, $\Psi\le N_m$ and $\Gamma\le2a$, and with $\mathfrak m_\tau$,
$\mathfrak m_\tau'$ in the place of $\log|\vartheta|$, $\log\kappa$ there. This reading is ours: in the
source $\mathfrak m_\tau$ and $\mathfrak m_\tau'$ are by definition the two majorants
$(N_m-1)\,d_{1,\tau}\log(2a)$ and $\frac{N_m(N_m-1)}2\log(2a)$ (Lean: \texttt{logc}, \texttt{logxi} in
\texttt{lastSide\_of\_logT\_le}).
\item \emph{The eigenvalues of a lift.} The right side of \eqref{sa:fact} vanishes for
$\mathbf u\in U_0$. There are a finite nonempty index set $\mathcal I$, $K$-subspaces
$\bar U_p\subseteq\bar U$ $(p\in\mathcal I)$ with $\sum_p\bar U_p=\bar U$, and positive numbers
$\varkappa_v$ and $\lambda_{j,p}(v)$ such that at every $v$
\begin{equation}\label{D:257}
 \langle\iota\mathbf x,\iota\mathbf x'\rangle_{Y,v}
 =\varkappa_v\prod_{j\in\mathcal J}\lambda_{j,p}(v)\ \langle\mathbf x,\mathbf x'\rangle_{\bar U,v}
 \qquad(\mathbf x\in\bar U_p,\ \mathbf x'\in\bar U).
\end{equation}
There are real numbers $l_{6,j}$ with $l_{6,j}\le\log\lambda_{j,p}(v)$ for all $j,p,v$, and
$\sum_jh_{X,j}\le\sum_jl_{6,j}$.
\item \emph{The periods.} For every $\tau\in\mathcal P$, $\langle\log\Omega_\tau\rangle=\log\Omega_E$.
\item \emph{The slope at the lower level.} There are real numbers $l_{3,j}$ $(j\in\mathcal J)$,
$\mathfrak F$ and $E_{\rm loc}$ with
\begin{equation}\label{D:258}
 \mu_0-\langle\log\varkappa\rangle\le\sum_{j\in\mathcal J}l_{3,j}+2\mathfrak F+E_{\rm loc}.
\end{equation}
Here $l_{3,j}$ is an upper bound for the normalised logarithmic covolume, divided by $\dim Y$, of the
lattice attached to the $j$th lift; $\mathfrak F$ is the sum of the local index terms; $E_{\rm loc}$ is
the sum of the local errors. This is the form of the source: its theorem has the bound by $l_{3,j}$ among
its hypotheses and the display as its conclusion (subsection 5).
\item \emph{The covolumes.} There is a real number $E_3$ with
\begin{equation}\label{D:259}
 \sum_{j\in\mathcal J}l_{3,j}\le\sum_{j\in\mathcal J}\sum_{\tau\in\mathcal P}p_j\,w_{j\tau}\log\Omega_E+E_3 .
\end{equation}
\item \emph{The passage between levels.} There are real numbers $\mathrm{de}\ge0$, $\mathrm{la}$, $H_{\max}$, $Q_0\ge1$ with
$\mu_{\max}(\mathcal R_0)\le H_{\max}-\langle\log\varkappa\rangle$, and such that, if
$\mathrm{de}\le\frac13$ and $H_*$ is a real number with $\mu_{\max}(\mathcal R_0)\le H_*$, then
\begin{equation}\label{D:260}
 \mu_{\rm CC}\le\mu_0+3\,\mathrm{de}\,\big(\mu_0+H_*+28\,\mathrm{la}\big)+28\,\mathrm{la} .
\end{equation}
Moreover $3\,\mathrm{de}\,Q_0\le1$,\ \
$\sum_jl_{3,j}+2\mathfrak F+E_{\rm loc}+H_{\max}+28\,\mathrm{la}\le Q_0$,\ \
$\mathrm{la}\le\lambda_0$ and $C_0\ge28$. Here $\langle\log\varkappa\rangle$ is the mean of the numbers
$\log\varkappa_v$ of (S3), as in (S5). (The source uses the same scalars $\varkappa_v$ for $\mathcal R_0$
and for $\mathcal R$; this explains why one number occurs at both levels. It is a comment, not a
hypothesis.)
\item \emph{The size of the local errors.} $E_3+E_{\rm loc}\le c_e\,m\,x^{\varepsilon}(B_0+N_c)$.
\end{enumerate}
\begin{enumerate}[label=\textup{(W\arabic*)},leftmargin=2.8em,itemsep=3pt,topsep=3pt]
\item \emph{Rows against hooks.} For all $j,\tau$:\quad
$\sum_ih(\nu_{j\tau i})-\log H_{j\tau}\le[\tau\in T]\,h(a_\tau)+\mathfrak e_{j\tau}$, with
$\mathfrak e_{j\tau}\le C_0\big(w_{j\tau}+S_{\rm c}+N_m^2+m(t_0+1)\eta\big)L_a$.
\item \emph{The total degree.} For all $j,\tau$:\quad
$n_{j\tau}=[\tau\in T]\,a_\tau+k\,w_{j\tau}+\mathfrak d_{j\tau}$, with
$|\mathfrak d_{j\tau}|\le C_0\big(w_{j\tau}+k(N_m+S_{\rm c})+m(t_0+1)\eta+[\tau\in T]\,k\,b'\big)$.
\end{enumerate}
\begin{enumerate}[label=\textup{(G)},leftmargin=2.8em,itemsep=3pt,topsep=3pt]
\item \emph{Sizes of the parameters.} $\mathrm{rd}_E\ge1$, $\log\mathrm{rd}_E\le L_{\rm sc}$,
$\log x\le L_{\rm sc}$; $m,f,k\ge1$, $k+m\le N_m$; $S_{\rm c},t_0,\eta,b',\lambda_0,\varrho,L\ge0$,
$\varrho\le N_m$; $0\le w_{j\tau}\le2b$; $|p_j-k|\le m$; $|\log\Omega_E|\le L_a$; and
\begin{gather}\label{D:261}
 m\lambda_0\le b,\quad N_m^2\le b,\quad m^3N_m^2L_{\rm sc}\le b,\quad m\,\eta\,N_m\,L_{\rm sc}\le b,\quad
 kN_m\le b,\\
 kS_{\rm c}\le w'',\quad m(t_0+1)\eta\le w'',\quad N_mw''\le b,\quad b'N_m\le a,\quad m\,k\,b'\le a,\\
 1\le C_0\le c_e,\qquad 0\le C_h\le C_0,\qquad C_hL\le C_0L_a .
\end{gather}
\end{enumerate}

\emph{Remarks on the list.} (1) \emph{\textup{(W1)} and \textup{(W2)} follow from ``Row by row''.} In
the source $\mathfrak e_{j\tau}$ and $\mathfrak d_{j\tau}$ are defined by equality in the first halves of
(W1) and (W2) (Lean: \texttt{eEN}, \texttt{dN}). Fix $j$ and $\tau$ and write $w=w_{j\tau}$. From its
inputs (o), (a)--(e) the paragraph ``Row by row'' of \S III.4 derives
\begin{equation}\label{D:262}
 \mathfrak e_{j\tau}\le(w+\Delta+B)\,L_a,\qquad|\mathfrak d_{j\tau}|\le k(\varsigma+\Delta)+w+B,
 \qquad\varsigma=\frac{k+(2j-1)(t_0+1)}2 ,
\end{equation}
where $\Delta\ge0$ is the width of the weight interval and $B$ bounds the extra degrees. (At a
distinguished place it gives $(w+B)L_a$ for the first.) Three facts are taken from the source (Lean:
the fields \texttt{constant\_weight\_sum\_bound}, \texttt{constant\_succ\_le}, \texttt{blk} of
\texttt{NumFacts}, and \texttt{Nm\_cast}): $2\Delta\le C_1S_{\rm c}$ for a constant $C_1$ with
$C_1+1\le C_0$; $B=m(t_0+1)\eta$; and $N_m=k+m(t_0+1)$. So $\Delta\le C_0S_{\rm c}$, because
$S_{\rm c}\ge0$, and $\varsigma\le N_m$, because $2j-1\le2m$. With $w\ge0$, $B\ge0$, $k\ge1$ and
$C_0\ge1$ from (G),
\begin{align}\label{D:263}
 w+\Delta+B&\le w+C_0S_{\rm c}+B\le C_0\big(w+S_{\rm c}+N_m^2+m(t_0+1)\eta\big),\\
 k(\varsigma+\Delta)+w+B&\le kN_m+C_0kS_{\rm c}+w+B\le C_0\big(w+k(N_m+S_{\rm c})+m(t_0+1)\eta\big).
\end{align}
Multiply the first line by $L_a\ge0$. These are the bounds of (W1) and (W2); the term
$[\tau\in T]\,k\,b'$ of (W2) is not needed. So (W1) and (W2) hold under the inputs (o), (a)--(e) of
``Row by row'' and the three facts. In the source this is \texttt{stepE\_num}, used in
\texttt{RecNums.ofNum}. The theorem below assumes (W1) and (W2) as stated.

(2) In the source $p_j=k+(2j-1-m)$ (Lean: \texttt{ClosedScale3.l3m}), so that $|p_j-k|\le m-1$.

(3) The condition $C_0\ge28$ stands in (S7) because the source has it there (Lean: the field
\texttt{constant\_ge\_twenty\_eight}). The proof uses it once, in Step 3, for
$28\,\mathrm{la}\le C_0\,m\,\lambda_0$. It is not an extra condition: $2(16+C_T+C_\Pi)\le C_h$ of (S2) and
$C_h\le C_0$ of (G) give $C_0\ge32$. Two other conditions are also used once each: $k+m\le N_m$ of (G)
only through $m\le N_m$ in Step 2, and $C_0\le c_e$ only for the last display of the theorem.

\subsection*{3.\quad The assembly}

\begin{theorem}\label{sa:thm}
Assume \textup{(S1)--(S8)}, \textup{(W1)}, \textup{(W2)} and \textup{(G)}. Then:
\begin{enumerate}[label=\textup{(\alph*)},leftmargin=2.4em,itemsep=2pt,topsep=3pt]
\item if $\Phi\ne0$, then $\ \mathrm{ht}(\Phi)\le\mathcal H-\langle\log\varkappa\rangle$;
\item with $\mathcal U_1=|T|\,a\,L_a$ and $\mathcal U_2=m\,x^{\varepsilon}(B_0+N_c)$,
\begin{equation}\label{D:264}
\begin{aligned}
 \mu_{\rm CC}-\langle\log\varkappa\rangle+\mathcal H-2\mathfrak F\ \le{}&
 \sum_{t\in T}\Big(m\big(h(a_t)-a_t\log\Omega_E\big)-a_t\log\mathcal T\Big)\\
 &+2C_0\,\mathcal U_1+\big(47C_0+7+c_e\big)\mathcal U_2 .
\end{aligned}
\end{equation}
\end{enumerate}
In particular, with $h=\mathcal H-\langle\log\varkappa\rangle$, $F_0=2\mathfrak F$ and
$C=\max(50C_0+10,\,60c_e)$,
\begin{equation}\label{D:265}
\begin{aligned}
 \mu_{\rm CC}+h-F_0\le{}&\sum_{t\in T}\Big(m\big(a_t\log a_t-a_t-a_t\log\Omega_E\big)-a_t\log\mathcal T\Big)\\
 &+C\,|T|\,a\,L_a+C\,m\,B_0\,x^{2\varepsilon}+C\,m\,N_c\,x^{2\varepsilon}.
\end{aligned}
\end{equation}
\end{theorem}

The last display is the inequality (\ref{III:sizefinal}) at the end of \S III.4. In it
$a_t\log a_t$ means $0$ when $a_t=0$, since $h(0)=0$; part (b) is written with $h(a_t)$. The constant $C$
is at least $1610$, because $C_0\ge32$. So the inequality has content only when $m$ is large compared
with $C$: otherwise the error term $C\,|T|\,a\,L_a$ can exceed the main term. The source states (b) with the
error $(50C_0+10)\,\mathcal U_1+60c_e\,\mathcal U_2$ (Lean: \texttt{RecNums.rhs}); this is weaker than
(b), since $47C_0+7+c_e\le55c_e$. The terms $h$ and $F_0$ of the final inequality are thus
$\mathcal H-\langle\log\varkappa\rangle$ and $2\mathfrak F$. The term $F_0$ enters in two places. In
(S5) it is added in the bound for the slope, and it is subtracted on the left. It also occurs in the
condition of (S7) that bounds $Q_0$ from below. The proof takes the rest of this subsection. It has four
steps: a bound for $\|\Phi\|_v$ at one place; the average over the places; the slope; the cancellation.

\medskip
\noindent\textbf{Step 1: a map that factors through a quotient.}

\textbf{Lemma (descent).} \emph{Fix a place $v$ and extend all spaces to $\C$ by the embedding at $v$.
Assume the first two sentences of \textup{(S3)}. Then}
\begin{equation}\label{D:266}
 \|\Phi\|_v^2\ \le\ \frac{\|\Pi\|_v^2\,\prod_{j\in\mathcal J}\|K_j\|_v^2}
 {\varkappa_v\,\min_{p\in\mathcal I}\prod_{j\in\mathcal J}\lambda_{j,p}(v)} .
\end{equation}
\emph{Proof.} Write $\mathfrak s_p=\varkappa_v\prod_j\lambda_{j,p}(v)$ and
$\mathfrak s_{\min}=\min_p\mathfrak s_p$. Let $\Xi$ be the operator on $\bar U\otimes\C$ with
$\langle\iota\mathbf x,\iota\mathbf x'\rangle_Y=\langle\Xi\mathbf x,\mathbf x'\rangle_{\bar U}$ for all
$\mathbf x,\mathbf x'$. It is Hermitian. The hypothesis says
$\langle\Xi\mathbf x,\mathbf x'\rangle=\mathfrak s_p\langle\mathbf x,\mathbf x'\rangle$ for
$\mathbf x\in\bar U_p$ and all $\mathbf x'$, that is, $\Xi\mathbf x=\mathfrak s_p\mathbf x$ on $\bar U_p$.
The spaces $\bar U_p\otimes\C$ span $\bar U\otimes\C$, so $\Xi$ is diagonalisable with all eigenvalues
among the $\mathfrak s_p$. Eigenspaces of a Hermitian operator for different eigenvalues are orthogonal.
Hence $\|\iota\mathbf x\|_Y^2=\langle\Xi\mathbf x,\mathbf x\rangle\ge\mathfrak s_{\min}\|\mathbf x\|_{\bar U}^2$
for every $\mathbf x$. Now let $\mathbf y\in Y\otimes\C$ and put $\mathbf x=\iota^{-1}\mathbf y$. Choose
$\mathbf u$ with $\bar{\mathbf u}=\mathbf x$ and $\|\mathbf u\|_U=\|\mathbf x\|_{\bar U}$; this is
possible by the definition of the quotient norm. By \eqref{sa:fact},
\begin{equation}\label{D:267}
 |\Phi(\varphi)(\mathbf y)|=\big|\Pi(K_m\cdots K_1\varphi)(\mathbf u)\big|
 \le\|\Pi\|\,\prod_j\|K_j\|\,\|\varphi\|\,\|\mathbf u\|_U
 \le\|\Pi\|\,\prod_j\|K_j\|\,\|\varphi\|\,\frac{\|\mathbf y\|_Y}{\sqrt{\mathfrak s_{\min}}} .
\end{equation}
The norm of $\Phi(\varphi)$ in the dual of $Y$ is the supremum of
$|\Phi(\varphi)(\mathbf y)|/\|\mathbf y\|_Y$. \qed

\noindent(Lean: \texttt{HeightBoundConstant}, proved as \texttt{height\_constant\_bound} and used in
\texttt{height\_chain\_atR}; there for a chain of matrix steps, as said in subsection 1.)

Assume now (S1), (S2), (S3) and $\Phi\ne0$. Then $\|\Phi\|_v>0$ at every $v$, because a nonzero entry of
the matrix of $\Phi$ stays nonzero under an embedding of $K$. Taking logarithms in the lemma and using
$l_{6,j}\le\log\lambda_{j,p}(v)$,
\begin{equation}\label{sa:place}
 \log\|\Phi\|_v^2\le h_\Pi+\sum_{j\in\mathcal J}\Sigma_j(v)-\log\varkappa_v-\sum_{j\in\mathcal J}l_{6,j}
 \qquad\text{for every }v.
\end{equation}

\noindent\textbf{Step 2: the average over the places.}

\textbf{Lemma (the height record).} \emph{Assume the numerical parts of \textup{(S1)} and
\textup{(S2)}, the identity \textup{(S4)}, and \textup{(G)}. Then}
\begin{equation}\label{D:269}
 h_\Pi\le h_A\qquad\text{and}\qquad\langle\Sigma_j\rangle\le h_{B,j}\quad(j\in\mathcal J).
\end{equation}
\emph{Proof.} From (G): $L_{\rm sc}\ge1$ since $x>1$; $L_a=\log x+L_{\rm sc}+\log\mathrm{rd}_E$, so
$1\le L_{\rm sc}\le L_a\le3L_{\rm sc}$; $0\le\log(2a)\le2L_a$; and $m\le N_m$.

\emph{The first inequality.} For $\tau\in T$ put $\bar D_\tau=S_{\rm c}+m\eta+2b'$, and for
$\tau\notin T$ put $\bar D_\tau=S_{\rm c}+m\eta$; then $d_{0,\tau}\le\bar D_\tau$ and
$d_{1,\tau}\le3\bar D_\tau+S_{\rm c}+\eta\le4S_{\rm c}+4m\eta+[\tau\in T]\,6b'$. Since
$(N_m-1)\log(2a)\ge0$ and the last bound is nonnegative,
\begin{equation}\label{D:270}
 2\,\mathfrak m_\tau\le\big(16N_m(S_{\rm c}+m\eta)+[\tau\in T]\,24\,b'N_m\big)L_a,
 \qquad\mathfrak m_\tau'\le N_m^2L_a .
\end{equation}
Next, $-\ell_\tau\log\mathcal T=(d_{0,\tau}-\lambda_{1,\tau})\log\mathcal T$ and
$0\le\log\mathcal T\le C_TmL_a$. For $\tau\in T$ the coefficient is
$-a_\tau+\frac{N_m-k}2+d_{0,\tau}\le-a_\tau+\frac{N_m}2+S_{\rm c}+m\eta+2b'$; for $\tau\notin T$ it is at
most $\frac{N_m-1}2-S_{\rm c}+S_{\rm c}+m\eta\le\frac{N_m}2+m\eta$. With $m\le N_m$ this gives
\begin{equation}\label{D:271}
 -\ell_\tau\log\mathcal T\le-[\tau\in T]\,a_\tau\log\mathcal T
 +C_T\Big(\tfrac{N_m^2}2+N_m(S_{\rm c}+m\eta)+[\tau\in T]\,2b'N_m\Big)L_a .
\end{equation}
Sum over the at most $f$ places $\tau$, of which $|T|$ are in $T$, and add $e_\Pi\le C_\Pi fN_m^2L_a$:
\begin{equation}\label{D:272}
 h_\Pi\le-\sum_{t\in T}a_t\log\mathcal T+\Big((24+2C_T)|T|\,b'N_m+(16+C_T)f(m\eta+S_{\rm c})N_m
 +\big(1+\tfrac{C_T}2+C_\Pi\big)fN_m^2\Big)L_a .
\end{equation}
Each of the three coefficients is at most $2(16+C_T+C_\Pi)\le C_h$. This is $h_\Pi\le h_A$.

\emph{The second inequality.} Average $\Sigma_j(v)$ over $v$. By (S4) the first sum becomes
$\sum_\tau(\sum_ih(\nu_{j\tau i})-n_{j\tau}\log\Omega_E)$, exactly. For the other terms, using
$|\mathcal P|\le f$, $0\le\mathfrak g_j\le N_m^2$, $1\le L_{\rm sc}\le L_a$ and
$L_g+\log\mathcal T\le(C_8+C_T)mL_{\rm sc}$:
\begin{align}\label{D:273}
 |\mathcal P|\,\mathfrak g_j\,\langle\log\Theta_1\rangle&\le C_{\rm th}fN_m^2L_a,\\
 |\mathcal P|\,\mathfrak g_j\log(8+32|\mathcal P|\mathfrak g_jv_{\rm h})&\le c_5fN_m^2L_a,\\
 \textstyle\sum_\tau\big(1+2\mathfrak g_j\log(2n_{j\tau}+2)\big)&\le(1+2c_6)fN_m^2L_a,\\
 2\langle\log\beta^{*}_j\rangle&\le2C_2(C_T+C_8)fm^3N_m^2L_{\rm sc}+2C_eC_8fN_m^2L_a,\\
 \log\mathrm{cl}_j+\log\mathrm{cl}_j'&\le2C_4fN_m^2L_a+2C_4(C_T+C_8)fm^3N_m^2L_{\rm sc}.
\end{align}
The sum of the coefficients of $fN_m^2L_a$ is $C_{\rm th}+c_5+2c_6+1+2C_4+2C_eC_8\le C_h$, and that of
$fm^3N_m^2L_{\rm sc}$ is $2(C_2+C_4)(C_T+C_8)\le C_h$. \qed

\noindent(Lean: \texttt{height\_bound\_at\_set}, applied through \texttt{HtNums.le\_hu}. A parallel
statement for the first inequality, \texttt{stepDiiiRecord}, is in the archive but is not used by the
final theorem.)

\emph{Proof of Theorem \ref{sa:thm}(a).} Average \eqref{sa:place} over $v$ and apply the lemma and the
last inequality of (S3):
\begin{equation}\label{D:274}
 \mathrm{ht}(\Phi)\le h_\Pi+\sum_j\langle\Sigma_j\rangle-\sum_jl_{6,j}-\langle\log\varkappa\rangle
 \le h_A+\sum_jh_{B,j}-\sum_jh_{X,j}-\langle\log\varkappa\rangle=\mathcal H-\langle\log\varkappa\rangle .
\end{equation}
\emph{This is the first place where $\log\Omega_E$ enters:} through (S4), each $h_{B,j}$ contains
$-\sum_\tau n_{j\tau}\log\Omega_E$. The hook products enter through (S3), as the subtracted terms
$h_{X,j}$. \qed

\medskip
\noindent\textbf{Step 3: the slope.} Put $\mu_0'=\mu_0-\langle\log\varkappa\rangle$ and
$\mu_C=\mu_{\rm CC}-\langle\log\varkappa\rangle$. By (S5),
\begin{equation}\label{sa:mu0a}
 \mu_0'\le\sum_jl_{3,j}+2\mathfrak F+E_{\rm loc},
\end{equation}
and then by (S6) and (S8),
\begin{equation}\label{sa:mu0}
 \mu_0'\le\sum_{j\in\mathcal J}\sum_{\tau\in\mathcal P}p_j\,w_{j\tau}\log\Omega_E+2\mathfrak F+E_b,\qquad
 E_b=E_3+E_{\rm loc}\le c_e\,m\,x^{\varepsilon}(B_0+N_c).
\end{equation}
\emph{This is the second place where $\log\Omega_E$ enters.} For the passage, $3\,\mathrm{de}\,Q_0\le1$
and $Q_0\ge1$ give $\mathrm{de}\le\frac13$. Apply (S7) with $H_*=H_{\max}-\langle\log\varkappa\rangle$ and
subtract the number $\langle\log\varkappa\rangle$ from both sides:
\begin{equation}\label{D:277}
 \mu_C\le\mu_0'+3\,\mathrm{de}\,\big(\mu_0'+H_{\max}+28\,\mathrm{la}\big)+28\,\mathrm{la} .
\end{equation}
By \eqref{sa:mu0a} and (S7), $\mu_0'+H_{\max}+28\,\mathrm{la}\le Q_0$; as $\mathrm{de}\ge0$, the middle
term is at most $3\,\mathrm{de}\,Q_0\le1$. Also $28\,\mathrm{la}\le28\lambda_0\le C_0m\lambda_0$. So
\begin{equation}\label{sa:pass}
 \mu_C\le\mu_0'+L_{\rm pass},\qquad L_{\rm pass}\le1+C_0\,m\,\lambda_0 .
\end{equation}

\noindent\textbf{Step 4: the cancellation.}

\textbf{Lemma (cancellation).} \emph{Assume \textup{(W1)}, \textup{(W2)}, \textup{(G)}, and let
$\mu_C,\mu_0',\mathfrak F,E_b,L_{\rm pass}$ be real numbers satisfying \eqref{sa:mu0} and
\eqref{sa:pass}. Then}
\begin{equation}\label{D:279}
 \mu_C+\mathcal H-2\mathfrak F\le\sum_{t\in T}\Big(m\big(h(a_t)-a_t\log\Omega_E\big)-a_t\log\mathcal T\Big)
 +2C_0\,\mathcal U_1+(47C_0+7+c_e)\,\mathcal U_2 .
\end{equation}
\emph{Proof.} Add \eqref{sa:pass}, \eqref{sa:mu0} and the definition of $\mathcal H$:
\begin{align}\label{D:280}
 \mu_C+\mathcal H-2\mathfrak F\le{}&L_{\rm pass}+E_b+h_A+E_K+E_X\\
 &+\sum_{j,\tau}\Big(p_jw_{j\tau}\log\Omega_E+\sum_ih(\nu_{j\tau i})-n_{j\tau}\log\Omega_E-\log H_{j\tau}\Big),
\end{align}
where $E_K=m\,C_hfN_m^2(L_a+m^3L_{\rm sc})$ and $E_X=m\,C_hf\,\varrho^2L$. By (W1) and (W2) the summand is
at most
\begin{equation}\label{D:281}
 [\tau\in T]\big(h(a_\tau)-a_\tau\log\Omega_E\big)+\mathfrak e_{j\tau}
 +(p_j-k)\,w_{j\tau}\log\Omega_E-\mathfrak d_{j\tau}\log\Omega_E .
\end{equation}
\emph{Here the two occurrences of the period cancel:} $+p_jw_{j\tau}\log\Omega_E$ from the slope against
$-k\,w_{j\tau}\log\Omega_E$ from the height. What is left of them is $(p_j-k)w_{j\tau}\log\Omega_E$. The
sum over $j\in\mathcal J$ and $\tau\in\mathcal P$ of the first bracket is
$m\sum_{t\in T}(h(a_t)-a_t\log\Omega_E)$, and $h_A$ contributes $-\sum_ta_t\log\mathcal T$. So the left
side is at most the main term plus
\begin{equation}\label{D:282}
 \mathcal E_1+\dots+\mathcal E_8,
\end{equation}
where
\begin{gather}\label{D:283}
 \mathcal E_1=L_{\rm pass},\qquad\mathcal E_2=E_b,\qquad\mathcal E_3=E_K,\qquad\mathcal E_4=E_X,\\
 \mathcal E_5=C_h\big(|T|b'N_m+f(m\eta+S_{\rm c})N_m+fN_m^2\big)L_a,\qquad
 \mathcal E_6=\textstyle\sum_{j,\tau}\mathfrak e_{j\tau},\\
 \mathcal E_7=\textstyle\sum_{j,\tau}|p_j-k|\,w_{j\tau}\,|\log\Omega_E|,\qquad
 \mathcal E_8=\textstyle\sum_{j,\tau}|\mathfrak d_{j\tau}|\,|\log\Omega_E| .
\end{gather}
Put $\mathcal V=f\,m\,b\,L_{\rm sc}$. From (G): $L_{\rm sc}\ge1$, $b\ge1$, $L_a\le3L_{\rm sc}$;
$S_{\rm c}\le kS_{\rm c}\le w''\le N_mw''\le b$ and $S_{\rm c}N_m\le b$; $m(t_0+1)\eta\le w''\le b$;
$b\le\mathcal V$; and
\begin{equation}\label{D:284}
 \mathcal U_2\ge m\,x^{\varepsilon}B_0=m\,L_{\rm sc}\,\mathcal V\ \ge\ m\mathcal V\ \ge\ \mathcal V .
\end{equation}
There are $m$ values of $j$ and at most $f$ values of $\tau$. Then:
\begin{align}\label{D:285}
 \mathcal E_1&\le1+C_0b\le(1+C_0)\mathcal V;\\
 \mathcal E_3&\le C_0\,mf\,(b\cdot3L_{\rm sc}+b)\le4C_0\mathcal V;\\
 \mathcal E_4&\le C_0\,mfN_m^2L_a\le3C_0\mathcal V&&(\text{by }C_hL\le C_0L_a,\ \varrho\le N_m);\\
 \mathcal E_5&\le C_0|T|aL_a+C_0f\big(3b+3bL_{\rm sc}+3bL_{\rm sc}\big)\le C_0\,\mathcal U_1+9C_0\mathcal V
   &&(\text{by }b'N_m\le a,\ m\eta N_mL_{\rm sc}\le b);\\
 \mathcal E_6&\le mf\,C_0\,(2b+b+b+b)\,3L_{\rm sc}=15C_0\mathcal V;\\
 \mathcal E_7&\le mf\cdot m\cdot2b\cdot3L_{\rm sc}=6m\mathcal V;\\
 \mathcal E_8&\le15C_0\mathcal V+m|T|\,C_0\,k\,b'L_a\le15C_0\mathcal V+C_0\,\mathcal U_1
   &&(\text{by }kN_m\le b,\ kS_{\rm c}\le b,\ mkb'\le a).
\end{align}
Finally $\mathcal E_2\le c_e\,\mathcal U_2$. The sum is at most
$2C_0\,\mathcal U_1+(47C_0+1)\mathcal V+6m\mathcal V+c_e\,\mathcal U_2
\le2C_0\,\mathcal U_1+(47C_0+7+c_e)\,\mathcal U_2$. \qed

\noindent(Lean: \texttt{divisor\_inequality\_explicit}, used through \texttt{divIneqRecord} and
\texttt{RecNums.divIneq}; it is \texttt{divisor\_inequality\_assembly}, the exact rearrangement, followed
by \texttt{error\_absorption}, the estimates of $\mathcal E_1,\dots,\mathcal E_8$. The source's constants
are $50C_0+10$ and $60c_e$.)

\emph{Proof of Theorem \ref{sa:thm}(b).} Steps 3 and 4. For the last display of the theorem:
$h(a_t)=a_t\log a_t-a_t$; $2C_0\le C$; $47C_0+7+c_e\le55c_e\le C$ because $1\le C_0\le c_e$; and
$\mathcal U_2=m\,x^{\varepsilon}(B_0+N_c)\le m\,x^{2\varepsilon}(B_0+N_c)$ because $x>1$. \qed

\medskip
\noindent\textbf{Where each term comes from.}
\begin{center}\small
\begin{tabular}{p{.27\textwidth}p{.66\textwidth}}
\toprule
Term of the inequality & Source in the proof\\
\midrule
$m\,h(a_t)$ & (W1): at a distinguished place the hook product cancels all rows of the kernel but the
first, once for each of the $m$ lifts\\
$-m\,a_t\log\Omega_E$ & (S4) and (W2): the part $[\tau\in T]\,a_\tau$ of the degree $n_{j\tau}$, times the
mean period, for each lift\\
$-a_t\log\mathcal T$ & (S2): the level factor of the last map, through
$\ell_\tau=\lambda_{1,\tau}-d_{0,\tau}$ and $\lambda_{1,\tau}=a_\tau-\frac{N_m-k}2$\\
cancelled: $k\,w_{j\tau}\log\Omega_E$ & $-k\,w_{j\tau}\log\Omega_E$ from (S4), (W2) in the height;
$+p_jw_{j\tau}\log\Omega_E$ from (S6) in the slope; remainder $\mathcal E_7$\\
$C\,|T|\,a\,L_a$ & $\mathcal E_5$ (the term $|T|\,b'N_m$ of the end pairing) and $\mathcal E_8$ (the term
$k\,b'$ of the degree error at distinguished places, times $|\log\Omega_E|$)\\
$C\,m\,B_0\,x^{2\varepsilon}$ & $\mathcal E_1$ (level passage), $\mathcal E_3$ (constants of the kernels,
class numbers, reference mean square), $\mathcal E_4$ (loss in the eigenvalue bound), the rest of
$\mathcal E_5$, $\mathcal E_6$ (rows against hooks), $\mathcal E_7$ (the periods cancel only up to
$|p_j-k|\le m$), the rest of $\mathcal E_8$; and the part of $\mathcal E_2$ with $B_0$\\
$C\,m\,N_c\,x^{2\varepsilon}$ & $\mathcal E_2$ only: the local errors $E_3+E_{\rm loc}$, by (S8)\\
\bottomrule
\end{tabular}
\end{center}
Among the terms of the second kind, $\mathcal E_7\le6m\mathcal V$ is the only one whose bound carries the
factor $m$; so it is the only one that needs the full factor $m$ in
$\mathcal U_2\ge mL_{\rm sc}\mathcal V$. It is the cost of $p_j\ne k$. Its bound is not the largest in
general: the bounds for $\mathcal E_6$ and $\mathcal E_8$ are $15C_0\mathcal V$, and $6m>15C_0$ only when
$m>2.5\,C_0$.

\subsection*{4.\quad The mean of the periods of the integral lines}

This subsection proves a formula for a mean of periods (the Proposition below). It does not prove (S4);
what is missing is listed at the end. \emph{In this subsection $\rho$ denotes an embedding of a field $M$
into $\C$, as in the paragraph ``The reference mean square'' of \S III.4, and $\tau$, $\tau'$ denote
weight places.}

Let $k_d=\Q(\sqrt{-d})$, embedded in $\C$. Its discriminant is $-d$, because $d$ is squarefree and
$d\equiv3\bmod4$ (the $d$ of the paper is a prime with $d\equiv7\bmod8$); Lemma A uses this. Let
$E\supset k_d$ be the field of the construction, of degree $[E:k_d]=f$. The weight places are the $f$
embeddings $\tau:E\to\C$ that extend the given embedding of $k_d$. So here $\mathcal P$ has exactly $f$
elements; subsection 3 used only $|\mathcal P|\le f$. Let $\mathfrak D$ be the different of $E$ over
$k_d$, and let the bar be complex conjugation on $E$.

\emph{A fact about the field $E$.} There is an integral ideal $\mathfrak b$ of $E$ with
$\mathfrak D=\mathfrak b\,\bar{\mathfrak b}$. This holds for the field of the paper,
$E=\Q(\zeta_N+\zeta_N^{-1},\sqrt{-q},\sqrt{-d})$, where $N=4n$, $n$ is a power of $2$, and $q$ is a prime
with $q\equiv3\bmod4$ that divides $d+1$. The proof has two steps.

\emph{First step: every prime of $E$ that divides $\mathfrak D$ lies over $2$ or over $q$.} (This is
(P2) below.) Put $F=\Q(\zeta_N+\zeta_N^{-1},\sqrt{-q})$. Then $E=k_dF$. Since $q\equiv3\bmod4$, the number
$\sqrt{-q}$ lies in $\Q(\zeta_q)$ (the quadratic Gauss sum). So $F\subseteq\Q(\zeta_{Nq})$, and $F$ is
unramified over $\Q$ at every prime number that does not divide $2q$. Let $\mathfrak Q$ be a prime of $E$
that does not lie over $2$ or $q$. An extension that is unramified at a prime stays unramified when the
ground field is enlarged; so $E=k_dF$ is unramified over $k_d$ at $\mathfrak Q$. A prime that is
unramified does not divide the different. (For $n=1$ the field $F$ is $\Q(\sqrt{-q})$, and only primes
over $q$ can divide $\mathfrak D$.) The Lean has this step as \texttt{diffSupport}.

\emph{Second step.} $2$ and $q$ split in $k_d$, because $d\equiv7\bmod8$ and $-d\equiv1\bmod q$. Complex
conjugation exchanges the two primes of $k_d$ over $2$, and the two over $q$. So a prime $\mathfrak Q$ of
$E$ that divides $\mathfrak D$ and its conjugate $\bar{\mathfrak Q}$ lie over different primes of $k_d$,
and $\bar{\mathfrak Q}\ne\mathfrak Q$. Conjugation maps $E$ to $E$ and $k_d$ to $k_d$; so
$\bar{\mathfrak D}=\mathfrak D$, and $\mathfrak Q$, $\bar{\mathfrak Q}$ occur in $\mathfrak D$ with the
same exponent. Take for $\mathfrak b$ the product of one prime power from each pair. (Lean:
\texttt{principal} in \texttt{KcDiff.lean}, from \texttt{diffSupport} and \texttt{kc\_below\_ne}; it is
used through \texttt{Arith.bb\_ex} to fill the field $\mathfrak b$\texttt{\_spec} of the record
\texttt{Host}.)

This $\mathfrak b$ is coprime to $\bar{\mathfrak b}$, as is the ideal of the source. There are other
integral choices when an exponent $e$ is larger than $1$: one may take $\mathfrak Q^{i}\bar{\mathfrak Q}^{e-i}$
with $0\le i\le e$ from a pair. The Proposition below holds for every $\mathfrak b$ with
$\mathfrak D=\mathfrak b\,\bar{\mathfrak b}$; its proof uses nothing else. The coprimality is not used in
this section.

The splitting of $2$ and $q$ in $k_d$ is sufficient for the existence of $\mathfrak b$. It is not
necessary. For a field $E$ with complex multiplication over an imaginary quadratic field, an integral
$\mathfrak b$ with $\mathfrak D=\mathfrak b\,\bar{\mathfrak b}$ exists if and only if every prime of
$\mathfrak D$ that is equal to its conjugate has an even exponent. (If $\bar{\mathfrak Q}=\mathfrak Q$,
the exponent of $\mathfrak Q$ in $\mathfrak b\,\bar{\mathfrak b}$ is twice its exponent in $\mathfrak b$.
Conversely, halve the exponents of these primes and treat the other primes in pairs as above.) So the
statement fails for some fields: for $E=\Q(\sqrt{-7},\sqrt{-3})$ over $\Q(\sqrt{-7})$ the different is the
prime over $3$, with exponent $1$, and it is its own conjugate.

Fix such a $\mathfrak b$ and put $\mathfrak a_0=\mathfrak b^{-1}$. Let $M\subset\C$ be a number field
that contains $\tau(E)$ for every $\tau$ and a field of definition of the elliptic curve of \S III.4,
with lattice $\Lambda$ normalised by $g_2^3-27g_3^2=1$. For an embedding $\rho$ of $M$, $A_\rho$ is the
area of the lattice of the conjugate curve, as in \S III.4.

Let $\mathcal M\subset M^{f}$ be the $\mathcal O_M$-module spanned by the vectors
$(\tau(\alpha))_\tau$, $\alpha\in\mathfrak a_0$, and let
\begin{equation}\label{D:286}
 \mathfrak l_\tau=\{z\in M:\ z\,\mathbf e_\tau\in\mathcal M\},
\end{equation}
where $\mathbf e_\tau$ is the $\tau$th coordinate vector. This is the integral line at the place $\tau$
(Lean: \texttt{ellT}). For a nonzero $z\in\mathfrak l_\tau$ the local squared period at $(\tau,\rho)$ is
$A_\rho/(\pi|\rho(z)|^2)$, and the mean of its logarithm is defined, with a correction for the index, by
\begin{equation}\label{sa:meandef}
 \overline{\log\Omega_\tau}=\frac1{[M:\Q]}\Big(\sum_{\rho}\log\frac{A_\rho}{\pi\,|\rho(z)|^2}
 +2\log\big[\mathfrak l_\tau:\mathcal O_Mz\big]\Big).
\end{equation}
(Lean: \texttt{meanLogOm}, \texttt{logMet}.) If $\mathfrak l_\tau=\mathcal O_Mz_\tau$ is principal, this
is the plain mean over $\rho$ of $\log\big(A_\rho/(\pi|\rho(z_\tau)|^2)\big)$.

\textbf{Proposition (the mean period).} \emph{The number \eqref{sa:meandef} does not depend on $z$, on
$\tau$, or on $M$, and}
\begin{equation}\label{D:288}
 \overline{\log\Omega_\tau}=\frac1{[M:\Q]}\sum_\rho\log\frac{A_\rho}{\pi}-\log\frac{\mathrm{rd}_E}{\sqrt d}.
\end{equation}
\emph{Under the inputs \textup{(I2)}, \textup{(I3)} of \S III.4, it equals
$\gamma+\log(2\pi/\mathrm{rd}_E)-c_d=\log\Omega_E$.}

The proof uses four statements about ideals. $\mathrm N$ is the absolute norm of ideals of $E$.

\textbf{Lemma A (the tower formula).} \emph{$|\mathrm{disc}\,E|=d^{\,f}\,\mathrm N(\mathfrak D)$.
Consequently $\frac1{[E:\Q]}\log\mathrm N(\mathfrak D)=\log\mathrm{rd}_E-\log\sqrt d$.}

\emph{Proof.} Cited: the different is multiplicative in towers,
$\mathfrak D_{E/\Q}=\mathfrak D\cdot\mathfrak D_{k_d/\Q}\mathcal O_E$, and the absolute norm of the
different of a number field is the absolute value of its discriminant (J.~Neukirch, \emph{Algebraic Number
Theory}, 1999, Ch.~III, (2.2)(i) and (2.9)). The norm of $\mathfrak D_{k_d/\Q}\mathcal O_E$ is
$\mathrm N(\mathfrak D_{k_d/\Q})^{[E:k_d]}=d^{f}$, since the discriminant of $k_d$ is $-d$. For the
second statement divide the logarithm by $[E:\Q]=2f$. \qed

\textbf{Lemma B (the integral line).} \emph{For a nonzero fractional ideal $\mathfrak c$ of $E$ let
$\mathcal M_{\mathfrak c}$ be the $\mathcal O_M$-module spanned by the vectors
$(\tau'(\alpha))_{\tau'}$, $\alpha\in\mathfrak c$. Then the line
$\{z\in M: z\,\mathbf e_\tau\in\mathcal M_{\mathfrak c}\}$ is the $\mathcal O_M$-module generated by
$\tau(\mathfrak c\,\mathfrak D)$. In particular
$\mathfrak l_\tau=\tau(\mathfrak a_0\mathfrak D)\mathcal O_M=\tau(\bar{\mathfrak b})\mathcal O_M$.}

\emph{Proof.} Let
$\mathfrak c^{*}=\{\alpha^{*}\in E:\mathrm{Tr}_{E/k_d}(\alpha^{*}\mathfrak c)\subseteq\mathcal O_{k_d}\}
=\mathfrak c^{-1}\mathfrak D^{-1}$ (the definition of the different), and let $\mathcal M^{*}$ be the
$\mathcal O_M$-module spanned by the vectors $(\tau'(\alpha^{*}))_{\tau'}$,
$\alpha^{*}\in\mathfrak c^{*}$. Use the bilinear form
$\langle\mathbf m,\mathbf m'\rangle=\sum_{\tau'}m_{\tau'}m'_{\tau'}$ on $M^f$. For $\alpha\in\mathfrak c$
and $\alpha^{*}\in\mathfrak c^{*}$ the pairing of their vectors is
$\sum_{\tau'}\tau'(\alpha\alpha^{*})=\mathrm{Tr}_{E/k_d}(\alpha\alpha^{*})\in\mathcal O_{k_d}\subseteq\mathcal O_M$.
(Every $\tau'$ restricts to the given embedding of $k_d$; so the map
$\alpha\mapsto(\tau'(\alpha))_{\tau'}$ is $k_d$-linear, and the sum over $\tau'$ is the trace to $k_d$.)
We claim that $\mathcal M_{\mathfrak c}$ is exactly the set of $\mathbf m$ with
$\langle\mathbf m,\mathbf m'\rangle\in\mathcal O_M$ for all $\mathbf m'\in\mathcal M^{*}$. It is enough
to check this after localising at each prime of $\mathcal O_M$. Above a prime $\mathfrak p$ of $k_d$ the
module $\mathfrak c_{\mathfrak p}$ is free over the local ring of $k_d$, with a basis
$\alpha_1,\dots,\alpha_f$, and $\mathfrak c^{*}_{\mathfrak p}$ has the dual basis
$\alpha_1^{*},\dots,\alpha_f^{*}$ for the trace form. The corresponding vectors $\mathbf a_i$,
$\mathbf a_i^{*}$ of $M^f$ satisfy $\langle\mathbf a_i,\mathbf a_{i'}^{*}\rangle=\delta_{ii'}$, so they
are two bases of $M^f$, dual to each other, and they generate the localisations of
$\mathcal M_{\mathfrak c}$ and $\mathcal M^{*}$. This proves the claim. Now $z\,\mathbf e_\tau$ pairs with
the vector of $\alpha^{*}$ to $z\,\tau(\alpha^{*})$. So $z\,\mathbf e_\tau\in\mathcal M_{\mathfrak c}$ if
and only if $z\,\tau(\alpha^{*})\in\mathcal O_M$ for all $\alpha^{*}\in\mathfrak c^{*}$, that is, if and
only if $z$ lies in the inverse of the ideal $\tau(\mathfrak c^{*})\mathcal O_M$, which is
$\tau(\mathfrak c\,\mathfrak D)\mathcal O_M$. \qed

\textbf{Lemma C (the index).} \emph{For nonzero $z\in\mathfrak l_\tau$,\quad
$[\mathfrak l_\tau:\mathcal O_Mz]=|\mathrm N_{M/\Q}(z)|\cdot\mathrm N(\bar{\mathfrak b})^{-[M:\Q]/[E:\Q]}$.}

\emph{Proof.} For a nonzero ideal $\mathfrak A$ of $\mathcal O_M$ and $z\in\mathfrak A$ one has
$[\mathfrak A:z\mathcal O_M]=|\mathrm N_{M/\Q}z|/\mathrm N_M(\mathfrak A)$, by multiplicativity of the
norm. By Lemma B, $\mathfrak A=\mathfrak l_\tau$ is the extension to $M$ of the ideal
$\tau(\bar{\mathfrak b})$ of $\tau(E)$, and the norm of an extended ideal is the norm raised to the power
$[M:\tau(E)]$. \qed

\textbf{Lemma D (the product formula).} \emph{For nonzero $z\in M$,\quad
$\sum_\rho\log|\rho(z)|^2=2\log|\mathrm N_{M/\Q}(z)|$.}

\emph{Proof.} $\mathrm N_{M/\Q}(z)=\prod_\rho\rho(z)$, the product over all embeddings. \qed

\noindent\emph{Proof of the Proposition.} By Lemmas C and D the terms with $z$ cancel:
\begin{equation}\label{D:289}
 2\log[\mathfrak l_\tau:\mathcal O_Mz]-\sum_\rho\log|\rho(z)|^2
 =-\frac{2[M:\Q]}{[E:\Q]}\log\mathrm N(\bar{\mathfrak b}) .
\end{equation}
In particular a unit, or any other change of $z$, does not change \eqref{sa:meandef}. Since
$\mathrm N(\mathfrak b)=\mathrm N(\bar{\mathfrak b})$ and $\mathfrak D=\mathfrak b\bar{\mathfrak b}$,
$2\log\mathrm N(\bar{\mathfrak b})=\log\mathrm N(\mathfrak D)$. By Lemma A,
\begin{equation}\label{D:290}
 \overline{\log\Omega_\tau}=\frac1{[M:\Q]}\sum_\rho\log\frac{A_\rho}\pi-\frac1{[E:\Q]}\log\mathrm N(\mathfrak D)
 =\frac1{[M:\Q]}\sum_\rho\log\frac{A_\rho}\pi-\log\frac{\mathrm{rd}_E}{\sqrt d}.
\end{equation}
The first term does not depend on $M$. Indeed, let $M_0\subseteq M$ be a field that contains the
coefficients $g_2,g_3$ of the curve. The number $A_\rho$ depends only on the restriction of $\rho$ to $M_0$,
and every embedding of $M_0$ extends to exactly $[M:M_0]$ embeddings of $M$. So the mean over the
embeddings of $M$ is the mean over the embeddings of $M_0$, and two fields $M$ are compared through their
compositum. This proves the first assertion.
By (I2) and (I3) of \S III.4 the mean of $\log A_\rho$ is $\log\pi+\gamma+\log(2\pi/\sqrt d)-c_d$. Subtracting, $\overline{\log\Omega_\tau}=\gamma+\log(2\pi/\mathrm{rd}_E)-c_d$, which is
$\log\Omega_E$ by its definition in \S I.7. \qed

\noindent So the constant $\Omega_E$ differs from the elliptic $\Omega$ by the factor
$\mathrm N(\mathfrak D)^{1/[E:\Q]}=\mathrm{rd}_E/\sqrt d$. Two things give this factor. The first is
Lemma B: the integral line at a weight place is generated by $\tau(\mathfrak a_0\mathfrak D)$, not by
$\tau(\mathfrak a_0)$. The second is the choice $\mathfrak a_0=\mathfrak b^{-1}$, for which
$\mathrm N(\mathfrak a_0\mathfrak D)^2=\mathrm N(\mathfrak D)$. For an arbitrary nonzero fractional ideal
$\mathfrak a_0$ the same proof gives
\begin{equation}\label{D:291}
 \overline{\log\Omega_\tau}=\frac1{[M:\Q]}\sum_\rho\log\frac{A_\rho}\pi
 -\frac2{[E:\Q]}\log\mathrm N(\mathfrak a_0\mathfrak D) ;
\end{equation}
with $\mathfrak a_0=\mathcal O_E$ the factor would be $(\mathrm{rd}_E/\sqrt d)^2$.
(Lean: \texttt{meanPeriodAt}, from \texttt{mean\_log\_period}, \texttt{br\_meanLogOm\_eq} and
\texttt{two\_log\_brN}; Lemma B is \texttt{lineR\_eq}, Lemma C is \texttt{index\_lineR}, Lemma D is
\texttt{br\_sum\_log\_norm\_sq}, Lemma A is \texttt{natAbs\_discr\_eq}. The source proves Lemma A from
the same two facts.)

\emph{What the Proposition proves, and what it does not.} It proves that the number
\eqref{sa:meandef}, which is defined by the ideal $\mathfrak a_0$ and the areas $A_\rho$, equals
$\log\Omega_E$, under (I2), (I3) and the fact about $E$. Input (S4) is a different statement. It concerns
the numbers $\Omega_\tau(v)$ of (S1), averaged over the archimedean places $v$ of $K$. To deduce (S4)
from the Proposition one needs three further statements. (P1) and (P3) are assumed, from the source,
and not proved here. (P2) was proved above:
\begin{enumerate}[label=\textup{(P\arabic*)},leftmargin=2.8em,itemsep=2pt,topsep=3pt]
\item the period that occurs in the block estimate of \S III.4 for the kernel conjugated by $\rho$, at
the weight place $\tau$, is $A_\rho/(\pi|\rho(z_\tau)|^2)$ for a generator $z_\tau$ of $\mathfrak l_\tau$
(Lean: \texttt{logOmAt}, \texttt{ExtraTie.om\_eq}). Generators need not exist over a given $M$. An example: $M=E$,
$q=3$, $d=23$, $N=4$, $8$ or $16$ (the class number of $E$ is $3$, $6$, $60$). There the ideal
$\tau(\bar{\mathfrak b})\mathcal O_E$ of Lemma B is not principal for any of the choices of $\mathfrak b$
coprime to $\bar{\mathfrak b}$. This was computed with Sage 10.9, with class groups computed in its mode
\texttt{proof=False}; so the example is correct under GRH and was not checked without it. Nothing in
this section depends on the example. It depends on the choice of $\mathfrak b$: for $N=8$ two of the
eight integral $\mathfrak b$, which are not coprime to their conjugates, give a principal line. The
existence of generators over a suitable larger field is the hypothesis \texttt{IsDgen} of \texttt{mean\_logMet\_eq};
\item every prime of $E$ that divides $\mathfrak D$ lies over $2$ or over $q$ (Lean:
\texttt{diffSupport}). This is the first step of the fact about $E$ above; it is proved there for the
field $E$ of the paper. With it the ideal $\mathfrak b$ exists;
\item the mean over the embeddings $\rho$ of $M$ can be replaced by the mean over the places $v$ of $K$:
the source takes $M\supseteq K$ and for $\log\Omega_\tau(v)$ an average, with weights of sum $1$, of the
numbers $\log\big(A_\rho/(\pi|\rho(z_\tau)|^2)\big)$, the weights having the same total for every $\rho$
(Lean: \texttt{kerSide\_of\_meansS}, \texttt{exists\_weights}).
\end{enumerate}

\subsection*{5.\quad The remaining inputs: statements and status}

Each statement below is assumed in this section. We give its form in the source and say what a proof
needs. The Lean names are collected in subsection 6.

\textbf{(S1), the chain bound.} The analytic content is proved in \S III.4: the bound for one block with
the period (``Several blocks, and the period''), the bound for one map (``Sums over classes''), and the
bound for the Gaussian constant by a mean square. Compare with the formula for $\log b_j(v)$ and with the
constant $S$ of (\ref{III:Sform}) there. $|\mathcal P|\,\mathfrak g_j$ is the number $g=2nNN'$ of variables,
with $N$, $N'$ the ranks of the pair as in \S III.4. $\beta^{*}_j(v)$ is the prefactor
$\beta\,c^{2[E:\Q]NN'}$, and $\mathrm{cl}_j,\mathrm{cl}_j'$ count classes. $\Theta_1(v)$ is a mean square
per variable. In \S III.4 the reference mean square at an embedding $\rho$ is
$\Theta_{0,\rho}=H_\rho^{fNN'}$, with $H_\rho$ the mean square of one elliptic theta function. It enters
$S^2$ once, with $\langle\log\Theta_{0,\rho}\rangle\le fNN'L_{\rm sc}$. In (S1) the same term is written
$|\mathcal P|\,\mathfrak g_j\log\Theta_1(v)$, with $\langle\log\Theta_1\rangle\le C_{\rm th}L_{\rm sc}$. So
$\Theta_1$ corresponds to $H_\rho$, not to $\Theta_{0,\rho}$. What is assumed: that the maps of the construction are
instances (item (vi) of the list in \S III.4), and the eight numerical bounds listed under (S1). Of
these, the bound for $\langle\log\Theta_1\rangle$ is the conclusion of the paragraph ``The reference mean
square'' of \S III.4 (conditional on (I1)--(I5)), divided by the number of variables; the bounds for
$\beta^{*}_j$ and for the numbers of classes are statements about the finite data of the construction
(coefficients, modulus, class numbers of Hermitian lattices of bounded level). A proof needs a bound for
class numbers of definite Hermitian lattices in terms of the mass and the level. The source's statements
were not read for this section, with one exception, the sign of $v_{\rm h}$. In the record of the source
$v_{\rm h}$ is a real number with no condition on its sign (Lean: the field \texttt{vv} of
\texttt{HtNums}, and the same in \texttt{KerSide}); the logarithm of the Lean is defined for every real
number, so the record needs none. Where the source builds the record, $v_{\rm h}=\mathrm{rd}_E/2$ (Lean:
\texttt{NumDatum.v}), and $v_{\rm h}\ge0$ is proved from $\mathrm{rd}_E\ge1$ inside the proof of the bound
with $c_5$ (\texttt{log\_eight\_leS}). The source's statement of the analytic bound for one map
(\texttt{LiftedMapEmbeddingLogBound}) has $v_{\rm h}>0$ among its hypotheses. So $v_{\rm h}\ge0$ is not a
hypothesis of the record, and it is a consequence of the construction.

\textbf{(S2), the pairing at the end.} The inequality $\|\Pi\|_v^2\le e^{h_\Pi}$, with $h_\Pi$ bounded
as in the first line of (S2) with $e_\Pi=0$, is the conclusion of the paragraph ``The pairing at the end
of the chain'' of \S III.4, with its own inputs; there the two terms are bounded by
$\log|\vartheta|\le\psi\,d_1\log\Gamma$ and $\log\kappa\le\frac{\Psi(\Psi-1)}2\log\Gamma$. The source
does not name $\log|\vartheta|$ and $\log\kappa$: its numbers $\mathfrak m_\tau$, $\mathfrak m_\tau'$ are
these two majorants with $\psi=N_m-1$, $\Psi=N_m$ and $2a$ in the place of $\Gamma$. The identification
of $d_{0,\tau}$, $d_{1,\tau}$, $\ell_\tau$, $\lambda_{1,\tau}$, $b_\tau$ with the $d_0$, $d_1$, $\ell$,
$\lambda_1$, $b_{\mathfrak t}$ of \S III.4 was confirmed in the Lean
(\texttt{lastSide\_of\_logT\_le}). The remaining conditions of (S2) express the exponent
$\ell_\tau=\lambda_{1,\tau}-d_{0,\tau}$ and the degrees $d_{0,\tau}$, $d_{1,\tau}$ through the weights of
the construction: at a distinguished place the first weight is $a_\tau$ up to the shift $(N_m-k)/2$. They
are statements about the weights only; this is the ``substitution of the actual weights'' of item (ii)
of the list in \S III.4. A proof needs the weights of the last two spaces of the chain as explicit
functions of the data. In the source $e_\Pi=0$ and $C_\Pi=0$.

\textbf{(S3), the eigenvalues.} This is the paragraph ``The eigenvalues of a lift'' of \S III.4, with its
hypotheses (E1)--(E7). The source's evaluation of the scalar is the following statement (Lean:
\texttt{scalar\_bound\_lowC}; among the declarations the final theorem depends on). Let
\begin{equation}\label{D:292}
 \lambda=D_0\cdot\mathrm{vol}(\mathcal F)\cdot
 \prod_\tau H(A_\tau^+)H(A_\tau^-)\cdot\frac{\|\phi\|^2}{D_{S^c}}\cdot\prod_\tau J_\tau\cdot|E_0|,
\end{equation}
where $D_0\ge1$ is a product of dimensions of representations; $\mathrm{vol}(\mathcal F)$ is the volume
of a fundamental set for the level subgroup of the definite group; $\|\phi\|^2$ is the squared
$L^2$-norm of the finite test function for a measure in which the rational lattice has covolume one;
$J_\tau$ is the integral, over the real unitary group at $\tau$, of the squared absolute value of a
matrix coefficient of a unit vector of the lowest type for the maximal compact subgroup; $E_0$ is an
Euler product over the prime ideals $\mathfrak P$ of $E$ outside a finite set, with parameters of absolute value
at most $(\mathrm N\mathfrak P)^{(r-1)/2-(1+t_0)/2}$, as in (E5); and $D_{S^c}$ is the normalising factor
outside that set. Assume the level-index identity (the index of the level subgroup times $\|\phi\|^2$
equals $1/D_S$), the lower bound $\big(w_E^{\,\varrho}\,\varrho!\big)^{-1}$ for the mass, with $w_E$ the
number of roots of unity of $E$, and the bound
$-\log J_\tau\le2+r^2\log\big(2\sqrt2\,r^2\varrho(d_\tau+2)\big)$. Then $\lambda>0$ and
\begin{equation}\label{D:293}
 \log\lambda\ \ge\ \sum_\tau\big(\log H(A_\tau^+)+\log H(A_\tau^-)\big)-8f\,\varrho^2\,L,\qquad
 L=\log\big(b\,\mathrm{rd}_E\,g_0g_1\,\varrho\big).
\end{equation}
This is the deduction written in \S III.4, with one difference. The mass bound in (\ref{III:massbd}) is
$\big((8f^2)^{\varrho}\varrho!\big)^{-1}$; here and in the Lean it is
$\big(w_E^{\,\varrho}\varrho!\big)^{-1}$. The second implies the first, because $w_E\le8f^2$. The reason:
write $\mathrm{eul}(u)$ for Euler's function of an integer $u\ge1$. The field $E$ contains
$\Q(\zeta_{w_E})$, which has degree $\mathrm{eul}(w_E)$ over $\Q$. So $\mathrm{eul}(w_E)$ divides
$[E:\Q]=2f$, and $\mathrm{eul}(w_E)\le2f$. Next, $u\le2\,\mathrm{eul}(u)^2$ for every $u\ge1$: the
quotient $\mathrm{eul}(u)^2/u$ is the product of the numbers $p^{e-2}(p-1)^2$ over the prime powers $p^e$
that divide $u$ exactly, and such a number is at least $\frac43$ for $p\ge3$ and at least $\frac12$ for
$p=2$. Hence $w_E\le2(2f)^2=8f^2$. The Lean has this inequality as a theorem, for the field $E$ of its
record \texttt{Host}: \texttt{host\_torsionOrder\_le}, which states
\texttt{Units.torsionOrder H.E} $\le8f^2$. The inner-product
formula itself is used in the following form (Lean: the hypothesis \texttt{InnerProductFormulaC} of
\texttt{term\_eval\_lowC}): for theta lifts $\theta(\phi_i,\mathfrak f_i)$, to the definite group, of
forms $\mathfrak f_i$ on the group $\mathsf G$ lifted from,
\begin{equation}\label{D:294}
 \big\langle\theta(\phi_1,\mathfrak f_1),\theta(\phi_2,\mathfrak f_2)\big\rangle
 =\mathrm{vol}(\mathcal F)\int_{\mathsf G(\mathbb A)}
 \langle\omega(g)\phi_1,\phi_2\rangle\ \overline{\langle g\cdot\mathfrak f_1,\mathfrak f_2\rangle}\ dg ,
\end{equation}
the integral being written as a sum over finitely many cosets at the finite places and an integral over
the real points, absolutely convergent. This is the inner-product formula in the convergent range, as an
integral of matrix coefficients; it is not the form with a special value of an $L$-function, and no
Siegel--Weil formula appears in the statement. For the classical statements see S.~Rallis,
\emph{$L$-functions and the oscillator representation}, Lecture Notes in Math.~1245 (1987) [Ral87] and
\emph{Injectivity properties of liftings associated to Weil representations}, Compositio Math.~52 (1984);
J.-S.~Li, \emph{Non-vanishing theorems for the cohomology of certain arithmetic quotients}, J.~reine
angew.~Math.~428 (1992) (the formula as an integral of matrix coefficients, in the stable range);
W.~T.~Gan, Y.~Qiu, S.~Takeda, \emph{The regularized Siegel--Weil formula (the second term identity) and
the Rallis inner product formula}, Invent.~Math.~198 (2014); S.~Yamana, \emph{$L$-functions and theta
correspondence for classical groups}, Invent.~Math.~196 (2014). All from memory; none was consulted. A
proof of (S3) needs: this formula for the actual test functions; the unramified computation of the local
integrals (the factors of $E_0$); the mass formula for definite unitary groups (for instance \cite{GHY01}; its formula is for maximal lattices, and local comparisons would be needed for the levels used here); the formal degree of the discrete series at the real places (Harish-Chandra); the transport to
the other embeddings (E7); and the description of the target by Hecke components (E6). The source has a
theorem that derives the inner-product formula (\texttt{innerProductFormula\_of\_hypotheses}) from four
hypotheses. Three are statements named \texttt{Published.CuspRapidDecay},
\texttt{WeilSiegelFormulaUnitary} and \texttt{SiegelEisensteinContinuation}. The fourth is the existence
of a choice of parameters (Lean: a \texttt{ParameterChoice'} exists for some \texttt{SizeData}$_2$). Each of
the three named statements has a closed theorem in the source. Their proofs
were not read. The dependency computation shows that the final theorem depends on \texttt{innerProductFormula\_of\_hypotheses}; the three named statements were not put to it.

\textbf{(S5), the slope at the lower level.} In the source this is a theorem
(\texttt{IndefSide.boundsMu}, from \texttt{slope\_bounds}). Its conclusion is the display of (S5), with
the sum of the $l_{3,j}$. Its hypotheses are the following. The first is the covolume bound: the
normalised logarithmic covolume of the lattice of the $j$th lift, divided by $\dim Y$, is at most
$l_{3,j}$ (Lean: the field \texttt{l3} of \texttt{IndefSide}; in the source the covolume is minus the
normalised slope). The others are all at finite places. For a finite place $y$ of $K$ write
$\mathrm{wt}(y)=\log\mathrm N(y)/[K:\Q]$, $\mathrm{ram}(y)$ for its ramification index and
$\mathrm{rc}(y)$ for its residue characteristic. There are a finite set $\mathcal S_{\rm fin}$ of places
containing the set $\mathcal S_{\rm tw}$ of places where a twist is made, and real numbers
$\mathrm{gd}(y)$, $\mathrm{ix}(y,j)$, $\mathrm{err}(y,j)\ge0$, $\mathrm{jx}(y,j)$, $\mathrm{et}(y,j)$,
such that: the valuation at $y$ of a Gram determinant of the quotient lattice is $0$ outside
$\mathcal S_{\rm fin}$ and at most $\dim Y\cdot\mathrm{gd}(y)$ on $\mathcal S_{\rm fin}$; on
$\mathcal S_{\rm tw}$ the sum over $j$ of the two local indices of the lattice of the $j$th lift (at $y$
and at $\bar y$) is at most
$\big(\sum_j(2\,\mathrm{ram}(y)\,\mathrm{ix}(y,j)+\mathrm{ram}(y)\,\mathrm{err}(y,j))
+2\,\mathrm{ram}(y)/\mathrm{rc}(y)^2\big)\dim Y$; a further local index is at most
$\mathrm{ram}(y)\,\mathrm{jx}(y,j)\dim Y+\mathrm{et}(y,j)$; outside $\mathcal S_{\rm tw}$, for each $j$, the
sum of these three local indices (the two at $y$ and $\bar y$ and the further one) vanishes, and so does
$\mathrm{ix}(y,j)$ (Lean: \texttt{IndefData.Untwisted}; it asks for the vanishing of the sum, not of each
index). Then (S5) holds with $\mathfrak F$ as follows, and with every $E_{\rm loc}$ that is at least
the sum written below:
\begin{align}\label{D:295}
 \mathfrak F&=\sum_j\sum_{y\in\mathcal S_{\rm tw}}\mathrm{wt}(y)\,\mathrm{ram}(y)\,\mathrm{ix}(y,j),\\
 E_{\rm loc}&=\sum_{y\in\mathcal S_{\rm fin}}\mathrm{wt}(y)\,\mathrm{gd}(y)
 +\sum_j\sum_{y\in\mathcal S_{\rm tw}}\mathrm{wt}(y)\Big(\mathrm{ram}(y)\big(\mathrm{err}(y,j)+\mathrm{jx}(y,j)\big)
 +\frac{\mathrm{et}(y,j)}{\dim Y}\Big)\\
 &\quad+\sum_{y\in\mathcal S_{\rm fin}}\frac{2\,\mathrm{wt}(y)\,\mathrm{ram}(y)}{\mathrm{rc}(y)^2}.
\end{align}
The display gives $E_{\rm loc}$ in the case of equality. In the source the four sums of $E_{\rm loc}$
are not computed: each is bounded by a free real number (Lean: \texttt{LG}, \texttt{E1}, \texttt{E2},
\texttt{E3}), and the conclusion has the sum of these majorants. The Lean \texttt{E3} is the last sum of
$E_{\rm loc}$. It is not the $E_3$ of (S6), which the Lean calls \texttt{El3}.
It is a statement of the geometry of numbers over $K$: the covolume of a lattice glued from local
lattices, compared with the covolumes of the lattices of the lifts through the scalar relations of (S3)
(which is where $\langle\log\varkappa\rangle$ enters) and through local indices. No period occurs in it.
The term $\mathfrak F$ is the term $F_0/2$ of Part II, \S6.7. The local index bounds are separate inputs;
those at the places of $\mathcal S_{\rm tw}$ are of the kind treated in \S III.5. The proof in the source
was not read.

\textbf{(S6), the covolumes.} In the source
$l_{3,j}=\sum_\tau p_jw_{j\tau}\log\Omega_E+C_3\,B_0\log\log x$ by definition, with a constant $C_3$
(Lean: \texttt{ClosedScale3.l3m}, \texttt{l3u}) and with $p_j=k+2j-1-m$. So (S6) holds with
$E_3=m\,C_3\,B_0\log\log x$ (\texttt{levelThree\_bound\_closed}). The content is therefore in the first
hypothesis of the theorem quoted under (S5), the covolume bound (\texttt{IndefSide.l3}): the normalised
logarithmic covolume of the lattice of the $j$th lift, divided by $\dim Y$, is at most
$\sum_\tau p_j\,w_{j\tau}\log\Omega_E$ up to $C_3\,B_0\log\log x$. This is the count of forms of Part II,
\S6.6 in its arithmetic form: the lattice is a space of algebraic automorphic forms of weight governed
by $w_{j\tau}$, its metric is the Petersson norm, and the period appears because integral forms are
normalised by the algebraic differential, with the same integral lines $\mathfrak l_\tau$ as in
subsection 4. \emph{That the constant here is the same $\Omega_E$ as in \textup{(S4)} is exactly what
makes the cancellation in Step 4 possible; it is assumed, not proved, in this section.} A proof needs the
covolume of the lattice of integral forms on a definite unitary group for the Petersson metric (mass
formula and dimension formula, as in \eqref{II:t6.4} of Part II), with the period normalisation of subsection 4.
The source derives it from a statement named \texttt{CovolumeBoundThirdPlaceSeries}
(\texttt{linked\_bundles\_l3\_of\_open\_alt}), whose hypotheses include statements on discrete series and
on Harish-Chandra modules; it was not read beyond its header. One of its hypotheses is a covolume bound
named \texttt{CovolumeBoundOfX}. The source proves that one (\texttt{topCovolSlim2\_use\_alt\_alt}) from
statements of algebraic and Arakelov geometry, which it takes as hypotheses: among them resolution of
singularities and principalisation in characteristic zero, a Poincar\'e--Lelong formula, an arithmetic
Hilbert--Samuel statement, a bigness statement named after Moriwaki, toroidal compactifications, and a
list named \texttt{PublishedHodgeBundleTheorems}. Only the headers of these two theorems were read. The
statements passed to this bound are listed, with their published sources, in the census of \S\ref{cs:sec}:
they include every row of its groups on abelian varieties and on algebraic and arithmetic geometry (groups~J
and~K).

\textbf{(S7), the passage between levels.} Two statements: the bound
$\mu_{\max}(\mathcal R_0)\le H_{\max}-\langle\log\varkappa\rangle$ with $H_{\max}=x^{C_H}b$ for a
constant $C_H$ (Lean: \texttt{IndefExtraR.fun0}), and the comparison of $\mu(\mathcal R^\vee)$ with
$\mu(\mathcal R_0^\vee)$ displayed in (S7) (\texttt{IndefExtraR.passage}). In the source $\mathrm{la}$ is
the logarithm of the order of a finite group, and $\mathrm{de}=8fr/\min(g_0,g_1)$ with $r=k+m-1$. The
number $Q_0$ is $\max(1,\,\cdot\,)$ of a majorant of the left side of its defining inequality: any upper
bound for $\sum_jl_{3,j}+2\mathfrak F+E_{\rm loc}$, plus $H_{\max}+28\,\mathrm{la}$. The condition
$3\,\mathrm{de}\,Q_0\le1$ then says that the level primes $g_0,g_1$ are at least $24\,f\,r\,Q_0$. (The
source proves an inequality of this form, $24\,f\,r\,Q\le\min(g_0,g_1)$ for its number \texttt{QOfR}, from
$g_0,g_1\ge(kb)^8$ and further size conditions: \texttt{hg\_of\_sizeG}, \texttt{SizeG}. Only the statement was read.) So the
passage costs a bounded amount because the level primes are large. The comparison is a statement about
two lattices of forms of different level and the trace map between them; the first bound is a bound for
the largest slope of a lattice of forms. Proofs were not read; the source reaches them through
statements on Fourier coefficients (\texttt{fittedFunction\_at\_zero\_close\_alt}).

\textbf{(S8), the size of the local errors,} and \textbf{(G)}. These are inequalities between the
parameters. (S8) contains the only occurrence of $N_c$. With $E_3=m\,C_3\,B_0\log\log x$ as above, and
with the bounds for $\mathrm{wt}(y)\,\mathrm{gd}(y)$, $\mathrm{err}$, $\mathrm{jx}$, $\mathrm{et}$
supplied by the local statements, it holds for large $x$. The part of $E_{\rm loc}$ with
$\mathrm{et}(y,j)$ is where the bound for Galois cohomology of Part II, \S6.8 is used (Lean:
\texttt{divinN\_of\_hypV}); this was read in the notes of another reader and not checked here. (G) is a
list of size conditions on the parameters of the chain, valid for $x\ge x_0$ in the source; among them
$|\log\Omega_E|\le L_a$ follows from $|c_d|\le4\log x$ (Theorem~\ref{thm:pair}(d)) for large $x$.

\textbf{(W1), (W2).} See remark (1) after the list of inputs: they follow from the conclusions of ``Row
by row'', which are conditional on the inputs (o), (a)--(e) there, and from three facts of the source.

\subsection*{6.\quad Assumed here, from the source}

The Lean statements corresponding to the inputs. ``Dependency: yes'' means that the name is among the
declarations the final theorem depends on, by the dependency computation; ``reached'' means that the name
is not in the lists computed, but is used in the proof or in the definition of a name that is, either
directly or through further declarations that are not in the lists. The route is given each time.
\begin{itemize}[leftmargin=1.6em,itemsep=2pt,topsep=3pt]
\item \texttt{divClausesR\_of} (dependency: yes): the assembly; its third clause is Theorem
\ref{sa:thm}(b) with the error $(50C_0+10)\,\mathcal U_1+60c_e\,\mathcal U_2$, its second clause is
Theorem \ref{sa:thm}(a) in the form of a majorant at each place.
\item (S1): fields \texttt{cb}, \texttt{nKcb} of \texttt{DivInputs0R} and \texttt{i2}--\texttt{i10},
\texttt{k2}, \texttt{k3} of \texttt{HtNums}: each matrix of the chain is a sum over classes of blocks
with squared norm at most a bound whose logarithm is at most $\Sigma_j(v)$.
\texttt{chain\_coupled\_bound} (yes) turns this into $\log\|K_j\|_v^2\le\Sigma_j(v)$.
\item (S2): field \texttt{nP} of \texttt{DivInputs0R} and \texttt{t1}--\texttt{t11}, \texttt{k1} of
\texttt{HtNums}.
\item (S3): fields \texttt{l1b}, \texttt{l1c}, \texttt{l6} of \texttt{IndefColsR} and \texttt{eig6} of
\texttt{DivInputs0R}; \texttt{scalar\_bound\_lowC} (yes). The source has a theorem named
\texttt{lamLow\_of\_classicalTheorems}; among its hypotheses are the Siegel--Weil formula for unitary
groups with its convergence, the doubling method, and the unramified doubling computation for
$\mathrm{GL}$ (\texttt{weilSiegelFormulaUnitary}, \texttt{siegelWeilUnitaryConvergent},
\texttt{doublingMethod}, \texttt{unramifiedDoublingGL}). By its name it gives the lower bound for the
scalars. That it is what fills the field \texttt{l6} was not verified.
\item {\sloppy (S4): field \texttt{i1} of \texttt{HtNums}; \texttt{meanPeriod} (yes), \texttt{mean\_log\_period}
(yes), \texttt{chowla\_selberg\_eta} (yes). Reached: \texttt{two\_log\_brN} and \texttt{br\_meanLogOm\_eq},
through \texttt{meanPeriodAt}, which is the body of \texttt{meanPeriod} and is not in the lists;
\texttt{index\_lineR}, through \texttt{br\_meanLogOm\_eq}; \texttt{lineR\_eq}, through
\texttt{index\_lineR}. The ideal $\mathfrak b$: \texttt{principal}, \texttt{Arith.bb\_ex}, field
$\mathfrak b$\texttt{\_spec} of \texttt{Host} (dependency: yes for \texttt{principal}, by a later run of the computation, which finds it in the
namespace \texttt{Formalism.HC.Der}).\par}
\item (S5): \texttt{IndefSide.boundsMu} (yes), \texttt{slope\_bounds} (yes), with the fields
\texttt{bi}, \texttt{bLG}, \texttt{biv}, \texttt{borderedMinorImage},
\texttt{twisted\_third\_sum\_le} of \texttt{DivInputs0R}, \texttt{on} and \texttt{l3} of
\texttt{IndefSide} and \texttt{untw} of \texttt{IndefExtraR}.
\item (S6): field \texttt{L3main} of \texttt{DivInputs0R}.
\item (S7): fields \texttt{level}, \texttt{passage}, \texttt{fun0} of \texttt{IndefExtraR}; fields
\texttt{de0}, \texttt{deQ}, \texttt{Q1}, \texttt{Qb}, \texttt{la\_le},
\texttt{constant\_ge\_twenty\_eight} of \texttt{DivInputs0R}.
\item (S8): field \texttt{errB} of \texttt{DivInputs0R}.
\item (W1), (W2), (G): fields \texttt{weight\_height\_sum\_le}, \texttt{weight\_sum\_eq}, \texttt{sE},
\texttt{s}$\delta$, \texttt{p1}--\texttt{p38}, \texttt{c1}--\texttt{c5} of \texttt{RecNums}; for remark
(1), \texttt{NumFacts}, \texttt{stepE\_num}, \texttt{RecNums.ofNum}.
\item {\sloppy Step 1: \texttt{height\_chain\_atR} (yes); \texttt{height\_constant\_bound} (reached directly: it
is used in the proof of \texttt{height\_chain\_atR}). Step 2:
\texttt{height\_bound\_at\_set} (yes). Step 4: \texttt{divisor\_inequality\_explicit} (yes),
\texttt{divIneqRecord} (yes), \texttt{error\_absorption} (yes); \texttt{divisor\_inequality\_assembly}
(reached through \texttt{divisor\_inequality}, which is not in the lists and is applied by
\texttt{divisor\_inequality\_explicit}).\par}
\item The identifications $h=\mathcal H-\langle\log\varkappa\rangle$, $F_0=2\mathfrak F$,
$\mu_{\rm CC}=\mu(\mathcal R^\vee)$ with the quantities of the final statement of the source:
\texttt{hTie\_of} (yes) and the hypotheses \texttt{tMu}, \texttt{tFo}, \texttt{tR} of
\texttt{wordFacts\_of\_clauses}. In the source the function $h(u)$ of the final statement is defined by
the formula for $\mathcal H-\langle\log\varkappa\rangle$ (\texttt{hExpl6}); so it already contains
$-\langle\log\varkappa\rangle$.
\end{itemize}

\subsection*{7.\quad What this section proves, and what it assumes}

\emph{Proved:} Theorem \ref{sa:thm}: the inequality (\ref{III:sizefinal}) at the end of \S III.4 follows from
(S1)--(S8), (W1), (W2), (G), with the constant $C=\max(50C_0+10,60c_e)$; the descent lemma; the two
lemmas on real numbers, with every constant; the table of the origin of each term. The Proposition of
subsection 4: the mean \eqref{sa:meandef} of the logarithms of the periods of the integral lines is
$\log\Omega_E$, given the class mean formula (I2), (I3) of \S III.4, the multiplicativity of the
different and the formula for its norm (cited), and the ideal $\mathfrak b$; Lemmas B, C, D are proved.
The existence of $\mathfrak b$ for the field $E$ of the paper, with (P2), on which it rests. Remark (1) after the list of
inputs: (W1), (W2) from the conclusions of ``Row by row'' and three facts of the source.

\emph{Assumed:} (S1)--(S8), (W1), (W2), (G). The input (S4) is not proved: it would follow from the
Proposition and (P1)--(P3), and (P1), (P3) are assumed from the source. Of the inputs, (S1), (S2), (S3),
(W1), (W2) correspond to conclusions of conditional paragraphs of \S III.4, each with inputs of its own
that are named there. For (S1) and (S2) the exact match of the numerical forms (the constants of (S1),
the weights in (S2)) is part of what is assumed. For (W1), (W2) the match is made in remark (1), given
the three facts quoted there.

\subsection*{8.\quad What remains}

\begin{enumerate}[label=\textup{(\roman*)},leftmargin=2.4em,itemsep=2pt,topsep=3pt]
\item The covolume bound with the period $\Omega_E$ of subsection 4 (the first hypothesis of the theorem
behind (S5), with the values of $l_{3,j}$ in (S6)): the statement that the lattice of integral forms of
the $j$th lift has normalised covolume at most $\sum_\tau p_jw_{j\tau}\log\Omega_E$ up to the stated
error. This is the largest gap: the cancellation of the period is proved here only as an identity
between two inputs.
\item The slope bound (S5) and its local index bounds, including the definition of $\mathfrak F$.
\item The passage between levels (S7).
\item The inner-product formula, the mass, the formal degrees and the local factors, for (S3).
\item The statements (P1) and (P3) on the periods of the actual kernels, and item (vi)
of the list of \S III.4.
\item The verification of (G) and (S8) for the parameters of Part II for large $x$, and the bounds for
the constants and the class numbers in (S1).
\end{enumerate}


\section{The images in general position: the trace-formula estimates}\label{tr:sec}

This section concerns the two smallness statements that the linear algebra of Part~II, \S6.9--6.10 consumes,
and the identity that connects them with the pairings of the chain of lifts. In the Lean development these
are three statements: \texttt{LinkS} (bounds for the Gram blocks of a family of maps; the trace statement),
\texttt{LinkIRR} (a count; the rank statement), and \texttt{BesselId} (the pairing at the end of the chain is
an inner product with the image of such a map; the pairing statement). The title needs a word of
explanation. The source does not prove that the images of the maps are nearly orthogonal. The images cannot
all be mutually orthogonal: by the trace statement their dimensions add up to more than the dimension of the
target (Remark~(1) of subsection~2). What the source bounds is this: every block
$B_{\mathsf w}^\dagger B_{\mathsf w'}$ with $\mathsf w\ne\mathsf w'$ has at most $(1+\eta_0)(1+\eta_1)$ times
the mean size that it has for subspaces in general position, where $\eta_0$ and $\eta_1$ are the numbers
of Statement~(S) of subsection~1, at most $\exp(-x^{\varepsilon/8})$. So the operator
$Q=\sum_{\mathsf w}B_{\mathsf w}B_{\mathsf w}^\dagger$ is close to its mean eigenvalue in Hilbert--Schmidt norm
(Remark~(1) of subsection~2). The operator $Q$ is not a sum of projections. \emph{Proved here:} Lemma~T1 and
Corollary~T2. This section deduces the dimension bound from the three trace, rank and pairing statements and
the size and subspace conditions in T2, with explicit constants. It also proves Lemma~T3, a statement about
finite sets and weights which is the mechanism of a second moment in the simplest case (a sum over twists of
$B_\chi^\dagger B_\chi$ is a count of pairs of representations, and its main term is a mass). \emph{Quoted
from the source:} the three statements themselves, transcribed in subsection~1, and the theorems from which
the source deduces them, listed in subsection~4 with their Lean names. The construction of the spaces and
maps and the arithmetic inputs are quoted; none of these is proved here. The assembly of (M) is in
\S\ref{as:sec}: it uses Corollary~T2, through Proposition~M4 there, and it reduces (M) to numbered
Hypotheses, which are not proved. The error term
that Part~II, \S6.10 takes from this section is recalled at the end of subsection~2. Below, \S1, \S2, \dots\
without a part number mean the subsections of this section.

{\sloppy Letters introduced in this section: $\mathsf W$, $\mathsf w$ (words); $\mathcal H$, $\mathcal C_{\mathsf w}$,
$\mathcal M_{\mathsf w}$, $\mathcal E_0$, $\mathcal E$, $\mathcal U_{\mathsf w}$ (spaces), with $\mathcal M_{\mathsf w,\mathbb L}$,
$\mathcal E_{\mathbb L}$ (spaces over a number field $\mathbb L$); $P_{\mathsf w}$, $B_{\mathsf w}$,
$K_{\mathsf w\mathsf w'}$, $\mathsf s$, $\mathfrak B_{\mathsf w}$, $\beta_{\mathsf w}$, $\Psi$, $\Psi_{\mathsf w}$, $\Lambda$, $\pi_{\mathcal E}$
(maps and pairings); $d_{\mathsf w}$, $d^{\rm ram}_{\mathsf w}$, $D_*$, $D_{\rm r}$ (dimensions);
$\eta_0,\eta_1,e_0,\delta_{\rm p},\delta_{\mathcal E},\varkappa,\alpha_{\rm ram},\alpha',\varepsilon_z,\varepsilon_C$ (small
numbers); $\mathsf G$, $C_{\mathsf G}$, $C_0$, $C_H$, $J_z$, $c_0(\mathsf w)$, $\lambda_{\mathsf w}$ (constants);
$\mathrm{Wg}$, $\mathcal T(\mathsf w)$, $\xi_i$, $\psi_i$, $\theta_j$, $e_j'$, $\Delta$ (in \S2, \S4, \S5); and, in \S3 only,
$\mathcal I,\mathcal I',\omega,\omega',G_i,\mathrm{Rep},A,\mathrm{inv},\Theta_\chi,B_\chi,\mathcal N,\mathcal N^\circ,\mathfrak N^\circ,\mathfrak m,\sigma,\ell_0,\varsigma_\iota$.
The letters $Q$, $D$, $S$, $e$, $\eta$, $\alpha$, $u$, $H$ have the meaning of Part~II, \S6.9--6.10: $Q$ is the sum of
the operators $B_{\mathsf w}B_{\mathsf w}^\dagger$, $D$ the dimension of the target, $S$ the total dimension of the
sources, $e$, $\eta$, $\alpha$ the three constants of the lemma of \S6.9, and $u$, $H$ the spread and the height
bound of \S6.10. From Corollary~T2 on, the share of directions lost at the level primes is written $\alpha_{\rm ram}$; it is one part of the constant $\alpha$ of the lemma, which in Corollary~T2 is $\alpha_{\rm ram}+\alpha'$, and it is not the $\alpha=u/(u+t)$ of \S6.10. The target space, called $H$ in \S6.9, is written $\mathcal H$ here. The letters $k,m,t_0,f,r$,
$g_0,g_1$, $q,d$, $b$ have the meaning of \S III.4 and \S III.5: $k$ is the initial rank, $m$ the number of lifts,
$t_0$ the integer parameter written $t$ in \S III.4 and $t_0$ in \S III.5 (the field \texttt{t} of the Lean record of ranks), $f=2n$ the number of weight places, $r=k+m-1$, $g_0,g_1$ are the two level primes, $q,d$ the auxiliary primes, and
$b=\exp(x^\varepsilon)$ with $0<\varepsilon<1/52$. $\|A\|_{\rm HS}^2=\operatorname{tr}(A^\dagger A)$ is the
squared Hilbert--Schmidt norm, and $A^\dagger$ the adjoint. Inner products are conjugate-linear in the
first variable and linear in the second. $V^\vee$ is the space of $\C$-linear forms on a complex space $V$.
The letter $\nu$ denotes a dimension where it
occurs, and $j$, $j'$, $l$, $\iota$ are indices. The exponential function is written $\exp$ in this section,
because $e$ is a constant of the lemma of \S6.9.\par}

{\sloppy \emph{Other letters.} Used in more than one place: $\varphi$ and $F$ (elements of
$\mathcal M_{\mathsf w}$ or $\mathcal M_{\mathsf w,\mathbb L}$, and of $\mathcal E_0$ or $\mathcal E_{\mathbb L}$);
$x_0$ (a threshold for $x$) and $c$, $c'$, $C$, $C'$ (positive constants), which are not the same at each
occurrence; $y$ (a real number in the proof of Lemma~T3, and $x^{\varepsilon/8}$ in \S4); $i$ (an index in
$c_i$, $\xi_i$, $\psi_i$, an element of $\mathcal I$ in \S3, and the imaginary unit in $iF$ in \S1); $v$ (a
finite place in (R1) and (R2), a positive real number in (R3)); $X$ (a tuple of vectors in (R2), the variable
of a polynomial in (R4)); $g$ (a group element in (R2), a prime in (R4)); $M$ (a rank, in (R2) and (R3)).
Used in one place only and explained there: $s$ (the inner mean of \S2, a real number; it is not the map
$\mathsf s$); $\mathcal X$, $\mathcal Z$, $\mathcal Y$, $h$ (in the proofs of Lemma~T1(d) and Corollary~T2);
$\Xi$, $\zeta$, $\zeta'$, $p$ (in the Remarks); $Y$, $\rho$, $\chi$ (in \S3); $\tau$, $K^{(\tau)}$, $K''$, $s_0$
(in (R1)); $E/E^+$, $J$, $Z$, $\Phi$, $\varrho_v$ (in (R2)). The number $e_0'$ of \S5 is the first of the
$e_j'$; it is not the $e_0$ of (S).\par}

{\sloppy \emph{The letters are local.} No letter of this section was changed to agree with the other sections of
Part~III, and several have another meaning there. The clashes most likely to confuse are these. $C_0$ is $3$
here; in \S III.7 it is a constant with $C_0\ge28$. $\mathcal H$ is the target space here; in \S III.7 it is a
height. $\mathsf G$ is the number $\exp(C_{\mathsf G}m^5fx^\varepsilon)$ here and in Part~II, \S6.10; in
\S III.7 it is a group. $\mathcal T(\mathsf w)$ is a level prime here, as $\mathcal T$ is in Part~II, \S6.7; in
\S III.6 $\mathcal T$ is a set. $\varsigma_\iota$ are minors here (\S3); in \S III.4 and \S III.7
$\varsigma$ is a number formed from the ranks. $\Lambda$ is a map here; in \S III.6 the $\Lambda_w$ are
lattices, and elsewhere $\Lambda$ is the von Mangoldt function. $\Xi$ is a random unitary map here; in
\S III.6 it is a set of primes and in \S III.7 an operator. $\Psi$, $\Psi_{\mathsf w}$ are unitary maps and
$\psi_i$ weights here; in \S III.4 $\Psi=\psi+1$ and $\psi$ are ranks, and in \S III.9 $\Psi$ is a character.
$\Delta$ is the step of the weights here (\S5); in \S III.4 it is the width of the weight interval. Also:
$\alpha$ is a constant here and a member of the family in \S III.6 (the word of this section; see \S1);
$\mathfrak B$ and $c_0$ occur in \S III.9, $\varkappa_v$ in \S III.7 and $\mathcal Y$ in \S III.5, with other
meanings.\par}

\subsection*{1.\quad The three statements in the source}

\textbf{The data.} At a scale $x$ the source works with a record (Lean: \texttt{D7}) which contains, among
many other things, the following. They are fields of the record. Their identification with spaces of
automorphic forms is made in the source by further hypotheses, which were not read (see \S6).
\begin{itemize}
\item A finite set $\mathsf W$ of \emph{words}. Each word $\mathsf w$ has a sequence of letters in
$\{0,1\}^{k-2}$, and distinct words have distinct sequences (\texttt{letter}, \texttt{letter\_inj}). A word
fixes, at each of the $f$ weight places, a weight $\xi=(\xi_0>\dots>\xi_{k-1})$ of a unitary group in $k$
variables; see \S5.
\item A complex inner product space $\mathcal H$ of finite dimension $D$ (\texttt{HH}, \texttt{Dall}), and a
real number $D_*>0$ (\texttt{Dstar}). The positivity of $D_*$ is not a field of the record. It follows from
two fields of other records: \texttt{Dstar\_eq} and \texttt{Dstar\_pos}.
\item For each word a complex inner product space $\mathcal C_{\mathsf w}$ of finite dimension
$d_{\mathsf w}$ (\texttt{Cu}, \texttt{duS}), a subspace $\mathcal M_{\mathsf w}\subset\mathcal C_{\mathsf w}$ of
dimension $d^{\rm ram}_{\mathsf w}$ (\texttt{Mram}, \texttt{du}), a linear map
$P_{\mathsf w}:\mathcal C_{\mathsf w}\to\mathcal H$ (\texttt{Pu}), and a number $c_0(\mathsf w)\ge0$
(\texttt{c0}, a product of two volumes divided by a dimension). When $c_0(\mathsf w)>0$ put
\begin{equation}\label{tr:K}
 B_{\mathsf w}=c_0(\mathsf w)^{-1/2}P_{\mathsf w},\qquad
 K_{\mathsf w\mathsf w'}=B_{\mathsf w}^\dagger B_{\mathsf w'}:\ \mathcal C_{\mathsf w'}\to\mathcal C_{\mathsf w}.
\end{equation}
When $c_0(\mathsf w)=0$ the definitions of the source give $K_{\mathsf w\mathsf w}=0$. In the notation of
Part~II, \S6.9 the maps $B_i$ are the $P_{\mathsf w}$, the spaces $V_i$ are the $\mathcal C_{\mathsf w}$, and the
weights are $c_i=1/c_0(\mathsf w)$ (the letter $c_i$ is that of \S6.9; the number $c_0(\mathsf w)$ of the
record is not one of the $c_i$); so $K_i=K_{\mathsf w\mathsf w}$ and
$Q=\sum_{\mathsf w}B_{\mathsf w}B_{\mathsf w}^\dagger$.
\item A complex inner product space $\mathcal E_0$ of dimension $D_{\rm r}$ (\texttt{EE}), and a
\emph{conjugate-linear} map $\mathsf s:\mathcal E_0\to\mathcal H$ which preserves norms (\texttt{sEE}). So
$\mathsf s(cF)=\bar c\,\mathsf s(F)$ for $c\in\C$, and $\mathsf s$ is injective. Its image
$\mathcal E=\mathsf s(\mathcal E_0)$ is a complex subspace of $\mathcal H$ of dimension $D_{\rm r}$. There is a
number $\delta_{\rm p}\le1$ with $D_{\rm r}=(1-\delta_{\rm p})D$ (\texttt{Drows}, \texttt{delp}, \texttt{D\_eq}).
The spaces $\mathcal E_0$ and $\mathcal E$ have the same elements up to $\mathsf s$, with conjugate actions of
the scalars. They are kept apart below.
\item For each word a pairing $\mathfrak B_{\mathsf w}:\mathcal M_{\mathsf w}\times\mathcal E_0\to\C$, linear over
$\C$ in each variable (\texttt{Bu}, \texttt{Bu\_linL}, \texttt{Bu\_linR}). In the source it stands for the
pairing at the end of the chain of $m$ lifts of Part~II, \S6.2, between the space at the start of the chain
and the space at its end (it is tied to the chain by the fields \texttt{perCol} and \texttt{PtilCol}, which
were not read). Both spaces carry structures over a number field $\mathbb L\subset\C$: there are
$\mathbb L$-vector spaces $\mathcal M_{\mathsf w,\mathbb L}$ and $\mathcal E_{\mathbb L}$, with lattices in them,
and maps $\mathcal M_{\mathsf w,\mathbb L}\to\mathcal M_{\mathsf w}$ and $\mathcal E_{\mathbb L}\to\mathcal E_0$,
linear over $\mathbb L$ (\texttt{jN}, \texttt{jEE}), under which an $\mathbb L$-basis of
$\mathcal M_{\mathsf w,\mathbb L}$ is a $\C$-basis of $\mathcal M_{\mathsf w}$ and an $\mathbb L$-basis of
$\mathcal E_{\mathbb L}$ is a $\C$-basis of $\mathcal E_0$ (\texttt{jN\_basis}, \texttt{jEE\_basis}: linear
independence and span). The second map goes to $\mathcal E_0$ and not to $\mathcal E$: composed with
$\mathsf s$ it would be conjugate-semilinear over $\mathbb L$, and so linear over $\mathbb L$ only if
$\mathbb L\subset\mathbb R$. In particular
\begin{equation}\label{tr:dimE}
 \dim_{\mathbb L}\mathcal E_{\mathbb L}=D_{\rm r}
\end{equation}
(\texttt{RRV\_finrank}). There is an $\mathbb L$-bilinear pairing
$\mathfrak B^{\rm alg}_{\mathsf w}:\mathcal M_{\mathsf w,\mathbb L}\times\mathcal E_{\mathbb L}\to\mathbb L$
(\texttt{Bk}). The slopes of Part~II, \S6.8--6.10 are those of these lattices.
\end{itemize}
\emph{Comparison with the letters of \S III.6.} The following was checked in the types of the Lean record.
The field $\mathbb L$ is the coefficient field $L_0$ of \S III.6 (\texttt{LL}, with the embedding
\texttt{embLL}). The words are the members $\alpha$ of \S III.6 (\texttt{Word}). The space
$\mathcal M_{\mathsf w,\mathbb L}$ is the space $U_\alpha$ that carries the glued lattice
$\Lambda^{\rm gl}_\alpha$ (\texttt{Nu}), and $\mathcal E_{\mathbb L}$ is the space $U'$ with the lattice
$\mathcal R$ (\texttt{RR}): the maps \texttt{jN} and \texttt{jEE} have the sources \texttt{(div.Nu u).V} and
\texttt{div.RR.V}. Not checked in the Lean: that $\mathfrak B^{\rm alg}_{\mathsf w}$ (\texttt{Bk}) is the
pairing $\langle\ ,\ \rangle_\alpha$ of \S III.6 (\texttt{perCol}), and that the level prime
$\mathcal T(\mathsf w)$ of \S2 is the prime $\mathrm{Lv}_\alpha$ of \S III.6; both sections describe them
in the same words. The letters of \S III.7 were not compared.
\emph{The index of the family is the word.} It is not a Hecke character: in the source the twists by Hecke
characters form another index (\texttt{Tup}), with one twist attached to each place (\texttt{muAt}). By our
reading it serves to glue the lattices (Part~II, \S6.8). The space $\mathcal E_0$ and the pairings
$\mathfrak B_{\mathsf w}$ do not depend on it: \texttt{Tup} does not occur in the types of \texttt{EE} and
\texttt{Bu}, nor in \texttt{LinkS\_at}, \texttt{Frame\_at} and \texttt{frame\_at\_of}.

\emph{Translates.} There is one record for each translate of Part~II, \S4.2, and each record has its own set
of words (Lean: the definition \texttt{data} in \texttt{QuadraticLogDerivCongr.lean}, whose family of index
sets is \texttt{fun u0 => (Fm.D x (M.pt u0)).what}$_2$\texttt{.Word}; there \texttt{u0} is the translate and \texttt{u}
the word). The weights of the table of Part~II, \S4.2 depend on the translate and not on the word
(\texttt{w u0 j}). The level prime, the height and the rank depend on the translate and on the word
(\texttt{T u0 u}, \texttt{h u0 u}, \texttt{d u0 u}).

\textbf{Statement (S)} (Lean: \texttt{LinkS\_at}, \texttt{LinkS}, with $C_0=3$). \emph{Put
$\mathsf G=\exp(C_{\mathsf G}\,m^5f\,x^{\varepsilon})$, where $C_{\mathsf G}>0$ is a constant of the family. There is
$x_0$ such that for every $x\ge x_0$ and every record of the family at the scale $x$ there are real numbers
$\eta_0,\eta_1$ and $e_0\ge0$ with}
\begin{equation}\label{tr:small}
 \eta_0\le \exp(-x^{\varepsilon/8}),\qquad \eta_1\le \exp(-x^{\varepsilon/8}),\qquad \sqrt{e_0}\le \exp(-x^{\varepsilon/8}),
\end{equation}
\emph{such that, if}
\begin{equation}\label{tr:proviso}
 C_0^{-1}\le d_{\mathsf w}\,\mathsf G^2/D_*\le C_0\qquad\text{for every word }\mathsf w,
\end{equation}
\emph{then}
\begin{align}
 \|K_{\mathsf w\mathsf w'}\|_{\rm HS}^2&\le(1+\eta_0)\,\frac{d_{\mathsf w}d_{\mathsf w'}}{D_*}
   \qquad(\mathsf w\ne\mathsf w'),\label{tr:S1}\\
 \|K_{\mathsf w\mathsf w}-1\|_{\rm HS}^2&\le e_0\,d_{\mathsf w},\label{tr:S2}\\
 D&\le(1+\eta_1)\,D_*.\label{tr:S3}
\end{align}
In the proof in the source the three numbers are
$\eta_0=b^{-6}$, $e_0=2C_0/\mathsf G+2\exp(-kf\log k)$ and $\eta_1=\exp(-c\,rf\log r)$ with a constant $c>0$
(\texttt{linkS\_concl}). So they are of size $\exp(-c'x^{\varepsilon})$, and \eqref{tr:small} is a weaker,
uniform way of recording this. When \eqref{tr:proviso} fails for a record, that proof takes
$\eta_0=\eta_1=e_0=0$, and the implication is empty (\texttt{linkS\_of\_standing\_hypotheses}). The source
proves (S) for families satisfying a further condition on the level primes of the words (\texttt{TuSize}).

\textbf{Statement (I)} (Lean: \texttt{LinkIRR}). \emph{There is $x_0$ such that for $x\ge x_0$ and every record
of the family at the scale $x$}
\begin{equation}\label{tr:I}
 0\le\delta_{\rm p}\le\exp\big(-x^{\varepsilon_z/7}\big),\qquad \varepsilon_z=\varepsilon/J_z,
\end{equation}
\emph{where $J_z\ge1$ is a constant of the family.} So $\mathcal E$ is almost all of $\mathcal H$.

\textbf{Statement (P)} (Lean: \texttt{BesselId}). \emph{For every record: \textup{(P1)} for
$\varphi\in\mathcal M_{\mathsf w,\mathbb L}$ and $F\in\mathcal E_{\mathbb L}$,
$\mathfrak B_{\mathsf w}(\varphi,F)=\mathfrak B^{\rm alg}_{\mathsf w}(\varphi,F)$, the right side being
read in $\C$ through the fixed embedding of $\mathbb L$, and $\varphi$, $F$ on the left being their images in
$\mathcal M_{\mathsf w}$ and $\mathcal E_0$; \textup{(P2)} there are a unitary map $\Psi$ of
$\mathcal E_0$ and, for each word, a complex number $\lambda_{\mathsf w}\ne0$ and a unitary map
$\Psi_{\mathsf w}$ of $\mathcal M_{\mathsf w}$, such that}
\begin{equation}\label{tr:P}
 \mathfrak B_{\mathsf w}(\varphi,F)=\lambda_{\mathsf w}\,\big\langle \mathsf s\Psi F,\ P_{\mathsf w}\Psi_{\mathsf w}\varphi\big\rangle_{\mathcal H}
 \qquad(\varphi\in\mathcal M_{\mathsf w},\ F\in\mathcal E_0).
\end{equation}
Both sides are linear in $\varphi$ and in $F$. For the right side and $F$: replace $F$ by $iF$. Then
$\Psi F$ is multiplied by $i$, since $\Psi$ is linear; $\mathsf s\Psi F$ is multiplied by $-i$, since
$\mathsf s$ is conjugate-linear; and the inner product is multiplied by $i$, since it is conjugate-linear in
its first variable. The map $\mathsf s$ cannot be omitted: for a complex subspace $\mathcal E$ of $\mathcal H$,
a pairing that is linear on $\mathcal M_{\mathsf w}\times\mathcal E$ cannot be equal to
$\lambda_{\mathsf w}\langle\Psi F,P_{\mathsf w}\Psi_{\mathsf w}\varphi\rangle$ with $\Psi$ unitary on $\mathcal E$, unless it is zero: the left side
would be multiplied by $i$ and the right side by $-i$. In the Lean, \texttt{sEE} has the type of a
conjugate-linear isometry from \texttt{EE} to the target, the pairing \texttt{besS} is
\texttt{inner (sEE F) (Pu \dots)}, and \texttt{BesselId} asserts
\texttt{Bu u phi F = a * besS u (U phi) (V F)} with linear isometric equivalences \texttt{V} of \texttt{EE} and
\texttt{U} of the space of $\varphi$. In the source (P) is a clause of the statement \texttt{LinkDIV\_at}. The
name refers to the realisation of $\langle\mathsf sF,P_{\mathsf w}\varphi\rangle$ in the source as a Bessel period
of $F$; in the form \eqref{tr:P} no integral appears.

\subsection*{2.\quad From the three statements to the hypotheses of Part~II, \S6.9}

The lemma of Part~II, \S6.9 has three hypotheses, with constants $e$, $\eta$, $\alpha$: a bound
$\sum_i\|K_i-1\|_{\rm HS}^2\le eS$ for the diagonal blocks; a bound
$\operatorname{tr}(Q^2)\le(1+\eta)(\operatorname{tr}Q)^2/D$ for the operator $Q$; and subspaces of total
codimension at most $\alpha S$. Its conclusion is that the images span all of the target but a part of
dimension at most $(\eta+4\alpha+4\sqrt e)D$. This is the form that the proof in \S6.9 gives. The display of
\S6.9 has $2\eta$ in place of $\eta$, which is weaker. This section uses the form with $\eta$ everywhere, and
so does the source (\texttt{FrameLA1}; see below). The next lemma obtains the first two hypotheses from \eqref{tr:S1}--\eqref{tr:S3}.

\textbf{Lemma T1.} \emph{Let $\mathcal H$ be a complex inner product space of dimension $D\ge1$, let
$\mathsf W$ be a finite nonempty set, and for $\mathsf w\in\mathsf W$ let
$B_{\mathsf w}:\mathcal C_{\mathsf w}\to\mathcal H$ be a linear map from an inner product space of dimension
$d_{\mathsf w}\ge1$. Put $K_{\mathsf w\mathsf w'}=B_{\mathsf w}^\dagger B_{\mathsf w'}$,
$Q=\sum_{\mathsf w}B_{\mathsf w}B_{\mathsf w}^\dagger$ and $S=\sum_{\mathsf w}d_{\mathsf w}$.}
\begin{enumerate}[label=\textup{(\alph*)},leftmargin=2.2em]
\item $\operatorname{tr}Q=\sum_{\mathsf w}\operatorname{tr}K_{\mathsf w\mathsf w}$ \emph{and}
$\operatorname{tr}(Q^2)=\sum_{\mathsf w,\mathsf w'}\|K_{\mathsf w\mathsf w'}\|_{\rm HS}^2$.
\item \emph{Let $D_*>0$, $\eta_0\ge-1$, $e_0\ge0$, and assume \eqref{tr:S1} and \eqref{tr:S2}. Then}
\begin{equation}\label{tr:cor}
 \operatorname{tr}(Q^2)\le(1+\eta_0)\frac{S^2}{D_*}+(1+\sqrt{e_0})^2S,\qquad
 \operatorname{tr}Q\ge(1-\sqrt{e_0})S,\qquad \sum_{\mathsf w}\|K_{\mathsf w\mathsf w}-1\|_{\rm HS}^2\le e_0S .
\end{equation}
\item \emph{Assume in addition \eqref{tr:S3}, and let $0<\varkappa\le\frac14$ be such that}
\begin{equation}\label{tr:kap}
 D_*\le\varkappa S,\qquad \eta_0\le\varkappa/2,\qquad\eta_1\le\varkappa/2,\qquad\sqrt{e_0}\le\varkappa/2 .
\end{equation}
\emph{Then}
\begin{equation}\label{tr:F1}
 \operatorname{tr}(Q^2)\le(1+5\varkappa)\,\frac{(\operatorname{tr}Q)^2}{D},\qquad e_0\le\tfrac1{64}.
\end{equation}
\item \emph{Under the hypotheses of \textup{(c)}, let $\mathcal U_{\mathsf w}\subset\mathcal C_{\mathsf w}$ be
subspaces with $\sum_{\mathsf w}(d_{\mathsf w}-\dim\mathcal U_{\mathsf w})\le\alpha S$, let
$\mathcal E\subset\mathcal H$ be a complex subspace of dimension $(1-\delta_{\mathcal E})D$ with $0\le\delta_{\mathcal E}<1$,
and let $\pi_{\mathcal E}$ be the orthogonal projection onto $\mathcal E$. Then}
\begin{equation}\label{tr:codim}
 \dim\mathcal E-\dim\sum_{\mathsf w}\pi_{\mathcal E} B_{\mathsf w}\,\mathcal U_{\mathsf w}
 \ \le\ \frac{5\varkappa+4\alpha+4\sqrt{e_0}}{1-\delta_{\mathcal E}}\,\dim\mathcal E .
\end{equation}
\end{enumerate}

\emph{Proof.} (a) The first identity is $\operatorname{tr}(BB^\dagger)=\operatorname{tr}(B^\dagger B)$. For the
second, $K_{\mathsf w\mathsf w'}^\dagger=K_{\mathsf w'\mathsf w}$, so
\begin{equation}\label{D:307}
 \operatorname{tr}\big(B_{\mathsf w}B_{\mathsf w}^\dagger B_{\mathsf w'}B_{\mathsf w'}^\dagger\big)
 =\operatorname{tr}\big(B_{\mathsf w'}^\dagger B_{\mathsf w}\,B_{\mathsf w}^\dagger B_{\mathsf w'}\big)
 =\operatorname{tr}\big(K_{\mathsf w\mathsf w'}^\dagger K_{\mathsf w\mathsf w'}\big)
 =\|K_{\mathsf w\mathsf w'}\|_{\rm HS}^2 ,
\end{equation}
and $\operatorname{tr}(Q^2)$ is the sum of these over all pairs.

(b) Let $K$ be an operator on a space of dimension $\nu$. By the Cauchy--Schwarz inequality for the
Hilbert--Schmidt inner product, $|\operatorname{tr}(K-1)|\le\sqrt\nu\,\|K-1\|_{\rm HS}$, so
$\re\operatorname{tr}K\ge\nu-\sqrt\nu\,\|K-1\|_{\rm HS}$; and $\|K\|_{\rm HS}\le\sqrt\nu+\|K-1\|_{\rm HS}$. With
\eqref{tr:S2} this gives $\operatorname{tr}K_{\mathsf w\mathsf w}\ge(1-\sqrt{e_0})d_{\mathsf w}$ (the trace is
real) and $\|K_{\mathsf w\mathsf w}\|_{\rm HS}^2\le(1+\sqrt{e_0})^2d_{\mathsf w}$. Sum over $\mathsf w$ and use
(a). For the pairs $\mathsf w\ne\mathsf w'$ use \eqref{tr:S1} and
$\sum_{\mathsf w\ne\mathsf w'}d_{\mathsf w}d_{\mathsf w'}\le S^2$; here $1+\eta_0\ge0$ is used.

(c) By \eqref{tr:S3}, $1+\eta_1\ge D/D_*>0$ and $1/D\ge1/((1+\varkappa/2)D_*)$. By (b) and \eqref{tr:kap},
$\operatorname{tr}Q\ge(1-\varkappa/2)S\ge0$, and, since $S\le\varkappa S^2/D_*$,
\begin{align}\label{D:308}
 \operatorname{tr}(Q^2)&\le\Big((1+\tfrac\varkappa2)+(1+\tfrac\varkappa2)^2\varkappa\Big)\frac{S^2}{D_*},\\
 (1+5\varkappa)\frac{(\operatorname{tr}Q)^2}{D}&\ge\frac{(1+5\varkappa)(1-\frac\varkappa2)^2}{1+\frac\varkappa2}\cdot\frac{S^2}{D_*}.
\end{align}
So it suffices that $(1+\frac\varkappa2)^2+(1+\frac\varkappa2)^3\varkappa\le(1+5\varkappa)(1-\frac\varkappa2)^2$.
Expanding both sides, the difference of the right side and the left side is
\begin{equation}\label{tr:poly}
 \varkappa\big(2-\tfrac{13}2\varkappa+\tfrac12\varkappa^2-\tfrac18\varkappa^3\big),
\end{equation}
which is positive for $0<\varkappa\le\frac14$, since then the bracket is at least $2-\frac{13}8-\frac1{512}>0$.
Finally $e_0\le\varkappa^2/4\le1/64$.

(d) By (b) and (c) the hypotheses of the lemma of Part~II, \S6.9 hold for the maps $B_{\mathsf w}$, with
all weights $c_i$ equal to $1$, with $e=e_0\le\frac14$, $\eta=5\varkappa$, and the given $\alpha$. So, by the form of
that lemma with $\eta$ which its proof gives, the orthogonal
complement $\mathcal Z$ of $\mathcal X=\sum_{\mathsf w}B_{\mathsf w}\mathcal U_{\mathsf w}$ in $\mathcal H$ has
dimension at most $(5\varkappa+4\alpha+4\sqrt{e_0})D$. Now $\pi_{\mathcal E}(\mathcal X)=(\mathcal X+\mathcal E^\perp)\cap\mathcal E$,
and the codimension of this space in $\mathcal E$ equals the codimension of $\mathcal X+\mathcal E^\perp$ in
$\mathcal H$, because $\mathcal X+\mathcal E^\perp=\pi_{\mathcal E}(\mathcal X)\oplus\mathcal E^\perp$ and
$\mathcal H=\mathcal E\oplus\mathcal E^\perp$. This codimension is at most that of $\mathcal X$, which is
$\dim\mathcal Z$. Finally $D=\dim\mathcal E/(1-\delta_{\mathcal E})$. \qed

\emph{Remarks.} (1) \emph{The blocks are not small, and the images are not nearly orthogonal.} For every
family, $\operatorname{tr}(Q^2)\ge(\operatorname{tr}Q)^2/D$, by the Cauchy--Schwarz inequality for the
eigenvalues of $Q$. If all $B_{\mathsf w}$ are isometries, then $K_{\mathsf w\mathsf w}=1$, and this and (a) give
\begin{equation}\label{D:310}
 \sum_{\mathsf w\ne\mathsf w'}\|K_{\mathsf w\mathsf w'}\|_{\rm HS}^2\ \ge\ \frac{S^2}D-S .
\end{equation}
The sum of the bounds $d_{\mathsf w}d_{\mathsf w'}/D$ over the same pairs is $(S^2-\sum_{\mathsf w}d_{\mathsf w}^2)/D$,
which exceeds the right side by at most $S$, that is, by the relative amount $D/S$ at most. By \eqref{tr:S3}, $1/D_*\le(1+\eta_1)/D$. So \eqref{tr:S1} and \eqref{tr:S3} assert that the squared norm of
every off-diagonal block is at most $(1+\eta_0)(1+\eta_1)$ times $d_{\mathsf w}d_{\mathsf w'}/D$. Up to this
factor and up to the relative error $D/S$, that is the least size that the blocks can have on average.
Statement (S) bounds $D$ by $D_*$ from above only; the bound in the other direction is in (R3) of \S4 and is
not used in this section. The number $d_{\mathsf w}d_{\mathsf w'}/D$ is also the mean size for subspaces in
general position: if $\Xi$ is a random unitary map of $\mathcal H$ for Haar measure and $\zeta,\zeta'$ are unit
vectors, the mean of $|\langle\zeta,\Xi\zeta'\rangle|^2$ is $1/D$, so the mean of $\|B^\dagger\Xi B'\|_{\rm HS}^2$ for
isometries $B,B'$ of ranks $\nu,\nu'$ is $\nu\nu'/D$. Orthogonal images would mean
$K_{\mathsf w\mathsf w'}=0$. Mutually orthogonal images have dimensions that add up to at most $D$. For
injective maps this is impossible when $S>D$. The maps of the source are not known to be injective, but
\eqref{tr:S2} bounds their ranks from below. The operator $K_{\mathsf w\mathsf w}-1$ is $-1$ on the kernel of
$B_{\mathsf w}$, so $\|K_{\mathsf w\mathsf w}-1\|_{\rm HS}^2\ge d_{\mathsf w}-\operatorname{rank}B_{\mathsf w}$, and
\eqref{tr:S2} gives $\operatorname{rank}B_{\mathsf w}\ge(1-e_0)\,d_{\mathsf w}$. Under the hypotheses of
Corollary~T2 below, $e_0\le\frac1{64}$, $S\ge x^{C_H}D_*\ge4D_*$ and $D\le\frac98D_*$. So the ranks add up to
at least $\frac{63}{64}S>D$, which is more than the dimension of the target, and the images cannot all be
mutually orthogonal. The Gram block estimates control their total squared overlap
and put $Q$ close to its mean eigenvalue in Hilbert--Schmidt norm. Precisely, for a Hermitian operator $Q$ on
$\mathcal H$ one has $\operatorname{tr}(Q^2)-(\operatorname{tr}Q)^2/D=\|Q-(\operatorname{tr}Q/D)\cdot1\|_{\rm HS}^2$,
so under the hypotheses of Lemma~T1(c) the bound \eqref{tr:F1} is
\begin{equation}\label{tr:QHS}
 \Big\|Q-\frac{\operatorname{tr}Q}D\cdot1\Big\|_{\rm HS}^2\ \le\ 5\varkappa\,\frac{(\operatorname{tr}Q)^2}D .
\end{equation}
The operator $Q=\sum_{\mathsf w}B_{\mathsf w}B_{\mathsf w}^\dagger$, with
$B_{\mathsf w}=c_0(\mathsf w)^{-1/2}P_{\mathsf w}$ for the maps of the source, is not the sum of the projections onto the images:
$B_{\mathsf w}B_{\mathsf w}^\dagger$ is a projection only if $B_{\mathsf w}$ is a partial isometry.
(2) \emph{The constant $5$.} The term $(1+\sqrt{e_0})^2S$ in \eqref{tr:cor} comes from the diagonal blocks. It
is the reason why \eqref{tr:F1} needs $D_*\le\varkappa S$: the total dimension $S$ of the sources must exceed
the dimension of the target by the factor $1/\varkappa$. A term of this size is present on average for maps in
general position: for isometries $B_{\mathsf w}$ composed with independent random unitary maps of
$\mathcal H$, the mean of $\operatorname{tr}(Q^2)$ is $S+(S^2-\sum_{\mathsf w}d_{\mathsf w}^2)/D$ by Remark~(1), while
$\operatorname{tr}Q=S$; so the mean of $\operatorname{tr}(Q^2)D/(\operatorname{tr}Q)^2-1$ is about $D/S$. This is
a statement about a mean. Lemma~T1(c) is an upper bound, and \eqref{tr:S1}--\eqref{tr:S3} do not bound
$\operatorname{tr}(Q^2)D/(\operatorname{tr}Q)^2-1$ from below. For example, if every $B_{\mathsf w}$ is a unitary
map of $\mathcal H$, then \eqref{tr:S1}--\eqref{tr:S3} hold with $D_*=D$ and $\eta_0=e_0=\eta_1=0$, and
$Q=|\mathsf W|\cdot1$, so this number is $0$, while $D_*/S=1/|\mathsf W|$. (Another example, with $d_{\mathsf w}=1$: the $p(p+1)$ lines spanned by the vectors of $p+1$ mutually unbiased
bases of $\C^p$, where such bases exist. Here $Q=(p+1)\cdot1$ and $D_*/S=1/(p+1)$.)
(3) The archive contains statements with these names. \texttt{CorBounds} is for (b).
\texttt{TraceSquareAndDiagBounds} is for (c), with $e_0\le\frac14$ where (c) has $\frac1{64}$.
\texttt{CompressTraceSquareBound} is a variant of (c), not of (d): the maps are first restricted to
subspaces of total codimension at most $\alpha S$, and the conclusion is the second-moment bound for the
restricted maps with $\eta=5\varkappa+3\sqrt\alpha$ (for $\varkappa\le\frac1{20}$, $0\le\alpha\le\frac14$). The
bound for the diagonal blocks of the restricted maps, with $e=2e_0$, is the conclusion of another statement,
\texttt{CompressDiagDefectBound} (for $\alpha\le\frac12$). The proofs above
are ours. The Lean proofs of these statements were not read. By the dependency computation the final theorem depends on
\texttt{CompressTraceSquareBound} and not on \texttt{CorBounds} or \texttt{TraceSquareAndDiagBounds}. The name
\texttt{CompressDiagDefectBound} was not among those queried.

\begin{altbox}
Corollary~T2, with Proposition~M4 of \S\ref{as:sec}, is this paper's route to the mean-slope inequality
(H8). The source does not state Corollary~T2; it proves (H8) from a theorem on lattices of its own
(\texttt{TheoremF}), with three further conditions (\S\ref{as:sec}, \S4).
\end{altbox}
\textbf{Corollary T2} (the hypotheses of Part~II, \S6.9 for the maps of the source, and for arbitrary
subspaces). \emph{Assume \textup{(S)} with \eqref{tr:proviso}, \textup{(I)} with $\delta_{\rm p}<1$, and
\textup{(P)}, for one record. Let $C_H>0$ and put $\varkappa=x^{-C_H}$. Let $\alpha_{\rm ram},\alpha'\ge0$, and for each
word let $\mathcal U_{\mathsf w}\subset\mathcal M_{\mathsf w}$ be a complex subspace. Assume further}
\begin{equation}\label{tr:side}
 |\mathsf W|\ \ge\ C_0\,\mathsf G^2x^{C_H},\qquad
 d_{\mathsf w}-d^{\rm ram}_{\mathsf w}\le\alpha_{\rm ram}\,d_{\mathsf w}\ \ (\mathsf w\in\mathsf W),\qquad
 \varkappa\le\tfrac14,\qquad \exp(-x^{\varepsilon/8})\le\varkappa/2 ,
\end{equation}
\begin{equation}\label{tr:sub}
 \sum_{\mathsf w}\big(d^{\rm ram}_{\mathsf w}-\dim\mathcal U_{\mathsf w}\big)\ \le\ \alpha'\sum_{\mathsf w}d_{\mathsf w}.
\end{equation}
\emph{For each word let $\beta_{\mathsf w}:\mathcal M_{\mathsf w}\to\operatorname{Hom}_{\mathbb L}(\mathcal E_{\mathbb L},\mathbb L)\otimes_{\mathbb L}\C$
be the complexification of the map $\varphi\mapsto\mathfrak B^{\rm alg}_{\mathsf w}(\varphi,\cdot)$ on
$\mathcal M_{\mathsf w,\mathbb L}$; here $\mathcal M_{\mathsf w,\mathbb L}\otimes_{\mathbb L}\C$ is identified with
$\mathcal M_{\mathsf w}$ as in \S1. The target is identified with the dual $\mathcal E_0^\vee$ of
$\mathcal E_0$ by the map $\mathcal E_{\mathbb L}\to\mathcal E_0$ of \S1. Then the sum of the spaces
$\beta_{\mathsf w}(\mathcal U_{\mathsf w})$ has codimension at most}
\begin{equation}\label{tr:final}
 \frac{5x^{-C_H}+4(\alpha_{\rm ram}+\alpha')+4\sqrt{e_0}}{1-\delta_{\rm p}}\ D_{\rm r}
\end{equation}
\emph{in $\mathcal E_0^\vee$, which has dimension $D_{\rm r}$ by \eqref{tr:dimE}. In particular, for
$\mathcal U_{\mathsf w}=\mathcal M_{\mathsf w}$ and $\alpha'=0$, the sum of the images of the $\beta_{\mathsf w}$ has
codimension at most $(5x^{-C_H}+4\alpha_{\rm ram}+4\sqrt{e_0})D_{\rm r}/(1-\delta_{\rm p})$.}

\emph{Proof.} First, $c_0(\mathsf w)>0$ for every word, so that $B_{\mathsf w}$ is defined. Indeed
$d_{\mathsf w}\ge1$ by the lower bound in \eqref{tr:proviso}, and $e_0<1$ by \eqref{tr:small}; if
$c_0(\mathsf w)=0$, then $K_{\mathsf w\mathsf w}=0$ and \eqref{tr:S2} would read $d_{\mathsf w}\le e_0d_{\mathsf w}$.
(This is \texttt{unfolding\_constant\_pos\_of\_diagonal} in the source.) By the lower bound in
\eqref{tr:proviso}, $d_{\mathsf w}\ge D_*/(C_0\mathsf G^2)$, so
$S\ge|\mathsf W|D_*/(C_0\mathsf G^2)\ge D_*/\varkappa$. With \eqref{tr:small} and the last condition of
\eqref{tr:side}, \eqref{tr:kap} holds; and $D\ge1$, because \eqref{tr:S2} with $e_0<1$ gives $B_{\mathsf w}\ne0$. (If $\eta_0<-1$, \eqref{tr:S1} forces $|\mathsf W|=1$, which the first
condition of \eqref{tr:side} excludes.) By the second condition of \eqref{tr:side} and by \eqref{tr:sub},
$\sum_{\mathsf w}(d_{\mathsf w}-\dim\mathcal U_{\mathsf w})\le(\alpha_{\rm ram}+\alpha')S$.

\emph{The estimate in $\mathcal H$.} The space $\mathcal E=\mathsf s(\mathcal E_0)$ is a complex subspace of
$\mathcal H$, because $c\,\mathsf s(F)=\mathsf s(\bar cF)$. Its dimension is $D_{\rm r}=(1-\delta_{\rm p})D$,
because $\mathsf s$ is injective and maps a basis of $\mathcal E_0$ to a basis of $\mathcal E$. Let $\pi_{\mathcal E}$ be
the orthogonal projection of $\mathcal H$ onto $\mathcal E$. Apply Lemma~T1(d) with the subspaces
$\Psi_{\mathsf w}\mathcal U_{\mathsf w}\subset\mathcal M_{\mathsf w}\subset\mathcal C_{\mathsf w}$, which have the
dimensions of the $\mathcal U_{\mathsf w}$, with $\alpha_{\rm ram}+\alpha'$ in place of $\alpha$, and with
$\delta_{\mathcal E}=\delta_{\rm p}$. Since $P_{\mathsf w}$ and $B_{\mathsf w}$ differ by a positive factor, the
space $\mathcal Y=\sum_{\mathsf w}\pi_{\mathcal E} P_{\mathsf w}\Psi_{\mathsf w}\mathcal U_{\mathsf w}$ has codimension at most
\eqref{tr:final} in $\mathcal E$. This step uses only the subspace $\mathcal E$ of $\mathcal H$, and neither
$\mathsf s$ nor the scalars of $\mathcal E_0$.

\emph{The passage to the dual.} Let $\Lambda:\mathcal H\to\mathcal E_0^\vee$ send $h$ to the form
\begin{equation}\label{D:315}
 \Lambda(h):\ F\longmapsto\langle\mathsf s\Psi F,\,h\rangle_{\mathcal H}\qquad(F\in\mathcal E_0).
\end{equation}
The form $\Lambda(h)$ is $\C$-linear in $F$, by the computation after \eqref{tr:P}; so it lies in
$\mathcal E_0^\vee$. The map $\Lambda$ is $\C$-linear in $h$. Its kernel is the set of $h$ orthogonal to
$\mathsf s\Psi(\mathcal E_0)=\mathsf s(\mathcal E_0)=\mathcal E$, that is, $\mathcal E^\perp$. So
$\Lambda=\Lambda\circ\pi_{\mathcal E}$, the restriction of $\Lambda$ to $\mathcal E$ is injective, and since
$\dim\mathcal E=D_{\rm r}=\dim\mathcal E_0^\vee$ it maps $\mathcal E$ isomorphically onto $\mathcal E_0^\vee$.
Hence $\Lambda$ maps $\mathcal H$ onto $\mathcal E_0^\vee$, and $\Lambda(\mathcal Y)$ has in $\mathcal E_0^\vee$
the codimension of $\mathcal Y$ in $\mathcal E$. By \eqref{tr:P}, the map
$\varphi\mapsto\mathfrak B_{\mathsf w}(\varphi,\cdot)$ from $\mathcal M_{\mathsf w}$ to $\mathcal E_0^\vee$ is
$\lambda_{\mathsf w}\,\Lambda P_{\mathsf w}\Psi_{\mathsf w}$, a $\C$-linear map. The image of $\mathcal U_{\mathsf w}$
under it is $\Lambda\pi_{\mathcal E} P_{\mathsf w}\Psi_{\mathsf w}\mathcal U_{\mathsf w}$, because $\lambda_{\mathsf w}\ne0$ and the
image is a complex subspace. So the sum of these images is $\Lambda(\mathcal Y)$.

\emph{The algebraic model.} The map $\mathcal E_{\mathbb L}\to\mathcal E_0$ is linear over $\mathbb L$ and maps
an $\mathbb L$-basis to a $\C$-basis. So restriction of forms along it is a $\C$-linear isomorphism of
$\mathcal E_0^\vee$ with $\operatorname{Hom}_{\mathbb L}(\mathcal E_{\mathbb L},\mathbb L)\otimes_{\mathbb L}\C$.
Under it and under $\mathcal M_{\mathsf w,\mathbb L}\otimes_{\mathbb L}\C=\mathcal M_{\mathsf w}$, the map
$\varphi\mapsto\mathfrak B_{\mathsf w}(\varphi,\cdot)$ is $\beta_{\mathsf w}$: both are $\C$-linear, and by (P1)
they agree on the image of $\mathcal M_{\mathsf w,\mathbb L}$, which spans $\mathcal M_{\mathsf w}$. \qed

\emph{On the two scalar actions.} Every space in this proof is a complex space, and every map in it is
$\C$-linear, except $\mathsf s$. The conjugate-linear map $\mathsf s$ occurs only inside the inner product,
in its conjugate-linear variable. The constants are those of Lemma~T1(d).

The corollary is stated for arbitrary subspaces because Part~II, \S6.10 needs it so: there the rank bound is
applied after discarding directions, for every level of the slope filtration. The subspaces that occur are
complexifications of $\mathbb L$-subspaces of $\mathcal M_{\mathsf w,\mathbb L}$.

\emph{The side conditions in the source.} For the second condition of \eqref{tr:side}: each word has a level
prime $\mathcal T(\mathsf w)\ge b^{m/2}$, and
$d^{\rm ram}_{\mathsf w}=d_{\mathsf w}(1-1/(\mathcal T(\mathsf w)-1))^f$ (\texttt{RamRank\_at}); by Bernoulli's
inequality this gives $\alpha_{\rm ram}=2f/b$ for $b\ge3$ (\texttt{ramShare\_at\_of}). The number $\mathcal T(\mathsf w)$ is
the level $\mathcal T$ of Part~II, \S6.7, which is said there to be attached to one choice of lattice data. In
the source it depends on the translate and on the word (\texttt{T u0 u}). The first condition of
\eqref{tr:side} is \texttt{ChCode\_at}. In the source the number of words is fixed at that size: with $C_0=3$,
\begin{equation}\label{tr:words}
 |\mathsf W|=\big\lceil 3\,\mathsf G^2x^{C_H}\big\rceil,\qquad\text{so that, under \eqref{tr:proviso},}\qquad
 x^{C_H}\ \le\ \frac S{D_*}\ \le\ 9x^{C_H}+3
\end{equation}
(the equality is a clause in \texttt{DesignS.lean}, where the exponent is the one passed to
\texttt{ChCode\_at}; the two bounds are ours, from $D_*/(3\mathsf G^2)\le d_{\mathsf w}\le3D_*/\mathsf G^2$).
Since the words have $k-2$ letters, this requires $2^{k-2}\ge|\mathsf W|$. The source has
$\log|\mathsf W|\le(k-3)\log2$ (a clause of \texttt{Win\_at}) and the size condition
$8C_{\mathsf G}\,m^5f\,x^\varepsilon\le k$ (a clause of \texttt{SizeKRel}), with
$m\approx x^{1/2-\varepsilon/16}$, $n\approx x^{2+8\varepsilon}$ and $k\approx x^{9/2+10\varepsilon}$, each up to a
factor $x^{\pm\varepsilon/100}$ (\texttt{SizeM}, \texttt{SizeN}, \texttt{SizeK}). The rounded sizes
$m\approx\sqrt x$, $n\approx x^2$, $k\approx x^{9/2}$ of \S I.8 do not satisfy this condition: they give
$m^5f\,x^\varepsilon=2x^{9/2+\varepsilon}$. By the dependency computation the final theorem depends on
\texttt{RamRank\_at}, \texttt{ramShare\_at\_of} and \texttt{ChCode\_at}. The names \texttt{Win\_at},
\texttt{SizeKRel}, \texttt{SizeM}, \texttt{SizeN} and \texttt{SizeK} were not among those queried.

\textbf{How the source uses the corollary.} The source does not state Corollary~T2. It applies a theorem on
lattices (\texttt{TheoremF}). By the names of the declarations, its
proof uses the lemma of Part~II, \S6.9 in the form \texttt{FrameLA1}, which is stated for arbitrary
subspaces and has $\eta$ and not $2\eta$ (the
chain of names from \texttt{frame\_at\_of} to \texttt{frame\_linear\_algebra\_first} was followed; the proof
terms were not read). The conclusion is the mean-slope inequality of Part~II, \S6.10, for each translate,
with the error term
\begin{equation}\label{tr:frame}
 \frac1{1-\delta_{\rm p}}\Big(4u\log\big(1+\tfrac Hu\big)+(\eta+4\sqrt e)H\Big),\qquad
 \eta=5x^{-C_H}+3\sqrt{\alpha_{\rm ram}},\quad e=2e_0,
\end{equation}
and with $-\log|\mathsf W|$ as a separate term of the lower bound, outside the bracket and not multiplied
by $(1-\delta_{\rm p})^{-1}$ (\texttt{Frame\_at}, \texttt{frame\_at\_of}). Here $u$ and $H$ are the
spread and the height bound of Part~II, \S6.10, and $\alpha_{\rm ram}$ is any number with $\alpha_{\rm ram}\le\frac14$ for which the
second condition of \eqref{tr:side} holds (\texttt{de} in the Lean); the statement also requires $x^{-C_H}\le\frac1{20}$. The values of $\eta$ and $e$ in
\eqref{tr:frame} are the ones that the proof of \texttt{frame\_at\_of} supplies. The source then takes
$\alpha_{\rm ram}=e^{-2x^{\varepsilon/8}}$, so that $\sqrt{\alpha_{\rm ram}}=\exp(-x^{\varepsilon/8})$, and uses
$\sqrt{e_0}\le \exp(-x^{\varepsilon/8})$ from \eqref{tr:small} (\texttt{matRealF\_of}). For all sufficiently
large $x$, $(3+4\sqrt2)\,\exp(-x^{\varepsilon/8})x^{C_H}\le1$. Since $\eta=5x^{-C_H}+3\sqrt{\alpha_{\rm ram}}$ and $e=2e_0$,
this gives
\begin{equation}\label{tr:six}
 (\eta+4\sqrt e)\,x^{C_H}\le6,\qquad\text{so that}\qquad(\eta+4\sqrt e)H\le6b\quad\text{for }H=x^{C_H}b
\end{equation}
(\texttt{FrameF}, \texttt{frameF\_of}). Indeed $(\eta+4\sqrt e)x^{C_H}=5+(3\sqrt{\alpha_{\rm ram}}+4\sqrt{2e_0})x^{C_H}$, and
the bracket is at most $(3+4\sqrt2)\exp(-x^{\varepsilon/8})$. The weaker bound
$\exp(-x^{\varepsilon/8})\le x^{-C_H}/2$, which is the last condition of \eqref{tr:side}, would not suffice: it
leaves $(3+4\sqrt2)/2>1$ for the terms beyond $5$. In the Lean the hypothesis of \texttt{frameF\_of} is
$(3\sqrt{\texttt{de}}+4\sqrt{\texttt{eB}})\,x^{C_H}\le1$ with $0\le\texttt{de}\le\frac14$ and
$0\le\texttt{eB}\le\frac14$ (\texttt{FrameRealF}), where \texttt{de} is $\alpha_{\rm ram}$ and \texttt{eB} is an upper
bound for $e$; \texttt{matRealF\_of} proves it for $x\ge18+|a_0|+|a_1|$. Here $a_0$ is the threshold in its
hypothesis (\texttt{MatSF}), and $a_1$ is a number such that $x^{C_H+1}\le\exp(x^{\varepsilon/8})$ for all
$x\ge\max(a_1,1)$. For such $x$ the proof deduces $\exp(-x^{\varepsilon/8})x^{C_H}\le\frac1{18}$, and it
uses the cruder constant $11$ in place of $3+4\sqrt2$. So the threshold is not $18$; it depends on
$\varepsilon$ and $C_H$. The exponent $C_H$ is the one in the bound
$H\le x^{C_H}b$ of the source (a clause of \texttt{LinkDIV\_at}).

\emph{What is known about the use of these declarations.} \texttt{blockReal\_at\_of} is in the dependency lists
as a declaration that the final theorem depends on. \texttt{frame\_at\_of} and \texttt{matRealF\_of} were first reached by reading proof terms from \texttt{blockReal\_at\_of}, and a later dependency computation confirms that the final theorem depends on both. The chain of proof terms is this. \texttt{blockReal\_at\_of} is used only in the proof of
\texttt{matKF\_of}. \texttt{matKF\_of} is used in four places. In each, its value is an argument of
\texttt{frameRealF\_of''}, and the value of \texttt{frameRealF\_of''} is an argument of \texttt{frameF\_of}.
\texttt{frameRealF\_of''} uses \texttt{matRealF\_of} through \texttt{frameRealF\_of'}, and
\texttt{frameF\_of} uses \texttt{frame\_at\_of}. The same computation shows that the final theorem depends on \texttt{LinkIRR}, \texttt{TheoremF}, \texttt{FrameLA1},
\texttt{frame\_linear\_algebra\_first}, \texttt{Frame\_at}, \texttt{FrameF}, \texttt{matSF\_of},
\texttt{ChCode\_at} and \texttt{CodeUpF}, and does not depend on \texttt{MeanBoundSecondForm}. The statement \texttt{FrameLA} and its
proof \texttt{frame\_linear\_algebra\_bound} are in the lists as not used by the final theorem, as Part~II,
\S6.9 notes. The final theorem also depends on \texttt{matKF\_of}, on
\texttt{unfolding\_constant\_pos\_of\_diagonal} and on \texttt{TheoremFSharp}.

\textbf{The error term and the two means.} Part~II, \S6.9--6.10 states these points; here they are given
with the names of the source. \emph{The error.} For each translate, the mean-slope inequality of the source
has the form
\begin{equation}\label{tr:err}
 (\text{mean slope of the target})\ \ge\ s-\log|\mathsf W|
 -\frac1{1-\delta_{\rm p}}\Big(4u\log\big(1+\tfrac Hu\big)+(\eta+4\sqrt e)H\Big),
\end{equation}
with $\eta=5x^{-C_H}+3\sqrt{\alpha_{\rm ram}}$ and $e=2e_0$ as in \eqref{tr:frame}. Here $s$ is the mean over the words,
weighted by the ranks, of the slopes of the sources minus their heights (\texttt{sFrame}), and the target is
the dual lattice (\texttt{muDual}); this is the statement \texttt{Frame\_at}. The factor $(1-\delta_{\rm p})^{-1}$
multiplies the error bracket only. It is there because the lattice of the target spans a space of dimension
$D_{\rm r}$, not $D$, and the second moment is taken in $\mathcal H$; this is where (I) is used, as in
Lemma~T1(d). The term $-\log|\mathsf W|$ is separate. It is at most $(k-3)\log2$ in absolute value
(\texttt{Win\_at}), which is negligible against $b$. The archive has a statement with the same hypotheses and
no such term (\texttt{TheoremFSharp}). In $\eta$ only $\sqrt{\alpha_{\rm ram}}$ is exponentially small. The power of $x$
comes from the number of words \eqref{tr:words}: by Lemma~T1(c) a bound $5\varkappa$ for $\eta$ needs
$D_*\le\varkappa S$, and $D_*/S\ge1/(9x^{C_H}+3)$. So no bound for $\eta$ better than a negative power of $x$
follows from \textup{(S)} by Lemma~T1. This is enough, by \eqref{tr:six}. We do not assert that the number
$\operatorname{tr}(Q^2)D/(\operatorname{tr}Q)^2-1$ of the family is of that size: it was not bounded from
below (Remark~(2)). \emph{The two means.} Each translate of Part~II, \S4.2 has its own record, with its own
set of words (\S1). The mean over the words is an inner mean, weighted by the ranks, inside \eqref{tr:err}
for one translate: the number $s$ is, in the file \texttt{MeanBoundSecondForm.lean} (source [7] of Part~II),
\begin{center}\texttt{def s (u0 : U0) : }$\mathbb R$\texttt{ := wmean (D.d u0) (fun u => D.muN u0 u - D.h u0 u)}.\end{center} That file defines the
record and its hypotheses, among them the slope inequality \eqref{tr:err} for each translate
(\texttt{Data.SlopeIneq}; the final theorem depends on it). It also states a bound for the mean over the
translates, with an error \texttt{ErrX} that contains the mean of $\log|\mathsf W|$
(\texttt{MeanBoundSecondForm}). The final theorem does not depend on that statement. The route of the final
theorem keeps the translate fixed: \texttt{each\_unit\_bound\_statement}, in the file
\texttt{LevelOfCongruence.lean} (source [28] of Part~II), has the conclusion
\begin{center}\texttt{D.ASK u0 $\le$ D.Main u0 + ErrX5at D C Xpp u0}\quad for every \texttt{u0},\end{center}
and $\log|\mathsf W|$, for the set of words of the translate \texttt{u0}, is a term of \texttt{ErrX5at}. The
final theorem depends on \texttt{each\_unit\_bound\_statement}; its two uses in the archive, in the files
\texttt{ChTop.lean} (source [5]) and \texttt{PChain.lean}, are again for one translate. So the bound of the source holds for each
translate, and the mean over the translates in (M) is an outer mean, taken afterwards.
The inner mean is not visible in (M). The ``two moments'' of Part~II, \S6.9--6.10 are what the lemma there needs.
What the source assumes are the bounds \eqref{tr:S1}--\eqref{tr:S3} for the Gram blocks; the two moments
follow from them by Lemma~T1(b),(c).

\subsection*{3.\quad The mechanism in a finite model}

\begin{illbox}
This subsection, with Lemma~T3, is a model. It is not a step of the source's proof and is not used in the
proof.
\end{illbox}

Statement (S) is proved in the source by expanding each block $K_{\mathsf w\mathsf w'}$ as a sum of terms and
computing its squared norm as a sum over pairs of terms (see \S4). The simplest instance of this kind of
computation is the following, for theta lifts between two definite groups as in Part~II, \S6.1--6.2. It is
a statement about finite sets.

\textbf{The setting.} Let $\mathcal I$ and $\mathcal I'$ be nonempty finite sets with positive weights $\omega$ and
$\omega'$ (classes of lattices $L_i$ and $L'_{i'}$, with $\omega(i)=1/|\mathrm{Aut}\,L_i|$). For $i\in\mathcal I$
let $G_i$ be a finite group of order $1/\omega(i)$. For $i\in\mathcal I$, $i'\in\mathcal I'$ let
$\mathrm{Rep}(i,i')$ be a finite set on which $G_i$ acts freely (the isometric
embeddings $\rho:L_i\to L'_{i'}$, with $\sigma\in\mathrm{Aut}\,L_i$ acting by $\rho\mapsto\rho\sigma$). Let $A$ be a
finite abelian group and $\mathrm{inv}:\mathrm{Rep}(i,i')\to A$ maps with
$\mathrm{inv}(\rho\sigma)=\mathrm{inv}(\rho)$ for $\sigma\in G_i$. For a character $\chi$ of $A$ put
\begin{equation}\label{tr:theta}
 \Theta_\chi(i,i')=\sum_{\rho\in\mathrm{Rep}(i,i')}\chi(\mathrm{inv}(\rho)),\qquad
 (B_\chi f)(i')=\sum_{i\in\mathcal I}\omega(i)\,\Theta_\chi(i,i')\,f(i).
\end{equation}
So $\Theta_1(i,i')$ is the number of representations of $L_i$ by $L'_{i'}$, and $\Theta_\chi$ is a
representation number twisted by a character. The map $B_\chi$ goes from $L^2(\mathcal I,\omega)$ to
$L^2(\mathcal I',\omega')$, where $\langle f,g\rangle=\sum_i\omega(i)\overline{f(i)}g(i)$, conjugate-linear in the first variable as
everywhere in this section. An operator $Y$ on
$L^2(\mathcal I,\omega)$ has a kernel, also written $Y$, with
$(Yf)(i_1)=\sum_{i_2}\omega(i_2)Y(i_1,i_2)f(i_2)$; the kernel of the identity is $1/\omega(i_1)$ for $i_1=i_2$
and $0$ otherwise, and
$\|Y\|_{\rm HS}^2=\sum_{i_1,i_2}\omega(i_1)\omega(i_2)|Y(i_1,i_2)|^2$. For $i_1,i_2\in\mathcal I$ put
\begin{equation}\label{tr:N}
 \mathcal N(i_1,i_2)=\sum_{i'\in\mathcal I'}\omega'(i')\,\#\big\{(\rho_1,\rho_2)\in\mathrm{Rep}(i_1,i')\times\mathrm{Rep}(i_2,i'):\ \mathrm{inv}(\rho_1)=\mathrm{inv}(\rho_2)\big\},
\end{equation}
and let $\mathcal N^{\circ}(i_1,i_2)$ be the same sum restricted to the pairs that are \emph{not} of the form
$(\rho,\rho\sigma)$ with $i_1=i_2$ and $\sigma\in G_{i_1}$.

\textbf{Lemma T3.} \emph{\textup{(a)} The operator $\sum_\chi B_\chi^\dagger B_\chi$ on $L^2(\mathcal I,\omega)$, the
sum being over all characters of $A$, has the kernel $|A|\,\mathcal N$.}

\emph{\textup{(b)} Assume that there is a number $\mathfrak m$ with}
\begin{equation}\label{tr:mass}
 \sum_{i'\in\mathcal I'}\omega'(i')\,|\mathrm{Rep}(i,i')|=\mathfrak m\qquad\text{for every }i\in\mathcal I.
\end{equation}
\emph{Then $\sum_\chi B_\chi^\dagger B_\chi=|A|\,(\mathfrak m\cdot1+\mathfrak N^{\circ})$, where $\mathfrak N^{\circ}$ is the
operator with kernel $\mathcal N^\circ\ge0$, and}
\begin{equation}\label{tr:HS}
 \Big\|\sum_\chi B_\chi^\dagger B_\chi-|A|\,\mathfrak m\cdot1\Big\|_{\rm HS}^2=|A|^2\sum_{i_1,i_2\in\mathcal I}\omega(i_1)\omega(i_2)\,\mathcal N^{\circ}(i_1,i_2)^2 .
\end{equation}
\emph{\textup{(c)} Under \eqref{tr:mass} and $\mathfrak m>0$, the family of the adjoint maps
$B_\chi^\dagger:L^2(\mathcal I',\omega')\to L^2(\mathcal I,\omega)$ satisfies the second-moment hypothesis of
Part~II, \S6.9: with $Q=\sum_\chi B_\chi^\dagger B_\chi$,}
\begin{equation}\label{tr:eta}
 \operatorname{tr}(Q^2)\le(1+\eta)\frac{(\operatorname{tr}Q)^2}{|\mathcal I|},\qquad
 \eta=\frac1{\mathfrak m^2|\mathcal I|}\sum_{i_1,i_2}\omega(i_1)\omega(i_2)\,\mathcal N^{\circ}(i_1,i_2)^2 .
\end{equation}

\emph{Proof.} (a) The adjoint of $B_\chi$ is
$(B_\chi^\dagger g)(i)=\sum_{i'}\omega'(i')\overline{\Theta_\chi(i,i')}g(i')$, as one sees by writing out
$\langle B_\chi f,g\rangle=\langle f,B_\chi^\dagger g\rangle$. So $B_\chi^\dagger B_\chi$ has the kernel
$\sum_{i'}\omega'(i')\overline{\Theta_\chi(i_1,i')}\Theta_\chi(i_2,i')$. Insert \eqref{tr:theta} and sum over $\chi$:
by the orthogonality of the characters of $A$,
$\sum_\chi\overline{\chi(\mathrm{inv}(\rho_1))}\chi(\mathrm{inv}(\rho_2))$ is $|A|$ if
$\mathrm{inv}(\rho_1)=\mathrm{inv}(\rho_2)$ and $0$ otherwise. This gives $|A|\,\mathcal N(i_1,i_2)$.

(b) Let $i_1=i_2=i$. The pairs $(\rho,\rho\sigma)$ with $\rho\in\mathrm{Rep}(i,i')$ and $\sigma\in G_i$ are
distinct for distinct $(\rho,\sigma)$, because the action is free; they satisfy
$\mathrm{inv}(\rho)=\mathrm{inv}(\rho\sigma)$; and their number is $|\mathrm{Rep}(i,i')|/\omega(i)$. By
\eqref{tr:mass} their contribution to $\mathcal N(i,i)$ is $\mathfrak m/\omega(i)$, which is the kernel of
$\mathfrak m\cdot1$. The rest is $\mathcal N^\circ$, by definition, and $\mathcal N^\circ\ge0$ because it is a
weighted count. \eqref{tr:HS} is the formula for the Hilbert--Schmidt norm of an operator with a kernel.

(c) For a Hermitian operator $Q$ on a space of dimension $\nu$ and any real $y$,
$\operatorname{tr}(Q^2)-(\operatorname{tr}Q)^2/\nu=\|Q-\frac{\operatorname{tr}Q}\nu\|_{\rm HS}^2\le\|Q-y\|_{\rm HS}^2$.
Take $\nu=|\mathcal I|$ and $y=|A|\,\mathfrak m$ and use \eqref{tr:HS}. Also
$\operatorname{tr}Q=|A|(\mathfrak m|\mathcal I|+\operatorname{tr}\mathfrak N^{\circ})\ge|A|\,\mathfrak m|\mathcal I|$,
because $\operatorname{tr}\mathfrak N^{\circ}=\sum_i\omega(i)\mathcal N^\circ(i,i)\ge0$. Divide. \qed

\emph{What the lemma shows.} The sum over the twists is a count of pairs of representations subject to one
condition, $\mathrm{inv}(\rho_1)=\mathrm{inv}(\rho_2)$: this is the form a relative trace formula takes when
both groups are definite. The pairs that are forced, $(\rho,\rho\sigma)$, give a multiple of the identity, and
the multiple is the mean number of representations \eqref{tr:mass}. The other pairs are coincidences, and
$\eta$ is their normalised mean square. Part (b) uses three hypotheses, and each is needed. The first is that
the action of $G_i$ is free: with stabilisers of order $2$ the identity of (b) fails. The second is the
invariance $\mathrm{inv}(\rho\sigma)=\mathrm{inv}(\rho)$: without it the forced pairs do not all satisfy the
condition, and the identity of (b) and \eqref{tr:HS} fail. (The number $\mathcal N^\circ$ is a restricted count
and stays nonnegative; it is the kernel $\mathcal N-\mathfrak m/\omega$ on the diagonal that can become
negative.) The third is \eqref{tr:mass}, which is not formal. For lattices it is Siegel's theorem: if
$\mathcal I'$ is a full genus and the $L_i$ belong to one genus, the left side of \eqref{tr:mass} is the mass
of the genus $\mathcal I'$ times a product of local densities, which depend only on the genus of $L_i$
(C.~L.~Siegel, \emph{\"Uber die analytische Theorie der quadratischen Formen}, Ann.\ of Math.\ 36 (1935); in
the adelic form, A.~Weil, \emph{Sur la formule de Siegel dans la th\'eorie des groupes classiques}, Acta
Math.\ 113 (1965); both cited from memory).

\emph{An example.} Let $L_1,L_2$ be the binary lattices with Gram matrices
$\big(\begin{smallmatrix}2&1\\1&12\end{smallmatrix}\big)$ and $\big(\begin{smallmatrix}4&1\\1&6\end{smallmatrix}\big)$
(twice the forms $X_1^2+X_1X_2+6X_2^2$ and $2X_1^2+X_1X_2+3X_2^2$ of discriminant $-23$, which form one genus and represent the
two classes under $\mathrm{GL}_2(\Z)$), with automorphism groups of orders $4$ and $2$, and let $\mathcal I'$
consist of the lattice $\Z^5$, which is alone in its genus, with $\omega'=1/(2^5\,5!)$. A count gives
$|\mathrm{Rep}(L_1,\Z^5)|=|\mathrm{Rep}(L_2,\Z^5)|=1920$, so \eqref{tr:mass} holds with $\mathfrak m=\frac12$; for
$\Z^6$ both numbers are $19200$. (For two forms of determinant $60$ in different genera,
$\big(\begin{smallmatrix}2&0\\0&30\end{smallmatrix}\big)$ and $\big(\begin{smallmatrix}6&0\\0&10\end{smallmatrix}\big)$, the
numbers for $\Z^5$ are $4480$ and $13440$.) Let $\varsigma_0,\dots,\varsigma_9$ be the ten $2\times2$ minors of
the pair of vectors $\rho$, in lexicographic order. Take for $\mathrm{inv}(\rho)$ the $55$ products
$\varsigma_\iota\varsigma_{\iota'}$ modulo a prime $\ell_0$; they are invariant under $G_i$, and
$A=(\Z/\ell_0)^{55}$. For $\ell_0=3$ one finds
$\mathcal N^\circ=\big(\begin{smallmatrix}2&0\\0&3\end{smallmatrix}\big)$. Then the true value
$\operatorname{tr}(Q^2)|\mathcal I|/(\operatorname{tr}Q)^2-1$ is $\frac19$, and the bound \eqref{tr:eta} gives
$\eta=5$. For $\ell_0=5$ and $\ell_0=11$ there are no coincidences,
$\mathcal N^\circ=0$, and $\sum_\chi B_\chi^\dagger B_\chi=|A|\,\mathfrak m$ exactly. With the single invariant
$\mathrm{inv}(\rho)=\varsigma_0^2+2\varsigma_9^2+3\varsigma_1\varsigma_4$ modulo $5$ (so $|A|=5$) the
coincidences dominate: $\mathcal N^\circ\approx\big(\begin{smallmatrix}242&173\\173&209\end{smallmatrix}\big)$. So the family of
twists must be large compared with the number of representations for the count $\mathcal N^\circ$ to be small. The
identities \eqref{tr:HS}, \eqref{tr:eta} were also checked numerically on random finite models.

\subsection*{4.\quad The statements behind \textup{(S)}, \textup{(I)} and \textup{(P)} in the source}

These are assumed here. For each we give the statement that was read, its Lean name, the kind of theorem it
is, and the parameter that produces the saving. References to the literature are from memory.

\textbf{(R1) The expansion of a block and its four parts} (\texttt{ClKppOff}, \texttt{ClKppDiag}, \texttt{ClOff},
\texttt{ClDiag}; they give (S) by \texttt{linkS\_concl}). Each block is an absolutely convergent sum
$K_{\mathsf w\mathsf w'}=\sum_\tau K^{(\tau)}_{\mathsf w\mathsf w'}$ over a set of \emph{terms} $\tau$ (\texttt{Term}, \texttt{aTerm}). A term has an
invariant in the totally real field (\texttt{Gterm}) and may be the identity term, or \emph{generic}. Write
$K''_{\mathsf w\mathsf w'}$ for the sum of the terms with invariant $0$ other than the identity term. The four
statements are: $K''_{\mathsf w\mathsf w'}=0$ for $\mathsf w\ne\mathsf w'$;
$\|K''_{\mathsf w\mathsf w}\|_{\rm HS}^2\le d_{\mathsf w}\exp(-kf\log k)$;
\begin{equation}\label{tr:four}
 \|K_{\mathsf w\mathsf w'}-K''_{\mathsf w\mathsf w'}\|_{\rm HS}^2\le(1+b^{-6})\frac{d_{\mathsf w}d_{\mathsf w'}}{D_*}\ \ (\mathsf w\ne\mathsf w'),
 \qquad
 \|K_{\mathsf w\mathsf w}-K''_{\mathsf w\mathsf w}-1\|_{\rm HS}^2\le\mathsf G\,\frac{d_{\mathsf w}^2}{D_*}.
\end{equation}
With the upper bound in \eqref{tr:proviso}, $\mathsf G\,d_{\mathsf w}^2/D_*\le(C_0/\mathsf G)\,d_{\mathsf w}$: this is
\eqref{tr:S2}, and it shows where $e_0$ comes from. \emph{The diagonal bound is small because each
$\mathcal C_{\mathsf w}$ is a small part of $\mathcal H$}, of relative dimension about $\mathsf G^{-2}$, while
the constant in the second bound of \eqref{tr:four} is only $\mathsf G$. This is an expansion of the kind of
Petersson's formula for the inner products of Poincar\'e series, here for the Bessel periods of a unitary
group in $2(\frac m2-1)+k+1$ variables along a unitary group in $k$ variables; this number is $k+m-1$, the $r$ of \S III.4 (the periods of W.~T.~Gan,
B.~Gross and D.~Prasad, Ast\'erisque 346 (2012)); it is a relative trace formula in the sense of Jacquet. In
the source the squared norm of a part of the sum is expressed as a sum, over \emph{pairs} of terms, of inner
products of vectors attached to the terms (\texttt{pairSum}; compare \eqref{tr:N}), and the
pairs are divided by a rank $s_0\ge1$ and by being ``scalar'' or not. The pairs with $s_0=1$ that are scalar
and generic give the main term. The statements read are bounds for the classes of pairs:
\texttt{RowNGP} (non-generic terms: at most
$2^{-(k-m-3)}\exp(Cm^5f\log b)\,d_{\mathsf w}d_{\mathsf w'}/D_*$) and \texttt{RowMainEqP} (the main pairs of a diagonal
block: at most $\exp(Cm^5f\log b)\,d_{\mathsf w}^2/D_*$). Both are deduced in the source from (R2). The saving in
$\eta_0=b^{-6}$ comes from $k$ and from the level primes: the non-generic and the higher-rank pairs carry
factors $2^{-k}$ and $\exp(-ckf\log k)$ with $k\ge C'm^5fx^\varepsilon$, and the factor
$\prod_v(1+3/|\mathbb F_v|)$ over the places above $g_0,g_1$ is $1+O(b^{-7})$ when $g_0,g_1\ge(kb)^8$
(\texttt{NumPrice}). How these are combined in the proof of the first bound of \eqref{tr:four} was not read.

\textbf{(R2) The Siegel--Weil formula in the convergent range} (\texttt{SiegelWeilUnitaryConvergent}; the theorem
\texttt{siegelWeilUnitaryConvergent} is among the declarations the final theorem depends on). \emph{Let $E/E^+$ be
a CM extension, $J$ a nondegenerate Hermitian matrix of size $M$ over $E$, and $\nu$ an integer with $2\nu<M$.
The quotient of the adelic points of the unitary group of $J$ by its rational points has finite volume, and
for every nondegenerate Hermitian $Z$ of size $\nu$ and every factorisable Schwartz function
$\Phi=\Phi_\infty\otimes\bigotimes_v\Phi_v$, the mean over this quotient of
$\sum_{X}\Phi(gX)$, the sum being over the $\nu$-tuples $X$ of rational vectors with Gram matrix $Z$, equals
$\varrho_\infty(Z)\prod_v\varrho_v(Z)$, where $\varrho_v$ is the density at $Z$ of the image of the measure $\Phi_v\,dX$ under the
Gram map.} (The description of the theta sum is our reading of the names; its definition was not read.)
This is Siegel's theorem in Weil's form, for a unitary group, in a range of absolute convergence; for
unitary groups beyond that range see A.~Ichino, \emph{A regularized Siegel--Weil formula for
unitary groups}, Math.\ Z.\ 247 (2004), and S.~Yamana, \emph{On the Siegel--Weil formula: the case of singular
forms}, Compos.\ Math.\ 147 (2011). It plays the role of \eqref{tr:mass}: it evaluates the mean number of
representations, and so the main pairs, by local densities. It contains no saving by itself.

\textbf{(R3) The count of $\mathcal H$} (\texttt{CountAllPointsClass}; it gives \eqref{tr:S3}).
\emph{$|D-D_*|\le\exp(-c\,rf\log r)\,D_*$, where $D$ is the multiplicity of a discrete series representation,
at a fixed level, in the automorphic spectrum of the unitary group in $r$ variables, and $D_*$ is the product
of its formal degree, the volume of the quotient and the value at $1$ of the idempotent defining the level,}
provided the Harish-Chandra parameter has all gaps at least $Cr^{19}f^5\log(b\,\mathrm{rd}_E)$, and under two
conditions on the sizes of $f$ and $b$. (The description of $D$ and $D_*$ is our reading of the definitions
\texttt{Dcount} and \texttt{Dstar}.) This is a limit
multiplicity statement with a rate (in the tradition of D.~L.~de~George and N.~R.~Wallach, Ann.\ of Math.\ 107
(1978), and L.~Clozel, Invent.\ Math.\ 83 (1986)); $D_*$ is the identity term of the trace formula. The saving
comes from the gaps of the weights. The volume in $D_*$ is a mass. The mass formula of the source is
\texttt{siegel\_mass\_ratio}: the quotient of the masses in ranks $M$ and $M-1$ is
$((M-1)!\,(v/\pi)^M)^f(2v)^{-f/2}$ times a further factor (here $v$ is a positive real number, not a
place, and $\pi$ is the number), which is the recurrence displayed before \eqref{II:t6.4} in
Part~II, \S6.6 (H.~Braun, Abh.\ Math.\ Sem.\ Hamburg 14 (1941); W.~T.~Gan, J.~Hanke and J.-K.~Yu, Duke Math.\ J.\
107 (2001)). Here we read the source's \texttt{vInv}, \texttt{fDeg}, \texttt{lValEta} as $v$, $f$ and that factor;
their definitions, and the place where the formula is applied to $D_*$, were not read.

\textbf{(R4) The count behind \textup{(I)}} (\texttt{closed\_linkIRR\_of\_join\_count}, \texttt{delp\_of\_allRows}).
The space $\mathcal H$ has a basis each of whose $D$ elements carries, at each of the places above $g_0$ and
above $g_1$, a \emph{label}: a monic polynomial of degree $r$ over $\mathbb F_{g_\iota}$ ($\iota=0,1$) with nonzero constant term.
A tuple of labels at the places above one $g_\iota$ is \emph{good} if the polynomials are squarefree, their roots
satisfy two conditions of non-coincidence, for every $j$ with $1\le j\le r/2$ some place has a label with no
factor of degree $j$, and some place has a label with an irreducible factor of prime degree $j'$ with
$r/2<j'\le r-3$ (\texttt{IsGood}). $D_{\rm r}$ is the number of basis elements whose tuples are good at both
primes. The proof of \texttt{delp\_of\_allRows} contains the step
\begin{equation}\label{tr:Wg}
 \delta_{\rm p}\le\frac{1+\varepsilon_C}{1-\varepsilon_C}\big(\mathrm{Wg}(g_0)+\mathrm{Wg}(g_1)\big),\qquad \varepsilon_C=\exp(-c\,rf\log r),
\end{equation}
\begin{equation}\label{D:327}
 \mathrm{Wg}(g)=\frac{3f}g+\frac{23f\log g}{\sqrt g}+2fm\Big(\frac1{4\min(q,d)}+\frac{1+\log r}{g-1}\Big)
      +\frac r2\Big(\frac78\Big)^f+\exp\Big(-\frac f{5\log r}\Big),
\end{equation}
where $q,d$ are the auxiliary primes of \S III.1. The statement \texttt{delp\_of\_allRows} has as a
hypothesis that the right side of \eqref{tr:Wg} is at most a number of the source (\texttt{E7}), and
as its conclusion $0\le\delta_{\rm p}\le\exp(-x^{\varepsilon_z/7})$; \eqref{tr:Wg} itself is not a statement of the
source. Two inputs of that proof: the count (R3), applied to the part of
$\mathcal H$ with prescribed labels, whose expected dimension is $D_*$ times a product over the places of
local weights; and the fact that the local weight of a squarefree label is $g^{-r}(1+O(r/g))$, the labels
that are not squarefree having total weight at most $3/g$ (\texttt{admLq}, \texttt{lawAt}), so that the labels at each place are distributed, up to these errors, like uniformly random monic polynomials of degree $r$ with nonzero constant term. The latter
rests on the parametrisation of the irreducible characters of $\mathrm{GL}_r(\mathbb F_g)$ by J.~A.~Green,
\emph{The characters of the finite general linear groups}, Trans.\ Amer.\ Math.\ Soc.\ 80 (1955), which the
source takes as a hypothesis \texttt{Green} (the cuspidal characters of $\mathrm{GL}_j(\mathbb F_g)$ correspond
to the monic irreducible polynomials of degree $j$ with nonzero constant term and have degree
$\prod_{i=1}^{j-1}(g^i-1)$; for $j=1$ these are the $g-1$ polynomials $X-a$ with $a\ne0$, and the polynomial
$X$ is excluded; Lean: \texttt{IsPrimeLab}). The terms of $\mathrm{Wg}(g)$
are then statements about the factorisation of uniformly random monic polynomials with nonzero constant
term, taken at $f$ independent places. The
bound \eqref{tr:I} follows from \eqref{tr:Wg} and the sizes $g_0,g_1\ge\exp(8x^{\varepsilon_z})$, $f\ge x^2$,
$f,m,r\le x^{C}$ and $\min(q,d)\ge\exp(x^{\varepsilon_z/6})$ (\texttt{SizesAt}). \emph{The term that limits the
bound is $fm/(2\min(q,d))$}; the others are much smaller. The good tuples are those for which the source can
prove that an attached residual Galois representation is irreducible; this use was not read. By the
dependency computation the final theorem depends on \texttt{delp\_of\_allRows} and on \texttt{Green}.

\textbf{(R5) Uniqueness of Bessel models, and decay} (\texttt{besselMapUniqGap}, \texttt{wuDecayOfGap},
\texttt{adm13\_of\_kernel\_join}). \emph{For a discrete series representation of a real unitary group in $r$ variables whose Harish-Chandra
parameter has all gaps at least $120r^3$, a representation of the unitary group in $k$ variables of a given
type, and a generic character, any two Bessel maps on the smooth vectors are proportional.} (The source
states this for $k+2j+1$ variables with any $j\ge0$; here $j=\frac m2-1$, which gives $r$ variables.) Compare Theorem~A of D.~Jiang,
B.~Sun and C.-B.~Zhu, \emph{Uniqueness of Bessel models: the archimedean case}, Geom.\ Funct.\ Anal.\ 20 (2010):
it gives multiplicity at most one for Bessel models at an archimedean place. The statement of the source has
in addition a hypothesis on the gaps. Theorem~A does not give the quantitative decay used here. The decay is
a further statement of the source: from the uniqueness it
deduces an integrability statement for a function attached to a word (\texttt{WuDecay}; the
name refers to this function), under the condition that $r\ge1801$ and that consecutive weights $\theta_j$ of \S5 differ
by at least $16r^{19}$ (\texttt{GapHyp}); by our reading this serves the convergence of the expansion in (R1).
Statement (P2) is
proved by \texttt{adm13\_of\_kernel\_join}, from hypotheses that were not read.

\textbf{Where the exponents come from.}
\begin{center}\small
\begin{tabular}{lp{.26\textwidth}p{.48\textwidth}}
\toprule
number & value in the source & large parameter\\
\midrule
$\eta_0$ & $b^{-6}$ & $k$, and $g_0,g_1\ge(kb)^8$\\
$e_0$ & $2C_0/\mathsf G+2\exp(-kf\log k)$ & $\mathsf G=\exp(C_{\mathsf G}m^5fx^\varepsilon)$: the ratio $D_*/d_{\mathsf w}\approx\mathsf G^2$\\
$\eta_1$ & $\exp(-c\,rf\log r)$ & the rank $r\ge k$ and the degree $f$\\
$\alpha_{\rm ram}$ & $2f/b$ & the level primes $\mathcal T(\mathsf w)\ge b^{m/2}$\\
$\delta_{\rm p}$ & $fm/(2\min(q,d))$: leading term in the upper bound for $\delta_{\rm p}$ & the auxiliary primes $q,d\ge\exp(x^{\varepsilon_z/6})$\\
bound for $\eta$ & $5x^{-C_H}+3\sqrt{\alpha_{\rm ram}}$ & the number of words, $|\mathsf W|=\lceil3\mathsf G^2x^{C_H}\rceil$\\
\bottomrule
\end{tabular}
\end{center}
The exponent $x^{\varepsilon/8}$ in \eqref{tr:small} is not the true size of any of these numbers. It is
chosen in the source so that one bound, $\exp(-y)$ with $y=x^{\varepsilon/8}$ and $x^\varepsilon=y^8$, covers
$\eta_0$, $\eta_1$ and $\sqrt{e_0}$ (\texttt{eF\_nonneg\_and\_le\_exp\_neg}). The exponent $\varepsilon_z/7$ in
\eqref{tr:I} is just below the $\varepsilon_z/6$ of the auxiliary primes. The last row is the bound for the constant $\eta$ of Part~II, \S6.9 proved in the
source (\texttt{frame\_at\_of}); see the paragraph on the error term at the end of \S2 for what it does and does not say about $\eta$.

\subsection*{5.\quad The interlacing of the weights}

At each weight place the source fixes rational numbers
\begin{equation}\label{tr:inter}
 \xi_0>\psi_0>\xi_1>\psi_1>\dots>\psi_{k-2}>\xi_{k-1}>-\theta_0,\qquad \theta_0<e_0'<\theta_1<e_1'<\dots<\theta_{m-1}<e_{m-1}',
\end{equation}
with $\xi_i\in\Z+\frac12$ and $\psi_i,\theta_j,e_j'\in\Z$ (\texttt{HCParam}; the condition is named \texttt{Interlace}; we
write $\xi_i,\psi_i,\theta_j,e_j'$ for the numbers called \texttt{u}, \texttt{v}, \texttt{b}, \texttt{e} there; $e_0'$ is not the
$e_0$ of (S)). The first chain is
the condition under which the representation of the unitary group in $k-1$ variables with parameter $\psi$
occurs in the restriction of the one in $k$ variables with parameter $\xi$: the classical branching rule, here for parameters
shifted by half the sum of the positive roots. In this section it is used in one place, the definition of
the family: \emph{the words are the weights $\xi$ that interlace one fixed $\psi$}. Precisely, consecutive $\psi_l$
differ by $4\Delta$ for a fixed $\Delta$, the numbers $\xi_0$ and $\xi_{k-1}$ are the same for all words, and for
$0\le l\le k-3$
\begin{equation}\label{D:329}
 \xi_{l+1}=\psi_l-\tfrac12-\Delta\cdot(\text{the }l\text{th letter of the word}),
\end{equation}
and for both values of the letter this number lies between $\psi_{l+1}$ and $\psi_l$, since \eqref{tr:inter} is
part of the data for every word (\texttt{cloud\_gap}, \texttt{u\_letter}, \texttt{hcW\_v}, \texttt{hcW\_first},
\texttt{hcW\_last}). So two words differ in the weight $\xi$ at the positions
where their letters differ. That the bounds \eqref{tr:four} for $\mathsf w\ne\mathsf w'$ use this (through the
archimedean factors of the terms) is our reading of the definitions; the proofs were not read. The second
chain of \eqref{tr:inter}, with gaps at least $(t_0+3)/2$ (\texttt{gap\_be}, \texttt{gap\_eb}), concerns the weights along the chain of lifts and
belongs to the divisibility and size estimates (\S III.4, \S III.5).

\subsection*{6.\quad Assumed here, from the source}

\begingroup\sloppy
\begin{enumerate}[label=\textup{(\roman*)},leftmargin=2.4em]
\item \emph{Statement \textup{(S)}}: \texttt{LinkS} with $C_0=3$, proved by \texttt{link\_alt\_alt\_bis} from
\texttt{theoremS\_transport\_link\_zero}, for families with \texttt{TuSize}. Its inputs, as hypotheses of
\texttt{link\_alt\_alt\_bis}: the decay \texttt{WuDecayOfGap}; a finiteness statement \texttt{RowsFinDim}; the
bounds \texttt{TypeP}, \texttt{RowNGP}, \texttt{RowMainEqP}, \texttt{LoopRowS}; the off-diagonal bound
\texttt{TheoremSSelected}; the count \texttt{CountAllPointsClass}; and compatibility conditions. Only the
statements were read, and for \texttt{RowNGP}, \texttt{RowMainEqP} the deduction from (R2) by name.
\item \emph{Statement \textup{(I)}}: \texttt{LinkIRR}, proved by \texttt{closed\_linkIRR\_of\_join\_count} from
\texttt{Green} and a count hypothesis (\texttt{CountAllRows}).
\item \emph{Statement \textup{(P)}}: \texttt{BesselId}, the conjunction of \texttt{BuRational} and
\texttt{adm13}; the latter by \texttt{adm13\_of\_kernel\_join}.
\item \emph{The side conditions of Corollary~T2}: the proviso \eqref{tr:proviso} (\texttt{ChProviso\_at}) and
the number of words (\texttt{ChCode\_at}), which are conditions on the family; the formula for
$d^{\rm ram}_{\mathsf w}$ (\texttt{RamRank\_at}) and $\mathcal T(\mathsf w)\ge b^{m/2}$ (\texttt{TuGe}). That
$c_0(\mathsf w)>0$ is not assumed: it follows from (S), see the proof of Corollary~T2.
\item \emph{The identification of the data of \S1} with spaces of automorphic forms, with Poincar\'e series
and with the pairing of the chain of lifts; and the passage from Corollary~T2 to the mean-slope inequality
(\texttt{TheoremF}, \texttt{frame\_at\_of}), which uses the slope filtration of Part~II, \S6.10. (A
passage from Corollary~T2 to the mean-slope inequality is derived in \S\ref{as:sec}, Proposition~M4, from
the rank bound and a cited fact on the filtration by slopes; it is not the passage of the source.) For what is
known about the dependence of the final theorem on these two declarations see the end of \S2.
\item (R1)--(R5) of \S4.
\end{enumerate}

\textbf{What this section proves, and what it assumes.} This section deduces the dimension bound from
the three trace, rank and pairing statements (S), (I), (P) and the size and subspace conditions
\eqref{tr:side}, \eqref{tr:sub} in T2. \emph{Proved outright:} Lemma~T1, including the
polynomial inequality \eqref{tr:poly}; Corollary~T2, for arbitrary subspaces, as that deduction; Lemma~T3;
and the arithmetic of the example of \S3 (by computation). \emph{Quoted from the
source without proof:} (S), (I), (P) and everything in \S4 and \S6. So the construction of the spaces and
maps and the arithmetic inputs are quoted. The assembly of (M) from Corollary~T2 and the other inputs is
in \S\ref{as:sec}; it is a reduction of (M) to numbered Hypotheses, not a proof of (M). Lemma~T3 is a model. It is not a step of
the source's proof: there the family is the set of words, the terms are attached to a Bessel subgroup, and
the larger group is, by our reading, not definite (the parameter \texttt{sigY0} of the source is the
signature $(k+\frac m2-1,\frac m2)$; it is the signature of the space \texttt{Y 0} of the record at every
embedding, by the field \texttt{baseSpace\_signature\_eq}). The statements \texttt{FrameLA}, \texttt{frame\_linear\_algebra\_bound},
\texttt{ThmSFree} and \texttt{KernelFree} of the archive are not used by the final theorem; the declarations
\texttt{LinkS}, \texttt{LinkS\_at}, \texttt{BesselId}, \texttt{blockReal\_at\_of}, \texttt{link\_alt\_alt\_bis},
\texttt{theoremS\_transport\_link\_zero}, \texttt{closed\_linkIRR\_of\_join\_count}, \texttt{rowNGP\_of\_siegelWeil},
\texttt{rowMainEqP\_of\_siegelWeil}, \texttt{siegelWeilUnitaryConvergent}, \texttt{adm13\_of\_kernel\_join},
\texttt{wuDecayOfGap}, \texttt{besselMapUniqGap} and \texttt{siegel\_mass\_ratio} are. So are, by a later
run of the same computation, \texttt{LinkIRR}, \texttt{TheoremF}, \texttt{FrameLA1},
\texttt{frame\_linear\_algebra\_first}, \texttt{Frame\_at}, \texttt{FrameF}, \texttt{frameF\_of},
\texttt{frame\_at\_of}, \texttt{matRealF\_of}, \texttt{matSF\_of}, \texttt{ChCode\_at}, \texttt{CodeUpF},
\texttt{RamRank\_at}, \texttt{ramShare\_at\_of}, \texttt{delp\_of\_allRows}, \texttt{Green},
\texttt{matKF\_of}, \texttt{unfolding\_constant\_pos\_of\_diagonal}, \texttt{TheoremFSharp},
\texttt{RRV\_finrank}, \texttt{jEE\_basis} (the field of \texttt{D7}) and
\texttt{CompressTraceSquareBound}. The final theorem does not depend on \texttt{CorBounds},
\texttt{TraceSquareAndDiagBounds} or \texttt{MeanBoundSecondForm}. Several declarations of the archive are named \texttt{offDiag}. The one that concerns this
section is \texttt{GlobalTheory.CarrierGap.offDiag}, which proves the off-diagonal bound
\texttt{TheoremSSelected} of (i) from the Siegel--Weil statements and \texttt{LoopRowS}. The final theorem
depends on it, and also on \texttt{each\_unit\_bound\_statement} and on \texttt{Data.SlopeIneq} (\S2).

\textbf{What remains.}
\begin{enumerate}[label=\textup{(\roman*)},leftmargin=2.4em]
\item The definition of $\mathcal H$, $\mathcal C_{\mathsf w}$ and $P_{\mathsf w}$ in classical terms (the group, the
level, the discrete series, the Poincar\'e series), and of the terms of the expansion (R1).
\item A proof of the first bound of \eqref{tr:four}: the main term with constant $1$ from (R2), and the bounds
for the other pairs. This is the centre of (S).
\item A proof of (P): why the pairing of the chain is an inner product with $P_{\mathsf w}\varphi$.
\item The count (R3) with its rate, and the distribution of the labels in (R4).
\item A run of the dependency computation for the names that were not queried: \texttt{Win\_at},
\texttt{SizeKRel}, \texttt{SizeM}, \texttt{SizeN}, \texttt{SizeK}, \texttt{SizesAt},
\texttt{CompressDiagDefectBound}.
\item Whether the number $\operatorname{tr}(Q^2)D/(\operatorname{tr}Q)^2-1$ of the family of the source is of the
size $x^{-C_H}$ or smaller; only the upper bound is known (the paragraph on the error term at the end of \S2).
\end{enumerate}
\endgroup


\section{The bound for Galois cohomology}\label{gc:sec}

This section concerns the bound \eqref{II:shabound} of Part~II, \S6.8 (Part~I, \S I.7, ``Denominators''): a bound
for the lengths of certain Galois cohomology groups of the field $E$, which the announcement \cite{Alp26}
calls a weakening of Vandiver's conjecture. It is the part of the proof that uses the most mathematics from
outside analytic number theory, and no proof of it is given here. The section states it exactly, shows how it
is used, describes the structure of its proof in the Lean development with the published theorems that
correspond to each step, and works out the simplest case.

\emph{Proved outright.} Lemma~G1 (lengths, orders and indices); Proposition~G2 (how the bound gives the
estimate for the denominators, given the inequality \texttt{LinkN\_at} of the source).

\emph{Proved from theorems of books that we consulted.} Proposition~G3: for the field $\Q$ the groups
vanish when a class group vanishes; with its corollaries for odd twists (Herbrand's theorem), for the twists
with $(\ell-1)\mid 2e$, and under Vandiver's conjecture. The theorems used are in Washington's book
\cite{Was82}; their statements and numbers were checked in a copy of the first edition. Lemma~G4: the class
number of $\Q(\sqrt\ell)$ is less than $\ell$; it is proved from Minkowski's theorem on lattice points
\cite[Ch.~I, (4.4)]{Neu99}. The numerical statements of \S5 are computed exactly.

\emph{Proved from the standard statements of class field theory.} The vanishing of
$\mathrm{Sha}_5(\Q(i),4)$, at the end of \S5. The proof is for this one field and this one twist.

\emph{Taken from the source without proof.} Theorem~G (the bound itself, Lean: \texttt{VWinAt} for the
trivial character); the inequality \texttt{LinkN\_at}; and every statement of \S4.

\emph{Citations.} They are of three kinds. (1) The books \cite{Was82}, \cite{Neu99} and \cite{NSW00} were
read in local copies (\cite{Was82} and \cite{NSW00} in the first edition); the test group G19 compared each
statement cited from them with the page. (2) The articles \cite{Sou03}, \cite{Kat04}, \cite{Kur92},
\cite{BC04}, \cite{Car12}, \cite{Car14}, \cite{CH13}, \cite{BS23} and the book \cite{Rub00} were not available to us.
Where a theorem or a section of one of them is named, the wording was compared with the publication by an
outside reviewer, and we say no more than his report says. (3) Everything else is cited from memory and is
marked so: \cite{Sou79}, the theorem of Voevodsky and Rost, \cite{MW84}, \cite{FW79}, the work of Harris and
Taylor and of Shin, and the last column of the table of \S4.

Letters introduced in this section: $\ell$ (a prime, as in \S III.5), $e$ (the twist; Part~II, \S6.8 writes
$a$), $t$ (a number of twists; in the construction it is the parameter written $t_0$ in \S III.5, and $t$
in \S III.4; Part~II, \S6.8 does not name it and writes the bound $\lfloor x^{400}\rfloor$), $V_0$, $S_\ell$, $G_{E,\ell}$, $G_v$, $W_e$, $\mathrm{Sha}_\ell(E,e)$,
$h_{\ell,e}$, $h^{\pm}_{\ell,e}$, $\jmath$ (complex conjugation), $\psi$ (a Dirichlet character),
$\mathfrak B$, $\eta_j(w)$, $H_j$, $y=x^{\varepsilon}$, $V$, $e_O$, $f_O$, $\kappa$, $c_\psi$, $x'$; in \S4:
$\mathcal K$, $\mathcal R$, $R$, $R'$, $c_0$, $\vartheta$, $\varrho$; and in \S5: $A$, $A^{(i)}$,
$X_\infty$, $H_\infty$, $\Gamma$, $\Pi$, $\omega$, $F$, $h$, $\mathfrak d$, $\mathfrak a$, $\alpha$, $T$, $L$,
$j$, $N$, $c$, $\delta$, $\lambda$, $z$, $\phi$, $a_1$, $a_2$, $\mathrm{rec}$, $u$, $\epsilon$, $\beta$. The
letters $x,\varepsilon,b,m,f,q,d,n$ have their meaning from Parts I and II; in particular
$b=\exp(x^{\varepsilon})$.

\emph{The letters are local.} The letters of this section have the meanings given here in this section
only. The clashes most likely to confuse are these. $n$ is the power of $2$ with $[E:\Q]=4n$; \S III.5
takes its $n$ from the record of the construction (\texttt{ord}), and \S1 says how the order of the frame
is related to the $n$ of Theorem~G; Soul\'e's twist, which he writes $n$, is written $n'$ in \S4. $t$ is a
number of twists, the $t_0$ of \S III.5; in \S III.5 the letter $t$ is an index. $e$ is the twist; $e(w)$
and $e_O$ are ramification indices (the $e_w$ of \S III.6); $e$ is never the base of the natural logarithm
here. $L_0$ is the base field of a tower, in \S4 only; the coefficient field, which is $L_0$ in \S III.6, is
called $K$ here, and $\mathcal K$ is an imaginary quadratic field. $\psi$ is a Dirichlet character and
$\Psi=\psi\chi_\ell^{e}$ a character; in \S III.2 $\psi(u,\chi)$ is a sum over prime powers, in \S III.4 and
\S III.7 $\psi$ and $\Psi$ are ranks, and in \S III.8 they are weights and unitary maps. $C$ is the
constant of \eqref{gc:link}, or an exponent in the table of \S4; it is not the $C$ of another section, and
it is not one of $c$ (a homomorphism, in \S5), $c_0$, $c_\psi$. Inside the section four letters have two or
more uses. $T$ is a finite group in Lemma~G1 and a real number in Lemma~G4. $N$ is a level or a conductor
in \S3 and \S4, the norm $N(w)$, $N(\mathfrak a)$ of a prime or an ideal, and the norm map of class groups
in the proof of the corollaries. $j$ numbers the maps of the chain in \S2 and is the map that extends
ideals in \S5. $\epsilon=\pm1$ is a sign in Step~4 of \S5; it is not $\varepsilon$.

\subsection*{1.\quad The statement in the source}

\textbf{The groups.} Let $E$ be a number field, $\ell$ a prime and $e$ an integer. Let $S_\ell$ be the set of
primes of $E$ above $\ell$, and $G_{E,\ell}$ the Galois group over $E$ of the largest extension of $E$ in which no
finite prime outside $S_\ell$ ramifies. Let $\chi_\ell\colon G_{E,\ell}\to\Z_\ell^\times$ be the cyclotomic
character, and let $W_e=\Q_\ell/\Z_\ell(-e)$: the discrete group $\Q_\ell/\Z_\ell$ on which $\sigma$ acts by
multiplication by $\chi_\ell(\sigma)^{-e}$. For $v\in S_\ell$ let $G_v\subset G_{E}$ be a decomposition group at
$v$, mapped to $G_{E,\ell}$. Put
\begin{equation}\label{gc:sha}
\begin{aligned}
 \mathrm{Sha}_\ell(E,e)&=\ker\Big(H^1(G_{E,\ell},W_e)\longrightarrow\prod_{v\in S_\ell}H^1(G_v,W_e)\Big),\\
 h_{\ell,e}&=\operatorname{length}_{\Z_\ell}\mathrm{Sha}_\ell(E,e)\ \in\ \{0,1,2,\dots\}\cup\{\infty\}.
\end{aligned}
\end{equation}
Cohomology is that of continuous cochains: $H^1$ is the group of continuous maps $\xi$ with
$\xi(\sigma\sigma')=\xi(\sigma)+\sigma\cdot \xi(\sigma')$, modulo the maps $\sigma\mapsto \sigma\cdot w-w$. So
$\mathrm{Sha}_\ell(E,e)$ consists of the
classes that are unramified outside $\ell$ and trivial at every prime above $\ell$, as in Part~II, \S6.8. The
length of a $\Z_\ell$-module is the length of a longest chain of submodules; a finite module of length $h$
has $\ell^{h}$ elements, and the length is $\infty$ if there is no finite bound.
In the Lean: \texttt{LData} (a choice of a place of $\overline E$ above each $v\in S_\ell$ and of the
character $\chi_\ell$ on $G_{E,\ell}$), \texttt{Sha} and \texttt{shaLen}
(\texttt{Formalism/GaloisCohomology/Statement.lean}, lines 54, 75, 78); the value $\infty$ is written $\top$.
The length does not depend on the choices; this is proved in the source (\texttt{shaLen\_eq\_}\allowbreak\texttt{shaInv}).

\textbf{The statement} $\mathrm{VVle}(E,t,V_0)$, for an integer $t\ge1$ and a real number $V_0$ (Lean: \texttt{VVle}, in the namespace
\texttt{UpperTheory.Top}; \texttt{ExistsPageN.lean}:12).
\emph{For every algebraic closure of $E$ and every choice of the data above, for every prime $\ell$:}
\begin{equation}\label{gc:vvle}
 h_{\ell,e}<\infty\quad(1\le e<t),\qquad\text{and}\qquad
 \sum_{\ell\in S}\log\ell\sum_{1\le e<t}h_{\ell,e}\ \le\ V_0\quad\text{for every finite set $S$ of primes.}
\end{equation}
(The first clause is needed: in the Lean sum an infinite length is counted as $0$.) If
$\mathrm{VVle}(E,t,V_0)$ holds, $t'\le t$ and $V_0\le V_0'$, then $\mathrm{VVle}(E,t',V_0')$ holds
(\texttt{VVle.mono}).

\textbf{The field.} For $n\ge1$ a power of $2$ and primes $q\equiv3\pmod 8$, $d\equiv 7\pmod{16}$, let
$E=\Q(\zeta_{4n}+\zeta_{4n}^{-1},\sqrt{-q},\sqrt{-d})$. It is the subfield of a cyclotomic
field fixed by the automorphisms $\zeta\mapsto\zeta^{\alpha}$ with $\alpha\equiv\pm1\pmod{4n}$ and
$\big(\tfrac \alpha q\big)=\big(\tfrac \alpha d\big)=1$ (Lean: \texttt{IsChiField} for the trivial character,
line 19). It is a field with complex multiplication, of degree $4n$ over $\Q$, and its Galois group over $\Q$ is
$\Z/n\times(\Z/2)^2$.

\textbf{Theorem G} \quoted. (Lean: \texttt{VWinAt} for the trivial character and the exponent $401$,
lines 44--49 of the same file; proved in the source as \texttt{vwin\_of\_row}.) \emph{Let
$0<\varepsilon<1/52$, and for $x\ge1$ let $n(x)$ be a power of $2$ with
$x^{2+8\varepsilon}\le n(x)\le 2x^{2+8\varepsilon}$. There are $J\ge1$ and $X_0$ with the
following property. Let $X\ge X_0$, put $Y_0=X^{1/J}$, and let $q\equiv 3\pmod 8$ and $d\equiv7\pmod{16}$ be
primes with
\begin{equation}\label{D:332}
 \tfrac12Y_0^{\varepsilon/5}<\log q\le Y_0^{\varepsilon/5},\qquad
 \tfrac12Y_0^{\varepsilon/5}\log^2Y_0<\log d\le Y_0^{\varepsilon/5}\log^2Y_0,\qquad q\mid d+1,
\end{equation}
and such that every odd prime $p\le Y_0^{\varepsilon/5}$ divides $d+1$. Then there is $Y$ with
$Y_0\le Y$ and $Y^2\le X$ such that, for every $x\in[Y,Y^2]$ and for the field $E$ with $n=n(x)$,}
\begin{equation}\label{gc:G}
 \mathrm{VVle}\big(E,\ \lfloor x^{401}\rfloor,\ \sqrt b\,\big)\ \text{holds},\qquad b=\exp(x^{\varepsilon}).
\end{equation}
The congruences and sizes required of $(q,d)$ are those of the primes of Estimate~B (Lean:
\texttt{WinPair}). The interval $[Y,Y^2]$ is the ``window''; its origin is explained in \S4.

\textbf{From Theorem G to the bound of Part~II.} The bound \eqref{II:shabound} of Part~II has $\lfloor x^{400}\rfloor$
where \eqref{gc:G} has $\lfloor x^{401}\rfloor$. The source passes from one to the other by a change of scale
(Lean: \texttt{winPairs},
\texttt{WinPairs.lean}:5--81). Let $c_\psi\ge1$ be the constant of the
source such that the order of the frame of the construction is $c_\psi$ times the power of $2$ of
Theorem~G (the constant of \texttt{fieldChi}; we did not evaluate it for the trivial character). Let $\kappa\ge1$ be the
number with $\kappa^{2+8\varepsilon}=c_\psi$. Theorem~G is applied at a point $x'$, and the construction
works at the scale $x=\kappa x'$. Then $x^{2+8\varepsilon}=c_\psi\,x'^{\,2+8\varepsilon}$, so the bounds for
$n(x')$ become $x^{2+8\varepsilon}\le c_\psi\,n(x')\le2x^{2+8\varepsilon}$. The proof takes $x'$ so large
that $\kappa^{400}\le x'$. Hence
\begin{equation}\label{D:334}
 \lfloor x^{400}\rfloor=\lfloor\kappa^{400}x'^{\,400}\rfloor\le\lfloor x'^{\,401}\rfloor,\qquad
 \exp(x'^{\,\varepsilon})\le\exp(x^{\varepsilon}),
\end{equation}
and \eqref{gc:G} at $x'$ gives $\mathrm{VVle}(E,\lfloor x^{400}\rfloor,\sqrt b)$ with $b=\exp(x^\varepsilon)$,
by \texttt{VVle.mono}. This is the field \texttt{bound} of the record \texttt{WinDatum}
(\texttt{TwistForm.lean}:25). The window $[Y,Y^2]$ for $x'$ becomes a window $[Z,Z^{3/2}]$ for $x$, with
$Z=\kappa Y$; this uses $\kappa^{3}\le Y$. So the remark of Part~II, \S6.8, that the source states the bound
for twists below $x^{401}$ at a slightly different scale, is accurate.

The number of twists of the construction (Lean: \texttt{R.t}) is not a free parameter. The condition
\texttt{RankVal} (\texttt{PairRowIffQdRows.lean}:28--32) fixes it by an equation in the other parameters, and
\texttt{eventually\_rank\_le\_}\allowbreak\texttt{rpow\_of\_lt\_one\_div\_fiftytwo} proves that it is at most $x^{400}$ for
$x\ge2^{100}$. The theorem \texttt{supplyOfAdmOf} then applies \texttt{VVle.mono} once more. The family of
the construction receives the result in the form \texttt{HypV}: for large $x$,
$\mathrm{VVle}(E,t,\sqrt b)$ holds for its fields $E$ and its number of twists $t$
(\texttt{DivinNOfHypV.lean}:10). The declarations named in this paragraph and in the preceding one are
among those the final theorem depends on.

\subsection*{2.\quad From lengths to indices of lattices}

\textbf{Lemma G1.} \emph{\textup{(i)} Let $T$ be a finite abelian group and $h_\ell$ the length of its
$\ell$-primary part as a $\Z_\ell$-module. Then $\log|T|=\sum_\ell h_\ell\log\ell$.
\textup{(ii)} Let $O$ be the ring of integers of a finite extension of $\Q_\ell$, with ramification index
$e_O$ and residue degree $f_O$. If $T$ is an $O$-module of finite length, then
$\operatorname{length}_{\Z_\ell}T=f_O\operatorname{length}_OT$ and $|T|=\ell^{f_O\operatorname{length}_OT}$.
If $T_0$ is a finite $\Z_\ell$-module, then $\operatorname{length}_O(O\otimes_{\Z_\ell}T_0)=e_O\operatorname{length}_{\Z_\ell}T_0$.
\textup{(iii)} Let $M'\subset M$ be lattices of the same rank over the ring of integers $O_K$ of a number
field $K$. Then
$\log[M:M']=\sum_w\operatorname{length}_{O_{K,w}}\big((M/M')_w\big)\log N(w)$,
the sum over the primes $w$ of $K$; here $N(w)$ is the norm and the subscript $w$ is the completion.}

\emph{Proof.} (i) A finite $\Z_\ell$-module of length $h$ has a chain of submodules with $h$ successive
quotients $\Z/\ell$, so it has $\ell^{h}$ elements; and $T$ is the product of its primary parts. (ii) The
simple $O$-module is the residue field, which has $\ell^{f_O}$ elements and length $f_O$ over $\Z_\ell$.
For $T_0=\Z/\ell^{s}$ the module $O/\ell^{s}O$ has length $e_Os$, because $\ell O$ is the $e_O$th power of the
maximal ideal; a general $T_0$ is a product of such groups. (iii) $M/M'$ is the product of its completions,
by the Chinese remainder theorem, and a module of length $h$ over $O_{K,w}$ has $N(w)^{h}$ elements.\qed

By (ii), $\operatorname{length}_O(O\otimes T_0)/e_O$ does not depend on $O$. The Lean uses this normalisation
when it enlarges the coefficients to treat a group character by character (\texttt{Chr.nrm}: the length
divided by the length of $O/\ell O$, which is $e_O$; Lean: \texttt{eO}).

\textbf{The inequality of the source.} The construction has a coefficient field, a number field which we
call $K$ here (Lean: \texttt{LL}), and for each map of the chain, numbered $j=1,\dots,m$, and each prime $w$ of
$K$ a rational number $\eta_j(w)\ge0$ (Lean: \texttt{etaJ}). A prime $w$ is called \emph{twisted} if the
character twist chosen at $w$ is not the base one (Lean: $\mu_w\ne\mu_0$). Put
\begin{equation}\label{gc:eta}
 H_j=\sum_{w\ \mathrm{twisted}}\frac{\log N(w)}{[K:\Q]}\;e(w)\,\eta_j(w),
\end{equation}
where $e(w)$ is the ramification index of $w$ over $\Q$ and the sum is over a finite set (Lean:
\texttt{etaSum}, \texttt{GlobalTheoremInputs.lean}:36). The sum runs over twisted $w$ only.

\emph{What is known about the meaning of $\eta_j(w)$.} The number $\eta_j(w)$ of \eqref{gc:eta} is the
field \texttt{etaJ} of the record of data of the final theorem (\texttt{DataRows.lean}:162). The source
also has a number with the same name for a family of lattices (\texttt{DefFam.etaJ},
\texttt{LiftLat.lean}:147). For that number the meaning is clear. Times $e(w)$ it is the length of a
quotient of two lattices at $w$: the definition is \texttt{BlockDat.etaW} (\texttt{LiftLat.lean}:116), the
length over the localisation at $w$ of $A/(A\cap A')$ for two lattices $A$, $A'$ (\texttt{locLen}), and the
relation $\eta\cdot e(w)=\texttt{etaW}$ is \texttt{eta\_units} (\texttt{LiftLatLemmas.lean}:46). The final
theorem depends on \texttt{etaW}, \texttt{locLen} and \texttt{eta\_units}.
\emph{It is not established here that the field \texttt{etaJ} of the record equals the number of the
family at the twisted primes.} An earlier version of this section cited the declarations \texttt{TieEta}
(\texttt{Statement.lean}:90) and \texttt{etaN\_eq\_etaJ} (\texttt{CoverAct.lean}:108) for this. By the
dependency computation of the test group G19, both are in the archive and the final theorem depends on
neither. So they establish nothing about the $\eta_j(w)$ that the final theorem uses. The statement that
is missing is this: for the record of data of the final theorem and every twisted $w$,
$e(w)\,\eta_j(w)$ equals \texttt{etaW} at $w$, a length of $A/(A\cap A')$. The archive contains statements of this
form (the field \texttt{etaJ} of the record \texttt{NumTie}, \texttt{RowSideTie.lean}:244, and the theorem
\texttt{etaJ\_join}); whether the final theorem depends on them was not computed. The record itself only
says that $\eta_j(w)$ is a sum of local numbers (its field \texttt{etaJ\_tie}). To decide the question one
has to read the proof of \texttt{LinkN} (\texttt{linkN\_of\_twGood\_}\allowbreak\texttt{galoisPreconditions},
which the final theorem depends on). Until this is done, the following is an expectation and not a
statement of this paper: that $e(w)\eta_j(w)$ is the length at $w$ of a quotient $M/M'$ of two lattices over
$O_K$; that, by Lemma~G1(iii), $[K:\Q]\,H_j$ is the part of $\log[M:M']$ that comes from the twisted primes;
and that $\eta_j(w)$ is the defect of the index of lattices of Part~II, \S6.8. No proof in this paper uses
it. Proposition~G2 and the chain use only the inequality \texttt{LinkN\_at}.

\texttt{LinkN\_at} (line 40) \quoted. \emph{If $\sum_\ell\log\ell\sum_{1\le e<t}h_{\ell,e}$ converges, then}
\begin{equation}\label{gc:link}
 H_j\ \le\ d_j\Big(C\big(f\,m\,b\,x^{2\varepsilon}+(\log x)^2N_c\big)+\sum_{\ell}\log\ell\sum_{1\le e<t}h_{\ell,e}\Big).
\end{equation}
Here $C$ is a constant, $d_j$ is a dimension (Lean: \texttt{dZ0}; the statement \texttt{SmallLevelRankEq}
says that $d_j$ equals a number $D_0$ not depending on $j$), $f\ge1$ and $m\ge1$ are integers of the
construction ($m$ is the number of maps), and $N_c\ge0$ (Lean: \texttt{budgetNc}). The statement
\texttt{LinkN} asserts \eqref{gc:link} for all members of the family, for large $x$, with one $C$, assuming
the first clause of \eqref{gc:vvle}. It is proved in the source; see \S6(ii).

\textbf{Proposition G2} (Lean: \texttt{divinN\_at}, \texttt{DivinNOfHypV.lean}:76). \emph{Let $x>1$,
$y=x^{\varepsilon}$, $b=\exp(y)$, $f,m\ge1$, $N_c\ge0$, $D_0\ge0$, $C$ real, $B_0=f\,m\,b\,y$. Suppose
$(\log x)^2\le y$, that \eqref{gc:link} holds with $d_j=D_0$, and that
$V=\sum_\ell\log\ell\sum_{1\le e<t}h_{\ell,e}\le b$. Then}
\begin{equation}\label{gc:div}
 H_j\ \le\ D_0\,(|C|+1)\,x^{\varepsilon}\,(B_0+N_c).
\end{equation}

\emph{Proof.} Since $f\,m\,b\,y^{2}=yB_0$ and $(\log x)^2N_c\le yN_c$, the first term in the bracket of
\eqref{gc:link} is at most $|C|\,y\,(B_0+N_c)$. Since $f,m,y\ge1$, we have $V\le b\le B_0\le y(B_0+N_c)$.\qed

\eqref{gc:div} is the statement \texttt{DIVinN\_at}; it is the bound for the denominators that the rest of
the chain uses. Theorem~G gives $V\le\sqrt b$, and $\sqrt b\le b$. So at this step a bound by $b$ would
suffice. The square root is not used anywhere else. In the source the statement \texttt{HypV} is used by one
theorem only, \texttt{divinN\_of\_hypV} (its single use is in \texttt{DesWin.lean}:29), and that theorem
replaces $\sqrt b$ by $b$ at once, by $\sqrt b\le b$ for $b\ge1$. Besides the bound it takes from
\texttt{HypV} the finiteness of the lengths and the convergence of the sum (\texttt{vvRecord}). So a bound by
$b$ in Theorem~G would suffice for the deduction.

\subsection*{3.\quad Why the statement is of Vandiver type}

Let $\ell$ be odd and let $\jmath$ be complex conjugation on $E$. It acts on $\mathrm{Sha}_\ell(E,e)$ as an
involution, and the Lean splits the group into the part where $\jmath=(-1)^{e}$ and the part where
$\jmath=(-1)^{e+1}$, of lengths $h^{+}_{\ell,e}$ and $h^{-}_{\ell,e}$; $h_{\ell,e}=h^+_{\ell,e}+h^-_{\ell,e}$
(\texttt{shaLen\_eq\_add}, proved there by linear algebra). Let $E$ be abelian over $\Q$ and
$\ell\nmid[E:\Q]$. In the classical theory the group then splits, after enlarging the coefficients to a ring
$O$, into parts indexed by the Dirichlet characters $\psi$ of $\mathrm{Gal}(E/\Q)$; the part of $\psi$ lies
in the plus part if $\psi(-1)=(-1)^{e}$ and in the minus part otherwise. The source does not state a
decomposition. It proves an inequality (\texttt{HV.redPrime}): for $\ell\ne2$ and $\ell\nmid[E:\Q]$,
\begin{equation}\label{D:338}
 e_O\,h^{+}_{\ell,e}\ \le\ \sum_{\psi(-1)=(-1)^{e}}\operatorname{length}_O\big(\text{the group of }\psi\big),
\end{equation}
the sum over the characters $\psi$ of $\mathrm{Gal}(E/\Q)$ with values in $O$. The group of $\psi$ is defined
over $\Q$, with coefficients twisted by $\psi$ (Lean: \texttt{Vlen}, the length of \texttt{Chr.V'} at the
level $\ell N$, $N$ the conductor of $E$). For the other $\ell$ there is a similar inequality with a loss
term (\texttt{HV.redAll}). For the field $E$ of Theorem~G every odd $\ell$ is prime to $[E:\Q]=4n$.

By the duality theorem of Poitou and Tate, $\mathrm{Sha}_\ell(E,e)$ has the same length as the kernel of
$H^2(G_{E,\ell},\Z_\ell(e+1))\to\prod_{v\mid\ell}H^2(G_v,\Z_\ell(e+1))$. The duality is
\cite[\S8.6, Thm.~8.6.8]{NSW00}, and the long exact sequence that contains it is Theorem~8.6.13 there; these
numbers are those of the first edition, which was read. Theorem~8.6.8 is stated there for finitely generated
modules and Theorem~8.6.13 for finite modules. For the statement used here, with the coefficients
$\Q_\ell/\Z_\ell(-e)$ and $\Z_\ell(e+1)$, one passes to limits, and we did not write this out. In the Lean this equality of
lengths is the statement \texttt{PTShaAt} (\texttt{Statement.lean}:101); it is stated under the condition
that $\ell\ne2$ or $E$ has no real place. For $e\ge1$ the group $H^2(G_{E,\ell},\Z_\ell(e+1))$ is finite
(Soul\'e \cite{Sou79}, from memory)
and, for odd $\ell$, isomorphic to $K_{2e}(O_E)\otimes\Z_\ell$, where $K_{2e}$ is Quillen's $K$-group (the
conjecture of Quillen and Lichtenbaum, a consequence of the theorem of Voevodsky and Rost; from memory). Take
$E=\Q$. For odd $e$ the order of $K_{2e}(\Z)\otimes\Z_\ell$ is the power of $\ell$ in the numerator of
$B_{e+1}/(e+1)$: this is the minus part, and it is governed by a value of $\zeta$ (Mazur and Wiles
\cite{MW84}, from memory). For even $e$ the group is the plus part. Vandiver's conjecture says that $\ell$ does not divide
the class number of $\Q(\zeta_\ell+\zeta_\ell^{-1})$. Kurihara \cite{Kur92} relates it to the vanishing
of the groups $K_{4s}(\Z)$, $s\ge1$: his Proposition~3.7 proves one implication, and his Remark~3.11 gives the
converse under the comparison of $K$-theory with cohomology that is used there (compared with the article by
a reviewer, not by us). An equivalence for the
integral groups also needs that comparison and a treatment of the $2$-primary component. We do not use an
equivalence. Corollary~(c) of \S5 proves the sufficient implication that is used here: under Vandiver's
conjecture, $h_{\ell,e}(\Q)=0$ for all odd $\ell$ and all even $e\ge2$.

Theorem~G is weaker than a statement of vanishing in two ways: it bounds a sum weighted by $\log\ell$, and
the bound $\sqrt b=\exp(x^{\varepsilon}/2)$ is very large. It is stronger in one way: the field is $E$, of
degree about $4x^{2+8\varepsilon}$, and the twists go up to $x^{401}$. Theorem~G asserts no vanishing for
the plus part of $E$. For the field $\Q$ some plus parts do vanish: those with $(\ell-1)\mid2e$
(Corollary~(b) of \S5; this is the statement \texttt{PropZ} of the source under the identification
discussed in \S5; for example $\ell=5$, $e=2$), and all of
them under Vandiver's conjecture (Corollary~(c)).

\subsection*{4.\quad The structure of the proof in the source}

What follows was obtained by reading statements, and the proof terms that connect them, downward from
\texttt{vwin\_of\_row}. The long proofs were not read. Sources in the literature are cited as said in the
paragraph ``Citations'' at the head of the section: from memory, unless a theorem or a section is named.

Two dependency computations were made for this section, by the test groups G12 and G19 (files
\texttt{outputs/depcheckG12.out} and \texttt{outputs/depcheckG19.out} in
\texttt{zeta-paper\_tests/displays/G12} and \texttt{G19}). They give the same number of declarations that
the final theorem depends on ($1\,478\,198$) and the same three axioms. The version of this section that
was tested prints $133$ Lean names. G12's run covers $89$ of them: $84$ are among the declarations the final
theorem depends on, and five are not. G19's run covers $35$ of the other $44$: $33$ are among them, and two
are not. The remaining nine are not names of declarations (local names, fields of records and the three
axioms). The seven exceptions are these. \texttt{OneChar},
\texttt{Adm3}, \texttt{Adm2}, \texttt{Sp4In} and \texttt{KatoUse} are in the archive but the final theorem
does not depend on them. It uses variants: \texttt{OneCharU\_s} (\texttt{NumberU3US.lean}:49) in place of
\texttt{OneChar}, with the conditions \texttt{Adm3U\_s} and \texttt{Adm2A} in place of \texttt{Adm3} and
\texttt{Adm2}; and \texttt{SymplecticFourImageDense} in place of \texttt{Sp4In}. That the variants are
used is in G19's run and in a run of the authors, not in G12's. The other two exceptions are \texttt{TieEta}
and \texttt{etaN\_eq\_etaJ}; see \S2. Names first printed in the present version (\texttt{NumTie},
\texttt{etaJ\_join}, \texttt{etaJ\_tie}, \texttt{DefFam.etaJ}) were not tested.

\textbf{The top.} \texttt{vwin\_of\_row} takes one hypothesis, \texttt{Row0b}
(\texttt{WeightsConjBis.lean}:69): \emph{for every field $E$
with complex multiplication, every $s\ge1$, every regular algebraic cuspidal automorphic representation
$\pi$ of $\mathrm{GL}_s$ over $E$ that is conjugate self-dual up to a twist, every $\ell$ and every
isomorphism $\iota\colon\overline{\Q}_\ell\simeq\C$, there is a
continuous semisimple representation of $\mathrm{Gal}(\overline E/E)$ on $\overline{\Q}_\ell^{\,s}$, with
predicted weights and the local comparisons in the record \texttt{PropsB}} (lines 54--67). The record says
the following. The representation is continuous and semisimple. At the primes not above $\ell$ it is
compared with $\pi$ in a weakened sense (\texttt{AwayB}). For each embedding $\tau$ of $E$ a multiset of
integers is given, and it equals the multiset computed from the infinitesimal character of $\pi$; this
equality, and a symmetry under complex conjugation, are stated without a condition. They prescribe the
multiset. They do not say that it is the multiset of the jumps of the representation. The assertions about
the actual jumps are in the part \texttt{AboveB} (lines 39--52). They are made with respect to a local
comparison datum (Lean: \texttt{LocBridge}, \texttt{Pst.Induces}), and in two cases only: where the local
component of $\pi$ is spherical (the representation is then crystalline and de~Rham there, with the predicted
jumps), and in a specified local case (the conditions \texttt{TmW} and \texttt{Sqf} with \texttt{IsQpBig},
or \texttt{Tor}, together with \texttt{mono0}; we did not transcribe them). The classical comparison is
with the theorem stated in the introduction of Chenevier and Harris \cite{CH13} and with the compatibility of
local and global correspondences of Caraiani \cite{Car12}, \cite[Thm.~1.1]{Car14} (these by the reviewer's
report); they build on the work of Harris and Taylor and of Shin (from memory). The normalisations and the hypotheses on the two sides have to
be specified before an exact dictionary can be claimed. We have not done this. One difference is clear from the
Lean itself. The local statement that enters (\texttt{RecTExists}: a map from irreducible representations of
$GL_n$ of a local field to Weil--Deligne representations satisfying the laws \texttt{IsRecT}) is much weaker than
the local Langlands correspondence. The map by which the Lean proves it is built on the supercuspidal representations in
\texttt{cuspEx} and extended to all irreducible representations through segments (\texttt{recTExists}). On a
restricted class of supercuspidal representations (\texttt{GoodCusp}: the characters,
and over $\Q_p$ the twists of representations with a vector fixed by the first congruence subgroup) the map is
given by another construction of the Lean (\texttt{goodMap}; we did not compare it with the correspondence).
To every other supercuspidal representation of $GL_n$ at a place $w$ it assigns one fixed irreducible
parameter of dimension $n$, depending on $w$ and $n$ only, twisted by a real power of the norm; the Lean obtains this parameter
from a lemma it names \texttt{junkEx}. That map satisfies all the laws. So the laws do not single out the
correspondence, and what the Lean calls local Langlands here is a weak form of it. In the Lean \texttt{Row0b} is
a theorem proved in the archive
(\texttt{galoisLocalGlobalStrong\_}\allowbreak\texttt{recParameter}). The theorem \texttt{vwin\_of\_row}
depends only on the three standard axioms of Lean (\texttt{propext}, \texttt{Classical.choice},
\texttt{Quot.sound}), by the dependency computation. So in the Lean nothing on this route is left as an assumption: \texttt{Row0b}, the one hypothesis of
\texttt{vwin\_of\_row}, is itself a theorem of the archive.

All other inputs named ``literature'' on this route are theorems proved in the archive. They are passed
through eight structures whose fields are hypotheses. Each structure is built by a term, as follows.
\begin{enumerate}[label=\textup{(\arabic*)},leftmargin=2.4em,itemsep=1pt,topsep=3pt]
\item \texttt{OpenImageCondition}: one field, \texttt{Row0b}, filled in \texttt{vwin\_of\_row}. The field is
never read: the declarations the final theorem depends on contain no projection and no case analysis of
this structure. In its file 62 theorems take it as an argument and ignore it.
\item \texttt{LiteratureRankNHypothesis}: four fields (local, automorphic, cohomological, complex
multiplication), filled by the next four structures. The last field is not read.
\item \texttt{LitLocAlt}: 18 fields, filled by the term \texttt{local\_inputs\_from\_literature}, each
field by a theorem of the archive; 11 are read.
\item \texttt{LiteratureAutomorphyHypothesis}: 15 fields; 13 are filled by theorems of the archive, one by
the theorem \texttt{baseChangeN}, and one by the hypothesis \texttt{Row0b}; 10 are read.
\item \texttt{LitCoh}: 12 fields, filled by \texttt{galois\_cohomology\_from\_literature}; all are read.
\item \texttt{LiteratureCMHypothesis}: 11 fields, filled by
\texttt{cm\_abelian\_inputs\_from\_literature}; none is read.
\item \texttt{HypothesesH'}: 15 fields with those of its parent; all but one are read. It is filled by
the term \texttt{published\_hypotheses\_from\_patching\_alt}.
\item \texttt{U3.CohInp}: 5 fields, filled by \texttt{Core0.cohInp\_alt}; all are read.
\end{enumerate}
The propositions passed as arguments further down (\texttt{LemmaK}, \texttt{CountI} to \texttt{CountV},
\texttt{CountZ}, \texttt{KatoTwoW}, \texttt{ThmBab}, \texttt{ThmBcFree}, \texttt{PsiBound},
\texttt{FerreroWashington}, \texttt{SymplecticFourImageDense} and others) are each supplied by a theorem
that the final theorem depends on. 

The structure \texttt{OpenLiteratureStatements} is not on this route. It belongs to the other half of
the final theorem, which proves \texttt{LinkN}. It has three fields, and all three are read. Two of
them are filled with the hypothesis \texttt{Row0b}; their names are \texttt{caraianiLocalGlobalDual} and
\texttt{galoisRepOfPolarizedPairLit}. The third is filled with the theorem \texttt{recTExists}; its name
is \texttt{localLanglandsExists}.

\textbf{The decomposition} (\texttt{vvle\_of\_parts}, \texttt{vwin\_of\_counts}). The sum in
\eqref{gc:vvle} is split into the prime $2$, the minus parts and the plus parts for odd $\ell$. Let
$\mathfrak B$ be a finite set of odd primes depending only on the character (Lean: \texttt{Bchi}); for
$\zeta$ the definitions give $\mathfrak B=\{3,7\}$. In the table, $N$ is the conductor of $E$.

\begin{center}\small
\begin{tabular}{@{}p{0.19\linewidth}p{0.43\linewidth}p{0.32\linewidth}@{}}
\toprule
Case & Statement in the Lean, in classical terms & \raggedright Related published work (from memory)\tabularnewline
\midrule
\raggedright $\ell=2$; and $\ell\in\mathfrak B$, plus part
& \raggedright \texttt{ThmBab}: on the window, $\sum_{e<t}h_{\ell,e}\le\exp(x^{4\varepsilon/5})$ for $\ell=2$ over $E$,
and for $\ell\in\mathfrak B$ over bounded layers above a fixed field (layers in the sense of ``The window'' below)
& \raggedright Ferrero and Washington \cite{FW79} ($\mu=0$); Iwasawa's theory of $\Z_\ell$-extensions; Soul\'e \cite{Sou03}\tabularnewline
\raggedright Minus part, $\ell$ odd
& \raggedright \texttt{ThmBcFree}, for $n\ge4$ and $t\ge2$: $\sum_\ell\log\ell\sum_{e<t}h^-_{\ell,e}\le3\,[E:\Q]\,t^2\log(tN)$,
plus a term for the odd primes dividing the degree of the field of the character $\psi$ of the source
(Lean: \texttt{M.deg}; none for $\zeta$). From \texttt{PsiBound}: the length of the part of $\psi$
is at most the valuation of $N\prod_{p\le e+2}p\cdot B_{e+1,\psi}$, if $\ell\nmid[E:\Q]$ and
$\psi(-1)=(-1)^{e+1}$; the hypotheses are stated after the table
& \raggedright Stickelberger; Mazur and Wiles \cite{MW84}; Euler systems of Gauss sums (Kolyvagin; Rubin \cite[\S2.2]{Rub00})\tabularnewline
\raggedright Plus part, $\ell\notin\mathfrak B$, $(\ell-1)\nmid2e$ (Case I)
& \raggedright \texttt{CountI}, from \texttt{OneCharU\_s}: under the hypotheses stated after the table, the length is at most
$2\sum_{\rho\in\mathcal R}v_\ell\big(\operatorname{tr}\rho(\sigma)-\Psi(\sigma)-1-\Psi(\sigma)^{-1}\big)$,
$\Psi=\psi\chi_\ell^{e}$, for a finite set $\mathcal R$ of three-dimensional representations $\rho$ of the
Galois group of an imaginary quadratic field
& \raggedright congruences with Eisenstein series (Ribet; Mazur and Wiles; Bella\"iche and Chenevier \cite{BC04});
modularity lifting (Taylor and Wiles); Kato \cite{Kat04}\tabularnewline
\raggedright Plus part, $\ell\notin\mathfrak B$, $(\ell-1)\mid2e$, $\psi=1$ (Case Z)
& \raggedright \texttt{PropZ}: for $\ell\ge5$ and even $e\ge2$ the group over $\Q$ vanishes
& \raggedright class numbers of $\Q$ and $\Q(\sqrt\ell)$; see \S5, Corollary (b), and the identification discussed there\tabularnewline
\raggedright Plus part, the remaining $\psi$ with $(\ell-1)\mid2e$ (Cases II, III)
& \raggedright \texttt{CountII}, \texttt{CountIII}: the length times $\log\ell$ is at most $(\ell Ne)^{C}$; from a statement on
automorphic forms on $\mathrm{GSp}_4$ (\texttt{SymplecticFourImageDense}). Case III is empty for $\zeta$
& \raggedright not identified\tabularnewline
\raggedright Plus part, $\ell\in\mathfrak B$ (Cases II, IV, V)
& \raggedright \texttt{SumIn}: reduction to the layers of the first row
& \raggedright as in the first row\tabularnewline
\bottomrule
\end{tabular}
\end{center}

The last column names related work. It does not say that the statement of the Lean is a published theorem;
for that question see the three lists at the end of the section. \texttt{CountV} has the shape of
\texttt{CountII}, for Case~V with a further condition that we did not transcribe.

\textbf{The row for the minus part, with its hypotheses} (\texttt{PsiBound},
\texttt{ZeroMemUnrHom.lean}:97) \quoted. Let $E$ be an abelian field with complex multiplication, and $N\ge1$
a level with a homomorphism $(\Z/N)^\times\to\mathrm{Gal}(E/\Q)$ compatible with the action on roots of unity
(Lean: \texttt{IsCycAct}). Let $\ell$ be an odd prime with $\ell\nmid[E:\Q]$. Let $O$ be a discrete valuation
ring, finite over $\Z_\ell$, inside an algebraic closure of $\Q_\ell$, and let $v_O$ be its additive
valuation, which is $1$ on a generator of the maximal ideal. Let $e\ge1$ and let $\psi$ be a character of
$\mathrm{Gal}(E/\Q)$ with values in $O^\times$ and $\psi(-1)=(-1)^{e+1}$. Then
\begin{equation}\label{D:339}
 \operatorname{length}_O\big(\text{the group of }\psi\big)\ \le\ v_O\Big(N\prod_{p\le e+2}p\cdot B_{e+1,\psi}\Big),
\end{equation}
with the Bernoulli number of the primitive character. Both sides are in units of length over $O$; on $\Z_\ell$
we have $v_O=e_O\,v_\ell$. So \texttt{PsiBound} is an upper bound for a cohomology length in terms of the
valuation of a Bernoulli number with explicit integral factors. It is not a bound for the valuation. For
$\psi=1$ the field $E$ is still a field with complex multiplication, so the statement does not apply to
$E=\Q$. To use it over $\Q$ one needs a descent, for example from the part of the trivial character over
$\Q(i)$; we did not carry this out. Rubin \cite[\S2.2]{Rub00} gives general bounds from Euler systems under
explicit hypotheses (by the reviewer's report). The specialisation to the bound by a Bernoulli number is not supplied here.

\textbf{The row for Case I, with its hypotheses} (\texttt{OneCharU\_s}, \texttt{NumberU3US.lean}:49) \quoted. Let $\mathcal K$ be an
imaginary quadratic field, $\ell$ a prime, $e$ an integer, $N$ a level and $\psi$ a Dirichlet character
modulo $N$ with values in the coefficient ring. Assume the twelve conditions of the record \texttt{GD}
(\texttt{Counts.lean}:9); among them: $\ell\ne2$; $e\ge1$; $N$ is even; $\ell$ does not divide the order of
$\psi$; $\ell$ does not divide the discriminant of $\mathcal K$ times its class number; the discriminant of
$\mathcal K$ is prime to $N$; $\psi(-1)=(-1)^e$; and $(\ell-1)\nmid2e$. Fix an isomorphism
$\iota\colon\overline{\Q}_\ell\simeq\C$. Call a three-dimensional
representation $\rho$ of $\mathrm{Gal}(\overline{\mathcal K}/\mathcal K)$ over $\overline{\Q}_\ell$
\emph{admissible} if it satisfies
\begin{equation}\label{D:340}
 \texttt{Adm3U\_s}(\iota,e,N,\rho)\qquad\text{or}\qquad\texttt{Adm2A}\big(\iota,e,N\,|\mathrm{disc}\,\mathcal K|,\rho\big).
\end{equation}
The same $\iota$, $e$ and levels are used in the hypothesis and in the conclusion. Let
$\sigma\in\mathrm{Gal}(\overline{\mathcal K}/\mathcal K)$ be such that
$\operatorname{tr}\rho(\sigma)\ne\Psi(\sigma)+1+\Psi(\sigma)^{-1}$ for \emph{every} admissible $\rho$. Then
the length of the group of $\psi$ is finite, and there is a finite set $\mathcal R$ of admissible
representations, no two of them isomorphic, such that the normalised length is at most
$2\sum_{\rho\in\mathcal R}v_\ell\big(\operatorname{tr}\rho(\sigma)-\Psi(\sigma)-1-\Psi(\sigma)^{-1}\big)$.
The condition on $\sigma$ makes every term of the sum finite. For the statement \texttt{OneChar}, which
has the same shape, the conditions \texttt{Adm3} and \texttt{Adm2} say in classical terms: $\rho$ is
irreducible, automorphic, conjugate self-dual, of weights $(e,0,-e)$; or $\rho=W\oplus1$ with $W$ attached to
a newform of weight $2e+1$. Whether \texttt{Adm3U\_s} and \texttt{Adm2A} differ from them in substance was
not examined. Bella\"iche and Chenevier \cite[Thm.~1.1]{BC04} construct an extension in a Selmer group for
their characters (by the reviewer's report). This illustrates the method of congruences. A uniform upper bound for a length is a
further statement.

\textbf{Further inputs, all proved in the archive.} \texttt{LemmaK} (\texttt{LemmaK.lean}:13): for
every prime $\ell$, every totally complex field $L$
containing the $\ell$th roots of unity (the fourth roots if $\ell=2$) and every $e\ge1$, the group
$H^2(G_{L,\ell},\Z_\ell(e+1))$ is finite and the logarithm of its order is at most
$\exp\big((2e[L:\Q])^{C}\log(2|\mathrm{disc}\,L|)\big)$. So \texttt{LemmaK} bounds the whole group, for
every $\ell$, and it includes $\ell=2$. Finiteness is due to Soul\'e (\cite{Sou79}, from memory). A
quantitative bound away from small primes is in \cite{Sou03}: Theorem~2 there discards the torsion at small
primes, and Theorem~3 bounds the groups $H^2(O_L[1/p],\Z_p(n'))$ in a product over the primes $p\ge4n'-1$,
$p\ge[L:\Q]+2$ only. Here $n'$ is the twist, which Soul\'e writes $n$; it is not the $n$ of this section.
These two statements of \cite{Sou03} were compared with the article by a reviewer, not by us. We have not
identified a published estimate for \texttt{LemmaK} in its full small-prime range. The proof of
\texttt{LemmaK} in the Lean names a statement on the rank of the homology
of $\mathrm{GL}_s(O_L)$ (\texttt{BOdd}; Borel's theorem). \texttt{FerreroWashington}: for an abelian field
$L$ containing those roots of unity, the $\ell$-rank of the class group of $L(\zeta_{\ell^{r}})$ is bounded
independently of $r$; it is proved in the archive. In the cyclotomic $\Z_\ell$-tower, a bounded $\ell$-rank of
the class groups is equivalent to $\mu=0$ \cite[Prop.~13.23]{Was82} (first edition, p.~284); so this is the
theorem of \cite{FW79} (from memory) in its form with bounded rank. The theorems of Poitou and Tate, Tate's local duality and
Euler characteristic, the formula of Greenberg and Wiles, and the local theorems of Bloch and Kato
are collected in the structure \texttt{LitCoh}. \texttt{KatoTwoW} (\texttt{KatoTwoW.lean}:10--14): let $W$ be
a two-dimensional representation of $\mathrm{Gal}(\overline{\Q}/\Q)$ over $\overline{\Q}_\ell$ and $\theta$ a
character, such that the twist of $W$ by $\chi_\ell^{-e}$ is attached to a newform of weight $2e+1$
(the attachment hypothesis), and assume the hypotheses \texttt{Used} and \texttt{WitAQ} of the source, which
we did not transcribe. Then the two rational Selmer spaces of Bloch and Kato, of $W$ and of
$W\otimes\theta$, vanish: every cocycle with values in the vector space over $\overline{\Q}_\ell$ that
satisfies the local conditions is a coboundary. Kato \cite[\S14.1 and Thm.~14.2(1)]{Kat04} proves that the
discrete Selmer groups are finite at the critical twists other than the central one (by the reviewer's
report). The two statements are
of different kinds: the Lean statement is about vector spaces, and it bounds no integral length. In our
reading the relevant point is $s=e+1$, on the edge of the critical strip.

\textbf{The window.} The statement \texttt{ABound} (\texttt{SecondBound.lean}:19--25; line 25 is its second
clause) explains it. Let
$\ell=2$ or $\ell\in\mathfrak B$, let $L_0$ be a fixed base field and $R\ge0$ an integer. Put
\begin{equation}\label{D:341}
 j_0=3+\lfloor\log_\ell\max(1,R)\rfloor
\end{equation}
(Lean: \texttt{jZero}; the logarithm of the Lean is $0$ at $R=0$). The first clause
of \texttt{ABound} says: there are an integer $c_0\ge0$ such that $L_0$ contains a primitive root of unity
of order $\ell^{c_0}$, an integer $R'$ with $0\le R'\le R$, and points $s_1,\dots,s_{R'}$ of a field containing
$\Q_\ell$, with a non-archimedean absolute value $|\cdot|_\ell$ extending that of $\Q_\ell$, with this
property. We call the $s_i$ the \emph{exceptional centres}. Let $e\ge1$, let $\vartheta\ge0$ be an integer,
and let $L_1$ be a \emph{layer of degree $\ell^{\vartheta}$ over $L_0$} in the sense of the Lean
(\texttt{IsLayer}, \texttt{StmtFields.lean}:47): a number field $L_1=L_0(\zeta_{\ell^{k}})$ for some $k\ge0$,
with $[L_1:L_0]=\ell^{\vartheta}$. Let $h_{\ell,e}$ be the length \eqref{gc:sha} for the field $L_1$. Then
for every integer $\varrho\ge0$
\begin{equation}\label{gc:abound}
 h_{\ell,e}\ \le\ R'\,\big(\max(\vartheta,\,j_0)+c_0+\varrho\big)
 \qquad\text{whenever } |s_i-e|_\ell\ge\ell^{-\varrho}\ \text{for all } i .
\end{equation}
If $L_0$ contains the $\ell$th roots of unity (the fourth roots if $\ell=2$), a layer in this sense is
the field of degree $\ell^{\vartheta}$ over $L_0$ in the cyclotomic $\Z_\ell$-extension of $L_0$; this is
the case in the uses \texttt{TwoIn} and \texttt{FixedIn}. The statement \texttt{ABound} itself does not
assume it, since $c_0$ may be $0$. Without it a layer need not lie in the cyclotomic $\Z_\ell$-extension:
for $\ell=2$ and $L_0=\Q$ the field $\Q(i)$ is a layer with $\vartheta=1$, and it is not in the cyclotomic
$\Z_2$-extension of $\Q$.
A second clause (\texttt{SecondBound}) compares a layer of degree $\ell^{\vartheta}$ with a layer of
degree $\ell^{j_0}$, with a loss $R\cdot\max(\vartheta-j_0,\,0)$; we do
not transcribe it. By \eqref{gc:abound} the lengths are small unless $e$ is $\ell$-adically very close to
one of at most $R$ exceptional centres. The statement defines its conclusion by the distances from $e$ to
the centres. It does not say what the centres are. Iwasawa's theory gives a statement of this shape when
$\mu=0$: a module with $\lambda\le R$, and the zeros of its characteristic power series in the place of the
centres. The Lean statement does not identify the centres with such zeros in the coordinate $e$, and we did
not check this in the proof. For the exceptional $e$ only \texttt{LemmaK} is available, and it is useful
only if $e$ is at most a small power of $x$. The window is an interval of $x$ in which no exceptional $e$
lies between that small power and $x^{401}$; the number $J$ of Theorem~G grows with $R$. This paragraph is
our reading of the statements \texttt{ABound}, \texttt{FirstBound} and \texttt{ThmBab}; the proof was not read.

\textbf{Three claims of earlier notes, compared with the Lean statements.}
\emph{(1) ``The minus part is a Stickelberger index bound.''} The statement is a bound of a length by the
valuation of a generalised Bernoulli number, as in the table; summed over $\ell$ it is bounded by
$[E:\Q]\log\big(16N^{e+2}(e+1)!\big)$ (\texttt{PrimeBound}). Its proof passes through a reciprocity
statement (\texttt{Rec}) deduced from the duality of Poitou and Tate, and the files that prove it are named
after Kolyvagin's derivative construction and Jacobi sums. A cruder bound of Stickelberger type
(\texttt{StickBound}) is used only for $\ell\mid[E:\Q]$.
\emph{(2) ``The plus part uses congruences on $U(3)$ in the manner of Ribet.''} Confirmed for Case~I at the
level of the statement \texttt{OneCharU\_s}. The inequality bounds a cohomology group from above by a congruence
number. Ribet's method produces classes; an upper bound of this kind needs in addition a comparison of
a deformation ring with a Hecke algebra, and the Lean names such inputs (\texttt{adjoint\_lifting\_of\_cohomology},
\texttt{published\_hypotheses\_from\_patching\_alt}).
\emph{(3) ``Prismatic $F$-crystals are needed because the local fields at $2$ have unbounded degree.''}
The archive contains the statement \texttt{BhattScholzeIntegral}: every lattice in a crystalline
representation of the Galois group of any finite extension of $\Q_\ell$ comes from a prismatic $F$-crystal.
Bhatt and Scholze \cite[Thms.~1.2 and~5.6]{BS23} give the equivalence of prismatic $F$-crystals with the
finite free $\Z_p$-lattices whose rationalisations are crystalline (by the reviewer's report). By the dependency computation it belongs to the half that
proves \texttt{LinkN}, not to the proof of
Theorem~G. The completion of $E$ at a prime above $2$ has ramification index $n$ and
residue degree $2$, so its degree $2n$ is unbounded; the theory of Fontaine and Laffaille needs an
unramified field and small weights. So the claim is consistent with the statements. That the result is
\emph{needed} cannot be read off from statements.

\subsection*{5.\quad The simplest case: the field $\Q$, and $\Q(i)$}

\begin{illbox}
The whole of this subsection, to the end of the computation for $\Q(i)$, is an illustration: the field
$\Q$ and one twist over $\Q(i)$, where the groups can be computed from class numbers. The proof uses the
general bound of Theorem~G, with the quoted comparison inequality, in place of these special computations.
\end{illbox}

It is not used by the source. Let $\ell$ be an odd prime. Let $A$ be the
$\ell$-primary part of the class group of $\Q(\zeta_\ell)$, let $\Pi=\mathrm{Gal}(\Q(\zeta_\ell)/\Q)$, let
$\omega\colon\Pi\to\Z_\ell^\times$ be the character with $\omega(\varpi)\equiv\alpha\pmod\ell$ when
$\varpi\zeta_\ell=\zeta_\ell^{\alpha}$, and let $A^{(i)}$ be the subgroup of $A$ on which $\Pi$ acts by
$\omega^{i}$. References to \cite{Was82} are to the first edition, which we consulted; there $A^{(i)}$ is
written $A_i$.

\textbf{Proposition G3.} \emph{Let $e\ne0$. If $A^{(i)}=0$ for $i\equiv-e\pmod{\ell-1}$, then
$\mathrm{Sha}_\ell(\Q,e)=0$.}

\emph{Proof.} Let $F_\infty$ be the field generated over $\Q$ by all roots of unity of $\ell$-power order,
and let $H_\infty\subset G_{\Q,\ell}$ be the subgroup fixing
$F_\infty$; then $\chi_\ell$ identifies $G_{\Q,\ell}/H_\infty$ with $\Z_\ell^\times$, and $H_\infty$ acts trivially on $W_e$.
The sequence of inflation and restriction reads
\begin{equation}\label{D:343}
 0\to H^1(\Z_\ell^\times,W_e)\to H^1(G_{\Q,\ell},W_e)\to\mathrm{Hom}(H_\infty,W_e).
\end{equation}
The group $\Z_\ell^\times$ has a topological generator $\gamma_0$, so $H^1(\Z_\ell^\times,W_e)=W_e/(\gamma_0^{-e}-1)W_e$.
This is $0$: $\gamma_0^{-e}-1$ is a nonzero element of $\Z_\ell$ and $\Q_\ell/\Z_\ell$ is divisible. So a class $z$
in $\mathrm{Sha}_\ell(\Q,e)$ is determined by its restriction $\phi\colon H_\infty\to W_e$, a continuous
homomorphism with $\phi(\sigma\sigma'\sigma^{-1})=\chi_\ell(\sigma)^{-e}\phi(\sigma')$ for $\sigma\in G_{\Q,\ell}$,
$\sigma'\in H_\infty$. Since $z$ is trivial on a decomposition group
$G_\ell$ at $\ell$, $\phi$ vanishes on $G_\ell\cap H_\infty$, and by the last formula on its conjugates. In particular $\phi$ is
unramified at the primes above $\ell$; it is unramified elsewhere by definition of $G_{\Q,\ell}$. So $\phi$
factors through $X_\infty$, the Galois group over $F_\infty$ of its largest unramified abelian extension of
$\ell$-power order. Restricted to $\Pi\subset\Z_\ell^\times$ the formula says
$\phi(\varpi\cdot z')=\omega(\varpi)^{-e}\phi(z')$, so $\phi$ factors through the part $X_\infty^{(i)}$ on
which $\Pi$ acts by $\omega^{i}$, $i\equiv-e$. Let $\Gamma=\mathrm{Gal}(F_\infty/\Q(\zeta_\ell))$, with
generator $\gamma$. One prime of $\Q(\zeta_\ell)$ ramifies in $F_\infty$, and it ramifies totally; hence
$X_\infty/(\gamma-1)X_\infty\cong A$ \cite[Prop.~13.22, at the bottom level]{Was82}. The proposition does not
state that the isomorphism commutes with $\Pi$. It does. Let $H$ be the largest unramified abelian extension
of $\Q(\zeta_\ell)$ of $\ell$-power degree. The isomorphism is the restriction
$X_\infty\to\mathrm{Gal}(H/\Q(\zeta_\ell))$, which has kernel $(\gamma-1)X_\infty$, followed by the inverse
of the Artin map $A\to\mathrm{Gal}(H/\Q(\zeta_\ell))$. An element $\varpi$ of
$\mathrm{Gal}(F_\infty/\Q)=\Z_\ell^\times$ acts on both Galois groups by conjugation with a lift of
$\varpi$, and restriction commutes with this. The Artin map sends $\varpi\mathfrak a$ to the conjugate by
$\varpi$ of the image of $\mathfrak a$. So
$X^{(i)}_\infty/(\gamma-1)=A^{(i)}=0$, and $X^{(i)}_\infty=0$ by Nakayama's lemma for compact modules over
$\Z_\ell[[\Gamma]]$ \cite[Lemma~13.16]{Was82}. Hence $\phi=0$.\qed

\textbf{Lemma G4.} \emph{For every prime $\ell$ the class number of $\Q(\sqrt\ell)$ is less than $\ell$.}

\emph{Proof.} Put $F=\Q(\sqrt\ell)$ and write $h$ for its class number. Let $\mathfrak d$ be the
discriminant of $F$: $\mathfrak d=\ell$ if $\ell\equiv1\pmod4$, $\mathfrak d=4\ell$ if $\ell\equiv3\pmod4$,
and $\mathfrak d=8$ if $\ell=2$. In each case $\tfrac12\sqrt{\mathfrak d}\le\sqrt\ell$. Write $\alpha'$ for
the conjugate of $\alpha\in F$, and $N(\mathfrak a)$ for the norm of an ideal.

\emph{Minkowski's bound: every ideal class of $F$ contains an integral ideal of norm at most
$\tfrac12\sqrt{\mathfrak d}$.} This is \cite[Ch.~I, \S6, Exercise~3]{Neu99} for a field of degree $2$ with two
real places. It is an exercise there, so we give the proof from two statements of the text, (5.2) and
(4.4). Let
$\mathfrak a$ be a nonzero integral ideal. Its image under $\alpha\mapsto(\alpha,\alpha')$ is a lattice in
$\mathbb R^2$ of covolume $\sqrt{\mathfrak d}\,N(\mathfrak a)$ \cite[Ch.~I, (5.2)]{Neu99}. For $T>0$ the set
$|x|+|y|\le T$ is convex and symmetric, of area $2T^2$. Minkowski's lattice point theorem
\cite[Ch.~I, (4.4)]{Neu99} says that it contains a nonzero point of the lattice if
$2T^2>4\sqrt{\mathfrak d}\,N(\mathfrak a)$. The inequality is strict there. (The case of equality for
compact sets is the last sentence of Exercise~2 of \S4 there, without proof; we do not use it.) Take
$T^2=2\sqrt{\mathfrak d}\,N(\mathfrak a)$. For every $T'>T$ the theorem gives a nonzero lattice point with
$|x|+|y|\le T'$. The lattice has only finitely many points with $|x|+|y|\le 2T$. So one of them, nonzero,
serves for a sequence of numbers $T'$ decreasing to $T$. We get
$\alpha\in\mathfrak a$, $\alpha\ne0$, with $|\alpha|+|\alpha'|\le T$, hence
\begin{equation}\label{D:344}
 |\alpha\alpha'|\ \le\ \tfrac14\big(|\alpha|+|\alpha'|\big)^2\ \le\ \tfrac14T^2=\tfrac12\sqrt{\mathfrak d}\,N(\mathfrak a).
\end{equation}
Now let an ideal class be given, and let $\mathfrak a$ be an integral ideal in the inverse class. The ideal
$(\alpha)\mathfrak a^{-1}$ is integral, because $\alpha\in\mathfrak a$. It lies in the given class, and its
norm is $|\alpha\alpha'|/N(\mathfrak a)\le\tfrac12\sqrt{\mathfrak d}$.

\emph{Counting: for $j\ge1$, the number of integral ideals of $F$ of norm $j$ is at most $j$.} An ideal of norm
$j$ is in one way a product, over the primes $p$ dividing $j$, of ideals whose norms are the powers of $p$
in $j$. So let $j=p^{k}$. If $p$ splits, $(p)=\mathfrak p\mathfrak p'$, the ideals of norm $p^k$ are the
$k+1$ ideals $\mathfrak p^{k_1}\mathfrak p'^{\,k_2}$ with $k_1+k_2=k$. If $p$ is inert, there is one such ideal
for even $k$ and none for odd $k$. If $p$ ramifies, $(p)=\mathfrak p^2$, there is one, $\mathfrak p^{k}$. In
each case the number is at most $k+1$, the number of divisors of $p^k$. Hence the number of ideals of norm
$j$ is at most the number of positive divisors of $j$, which is at most $j$.

\emph{Conclusion.} Every ideal class contains an integral ideal of norm $j\le\tfrac12\sqrt{\mathfrak d}\le\sqrt\ell$.
The sum of the integers $j$ with $1\le j\le\sqrt\ell$ is at most $\tfrac12\sqrt\ell\,(\sqrt\ell+1)$. So
\begin{equation}\label{D:345}
 h\ \le\ \sum_{1\le j\le\sqrt\ell}j\ \le\ \frac{\sqrt\ell\,(\sqrt\ell+1)}2=\frac{\ell+\sqrt\ell}2\ <\ \ell ,
\end{equation}
because $\sqrt\ell<\ell$ for $\ell>1$. This holds for every prime $\ell$, including $2$ and $3$.\qed

\emph{Numerical check.} The test group G12 computed the class numbers of
$\Q(\sqrt\ell)$ for all $17\,984$ primes $\ell<2\cdot10^5$ by cycles of reduced forms (script
\texttt{t04\_classno.py} in \texttt{zeta-paper\_tests/displays/G12}): $h<\sqrt\ell$ in every case. The test
group G19 reproduced this with PARI (the largest class number is $153$, at $\ell=176401$; PARI's values
were not certified), and for the $9592$ primes $\ell<10^5$ it checked every inequality of the proof.

\textbf{Corollaries.} \emph{\textup{(a)} Let $e\ge1$ be odd. If $\ell$ does not divide the numerator of
$B_{e+1}/(e+1)$, then $\mathrm{Sha}_\ell(\Q,e)=0$. \textup{(b)} Let $e\ge2$ be even and $(\ell-1)\mid2e$. Then
$\mathrm{Sha}_\ell(\Q,e)=0$. \textup{(c)} Let $e\ge2$ be even. If $\ell$ does not divide the class number of
$\Q(\zeta_\ell+\zeta_\ell^{-1})$, then $\mathrm{Sha}_\ell(\Q,e)=0$.}

\emph{Proof.} We use this fact: if $F\subset\Q(\zeta_\ell)$ is a subfield, then the $\ell$-part of the class
group of $F$ is the part of $A$ fixed by $\mathrm{Gal}(\Q(\zeta_\ell)/F)$. Indeed, write
$L=\Q(\zeta_\ell)$ in this proof, let $j\colon\mathrm{Cl}(F)\to\mathrm{Cl}(L)$ be the map that extends
ideals, and $N\colon\mathrm{Cl}(L)\to\mathrm{Cl}(F)$ the norm. Then
\begin{equation}\label{D:346}
 N\circ j=[L:F]\ \text{ on }\mathrm{Cl}(F),\qquad j\circ N=\sum_{\sigma\in\mathrm{Gal}(L/F)}\sigma\ \text{ on }\mathrm{Cl}(L).
\end{equation}
The degree $[L:F]$ is prime to $\ell$. By the first identity $j$ is injective on the $\ell$-part. By the
second, a class $z$ of $A$ fixed by $\mathrm{Gal}(L/F)$ satisfies $j(N(z))=[L:F]\,z$, so it lies in the
image of $j$; and the image of $j$ is fixed by $\mathrm{Gal}(L/F)$.
(a) Here $i\equiv-e$ is odd. If $(\ell-1)\mid(e+1)$ then $i\equiv1$ and $A^{(1)}=0$
\cite[Prop.~6.16]{Was82}. Otherwise take $3\le i\le\ell-2$. By Herbrand's theorem \cite[Thm.~6.17]{Was82},
$A^{(i)}=0$ unless $\ell\mid B_{\ell-i}$; and $B_{\ell-i}/(\ell-i)\equiv B_{e+1}/(e+1)\pmod\ell$
by Kummer's congruence \cite[Cor.~5.14]{Was82}, since $\ell-i\equiv e+1\not\equiv0\pmod{\ell-1}$. (b) If
$(\ell-1)\mid e$ then $i\equiv0$ and
$A^{(0)}$ is the $\ell$-part of the class group of $\Q$, by the fact above. Otherwise $(\ell-1)/2$ divides
$e$ and is even, so $\ell\equiv1\pmod4$ and $i\equiv(\ell-1)/2$. The part of $A$ fixed by the squares in
$\Pi$ is $A^{(0)}\oplus A^{((\ell-1)/2)}$, and the field fixed by the squares is $\Q(\sqrt\ell)$. So
$A^{((\ell-1)/2)}$ is the $\ell$-part of the class group of $\Q(\sqrt\ell)$. It is $0$, because the class
number is less than $\ell$ (Lemma~G4). (c) For even $e$, $A^{(i)}$ lies in the part of $A$ fixed by complex
conjugation, which is the $\ell$-part of the class group of the real subfield.\qed

Kummer's congruence, as used in (a), was also tested numerically by the test group G12:
$B_r/r\equiv B_{r'}/r'\pmod\ell$ for the $2543$ triples $(\ell,r,r')$ with $\ell$ prime, $5\le\ell<200$,
$r$ and $r'$ even, $2\le r\le\ell-3$, $r<r'\le220$ and $r\equiv r'\pmod{\ell-1}$; in $14$ of them $\ell$
divides $B_r$ (script \texttt{t03\_bernoulli.py}, part G, in \texttt{zeta-paper\_tests/displays/G12}). The
test group G19 dropped the condition $r\le\ell-3$ and required $r\not\equiv0\pmod{\ell-1}$ instead: the
congruence holds for all $9738$ such triples.

\emph{What is proved here, and what the Lean states.} Everything proved in this subsection is about the
groups $\mathrm{Sha}_\ell(E,e)$ of \eqref{gc:sha}. Two identifications with the Lean are assumed and not
proved. The first is the transcription of \S1: that \texttt{Sha} and \texttt{shaLen} of the Lean are the
group and the length of \eqref{gc:sha}. It was compared at the level of the definitions (the same set
$S_\ell$, the same twist $\chi_\ell^{-e}$, discrete coefficients $\Q_\ell/\Z_\ell$, continuous cocycles, the
kernel of the restrictions to chosen decomposition groups; \texttt{shaLen\_eq\_}\allowbreak\texttt{shaInv}
removes the choices). Under it, the vanishing of $\mathrm{Sha}_5(\Q(i),4)$ proved below says that
\texttt{shaLen} is $0$ for $\Q(i)$, $\ell=5$, $e=4$. The second concerns \texttt{PropZ}, which says: for
$\ell\ge5$ and even $e\ge2$ with $(\ell-1)\mid2e$, the number \texttt{Chr.V} of the trivial character is $0$.
Corollary (b) is, for $\ell\ge5$, this statement, if the group \texttt{Chr.V}
of the Lean is identified with $\mathrm{Sha}_\ell(\Q,e)$ with enlarged coefficients. At the level of the
definitions this is so: \texttt{Chr.V} is the length of the group of classes that vanish on the
decomposition groups above $\ell$ and on all inertia groups (\texttt{ShaFamMono.lean}:30--32). The cohomology of the Lean
is continuous cohomology: \texttt{H1U} is defined by restricting the representation to the subgroup
(\texttt{RestrictionMemInvariantH1.lean}:35), and the group \texttt{H1} is the quotient of the continuous
cocycles by the coboundaries (\texttt{SeparatedDvrSecondCohomologyMem.lean}:42--66). A full identification of
the coefficients and of the local conditions of the Lean with those of \eqref{gc:sha} was not checked. So
Corollary (b) is a proof of the classical statement, and it is a proof of \texttt{PropZ} only under this
identification. The source proves \texttt{PropZ} independently (\texttt{X.propZ}). Corollary (a) is proved from Herbrand's
theorem and does not use \texttt{PsiBound}. As said in \S4, \texttt{PsiBound} with $\psi=1$ needs a field
with complex multiplication, so it gives a statement over $\Q$ only after a descent. The precise length over
$\Q$, $v_\ell(B_{e+1}/(e+1))$ for $(\ell-1)\nmid(e+1)$, is the theorem of Mazur and
Wiles (\cite{MW84}, from memory).

\textbf{The field $\Q(i)$.} For odd $\ell$, $\mathrm{Sha}_\ell(\Q(i),e)$ is the sum of the part on which
complex conjugation $\jmath$ acts as $+1$ and the part on which it acts as $-1$. The first is
$\mathrm{Sha}_\ell(\Q,e)$, and the second is the part of the character $\chi_4$ modulo $4$
(restriction and corestriction; not proved here). For odd $e$ the first is the minus part and the second
the plus part; for even $e$ it is the other way round. The minus part is governed by $B_{e+1}$ ($e$ odd)
and by $B_{e+1,\chi_4}$ ($e$ even), where
$B_{r,\chi}=f_\chi^{\,r-1}\sum_{\alpha=1}^{f_\chi}\chi(\alpha)B_r(\alpha/f_\chi)$, $f_\chi$ the conductor and
$B_r(X)$ the Bernoulli polynomial (Lean: \texttt{genB}). With this range of summation $B_{1,1}=+\frac12$, as
in the Lean. Exact computation gives the odd primes dividing the numerator of $B_{e+1,\chi}/(e+1)$:

\begin{center}\small
\begin{tabular}{@{}rll@{\qquad}rll@{}}
\toprule
$e$ & $\chi$ & odd prime factors & $e$ & $\chi$ & odd prime factors\\
\midrule
1, 3, 5, 7, 9 & $1$ & none & 2 & $\chi_4$ & none\\
11 & $1$ & $691$ & 4 & $\chi_4$ & $5$\\
13 & $1$ & none & 6 & $\chi_4$ & $61$\\
15 & $1$ & $3617$ & 8 & $\chi_4$ & $5\cdot277$\\
17 & $1$ & $43867$ & 10 & $\chi_4$ & $19\cdot2659$\\
19 & $1$ & $283\cdot617$ & 12 & $\chi_4$ & $5\cdot13\cdot43\cdot967$\\
21 & $1$ & $131\cdot593$ & 14 & $\chi_4$ & $47\cdot4241723$\\
23 & $1$ & $103\cdot2294797$ &&&\\
25 & $1$ & $657931$ &&&\\
27 & $1$ & $9349\cdot362903$ &&&\\
29 & $1$ & $1721\cdot1001259881$ &&&\\
31 & $1$ & $37\cdot683\cdot305065927$ &&&\\
\bottomrule
\end{tabular}
\end{center}

By Corollary (a), $\mathrm{Sha}_\ell(\Q,e)=0$ for every odd $\ell$ when $e\in\{1,3,5,7,9,13\}$, and for
every odd $\ell\ne691$ when $e=11$. Grant the description of the two parts given above. Then for $E=\Q(i)$
and $t=10$ the minus part of the sum in \eqref{gc:vvle}
has no contribution from odd $e$; for even $e\le8$ and odd $\ell$ the statement \texttt{PsiBound} of the source
bounds the contribution of $\ell$ by $v_\ell\big(4\prod_{p\le e+2}p\cdot B_{e+1,\chi_4}\big)$.

\textbf{A prime that divides the Bernoulli number while the group is zero.} Take $e=4$ and $\ell=5$. The
prime $5$ divides the numerator of $B_{5,\chi_4}/5=-\frac52$. The reason is an Euler factor. Let
$L_5(s,\chi_4\omega)$ be the $5$-adic $L$-function. For $r\equiv1\pmod4$ its value at $1-r$ is
$-(1-\chi_4(5)5^{r-1})B_{r,\chi_4}/r$ \cite[Thm.~5.11]{Was82}. At $r=1$ this is $0$, because $\chi_4(5)=1$.
All its values at integers are congruent modulo $5$ \cite[Cor.~5.13]{Was82}. So $5$ divides the numerator of
$B_{e+1,\chi_4}/(e+1)$ for every $e\equiv0\pmod4$, $e\ge4$. Consistently, the computation finds the factor
$5$ for every $e\equiv0\pmod 4$ with $e\le24$ and for no $e\equiv2\pmod4$ in that range.

The group $\mathrm{Sha}_5(\Q(i),4)$ is nevertheless $0$; in particular its part on which $\jmath$ acts as
$-1$ is $0$. We prove this. The proof uses local and global class field theory in their standard form and
one computation modulo $25$. It is for this field and this twist only. Put $E=\Q(i)$, $G=G_{E,5}$,
$W=W_4=\Q_5/\Z_5(-4)$ and $\chi=\chi_5$. For a subgroup $H$ write $\mathrm{Hom}(H,\mathbb F_5)$ for the continuous
homomorphisms.

\emph{Step 1: the invariants.} The prime $5$ splits in $E$. So $S_5=\{v,v'\}$ and $E_v=E_{v'}=\Q_5$. The
character $\chi$ maps $G$ onto $\Z_5^\times$, because $E$ is ramified only at $2$ and every subfield
$\ne\Q$ of the field of $5$-power roots of unity is ramified at $5$. It maps each of $G_v$, $G_{v'}$ onto
$\Z_5^\times$, because the cyclotomic character of $\Q_5$ is onto. An element $w\in W$ is fixed if and only
if $(u^{-4}-1)w=0$ for all $u\in\Z_5^\times$. Now $u^{4}\equiv1\pmod5$ for all $u$, and $2^4-1=15$ is not
divisible by $25$. So, for $H=G$, $G_v$ or $G_{v'}$,
\begin{equation}\label{D:347}
 H^0(H,W)=W[5],
\end{equation}
the subgroup of order $5$, and $H$ acts trivially on it. We identify $W[5]$ with $\mathbb F_5$ by $\frac15\mapsto1$.

\emph{Step 2: the connecting homomorphism.} The sequence $0\to W[5]\to W\xrightarrow{\,5\,}W\to0$ gives, for
the same $H$, the exact sequence
\begin{equation}\label{D:348}
 H^0(H,W)\xrightarrow{\ 5\ }H^0(H,W)\xrightarrow{\ \delta\ }\mathrm{Hom}(H,\mathbb F_5)\longrightarrow H^1(H,W) .
\end{equation}
The first map is $0$ by Step~1. So the kernel of $\mathrm{Hom}(H,\mathbb F_5)\to H^1(H,W)$ is $\delta(W[5])$, a
group of order $5$ generated by $c=\delta(\frac15)$. To compute $c$, lift $\frac15$ to $\frac1{25}\in W$:
$c(g)$ is the class of $g\cdot\frac1{25}-\frac1{25}=(\chi(g)^{-4}-1)/25$, that is,
\begin{equation}\label{D:349}
 c(g)=\frac{\chi(g)^{-4}-1}{5}\bmod 5 .
\end{equation}
The same formula holds on $G$ and on $G_v$, $G_{v'}$, so the restriction of $c$ is the generator there. For
$u\in\Z_5^\times$ put $\lambda(u)=(u^{4}-1)/5\bmod5$. It depends only on $u$ modulo $25$. It is a
homomorphism $\Z_5^\times\to\mathbb F_5$, because $(1+5x)(1+5y)\equiv1+5(x+y)\pmod{25}$. And
$c(g)=-\lambda(\chi(g))$, because $u^{-4}-1=-(u^4-1)u^{-4}$ and $u^{-4}\equiv1\pmod5$.

\emph{Step 3: a class killed by $5$.} Let $z\in\mathrm{Sha}_5(E,4)$ with $5z=0$. By the sequence of Step~2
continued one term, $z$ is the image of some $\phi\in\mathrm{Hom}(G,\mathbb F_5)$. The restriction of $z$ to $G_v$
is $0$, so the restriction of $\phi$ to $G_v$ lies in the kernel described in Step~2. The same holds at
$v'$. So there are $a_1,a_2\in\mathbb F_5$ with
\begin{equation}\label{D:350}
 \phi|_{G_v}=a_1\,c|_{G_v},\qquad \phi|_{G_{v'}}=a_2\,c|_{G_{v'}} .
\end{equation}

\emph{Step 4: reciprocity gives $a_1=a_2$.} Let $\mathrm{rec}$ be the reciprocity map of local class field
theory for $\Q_5$, from $\Q_5^\times$ to the Galois group of the largest abelian extension. In its action on
the roots of unity of $5$-power order, $\mathrm{rec}(5)$ is trivial, and $\mathrm{rec}(u)$ acts as $u^{\epsilon}$
for $u\in\Z_5^\times$, where $\epsilon=\pm1$ depends only on the normalisation of $\mathrm{rec}$. So at each
of the two places, for $u\in\Z_5^\times$ and an integer $k$,
\begin{equation}\label{D:351}
 c\big(\mathrm{rec}(5^ku)\big)=-\lambda(u^{\epsilon})=-\epsilon\,\lambda(u).
\end{equation}
Global class field theory says: the homomorphism $\phi$ factors through the largest abelian quotient of $G$,
and for every $\beta\in E^\times$ the sum over all places $w$ of $E$ of the values of $\phi$ on the local
reciprocity image of $\beta$ at $w$ is $0$. Take $\beta=2+i$, of norm $5$. At a finite place $w\nmid5$,
$\beta$ is a unit and $\phi$ is unramified, so the term is $0$. At the infinite place the term is $0$,
because $\C^\times$ is connected. There remain $v$ and $v'$. They correspond to the two square roots of $-1$
in $\Z_5$. Let $v$ be the place at which $i\equiv7\pmod{25}$; indeed $7^2=49\equiv-1$. At $v$,
$\beta\equiv9\pmod{25}$ is a unit, and
\begin{equation}\label{D:352}
 9^4=6561\equiv11\pmod{25},\qquad \lambda(9)=2 .
\end{equation}
At $v'$, $i\equiv-7$ and $\beta\equiv-5\pmod{25}$ is not a unit. Write $\beta=5/\bar\beta$ with
$\bar\beta=2-i$. At $v'$, $\bar\beta\equiv9\pmod{25}$, so $\beta=5u$ with $\lambda(u)=-\lambda(9)=-2$. The
relation is
\begin{equation}\label{D:353}
 0=a_1\cdot(-\epsilon)\cdot2+a_2\cdot(-\epsilon)\cdot(-2)=-2\epsilon\,(a_1-a_2)\quad\text{in }\mathbb F_5 .
\end{equation}
Hence $a_1=a_2$. (The unit $i$ gives no condition, as it must: $\lambda(7)=0$, because $7^4\equiv1\pmod{25}$.)

\emph{Step 5: $z=0$.} The homomorphism $\phi-a_1c\colon G\to\mathbb F_5$ is trivial on $G_v$ and on $G_{v'}$. So it
is unramified at the primes above $5$. It is unramified at all other primes, by the definition of $G$. The
field $E$ has no real place. So the extension of $E$ cut out by $\phi-a_1c$ is abelian and unramified
everywhere. The class number of $\Q(i)$ is $1$, so this extension is $E$ and $\phi=a_1c$. By Step~2 the image
of $c$ in $H^1(G,W)$ is $0$. Hence $z=0$.

\emph{Step 6: the whole group.} A continuous cocycle $G\to W$ has finite image, because $G$ is compact and
$W$ is discrete. So its values lie in the subgroup killed by some power of $5$, and that power kills its
class. If $z\ne0$ were a class in $\mathrm{Sha}_5(E,4)$ of order $5^{k}$, then $5^{k-1}z$ would be a nonzero
class in $\mathrm{Sha}_5(E,4)$ killed by $5$. Step~5 excludes this. So $\mathrm{Sha}_5(\Q(i),4)=0$.\qed

The congruences of Step~4 were recomputed by a script; the element $1+2i$ in place of $2+i$ gives the same
relation. The result agrees with Corollary~(b): if the description of the two parts given above is granted,
the part with $\jmath=+1$ is $\mathrm{Sha}_5(\Q,4)$, which is $0$ because $(5-1)\mid 2\cdot4$.

\emph{Why the Bernoulli number is divisible by $5$ and the group is $0$} (from memory; not proved here and not
used). By the duality of Poitou and Tate, $\mathrm{Sha}_5(E,4)$ is dual to the kernel of
$H^2(G,\Z_5(5))\to\prod_{v\mid5}H^2(G_v,\Z_5(5))$, and by Tate's local duality and Step~1 each local group has
order $5$. A consequence of the main conjecture \cite{MW84} should give the order $5^{\,v_5(B_{5,\chi_4}/5)}=5$
for the part of $H^2(G,\Z_5(5))$ with $\jmath=-1$. So the Bernoulli number would measure a global group that
maps injectively to the local groups. We have not compared this specialisation of the main conjecture with
\cite{MW84}. The proof above does not use it. The class number
of $\Q(\zeta_{20})$, which is $1$ (from memory), does not give the result by the method of Proposition~G3: two
primes of $\Q(\zeta_{20})$ lie above $5$ and both ramify in the cyclotomic $\Z_5$-tower, so
\cite[Prop.~13.22]{Was82} does not apply.

Three tests of the normalisations of the source, for the characters of conductor $1$, $4$, $8$ and
$r\le42$: $B_{r,\chi}\ne0$ exactly when $\chi(-1)=(-1)^{r}$ (except $\chi=1$, $r=1$, where
$B_{1,1}=\frac12$ in the normalisation above); the number
$N\prod_{p\le r+1}p\cdot B_{r,\chi}$ is an integer for $N\in\{f_\chi,8\}$ (and $N\,B_{r,\chi}$ alone is not, in $42$
cases); and its absolute value is at most $16N^{r+1}r!$ (the largest ratio found is $1/32$). For
$1\le e\le40$ the sum $\sum_\ell v_\ell\log\ell$ over the odd primes in the numerators above is at most
$2\log\big(16\cdot4^{e+2}(e+1)!\big)$, the bound that \texttt{PrimeBound} gives for $\Q(i)$.

\textbf{Sizes.} For the field of Theorem~G the bound for the minus part in the table is
$3\cdot4n\cdot t^2\log(tN)$ with $n\le2x^{2+8\varepsilon}$ and $t\le x^{401}$: a power of $x$, about
$x^{804}$. This upper bound is below $\sqrt b=\exp(x^{\varepsilon}/2)$ only for very large $x$: for
$\varepsilon=1/53$, from about $\log x=742$ on (script). The crossing illustrates a sufficient scale, obtained
from an upper bound. It gives no lower bound on the scale at which the actual cohomology satisfies
Theorem~G.

\subsection*{6.\quad Assumed here, from the source}

\begin{enumerate}[label=\textup{(\roman*)},leftmargin=2.4em]
\item \emph{Theorem G.} The statement \texttt{VWin} (\texttt{ExistsPageN.lean}:49). It is proved by
\texttt{vwin\_of\_row} from \texttt{Row0b}, and \texttt{Row0b} by
\texttt{galoisLocalGlobalStrong\_}\allowbreak\texttt{recParameter}. Both are among the declarations the
final theorem depends on, and \texttt{vwin\_of\_row} depends only on the three standard axioms.
\item \emph{The inequality \eqref{gc:link}.} The statement \texttt{LinkN}. Its proof in the source is in
the theorems \texttt{final\_alt\_alt} and \texttt{linkN\_of\_twGood\_}\allowbreak\texttt{galoisPreconditions}. Only
the statements were read. The proof names bounds at the primes $2$ and $3$
(\texttt{boundAtTwoThree\_}\allowbreak\texttt{linked\_of\_recT}) and a family of ``Galois preconditions''.
\item \emph{The fact that $d_j=D_0$} (\texttt{SmallLevelRankEq}). The meaning of $\eta_j(w)$ is \emph{not}
in this list of things taken from the source, because it is not established: the lattice definition
(\texttt{etaW}, \texttt{locLen}, \texttt{eta\_units}) is for the number of a family, and the declarations
\texttt{TieEta} and \texttt{etaN\_eq\_etaJ}, cited earlier for the tie with the field \texttt{etaJ} of the
record, are not among those the final theorem depends on (\S2). The identification of $\eta_j(w)$ with the
defect of the index of lattices of Part~II, \S6.8 is open. Nothing proved here uses it.
\item \emph{The identification of the data.} That the numbers \texttt{shaLen} in the record of the
construction are
the lengths $h_{\ell,e}$ of the field $E$ (\texttt{nd.sha\_eq}); and the passage from Theorem~G to
\texttt{HypV} described in \S1: the change of scale (\texttt{winPairs}), the bound for the number of twists
(\texttt{RankVal} and the theorem
\texttt{eventually\_rank\_le\_}\allowbreak\texttt{rpow\_of\_lt\_one\_div\_fiftytwo}) and
\texttt{supplyOfAdmOf}. The statements and the proof of \texttt{winPairs} were read; the other proofs were not.
\item \emph{Every row of the table of \S4,} the hypotheses of the row for Case~I, and the inputs listed
after it.
\end{enumerate}

\textbf{What this section proves, and what it assumes.} \emph{Proved here outright:} Lemma~G1 and
Proposition~G2. \emph{Proved here from theorems of \cite{Was82}, checked in the first edition}
(Propositions 6.16 and 13.22, Lemma 13.16, Theorem 6.17, Corollary 5.14): Proposition~G3 and the corollaries
(a), (b), (c); corollary (b) also uses Lemma~G4. Proposition 13.22 does not state the compatibility with
$\Pi$; the proof of Proposition~G3 supplies it. The divisibility of
$B_{e+1,\chi_4}/(e+1)$ by $5$ for $e\equiv0\pmod4$ follows from Theorem 5.11 and Corollary 5.13 of the same
book. \emph{Proved here from Minkowski's lattice point theorem} (\cite{Neu99}, Ch.~I, (4.4) with its strict
inequality, and (5.2); checked in a copy): Lemma~G4. \emph{Proved here from the standard statements of local and global class field theory}
(the reciprocity law for principal ideles, the action of the local reciprocity map of $\Q_5$ on roots of
unity, and the triviality of the Hilbert class field of $\Q(i)$), for this single field and twist:
$\mathrm{Sha}_5(\Q(i),4)=0$.
\emph{Not proved here:}
the description of the two parts of $\mathrm{Sha}_\ell(\Q(i),e)$ by characters; the remark on the order of
$H^2(G_{E,5},\Z_5(5))$ after the proof for $\Q(i)$; the two identifications with the Lean stated in \S5; and
the meaning of $\eta_j(w)$ as a length of a quotient of lattices (\S2). \emph{Computed exactly:}
the table of \S5 and the three tests. \emph{Taken from the source:}
Theorem~G, \eqref{gc:link}, items (i)--(v). \emph{Cited as published theorems:} the
correspondences in the last column of the table of \S4, and the theorems of Soul\'e, Voevodsky and Rost,
Mazur and Wiles, and Kurihara quoted in \S3. They are cited from memory, except where a theorem or a section
is named; a named place in an article was compared by a reviewer, not by us (see ``Citations'' at the head
of the section).

\textbf{What is in the literature and what is not.} We sort the statements of the source named in this
section into three lists. The source proves all of them again.

\emph{(A) Published, and cited with the place.} \texttt{FerreroWashington}: \cite{FW79} (from memory), in the
form with bounded rank by \cite[Prop.~13.23]{Was82} (read). The duality theorems in \texttt{LitCoh}: Poitou and Tate
\cite[\S8.6]{NSW00}; for Tate's local theorems, the formula of Greenberg and Wiles and the local theorems of
Bloch and Kato the attribution is from memory. \texttt{BhattScholzeIntegral}: \cite[Thms.~1.2 and~5.6]{BS23}
(compared by a reviewer). We did not compare the Lean statements with the published ones
word by word. \texttt{PropZ} does not belong to this list: it is not cited from a publication. Its classical
form is Corollary~(b), proved in \S5 from class numbers; that Corollary~(b) is \texttt{PropZ} depends on the
identification of \texttt{Chr.V} stated in \S5, which was not checked in full.

\emph{(B) A published counterpart exists, and it covers less or has not been matched.}
\texttt{PsiBound}: an upper bound for a cohomology length in terms of the valuation of a Bernoulli number
with explicit integral factors. It corresponds to a consequence of \cite{MW84} (from memory); general bounds
from Euler systems are in \cite[\S2.2]{Rub00}, and the specialisation is not supplied.
\texttt{Row0b}: a continuous semisimple representation, with predicted weights and the local comparisons in
\texttt{PropsB}; compare \cite{CH13}, \cite{Car12}, \cite[Thm.~1.1]{Car14}, with normalisations and
hypotheses still to be matched. \texttt{KatoTwoW}: the two rational Selmer spaces of Bloch and Kato vanish
under the attachment, \texttt{Used} and \texttt{WitAQ} hypotheses; Kato \cite[\S14.1, Thm.~14.2(1)]{Kat04}
proves finiteness of the discrete Selmer groups at the critical twists other than the central one.
\texttt{LemmaK}: Soul\'e gives finiteness (\cite{Sou79}, from memory) and a quantitative bound away from
small primes (\cite[Thms.~2 and~3]{Sou03}, compared by a reviewer). We have not identified a published estimate for
\texttt{LemmaK} in its full small-prime range. \texttt{OneCharU\_s}: the method of congruences is illustrated
by \cite[Thm.~1.1]{BC04}, which constructs a Selmer extension; the upper bound is in list (C).

\emph{(C) Appears to be new, or not identified in the literature.} Theorem~G itself (\texttt{VWin}); the
bounds for the plus
part, \texttt{CountI}, \texttt{OneCharU\_s}, \texttt{PropGK}, \texttt{CountII}, \texttt{CountIII},
\texttt{CountV}, \texttt{SumIn}, \texttt{SumOut}; the uniform statement \texttt{ABound} with the window,
\texttt{ThmBab}; \texttt{LinkN}; the statement \texttt{SymplecticFourImageDense} on $\mathrm{GSp}_4$; and
\texttt{LemmaK} in its range of small primes. Each of them is a theorem
of the archive that the final theorem depends on. Three of these names are not explained in the body.
\texttt{SumOut} says: for the odd primes $\ell\notin\mathfrak B$ together, the sum of the plus parts over
$E$, weighted by $\log\ell$, for the twists $e<t$, is at most $(Nt)^{C}$, with the level $N$ of the source. \texttt{PropGK} says: there is a
nonzero integer, of logarithm bounded by a power of $N(e+1)$ times $\log p$ for an auxiliary prime $p$ (and a factor that depends on the orbit), such
that for every $\ell\ne p$ twice its $\ell$-adic valuation bounds the sum of the normalised lengths over an
orbit of characters. \texttt{CountV} is described after the table of \S4. For \texttt{ThmBab} and
\texttt{SumIn} the table of \S4 names related work (Ferrero and Washington, Iwasawa, Soul\'e); that work is
not a counterpart of the statements. An upper bound for the plus part of the groups $\mathrm{Sha}_\ell(E,e)$, $e\ge1$ (groups of twisted
Galois cohomology; they are not ideal class groups), by a
congruence number of automorphic forms, uniform in the field and the twist, is not a published theorem as
far as we know. Every statement of novelty in this list is provisional: it is a judgment made without a
search of the literature, and it was not tested.

\textbf{What remains.}
\begin{enumerate}[label=\textup{(\roman*)},leftmargin=2.4em]
\item To read the proof of one row of the table completely. The minus part (\texttt{PsiBound}) is the natural
first choice: it concerns only abelian fields and Bernoulli numbers.
\item To state \texttt{OneCharU\_s} with its hypotheses on $\rho$ (\texttt{Adm3U\_s}, \texttt{Adm2A}) in
classical terms, to compare them with \texttt{Adm3} and \texttt{Adm2}, and to find how the nonvanishing
integer of \texttt{PropGK} is produced.
\item To identify the published counterpart of the rows for $\mathrm{GSp}_4$.
\item To establish the meaning of $\eta_j(w)$: to find in the proof of \texttt{LinkN}
(\texttt{linkN\_of\_twGood\_}\allowbreak\texttt{galoisPreconditions}) the statement that ties the field
\texttt{etaJ} of the record to the lattice length \texttt{etaW} at the twisted primes, and to check that
the final theorem depends on it. The declarations \texttt{TieEta} and \texttt{etaN\_eq\_etaJ} do not serve:
the final theorem does not depend on them. Until then the identification of $\eta_j(w)$ with the defect of
the index of lattices of Part~II, \S6.8 is not established.
\item To check the constants in the passage from the three parts to $\sqrt b$ (\texttt{vwin\_of\_joint}),
and to evaluate the constant $c_\psi$ of \S1 for the trivial character.
\item To prove the description of the parts of $\mathrm{Sha}_\ell(\Q(i),e)$ by characters, and to compare
the remark on $H^2(G_{E,5},\Z_5(5))$ in \S5 with \cite{MW84}.
\item To state \texttt{PsiBound} over $\Q$ by descent, and to match the normalisations of \texttt{Row0b} with
\cite{CH13,Car12,Car14}.
\end{enumerate}


\section{The objects of the construction}\label{ob:sec}

Parts~I and~II describe the objects behind the main inequality (M) informally, and \S III.4--III.9 each
work in an abstract setting: a chain of maps with Gaussian estimates, power series in formal-group
coordinates, lattices with a pairing, operators on a Hilbert space, cohomology groups. This section says
what the objects are. It defines the field $E$, the Hermitian spaces and lattices, the parameters with
their exact sizes, the spaces of forms, the theta kernels and the family of maps, all in classical
language. Then it gives a dictionary: for each abstract object of \S III.4--III.9, the object of the
construction it stands for, and how far this was checked in the source. \S\ref{as:sec} reduces (M) to twelve Hypotheses on the objects defined here, and the Riemann hypothesis
to these and a thirteenth; it uses the notation fixed here. No proof of (M) is given in this section or in that one.

\emph{Proved here.} The structure of the field $E$ (Lemma~O2). The splitting of the primes $p\nmid2qd$ in
$E$, from two cited theorems (Lemma~O3). The rule that sends a prime power to a cell of the table of
\S I.5, for the primes $p\nmid2qd$ (Lemma~O5) and for the ramified prime $d$ (Lemma~O22). The root
discriminant of $E$, from the conductor--discriminant formula, which is cited (Lemma~O6). Elementary
facts on Hermitian lattices, duals and slopes (Lemmas~O9 and~O11). The arithmetic of the ranks and
signatures along the chain (Lemma~O13). For a theta series of a pair of lattices: its convergence, its
Gaussian majorant, its Taylor coefficients, its equivariance, and the fact that it defines a linear map
between spaces of algebraic modular forms (Lemmas~O17 and~O18).

\emph{Cited.} Four theorems of algebraic number theory, (C1)--(C4) in \S1; the finiteness of the number of
classes in a genus, (C5), and the structure of lattices over $\mathfrak o_E$, in \S2; the classification
of the representations of the unitary groups, their dimensions and the branching rule, in \S4.

\emph{Assumed here, from the Lean development.} Seven numbered Hypotheses; together they contain every
statement of this section that is taken from the source. O23: the parameters exist, with the sizes of the
table of \S3. O12: the Hermitian spaces, lattices and levels exist. O24: the weights of the forms.
O19: the maps of the chain are given by kernels that are Taylor coefficients of entire functions. O20:
these Taylor coefficients are algebraic and, in the coordinates of \S III.5, integral. O21: the kernels
satisfy the hypotheses of the estimates of \S III.4. O25: the family of words, the pairing at the end of
the chain and its lattices. Nothing in this section proves that these data exist. The source asserts it,
as theorems, for $0<\varepsilon\le\varepsilon_0$ (\S3, ``Status in the source'').

\emph{The dictionary} is in \S7. Its third column says, row by row, whether an identification was verified
in the source, read from the types of the source only, or not verified. The letters fixed in this section
are listed in a table after \S8. Names of declarations of the Lean development occur only in \S9.

\subsection*{1.\quad The field $E$}

\textbf{Definition O1 (number fields; recalled).} A \emph{number field} $K$ is a field containing $\Q$ of
finite degree $n_0=[K:\Q]$. Its \emph{ring of integers} $\mathfrak o_K$ is the set of elements of $K$ that
are roots of monic polynomials with integer coefficients. It is a ring, and as an abelian group it is
free of rank $n_0$. Every nonzero ideal of $\mathfrak o_K$ is a product of prime ideals, uniquely up to
order. The \emph{norm} of a nonzero ideal $\mathfrak a$ is the number of elements of
$\mathfrak o_K/\mathfrak a$. For a prime number $p$ write $p\,\mathfrak o_K=\mathfrak p_1^{e_1}\cdots\mathfrak p_g^{e_g}$
with distinct prime ideals $\mathfrak p_i$. The field $\mathfrak o_K/\mathfrak p_i$ has $p^{f_i}$ elements;
$e_i$ is the \emph{ramification index} and $f_i$ the \emph{residue degree}, and $\sum_ie_if_i=n_0$. The
prime $p$ is \emph{unramified} in $K$ if all $e_i=1$, and \emph{ramified} otherwise. An \emph{embedding}
of $K$ is a field homomorphism
$K\to\C$; there are $n_0$ of them. The \emph{discriminant} $\mathrm{disc}\,K$ is $\det(\tau_i(\omega_j))^2$ for a
$\Z$-basis $\omega_1,\dots,\omega_{n_0}$ of $\mathfrak o_K$ and the embeddings $\tau_1,\dots,\tau_{n_0}$, and
the \emph{root discriminant} is $\mathrm{rd}_K=|\mathrm{disc}\,K|^{1/n_0}$. (These facts are in
\cite[Ch.~I, \S2--\S3 and \S8]{Neu99}.) The letters $e_i$, $f_i$, $g$ of this definition always carry an
index or stand in a triple $(e,f,g)$; without index, $e$ and $f$ are the numbers of ``The data'' below.

\emph{Example.} $K=\Q(i)$, $\mathfrak o_K=\Z[i]$, $\mathrm{disc}=-4$, $\mathrm{rd}_K=2$. Here
$5\,\Z[i]=(2+i)(2-i)$, two primes with $e_i=f_i=1$; $3\,\Z[i]$ is prime, with $f_1=2$; and
$2\,\Z[i]=(1+i)^2$, with $e_1=2$: the prime $2$ is ramified.

\textbf{Cited theorems.} For (C1) and (C2) the statements were read in a local copy of \cite{Neu99}; the
other references are from memory.
\begin{itemize}[leftmargin=2.6em,itemsep=1pt,topsep=2pt]
\item[(C1)] \emph{Cyclotomic fields} (\cite[Ch.~I, Propositions~(10.2) and~(10.3)]{Neu99}; for the Galois
group, \cite[Theorem~2.5]{Was82}). Let $Q\ge3$ and $\zeta_Q=e^{2\pi i/Q}$. The field $\Q(\zeta_Q)$ has degree
$\varphi(Q)$ and ring of integers $\Z[\zeta_Q]$. For each $a$ prime to $Q$ there is one automorphism
$\sigma_a$ with $\sigma_a(\zeta_Q)=\zeta_Q^a$, and $a\mapsto\sigma_a$ is an isomorphism of
$(\Z/Q)^\times$ onto the Galois group. A prime $p\nmid Q$ is unramified, and every prime above it has
residue degree equal to the order of $p$ in $(\Z/Q)^\times$.
\item[(C2)] \emph{Primes in Galois extensions} (\cite[Ch.~I, Proposition~(9.1) and \S9]{Neu99}). Let
$K\subset K'$ be number fields with $K'/K$ Galois, and $\mathfrak p$ a prime ideal of $\mathfrak o_K$. The
Galois group permutes the prime ideals of $\mathfrak o_{K'}$ that contain $\mathfrak p$ transitively.
Hence they all have the same $e$ and $f$ over $\mathfrak p$, and $e\le[K':K]$. Ramification indices
multiply in a tower of fields.
\item[(C3)] \emph{The conductor--discriminant formula} (\cite[Theorem~3.11]{Was82}; in general form
\cite[Ch.~VII, (11.9)]{Neu99}). Let $K\subset\Q(\zeta_Q)$ and let $\mathcal X$ be the group of Dirichlet
characters modulo $Q$ that are trivial on the subgroup of $(\Z/Q)^\times$ fixing $K$. Then
$|\mathrm{disc}\,K|=\prod_{\chi\in\mathcal X}\mathrm{cond}(\chi)$, where $\mathrm{cond}(\chi)$ is the conductor.
\item[(C4)] \emph{Gauss sums} (\cite[Ch.~2]{Dav00}). For a prime $q\equiv3\bmod4$ the sum
$\mathfrak g_q=\sum_{a\bmod q}(a/q)\zeta_q^a$ satisfies $\mathfrak g_q^2=-q$.
\end{itemize}

\textbf{The data.} Let $n=2^e$ with $e\ge2$, $N=4n$, and let $q\equiv3\bmod8$, $d\equiv7\bmod16$ be
distinct primes. Put $\chi_q(a)=(a/q)$, $\chi_d(a)=(a/d)$ (Legendre symbols), and
\begin{equation}\label{D:354}
 E=\Q\big(\zeta_N+\zeta_N^{-1},\ \sqrt{-q},\ \sqrt{-d}\big),\qquad k_d=\Q(\sqrt{-d}),\qquad f=2n .
\end{equation}
Let $H$ be the group of odd residue classes modulo $N$ with $a$ and $-a$ identified, written $[a]$, as in
Part~II, \S4.1. The group $H$ is written multiplicatively, as in Parts~I and~II; \S III.3, \S III.6 and
the source write it additively (the column $t[p^i]$ here is $t+ih_p$ in \eqref{lt:cell}).

\textbf{Lemma O2 (the field $E$).} \emph{Put $Q=Nqd$.}
\begin{enumerate}[label=\textup{(\alph*)},leftmargin=2.2em,itemsep=1pt,topsep=2pt]
\item \emph{$E\subset\Q(\zeta_Q)$, and $E$ is Galois over $\Q$ with abelian group $G$. The map}
\begin{equation}\label{D:355}
 (\Z/Q)^\times\to H\times\{\pm1\}\times\{\pm1\},\qquad a\mapsto\big([a],\chi_q(a),\chi_d(a)\big),
\end{equation}
\emph{induces an isomorphism of $G$ onto $H\times\{\pm1\}^2$. For $\sigma\in G$ we write
$\sigma=([\sigma],\varepsilon_q(\sigma),\varepsilon_d(\sigma))$. Then
$\sigma(\zeta_N+\zeta_N^{-1})=\zeta_N^a+\zeta_N^{-a}$ for $[a]=[\sigma]$,
$\sigma(\sqrt{-q})=\varepsilon_q(\sigma)\sqrt{-q}$ and $\sigma(\sqrt{-d})=\varepsilon_d(\sigma)\sqrt{-d}$.}
\item \emph{$[E:\Q]=4n$. The group $H$ is cyclic of order $n$, generated by $[5]$.}
\item \emph{Complex conjugation restricts to the element $c=([1],-1,-1)$ of $G$. Its fixed field is
$F=\Q(\zeta_N+\zeta_N^{-1},\sqrt{qd})$, of degree $2n$. Every embedding of $F$ is real and no embedding
of $E$ is real, and $E=F(\sqrt{-d})$. For every embedding $\tau$ of $E$ and every $y\in E$,
$\tau(c(y))=\overline{\tau(y)}$.}
\end{enumerate}

\noindent A field with the properties in (c) is called a \emph{CM field}, and $F$ is its maximal real
subfield.

\emph{Proof.} (a) By (C4), $\sqrt{-q}=\pm\mathfrak g_q\in\Q(\zeta_q)$ and $\sqrt{-d}\in\Q(\zeta_d)$; and
$\zeta_N,\zeta_q,\zeta_d$ are powers of $\zeta_Q$. So $E\subset\Q(\zeta_Q)$. By (C1) the automorphisms of
$\Q(\zeta_Q)$ are the $\sigma_a$. From the definition of $\mathfrak g_q$,
$\sigma_a(\mathfrak g_q)=\sum_b(b/q)\zeta_q^{ab}=(a/q)\mathfrak g_q$,
and likewise for $d$; and $\sigma_a(\zeta_N+\zeta_N^{-1})=\zeta_N^a+\zeta_N^{-a}$. Now
$\zeta_N^a+\zeta_N^{-a}=2\cos(2\pi a/N)$, and this equals $2\cos(2\pi/N)$ exactly when $a\equiv\pm1\bmod N$.
Hence $\sigma_a$ fixes the three generators of $E$ if and only if $a$ lies in the kernel $\Gamma$ of the
displayed map. By Galois theory $E$ is the fixed field of $\Gamma$, it is Galois over $\Q$, and
$G=(\Z/Q)^\times/\Gamma$ is isomorphic to the image. The map is onto, by the Chinese remainder theorem:
the classes of $a$ modulo $N$, $q$ and $d$ can be chosen independently. (b) $|G|=|H|\cdot4$ and
$|H|=\varphi(N)/2=n$. The class of $5$ has order $2^e=n$ in $(\Z/2^{e+2})^\times$ and $-1$ is not a power
of $5$ modulo $2^{e+2}$ (powers of $5$ are $\equiv1\bmod4$). So $[5]$ has order $n$ in $H$. (c) Complex
conjugation is $\sigma_{-1}$. Since $q\equiv d\equiv3\bmod4$, $\chi_q(-1)=\chi_d(-1)=-1$, and
$[-1]=[1]$. The element $\sqrt{qd}=-\sqrt{-q}\sqrt{-d}$ is fixed by $c$, so $F$ lies in the fixed field
of $c$. Conversely, an element $\sigma$ of $G$ that fixes $\zeta_N+\zeta_N^{-1}$ and $\sqrt{qd}$ has
$[\sigma]=[1]$ and $\varepsilon_q(\sigma)\varepsilon_d(\sigma)=1$, so $\sigma\in\{1,c\}$. By Galois theory
$F$ is the fixed field of $\{1,c\}$, and $[F:\Q]=4n/2=2n$. An embedding $\tau$ of $E$ is the restriction
of some $\sigma_a$, and $G$ is abelian, so $\tau c=c\tau$, which is the last assertion. It shows that
$\tau(F)\subset\mathbb R$ and, since $\tau(\sqrt{-d})=\pm i\sqrt d$, that $\tau(E)\not\subset\mathbb R$. \qed

\textbf{Lemma O3 (splitting of primes in $E$).} \emph{Let $p$ be a prime with $p\nmid2qd$, and let
$\sigma_p\in G$ be the image of $p$, so that $\sigma_p=([p],\chi_q(p),\chi_d(p))$. Then $p$ is
unramified in $E$. All primes of $\mathfrak o_E$ above $p$ have residue degree $f_p$, and there are $g_p$
of them, where}
\begin{equation}\label{ob:split}
 f_p=\text{order of }\sigma_p\text{ in }G=\max\big(\mathrm{ord}_H[p],\,o_p\big),\qquad g_p=\frac{4n}{f_p},\qquad
 o_p=\begin{cases}1,&\chi_q(p)=\chi_d(p)=1,\\2,&\text{otherwise.}\end{cases}
\end{equation}
\emph{Here $\mathrm{ord}_H[p]$ is the least $j\ge1$ with $p^j\equiv\pm1\bmod N$; it is a power of $2$.
The set of $\sigma\in G$ that map a given prime of $\mathfrak o_E$ above $p$ to itself is the subgroup
$\langle\sigma_p\rangle$ generated by $\sigma_p$.}

\noindent The automorphism $\sigma_p$ is the \emph{Frobenius} of $p$. So the splitting of $p$ in $E$ is
decided by $p$ modulo $N$ up to sign and by the two Legendre symbols. This is what makes the table of
\S I.5 a table: the cell of a prime power depends only on these three data (Lemma~O5).

\emph{Proof.} Put $G_Q=(\Z/Q)^\times$, the Galois group of $\Q(\zeta_Q)$, and let $\Gamma\subset G_Q$ be
the kernel of Lemma~O2(a). By (C1), $p$ is unramified in $\Q(\zeta_Q)$, and the primes above it have
residue degree the order $f'$ of $p$ in $G_Q$; let $g'$ be their number, so $f'g'=\varphi(Q)$. Let
$\mathfrak P$ be one of them, and let $G_{\mathfrak P}\subset G_Q$ be the group of the elements that map
$\mathfrak P$ to itself. By (C2), $G_Q$ permutes the $g'$ primes transitively, so
$|G_{\mathfrak P}|=\varphi(Q)/g'=f'$. The automorphism $\sigma_p$ satisfies
$\sigma_p(y)\equiv y^p\bmod p\,\Z[\zeta_Q]$ for every $y\in\Z[\zeta_Q]$: for
$y=\sum_ic_i\zeta_Q^i$ with $c_i\in\Z$, $y^p\equiv\sum_ic_i^p\zeta_Q^{ip}\equiv\sum_ic_i\zeta_Q^{ip}$.
So $\sigma_p$ maps $\mathfrak P$ into itself and lies in $G_{\mathfrak P}$. It has order $f'$. Hence
$G_{\mathfrak P}=\langle p\rangle$. Since ramification indices multiply in towers, $p$ is unramified in
$E$. Next, two primes of $\Q(\zeta_Q)$ above $p$ lie above the same prime of $E$ exactly when one is the
image of the other under an element of $\Gamma$, the Galois group of $\Q(\zeta_Q)$ over $E$: one direction
is clear, the other is (C2). The primes of $\Q(\zeta_Q)$ above $p$ are the $\gamma\mathfrak P$,
$\gamma\in G_Q$, and $\gamma\mathfrak P=\gamma'\mathfrak P$ exactly when $\gamma^{-1}\gamma'\in G_{\mathfrak P}$.
So $\gamma\mathfrak P$ and $\gamma'\mathfrak P$ lie above the same prime of $E$ exactly when
$\gamma'\in\gamma\,\Gamma G_{\mathfrak P}$; here $\Gamma G_{\mathfrak P}$ is a subgroup because $G_Q$ is
abelian. Hence the primes of $E$ above $p$ correspond to the cosets of $\Gamma G_{\mathfrak P}$ in $G_Q$,
their number is $[G_Q:\Gamma G_{\mathfrak P}]$, and the elements of $G=G_Q/\Gamma$ that fix the prime
$\mathfrak P\cap\mathfrak o_E$ are those of $\Gamma G_{\mathfrak P}/\Gamma=\langle\sigma_p\rangle$. Now
$[G_Q:\Gamma G_{\mathfrak P}]=|G|/|\langle\sigma_p\rangle|$. This gives $g_p=4n/\mathrm{ord}(\sigma_p)$,
and $f_p=4n/g_p$ by (C2) and $\sum_ie_if_i=4n$. Finally the order of $([p],\chi_q(p),\chi_d(p))$ is the
least common multiple of $\mathrm{ord}_H[p]$ and $o_p$. Both are powers of $2$, since $H$ has order
$2^e$, so it is their maximum. \qed

\emph{Remarks.} (1) \emph{Three classes of primes.} A prime $p\nmid2qd$ is called
\emph{split} if $\chi_d(p)=1$, \emph{signal} if $\chi_d(p)=-1$ and $\chi_q(p)=1$, and of \emph{inert
class} if $\chi_d(p)=\chi_q(p)=-1$ (these are the names of the source). Since $d\equiv3\bmod4$,
$\chi_d(p)=(-d/p)$: the split primes are those that split in $k_d$, and the other two classes are the
primes inert in $k_d$. These are the ``inert primes'' of Part~I.

(2) \emph{Primes of $F$.} \emph{A prime of $F$ above $p\nmid2qd$ stays prime in $E$ exactly when
$\sigma_p=c$, that is, when $\chi_q(p)=\chi_d(p)=-1$ and $p\equiv\pm1\bmod N$. For every other $p$, each
prime of $F$ above $p$ splits into two primes of $E$, interchanged by $c$.} Proof: let $\mathfrak P_E$ be
a prime of $E$ above $p$. By (C2) for $E/F$, whose Galois group is $\{1,c\}$, the primes of $E$ above
$\mathfrak P_E\cap F$ are $\mathfrak P_E$ and $c\,\mathfrak P_E$. So $\mathfrak P_E\cap F$ stays prime
exactly when $c$ maps $\mathfrak P_E$ to itself, that is, by Lemma~O3, when $c\in\langle\sigma_p\rangle$.
If $c=\sigma_p^j$, then $\varepsilon_d$ gives $\chi_d(p)^j=-1$, so $j$ is odd and $\chi_d(p)=-1$; likewise
$\chi_q(p)=-1$; and $[p]^j=[1]$ with $j$ odd forces $[p]=[1]$, because the order of $[p]$ is a power of
$2$. So $\sigma_p=c$. The converse is clear. \qed

(3) \emph{The ramified primes} are $2$, $q$ and $d$. The primes $2$ and $q$ do not occur in (M)
(Part~II, \S4.1). The prime $d$ does, and is treated in Lemma~O22. For $N=16$, $q=11$, $d=1319$ (the example
of \S8) the triples $(e,f,g)$ are $(4,2,2)$ at $2$, $(2,4,2)$ at $11$ and $(2,2,4)$ at $1319$ (computed with
PARI).

\textbf{Definition O4 (the CM type, places, cells).} The \emph{CM type} is the set
\begin{equation}\label{D:357}
 \Phi=\{\tau:E\to\C\ :\ \tau(\sqrt{-d})=+i\sqrt d\}.
\end{equation}
It has $f=2n$ elements, and every embedding of $E$ is in $\Phi$ or is the complex conjugate of exactly
one element of $\Phi$. The elements of $\Phi$ are the \emph{places} (the ``weight places'' of
\S III.4--III.8). Fix the \emph{base embedding} $\tau_{\rm b}\in\Phi$ with
$\tau_{\rm b}(\sqrt{-q})=+i\sqrt q$ and
$\tau_{\rm b}(\zeta_N+\zeta_N^{-1})=2\cos(2\pi/N)$. Every embedding is $\tau_{\rm b}\circ\sigma$ for one
$\sigma\in G$, and $\tau_{\rm b}\circ\sigma\in\Phi$ exactly when $\varepsilon_d(\sigma)=1$. A \emph{cell} is a
pair $(v,s)$ with $v\in\{0,1\}$ (the row) and $s\in H$ (the column). The cell of the place
$\tau_{\rm b}\circ\sigma$ is
\begin{equation}\label{ob:cell}
 (v,s),\qquad (-1)^v=\varepsilon_q(\sigma),\quad s=[\sigma].
\end{equation}
This is a bijection between $\Phi$ and the $2n$ cells of the table of \S I.5.

\textbf{Lemma O5 (the cell of a prime power).} \emph{Let $p\nmid2qd$ be a prime with $\chi_d(p)=-1$ and
let $i\ge1$ be odd. Let $\tau\in\Phi$ have the cell $(v,s)$. Then $\tau\circ\sigma_p^i\notin\Phi$,
and its complex conjugate $\tau\circ c\circ\sigma_p^i$ lies in $\Phi$ and has the cell}
\begin{equation}\label{ob:partner}
 \Big(v+\frac{1+\chi_q(p^i)}2\ \bmod2,\ \ s\cdot[p^i]\Big).
\end{equation}
\emph{In particular, from a place in row $0$ and column $t$ a signal prime ($\chi_q(p)=1$) leads to row
$1$ and column $t[p^i]$, and a prime of inert class leads to row $0$ and column $t[p^i]$. If
$\chi_d(p)=1$, then $\tau\circ\sigma_p^i\in\Phi$ for every $i$.}

\emph{Proof.} $\varepsilon_d(\sigma_p^i)=\chi_d(p)^i=-1$, so $\tau\circ\sigma_p^i$ sends $\sqrt{-d}$ to
$-i\sqrt d$. Composing with $c$ restores the sign. By Lemma~O2,
$c\,\sigma_p^i=([p^i],-\chi_q(p^i),+1)$. If $\tau=\tau_{\rm b}\circ\sigma$, the new place is
$\tau_{\rm b}\circ\sigma c\sigma_p^i$, with $\varepsilon_q=-\chi_q(p^i)(-1)^v$ and class $s[p^i]$. And
$-\chi_q(p^i)=(-1)^{(1+\chi_q(p^i))/2}$. \qed

\noindent For $v=0$, $s=t$ this is the rule of Part~I, \S I.5: the prime power $p^i$ goes to row
$\frac12(1+\chi_q(p^i))$ and column $t[p^i]$. (Parts~I and~II write $r=p^h$ for the prime power; here
$r$ is a rank and $h$ a Hermitian form, so we write $p^i$.) So the label $t\in T$ of (M) is a place in
row $0$, the place $(0,t)$, called a \emph{top place} below. The map
$\tau\mapsto\tau\circ c\circ\sigma_p$ is the ``Frobenius followed by complex conjugation'' of \S I.4.3 and
\S III.5.

\textbf{Lemma O22 (the cell of a power of $d$).} \emph{Let $E'=\Q(\zeta_N+\zeta_N^{-1},\sqrt{-q})$, so
that $E=E'(\sqrt{-d})$ and $E'$ is the fixed field of $I_d=\{1,([1],1,-1)\}\subset G$.}
\begin{enumerate}[label=\textup{(\alph*)},leftmargin=2.2em,itemsep=1pt,topsep=2pt]
\item \emph{The prime $d$ is unramified in $E'$, and every prime of $\mathfrak o_{E'}$ above $d$ is the
square of a prime of $\mathfrak o_E$: $d$ is ramified in $E$ with all $e_i=2$.}
\item \emph{The elements of $G$ whose restriction to $E'$ is the $i$th power of the Frobenius of $d$ in
$E'$ are the two elements $([d^i],\chi_q(d^i),\pm1)$. Exactly one of them, namely
$\sigma_{d,i}=([d^i],\chi_q(d^i),-1)$, carries $\Phi$ outside $\Phi$. For $\tau\in\Phi$ with the cell
$(v,s)$ and every $i\ge1$, odd or even, the place $\tau\circ c\circ\sigma_{d,i}$ lies in $\Phi$ and has
the cell \eqref{ob:partner} with $p=d$.}
\end{enumerate}

\emph{Proof.} (a) $E'\subset\Q(\zeta_{Nq})$ by Lemma~O2, and $d\nmid Nq$; so $d$ is unramified in $E'$ by
(C1) and (C2). In $\mathfrak o_E$ we have $d\,\mathfrak o_E=(\sqrt{-d}\,\mathfrak o_E)^2$, so every $e_i$ is
even, and $e_i\le[E:E']=2$ by (C2). (b) As in Lemma~O3, with $Nq$ in place of $Q$, the Frobenius of $d$
in $E'$ is the image $([d],\chi_q(d))$ of $d$ in $H\times\{\pm1\}$, the Galois group of $E'$. The
elements of $G$ above its $i$th power differ by $I_d$. The one with $\varepsilon_d=-1$ sends $\sqrt{-d}$
to $-\sqrt{-d}$. The rest is the proof of Lemma~O5. \qed

\noindent For $v=0$, $s=t$ this is the rule of Part~II, \S4.1 for the prime $d$: \emph{every} power
$d^i$ goes to row $\frac12(1+\chi_q(d^i))$ and column $t[d^i]$. In (M) the power $d^i$ has the
coefficient $\log d$, where a power $p^i$ of an inert prime has $2\log p$; this is the rule
$(1-\chi_d(p^i))\log p$ of Part~II, \S4.1, with $\chi_d(d)=0$. The coefficient is part of the definition
of the sums of (M); it is not derived here. The definition of these sums in the source has this form
(\S9, row~5): the Frobenius of $d$ is taken up to $I_d$, and the element outside $\Phi$ is used. Since
$q\equiv3\bmod4$ and $(-d/q)=1$ (a condition of \S3), $\chi_q(d)=-1$: odd powers of $d$ go to row $0$ and even powers to
row $1$.

\textbf{Lemma O6 (the discriminant).}
\begin{equation}\label{ob:rd}
 |\mathrm{disc}\,E|=2^{4((e+1)n-1)}\,(qd)^{2n},\qquad \mathrm{rd}_E=2^{\,e+1-1/n}\sqrt{qd}=\frac N2\,2^{-1/n}\sqrt{qd}.
\end{equation}
\emph{Proof.} By Lemma~O2 the characters in (C3) are the products $\psi\chi_q^{i}\chi_d^{i'}$ with
$i,i'\in\{0,1\}$ and $\psi$ a character modulo $N$ with $\psi(-1)=1$. The conductor of the product is
$\mathrm{cond}(\psi)\,q^{i}d^{i'}$. There are $2^{j-2}$ even characters modulo $2^j$ for $j\ge2$, so for
$3\le j\le e+2$ there are $2^{j-2}-2^{j-3}=2^{j-3}$ even characters of conductor exactly $2^j$, and one,
the trivial character, of conductor $1$. So $\prod_\psi\mathrm{cond}(\psi)$ is $2$ to the power
$\sum_{j=3}^{e+2}j\,2^{j-3}=(e+1)2^e-1$. Each $\psi$ occurs four times, and $q$ (resp.\ $d$) occurs in
$2n$ of the $4n$ characters. This gives the discriminant; take the $4n$th root. \qed

\noindent In Part~II, \eqref{II:t6.3} and in \S III.4, \S III.7 the number $\mathrm{rd}_E/2$ occurs under the name
$v$. Here $v$ is the row of a cell, and we always write $\mathrm{rd}_E/2$. By
\eqref{ob:rd} it is $\frac N4\,2^{-1/n}\sqrt{qd}$, a little less than $n\sqrt{qd}$.

\subsection*{2.\quad Hermitian spaces and lattices}

\textbf{Definition O7 (Hermitian spaces).} Write $\bar y=c(y)$ for $y\in E$. A \emph{Hermitian space} of
rank $r_1$ over $E$ is an $E$-vector space $\mathsf X$ of dimension $r_1$ with a map
$h:\mathsf X\times\mathsf X\to E$ that is $E$-linear in the second variable, satisfies
$h(y,y')=\overline{h(y',y)}$, and is \emph{nondegenerate}: if $h(y,y')=0$ for all $y'$, then $y=0$. In a
basis $e_1,\dots,e_{r_1}$ it is given by its \emph{Gram matrix} $J=(h(e_i,e_j))$, with
$\overline J^{\,t}=J$ and $\det J\ne0$. For $\tau\in\Phi$ the matrix $\tau(J)$ is a complex Hermitian
matrix (Lemma~O2(c)). It has $r^+_\tau$ positive and $r^-_\tau$ negative eigenvalues,
$r^+_\tau+r^-_\tau=r_1$; the pair $(r^+_\tau,r^-_\tau)$ is the \emph{signature at $\tau$}, and it does not
depend on the basis. The space is \emph{positive definite} if $r^-_\tau=0$ for all $\tau\in\Phi$,
\emph{definite} if at each $\tau$ one of $r^+_\tau,r^-_\tau$ is $0$, and \emph{indefinite at $\tau$} if
$r^+_\tau r^-_\tau\ne0$. We write $\mathsf X_\tau$ for the complex vector space $\C^{r_1}$ with the
Hermitian form $h_\tau$ of Gram matrix $\tau(J)$; a vector $y=\sum_iy_ie_i$ of $\mathsf X$ gives the
vector $y_\tau=(\tau(y_i))_i$ of $\mathsf X_\tau$, and an $E$-linear map $\lambda$ gives a complex linear
map $\lambda_\tau$, by applying $\tau$ to its matrix. The \emph{unitary group} is
\begin{equation}\label{D:361}
 U(\mathsf X)=\{\gamma:\mathsf X\to\mathsf X\ \text{$E$-linear and bijective}:\ h(\gamma y,\gamma y')=h(y,y')\ \text{for all }y,y'\}.
\end{equation}
It is the group of Part~II, \S6.1 (there $\overline g^{\,t}Jg=J$).

\emph{Examples.} (1) $\mathsf X=E^2$ with $J=\mathrm{diag}(1,\sqrt{qd})$. Since
$\tau_{\rm b}(\sqrt{qd})=-(i\sqrt q)(i\sqrt d)>0$, a place $\tau$ with the cell $(v,s)$ has
$\tau(\sqrt{qd})=(-1)^v\sqrt{qd}$. So the signature is $(2,0)$ at the places of row $0$ and $(1,1)$ at
those of row $1$. (2) A \emph{hyperbolic plane} is a Hermitian space with a basis $e,e'$ such that
$h(e,e)=h(e',e')=0$ and $h(e,e')=1$. Its Gram matrix has the eigenvalues $\pm1$: the signature is $(1,1)$
at every place.

\textbf{Definition O8 (lattices, duals, slopes).} An \emph{$\mathfrak o_E$-submodule} of $\mathsf X$ is a
subgroup $L$ with $aL\subset L$ for all $a\in\mathfrak o_E$. A \emph{lattice} in $\mathsf X$ is an
$\mathfrak o_E$-submodule generated by finitely many elements that span $\mathsf X$ over $E$. It is
\emph{free} if $L=\bigoplus_i\mathfrak o_Ee_i$ for a basis $(e_i)$ of $\mathsf X$; a lattice need not be
free (see the example below). For a nonzero ideal $\mathfrak a$, or a quotient $\mathfrak a\mathfrak b^{-1}$
of two ideals, $\mathfrak aL$ is the lattice of the finite sums $\sum a_iy_i$ with $a_i\in\mathfrak a$,
$y_i\in L$. The \emph{dual} of $L$ is
\begin{equation}\label{D:362}
 L^\#=\{y\in\mathsf X:\ h(y',y)\in\mathfrak o_E\ \text{for all }y'\in L\},
\end{equation}
again a lattice. $L$ is \emph{integral} if $L\subset L^\#$ and \emph{unimodular} if $L=L^\#$. Suppose
$\mathsf X$ is positive definite. Put $\mathsf X_{\mathbb R}=\bigoplus_{\tau\in\Phi}\mathsf X_\tau$, a real
vector space of dimension $4nr_1$ with the Euclidean norm $\|y\|^2=\sum_{\tau\in\Phi}h_\tau(y_\tau,y_\tau)$. The
map $y\mapsto(y_\tau)_\tau$ embeds $\mathsf X$ in it, and $L$ becomes a discrete subgroup of full rank. Its
\emph{covolume} $\mathrm{covol}(L)$ is the volume of a fundamental parallelepiped for this norm. Its
\emph{slope} is
\begin{equation}\label{ob:mu}
 \mu(L)=-\frac{2}{4n\,r_1}\,\log\mathrm{covol}(L)+\mathrm{cst}_E,\qquad
 \mathrm{cst}_E=-\log2+\frac{\log|\mathrm{disc}\,E|}{4n}=\log\frac{\mathrm{rd}_E}2 .
\end{equation}
This is \eqref{lt:mu} of \S III.6 for the field $E$, with the volume fixed as the Euclidean volume of
this norm. It is the slope \emph{for squared norms}: twice the $\mu$ of Part~II, \S6.8, up to the
constant. The same definition is used over any CM field in place of $E$.

\textbf{Lemma O9 (slopes).} \emph{Let $\mathsf X$ be positive definite of rank $r_1$.}
\begin{enumerate}[label=\textup{(\alph*)},leftmargin=2.2em,itemsep=1pt,topsep=2pt]
\item \emph{If $L=\bigoplus_i\mathfrak o_Ee_i$ for a basis $(e_i)$ with Gram matrix $J$, then
$\mu(L)=-\frac1{2nr_1}\sum_{\tau\in\Phi}\log\det\tau(J)$. In particular $\mu(\mathfrak o_E^{r_1})=0$ for
the form $h(y,y')=\sum_i\bar y_iy_i'$.}
\item \emph{If $L'\subset L$ are lattices, then $\mu(L')=\mu(L)-\frac{1}{2nr_1}\log[L:L']$.}
\item \emph{For every lattice $L$, $\mu(L^\#)=-\mu(L)$. In case \textup{(a)},
$L^\#=\bigoplus_i\mathfrak o_Ee_i^\#$ for the dual basis ($h(e_i,e_j^\#)=\delta_{ij}$), whose Gram matrix
is $J^{-1}$.}
\end{enumerate}
\emph{Proof.} (a) For $r_1=1$ and $J=1$: the covolume of $\mathfrak o_E$ in $\bigoplus_{\tau\in\Phi}\C$ is
$2^{-2n}|\mathrm{disc}\,E|^{1/2}$. Indeed, the $4n\times4n$ matrix with rows
$(\tau(\omega_j),\overline{\tau(\omega_j)})_{\tau\in\Phi}$ has determinant of absolute value
$|\mathrm{disc}\,E|^{1/2}$, and passing from $(\tau(y),\overline{\tau(y)})$ to
$(\re\tau(y),\mathrm{Im}\,\tau(y))$ divides the determinant by $2$ for each $\tau$. So
$\log\mathrm{covol}(\mathfrak o_E)=2n\,\mathrm{cst}_E$ and $\mu(\mathfrak o_E)=0$. In general choose at each
$\tau$ a matrix $A_\tau$ with $\tau(J)=\overline A_\tau^{\,t}A_\tau$. The map $y\mapsto A_\tau y$ carries $h_\tau$
to the standard form on $\C^{r_1}$. As a real linear map of $\C^{r_1}=\mathbb R^{2r_1}$ it has determinant
$|\det A_\tau|^2=\det\tau(J)$: for $r_1=1$ multiplication by a complex number $\alpha$ is a rotation
followed by the scaling by $|\alpha|$, of determinant $|\alpha|^2$; in general write $A_\tau$ as a
product of a unitary matrix, which preserves volume, and a triangular one. So
$\mathrm{covol}(L)=\mathrm{covol}(\mathfrak o_E)^{r_1}\prod_\tau\det\tau(J)$.
(b) $\mathrm{covol}(L')=[L:L']\,\mathrm{covol}(L)$.
(c) In case (a) write $y=\sum_iy_ie_i^\#$; then $h(e_i,y)=y_i$, so $y\in L^\#$ exactly when all
$y_i\in\mathfrak o_E$. If $e_j^\#=\sum_la_{lj}e_l$, then $\delta_{ij}=\sum_lJ_{il}a_{lj}$, so $(a_{lj})=J^{-1}$,
and the Gram matrix of the dual basis is $\overline{J^{-1}}^{\,t}JJ^{-1}=J^{-1}$. For a general lattice
we use two cited facts. First (Steinitz; from memory), every lattice in $\mathsf X$ has the form
$M=\mathfrak o_Ee_1\oplus\dots\oplus\mathfrak o_Ee_{r_1-1}\oplus\mathfrak be_{r_1}$ for a basis $(e_i)$ and
an ideal $\mathfrak b$; since $[\mathfrak b:\mathfrak a\mathfrak b]$ is the norm $N(\mathfrak a)$, this
gives $\mathrm{covol}(\mathfrak aM)=N(\mathfrak a)^{r_1}\mathrm{covol}(M)$ for an ideal $\mathfrak a$, and
then for a quotient of ideals, with the norm extended multiplicatively. Second
(\cite[Ch.~III, \S2]{Neu99}; from memory), the set $\mathfrak d^{-1}$ of the $a\in E$ with
$\mathrm{Tr}_{E/\Q}(a\,\mathfrak o_E)\subset\Z$ is the inverse of an ideal $\mathfrak d$ of norm
$|\mathrm{disc}\,E|$. Now put $T(y,y')=\mathrm{Tr}_{E/\Q}h(y,y')$. It is a symmetric $\Q$-bilinear form,
and $\frac12T(y,y)=\|y\|^2$, because the trace of an element of $F$ is twice the sum of its images under
the $\tau\in\Phi$. Let $L^{*}$ be the set of the $y$ with $\frac12T(y',y)\in\Z$ for all $y'\in L$. By the
definition of $\mathfrak d^{-1}$, and because $L$ is an $\mathfrak o_E$-module, $L^*$ is the set of $y$
with $h(L,y)\subset2\mathfrak d^{-1}$, which is $2\mathfrak d^{-1}L^\#$ (if $h(L,y)\subset\mathfrak a$, then
$\mathfrak a^{-1}y\subset L^\#$, so $y\in\mathfrak aL^\#$). So
$\mathrm{covol}(L^*)=2^{4nr_1}|\mathrm{disc}\,E|^{-r_1}\mathrm{covol}(L^\#)$. But $L^*$ is the dual of $L$
for the Euclidean inner product, and a discrete subgroup of full rank and its Euclidean dual have
covolumes with product $1$ (their bases are related by the inverse transpose). Hence
$\mathrm{covol}(L)\,\mathrm{covol}(L^\#)=(2^{-2n}|\mathrm{disc}\,E|^{1/2})^{2r_1}=\mathrm{covol}(\mathfrak o_E^{r_1})^2$,
and $\mu(L)+\mu(L^\#)=2\mu(\mathfrak o_E^{r_1})=0$. \qed

\noindent Part (b) is Lemma~L1 of \S III.6 for the field $E$; the form with lengths at places is
\eqref{lt:index}.

\emph{Example} (script \texttt{o3\_lattice.py}; here the field is $K=\Q(\sqrt{-23})$, in the role of $k_d$, with $4n$
replaced by $2$). Let $\omega=\frac12(1+\sqrt{-23})$, $\mathfrak o=\Z[\omega]$, and
$\mathfrak p=(2,\omega)$, a prime ideal of norm $2$ that is not principal. In $\mathsf X=k_d^2$ with
$h(y,y')=\bar y_1y_1'+\bar y_2y_2'$ let $L=\mathfrak o\oplus\mathfrak p$; it is not free. Then
$L^\#=\mathfrak o\oplus\bar{\mathfrak p}^{-1}$, and $\bar{\mathfrak p}^{-1}=\frac12\mathfrak p=\Z+\Z\frac\omega2$
because $\mathfrak p\bar{\mathfrak p}=(2)$. The covolumes of $\mathfrak o^2$, $L$, $L^\#$ are $\frac{23}4$,
$\frac{23}2$, $\frac{23}8$, and the slopes are $0$, $-\frac12\log2$, $+\frac12\log2$, in agreement with (b)
for the index $2$, and with (c).

\textbf{Definition O10 (classes, genus, automorphisms).} Two lattices $L,L'$ in $\mathsf X$ are
\emph{isometric} if $L'=\gamma L$ for some $\gamma\in U(\mathsf X)$; the \emph{class} of $L$ is the set of
such $L'$. The \emph{automorphism group} of $L$ is $\Gamma_L=\{\gamma\in U(\mathsf X):\gamma L=L\}$. For a
prime ideal $\mathfrak p$ of $\mathfrak o_F$ let $F_{\mathfrak p}$ be the completion of $F$ for the absolute
value $|y|_{\mathfrak p}=N(\mathfrak p)^{-j}$, where $\mathfrak p^j$ is the exact power of $\mathfrak p$ in
$y\,\mathfrak o_F$; it is constructed from $F$ as $\Q_p$ is from $\Q$ (Part~II, \S6.8), and its elements
with $|y|_{\mathfrak p}\le1$ form a ring $\mathfrak o_{\mathfrak p}$. Let $\mathsf X_{\mathfrak p}$ be the
$F_{\mathfrak p}$-vector space with the same basis as the $F$-vector space $\mathsf X$; the form $h$ and
the multiplication by $E$ extend to it. Let $L_{\mathfrak p}\subset\mathsf X_{\mathfrak p}$ be the set of
the finite sums $\sum a_iy_i$ with $a_i\in\mathfrak o_{\mathfrak p}$, $y_i\in L$. Two lattices $L,L'$ are
\emph{in the same genus} if for every $\mathfrak p$ there is a bijective map of $\mathsf X_{\mathfrak p}$ to
itself that commutes with $E$, preserves $h$, and carries $L_{\mathfrak p}$ to $L'_{\mathfrak p}$. A genus is
a union of classes. Two Hermitian spaces are \emph{isometric over $F_{\mathfrak p}$} if there is such a
map from the one $\mathsf X_{\mathfrak p}$ to the other.

\emph{Example} (with $K=\Q(\sqrt{-23})$ in place of $k_d$, and $\Q$ in place of $E$ and $F$). Let $\mathsf X=k_d$ with
$h(y,y')=\bar yy'$. A lattice is a quotient of ideals $\mathfrak a$, and $\mathfrak a^\#=\bar{\mathfrak a}^{-1}$.
The unimodular lattices are the $\mathfrak a$ with $\mathfrak a\bar{\mathfrak a}=\mathfrak o$, that is
$\mathfrak a=\mathfrak b\bar{\mathfrak b}^{-1}$. They form one genus: at each prime $\mathfrak b$ is
generated by one element $\beta$, and multiplication by $\beta/\bar\beta$ is an isometry there. Two of
them are isometric exactly when they differ by a principal ideal. The ideal class group of $k_d$ has
order $3$ and is generated by the class of $\mathfrak p$, so the genus has three classes, represented by
$\mathfrak o$, $\mathfrak p\bar{\mathfrak p}^{-1}$, $\bar{\mathfrak p}\mathfrak p^{-1}$. Each has the
automorphism group $\{\pm1\}$.

\textbf{Lemma O11 (finiteness of $\Gamma_L$).} \emph{If $\mathsf X$ is positive definite, $\Gamma_L$ is
finite.}

\emph{Proof.} $\Gamma_L$ preserves the Euclidean norm of Definition~O8 and the discrete set $L$. A
$\Z$-basis of $L$ can be sent only to vectors of $L$ of the same norms, and there are finitely many of
those. \qed

\noindent (C5) \emph{Finiteness of the class number} (Borel, \emph{Some finiteness properties of adele
groups over number fields}, 1963, Theorem~5.1; from memory): a genus consists of finitely many classes. So
for a positive definite space ``a function on the classes of a genus'' is a function on a finite set
$L_1,\dots,L_\kappa$, and the weights $1/|\Gamma_{L_i}|$ of Part~II, \S6.1 are defined. A \emph{level
structure} at a prime $\mathfrak p$ of $F$ is an additional datum attached to $L/\mathfrak pL$ (or to
$L^\#/L$), which isometries and automorphisms are then required to respect; classes, genera and the
groups $\Gamma_L$ of lattices with level structure are defined in the same way, and (C5) holds for them.

\subsection*{3.\quad The standing parameters}

The parameters of the construction are assumed to exist, with the sizes of the following table. The
middle column gives the condition under which the datum is constructed, and, after ``record:'', the
weaker condition that the estimates of \S III.4--III.9 use. The last column gives the rounded size used
in Part~I, \S I.5 and \S I.8, and in Part~II.

\textbf{Hypothesis O23 (the parameters; assumed here, from the Lean development).} \emph{There is
$\varepsilon_0$ with $0<\varepsilon_0<\frac1{52}$ such that for every $0<\varepsilon\le\varepsilon_0$ there
are constants $C$, $C_{\rm abs}$, $C_H$, $C_{\mathsf G}$, $C_{\rm c}$, $C_8$, $J\ge1$ and an unbounded set
of real numbers $x$, the admissible scales, with the following property. For every admissible $x$
there is $n=2^e$, and for every $u$ in the group $H$ of \S1 there are the integers, primes, sets and
numbers named in the first column of the table below, satisfying the conditions of its middle column.
The data of the rows above the row of $\eta$ are the same for all $u$.}

\noindent (M) needs the last sentence (Part~II, \S4). In the source it is expressed by statements that
tie the data of different translates (\S9, row~28); we did not verify them.

{\small\setlength{\LTleft}{0pt}\setlength{\LTright}{0pt}
\begin{longtable}{@{}p{.215\textwidth}p{.50\textwidth}p{.225\textwidth}@{}}
\toprule
Parameter and meaning & Condition (assumed in O23) & Parts I and II\\
\midrule\endhead
\bottomrule\endfoot
$\varepsilon$; $x$, the scale & $0<\varepsilon\le\varepsilon_0<\frac1{52}$; $x$ admissible & $0<\varepsilon<\frac1{52}$; $x$ in an interval $[Z,Z^{3/2}]$\\
$b$, $a$ & $b=\exp(x^{\varepsilon})$, $a=xb$ & the same\\
$X'$, $b'$, $w''$, $X''$, $\delta_1$ & $X'=\exp(x^{\varepsilon/3})$, $b'=a/X'$, $w''=b'/(xn)$, $X''=a/w''=xnX'$, $\delta_1=x^{-\varepsilon/8}$ & cut-off $\frac13\exp(x^{\varepsilon/3})$; the others not named\\
$n$, $N$, $f$: the table has $2$ rows and $n$ columns & $n=2^e$, $e\ge2$, $x^{2+8\varepsilon}\le n\le2x^{2+8\varepsilon}$; $N=4n$, $f=2n$. Record: $x^{2+8\varepsilon-\varepsilon/100}\le n\le x^{2+8\varepsilon+\varepsilon/100}$ & $n\approx x^{2+8\varepsilon}$ (\S I.5); $n\approx x^2$ (\S I.8)\\
$m$, the number of lifts & $m$ even, $m\ge2$, $\sqrt{x\delta_1}\le m\le2\sqrt{x\delta_1}$, that is $m\asymp x^{1/2-\varepsilon/16}$. Record: within a factor $x^{\pm\varepsilon/100}$ of $\sqrt{x\delta_1}$ & $m\approx\sqrt x$; $m=x^{1/2+O(\varepsilon)}$\\
$k$, the initial rank & $k$ even, $x^{9/2+10\varepsilon}\le k\le2x^{9/2+10\varepsilon}$; $m+3\le k$, $m\le f\le k$. Record: within $x^{\pm\varepsilon/100}$ & $k\approx x^{9/2}$\\
$t_0$, the odd parameter & $t_0$ odd, $t_0\ge k+m+4$, $t_0>3m$, $m^2\le N_m$; and
$r+t_0=8fr+4+\lceil x^{377}\rceil\,n_W$ with $n_W=8\lceil\frac{10}8r(f+f^2(k-1)^3x^{2\varepsilon}/m)\rceil$. So $t_0\asymp x^{797/2+(929/16)\varepsilon}$, and $t_0\le x^{400}$ for $x\ge2^{100}$ & $\tau_0$ odd (II, \S6.2); twists below $\lfloor x^{400}\rfloor$ (II, \S6.8)\\
$r$, $N_j$, $\varrho$: the ranks & $r=k+m-1$; $N_j=k+j(t_0+1)$ for $0\le j\le m$; $\varrho=r+t_0$ & $N_j$ (II, \S6.2)\\
$q$, $d$: the auxiliary primes; $Y$, the scale at which they are chosen & primes, $q\equiv3\bmod8$, $d\equiv7\bmod16$, $q\mid d+1$ (so $(-d/q)=1$), $q\ne3$, $q,d>4t_0r+1$; $x^{1/J}\le Y\le x$; every odd prime $p\le Y^{\varepsilon/5}$ divides $d+1$ (so $(-d/p)=1$); $\frac12Y^{\varepsilon/5}<\log q\le Y^{\varepsilon/5}$, $\frac12Y^{\varepsilon/5}\log^2Y<\log d\le Y^{\varepsilon/5}\log^2Y$. Record: $(-d/q)=1$; $(-d/p)=1$ for the odd primes $p$ up to a bound that is at least $11$ and at least $\lfloor x^{\varepsilon_z/5}\rfloor$; $x^{\varepsilon_z/5-\varepsilon_z/100}\le\log q,\log d\le x^{\varepsilon/5+\varepsilon/100}$, where $\varepsilon_z=\varepsilon/J$ & $q\mid d+1$; $\log q\approx Y^{\varepsilon/5}$, $\log d\approx Y^{\varepsilon/5}(\log Y)^2$\\
$g_0$, $g_1$: the level primes & distinct primes, $g_i\equiv1\bmod Nqd$, $(kb)^8\le g_i\le2(kb)^8$. Record: $(kb)^8\le g_i\le2g_j$, $\log g_i\le x^{\varepsilon+\varepsilon/100}$ & $\ell_g=\log(g_0g_1)$ (II, \S6.7)\\
$g_2$, $g_3$, $g_4$: three further level primes & primes, at most $(Nqd)^{C_8}$; $g_2\ne g_3$, both different from $g_0,g_1,g_4$; $g_2\equiv g_3\equiv1\bmod\mathrm{lcm}(15,Nqd)$, $g_4\equiv1\bmod Nqd$. So $g_0,\dots,g_4$ split completely in $E$ (Lemma~O3). Record: two further level primes, distinct, $\equiv1\bmod15$ & not mentioned\\
$H$; a translate $u$ & $H$ cyclic of order $n$ (Lemma~O2); $u$ runs over $H$; there is one set of data for each $u$ & the same\\
$T_2$, $\delta$: the thick columns & a set of integers $i\ge0$ with $4L_{\rm th}i<n$ and property (b) of Theorem~\ref{thm:T2} for $L=L_{\rm th}$, where $L_{\rm th}=10\lceil405\log x\rceil+2\lceil30/\varepsilon\rceil+22$; as a set of columns, $\{[5]^i:i\in T_2\}$; $\delta=|T_2|/n\ge x^{-\varepsilon}$ & $\delta\ge x^{-\varepsilon}$\\
$T$: the labels & $T\subset H$ as in Theorem~\ref{thm:T}, $\frac12\delta_1n\le|T|\le2\delta_1n$. Record: $x^{-\varepsilon/100}\delta_1n\le|T|\le2\delta_1n$ & $|T|/n\ge\frac12x^{-\varepsilon/8}$\\
$a_t$, $t\in T$: the orders & integers, $a\le a_t\le a\,(1+1/(kx^{13}))$ & $|a_t/b-x|\le1$\\
$b_t$, $t\in T$ & integers, $b'\le b_t\le b'\,(1+2n\delta_1x^{\varepsilon}\Lambda_G/k)$, with $x^\varepsilon\le\Lambda_G\le x^{\varepsilon+\varepsilon/100}$ & not named\\
$\eta$: the width of the weights & $t_0\le\eta\le b/(t_0^3fk^3)$, $m(t_0+1)\eta\le w''$, $6f(x+1)(\eta+1)\le b/x$ & $O(b/(nx^2))$\\
$w_{u,j}(v,s)$: the weight of the cell $(v,s)$ for the translate $u$ and the lift $j$ & integers. Thick cells $(1,u[5]^i)$, $i\in T_2$: $b\le w\le b+\eta+1$. Top cells $(0,t)$, $t\in T$: $b_t+b'/x\le w\le b_t+b'/x+\eta$. All other cells (thin): $w''\le w\le w''+\eta$ & thick: $b\le w\le b+O(b/(nx^2))$; others: $a_t/w\ge\frac13\exp(x^{\varepsilon/3})$\\
$\Delta$, $S_{\rm c}$ & integer $\Delta\ge N_mm^{10}$, $\log\Delta\le\frac12x^{\varepsilon/4}$; $S_{\rm c}=4k\Delta$ & not named\\
$\mathsf W_u$: the words & a finite set, $|\mathsf W_u|\le2^{k-3}$ & $\log|\mathsf W|\le(k-3)\log2$\\
$\mathrm{Lv}(\mathsf w)$: the level prime of a word & a prime $\pi\bar\pi$ with $\pi\in\mathfrak o_{k_d}$, $\pi\equiv1\bmod Nqd(g_0g_1)^2$; $\mathrm{Lv}(\mathsf w)\ge b^{m/2}$, $\log\mathrm{Lv}(\mathsf w)\le Cmx^\varepsilon$; different from $g_2$, $g_3$ & level $\mathcal T$ (II, \S6.7)\\
$B_0$, $N_c$, $L_a$; $B$ & definitions: $B_0=fmbx^\varepsilon$, $N_c=|T_2|kbx^{2\varepsilon}/n^2$, $L_a=\log(a\,\mathrm{rd}_E)$; $B=|T|\,m\,b\,\delta/2$, the unit of (M) & as in \eqref{III:sizefinal}; the same\\
Constants & $C$, $C_{\rm abs}$, $C_H$ (in the bound $x^{C_H}b$ of II, \S6.10), $C_0=3$, $C_{\mathsf G}$ (in $\mathsf G=\exp(C_{\mathsf G}m^5fx^\varepsilon)$), $C_{\rm c}$ (in O24), $C_8$, $J\ge1$; they do not depend on $x$, $u$, $\mathsf w$ & $C$ absolute\\
\end{longtable}}

\noindent The conditions on the numbers $\xi,\psi,\theta$ attached to the weights are in Hypothesis~O24.

\textbf{Remarks on the table: where Parts~I and~II are rounded or differ.} These are findings about
Parts~I and~II; none affects a proof there.
\begin{enumerate}[label=\textup{(\arabic*)},leftmargin=2.2em,itemsep=1pt,topsep=2pt]
\item \emph{Exponents.} With the exact sizes, the four errors of (M) are $x/m=x^{1/2+\varepsilon/16}$,
$xm/k=x^{-3-(161/16)\varepsilon}$, $m=x^{1/2-\varepsilon/16}$ and $k/n^2=x^{1/2-6\varepsilon}$, up to bounded
factors. Part~I, \S I.8 uses $m\approx\sqrt x$, $k\approx x^{9/2}$, $n\approx x^2$; the statement of
Estimate~A rounds $m$ and $k$ but not $n$. The conclusion $x^{1/2+C\varepsilon}$ is not affected.
\item \emph{The number of twists.} Part~II, \S6.8 states the bound \eqref{II:shabound} for the twists
with index below $\lfloor x^{400}\rfloor$; this is the form in which the source supplies the bound
(\S III.9). The construction uses the twists with index $1,\dots,t_0-1$, and $t_0\le x^{400}$ for
$x\ge2^{100}$ (a statement of the source, quoted in \S III.9), so \eqref{II:shabound} covers them.
\item \emph{The thick weights.} Part~II, \S4.2 has $b\le w\le b+O(b/(nx^2))$. Here
$b\le w\le b+\eta+1$, and $\eta+1\le b/(6f\,x(x+1))$; the two agree.
\item \emph{The rank example of Part~II, \S6.2.} Besides the four numerical conditions, the table has
$m\le2n\le k$. Since $n\ge4$, this excludes $k=6$, and the smallest triple allowed by the rank
conditions is $(m,k,t_0)=(2,8,15)$, with ranks $8,24,40$ (\S8); Part~II, \S6.2 uses this triple. In the
construction $t_0$ is also fixed by the equation for $r+t_0$ in the table, which makes it of size about
$x^{797/2}$; so this triple only illustrates the rank conditions.
\item \emph{The running example} $q=11$, $d=1319$, $N=16$ of \S I.5 has the right congruences, $q\mid d+1$,
$q\ne3$, and the smallest $N$ allowed; and the odd primes $3,5,7,11$ split in $k_d$. It violates
$q,d>4t_0r+1$ (at least $541$ even for the triple of Remark~(4), and much larger once the equation for $t_0$
is imposed; here $q=11$ is below even $541$, and $d=1319$ below the latter); the condition that the odd
primes up to a bound $\ge11$ split in $k_d$ whenever that bound is $13$ or more ($13$ is inert); the
condition of Theorem~\ref{thm:pair} that every odd prime up to $Y^{\varepsilon/5}$ divide $d+1$, as soon
as that bound is $7$ or more ($d+1=1320$ is prime to $7$); and all conditions of size.
(The example of the first version, $q=3$, $d=23$, also violated $q\ne3$, and there $5,7,11$ are inert.) Part~II, \S4.1 and \S I.5 both say that its parameters do not satisfy the conditions of (M).
No proof uses the example.
\item \emph{The hypothesis on $q$ and $d$.} The estimates assume only $(-d/q)=1$ (the record). The
construction has $q\mid d+1$, as Part~I has; it implies $(-d/q)=1$, and the converse fails (for example
$q=11$, $d=7$). Parts~I and~II use $q\mid d+1$ only to get $\chi_d(q)=1$.
\item \emph{Not in Parts~I and~II:} that the labels $t\in T$ are places of row $0$ (Lemma~O5); the numbers
$b_t$ and the weights of the top cells; three further level primes; the cut-off $X''$ and the number
$w''$; the parameter $J$. (Of these, $b_t$, $X''$, $w''$ and $J$ occur in \S III.4--III.9.) In the data
used by the final theorem $J$ is a constant $\ge1$ whose value is not specified.
\end{enumerate}

\textbf{Status in the source} (statements read, no proof read; the declarations are in \S9, rows 2 and
8--15). (i) \emph{Existence.} The source asserts, as theorems, that data with the properties of O23, O12,
O24 and O25 exist for $0<\varepsilon\le\varepsilon_0$ and all admissible $x$; one of these theorems is
among the declarations the final theorem depends on. (ii) \emph{The field.} The source does not define
$E$ by generators. On the route of the final theorem it is the subfield of a cyclotomic field fixed by
the residues $a$ with $a\equiv\pm1\bmod2^{e+2}$ and $\chi_q(a)=\chi_d(a)=1$; by Lemma~O2(a) this is the
field $E$ of \S1. This identification is our inference from the definitions, not a statement of the
source. (iii) \emph{The sets $T_2$ and $T$} are installed in the data by the construction of the route of
the final theorem: $T_2$ from the thick-set theorem (Theorem~\ref{thm:T2}), then $T$ from the random
choice (Theorem~\ref{thm:T}); both are among the declarations the final theorem depends on. (iv)
\emph{The constant $J$} is obtained existentially, with $J\ge1$, from the Galois-cohomology input
(\S III.9, Theorem~G); it is not fixed to $1$. The primes $q,d$ are chosen at a scale $Y$ with
$x^{1/J}\le Y\le x$; this is why the record's lower bound for $\log q$ has the exponent
$\varepsilon_z=\varepsilon/J$. (v) \emph{Not located in the source:} the number $B$ and the constants
$C_H$, $C_0$, $C_{\rm abs}$ of the last two rows; they are taken from Part~II and \S III.4--III.8.

\subsection*{4.\quad The spaces, the lattices and the chain of lifts}

\textbf{Hypothesis O12 (the spaces and lattices; assumed here, from the Lean development).} \emph{Fix
$x$, a translate $u\in H$ and the parameters of \textup{O23}. There are Hermitian spaces over $E$ as
follows.}
\begin{enumerate}[label=\textup{(\roman*)},leftmargin=2.4em,itemsep=1pt,topsep=2pt]
\item \emph{$\mathsf X_j$ for $0\le j\le m$, positive definite of rank $N_j$.}
\item \emph{$\mathsf X'$, the orthogonal complement in $\mathsf X_m$ of a vector $y_0$ with
$h(y_0,y_0)\ne0$; it is positive definite of rank $N_m-1$.}
\item \emph{$\mathsf Y_j$ for $0\le j\le m$, of rank $r+jt_0$ and of signature
$\big(r+jt_0-\lfloor\frac{m-j}2\rfloor,\lfloor\frac{m-j}2\rfloor\big)$ at every $\tau\in\Phi$, with an
$E$-linear map $\mathsf X_j\to\mathsf Y_j$ preserving the forms for $j<m$, and a bijective such map
$\mathsf X'\to\mathsf Y_m$. The space $\mathsf Y_0$ is the orthogonal sum of $\mathsf X_0$, a line
spanned by a vector $y$ with $h(y,y)=-\mathrm{Lv}$, and $\frac m2-1$ hyperbolic planes; here
$\mathrm{Lv}=\mathrm{Lv}(\mathsf w_0)$ for a base word $\mathsf w_0$.}
\item \emph{$\mathsf Z_j$ for $1\le j\le m$, of rank $r$ and signature $(k+j-1,m-j)$ at every
$\tau\in\Phi$, isometric to $\mathsf Y_0$ over $F_{\mathfrak p}$ for every prime $\mathfrak p$ of $F$.}
\item \emph{$\mathsf V_j$ for $1\le j\le m$, definite of rank $\varrho=r+t_0$, containing a free lattice
with a basis that is orthonormal for $h$ or for $-h$.}
\item \emph{For every word $\mathsf w\in\mathsf W_u$ and every $0\le j\le m$ there are a finite set
$\mathrm{Cl}_j(\mathsf w)$ of lattices in $\mathsf X_j$, and for
$L\in\mathrm{Cl}_j(\mathsf w)$ a finite group $\Gamma_L$.}
\item \emph{For $j=0$: there is a lattice $L_{\rm b}$ in $\mathsf X_0$ with $L_{\rm b}\subset L_{\rm b}^\#$,
$\mathrm{Lv}(\mathsf w)L_{\rm b}^\#\subset L_{\rm b}$ and
$[L_{\rm b}^\#:L_{\rm b}]=\mathrm{Lv}(\mathsf w)^{2f}$; the set $\mathrm{Cl}_0(\mathsf w)$ is a set of
representatives of the classes in the genus of $L_{\rm b}$ with a level structure at the primes above
$g_0$, $g_1$ and $\mathrm{Lv}(\mathsf w)$; and $\Gamma_L$ is the group of the automorphisms of $L$ that
respect it.}
\item \emph{For $1\le j\le m$: $\mathrm{Cl}_j(\mathsf w)$ is a set of representatives of the classes of a
set $\mathcal G_j(\mathsf w)$ of lattices with level structure in $\mathsf X_j$ that is a union of classes,
and $\Gamma_L$ is the group of the automorphisms of $L$ that respect the level structure.}
\end{enumerate}

\noindent \emph{The level structure at stage $0$.} A prime of $F$ above $g_0$ or $g_1$ splits into two
primes $\mathfrak P,\overline{\mathfrak P}$ of $E$ (Lemma~O3). The structure there contains a \emph{full
flag} in the vector space $L/\mathfrak PL$ over $\mathfrak o_E/\mathfrak P$: a chain of subspaces
$0\subset V_1\subset\dots\subset V_{N_0}=L/\mathfrak PL$ with $\dim V_i=i$; we did not unfold the rest. At a
prime above $\mathrm{Lv}(\mathsf w)$ it is a nonzero linear form on $L_{\rm b}^\#/L_{\rm b}$. We use no
more of it.

\noindent \emph{How O12 compares with the source} (\S9, row~16). In (i) the record requires
``definite''; that the spaces are positive definite is a further statement of the source
(\texttt{X\_posDef}, which is in the dependency lists). In (iv) the source has one statement for all $\mathfrak p$ together, which implies
the one here. Clause (vii) is not part of the record of a translate; it is a separate datum, assumed in
later statements of the source. Clause (viii) \emph{was not found in the source}: for $j\ge1$ the record
has the groups $\Gamma_L$ only as abstract finite groups, and a level datum only at the primes above
$\mathrm{Lv}(\mathsf w)$. We assume (viii) because Definition~O14 and the chain \eqref{ob:chain} use
$1/|\Gamma_L|$ as the inverse of an order of automorphisms. Only the spaces $\mathsf X_j$, $\mathsf X'$
and $\mathsf Y_0$ are used below; $\mathsf Z_j$ and $\mathsf V_j$ occur in statements quoted in
\S III.7--III.9.

\textbf{Lemma O13 (ranks and signatures).} \emph{Let $m,k$ be even, $t_0$ odd, $r=k+m-1$,
$N_j=k+j(t_0+1)$. Then: $r$ is odd and every $N_j$ is even; $N_m=(r+mt_0)+1$; the rank of $\mathsf Y_j$
exceeds $N_j$ by $m-1-j$ for $j\le m$; $N_0+1+2(\frac m2-1)=r$; the signature of $\mathsf Y_0$ in
\textup{(iii)} is $(k+\frac m2-1,\frac m2)$, which is that of $\mathsf Z_{m/2}$; and the signatures of
$\mathsf Y_m$ and $\mathsf Z_m$ are $(r+mt_0,0)$ and $(r,0)$.}

\emph{Proof.} Direct substitution. For instance $r+mt_0+1=k+m+mt_0=N_m$, and for $j=0$,
$\lfloor\frac m2\rfloor=\frac m2$ and $r-\frac m2=k+\frac m2-1$. \qed

\noindent So, under O12, $\mathsf Y_m$ and $\mathsf Z_m$ are definite, the orthogonal complement of
$\mathsf X_j$ in $\mathsf Y_j$ has rank $m-1-j$ and contains all the negative directions, and the
codimension of $\mathsf X'$ in $\mathsf X_m$ is one. (Script \texttt{o4\_ranks.py} checks these identities for $527$
triples.)

\textbf{Definition O14 (algebraic modular forms).} Let $\mathsf X$ be a positive definite Hermitian
space of rank $r_1$, and $\mathcal G$ a set of lattices in $\mathsf X$ that is a union of finitely many
classes (for instance a class, a genus, or a genus with a level structure). For $\tau\in\Phi$ let
$U(\mathsf X_\tau)$ be the group of the complex linear maps of $\mathsf X_\tau$ that preserve $h_\tau$; in
an orthonormal basis it is the group $U(r_1)$ of the unitary $r_1\times r_1$ matrices. Put
$U_\infty=\prod_{\tau\in\Phi}U(\mathsf X_\tau)$, a product of $f$ such groups; an element $\gamma$ of
$U(\mathsf X)$ gives the element $(\gamma_\tau)_\tau$ of $U_\infty$. A \emph{unitary representation} of
$U_\infty$ is a finite-dimensional complex vector space $W_\rho$ with an inner product, together with a
continuous homomorphism $\rho$ of $U_\infty$ into the group of the linear maps of $W_\rho$ that preserve
the inner product. It is \emph{irreducible} if $W_\rho\ne0$ and no subspace other than $0$ and $W_\rho$
is mapped to itself by all $\rho(g)$. A \emph{form of weight $\rho$ on $\mathcal G$} is a map
$\phi:\mathcal G\to W_\rho$ with
\begin{equation}\label{D:364}
 \phi(\gamma L)=\rho(\gamma)\,\phi(L)\qquad(\gamma\in U(\mathsf X),\ L\in\mathcal G).
\end{equation}
The forms make a complex vector space $\mathcal F(\mathcal G,\rho)$. If $L_1,\dots,L_\kappa$ represent the
classes in $\mathcal G$, then $\phi\mapsto(\phi(L_i))_i$ is an isomorphism of $\mathcal F(\mathcal G,\rho)$
onto $\bigoplus_iW_\rho^{\Gamma_{L_i}}$, the vectors fixed by the finite groups $\Gamma_{L_i}$. It carries
the inner product
\begin{equation}\label{D:365}
 \langle\phi,\phi'\rangle=\sum_{i=1}^\kappa\frac1{|\Gamma_{L_i}|}\,\big\langle\phi(L_i),\phi'(L_i)\big\rangle_{W_\rho} .
\end{equation}
It does not depend on the representatives: $\Gamma_{\gamma L}=\gamma\Gamma_L\gamma^{-1}$ has the order of
$\Gamma_L$, and $\rho(\gamma)$ preserves the inner product of $W_\rho$. For $W_\rho=\C$ with trivial
$\rho$ a form is a function on the finite set of classes, as in Part~II, \S6.1.

\emph{Examples.} (1) The group $U(1)$ is the circle $|g|=1$. For an integer $\lambda_1$,
$\rho(g)=g^{\lambda_1}$ on $\C$ is an irreducible unitary representation. (2) For integers $a\ge0$ and
$a'$, the group $U(2)$ acts on the homogeneous polynomials $P$ of degree $a$ in two variables by
$(\rho(g)P)(X)=\det(g)^{a'}P(Xg)$, $X=(X_1,X_2)$; this representation is irreducible of dimension $a+1$.
(3) Take $\Q(i)$ in place of $E$, one embedding, $\mathsf X=\Q(i)$ with $h=\bar yy'$, and $\mathcal G$ the
class of $L=\Z[i]$. Then $\Gamma_L=\{\pm1,\pm i\}$. For $\rho(g)=g^{\lambda_1}$ a form is a number
$\phi(L)$ with $i^{\lambda_1}\phi(L)=\phi(L)$: the space of forms has dimension $1$ if $4\mid\lambda_1$ and
$0$ otherwise, and $\langle\phi,\phi\rangle=\frac14|\phi(L)|^2$.

\textbf{Definition O15 (weights, and the interlacing condition).} We cite the classification of the
irreducible unitary representations of $U(r_1)$ (H.~Weyl; \cite{Wey25}; from memory). In such a
representation the diagonal matrices $\mathrm{diag}(g_1,\dots,g_{r_1})$ have a common basis of
eigenvectors, with eigenvalues $g_1^{\nu_1}\cdots g_{r_1}^{\nu_{r_1}}$ for lists of integers $\nu$. The
list that is largest in the lexicographic order is the \emph{highest weight}; it satisfies
$\lambda_1\ge\dots\ge\lambda_{r_1}$, it determines the representation, and every such list occurs. In
Example~(1) it is $(\lambda_1)$; in Example~(2) it is $(a+a',a')$. Multiplying a representation by
$\det(g)^{a'}$ adds $a'$ to every $\lambda_i$. The \emph{shifted weight} is the strictly decreasing list
$\lambda^\circ_i=\lambda_i+\frac{r_1+1}2-i$; its entries lie in $\Z$ for $r_1$ odd and in $\Z+\frac12$ for
$r_1$ even. The representation has dimension $\prod_{i<j}(\lambda^\circ_i-\lambda^\circ_j)/\prod_{j<r_1}j!$
(Weyl's dimension formula; Part~II, \S6.6); in Example~(2) this is $a+1$. A \emph{weight} for $\mathsf X$
is a choice of such a list $\lambda(\tau)$ at each $\tau\in\Phi$. Its representation of $U_\infty$ is the
product of the corresponding representations $\rho_\tau$ of the factors: the space $W_\rho$ has as a basis
the tuples $(w_\tau)_\tau$ of basis vectors of the spaces $W_{\rho_\tau}$, and $(g_\tau)_\tau$ acts by the
Kronecker product of the matrices $\rho_\tau(g_\tau)$. The \emph{branching rule} (Weyl; from memory): the
representation of $U(r_1-1)$ with shifted weight $\psi$ occurs in the restriction of the representation
of $U(r_1)$ with shifted weight $\xi$ exactly when the two lists \emph{interlace},
$\xi_0>\psi_0>\xi_1>\dots>\psi_{r_1-2}>\xi_{r_1-1}$, and then it occurs once. (For example the
representation of $U(3)$ with highest weight $(4,2,0)$, of dimension $27$, restricts to the sum of the
nine representations of $U(2)$ with highest weights $(\lambda'_1,\lambda'_2)$,
$4\ge\lambda'_1\ge2\ge\lambda'_2\ge0$, whose dimensions add up to $27$; script \texttt{o3\_lattice.py}.)
For a number $y$ and an integer $l\ge1$ write $\mathrm{bl}(y,l)$ for the list of $l$ numbers in
arithmetic progression with step $1$ and middle point $y$.

\textbf{Hypothesis O24 (the weights of the construction; assumed here, from the Lean development).}
\emph{Fix $x$, $u$ and a word $\mathsf w$. At each place $\tau\in\Phi$ there are numbers}
\begin{equation}\label{ob:inter}
 \xi_0>\psi_0>\xi_1>\dots>\psi_{k-2}>\xi_{k-1}>-\theta_0,\qquad
 \theta_0<w_{1}<\theta_1<w_{2}<\dots<\theta_{m-1}<w_{m},
\end{equation}
\emph{with $\xi_i\in\Z+\frac12$ and $\psi_i,\theta_j,w_j\in\Z$, and with all gaps $w_j-\theta_{j-1}$ and
$\theta_j-w_j$ at least $(t_0+3)/2$, such that \textup{(a)--(d)} hold.}
\begin{enumerate}[label=\textup{(\alph*)},leftmargin=2.2em,itemsep=1pt,topsep=2pt]
\item \emph{$w_j=w_j(\tau)$ is the weight $w_{u,j}(v,s)$ of the table of \S3 for the cell $(v,s)$ of
$\tau$. The numbers $\psi,\theta,w$ do not depend on the word; the numbers $\xi$ do.}
\item \emph{For $0\le j\le m$ the forms at stage $j$ have, at the place $\tau$, the highest weight
$\lambda^{(j)}(\tau)$ of rank $N_j$ with entries
$\lambda^{(j)}_i=\xi_{i-1}-\frac{k+1}2+i$ for $1\le i\le k$, followed, for $j'=1,\dots,j$, by $t_0+1$
entries equal to $-w_{j'}(\tau)+\frac k2+(2j'-1)\frac{t_0+1}2$.}
\item \emph{At a top place $(0,t)$, $t\in T$,}
\begin{equation}\label{ob:top}
 \xi_0=a_t+\tfrac{k-1}2,\qquad \xi_{k-1}=-b_t-\tfrac{k-1}2 .
\end{equation}
\item \emph{At the other places $S_{\rm c}\le\xi_0-\psi_0\le C_{\rm c}S_{\rm c}$, and at all places
$0\le\psi_i\le S_{\rm c}$.}
\end{enumerate}

\noindent The display \eqref{ob:inter} is \eqref{tr:inter} of \S III.8, where $w_{j+1}$ is written $e_j'$.
By (a), \emph{the weights of the table are the parameters $w_j$ of the chain}. Put
\begin{equation}\label{ob:Xpar}
 \{\xi_0,\dots,\xi_{k-1}\}\ \cup\ \bigcup_{i=1}^{j}\mathrm{bl}\big(-w_i(\tau),\,t_0+1\big),
\end{equation}
a set of $N_j=k+j(t_0+1)$ half-integers; the gaps in \eqref{ob:inter} make the blocks disjoint. The set
\eqref{ob:Xpar} is the \emph{parameter} of the forms at stage $j$. By a direct computation from (b), the
shifted weight of $\lambda^{(j)}(\tau)$ is \eqref{ob:Xpar} translated by $j(t_0+1)/2$. So
\eqref{ob:Xpar} is the shifted weight of the representation obtained from that of highest weight
$\lambda^{(j)}(\tau)$ by multiplying it by $\det^{-j(t_0+1)/2}$; for $j=0$ it is the shifted weight of
$\lambda^{(0)}(\tau)$ itself. \emph{Each lift adds
to the parameter a block of $t_0+1$ consecutive numbers centred at $-w_j$}; in the highest weight these
are $t_0+1$ equal entries. By (b) and \eqref{ob:top}, the highest weight at stage $0$ at a top place
begins with $\lambda_1=a_t$ and ends with $\lambda_k=-b_t$: \emph{the order $a_t$ is the length of the
first row of the weight at the top place $(0,t)$} (\S I.4.3). The first chain of \eqref{ob:inter} is the
interlacing condition of the branching rule for the ranks $k$ and $k-1$. Clause (b) transcribes the
list that the record attaches to stage $j$. The forms of the source are scalar functions on a finite set
lying over the classes; we did not trace how the source passes from the list to the representation
(\S9, row~18).

\textbf{The chain.} For a word $\mathsf w$ and $0\le j\le m$ let $V_j=V_j(\mathsf w)$ be the space of forms
on $\mathrm{Cl}_j(\mathsf w)$ (that is, on the classes of these lattices) of the weight
$(\lambda^{(j)}(\tau))_\tau$ of \textup{O24(b)}. The \emph{chain of lifts} is a sequence of linear maps
\begin{equation}\label{ob:chain}
 V_0\xrightarrow{\ K_0\ }V_1\xrightarrow{\ K_1\ }\cdots\xrightarrow{\ K_{m-1}\ }V_m,\qquad
 (K_j\phi)(L')=\sum_{L\in\mathrm{Cl}_j}\frac{1}{|\Gamma_L|}\,\epsilon(L,L')\,\mathbf K_j(L,L')\,\phi(L),
\end{equation}
where $\mathbf K_j(L,L')$ is a linear map between the spaces of values and $|\epsilon(L,L')|=1$. That
the construction provides such maps is part of Hypothesis~O19. This is
the form of Part~II, \S6.2 (the display defining $T_jf$ there), and the form assumed in \S III.4. The
numbering is that of \S III.4; \S III.7 writes $K_j:V_{j-1}\to V_j$. What $\mathbf K_j(L,L')$ is, is the
subject of \S5.

\subsection*{5.\quad The theta kernels}

\textbf{Definition O16 (the theta series of a pair of lattices).} Let $(\mathsf X,h)$, $(\mathsf X',h')$
be positive definite Hermitian spaces over $E$ of ranks $r_1$, $r_2$, and $L\subset\mathsf X$,
$L'\subset\mathsf X'$ lattices. (In this subsection $\mathsf X'$ is any such space, not the space of
Hypothesis~O12(ii).) Put
\begin{equation}\label{D:370}
 \mathrm{Hom}(L,L')=\{\lambda:\mathsf X\to\mathsf X'\ \text{$E$-linear}:\ \lambda(L)\subset L'\},
\end{equation}
a discrete subgroup of full rank in the $E$-vector space of the $E$-linear maps, which has dimension
$r_1r_2$. For $\tau\in\Phi$ let $\lambda_\tau:\mathsf X_\tau\to\mathsf X'_\tau$ be the complex linear map
induced by $\lambda$, and $\|\lambda\|^2=\sum_{\tau\in\Phi}\mathrm{tr}(\lambda_\tau^*\lambda_\tau)$, the adjoint
being taken for $h_\tau$ and $h'_\tau$. The variable is $z=(z_\tau)_{\tau\in\Phi}$, where $z_\tau$ is a complex
linear map $\mathsf X'_\tau\to\mathsf X_\tau$; in orthonormal bases $z_\tau$ is a complex
$r_1\times r_2$ matrix, and there are $fr_1r_2$ complex variables in all. Put
$\langle z,\lambda\rangle=\sum_\tau\mathrm{tr}(z_\tau\lambda_\tau)$ and
\begin{equation}\label{ob:theta}
 \Theta_{L,L'}(z)=\sum_{\lambda\in\mathrm{Hom}(L,L')}\exp\big(-\pi\|\lambda\|^2+\langle z,\lambda\rangle\big).
\end{equation}
For $E$ replaced by $\Q(i)$, $\Phi$ by one embedding and $L=L'=\Z[i]$ with $h=\bar yy'$ this is the
function $\Theta(z)$ of \S I.4.1, since $\mathrm{Hom}(L,L')=\Z[i]$. In this subsection $\lambda$ without
a lower index $1,2,\dots$ is an element of $\mathrm{Hom}(L,L')$, and $z$, $\nu$, $P$ have the meaning of
\S I.4.1.

\textbf{Lemma O17 (elementary properties).} \emph{Write $\|z\|^2=\sum_\tau\mathrm{tr}(z_\tau^*z_\tau)$.}
\begin{enumerate}[label=\textup{(\alph*)},leftmargin=2.2em,itemsep=1pt,topsep=2pt]
\item \emph{The series \eqref{ob:theta} converges absolutely and uniformly on bounded sets, so
$\Theta_{L,L'}$ is an entire function of the $fr_1r_2$ variables. Moreover}
\begin{equation}\label{ob:maj}
 |\Theta_{L,L'}(z)|\le S\,\exp\Big(\frac{\|z\|^2}{2\pi}\Big),\qquad
 S=\sum_{\lambda\in\mathrm{Hom}(L,L')}e^{-\pi\|\lambda\|^2/2}.
\end{equation}
\item \emph{Let $P$ be a polynomial in the matrix entries of the $\lambda_\tau$, written
$P(\lambda)=\sum_\nu P_\nu\lambda^\nu$ in the coordinates $\lambda_i$ dual to the coordinates $z_i$ of $z$,
so that $\langle z,\lambda\rangle=\sum_iz_i\lambda_i$. Let $\Theta_{L,L'}(z)=\sum_\nu c_\nu z^\nu$. Then}
\begin{equation}\label{ob:deriv}
 P(\partial)\Theta_{L,L'}(0)=\sum_\nu P_\nu\,\nu!\,c_\nu=\sum_{\lambda\in\mathrm{Hom}(L,L')}P(\lambda)\,e^{-\pi\|\lambda\|^2}.
\end{equation}
\item \emph{For $\gamma\in U(\mathsf X)$ and $\gamma'\in U(\mathsf X')$,
$\Theta_{\gamma L,\gamma'L'}(z)=\Theta_{L,L'}\big((\gamma_\tau^{-1}z_\tau\gamma'_\tau)_\tau\big)$.}
\end{enumerate}
\emph{Proof.} (a) The pairing $\langle z,\lambda\rangle$ is at most $\|z\|\,\|\lambda\|$ in absolute value,
by the Cauchy--Schwarz inequality for the Hilbert--Schmidt norm, and
$\|z\|\,\|\lambda\|\le\frac\pi2\|\lambda\|^2+\frac1{2\pi}\|z\|^2$. So the term of $\lambda$ is at most
$e^{-\pi\|\lambda\|^2/2}e^{\|z\|^2/(2\pi)}$ in absolute value. The sum $S$ converges because
$\|\cdot\|^2$ is a positive definite quadratic form on the real vector space spanned by the discrete group
$\mathrm{Hom}(L,L')$, so the number of $\lambda$ with $\|\lambda\|^2\le R$ is $O(R^{2nr_1r_2})$. (b) Expand
$e^{\langle z,\lambda\rangle}=\sum_\nu z^\nu\lambda^\nu/\nu!$; then $c_\nu=\frac1{\nu!}\sum_\lambda\lambda^\nu e^{-\pi\|\lambda\|^2}$,
the interchange being justified by (a). And $P(\partial)z^\nu|_{z=0}=P_\nu\nu!$. (c)
$\mathrm{Hom}(\gamma L,\gamma'L')=\gamma'\,\mathrm{Hom}(L,L')\,\gamma^{-1}$. The norm is unchanged, since
$\gamma_\tau$, $\gamma'_\tau$ are unitary, and
$\mathrm{tr}(z_\tau\gamma'_\tau\lambda_\tau\gamma_\tau^{-1})=\mathrm{tr}(\gamma_\tau^{-1}z_\tau\gamma'_\tau\lambda_\tau)$. \qed

\noindent The identity \eqref{ob:deriv} is \eqref{entry} of \S I.4.1: \emph{applying $P(\partial)$ at the
origin inserts the weight $P(\lambda)$ into the sum over the lattice}. After the change of variable
$z\mapsto\sqrt\pi z$ the bound \eqref{ob:maj} is the hypothesis $|F(z)|\le Se^{|z|^2/2}$ of the model at
the start of \S III.4.

\textbf{Lemma O18 (a theta series defines a map of forms).} \emph{Let $\mathcal G$, $\mathcal G'$ be sets of
lattices in $\mathsf X$, $\mathsf X'$ as in Definition~\textup{O14}, and let
$\rho$, $\rho'$ be unitary representations of $U_\infty$, $U'_\infty$ on $W_\rho$, $W_{\rho'}$. Let
$\mathbf P$ be a polynomial in the matrix entries of the $\lambda_\tau$ whose values are linear maps
$W_\rho\to W_{\rho'}$, such that}
\begin{equation}\label{ob:equiv}
 \mathbf P\big((\gamma'_\tau\lambda_\tau\gamma_\tau^{-1})_\tau\big)=\rho'(\gamma')\,\mathbf P(\lambda)\,\rho(\gamma)^{-1}
 \qquad(\gamma\in U_\infty,\ \gamma'\in U'_\infty).
\end{equation}
\emph{For $L\in\mathcal G$, $L'\in\mathcal G'$ put
$\mathbf K(L,L')=\mathbf P(\partial)\Theta_{L,L'}(0)=\sum_{\lambda\in\mathrm{Hom}(L,L')}\mathbf P(\lambda)e^{-\pi\|\lambda\|^2}$.
Then for $\phi\in\mathcal F(\mathcal G,\rho)$ the sum}
\begin{equation}\label{ob:lift}
 (K\phi)(L')=\sum_{L}\frac1{|\Gamma_L|}\,\mathbf K(L,L')\,\phi(L),
\end{equation}
\emph{over representatives $L$ of the classes in $\mathcal G$, does not depend on the representatives,
and $K\phi\in\mathcal F(\mathcal G',\rho')$. So $K$ is a linear map
$\mathcal F(\mathcal G,\rho)\to\mathcal F(\mathcal G',\rho')$. In
bases of $W_\rho$ and $W_{\rho'}$ each matrix entry of $\mathbf K(L,L')$ is a number
$P(\partial)\Theta_{L,L'}(0)$ as in \eqref{ob:deriv}.}

\emph{Proof.} By Lemma~O17(c) and \eqref{ob:equiv}, exactly as in the proof of (c),
$\mathbf K(\gamma L,\gamma'L')=\rho'(\gamma')\mathbf K(L,L')\rho(\gamma)^{-1}$ for $\gamma\in U(\mathsf X)$,
$\gamma'\in U(\mathsf X')$. Hence $\mathbf K(\gamma L,L')\phi(\gamma L)=\mathbf K(L,L')\phi(L)$, and
$|\Gamma_{\gamma L}|=|\Gamma_L|$; so each term depends only on the class of $L$. And
$(K\phi)(\gamma'L')=\rho'(\gamma')(K\phi)(L')$. \qed

\emph{Example} (over $\Q(i)$, as in Example~(3) after Definition~O14; numbers from \S I.4.1). Let
$L=L'=\Z[i]$, $\rho(g)=\rho'(g)=g^4$ and $\mathbf P(\lambda)=\lambda^4$; then \eqref{ob:equiv} holds. Here
$\Theta(0)=T_0=1.18034\ldots$ and $\mathbf K(L,L')=\sum_\lambda\lambda^4e^{-\pi|\lambda|^2}=T_4=0.14319\ldots$
($T_4/T_0=0.12131\ldots$), so $K\phi(L')=\frac14T_4\,\phi(L)$. In \eqref{ob:maj}, $S=2.01497\ldots$. With
$\mathbf P(\lambda)=\lambda^2$ and $\rho(g)=g^2$ the space of forms is $0$, and indeed
$\sum_\lambda\lambda^2e^{-\pi|\lambda|^2}=0$.

\noindent Lemma~O18 is the elementary part of the statement that the maps of the chain are well defined.
The lift of Part~II, \S6.1 is conjugate-linear in $\phi$; \eqref{ob:lift} is linear, in the convention of
Part~II, \S6.2. \S I.4.3 describes the entries of the polynomials $\mathbf P$ of the construction as
products of minors of the matrices $\lambda_\tau$, of degree about $a_t$ at a top place $(0,t)$; this is
not proved or used here.

\textbf{Hypothesis O19 (the maps of the chain; assumed here, from the Lean development).}
\emph{For each word $\mathsf w$ and each $0\le j<m$ there are, for every $L\in\mathrm{Cl}_j(\mathsf w)$ and
$L'\in\mathrm{Cl}_{j+1}(\mathsf w)$, an entire function $F_{L,L'}$ of the variable $z$ of
Definition~\textup{O16} $(fN_jN_{j+1}$ complex variables$)$ and a number $\epsilon(L,L')$ of absolute
value $1$, and there is a polynomial $\mathbf P_j$ satisfying \eqref{ob:equiv} for the representations of
the weights of \textup{O24(b)} at the stages $j$ and $j+1$, such that}
\begin{equation}\label{D:376}
 \mathbf K_j(L,L')=\mathbf P_j(\partial)F_{L,L'}(0),
\end{equation}
\emph{and such that \eqref{ob:chain} defines a linear map $K_j:V_j\to V_{j+1}$.}

\noindent \emph{Remark (the relation to \eqref{ob:theta}; not part of the Hypothesis).} In the classical
picture $F_{L,L'}$ is the theta series of \eqref{ob:theta} after a normalisation by periods
$\Omega_\tau>0$, namely
$F_{L,L'}(z)=\Theta_{L,L'}\big((\Omega_\tau^{-1/2}z_\tau)_\tau\big)/\Theta_{L,L'}(0)$, and then $K_j$ is a map \eqref{ob:lift} up to the unit scalars and a normalisation; its
equivariance is Lemmas~O17(c) and~O18. The source does not define its maps by the formula
\eqref{ob:theta}. It specifies them by a record of abstract data (\S9, rows 18 and 20), with statements
that tie these data to theta functions on abelian varieties with complex multiplication by $E$, of which
\eqref{ob:theta} is the classical expression. We did not verify this tie. Also, the source normalises by
linear maps applied before and after the kernel, not by one constant. Three properties of the
kernels are used later; O20 and O21 state the first two.

\textbf{Hypothesis O20 (algebraicity and integrality; assumed here, from the Lean development).}
\emph{Let $1\le j\le m$, $L\in\mathrm{Cl}_{j-1}(\mathsf w)$, $L'\in\mathrm{Cl}_j(\mathsf w)$, and let
$F=F_{L,L'}$ be the function of \textup{O19} $($for the index $j-1$ there$)$.}
\begin{enumerate}[label=\textup{(\alph*)},leftmargin=2.2em,itemsep=1pt,topsep=2pt]
\item \emph{There is a number field $L_0\supset E$, Galois over $\Q$ and without real embeddings, the same
for all $j,L,L'$ and all words, such that after a linear change of the variables $z_\tau$ with
coefficients in $L_0$ all Taylor coefficients of $F$ at $0$ lie in $L_0$.}
\item \emph{Let $\ell\ge5$ be a prime with $\chi_d(\ell)=-1$ and let $\sigma_\ell$ have order $d'$ in $G$.
Suppose $c$ is not a power of $\sigma_\ell$. Fix a place $\tau_1\in\Phi$, and an embedding of $L_0$ into an
algebraic closure of $\Q_\ell$; let $\mathrm{ord}_\ell$ be the valuation there with
$\mathrm{ord}_\ell(\ell)=1$. Call the place $\tau_1\circ(c\,\sigma_\ell)^{\varsigma}$ the place of the
\emph{site} $\varsigma\in\Z/d'$. Restrict the series of \textup{(a)} to the variables at the places of
the $d'$ sites, the others being set to $0$, and write these variables as $Z_{\varsigma,i,i'}$ with
$i<N_{j-1}$ and $i'<N_j$. Then there are bases of $L$ and $L'$ at the primes above $\ell$ and a number
$\epsilon_\ell$ with $\mathrm{ord}_\ell(\epsilon_\ell)=0$ such that, after the change to these bases
and the multiplication of the variables by $\epsilon_\ell$, the restricted series is a function
$\Theta$ of the form \eqref{dv:theta} of \S III.5, with $d=d'$, $N_0=N_{j-1}$, $N_1=N_j$: a power series
$U$ whose coefficients have $\mathrm{ord}_\ell\ge0$, composed with the map $g$ of \S III.5, times the
exponential of a sum of values of a series $\beta$ that satisfies hypothesis \textup{(B)} of
\S III.5, \S3.}
\end{enumerate}

\noindent Part~(a) is not a transcription of one statement of the source; the field $L_0$ and its two
properties are in the record. Part~(b) is the standing assumption of \S III.5, with its variables made
explicit: the $N_{j-1}\times N_j$ matrix variables at the $d'$ places of one orbit of
$\tau\mapsto\tau\circ c\circ\sigma_\ell$. To read it one needs \S III.5, \S1--3, where $g$, $\beta$ and (B)
are defined; \S III.5 writes $s$ for the site, $Z_{s,q,c}$ for the variables and $v$ for the valuation.
The order $d'$ is even, because $\sigma_\ell^{d'}=1$ gives $\chi_d(\ell)^{d'}=1$. The $d'$ places are
distinct: if $(c\sigma_\ell)^i=1$ with $i$ odd, then $c=\sigma_\ell^{-i}$, which is excluded; so $i$ is
even, and then $\sigma_\ell^i=1$ and $d'\mid i$. By Remark~(2) after Lemma~O3, $c$ is a power of
$\sigma_\ell$ only when $\sigma_\ell=c$; then the orbit has one element, and the source treats these primes
by a separate statement with one site, which \S III.5 does not cover. We did not locate in the source
the passage from its statement about $F$ to hypothesis~(B).

\textbf{Hypothesis O21 (Gaussian size and periods; assumed here, from the Lean development).}
\emph{Let $j$, $L$, $L'$ and $F=F_{L,L'}$ be as in \textup{O19}.}
\begin{enumerate}[label=\textup{(\alph*)},leftmargin=2.2em,itemsep=1pt,topsep=2pt]
\item \emph{For every embedding of $L_0$ in $\C$ the function $F$, with its Taylor coefficients embedded,
satisfies $|F(\sqrt\pi\,z)|\le S\,e^{\|z\|^2/2}$ for all $z$, with a constant $S$ of the form
\eqref{III:Sform} of \S III.4, in which $N=N_j$, $N'=N_{j+1}$, $g=2nN_jN_{j+1}$, and $\beta>0$, the modulus $c\ge1$ and
$\Theta_0>0$ are the numbers so named there.}
\item \emph{There are numbers $\Omega_\tau(v)>0$, for $\tau\in\Phi$ and $v$ an embedding of $L_0$ in $\C$
up to complex conjugation, which are the periods in \textup{(S1)} of \S III.7, and for every $\tau$ the
mean of $\log\Omega_\tau(v)$ over $v$ is $\log\Omega_E$, with $\Omega_E$ defined by Part~II, \textup{\eqref{II:t6.3}}.
This is input \textup{(S4)} of \S III.7.}
\end{enumerate}

\noindent By Lemma~O17(a) the series \eqref{ob:theta} is an entire function with a Gaussian majorant;
this is the setting of the block estimate of \S III.4 for one pair $(L,L')$. O21 is what \S III.4 needs
beyond it: item (vi) of the list at the end of \S III.4, and (S4). In (b) the letter $v$ is the place of
\S III.4 and \S III.7, not a row. In the record $\log\Omega_E$ is defined by the formula of \eqref{II:t6.3}, and
the statement about the mean occurs in the source as a hypothesis of two theorems (\S9, row~23). O21 is
not proved here.

\noindent \emph{Equivariance}, the third property: Lemma~O17(c) and Lemma~O18, proved here for kernels
of the form \eqref{ob:theta}. For the kernels of the construction it is contained in O19, in the
assertion that \eqref{ob:chain} maps $V_j$ to $V_{j+1}$.

\subsection*{6.\quad The family, the pairing and what is averaged}

Fix $x$ and a translate $u\in H$. The data of \S3--5 for this $u$ form what \S III.6 and \S III.8 call a
\emph{record}. The record contains a finite family of chains, indexed by words. Everything in this
subsection that is not a definition is collected in one Hypothesis.

\textbf{Hypothesis O25 (the family and the pairing; assumed here, from the Lean development).}
\emph{Fix $x$ and $u$, and the data of \textup{O23}, \textup{O12}, \textup{O24}, \textup{O19},
\textup{O20(a)}.}
\begin{enumerate}[label=\textup{(\alph*)},leftmargin=2.2em,itemsep=1pt,topsep=2pt]
\item \emph{Words. A word $\mathsf w\in\mathsf W_u$ has $k-2$ letters in $\{0,1\}$, and different words
have different letters. At every place, consecutive $\psi_l$ differ by $4\Delta$, the numbers $\xi_0$ and
$\xi_{k-1}$ do not depend on the word, and for $0\le l\le k-3$}
\begin{equation}\label{D:377}
 \xi_{l+1}=\psi_l-\tfrac12-\Delta\cdot(\text{the }l\text{th letter of }\mathsf w).
\end{equation}
\emph{Each word has its own level prime $\mathrm{Lv}(\mathsf w)$, lattices $\mathrm{Cl}_j(\mathsf w)$ and
chain \eqref{ob:chain}. The weights $w_{u,j}$, the sets $T$, $T_2$ and the orders $a_t$ do not depend on
the word.}
\item \emph{The source space. There is a subspace $\mathcal C_{\mathsf w}\subset V_0(\mathsf w)$, and for
each of the $f$ primes $\mathfrak p$ of $F$ above $\mathrm{Lv}(\mathsf w)$, with residue field
$\mathbb F_{\mathfrak p}$ of $\mathrm{Lv}(\mathsf w)$ elements, a linear action of the group
$\mathbb F_{\mathfrak p}^\times$ on $\mathcal C_{\mathsf w}$, written $\alpha\cdot\phi$. Put}
\begin{equation}\label{D:378}
 \mathcal M_{\mathsf w}=\{\phi\in\mathcal C_{\mathsf w}:\ \textstyle\sum_{\alpha\in\mathbb F_{\mathfrak p}^\times}\alpha\cdot\phi=0
 \ \text{for all }\mathfrak p\mid\mathrm{Lv}(\mathsf w)\}.
\end{equation}
\emph{Then
$\dim\mathcal M_{\mathsf w}=\big(1-\frac1{\mathrm{Lv}(\mathsf w)-1}\big)^f\dim\mathcal C_{\mathsf w}$.}
\item \emph{The period. Let $\mathcal A$ be the space of forms on the classes of a genus of lattices in
$\mathsf X'$ $(\cong\mathsf Y_m)$, of a weight determined by $\psi,\theta,w$. There is a linear map
$\mathrm{res}:V_m(\mathsf w)\to\mathcal A$, and the \emph{period} is the bilinear form}
\begin{equation}\label{ob:per}
 \mathrm{Per}_{\mathsf w}(\phi,\phi')=\Big\langle\ \mathrm{res}\big(K_{m-1}\cdots K_0\,\phi\big),\ \phi'\ \Big\rangle
 \qquad(\phi\in\mathcal M_{\mathsf w},\ \phi'\in\mathcal A),
\end{equation}
\emph{the pairing being the finite sum of Definition~\textup{O14}, made bilinear as in \eqref{ob:lift}.}
\item \emph{The target. There are a complex inner product space $\mathcal E_0$ of finite dimension
$D_{\rm r}$, the same for all words, a linear map $\iota:\mathcal A\to\mathcal E_0$, and a subspace
$\mathcal A_0\subset\mathcal A$. Let $\mathcal A_1$ be the orthogonal complement of
$\mathcal A_0\cap\ker\iota$ in $\mathcal A_0$, so that $\iota$ is injective on $\mathcal A_1$, and for
$y\in\mathcal E_0$ let $\mathrm{sec}(y)$ be the element of $\mathcal A_1$ whose image under $\iota$ is the
orthogonal projection of $y$ to $\iota(\mathcal A_1)$.}
\item \emph{Rationality. Define the pairing of the word by}
\begin{equation}\label{ob:B}
 \mathfrak B_{\mathsf w}:\mathcal M_{\mathsf w}\times\mathcal E_0\to\C,\qquad
 \mathfrak B_{\mathsf w}(\phi,y)=\mathrm{Per}_{\mathsf w}\big(\phi,\mathrm{sec}(y)\big).
\end{equation}
\emph{Let $L_0$ be the field of \textup{O20(a)}, with a fixed embedding in $\C$. There are
$L_0$-vector spaces $U_{\mathsf w}\subset\mathcal M_{\mathsf w}$ and $U'\subset\mathcal E_0$, each spanned
over $L_0$ by a basis of the complex space, such that $\mathfrak B_{\mathsf w}$ takes values in $L_0$ on
$U_{\mathsf w}\times U'$.}
\item \emph{Lattices. $U'$ contains a lattice $\mathcal R$ over $\mathfrak o_{L_0}$, the same for all
words, and $U_{\mathsf w}$ contains lattices
$\Lambda^{\rm div}_{\mathsf w}\subset\Lambda^{\rm src}_{\mathsf w}$. For every embedding $v$ of $L_0$ in $\C$, the complex vector spaces with the same bases as $U_{\mathsf w}$
and $U'$, the coefficients being embedded by $v$, carry Hermitian inner products.}
\end{enumerate}

\noindent \emph{Comments.} (a) So \emph{the words are the weights of rank $k$ that interlace one fixed
weight $\psi$ of rank $k-1$}: $\xi_{l+1}$ lies between $\psi_{l+1}$ and $\psi_l$ (\S III.8, \S5). The prime
sum $P_{u,t,j}$ of Part~II, \eqref{II:t4.1} depends only on $w_{u,j}$, $T$, $T_2$, $a_t$, hence not on the word. A
word is the \emph{member} $\alpha$ of \S III.6. (b) These are the spaces $\mathcal C_{\mathsf w}$,
$\mathcal M_{\mathsf w}$ of \S III.8. The action rescales the linear form of the level structure at
$\mathfrak p$; this description is read from the types of the source. The dimension formula is a
statement of the source on which the final theorem depends (\S9, row~24). By a statement that the source
uses as a hypothesis, $\mathcal C_{\mathsf w}$ is all of $V_0(\mathsf w)$. (c) \eqref{ob:per} is the
``single period integral'' of Part~II, \S6.2. Our reading of $\mathrm{res}$, which is not in the source:
a lattice $L^\circ$ of $\mathsf X'$ gives the lattice $L^\circ\oplus\mathfrak o_Ey_0$ of $\mathsf X_m$, the
form is evaluated there, and the value is projected to the weight of $\mathcal A$. The source gives the
period as a composite of three linear maps. The variants of $\mathcal A$ with the twists of Part~II,
\S6.8 are not needed here. (d) $\mathcal E_0$ is a space of forms on the unitary group of $\mathsf Y_0$,
which is not compact; we give no classical description of it or of $\iota$. (e) The restriction of
$\mathfrak B_{\mathsf w}$ to $U_{\mathsf w}\times U'$ is the pairing $\langle\ ,\ \rangle_\alpha$ of \S III.6
and $\mathfrak B^{\rm alg}_{\mathsf w}$ of \S III.8. That these two are the same pairing was not checked in the
source (\S III.8, \S2); the derivation of the mean-slope inequality in \S\ref{as:sec} (Hypothesis (H8) there)
assumes it. (f) $\Lambda^{\rm src}_{\mathsf w}$ is the \emph{source
lattice} of \S III.6 (the preimage of the first lattice of the chain) and $\Lambda^{\rm div}_{\mathsf w}$
its sublattice defined there; the \emph{glued lattice}
$\Lambda^{\rm gl}_{\mathsf w}=I_{\mathsf w}\Lambda^{\rm div}_{\mathsf w}$ is defined in \eqref{lt:N}.

\textbf{Slopes over $L_0$.} For a lattice in $U_{\mathsf w}$ or $U'$ the slope $\mu$ is defined by
\eqref{lt:mu}, as in Definition~O8 with $L_0$ in place of $E$: the covolume is taken for the Euclidean
norm that is the sum, over the embeddings of $L_0$ up to conjugation, of the Hermitian norms of
\textup{O25(f)}. We did not identify these norms with those of Definition~O14. The \emph{target lattice}
of the pairing is
\begin{equation}\label{D:381}
 \mathcal R^\vee=\{y^\vee:U'\to L_0\ \text{$L_0$-linear}:\ y^\vee(\mathcal R)\subset\mathfrak o_{L_0}\},
\end{equation}
a lattice in the space of the $L_0$-linear maps $U'\to L_0$, with the dual norms. The divisibility
statement of \S III.6 is that $\phi\mapsto\mathfrak B_{\mathsf w}(\phi,\cdot)$ maps
$\Lambda^{\rm gl}_{\mathsf w}$ into $\mathcal R^\vee$. The slope $\mu(\mathcal R^\vee)$ is the number
$\mu_{\rm CC}$ of Part~II, \S6.7.

\textbf{What is averaged.} For a word let $d_{\mathsf w}=\mathrm{rk}\,\Lambda^{\rm gl}_{\mathsf w}$ and let
$h_{\mathsf w}$ be the bound for the height of $\mathfrak B_{\mathsf w}$ (the $h$ of Part~II, \S6.7; with
its index it is not the Hermitian form). The
\emph{inner mean}, for one translate $u$, is the mean over the words weighted by the ranks:
\begin{equation}\label{ob:inner}
 s_u=\frac{\sum_{\mathsf w\in\mathsf W_u}d_{\mathsf w}\,\big(\mu(\Lambda^{\rm gl}_{\mathsf w})-h_{\mathsf w}\big)}{\sum_{\mathsf w\in\mathsf W_u}d_{\mathsf w}} .
\end{equation}
The slope inequality for many maps (Part~II, \S6.9--6.10; \S III.8) compares $s_u$ with
$\mu(\mathcal R^\vee)$. Given the Hypotheses of \S\ref{as:sec} it gives, for each $u$, an upper bound for
$\sum_{t\in T}\sum_{j=1}^m(P_{u,t,j}-M_{u,t,j})$ (Proposition~M7 there). The \emph{outer mean} is the mean $\frac1n\sum_{u\in H}$ of
these bounds over the $n$ translates; it is the left side of (M). Only the outer mean removes the
dependence on the residue class modulo $N$ (\S I.6(c)). The source sums the bounds over the
translates of a subgroup, which for $\zeta$ is all of $H$ (\S\ref{as:sec}, \S2, item (5) and \S9, Remark~(5)).

\subsection*{7.\quad The dictionary}

Each row gives an abstract object of \S III.4--III.9, in the notation of that section, the object of the
construction, a status, and the rows of the table of \S9 where the declarations of the source are named.
\textbf{V}: verified in the source, in the way stated. \textbf{T}: read from the types and statements of
the record only; no proof was followed. \textbf{N}: not verified. A record of the source is a list of
data and of properties that the data are required to have; ``V'' for a field of the record means that
the property is part of the record, not that the record used by the final theorem was inspected. The
level is called $\mathcal T$ in Part~II, \S6.7 and \S III.7--III.8, and $\mathrm{Lv}_\alpha$ in \S III.6;
in \S III.5--III.6 the letter $\mathcal T$ is the set of labels, which is $T$ here.

{\small\setlength{\LTleft}{0pt}\setlength{\LTright}{0pt}
\begin{longtable}{@{}p{.22\textwidth}p{.39\textwidth}p{.25\textwidth}p{.06\textwidth}@{}}
\toprule
Abstract object & Object of the construction & Status & \S9\\
\midrule\endhead
\bottomrule\endfoot
\multicolumn{4}{@{}l}{\emph{\S III.4 (size of the kernels)}}\\
field $K$, places $v$ & $L_0$ and its embeddings up to conjugation & T & 21\\
$V_j$, norms $\|\cdot\|_{j,v}$ & forms on $\mathrm{Cl}_j(\mathsf w)$ of the weight of O24(b), with the norm of Definition~O14 & V: fields of the record & 18\\
classes $p\in I_j$, $w_j(p)$ & $L\in\mathrm{Cl}_j$, $1/|\Gamma_L|$ & V for the formula $w_j=1/|\Gamma_L|$; that $\Gamma_L$ is a group of automorphisms: stage $0$ by O12(vii), stages $\ge1$ N & 16, 18\\
$K_j$; $\epsilon_{pq}$, $T_{pq}$ & the lift \eqref{ob:chain}; $\epsilon(L,L')$, $\mathbf K_j(L,L')$ & V: the definition has this shape; $T_{pq}$ is a composite of three maps; $|\epsilon|=1$ for $j<m$ & 18\\
$T_{pq}$ is a block $P(\partial)F$ with Gaussian $F$ & Hypotheses O19, O21 & N & 20, 23\\
$P=K_{m-1}\cdots K_0$ & the composite in \eqref{ob:per} & V: definition & 18\\
$\Omega_\tau$; $\Omega_E$ & periods of O21(b); Part~II, \eqref{II:t6.3} & N & 23\\
distinguished places $\mathfrak t$, $a_{\mathfrak t}$, $b_{\mathfrak t}$ & top places $(0,t)$, $a_t$, $b_t$, by \eqref{ob:top} & V: field of the record & 7\\
level $\mathcal T$ & $\mathrm{Lv}(\mathsf w)$ & V: see (b) & 24\\
\multicolumn{4}{@{}l}{\emph{\S III.5 (divisibility at inert primes)}}\\
$\ell$; $K$, $v$ & a prime $\ell\ge5$ with $\chi_d(\ell)=-1$; an algebraic closure of $\Q_\ell$, $\mathrm{ord}_\ell$ & V: field & 22\\
$d$, the sites $\Z/d$ & the order $d'$ of $\sigma_\ell$; the site $\varsigma$ is the place $\tau_1(c\sigma_\ell)^\varsigma$ & N: stated in the source only in a declaration that the final theorem does not depend on (see (d)) & 22\\
$N_0$, $N_1$ & $N_{j-1}$, $N_j$ & V: field & 22\\
$Z_{s,q,c}$; the frame $\xi$ & $Z_{\varsigma,i,i'}$: matrix variables at the places of one orbit; change of basis to bases of $L$, $L'$ at $\ell$ & T & 22\\
$\Theta$, hypothesis (B) & O20(b) & N & 22\\
top sites $\mathit{TT}$ & sites whose place is a top place $(0,t)$, $t\in T$ & N & 22\\
shape, labels, Theorem~J's frame & column lengths of the weights; not described here & N & --\\
cell of a site & Lemma~O5: \eqref{ob:partner} & V: proved here; the definition in the source read & 5\\
\multicolumn{4}{@{}l}{\emph{\S III.6 (from derivatives to lattices)}}\\
$L_0$, a place $w$ & coefficient field, its primes & V: field & 21\\
translate $u$ & $u\in H$; thick cells $(1,u[5]^i)$ & V: definition & 4\\
member $\alpha$ & word $\mathsf w$ & V: fields & 24\\
$U_\alpha$, $\Lambda^{\rm src}_\alpha$, $\Lambda^{\rm div}_\alpha$, $\Lambda^{\rm gl}_\alpha$ & $U_{\mathsf w}$ and its lattices & T & 26\\
$U'$, $\mathcal R$ & $U'\subset\mathcal E_0$, its lattice & V: types & 26\\
$\langle\ ,\ \rangle_\alpha$ & $\mathfrak B_{\mathsf w}$ on $U_{\mathsf w}\times U'$, \eqref{ob:B} & V, with one hypothesis not discharged: see (a) & 25, 33\\
$\mathrm{Lv}_\alpha$ & $\mathrm{Lv}(\mathsf w)$ & V: see (b) & 24\\
$\mathrm{rk}^*_\alpha$ & $\dim\mathcal C_{\mathsf w}$ & V: definition & 24\\
$\mathrm{wt}_j$ on cells & $w_{u,j}(v,s)=w_j(\tau)$ & V: definition & 7\\
classes of primes & Remark~(1) after Lemma~O3 & V: definition & 4\\
gain, $\phi$, err, $\tau(w)$, io & numbers attached to places of $L_0$; not described here & T & 29\\
\multicolumn{4}{@{}l}{\emph{\S III.7 (arithmetic inputs of the size bound)}}\\
$K$; $\mathcal P$; $T\subseteq\mathcal P$ & $L_0$; $\Phi$; the top places & T & 4, 21\\
$V_0,\dots,V_m$, $K_j$ & the chain \eqref{ob:chain}, renumbered & T & 18\\
$U$, $U_0$, $\bar U$ & $\mathcal A$; $\mathcal A_0\cap\ker\iota$; the subspace $\mathcal A_1$ of O25(d) & T & 26\\
$Y$, $\iota$ & $\mathcal E_0$; $\iota$ of O25(d) & T & 26\\
$\Pi$; the map $\Phi$ with \eqref{sa:fact} & $\mathrm{Per}_{\mathsf w}$; $\mathfrak B_{\mathsf w}$ (not the CM type); \eqref{sa:fact} is \eqref{ob:B} & T: see (a) & 25\\
$\mathcal R$, $\mu_{\rm CC}$ & $\mathcal R$, $\mu(\mathcal R^\vee)$ & T: the equality is a hypothesis in the source & 26\\
$\mathcal R_0$, $\mu_0$ & a lattice at lower level; not described here & N & 26\\
$\nu_{j\tau i}$, $w_{j\tau}$, $H_{j\tau}$, $p_j$ & degrees and hook products of the weights of O24(b) & N & 30\\
$\mathcal T$; $\Omega_\tau(v)$; $v_{\rm h}$ & $\mathrm{Lv}(\mathsf w)$; periods; $\mathrm{rd}_E/2$ & V; N; V & 4, 23, 24\\
\multicolumn{4}{@{}l}{\emph{\S III.8 (trace-formula estimates)}}\\
$\mathsf W$, $\mathsf w$ & $\mathsf W_u$ & V: fields & 24\\
$\mathcal C_{\mathsf w}$, $\mathcal M_{\mathsf w}$ & O25(b) & V: definition & 24\\
$\mathcal H$, $P_{\mathsf w}$, $c_0(\mathsf w)$ & a space of forms containing a copy of $\mathcal E_0$; not described here & N & 31\\
$\mathcal E_0$, $\mathsf s$, $D_{\rm r}$ & $\mathcal E_0$ of O25(d) & V: types & 26\\
$\mathfrak B_{\mathsf w}$; $\mathfrak B^{\rm alg}_{\mathsf w}$ & \eqref{ob:B}; its restriction to $U_{\mathsf w}\times U'$ & V: see (a) & 25\\
$\mathbb L$ & $L_0$ & V: types & 21\\
$\mathcal T(\mathsf w)$ & $\mathrm{Lv}(\mathsf w)$ & V: see (b) & 24, 28\\
$\xi,\psi,\theta,e'$ & \eqref{ob:inter}, $e'_j=w_{j+1}$ & V: definition & 6\\
\multicolumn{4}{@{}l}{\emph{\S III.9 (Galois cohomology)}}\\
$E$; twists $1\le e<t$ & $E$; $t=t_0$ & V: definition & 32\\
$h_{\ell,e}$ & the record's number, equal to the length when it is finite & V: field & 32\\
$K$; twisted $w$ & $L_0$; a place of $L_0$ with a nontrivial twist & V: types & 32\\
$\eta_j(w)$ & a number of the record, a sum of local terms; see (c) & V: field; meaning T & 32\\
\end{longtable}}

\noindent\textbf{Points that earlier sections left open.}

(a) \emph{The pairing.} \S III.6 names the pairing $\langle\ ,\ \rangle_\alpha$ by the period of the chain,
and \S III.8 names $\mathfrak B^{\rm alg}_{\mathsf w}$ by a bilinear form with values in $L_0$; they left open
whether these are the same. In the source there are three objects. The period
$\mathrm{Per}_{\mathsf w}(\phi,\phi')$ of \eqref{ob:per} is defined for $\phi'\in\mathcal A$. The form
$\mathfrak B_{\mathsf w}$ on $\mathcal M_{\mathsf w}\times\mathcal E_0$ is \emph{defined} from it by
\eqref{ob:B}; this is a property required of the family of records. And the $L_0$-valued form is tied to
$\mathfrak B_{\mathsf w}$ by the statement that $\mathfrak B_{\mathsf w}$ takes the same values on
$U_{\mathsf w}$ and $U'$; this statement is a clause of the hypothesis (P) of \S III.8. So the two
pairings agree, given \eqref{ob:B} and (P). In the statement that \S III.6 uses for the divisibility, the
equality of the $L_0$-valued form with the bilinear form of that statement is one of its hypotheses; we
did not find where it is discharged. Also, \eqref{ob:B} is the equation \eqref{sa:fact} of \S III.7:
$\Pi\circ K=\Phi\circ\iota$ on the complement of the kernel of $\iota$.

(b) \emph{The level.} The level $\mathcal T$ of Part~II, \S6.7, $\mathrm{Lv}_\alpha$ of \S III.6 and
$\mathcal T(\mathsf w)$ of \S III.8 are the same number, the prime $\mathrm{Lv}(\mathsf w)$ of the word. This
was checked in the statements: the lattice statement and the slope statement behind \S III.6, the main
term of the size bound, and the table of data for each translate used by \S III.8 all refer to the same
field of the record.

(c) \emph{The number $\eta_j(w)$ of \S III.9.} In the record it is a sum of local terms, each a quotient
of two lengths of modules over a ring of $\ell$-adic integers; it vanishes at the places without twist.
The description belongs to \S III.9 and is not repeated here. It does not settle the question of
\S III.9, \S2, whether this number equals the number defined there for the family of twists; the
inequality \eqref{gc:link} is stated in the source for the field of the record, so the assembly does not
need the answer. (This $\eta_j(w)$ is not the width $\eta$ of \S3.)

(d) \emph{Sites and cells for Theorem~J.} \S III.5 and \S III.6 assumed the identification of the sites
of an orbit with places. The source states it: for each prime $\ell$, lift $j$, word and pair of classes
to which the local estimate is applied, there are a Frobenius element $\sigma_\ell$, of order $d'$, and a
set of base places such that (base place, site $\varsigma$) $\mapsto$ the place
$\tau_1(c\sigma_\ell)^\varsigma$ is a bijection onto $\Phi$; and the power series to which Theorem~J is
applied is the expansion of the kernel in the variables at these places, in bases of the two lattices at
$\ell$, up to a unit constant. By Lemma~O5 the site $\varsigma$ of a top place $(0,t)$ leads, for odd $i$,
to the cell $\mathrm{cell}(\ell,t,i)$ of \eqref{lt:cell}. The declaration of the source that states this (\S9, row 22) is not one that the final theorem depends on, by the
dependency computation; in the record itself $d'$ is only required to be even and $\ge2$. The frames,
labels and shapes of Theorem~J remain unverified. In the source the primes with one site are those of
inert class whose Frobenius has a component of odd order in $H$; as $H$ is a $2$-group this means
$[\ell]=[1]$, that is $\sigma_\ell=c$, in agreement with O20(b).

\emph{What remains open in the dictionary:} every row marked N; the rows of \S III.7 as a whole (we
matched types, not the eleven inputs); the space $\mathcal H$ and the maps $P_{\mathsf w}$ of \S III.8;
and the question in (c).

\subsection*{8.\quad A toy instance}

\begin{illbox}
This subsection computes the objects for a small $N$ and the primes of the running example. Its parameters
do not satisfy the conditions of the construction, and it is not used in the proof.
\end{illbox}

Take $N=16$, $q=11$, $d=1319$, as in \S I.5; so $e=2$, $n=4$, $f=8$. This is the smallest $N$ allowed, and
the primes violate conditions of O23 (Remark~(5) of \S3). The computations are in the
scripts \texttt{o1\_field.gp}, \texttt{o1b\_table.gp}, \texttt{o2\_cells.py},
\texttt{o4\_ranks.py}; those for the example of the first version, $q=3$, $d=23$, carry the suffix
\texttt{\_q3\_d23}.

\emph{The field.} $E=\Q(\zeta_{16}+\zeta_{16}^{-1},\sqrt{-11},\sqrt{-1319})$ has degree $16$ and no real
embedding. Its discriminant is $2^{44}11^8\,1319^8$, as \eqref{ob:rd} predicts ($4((e+1)n-1)=44$), and
$\mathrm{rd}_E=2^{11/4}\sqrt{14509}=810.3100\ldots$, $\mathrm{rd}_E/2=405.1550\ldots$. We did not compute its
class group. For the field of the first version, with $q=3$, $d=23$, the class group is cyclic of order
$60$ (PARI, certified by \texttt{bnfcertify}); so lattices over $\mathfrak o_E$ need not be free. The group
$H$ is $\{[1],[5],[7],[3]\}=\{[5]^0,[5]^1,[5]^2,[5]^3\}$.

\emph{Splitting.} For the $76$ primes $3\le p\le400$, $p\ne11$, the factorisation of $p\,\mathfrak o_E$ computed
by PARI agrees with \eqref{ob:split}. The same holds for $(N,q,d)=(16,3,23)$ ($75$ primes), $(32,11,7)$
($27$ primes) and $(16,11,23)$ ($43$ primes). As a control, replacing the order of $[p]$ in
$H$ by the order of $p$ in $(\Z/16)^\times$ gives five wrong predictions among the $76$. The first
primes:
\begin{center}\small
\begin{tabular}{r|ccccccccccccccc}
\toprule
$p$ & $3$&$5$&$7$&$13$&$17$&$19$&$23$&$29$&$31$&$37$&$41$&$43$&$47$&$53$&$59$\\
\midrule
$[p]$ & $[3]$&$[5]$&$[7]$&$[3]$&$[1]$&$[3]$&$[7]$&$[3]$&$[1]$&$[5]$&$[7]$&$[5]$&$[1]$&$[5]$&$[5]$\\
$\chi_{11}(p)$ & $+$&$+$&$-$&$-$&$-$&$-$&$+$&$-$&$+$&$+$&$-$&$-$&$+$&$+$&$+$\\
$\chi_{1319}(p)$ & $+$&$+$&$+$&$-$&$-$&$+$&$-$&$-$&$+$&$-$&$-$&$+$&$-$&$+$&$-$\\
$f_p$ & $4$&$4$&$2$&$4$&$2$&$4$&$2$&$4$&$1$&$4$&$2$&$4$&$2$&$4$&$4$\\
$g_p$ & $4$&$4$&$8$&$4$&$8$&$4$&$8$&$4$&$16$&$4$&$8$&$4$&$8$&$4$&$4$\\
class & sp&sp&sp&i&i&sp&s&i&sp&s&i&sp&s&sp&s\\
\bottomrule
\end{tabular}
\end{center}
\noindent (s: signal; i: inert class; sp: split.) The prime $31$ splits completely. The prime $17$ has
$\sigma_{17}=c$: the primes of $F$ above $17$ stay prime in $E$ (Remark~(2) after Lemma~O3; the Remark
was tested on $2400$ primes in $8$ fields), and the orbit of a place under
$\tau\mapsto\tau c\sigma_{17}$ has one element. Among the primes below $10^5$ the three classes have
proportions $0.502$, $0.250$, $0.248$, and the eight cells between $0.061$ and $0.064$ each. The
ramified primes are in Remark~(3) after Lemma~O3.

\emph{The table.} By Lemma~O5 with $t=[1]$, a prime $p\ne1319$ with $\chi_{1319}(p)=-1$ goes to row
$\frac12(1+\chi_{11}(p))$ and column $[p]$. By Lemma~O22 the prime $1319$ goes to row
$\frac12(1+\chi_{11}(1319))=0$ and column $[1319]=[7]$, and $1319^2$ to row $1$ and column $[1]$. This
reproduces the table of \S I.5: row $0$ contains $17\,|\,13,29,61\,|\,-\,|\,41$ and row $1$ contains
$47\,|\,-\,|\,37,59\,|\,23$ in the columns $[1],[3],[5],[7]$. The rule \eqref{ob:partner} was compared with
the rule of \S I.5 for all four labels $t$, all primes below $4000$ with $\chi_{1319}(p)=-1$ and exponents
$1,3,5$: $3288$ agreements, no disagreement; without the conjugation $c$ the image is never a place. Two
orbits: for $p=23$ (signal), $c\sigma_{23}=([7],-1,+1)$ has order $2$, and the orbit of the top place
$(0,[1])$ is $\{(0,[1]),(1,[7])\}$, two sites. For $p=13$ (inert class), $c\sigma_{13}=([3],+1,+1)$ has
order $4$, and the orbit is $(0,[1]),(0,[3]),(0,[7]),(0,[5])$: four sites, all in row $0$.

\emph{Ranks.} With $n=4$ the smallest triple allowed by the rank conditions is $(m,k,t_0)=(2,8,15)$ (see Remark~(4) of \S3). Then $r=9$,
$\varrho=24$, the definite ranks are $N_0,N_1,N_2=8,24,40$, the spaces $\mathsf Y_0,\mathsf Y_1,\mathsf Y_2$
have ranks $9,24,39$ and signatures $(8,1),(24,0),(39,0)$, and $\mathsf Z_1,\mathsf Z_2$ have signatures
$(8,1),(9,0)$. A word has $6$ letters and there are at most $2^5=32$ words. The kernel of the first lift
is a function of $f\,N_0N_1=8\cdot8\cdot24=1536$ complex variables.

\subsection*{Letters of this section}

The letters fixed here agree with Part~II, \S4 and \S6 and with \S III.6, except as stated in the last
column. A letter listed with two meanings keeps both; the index or the subsection tells them apart.

{\small\setlength{\LTleft}{0pt}\setlength{\LTright}{0pt}
\begin{longtable}{@{}p{.25\textwidth}p{.37\textwidth}p{.31\textwidth}@{}}
\toprule
Letters & Meaning & Elsewhere; second meanings\\
\midrule\endhead
\bottomrule\endfoot
$x$, $\varepsilon$, $b$, $a$, $X'$, $X''$, $b'$, $w''$, $\delta_1$, $\eta$ & the scale and the cut-offs (\S3) & $\eta_j(w)$ of \S III.9 is another number; $a$ is also a residue class in \S1\\
$E$, $F$, $k_d$, $E'$, $\mathfrak o_E$, $N$, $n=2^e$, $f=2n$, $Q=Nqd$, $q$, $d$, $\chi_q$, $\chi_d$ & the field (\S1) & $e_i$, $f_i$, $f_p$, $g_p$, $(e,f,g)$: ramification data; $d'$: order of $\sigma_\ell$ (the $d$ of \S III.5)\\
$G$, $H$, $c$, $\sigma_p$, $\sigma_{d,i}$, $I_d$, $\varepsilon_q,\varepsilon_d$, $\Phi$, $\tau_{\rm b}$, $\mathrm{rd}_E$ & Galois group, conjugation, Frobenius, CM type, base embedding & $\tau_0$ in Part~II, \S6.2 is our $t_0$; $v=\mathrm{rd}_E/2$ in Part~II, \eqref{II:t6.3}; $\Phi$ is also a map in \S III.7; $c$ is also the modulus in \eqref{III:Sform}\\
$(v,s)$, $t\in T$, $u\in H$, $T_2$, $\delta$, $L_{\rm th}$ & cell (row, column), label, translate, thick columns & $s$ is a site and $v$ a valuation or a place in \S III.4--III.7; $L_{\rm th}$ is $L$ in Theorem~\ref{thm:T2}\\
$m$, $k$, $t_0$, $r$, $N_j$, $\varrho$ & the ranks (\S3) & $t_0$ is $\tau_0$ in Part~II and $t$ in \S III.4, \S III.9; $r_1,r_2$: rank of a general space\\
$\mathsf X_j$, $\mathsf X'$, $\mathsf Y_j$, $\mathsf Z_j$, $\mathsf V_j$; $h$, $J$, $y_0$ & Hermitian spaces of O12; form, Gram matrix & sans-serif because $X'$, $X''$ are cut-offs and $V_j$ spaces of forms; $J\ge1$ is also the constant of O23; $h_{\mathsf w}$ is a height\\
$L$, $L^\#$, $\Gamma_L$, $\mathrm{Cl}_j$, $\mu$, $\kappa$ & lattice, dual, automorphisms, classes, slope, number of classes & $L_0$ is a field, $L_a$ a number\\
$\rho$, $W_\rho$, $\mathcal F(\mathcal G,\rho)$, $\lambda_i$, $\lambda^\circ_i$, $\lambda^{(j)}(\tau)$ & representation, forms, highest and shifted weight & $\lambda$ without index: element of $\mathrm{Hom}(L,L')$ (\S5); $\lambda_w$: \eqref{lt:lam}\\
$\xi,\psi,\theta$, $w_j(\tau)=w_{u,j}(v,s)$, $a_t$, $b_t$, $\Delta$, $S_{\rm c}$ & weights (O24) & as in \S III.8, Part~II; $S$ alone is the constant of \eqref{ob:maj}\\
$V_j$, $K_j$, $\mathbf K_j$, $\epsilon$, $\mathbf P$, $F_{L,L'}$, $\Theta_{L,L'}$, $z$, $\nu$, $P$ & chain and kernels (\S4--5) & as in \S III.4, \S I.4.1; $K$ is a number field in \S1; $P_{u,t,j}$, $M_{u,t,j}$ are the sums of (M)\\
$\ell$, $d'$, $\varsigma$, $\mathrm{ord}_\ell$, $\epsilon_\ell$, $\Omega_\tau(v)$, $L_0$ & O20, O21 & $\varsigma$ is $s$ in \S III.5\\
$g_0,\dots,g_4$, $\mathrm{Lv}(\mathsf w)$ & level primes & $\mathrm{Lv}(\mathsf w)$ is $\mathcal T$ in Part~II, \S6.7; $\mathfrak g_q$ is a Gauss sum\\
$\mathsf w\in\mathsf W_u$, $\mathcal C_{\mathsf w}$, $\mathcal M_{\mathsf w}$, $\mathcal A$, $\mathcal A_0$, $\mathcal A_1$, $\mathrm{Per}_{\mathsf w}$, $\mathcal E_0$, $\iota$, $\mathrm{sec}$, $\mathfrak B_{\mathsf w}$, $U_{\mathsf w}$, $U'$, $\mathcal R$, $\Lambda^{\rm src}_{\mathsf w}$, $\Lambda^{\rm div}_{\mathsf w}$, $\Lambda^{\rm gl}_{\mathsf w}$, $d_{\mathsf w}$, $h_{\mathsf w}$, $s_u$ & the family and the pairing (\S6) & $\mathsf w$ is the member $\alpha$ of \S III.6; $U(\mathsf X)$ is a unitary group\\
\end{longtable}}

\subsection*{9.\quad Correspondence with the Lean development}

{\footnotesize\sloppy Files are given by short names. In the folder
\path{lean/src/Formalism/LocalRingCardinality/}: \texttt{Params}, \texttt{DataHost}, \texttt{DataSpaces},
\texttt{DataLevel}, \texttt{DataForms}, \texttt{DataCoeff}, \texttt{DataKernels}, \texttt{DataRows},
\texttt{DataBu}, \texttt{DataTies}, \texttt{DataArith}, \texttt{DataTheta}, \texttt{DataJet},
\texttt{DataParams}$_2$ (abbreviated DP$_2$), \texttt{PosOfIsomEmb},
and \texttt{D7} for the file \path{RhoDetPrimeEqMuTildeMulChiPow.lean}. In
\path{lean/src/Formalism/SharpRestriction/}: \texttt{ElevenLe}, \texttt{FamCore}, \texttt{ParConsts},
\texttt{GlobalTheoremInputs} (GTI); \texttt{DataS} in \path{CuspSpanNormBounds/}; in
\path{GlobalSchurBounds/}: \texttt{TieGlobal}, \texttt{StandardLattice}, \texttt{QLDC}
for \path{QuadraticLogDerivCongr.lean}, and \texttt{JTC} for the file
\path{IsDiscriminantTwelfthRootOfIdealLattice.lean} in \path{GeneralLinearActions/}; in
\path{RamificationFrobenius/TotientGaloisClasses/}: \texttt{CEM} for \path{ConjEMemHostFIff.lean}; in
\path{LevelLetterCocycles/}: \texttt{DataN}; in \path{InertJacobiLevels/}: \texttt{HostFam}. Elsewhere
under \path{lean/src/Formalism/}: \texttt{FPNS} is
\path{SizeWindowBrackets/FirstPrimeNeSecondPrime.lean}; \texttt{RecBridge}, \texttt{Numbers},
\texttt{Inputs} are in \path{LocalToGlobal/Divisibility/}; \texttt{TiedCH} in
\path{LocalToGlobal/Columns/}. Under \path{lean/src/Development/}: \texttt{PRQ} is
\path{PairRowIffQdRows.lean} and \texttt{TSH} is \path{ThinScalesIffHostBound.lean} in
\path{AbelianVarietyTheory/Definitions/GlobalTwisting/}; \texttt{EFK} is \path{ExistsFamPOfKits.lean} in
\path{AbelianVarietyTheory/Proofs/GlobalInterpolation/}; in \path{AutomorphicArithmetic/}: \texttt{EHW} is
\path{PrimeDiscriminants/ExistsHostWithPrimesCyclicTwoPower.lean}, \texttt{HOF} and \texttt{SOA} are
\path{hostOfFieldOf.lean} and \path{supplyOfAdmOf.lean} in \path{ShortNorms/}, \texttt{EPN} and
\texttt{TwistForm} are \path{ExistsPageN.lean} and \path{TwistForm.lean} in \path{CharacterBoxes/}.\par}

The last column says whether the final theorem
\texttt{riemannHypothesis} depends on the declaration: \emph{yes} if it is in the lists of the dependency
computation; \emph{ref} if it is not in the lists but occurs in the statement or definition of a
declaration that is (\texttt{LinkG}, \texttt{LinkS}, \texttt{LinkDIV}, \texttt{BesselId},
\texttt{Frame\_at}, \texttt{hostOfFieldOf}); \emph{no} if the lists say that it is not used; \emph{?} if
it is in no list. Only statements and definitions were read; no proof was read.

{\footnotesize\setlength{\LTleft}{0pt}\setlength{\LTright}{0pt}
\begin{longtable}{@{}p{.035\textwidth}p{.235\textwidth}p{.56\textwidth}p{.085\textwidth}@{}}
\toprule
 & This section & Declarations in the source (file:line) & Dep.\\
\midrule\endhead
\bottomrule\endfoot
1 & conditions on $q,d$; Lemmas O2, O3, O6 as required properties of $E$ & \texttt{HostHyp} (\path{FrobeniusOrbits/Statement.lean}:13); \texttt{HostRaw}, \texttt{HostDer} with \texttt{finrank\_E}, \texttt{frob}, \texttt{cmType}, \texttt{rd\_ge}, \texttt{rd\_le} (CEM:12, 87); \texttt{q\_ne\_three}, \texttt{qd\_big}, \texttt{split\_small} (DataHost:52--57) & ref\\
2 & the field on the route of the final theorem; the data of O23 on that route & \texttt{IsChiField} (EPN:19); \texttt{hostOfFieldOf}, concluding \texttt{HostGood} (HOF:158--164); \texttt{HostGood} (TSH:101--109), \texttt{HostVal}, \texttt{HostArith} (PRQ:78--92); \texttt{WinPair} (EPN:27) & yes; yes; ref; ref; yes\\
3 & a second construction, with $n=2^e$ and the further level primes $31$, $61$ & \texttt{cyclicFrame}; \texttt{exists\_host\_with\_primes\_cyclic\_two\_power}, \texttt{exists\_host\_with\_level\_prime} (EHW:53, 87; the first is used, the two existence statements are not) & yes, no, no\\
4 & classes of primes; top, thick, thin places; $\mathrm{rd}_E/2$; the translate & \texttt{primeClass}, \texttt{topPlaces}, \texttt{thickPlaces}, \texttt{IsThin}, \texttt{Host.v} (DataHost:33, 129, 133, 136, 143); \texttt{Host.u0} & ref\\
5 & Lemmas O5, O22; the sums of (M) & \texttt{partner}, \texttt{askD}, \texttt{ASK}, \texttt{MAIN} (FPNS:19, 79, 84, 89); \texttt{HostFacts.inertia\_d} (FPNS:59) & yes; the last ?\\
6 & rank conditions; \eqref{ob:inter} & \texttt{Ranks}, \texttt{HCParam} (Params:6, 50) & ref\\
7 & scale; $a_t$, $b_t$; weights of the cells & \texttt{Scale}, \texttt{SizeData}$_2$ with \texttt{aT}, \texttt{bT}, \texttt{wj}, \texttt{aT\_win}, \texttt{top\_u}, \texttt{win\_thick}, \texttt{win\_top}, \texttt{win\_thin}, \texttt{cloud\_gap}, \texttt{cloud\_v} (DP$_2$:9, 73--160) & ref\\
8 & O23, ``record'' conditions & \texttt{SizeN}, \texttt{SizeM}, \texttt{SizeK}, \texttt{SizeKRel}, \texttt{SizeT}, \texttt{SizeQD}, \texttt{SizeTT}, \texttt{SizeEta}, \texttt{SizeDel}, \texttt{SizeG}, \texttt{SizeWords}, collected in \texttt{SizesAt} (ElevenLe:24--140) & yes (\texttt{SizesAt})\\
9 & O23, conditions of the construction; $t_0\le x^{400}$ & \texttt{RankVal}, \texttt{TOv}, \texttt{nWv}, \texttt{Lwv} (PRQ:17--32); \texttt{eventually\_rank\_le\_rpow\_of\_lt\_one\_div\_fiftytwo} (\path{CharacterBoxes/}). \texttt{RankVal} occurs in the statement of \texttt{hostOfFieldOf} & ref; ?\\
10 & O23, further conditions & \texttt{Win\_at}, \texttt{Count\_at}, \texttt{Nodes\_at} (QLDC:231, 167, 211) & yes; yes; yes\\
11 & the level primes $g_2,g_3,g_4$ & \texttt{pde}, \texttt{pde'}, \texttt{pde\_mod}, \texttt{levelPrimes} (DataHost:76--84); in \texttt{HostGood}: \texttt{pde}, \texttt{pde'} $\equiv1$ mod \texttt{lcm 15 (QS*q*d)}, \texttt{q0} (TSH:107--109) & ref\\
12 & $T_2$ and $T$ installed & \texttt{T2 := T} from \texttt{hof\_thick} (HOF:116); \texttt{\{ D with TT := TT \}} from \texttt{topsChoiceH} (HOF:221--223); \texttt{hostThickExists} & yes\\
13 & the constant $J$; the scale $Y$ & \texttt{Scale.J}, \texttt{ez} (DP$_2$:16--23); \texttt{winPairs} in \texttt{supplyOfAdmOf} (SOA:29); \texttt{WinDatum.lo}, \texttt{up} (TwistForm:19--20); \texttt{PairWin} (PRQ:66) & ref; yes; ?; ref\\
14 & existence of the data for $0<\varepsilon\le\varepsilon_0$ & \texttt{exists\_famP\_of\_kits} (EFK:159); \texttt{designS} (\path{CharacterBoxes/DesignS.lean}:21); \texttt{topDesWin} & ?; yes; yes\\
15 & the twists & \texttt{WinDatum.bound}, with \texttt{VVle E} $\lfloor x^{400}\rfloor$ (TwistForm:26; EPN:12--17); \texttt{LinkV\_at} (sum over \texttt{Ico 1 R.t}) & ?; no\\
16 & Hypothesis O12 & \texttt{Spaces} with \texttt{X\_def} (definite), \texttt{sigYD} (DataSpaces:77, 87, 75); \texttt{Y\_sigval}, \texttt{clsLat} (D7:334, 162); \texttt{X\_posDef} (PosOfIsomEmb:173); \texttt{baseSpace\_isometric\_graded} (DataSpaces:126); stage $0$: \texttt{ColLevel} (DataLevel:94), entering as \texttt{LevelData} (TiedCH:15); all stages, at $\mathrm{Lv}(\mathsf w)$: \texttt{MemberForms.lv}, \texttt{fv} (DataForms:109--116) (\texttt{X\_posDef}, \texttt{ColLevel}: yes) & ref\\
17 & the record & \texttt{D7} (D7:76) & ref\\
18 & the chain \eqref{ob:chain}; O24(b) & \texttt{ColumnChain} with \texttt{cls}, \texttt{stab}, \texttt{wt}, \texttt{K}, \texttt{Kmap}, \texttt{rhoDet}, \texttt{rhoDet'}, \texttt{unf}, \texttt{sharp}, \texttt{divDe}, \texttt{col} (DataKernels:14--76); \texttt{kerOpM} (\path{ThetaLiftIndependence/Statement.lean}:47); \texttt{MemberForms} with \texttt{lam}, \texttt{lam\_eq}, \texttt{Fi}, \texttt{Fi\_zero}, \texttt{ipC} (DataForms:155--165, 227); \texttt{rhoVec} (DataCoeff:39); \texttt{wG} (Params:75) & ref\\
19 & the parameter \eqref{ob:Xpar} & \texttt{TPar.Xpar}, \texttt{Dpar}, \texttt{Interlace} (DataTheta:545, 538, 529) & no; yes; yes\\
20 & O19 in automorphic form & \texttt{ThetaX}, \texttt{TieBm} (DataTheta:569, 747) & yes, no\\
21 & the field $L_0$ of O20(a) & \texttt{D7.LL}, \texttt{galLL}, \texttt{tcLL}, \texttt{embE}, \texttt{places}, \texttt{plEq} (D7:106--138) & ref\\
22 & O20(b); sites & \texttt{ThetaJet}, \texttt{JetLoc} (DataJet:17--57) with \texttt{l}, \texttt{l\_inert}, \texttt{N0\_eq}, \texttt{Th}, \texttt{d\_even}; \texttt{letterOf}, \texttt{Pos} with \texttt{d\_eq}, \texttt{bij}, \texttt{JetTorsionCompatibility}, \texttt{chg}, \texttt{LocBasisW}, \texttt{jetTau}, \texttt{IsComb} (JTC:399--485); \texttt{BHyp} (\path{GaussJetBounds/Statement.lean}:85); \texttt{HNeZero} (DataHost:198); \texttt{TT} (\path{IncomingBounds/IncomingBound.lean}:101) (\texttt{JetTorsionCompatibility}: no; \texttt{letterOf}: the declaration in that file is not used, another of the same name is; \texttt{ThetaJet}, \texttt{JetLoc}, \texttt{BHyp}: yes; \texttt{Pos}: no; the rest: ?) & ref\\
23 & O21; periods & \texttt{SizeData$_2$.logOm} (DP$_2$:224); \texttt{logOmAt} (two definitions under \path{Development/AbelianVarietyTheory/BlockTheory/DivisorBase/}); hypothesis \texttt{hmean} of two theorems under \path{Development/AbelianVarietyTheory/Proofs/DivisorProofs/} & ref; ?; ?\\
24 & words; $\mathrm{Lv}(\mathsf w)$, $\mathcal C_{\mathsf w}$, $\mathcal M_{\mathsf w}$ & \texttt{Word}, \texttt{hcW}, \texttt{T\_norm}, \texttt{s\_Cu}, \texttt{duS} (D7:96, 144, 233, 345); \texttt{letter}, \texttt{letter\_inj}, \texttt{u\_letter}, \texttt{cloud\_gap} (DataS:22--29); \texttt{wordBits} (FamCore:12); \texttt{ColumnsAtT} with \texttt{T}, \texttt{Mu0}, \texttt{Mram} (DataKernels:104--150); \texttt{RamRank\_at}; \texttt{Mu0Tie} (StandardLattice:458) (the last two: yes; ?) & ref\\
25 & \eqref{ob:per}, \eqref{ob:B} & \texttt{PeriodMap}, \texttt{PtilCol} (DataKernels:163--195); \texttt{perCol\_base\_eq\_PtilCol}, \texttt{BuRational} (D7:169, 364); \texttt{Rows.BuDefQ} with \texttt{Bu}, \texttt{sec} (DataBu:18--32) in \texttt{Ties} (DataTies:38); \texttt{BesselId} (GTI:65) & ref; the last yes\\
26 & $\mathcal E_0$, $\mathcal A$, $\iota$, $\mathcal R$, the lattices & \texttt{Rows.Data} with \texttt{EE}, \texttt{Amb}, \texttt{iota}, \texttt{Usp}, \texttt{Qsp}, \texttt{EE0}, \texttt{D0} (DataRows:36--194); \texttt{DivData} with \texttt{RR}, \texttt{Nu}, \texttt{Bk}, \texttt{cc} (DataArith:139--166); \texttt{RR.muDual} (DataArith:98), hypothesis \texttt{tMu} (RecBridge:193); \texttt{jEE}, \texttt{sEE}, \texttt{Drows} (D7:185, 307); \texttt{MramL}, \texttt{NuLat} (DP$_2$:134) & ref\\
27 & \eqref{ob:mu}, \eqref{ob:inner} & \texttt{muDef} (\path{HarderNarasimhan/StatementPl.lean}:64); \texttt{HLat}, \texttt{sFrame} (DataArith:61, 131); \texttt{Frame\_at} (QLDC:188) & yes; ref; yes\\
28 & one record for each $u$; ties between them & \texttt{FamCore} (FamCore:71), \texttt{FamP} (ParConsts:31); \texttt{TiePts}, \texttt{TieLocal}, \texttt{ChainTied} (TieGlobal:39, 63, 97) & yes; yes\\
29 & numbers of \S III.6 & \texttt{gain}, \texttt{F}, \texttt{E}, \texttt{tau}, \texttt{io} (DP$_2$:93--97) & ref\\
30 & numbers of \S III.7 & \texttt{NumDatum} (Numbers:15), \texttt{RecNums} (Inputs:18) & yes\\
31 & objects of \S III.8 & \texttt{SData} with \texttt{HH}, \texttt{Pu} (DataS:13, 33, 43) & ref\\
32 & objects of \S III.9 & \texttt{host.E}, \texttt{R.t}; \texttt{NData} with \texttt{shaLen}, \texttt{sha\_eq}, \texttt{muAt}, \texttt{etaLoc} (DataN:55--219); \texttt{etaJ}, \texttt{etaJ\_tie} (D7:283) & ref\\
33 & \S7(a), last part & \texttt{wordFacts\_of\_clauses}, hypothesis \texttt{hBk} (RecBridge:179) & yes\\
34 & a small record & \texttt{Inhab.Rt}: $N=16$, $k=8$, $m=2$, $t_0=2s+15$ (HostFam:14) & no\\
\end{longtable}}

\emph{Remarks.} (1) \emph{The field.} A record contains a field with required properties (row~1): it is
Galois over $\Q$ of degree $4n$, with group $\{\pm1\}^2\times A$ for an abelian group $A$ of order $n$ that
maps onto $H$; it lies in a cyclotomic field of conductor $Q_Sqd$ with $2^{e+2}\mid Q_S$; it contains
square roots of $-q$ and $-d$; its Frobenius elements are as in Lemma~O3. These properties alone do not
determine the field. This generality serves the Dirichlet $L$-functions. The final theorem is the case
of the trivial character (\path{Formalism/GammaFactorNonvanishing/CritLine.lean}:29), and the hosts on its
route come from \texttt{hostOfFieldOf} with \texttt{IsChiField} (row~2), which says that the field is
fixed by the residues $a\equiv\pm1\bmod2^{e+2}$ with $\psi(a)=1$, $(a/q)=(a/d)=1$; for $\psi=1$ this is
the group $\Gamma$ of Lemma~O2, so $A=H$ and $Q_S=N$. This last step is our inference. (2) \emph{The
chain.} In \texttt{ColumnChain} the block $\mathbf K_j(L,L')$ is the composite
\texttt{divDe}\,$\circ$\,\texttt{unf}\,$\circ$\,\texttt{sharp}, the weight is \texttt{wt}${}=1/{}$\texttt{card (stab)},
and $|\epsilon|=1$ is a field for $i<m$ only. The forms are scalar functions on \texttt{cls i}\,$\times$\,\texttt{Y i};
the space $V_j$ is the subspace \texttt{Fi j}, and \texttt{Fi 0 = Mu0}. (3) \emph{The weights.}
\texttt{lam\_eq} gives $\lambda_i=u_i-\rho_k(i)$ for $i<k$, with \texttt{rhoVec} of rank $k$, and
$-$\texttt{wG} on the blocks; O24(b) is this formula, and the translation by $j(t_0+1)/2$ after
\eqref{ob:Xpar} follows from it. (4) \emph{The small record} of row~34 has sizes that cannot satisfy O23
for large $x$. It stands in a part of the source named \texttt{Inhab}; by that name it serves to show
that the type of records is not empty. We did not verify this. (5) The source describes the maps in
adelic language (\texttt{ThetaX}): cusp forms on the groups $U(\mathsf X_j)$ and $U(\mathsf Y_j)$, theta
lifts between $U(\mathsf Y_{j+1})$ and $U(\mathsf Y_j)$, and periods for the pairs
$(\mathsf X_j,\mathsf Y_j)$. Definition~O14 is the classical form for the definite groups. At the archimedean places
$U(\mathsf Y_j)$ is not compact for $0\le j\le m-2$, since the negative index of $\mathsf Y_j$ is
$\lfloor(m-j)/2\rfloor$; both $\mathsf Y_{m-1}$ and $\mathsf Y_m$ are positive definite. We gave no classical form
for the non-compact ones. (6) In O20(b) the
source has \texttt{Th = atLin (c}$\,\cdot\,$\texttt{chg \dots) F} with a constant $c$ of norm $1$; this is the
$\epsilon_\ell$ of O20(b). The names \texttt{Loc.AdmB} and \texttt{c0}, used in an earlier version of this
table, were not found in the source.

\emph{What remains.} (i) A classical description of the forms on the non-compact groups
$U(\mathsf Y_j)$, of the space $\mathcal E_0$ and of the maps $\iota$, $\mathrm{res}$. (ii) The level
structure at the stages $j\ge1$, that is, clause (viii) of O12. (iii) The identification of the functions
$F_{L,L'}$ of O19 with normalised theta series \eqref{ob:theta}, and proofs of O19--O21. (iv) The rows
marked N in \S7. (v) The dependency of the final theorem on the declarations marked ``?''. (vi) The ties
between the data of different translates (O23, last sentence).


\section{The proof of (M) from its inputs}\label{as:sec}

\emph{The plan of the proof of \textup{(M)}, in plain words.} The main inequality (M) says that a weighted
sum over prime powers, $P$, is at most a main term, $M$, plus an error. This section does not prove (M).
It \emph{reduces} (M), at one scale $x$ that is large in terms of the constants involved, to twelve numbered Hypotheses, (H0)--(H11), on the objects of
\S\ref{ob:sec}. A thirteenth Hypothesis, (H12), says at which scales and for which auxiliary primes the
twelve hold; with it the section reduces the Riemann hypothesis itself. None of the thirteen is proved
outright in this paper. The reduction of (M) is complete, and it is elementary; it is for one scale $x$
(see \S7 for the scales). Here is its course, with the section that supplies each Hypothesis.
\begin{enumerate}[label=\textup{\arabic*.},leftmargin=2.0em,itemsep=1pt,topsep=2pt]
\item \emph{The objects.} For each translate $u$, \S\ref{ob:sec} provides a finite family of linear maps
$\beta_{\mathsf w}$, one for each ``word'' $\mathsf w$. Each maps a lattice $\Lambda^{\rm gl}_{\mathsf w}$,
the source, into one fixed lattice $\mathcal R^\vee$, the target. The sizes of all parameters are (H0);
the two sets of columns $T_2$, $T$ in it are constructed in \S III.3.
\item \emph{The maps are divisible.} The source is a large lattice: its slope exceeds that of a known
lattice by a gain, and the gain is the prime sum $P$, up to losses. These are (H6) and (H7), from
\S III.6, which rests on the local estimate of \S III.5. The known lattice is controlled by (H5).
\item \emph{The maps are small.} The size of $\beta_{\mathsf w}$ at the complex places is bounded by an
explicit main term. This is (H1), from \S III.4 and \S III.7; the bound for denominators of \S III.9
enters it. Three companions, (H2), (H3), (H4), compare the levels and the heights of different words;
(H4) contains the count of forms of Part~II, \S6.6.
\item \emph{A divisible map cannot be small.} For one injective map this is the slope inequality, proved
below (Lemma~M3). For the whole family it is the mean-slope inequality (H8), which Proposition~M4 below
derives from the rank bound of \S III.8; (H9), (H10), (H11) are its side conditions.
\item \emph{The cancellation.} Insert 2 and 3 into 4. The slope of the target cancels, and what remains
is $P\le M+{}$errors for each translate (Lemma~M5, Proposition~M7). The errors are the four errors of (M)
(Lemma~M8, Theorem~M). The mean over the translates is then immediate (Corollary~M9).
\item \emph{The scales.} Items 1--5 are for one scale $x$. Parts~I and~II need (M) at every $x$ of
intervals $[Z,Z^{3/2}]$ with $Z$ arbitrarily large. At each such $x$ they need it with auxiliary primes
$q$, $d$ for which Estimates~B and~C hold at that $x$, and with constants that do not depend on $x$. That
such data exist is Hypothesis (H12), in \S7. Given (H12), Parts~I and~II apply at each $x$, and the
Riemann hypothesis follows (Corollary~M11).
\end{enumerate}
The plan is given in more detail in \S1, with a table of the Hypotheses and the status of each.

\emph{The result of the section.} The Riemann hypothesis is reduced to the Hypotheses (H0)--(H12) and
the cited theorems; none of the Hypotheses is proved outright in this paper.

\emph{What is proved, and what is not.} The section proves two implications. Theorem~M says that
\textup{(H0)--(H11)} imply \textup{(M)} for one translate at one scale, with all constants, provided
the scale is large in terms of those constants (the condition \eqref{as:big}), and
Corollary~M9 gives the mean over the translates, which is the (M) of Parts~I and~II. Corollary~M11 says
that \textup{(H12)} implies the Riemann hypothesis; (H12) contains (H0)--(H11) at each scale of its
intervals. So (M) and the Riemann hypothesis are reduced to the Hypotheses; they are not proved. None of
(H0)--(H12) is proved outright in this paper. Some are derived, here or in an earlier section, from
statements that the paper takes from the source without proof: (H0), except that every translate has a
word, from Hypothesis~O23 of \S\ref{ob:sec}, and that exception from a condition quoted in \S III.8;
(H1) in \S III.7; (H6)(i),(ii) in \S III.6, and (H6)(iii) in Lemma~M10 below; (H8) from \S III.8, by
Proposition~M4; (H10) and (H11) from statements quoted in \S III.8. One, (H7), is quoted in \S III.6.
The others, (H2), (H3), (H5), (H9), the third part of (H4), and (H12), are assumed here, from the Lean
development. Of (H4) the lower bound is quoted in Part~II in shape, with the size of its allowance
assumed, and the upper bound is a condition of O23. Of (H12) the paper proves one thing: auxiliary
primes with the properties of its clause (b) exist, for every large scale $Y$ and every $x$ in $[Y,Y^J]$
(\S III.1). Estimate~C for such primes is proved in \S III.2; it is used in Corollary~M11 and is not part of
(H12). The section does not prove that the data of \S\ref{ob:sec} exist, nor that they exist for the
scales and the primes of (H12).

\emph{Proved here.} The slope inequality for one injective map of lattices (Lemma~M3). The mean-slope
inequality for a family of maps, given a bound for the rank of the span of the images
(Proposition~M4); its proof uses three cited facts on the filtration of a lattice by slopes. An upper
bound for the largest slope of the glued lattice, from three statements of the source and an
inclusion that is part of the data (Lemma~M10). The
cancellation of the main terms (Lemma~M5). An elementary inequality for the function
$u\log(1+H/u)$ (Lemma~M6). The inequality for one translate with an explicit error (Proposition~M7). The
size of that error (Lemma~M8). Theorem~M and its corollary for the mean over the translates. The
deduction of the Riemann hypothesis from (H12), by Theorem~M, Part~II and \S III.2 (Corollary~M11).

\emph{Cited.} The existence and three properties of the filtration of a lattice by slopes, (C6) in \S3
(Stuhler 1976, Grayson 1984; from memory). It is used for Proposition~M4 only. The proof of
Corollary~M11 uses, through Part~II and \S III.2, the theorems cited there: the prime number theorem for
arithmetic progressions with a fixed modulus, Chebyshev's bound, the continuation of $\zeta$ with its
expansion at $s=1$, quadratic reciprocity (Part~II, \S4.1), the three theorems (Z1)--(Z3) on
Dirichlet $L$-functions, the functional equation of $L(s,\chi)$ and Dirichlet's theorem $L(1,\chi)\ne0$.
Its last step uses the functional equation of $\zeta$.

\emph{Assumed here, from the Lean development.} The Hypotheses (H0)--(H11) of \S4 and (H12) of \S7, with
the status that the table of \S1 gives for each.

\emph{Letters.} The notation of \S\ref{ob:sec} is used (see the table ``Letters of this section'' there):
the scale $x,\varepsilon,b,a,b',X',X''$; the field $E$ with $\mathrm{rd}_E$, the group $H$ of order $n$,
$N=4n$, $f=2n$; the parameters $m,k$; the sets $T$, $T_2$ and $\delta=|T_2|/n$; the orders $a_t$; the
weights $w_{u,j}$; a translate $u$; a word $\mathsf w\in\mathsf W_u$ with its level prime
$\mathrm{Lv}(\mathsf w)$; the coefficient field $L_0$; the spaces $U_{\mathsf w}$, $U'$; the lattices
$\Lambda^{\rm src}_{\mathsf w}$, $\Lambda^{\rm div}_{\mathsf w}$, $\Lambda^{\rm gl}_{\mathsf w}$ and
$\mathcal R^\vee$; the numbers $B_0$, $N_c$, $L_a$, $B$; the level primes $g_0,g_1$. From \S III.6 we use
$G_{\rm top}$, $F_0$, $E_{\rm err}$, $\mathrm{sg}$, $\tau(w)$, $\mathrm{io}(w)$, $\lambda_w$, $e_w$, for a
place $w$ of $L_0$ (a place is written $w$, a word $\mathsf w$). From \S III.8 we use $\delta_{\rm p}$,
$\alpha_{\rm ram}$, $e_0$. The letters introduced here are: $P_u$, $M_u$ (the two sides for one
translate); $A=\sum_{t\in T}a_t$; $\ell_g=\log(g_0g_1)$ and $\ell_{\mathsf w}=\log\mathrm{Lv}(\mathsf w)$;
$S$ (a sum, in \S4 and \S5 only); $\Lambda^{0}_{\mathsf w}$ (the first lattice of the chain);
$d_{\mathsf w}$, $h_{\mathsf w}$; $\mu_{\rm CC}$, $\mu^{\min}_{\rm CC}$; $s_u$ (as in \eqref{ob:inner});
$\mathrm{sp}_u$ (the spread); $I_u$; $\mathsf R$ (a bound for a range of slopes); $\eta_u$,
$\mathsf e_u$ (the two numbers of (H8)); the constants $c_1,\dots,c_7$, $C_{\rm abs}$, $e_L$, $C_H$,
$e_K$, $K$; the error terms $K_0$, $\mathrm{LoP}$, $\mathrm{MLo}$, $V_{\rm var}$, $U_{\rm sp}$,
$\mathrm{Err}_u$, $\mathfrak E$; and in \S3 a number field $K$ with a set $\Phi_K$ of embeddings, lattices
$\Lambda$ and $Q$ of ranks $r_1$ and $D$, maps $\phi$, $\beta$, numbers $\eta_*$, $\delta_*$, $\alpha_*$, a
level $\vartheta$, and $\mathrm{ht}$, $\mu_{\max}$, $\mu_{\min}$, $\deg$; in (C6) and in the proof of
Proposition~M4 the steps $Q_i$ of a filtration with slopes $\mu_i$, the slopes $y_j$ with multiplicity,
the step $F$, the count $\mathcal N(\vartheta)$, and for each word $V^\circ_{\mathsf w}$,
$\Lambda'_{\mathsf w}$, $\Lambda''_{\mathsf w}$, $r_{\mathsf w}$; in Lemma~M10 ideals $\mathfrak a$, $\mathfrak b$ and
integers $\nu_w$; in \S2 the smallest weight $w_{\min}$; and in \S7 the numbers $\varepsilon_1$, $J$,
$Z$, $Y$, $y$, $C_*$, $C_\varepsilon$, $x_1(\varepsilon)$.

\emph{Letters with two meanings.} The following could not be avoided without departing from the sections
on which this one builds; the context separates the meanings. $K$ is a number field in \S3 and a constant
from \S6 on; $K_0$ is an error term (in \S\ref{ob:sec} the $K_j$ are the maps of the chain). $Q$ is a
lattice in \S3 (in \S\ref{ob:sec}, \S1, $Q=Nqd$). $\Lambda$, with or without indices, is a lattice,
except $\Lambda_{\rm c}$ in Lemma~M5, which is a real number. $\varrho$ in Lemma~M5 is a positive real
number (in \S\ref{ob:sec} it is a rank). There are five letters $E$: the field $E$; the level correction
$E_{\rm err}$ (abbreviated $E_{\mathsf w}$ in the proof of Proposition~M7); the error $\mathrm{Err}_u$;
the bracket $\mathfrak E$; and $\mathcal E$, the error of (M) in Part~II. The letter $e$ with an index is
a ramification index ($e_w$) or a constant ($e_L$, $e_K$, $e_0$); the number of (H8) is written in
sans-serif, $\mathsf e_u$. $\phi$ is a linear map in \S3 and an element of $U_{\mathsf w}$ in \S4.
$\kappa_1,\kappa_2,\kappa_3$ and $z$, $z_D$, $z_L$ are abbreviations in the proof of Lemma~M8 only. The
spread is called $u$ in Part~II, \S6.10, and the bound $\mathsf R$ is called $H$ there; here $u$ is the
translate and $H$ is the group. The level $\vartheta$ in the proof of Proposition~M4 is called $t$ in
Part~II, \S6.10; here $t$ is a label. The letter $v$ is the row of a cell; the number
$v=\mathrm{rd}_E/2$ of Part~II, \eqref{II:t6.3} is written $\mathrm{rd}_E/2$. A prime power is written $p^i$, as
in \S\ref{ob:sec}; Parts~I and~II write $r=p^h$. The letter $\tau$ alone is an embedding in $\Phi_K$
(\S3); with a place as argument, $\tau(w)$ is the number of \S III.6. $I_{\mathsf w}$ is the ideal of
\eqref{lt:N}; $I_u$ is the loss sum of \S4. $D$ is the rank of the target in Proposition~M4; in the
application it is $D_{\rm r}$ (in \S III.8, $D$ is the dimension of another space). $F$ is a step of a
filtration in the proof of Proposition~M4 (in \S\ref{ob:sec} it is the real subfield of $E$). The letter
$w$ alone is a place of $L_0$; with the indices $u,j$ it is a weight. $y_j$ is a slope in \S3, and $y$ is
a real number in \S7. $Y$ in \S7 is the scale at which the primes are chosen; Part~II, \S3 writes $Y$
for the left end of an interval, which is $Z$ here. The unit $B$ of (M) is not the bound called $B$ in
Theorem~\ref{thm:C}.

\subsection*{1.\quad The plan of the proof of (M)}

Fix a translate $u$. The construction gives, for each word $\mathsf w$ of the translate, a chain of $m$
maps and, at its end, a pairing. The pairing is a linear map $\beta_{\mathsf w}$ from a space with a
lattice $\Lambda^{\rm gl}_{\mathsf w}$, the \emph{source}, to a space with a lattice $\mathcal R^\vee$, the
\emph{target}. The target is the same for all words. Two things are known about $\beta_{\mathsf w}$.

\emph{It is divisible.} The map $\beta_{\mathsf w}$ carries $\Lambda^{\rm gl}_{\mathsf w}$ into
$\mathcal R^\vee$, and $\Lambda^{\rm gl}_{\mathsf w}$ is large: it is the first lattice of the chain,
enlarged at each prime by the amount by which the values of the pairing are divisible there. In the
language of slopes (\S3; the slope of a lattice is minus the logarithm of its covolume, per unit of
rank), the slope of $\Lambda^{\rm gl}_{\mathsf w}$ exceeds that of the first lattice by a \emph{gain}
$G_{\rm top}+F_0$, up to losses (H6). The number $G_{\rm top}$ is, up to losses, the prime sum
$P_u=\sum_{t,j}P_{u,t,j}$ of (M) (H7).

\emph{It is small.} The size of $\beta_{\mathsf w}$ at the complex places, measured by a number
$h_{\mathsf w}$, is bounded above: $h_{\mathsf w}$ plus the slope of the target is at most a main term
$\sum_tm(a_t\log a_t-a_t-a_t\log\Omega_E)-A\log\mathrm{Lv}(\mathsf w)$, up to errors (H1).

\emph{A map that is divisible cannot be small.} For one injective map of lattices of the same rank this
is the inequality $\mu(\text{source})\le\mu(\text{target})+(\text{size})$ of Part~II, \S6.8
(Lemma~M3). Here the sources have smaller rank than the target, and one map controls only the largest
slope of a sublattice of the target. The remedy is to use all words together. Their images span almost
all of the target (the rank bound of \S III.8), and from this one gets the inequality for the
\emph{mean} over the words: $s_u\le\mu(\mathcal R^\vee)+{}$errors, where $s_u$ is the mean of
$\mu(\Lambda^{\rm gl}_{\mathsf w})-h_{\mathsf w}$ (Proposition~M4; (H8)).

\emph{The cancellation.} Insert the gain and the size bound into $s_u\le\mu(\mathcal R^\vee)+\dots$. The
slope of the target occurs on both sides and cancels. So does $F_0$. What remains is
$G_{\rm top}\le\sum_tm(a_t\log a_t-a_t-a_t\log\Omega_E)-A\log\mathrm{Lv}(\mathsf w)+E_{\rm err}(\mathsf w)+{}$errors.
The level prime is bounded below by a count of forms (H4), and $E_{\rm err}$ is a level correction (H2).
With these, the right side becomes $\sum_{t,j}M_{u,t,j}$ exactly: this is Lemma~M5. Adding (H7) gives
$P_u-M_u\le{}$errors for the translate $u$ (Proposition~M7).

\emph{The errors.} Each error is bounded by a constant times a power $x^{O(\varepsilon)}$ times one of
$|T|\,a$, $|T|\,a\,m^2/k$, $m\,B_0$, $m\,N_c$ (Lemma~M8). After division by the unit $B=|T|mb\delta/2$ these
are the four errors $x/m$, $xm/k$, $m$, $k/n^2$ of (M) (Theorem~M). The mean over the translates is then
taken; it needs nothing more (Corollary~M9).

\emph{The scales.} All this is for one scale $x$, one pair of auxiliary primes $q$, $d$ and one set of
data. The deduction of the Riemann hypothesis in Parts~I and~II needs the bound for $G(x)$ at every $x$
of intervals $[Z,Z^{3/2}]$ with $Z$ arbitrarily large. At each such $x$ it needs data for which (M)
holds, with primes for which Estimates~B and~C hold at that $x$, and with constants that are the same
for all $x$. Hypothesis (H12) of \S7 asserts that such data exist. Corollary~M11 then goes through
Parts~I and~II step by step and obtains the Riemann hypothesis from (H12).

\textbf{The Hypotheses and their status.} The last column uses three terms. \emph{Proved from named
statements}: derived, in the place named, from statements that are listed there and that the paper
takes from the source without proof. \emph{Quoted}: stated in the section named, without proof.
\emph{Assumed}: taken from the Lean development; no section of this paper treats it. A fourth status,
\emph{proved in the paper} without qualification, does not occur: no Hypothesis is proved outright.

{\small\setlength{\LTleft}{0pt}\setlength{\LTright}{0pt}
\begin{longtable}{@{}p{.09\textwidth}p{.20\textwidth}p{.40\textwidth}p{.23\textwidth}@{}}
\toprule
Hyp. & What it says & Where it comes from & Status\\
\midrule\endhead
\bottomrule\endfoot
(H0), except the first half of (g) & sizes of the parameters; densities of $T_2$, $T$ & the conditions of Hypothesis~O23 (table of \S\ref{ob:sec}, \S3), by a direct check for $x$ large in terms of $\varepsilon$. Sets with the densities of (h) are constructed in \S III.3 & proved here from Hypothesis~O23, which is assumed\\
(H0)(g), first half & every translate has a word: $1\le|\mathsf W_u|$ & the first condition of \eqref{tr:side}, which is quoted in \S III.8 and is used for (H8) and (H11) as well & proved from the named statement \eqref{tr:side}\\
(H1) & the size of the pairing: \eqref{III:sizefinal} & \S III.4, \S III.7: Theorem~\ref{sa:thm}, from eleven inputs, which are quoted. The bound for denominators of \S III.9 (Proposition~G2) enters through one of them & proved from named statements (for $C_{\rm abs}\ge2$)\\
(H2) & the level correction $E_{\rm err}$ & the Lean development: a statement on the local exponents at the level primes & assumed\\
(H3) & heights of two words compared & the Lean development: a comparison of one explicit number for two words & assumed\\
(H4) & the level prime and the count of forms & lower bound: Part~II, \S6.6--6.7, second line of \eqref{II:t6.5}, where the allowance $E_{\rm cnt}$ has no stated size; its size here is from the Lean development. Upper bound: a condition of Hypothesis~O23. Third part: the Lean development & mixed. Lower bound: shape quoted in Part~II, size of $E_{\rm cnt}$ assumed. Upper bound from O23. Third part assumed\\
(H5) & the first lattice of the chain & the Lean development: the covolume of the lattice of forms at stage $0$ & assumed\\
(H6)(i) & the pairing is integral on the glued lattice & \S III.6, Theorem~L4(i), from the statement (Val); (Val) rests on \S III.5 (Theorem~J, under hypothesis (B)) & proved from named statements; used only for (H8)\\
(H6)(ii) & lower bound for the slope of the glued lattice & \S III.6: \eqref{lt:slope} from the statement (Lat), then \eqref{lt:Lo} from the statements (a)--(d) of \S III.6, \S2 & proved from named statements\\
(H6)(iii) & upper bound for its largest slope & Lemma~M10 of this section, from (Lat), two more statements of the source, and the inclusion $\Lambda^{\rm div}_{\mathsf w}\subset\Lambda^{\rm src}_{\mathsf w}$ of O25(f) & proved here from named statements\\
(H7) & the gain against the prime sum; the loss $I_u$ & \S III.6, \eqref{lt:loss}; the splitting and three of the estimates are proved there, with Theorem~\ref{thm:T} & quoted\\
(H8) & the mean-slope inequality & Proposition~M4 of this section, from: the rank bound (Corollary~T2 of \S III.8, from the statements (S) with \eqref{tr:proviso}, (I), (P), and the conditions \eqref{tr:side} with $\alpha_{\rm ram}\le\frac9{16}$, \eqref{tr:six}, $C_H>0$); (H6)(i); the statement (8b), which is Theorem~\ref{sa:thm}(a); the identification of the pairings of \S III.6, \S III.7 and \S III.8 with one map $\beta_{\mathsf w}$ (Comment~(e) to O25; not checked in the source); and (C6), cited & proved here from named statements\\
(H9) & the range of the slopes & the Lean development; in the source it rests on the bound for denominators and on the statement behind (H2) & assumed\\
(H10) & $\delta_{\rm p}\le\frac12$ & Statement~(I) of \S III.8, \eqref{tr:I}, which is quoted there & proved from the quoted statement (I), for $x$ at least its threshold $x_0$\\
(H11) & ranks are positive & \S III.8: the lower bound in \eqref{tr:proviso} and the second condition of \eqref{tr:side} & proved from named statements\\
(H12) & supply on windows: data with (H0)--(H11) at every $x$ of intervals $[Z,Z^{3/2}]$, with constants independent of $x$ and with primes $q$, $d$ good at $x$ (\S7) & the Lean development. \S III.1 proves that primes as in its clause (b) exist at each large scale. Estimate~C for such primes is proved in \S III.2; it is not part of (H12) & assumed; \S III.1 proves that primes as in clause (b) exist\\
\end{longtable}}

\noindent The statements of the source behind each row have their own status, given in the section named
and in the table at the start of Part~III. \emph{The count.} The table has sixteen rows for the
thirteen Hypotheses. No row is proved outright. Nine rows are proved from named statements: (H0) except
the first half of (g), that half, (H1), (H6)(i), (H6)(ii), (H6)(iii), (H8), (H11), and (H10), which
needs one quoted statement only. One row is quoted: (H7). Five rows are assumed: (H2), (H3), (H5), (H9),
(H12). One row is mixed: (H4). For (H2), (H3), (H5), (H9) and the lower bound of (H4) we record, from the
statements of the source (their proofs were not read), what kind of statement each is: (H3) is
elementary in kind; (H2) and (H9) rest on local statements at the level primes, and (H9) also on the
bound for denominators; the lower bound in (H4), and (H5), are deep: they contain the count of forms and
the covolume of a lattice of forms.

\subsection*{2.\quad The statement of (M)}

\textbf{Definition M1 (the two sides).} Fix $x$, the field $E$ and the parameters of \S\ref{ob:sec}
(Hypothesis~O23). A \emph{cell} is a pair $(v,s)$ with $v\in\{0,1\}$ and $s\in H$; there are $2n$ cells.
For a translate $u\in H$ and a lift $1\le j\le m$ let $w_{u,j}(v,s)>0$ be the weight of the cell. For an
odd prime power $p^i$ with $p\ne q$ and a label $t\in T$ put
$\mathrm{cell}(p^i,t)=\big(\tfrac12(1+\chi_q(p^i)),\,t[p^i]\big)$. This is the cell \eqref{ob:partner}
for $v=0$, $s=t$: by Lemma~O5 for a prime $p\nmid2qd$ with $\chi_d(p)=-1$ and odd $i$, and by Lemma~O22
for the powers of $d$. All other prime powers have the coefficient $1-\chi_d(p^i)=0$ in the sum below.
For $t\in T$ and $1\le j\le m$,
\begin{equation}\label{as:P}
 P_{u,t,j}=\sum_{\substack{p\ \mathrm{odd\ prime}\\ p\ne q}}\ \sum_{i\ge1}\big(1-\chi_d(p^i)\big)\log p\,
 \Big(\frac{a_t}{p^i}-w_{u,j}\big(\mathrm{cell}(p^i,t)\big)\Big)_+ ,
\end{equation}
where $(\,\cdot\,)_+$ is the positive part, $\max(\,\cdot\,,0)$; and
\begin{equation}\label{as:Mdef}
 M_{u,t,j}=a_t\Big(\log a_t-1-\gamma+c_{\chi_d}-\langle\log w_{u,j}\rangle\Big),\qquad
 \langle\log w_{u,j}\rangle=\frac1{2n}\sum_{(v,s)}\log w_{u,j}(v,s),
\end{equation}
where $\gamma$ is Euler's constant and $c_{\chi_d}=\frac{L'}{L}(1,\chi_d)$. These are \eqref{II:t4.1} and \eqref{II:t4.2} of
Part~II. \emph{Convention:} if $a_t=0$, then $a_t\log a_t$ means $0$; so $M_{u,t,j}=0$, and also
$P_{u,t,j}=0$. (Hypothesis (H0)(c) below allows $a_t=0$. In the data of Hypothesis~O23,
$a_t\ge a>0$.) The sum \eqref{as:P} is finite: the term of $p^i$ vanishes when $p^i\ge a_t/w_{\min}$ for the
smallest weight $w_{\min}$. Put
\begin{equation}\label{D:385}
 P_u=\sum_{t\in T}\sum_{j=1}^mP_{u,t,j},\qquad M_u=\sum_{t\in T}\sum_{j=1}^mM_{u,t,j},\qquad
 A=\sum_{t\in T}a_t .
\end{equation}

\emph{Example.} For $q=11$, $d=1319$, $N=16$, $T=T_2=\{[1]\}$, $m=1$, $a=10000$ and the weights of Part~II,
\S4.2 (one weight $10$, seven weights $0.05$), Part~II, \S4.2--4.3 computes $P_u=84701.0317$ and
$M_u=84463.3942$ for $u=[1]$, and the four differences $P_u-M_u$ are $237.6375$, $-300.2570$, $-176.7314$,
$229.1329$.

\textbf{The form of (M) that the end of the proof uses.} Put
\begin{equation}\label{as:E}
 \mathfrak E=|T|\,a\,\Big(1+\frac{m^2}k\Big)+m\,(B_0+N_c),\qquad a=xb,\quad B_0=2n\,m\,b\,x^{\varepsilon},\quad
 N_c=\frac{|T_2|\,k\,b\,x^{2\varepsilon}}{n^2}.
\end{equation}
The statement is: \emph{there are real numbers $C_{\rm abs},e_L\ge0$ and $\varepsilon_1>0$ such that for
every $0<\varepsilon<\varepsilon_1$ and every admissible scale $x\ge x_0(\varepsilon)$ the data of
\S\ref{ob:sec} exist, with $\delta=|T_2|/n\ge x^{-\varepsilon}$ and $|T|\ge\frac12x^{-\varepsilon/8}n$,
and for every translate $u\in H$}
\begin{equation}\label{as:M}
 P_u-M_u\ \le\ x^{(C_{\rm abs}+e_L+3)\varepsilon}\ \mathfrak E .
\end{equation}
This is a statement for each translate. Theorem~M proves \eqref{as:M} from the Hypotheses of \S4. The
numbers $C_{\rm abs}$ and $e_L$ are the exponents in (H1) and (H7); the source calls the first $e_D$.
The two conditions on the densities are (H0)(h). Theorem~M does not prove that the data exist, and it
does not say which scales are admissible; that is the subject of Hypothesis (H12) in \S7. The statement
of the source contains further conditions on the sizes, among them $m\ge x^{1/2-\varepsilon}$ and
$k\le n^2x^{1/2-\varepsilon}$. These two are not among (H0)--(H11), and Theorem~M does not use them. They
are part of (H12), clause (e). They are used after (M), in the passage to the bound for $G(x)$
(Part~II, \S3 and \S5; Corollary~M11), where the four errors must be at most $x^{1/2+O(\varepsilon)}$.

\emph{How the end of the proof uses it.} Part~II, \S5 expresses $P_{u,t,j}-M_{u,t,j}$ as a sum, over the
cells and over the characters of $E$, of the functions $G_\chi$ at the points $a_t/w$. The sum of these
expressions over the translates $u$ removes the characters that are not trivial on $H$. The source sums
\eqref{as:M} over the translates, and only the sum is used. So \eqref{as:M} is used through its mean
over the translates, which is the form of Part~II.

\textbf{Comparison with the display (M) of Part~II, \S4.3.} Part~II states
$\frac1n\sum_{u}\sum_{t,j}(P_{u,t,j}-M_{u,t,j})\le\mathcal E$ with
$\mathcal E/B\le x^{C\varepsilon}(x/m+xm/k+m+k/n^2)$, $B=|T|mb\delta/2$. The two statements compare as
follows.
\begin{enumerate}[label=\textup{(\arabic*)},leftmargin=2.2em,itemsep=1pt,topsep=2pt]
\item \emph{Each translate.} The source proves \eqref{as:M} for every $u$. The mean over $u$ follows, with
$\mathcal E=x^{(C_{\rm abs}+e_L+3)\varepsilon}\mathfrak E$ (Corollary~M9). The display (M) of Part~II is
the mean.
\item \emph{The unnormalised error.} The bracket in the display after (M) in Part~II, \S4.3, with
$h=|T|$, $h_2=|T_2|$ and the scale $a=xb$, is $\mathfrak E$: indeed
$h_2kbx^{\varepsilon}/(n\cdot nx^{-\varepsilon})=N_c$. It is the bracket of the statement of the source
that the final theorem uses (\S9).
\item \emph{The constant.} After division by $B$, \eqref{as:M} gives
$(P_u-M_u)/B\le8\,x^{(C_{\rm abs}+e_L+3+17/8)\varepsilon}(x/m+xm/k+m+k/n^2)$ (Theorem~M). When
$8\le x^{7\varepsilon/8}$, which follows from \eqref{as:big} below, this is (M) with
$C=C_{\rm abs}+e_L+6$. The numbers $C_{\rm abs},e_L$ do not depend on $x$ or on $\varepsilon$. On the
route of the final theorem $e_L=0$, so that $C=C_{\rm abs}+6$. The number $C_{\rm abs}$ is a free
parameter $\ge0$ of the statements of the source that we read; the value with which the final theorem
uses them was not determined (\S9, Remark~(2)). In \eqref{III:sizefinal} it is $2$, and then $C=8$.
\item \emph{The orders.} Part~II has $|a_t/b-x|\le1$. The data of Hypothesis~O23 have
$a\le a_t\le a(1+1/(kx^{13}))$, which is stronger. The proof below uses only $0\le a_t\le2a$.
\item \emph{The translates that are used.} The source uses \eqref{as:M} for $u$ in the subgroup of $H$
generated by one element, which must contain $T_2$. In the setting of this paper the element is $[5]$.
The group $H$ is cyclic of order $n$ and $[5]$ generates it, because the odd residues modulo $4n=2^{e+2}$
are the numbers $\pm5^i$; so the subgroup is all of $H$, and the mean is $\frac1n\sum_{u\in H}$, as in
Part~II. In the source this was confirmed for the data whose group is cyclic of order $2^e$: there the
generator is the class of $5$ and the subgroup has index $1$. For the zeta function the statements of
the source require index $1$ for every set of data that they admit, so the mean is over all of $H$
(\S9, Remark~(5)).
\end{enumerate}

\subsection*{3.\quad Slopes: one map, and many maps}

This subsection recalls Part~II, \S6.8--6.10 over a number field and proves the two inequalities that
the assembly uses. Part~II, \S6.8 proves the first for lattices over $\Z$ and says that the form used in
the chain is quoted. It is proved here.

\textbf{Definition M2 (lattices with norms; slopes; heights).} Let $K$ be a number field of degree
$n_0$ without real embeddings, $\mathfrak o_K$ its ring of integers, and $\Phi_K$ a set of $n_0/2$
embeddings $K\to\C$ containing one of each pair of complex conjugates. (In the assembly $K$ is the
coefficient field $L_0$.) A \emph{lattice with norms} is a $K$-vector space $V$ of finite dimension $r_1$,
a finitely generated $\mathfrak o_K$-submodule $\Lambda\subset V$ that spans $V$, and for each $\tau\in\Phi_K$ a
positive definite Hermitian form (Definition~O7) on the complex vector space $V_\tau$, with norm
$\|\cdot\|_\tau$. Here $V_\tau$ is the complex vector space with the same basis as $V$, the coefficients
being embedded in $\C$ by $\tau$ (in symbols, $V\otimes_{K,\tau}\C$). Then
$\|y\|^2=\sum_{\tau\in\Phi_K}\|y_\tau\|_\tau^2$, for $y=(y_\tau)_\tau$, is a Euclidean norm on the real
vector space $V_{\mathbb R}=\bigoplus_\tau V_\tau$ of dimension $n_0r_1$. The space $V$ is embedded in
$V_{\mathbb R}$ by its images in the $V_\tau$, and $\Lambda$ is a lattice in $V_{\mathbb R}$. The
\emph{rank} of $\Lambda$ is $r_1$. For $r_1\ge1$ the \emph{slope} and the \emph{degree} are
\begin{equation}\label{as:mu}
 \mu(\Lambda)=-\frac2{n_0\,r_1}\log\mathrm{covol}(\Lambda)+\mathrm{cst}_K,\qquad\deg \Lambda=r_1\,\mu(\Lambda),\qquad
 \mathrm{cst}_K=-\log2+\frac{\log|\mathrm{disc}\,K|}{n_0},
\end{equation}
as in \eqref{ob:mu}; and $\deg0=0$. A submodule $\Lambda'\subset \Lambda$ is a lattice with norms in the $K$-span
$V'$ of $\Lambda'$, with the restricted forms. It is \emph{saturated} if $\Lambda'=\Lambda\cap V'$. For a saturated $\Lambda'$
the quotient $\Lambda/\Lambda'$ is a lattice with norms in $V/V'$: at each $\tau$ the space $V_\tau/V'_\tau$ is
identified with the orthogonal complement of $V'_\tau$ in $V_\tau$. Put
\begin{equation}\label{D:389}
 \mu_{\max}(\Lambda)=\max_{0\ne \Lambda'\subset \Lambda}\mu(\Lambda'),\qquad
 \mu_{\min}(\Lambda)=\min_{\Lambda'\ne \Lambda\ \mathrm{saturated}}\mu(\Lambda/\Lambda') ,
\end{equation}
over submodules $\Lambda'$; the extrema exist by (C6) below. So $\mu_{\min}(\Lambda)\le\mu(\Lambda)\le\mu_{\max}(\Lambda)$
(take $\Lambda'=0$ and $\Lambda'=\Lambda$). For a $K$-linear map $\phi:V\to V'$ between spaces with norms, let
$\phi_\tau:V_\tau\to V'_\tau$ be the complex linear map with the same matrix as $\phi$, and let
$\|\phi\|_\tau=\max\{\|\phi_\tau(y)\|_\tau:\ y\in V_\tau,\ \|y\|_\tau\le1\}$ be its \emph{operator norm}. The \emph{height} of $\phi\ne0$ is
\begin{equation}\label{as:ht}
 \mathrm{ht}(\phi)=\frac2{n_0}\sum_{\tau\in\Phi_K}\log\|\phi\|_\tau^2 ,
\end{equation}
the mean over the $n_0/2$ places of the logarithm of the squared norm (Part~II, \S6.7).

\emph{Example.} $K=\Q(i)$, $n_0=2$, $\Phi_K=\{\mathrm{id}\}$, $\mathrm{cst}_K=-\log2+\frac12\log4=0$. In
$V=K^2$ with the standard form let $\Lambda=\Z[i]e_1\oplus\Z[i]\,2e_2$. Its covolume in $\C^2=\mathbb R^4$ is
$4$, so $\mu(\Lambda)=-\frac12\log4=-\log2$ and $\deg \Lambda=-2\log2$. The submodule $\Z[i]e_1$ has slope $0$, and
the quotient by it is $\Z[i]\,2e_2$, of slope $-\log4$. Here $\mu_{\max}(\Lambda)=0$, $\mu_{\min}(\Lambda)=-2\log2$,
and $0+(-2\log2)=\deg \Lambda$. The map $\phi=\mathrm{diag}(2+i,3)$ of $\Z[i]^2$ into itself has
$\mathrm{ht}(\phi)=\log9$ and image of index $5\cdot9=45$.

\textbf{Lemma M3 (one injective map; proved here).} \emph{Let $\Lambda\subset V$ and $\Lambda'\subset V'$ be lattices
with norms, $\Lambda$ of rank $r_1\ge1$, and $\phi:V\to V'$ an injective $K$-linear map with $\phi(\Lambda)\subset \Lambda'$.
Then}
\begin{equation}\label{as:one}
 \mu(\Lambda)\ \le\ \mu(\phi(\Lambda))+\mathrm{ht}(\phi)\ \le\ \mu_{\max}(\Lambda')+\mathrm{ht}(\phi),
\end{equation}
\emph{where $\phi(\Lambda)$ is the submodule of $\Lambda'$ with the norms restricted from $V'$. Moreover,
for every $K$-linear map $\phi:V\to V'$, injective or not, and all $K$-subspaces $V_1\subset V$,
$V_1'\subset V'$ with $\phi(V_1)\subset V_1'$: the induced map $\bar\phi:V/V_1\to V'/V_1'$ satisfies,
for the quotient norms, $\|\bar\phi\|_\tau\le\|\phi\|_\tau$ for every $\tau$.}

\emph{Proof.} Fix $\tau$. The map $\phi_\tau$ is an injective complex linear map from $V_\tau$, of
dimension $r_1$, onto its image. In orthonormal bases it is an $r_1\times r_1$ complex matrix $A$. Let $A^*$ be its conjugate transpose.
The \emph{singular values} $\sigma_1,\dots,\sigma_{r_1}\ge0$ of $A$ are the square roots of the
eigenvalues of the Hermitian matrix $A^*A$. If $y$ is a unit eigenvector of $A^*A$ for $\sigma_i^2$, then
$\|Ay\|^2=y^*A^*Ay=\sigma_i^2$; so $\sigma_i\le\|\phi\|_\tau$. And $|\det A|^2=\det(A^*A)=\prod_i\sigma_i^2$.
As a real linear map of $\mathbb R^{2r_1}$, $A$ multiplies volumes by
$|\det A|^2\le\|\phi\|_\tau^{2r_1}$. The Euclidean spaces are orthogonal sums
over $\tau$, so $\mathrm{covol}(\phi(\Lambda))\le\mathrm{covol}(\Lambda)\prod_\tau\|\phi\|_\tau^{2r_1}$. By
\eqref{as:mu},
\begin{equation}\label{D:392}
 \mu(\phi(\Lambda))\ \ge\ \mu(\Lambda)-\frac2{n_0r_1}\sum_\tau2r_1\log\|\phi\|_\tau=\mu(\Lambda)-\mathrm{ht}(\phi).
\end{equation}
This is the first inequality of \eqref{as:one}; the second is the definition of $\mu_{\max}$. For the
last assertion, $\bar\phi_\tau$ is the composite of the inclusion of $V_{1,\tau}^\perp$ in $V_\tau$, of
$\phi_\tau$, and of the orthogonal projection of $V'_\tau$ onto $(V'_{1,\tau})^\perp$; the first and the
third have norm at most $1$. Injectivity is not used for this assertion. \qed

\noindent For $\phi=\mathrm{diag}(2+i,3)$ of the example, \eqref{as:one} reads
$0\le-\frac12\log45+\log9$, that is, $45\le81$: \emph{the index is bounded by the size}. If $\Lambda'$ has the
rank of $\Lambda$, then $\mu(\phi(\Lambda))=\mu(\Lambda')-\frac2{n_0r_1}\log[\Lambda':\phi(\Lambda)]$ (Lemma~L1 of \S III.6), and
\eqref{as:one} is the second display of Part~II, \S6.8.

\textbf{(C6) The filtration by slopes} (cited: U.~Stuhler \cite{Stu76} for $K=\Q$; D.~Grayson
\cite[\S1]{Gra84} for lattices over the integers of a number field; from memory: no local copy of either
paper was consulted, so no theorem numbers are given). \emph{Let $\Lambda$ be a
lattice with norms of rank $r_1\ge1$.}
\begin{enumerate}[label=\textup{(\roman*)},leftmargin=2.4em,itemsep=1pt,topsep=2pt]
\item \emph{The slopes of the nonzero submodules of $\Lambda$ are bounded above and their maximum is attained;
the minimum defining $\mu_{\min}(\Lambda)$ is attained.}
\item \emph{There are saturated submodules $0=\Lambda_0\subsetneq \Lambda_1\subsetneq\dots\subsetneq \Lambda_l=\Lambda$ such that
the numbers $\mu_i=\mu(\Lambda_i/\Lambda_{i-1})$ satisfy $\mu_1>\mu_2>\dots>\mu_l$, and
$\mu_{\max}(\Lambda/\Lambda_i)=\mu_{i+1}$ for $0\le i<l$. Moreover $\mu_l=\mu_{\min}(\Lambda)$.}
\item \emph{For a saturated submodule $\Lambda'$: $\deg \Lambda=\deg \Lambda'+\deg(\Lambda/\Lambda')$.}
\end{enumerate}
A chain as in (ii) is unique (this too is in the papers cited); it is called \emph{the filtration by
slopes} of $\Lambda$. Uniqueness is not used below: any chain as in (ii) serves.
Part (iii) is elementary: a $\Z$-basis of $\Lambda'$ extends to a $\Z$-basis of $\Lambda$ because $\Lambda/\Lambda'$ has no
torsion, the volume of the fundamental parallelepiped is the volume of its base in $V'_{\mathbb R}$
times the volume of its projection to the orthogonal complement, and $\mathrm{cst}_K\cdot r_1$ is additive
in $r_1$. In the example above the filtration is $0\subset\Z[i]e_1\subset \Lambda$, with $\mu_1=0$,
$\mu_2=-2\log2$. The \emph{slopes of $\Lambda$ with multiplicity} are the $r_1$ numbers $y_1\ge\dots\ge y_{r_1}$ in
which $\mu_i$ occurs $\mathrm{rk}(\Lambda_i/\Lambda_{i-1})$ times. By (iii), $\sum_jy_j=\deg \Lambda$; and
$y_j\ge\mu_{\min}(\Lambda)$ by (ii).

\begin{altbox}
\textbf{Proposition M4 (many maps: the mean-slope inequality from a rank bound; proved here from (C6)).}
\emph{Let $\mathsf W$ be a finite nonempty set. For $\mathsf w\in\mathsf W$ let $\Lambda_{\mathsf w}\subset V_{\mathsf w}$
be a lattice with norms of rank $d_{\mathsf w}\ge1$, let $Q\subset V'$ be a lattice with norms of rank
$D\ge1$, and let $\beta_{\mathsf w}:V_{\mathsf w}\to V'$ be $K$-linear with
$\beta_{\mathsf w}(\Lambda_{\mathsf w})\subset Q$. Let $h_{\mathsf w}$ be real numbers with
$\mathrm{ht}(\beta_{\mathsf w})\le h_{\mathsf w}$ whenever $\beta_{\mathsf w}\ne0$. Put}
\begin{equation}\label{as:s}
 s=\frac{\sum_{\mathsf w}d_{\mathsf w}\big(\mu(\Lambda_{\mathsf w})-h_{\mathsf w}\big)}{\sum_{\mathsf w}d_{\mathsf w}},\qquad
 \mathrm{sp}=\max_{\mathsf w}\big(\mu_{\max}(\Lambda_{\mathsf w})-h_{\mathsf w}\big)-s\ \ (\ge0).
\end{equation}
\emph{Assume the following} rank bound\emph{, with numbers $\eta_*\ge0$ and $0\le\delta_*<1$: for every
$0\le\alpha_*\le1$ and every choice of $K$-subspaces $V^\circ_{\mathsf w}\subset V_{\mathsf w}$ with
$\sum_{\mathsf w}(d_{\mathsf w}-\dim V^\circ_{\mathsf w})\le\alpha_*\sum_{\mathsf w}d_{\mathsf w}$, the subspace
$\sum_{\mathsf w}\beta_{\mathsf w}(V^\circ_{\mathsf w})$ of $V'$ has codimension at most
$(\eta_*+4\alpha_*)D/(1-\delta_*)$. Let $\mathsf R>0$ with $s-\mu_{\min}(Q)\le\mathsf R$. Then}
\begin{equation}\label{as:frame}
 \mu(Q)\ \ge\ s-\frac1{1-\delta_*}\Big(4\,\mathrm{sp}\,\log\Big(1+\frac{\mathsf R}{\mathrm{sp}}\Big)+\eta_*\,\mathsf R\Big),
\end{equation}
\emph{where $\mathrm{sp}\log(1+\mathsf R/\mathrm{sp})$ means $0$ if $\mathrm{sp}=0$.}

\noindent This is the inequality \eqref{tr:err} of \S III.8, without the term $-\log|\mathsf W|$. It
replaces \eqref{as:one} when the sources are smaller than the target. It will be used for (H8). The
numbers $\alpha_*$ and $\delta_*$ are called $\alpha'$ and $\delta_{\rm p}$ in Corollary~T2 of \S III.8; in
Part~II, \S6.9--6.10 the first is $\alpha$.

\emph{Proof.} \emph{Step 1: the filtration of the target.} Let $0=Q_0\subsetneq\dots\subsetneq Q_l=Q$ and
$\mu_1>\dots>\mu_l$ be as in (C6)(ii), and $y_1\ge\dots\ge y_D$ the slopes of $Q$ with multiplicity. Fix
$\vartheta\ge0$. Let $i$ be the largest index with $\mu_i\ge s-\vartheta$, or $i=0$ if there is none, and put $F=Q_i$. The
number of $j$ with $y_j<s-\vartheta$ is
\begin{equation}\label{D:395}
 \mathcal N(\vartheta)=D-\mathrm{rk}\,F .
\end{equation}
If $F\ne Q$, then $\mu_{\max}(Q/F)=\mu_{i+1}<s-\vartheta$.

\emph{Step 2: few directions are discarded.} Let $V'_F$ be the $K$-span of $F$. For each word put
$V^\circ_{\mathsf w}=\beta_{\mathsf w}^{-1}(V'_F)$, $\Lambda'_{\mathsf w}=\Lambda_{\mathsf w}\cap V^\circ_{\mathsf w}$, a saturated
submodule, and $r_{\mathsf w}=d_{\mathsf w}-\dim V^\circ_{\mathsf w}$, the rank of
$\Lambda''_{\mathsf w}=\Lambda_{\mathsf w}/\Lambda'_{\mathsf w}$. We claim
\begin{equation}\label{as:count}
 (\mathrm{sp}+\vartheta)\sum_{\mathsf w}r_{\mathsf w}\ \le\ \mathrm{sp}\sum_{\mathsf w}d_{\mathsf w}.
\end{equation}
Let $\mathsf w$ be a word with $r_{\mathsf w}\ge1$. Then $\beta_{\mathsf w}\ne0$ and $F\ne Q$. The map
$\beta_{\mathsf w}$ induces an injective map $\bar\beta:V_{\mathsf w}/V^\circ_{\mathsf w}\to V'/V'_F$, which
carries $\Lambda''_{\mathsf w}$ into $Q/F$. By Lemma~M3, applied to $\bar\beta$, and by its last assertion,
\begin{equation}\label{D:397}
 \mu(\Lambda''_{\mathsf w})\le\mu_{\max}(Q/F)+\mathrm{ht}(\bar\beta)\le(s-\vartheta)+h_{\mathsf w}.
\end{equation}
By (C6)(iii) and the definition of $\mu_{\max}$,
\begin{equation}\label{D:398}
 d_{\mathsf w}\,\mu(\Lambda_{\mathsf w})=\deg \Lambda'_{\mathsf w}+\deg \Lambda''_{\mathsf w}
 \le(d_{\mathsf w}-r_{\mathsf w})\,\mu_{\max}(\Lambda_{\mathsf w})+r_{\mathsf w}\,(s-\vartheta+h_{\mathsf w}).
\end{equation}
Since $\mu_{\max}(\Lambda_{\mathsf w})\le s+\mathrm{sp}+h_{\mathsf w}$ by \eqref{as:s},
\begin{equation}\label{D:399}
 d_{\mathsf w}\big(\mu(\Lambda_{\mathsf w})-h_{\mathsf w}\big)\le d_{\mathsf w}(s+\mathrm{sp})-r_{\mathsf w}(\mathrm{sp}+\vartheta).
\end{equation}
This holds also when $r_{\mathsf w}=0$, because $\mu\le\mu_{\max}$. Sum over the words. By \eqref{as:s} the
left side becomes $s\sum_{\mathsf w}d_{\mathsf w}$, and \eqref{as:count} follows.

\emph{Step 3: the rank bound.} Suppose $\mathrm{sp}+\vartheta>0$ and put $\alpha_*=\mathrm{sp}/(\mathrm{sp}+\vartheta)$, a number in $[0,1]$. By
\eqref{as:count}, $\sum_{\mathsf w}(d_{\mathsf w}-\dim V^\circ_{\mathsf w})\le\alpha_*\sum_{\mathsf w}d_{\mathsf w}$. The
space $\sum_{\mathsf w}\beta_{\mathsf w}(V^\circ_{\mathsf w})$ lies in $V'_F$, whose codimension is
$\mathcal N(\vartheta)$. By the rank bound,
\begin{equation}\label{as:Nt}
 \frac{\mathcal N(\vartheta)}D\ \le\ \frac1{1-\delta_*}\Big(\eta_*+\frac{4\,\mathrm{sp}}{\mathrm{sp}+\vartheta}\Big)\qquad(\mathrm{sp}+\vartheta>0).
\end{equation}

\emph{Step 4: integration.} For a real number $y\ge s-\mathsf R$,
$(s-y)_+=\int_0^{\mathsf R}[\,y<s-\vartheta\,]\,d\vartheta$, where the bracket is $1$ if the inequality holds and $0$
otherwise. All $y_j$ are at least $\mu_{\min}(Q)\ge s-\mathsf R$. So, by (C6)(iii),
\begin{equation}\label{D:401}
 \mu(Q)=\frac1D\sum_jy_j\ \ge\ s-\frac1D\sum_j(s-y_j)_+=s-\int_0^{\mathsf R}\frac{\mathcal N(\vartheta)}D\,d\vartheta .
\end{equation}
This is \eqref{II:t6.7} of Part~II. Insert \eqref{as:Nt} and use
$\int_0^{\mathsf R}\frac{\mathrm{sp}}{\mathrm{sp}+\vartheta}\,d\vartheta=\mathrm{sp}\log(1+\mathsf R/\mathrm{sp})$ for
$\mathrm{sp}>0$; for $\mathrm{sp}=0$ the integrand vanishes for $\vartheta>0$. \qed
The source proves the mean-slope inequality from a theorem on lattices of its own (\texttt{TheoremF});
the proof above is not a translation of it (\S9, Remark~(4)).
\end{altbox}

\emph{Example of Step 4.} Slopes $2,3,5$ and $s=4$: $\mathcal N(\vartheta)=2$ for $0<\vartheta<1$, $1$ for $1\le \vartheta<2$,
$0$ after. The integral is $3$, and the bound is $\frac{10}3\ge4-1$.

\emph{What Part~II quotes and what is proved here.} Part~II, \S6.10 quotes ``the inequality for the
minimal slope of the image of a map, and properties of the filtration of a lattice by slopes''. The first
is Lemma~M3 applied to the quotients, in Step~2. The second is (C6), which is cited. The rank bound is
the conclusion of Corollary~T2 of \S III.8; in Proposition~M4 it is a hypothesis. So the step that
Part~II, \S6.10 marks as quoted is derived here from one cited fact, (C6), given the rank bound. The
derivation is ours. The source reaches the mean-slope inequality by a theorem on lattices of its own,
whose hypotheses are not a rank bound for arbitrary subspaces (\S9, Remark~(4)). The proof was tested on
lattices of rank $2$ over $\Z$ (\S8).

\subsection*{4.\quad The Hypotheses}

Fix $\varepsilon$, a scale $x$, and a translate $u\in H$. The data of \S\ref{ob:sec} for this $u$ are
assumed given: the parameters (Hypothesis~O23), the spaces and lattices (O12), the weights (O24), the
maps of the chain (O19), and the family of words with its pairing (O25). Besides the parameters they
provide the following objects; the first column of the table in \S9 gives their names in the source.
\begin{itemize}[leftmargin=1.6em,itemsep=1pt,topsep=2pt]
\item For each word $\mathsf w\in\mathsf W_u$: the glued lattice $\Lambda^{\rm gl}_{\mathsf w}$ of
\eqref{lt:N}, in the $L_0$-vector space $U_{\mathsf w}$ of O25(e). It is a lattice with norms in the
sense of Definition~M2 for $K=L_0$, with the norms of O25(f), and its rank is $d_{\mathsf w}$. The linear
map $\beta_{\mathsf w}:U_{\mathsf w}\to\mathrm{Hom}_{L_0}(U',L_0)$,
$\beta_{\mathsf w}(\phi)=\mathfrak B_{\mathsf w}(\phi,\cdot)$, of \eqref{ob:B}. A real number
$h_{\mathsf w}$. The level prime $\mathrm{Lv}(\mathsf w)$. The \emph{first lattice of the chain}
$\Lambda^0_{\mathsf w}$, a lattice with norms (it is the lattice of the forms at stage $0$; its slope is
called $\mu_{0,\alpha}$ in \eqref{lt:Lo}). And real numbers $E_{\rm err,j}(\mathsf w)$, $1\le j\le m$,
with sum $E_{\rm err}(\mathsf w)$, as in \eqref{lt:totals}.
\item The target $\mathcal R^\vee$, a lattice with norms in $\mathrm{Hom}_{L_0}(U',L_0)$, of rank
$D_{\rm r}$. Write $\mu_{\rm CC}=\mu(\mathcal R^\vee)$ and $\mu_{\rm CC}^{\min}=\mu_{\min}(\mathcal R^\vee)$.
\item Real numbers $G_{\rm top}$, $F_0$ as in \eqref{lt:totals}, and
$I_u=\sum_w\lambda_w\,e_w\,\mathrm{io}(w)\ge0$, the sum over the places $w$ of $L_0$ of the numbers of \S III.6.
\item A real number $\delta_{\rm p}$, as in \S III.8.
\end{itemize}
\emph{The number $h_{\mathsf w}$.} The proofs of this section use only that $h_{\mathsf w}$ is a real
number satisfying the Hypotheses. In the source it is one explicit number: in the notation of \S III.7 it
is $\mathcal H-\langle\log\varkappa\rangle$, the $h$ of Theorem~\ref{sa:thm}, formed for the word
$\mathsf w$. The Hypotheses (H1), (H3), (H9) and the statement (8b) below all speak of this same number,
and (H3) has content only for it.

Put $\ell_g=\log(g_0g_1)$, and let $\Omega_E>0$ be the number with
$\log\Omega_E=\gamma+\log(2\pi/\mathrm{rd}_E)-c_{\chi_d}$, as in \eqref{II:t6.3} of Part~II. In this section
$\Omega_E$ is this number and nothing else; that it is a period is part of the proof of (H1), in
\S III.7. Let $s_u$ and $\mathrm{sp}_u$ be the numbers \eqref{as:s} for
$\Lambda_{\mathsf w}=\Lambda^{\rm gl}_{\mathsf w}$ and the given $h_{\mathsf w}$; so $s_u$ is the inner
mean \eqref{ob:inner}. Both are formed from means over the words with the weights $d_{\mathsf w}$. These
means are defined because the set of words is not empty, by (H0)(g), and every $d_{\mathsf w}$ is
positive, by (H11). Let $c_1,\dots,c_7\ge0$, $C_{\rm abs}\ge0$, $e_L\ge0$, $C_H\ge0$, $e_K\ge0$ be real
numbers. They are the constants of the Hypotheses. In the source they do not depend on $x$, on $u$ or on
the word.

\medskip
\noindent\textbf{Hypothesis (H0) (sizes).} \emph{The numbers $m\ge1$ and $n\ge1$ are integers, $T$ is not
empty, and}
\begin{enumerate}[label=\textup{(\alph*)},leftmargin=2.2em,itemsep=1pt,topsep=2pt]
\item $x\ge1$, $0<\varepsilon\le1$, $b\ge1$, $a=xb$, $\log x\le x^{\varepsilon}$;
\item $m\le x^{1/2+\varepsilon/2}$,\quad $xm\le k$,\quad $3\le k\le a$,\quad $k\le x^{e_K}$;
\item $0\le a_t\le2a$ \emph{for $t\in T$};
\item $0\le L_a\le2x^{\varepsilon}$,\quad $b'\ge0$,\quad $k\,b'\le a$,\quad $0\le X''\le a$;
\item $w_{u,j}(v,s)>0$ \emph{for all cells and lifts, and} $\mathrm{rd}_E>0$;
\item $B_0=2n\,m\,b\,x^{\varepsilon}$ \emph{and} $N_c=|T_2|\,k\,b\,x^{2\varepsilon}/n^2$;
\item $1\le|\mathsf W_u|$ \emph{and} $\log|\mathsf W_u|\le(k-3)\log2$;
\item $|T_2|\ge x^{-\varepsilon}n$ \emph{and} $|T|\ge\frac12x^{-\varepsilon/8}n$.
\end{enumerate}
\emph{Status.} For the parameters of Hypothesis~O23 (the table of \S\ref{ob:sec}, \S3) and $x$ large in
terms of $\varepsilon$, the first sentence, (a)--(f), the second half of (g) and (h) hold; this is a
direct check from the conditions of that table. For instance $k\le2x^{9/2+10\varepsilon}\le a$ and
$kb'=ka/X'\le a$ because $X'=\exp(x^{\varepsilon/3})$; $X''=xnX'\le xb$; $L_a=\log(xb\,\mathrm{rd}_E)\le2x^{\varepsilon}$
because $\log b=x^{\varepsilon}$ and $\log\mathrm{rd}_E\le x^{\varepsilon/2}$ by \eqref{ob:rd} and the sizes
of $q,d$; $a_t\le a(1+1/(kx^{13}))$; the weights are at least $w''>0$; and $|\mathsf W_u|\le2^{k-3}$. Among
these conditions the one that needs the largest $x$ is $kb'\le a$, that is, $k\le\exp(x^{\varepsilon/3})$.
Part (h) is a condition of O23 (the rows of $T_2$ and $T$); that sets with these densities exist is
proved in \S III.3 (Theorems~\ref{thm:T2}(a) and~\ref{thm:T}). The first half of (g), that the translate
has at least one word, is not among the conditions of O23. It follows from the first condition of
\eqref{tr:side}, $|\mathsf W|\ge C_0\mathsf G^2x^{C_H}$, which \S III.8 quotes: there $C_0=3$,
$\mathsf G\ge1$ and $x\ge1$, so $|\mathsf W|\ge3$. The status of (H8) and of (H11) rests on
\eqref{tr:side} in any case. So this half of (g) has the status of (H11): it is proved from a named
statement. In the source it is a hypothesis of the statement that corresponds to Theorem~M (\S9); which
statement of the source supplies it on the route of the final theorem was not followed. It is
used wherever a mean over the words is formed. Part (h) is used only for the division by $B$ in
Theorem~M. That data with these sizes exist is Hypothesis~O23; it is not proved in this paper.

\medskip
\noindent\textbf{Hypothesis (H1) (the size of the pairing).} \emph{For every word $\mathsf w$,}
\begin{equation}\label{as:H1}
 \mu_{\rm CC}+h_{\mathsf w}-F_0\ \le\ S-A\log\mathrm{Lv}(\mathsf w)+c_1\,|T|\,a\,L_a+c_1\,m\,(B_0+N_c)\,x^{C_{\rm abs}\varepsilon},
\end{equation}
\emph{where} $S=\sum_{t\in T}m\big(a_t\log a_t-a_t-a_t\log\Omega_E\big)$.

\noindent
\emph{Status.} For $C_{\rm abs}=2$ this is \eqref{III:sizefinal}, for the level
$\mathcal T=\mathrm{Lv}(\mathsf w)$ and $h=h_{\mathsf w}$; for $C_{\rm abs}\ge2$ it follows from that
case. \S III.7 (Theorem~\ref{sa:thm}) proves \eqref{III:sizefinal} from the inputs (S1)--(S8), (W1), (W2),
(G) there, which are quoted. So (H1) is proved in the paper from these named inputs, given the
identification of the objects of \S III.7 with those of \S\ref{ob:sec} (the dictionary of
\S\ref{ob:sec}, \S7, whose rows for \S III.7 are not verified). In the source $C_{\rm abs}$ is a free
parameter $\ge0$ of the statement; the value with which the final theorem uses it was not determined
(\S9, Remark~(2)).

\emph{How the bound for denominators enters.} Proposition~G2 of \S III.9 gives \eqref{gc:div} from the
bound \eqref{II:shabound} for Galois cohomology. It enters (H1) through the input (S8), the size of the
local errors. In the source (S8) is an inequality whose left side is a sum over the lifts, and each
summand contains a bound for the left side of \eqref{gc:div}. We read the statements that show this
(\S9, Remark~(6)). One link was not followed: that the bound which occurs in (S8) is, on the route of
the final theorem, the right side of \eqref{gc:div}. The bound for denominators enters (H9) as well; see
there.

\medskip
\noindent\textbf{Hypothesis (H2) (the level correction).} \emph{For every word and every lift $j$,
$|E_{\rm err,j}(\mathsf w)-A\,\ell_g|\le c_2B_0$. Hence
$|E_{\rm err}(\mathsf w)-m\,A\,\ell_g|\le c_2\,m\,B_0$.} Assumed here, from the Lean development. It is the
``local level correction'' of Part~II, \S6.7. In the source it is a statement on the local error
exponents: their weighted sum is $A\log(g_0g_1)$ up to a multiple of $B_0$, and the part away from
$g_0,g_1$ is at most a multiple of $B_0$. It needs the local exponents at the level primes; no section
of this paper treats it.

\medskip
\noindent\textbf{Hypothesis (H3) (the heights of two words).} \emph{For all words $\mathsf w,\mathsf w'$,}
\begin{equation}\label{D:403}
 \Big|\big(h_{\mathsf w}+A\log\mathrm{Lv}(\mathsf w)\big)-\big(h_{\mathsf w'}+A\log\mathrm{Lv}(\mathsf w')\big)\Big|
 \le c_3\big(|T|\,a\,L_a+m\,B_0\big).
\end{equation}
Assumed here, from the Lean development. By (H1) the number $h_{\mathsf w}+A\log\mathrm{Lv}(\mathsf w)$ is
bounded above by a quantity that does not depend on the word; (H3) says that it varies little. In the
source it compares the explicit number $h_{\mathsf w}$ for two words. It is elementary in kind: two
explicit formulas are compared. We did not read its proof.

\medskip
\noindent\textbf{Hypothesis (H4) (the level prime and the count of forms).} \emph{For all words
$\mathsf w,\mathsf w'$,}
\begin{equation}\label{as:H4}
 \sum_{j=1}^m\Big(\langle\log w_{u,j}\rangle+\log\frac{\mathrm{rd}_E}{2\pi}\Big)+m\,\ell_g
 -c_4\Big(L_a+\frac{m^2L_a}k\Big)\ \le\ \log\mathrm{Lv}(\mathsf w)\ \le\ c_4\,m\,x^{\varepsilon},
\end{equation}
\emph{and $|\log\mathrm{Lv}(\mathsf w)-\log\mathrm{Lv}(\mathsf w')|\le c_4\log k$.} \emph{Status.} The
lower bound has the shape of the second line of \eqref{II:t6.5} of Part~II, with $E_{\rm cnt}=c_4(L_a+m^2L_a/k)$;
Part~II, \S6.6 derives the shape from the mass formula, which is quoted. Part~II leaves the allowance
$E_{\rm cnt}$ unspecified and says, in \S6.10, that its size is part of the quoted statements. The size
$c_4(L_a+m^2L_a/k)$ is not in Part~II; it is assumed here, from the Lean development. The lower bound is a
deep statement: it contains the count of forms. The upper bound is a condition of Hypothesis~O23 (the row of $\mathrm{Lv}(\mathsf w)$ in its
table), with $c_4$ at least the constant $C$ there. The third part is assumed here, from the Lean
development; we did not read how the source proves it.

\medskip
\noindent\textbf{Hypothesis (H5) (the first lattice of the chain).} \emph{For every word,}
\begin{equation}\label{D:405}
 -c_5\,|T|\,a\,L_a\Big(1+\frac mk\Big)\le\mu(\Lambda^0_{\mathsf w})\qquad\text{and}\qquad
 \mu_{\max}(\Lambda^0_{\mathsf w})\le c_5\,|T|\,a\,L_a .
\end{equation}
Assumed here, from the Lean development. It is a statement on the covolume of the lattice of forms at
stage $0$. It is deep; in the source it is deduced from statements of the kind that give the lower bound
of (H4).

\medskip
\noindent\textbf{Hypothesis (H6) (the glued lattice).} \emph{For every word:
\textup{(i)} $\beta_{\mathsf w}(\Lambda^{\rm gl}_{\mathsf w})\subset\mathcal R^\vee$;
\textup{(ii)}}
\begin{equation}\label{as:H6}
 \mu(\Lambda^{\rm gl}_{\mathsf w})\ \ge\ \mu(\Lambda^0_{\mathsf w})+G_{\rm top}+F_0-E_{\rm err}(\mathsf w)-\mathrm{Lo}_{\mathsf w},
 \qquad
 \mathrm{Lo}_{\mathsf w}=c_6I_u+c_6|T|\,m\Big(\frac ak+b'\Big)\log\mathrm{Lv}(\mathsf w)+c_6X''+1;
\end{equation}
\emph{\textup{(iii)}
$\mu_{\max}(\Lambda^{\rm gl}_{\mathsf w})\le\mu_{\max}(\Lambda^0_{\mathsf w})+G_{\rm top}+F_0-E_{\rm err}(\mathsf w)+c_6\,m\,B_0$.}

\noindent\emph{Status.} Part (i) is Theorem~L4(i) of \S III.6, proved there from the statement (Val).
The statement (Val) is where the local estimate of \S III.5 (Theorem~J, under hypothesis (B)) is used.
Part (i) is \emph{not used} in the proof of Proposition~M7 or of Theorem~M. It is listed because the
derivation of (H8) from Proposition~M4 uses it; see (8a) below. Part (ii) is \eqref{lt:Lo}; \S III.6
proves \eqref{lt:slope} from the statement (Lat) (Theorem~L4(iii)) and passes to \eqref{lt:Lo} with the
statements (a)--(d) of \S III.6, \S2. These are four items; they contain at least eight statements of
the source, all assumed there. Part (iii) is not treated in \S III.6. It
is proved in Lemma~M10 (\S4a below), from (Lat) and two further statements of the source, and in a
stronger form, without the term $c_6mB_0$. The term is kept in (iii) only to match the source. So each
part is proved in the paper from named statements of the source; none is proved outright.

\medskip
\noindent\textbf{Hypothesis (H7) (the losses).}
\begin{equation}\label{as:H7}
 P_u-G_{\rm top}\ \le\ \mathrm{LoP}\qquad\text{and}\qquad I_u\ \le\ \mathrm{LoP},\qquad
 \mathrm{LoP}=c_7\,|T|\,m\,b\,x^{e_L\varepsilon}.
\end{equation}
\emph{Status.} This is \eqref{lt:loss} of \S III.6 (there $e_L=0$), quoted there. On the route of the
final theorem the source has $e_L=0$ as well (\S9, Remark~(2)); the exponent is kept here because the
statement of the source that corresponds to Theorem~M carries it. \S III.6 proves the
splitting of the left side by classes of primes and three of the estimates; one of them uses
Theorem~\ref{thm:T}. The identity between $P_u$ and the sum that the source bounds is Lemma~L6 of
\S III.6 with a statement of the source that was not read (\S9).

\medskip
\noindent\textbf{Hypothesis (H8) (the mean-slope inequality).} \emph{There are real numbers
$\eta_u,\mathsf e_u\ge0$ with $(\eta_u+4\sqrt{\mathsf e_u})\,x^{C_H}\le6$ such that for every $\mathsf R>0$ with
$s_u-\mu^{\min}_{\rm CC}\le\mathsf R$,}
\begin{equation}\label{as:H8}
 \mu_{\rm CC}\ \ge\ s_u-\log|\mathsf W_u|-\frac1{1-\delta_{\rm p}}
 \Big(4\,\mathrm{sp}_u\log\Big(1+\frac{\mathsf R}{\mathrm{sp}_u}\Big)+(\eta_u+4\sqrt{\mathsf e_u})\,\mathsf R\Big).
\end{equation}
\emph{Status.} This is \eqref{tr:err} with \eqref{tr:six}. It follows from Proposition~M4 and the
following three statements, by taking $K=L_0$, $\Lambda_{\mathsf w}=\Lambda^{\rm gl}_{\mathsf w}$,
$Q=\mathcal R^\vee$, $\delta_*=\delta_{\rm p}$, so that $s=s_u$ and $\mathrm{sp}=\mathrm{sp}_u$:
\begin{enumerate}[label=\textup{(8\alph*)},leftmargin=2.6em,itemsep=1pt,topsep=2pt]
\item (H6)(i); it is the hypothesis $\beta_{\mathsf w}(\Lambda_{\mathsf w})\subset Q$ of Proposition~M4;
\item $\mathrm{ht}(\beta_{\mathsf w})\le h_{\mathsf w}$ whenever $\beta_{\mathsf w}\ne0$ (this is what ``$h$ is
an upper bound for the height of the pairing map'' means in Part~II, \S6.7). This is not an extra
assumption of the source: it is part~(a) of Theorem~\ref{sa:thm} of \S III.7, where the map is called
$\Phi$ and $h_{\mathsf w}$ is $\mathcal H-\langle\log\varkappa\rangle$. So it has the status of (H1): it
is proved there from the inputs (S1)--(S8), (W1), (W2), (G), given the same identification of the map
and of the norms;
\item the rank bound of Proposition~M4 with
$\eta_*=5x^{-C_H}+4\alpha_{\rm ram}+4\sqrt{e_0}$, which is the conclusion \eqref{tr:final} of
Corollary~T2 with $\alpha'=\alpha_*$. (A $K$-subspace $V^\circ_{\mathsf w}$ of $U_{\mathsf w}$ gives the complex
subspace $\mathcal U_{\mathsf w}$ of $\mathcal M_{\mathsf w}$ with the same basis, the codimension of
$\sum\beta_{\mathsf w}(V^\circ_{\mathsf w})$ over $L_0$ is the complex codimension of
$\sum\beta_{\mathsf w}(\mathcal U_{\mathsf w})$, and \eqref{tr:sub} holds because
$d_{\mathsf w}=\dim\mathcal M_{\mathsf w}\le\dim\mathcal C_{\mathsf w}$. In \S III.8 the letter $d_{\mathsf w}$ is
$\dim\mathcal C_{\mathsf w}$ and our $d_{\mathsf w}$ is $d^{\rm ram}_{\mathsf w}$.)
\end{enumerate}
Indeed \eqref{as:frame} is then \eqref{as:H8} without the term $\log|\mathsf W_u|$, which is $\ge0$ by
(H0)(g), and with $\eta_*$ in place of $\eta_u+4\sqrt{\mathsf e_u}$; and for the values
$\eta_u=5x^{-C_H}+3\sqrt{\alpha_{\rm ram}}$, $\mathsf e_u=2e_0$ of \eqref{tr:frame},
$\eta_*\le\eta_u+4\sqrt{\mathsf e_u}$ as soon as $\alpha_{\rm ram}\le\frac9{16}$. The bound
$(\eta_u+4\sqrt{\mathsf e_u})x^{C_H}\le6$ is \eqref{tr:six}. Corollary~T2 is proved in \S III.8 from the
statements (S) with \eqref{tr:proviso}, (I), (P) and the side conditions \eqref{tr:side}, which are
quoted. It requires $C_H>0$, and $x^{-C_H}\le\frac14$ (the third condition of \eqref{tr:side}). So the
derivation of (H8) needs $C_H>0$; Theorem~M itself holds for $C_H\ge0$. Proposition~M4 also needs
$\mathsf W_u\ne\emptyset$ and $d_{\mathsf w}\ge1$, which are (H0)(g) and (H11). \emph{One map, three
sections.} Proposition~M4 needs one map $\beta_{\mathsf w}$ for which (8a), (8b) and (8c) hold at once.
But (8a) comes from \S III.6, for the pairing $\langle\ ,\ \rangle_\alpha$ there; (8b) from \S III.7, for
the map $\Phi$ there; and (8c) from \S III.8, for the pairing $\mathfrak B^{\rm alg}_{\mathsf w}$ there.
That the three are the map $\beta_{\mathsf w}$ of \eqref{ob:B} is Comment~(e) to O25 and the dictionary
of \S\ref{ob:sec}, \S7. It is a further assumption: \S III.8, \S2 records that the identification of
$\mathfrak B^{\rm alg}_{\mathsf w}$ with $\langle\ ,\ \rangle_\alpha$ was not checked in the source, and
Comment~(e) says the same. The
conditions of Corollary~T2 are not derived from (H0)--(H11); they are further named statements. In sum,
(H8) is proved here from named statements: (S), (I), (P), \eqref{tr:proviso}, \eqref{tr:side} with
$\alpha_{\rm ram}\le\frac9{16}$, \eqref{tr:six}, $C_H>0$, (Val), the inputs of Theorem~\ref{sa:thm}, the
identification of the three maps (Comment~(e) to O25; not checked in the source), and (C6), which is
cited.

The source does not go through Corollary~T2 and Proposition~M4; it proves \eqref{as:H8} from a theorem
on lattices of its own (\S9, Remark~(4)). When \eqref{as:H8} is taken from the source in that way, it
comes with three further conditions: $\mathsf e_u\le\frac14$, $\alpha_{\rm ram}\le\frac14$ and
$x^{-C_H}\le\frac1{20}$ (\S III.8, after \eqref{tr:frame}). They are not needed in the proofs of this
section and are not part of (H8).

\medskip
\noindent\textbf{Hypothesis (H9) (the range of the slopes).} $s_u-\mu^{\min}_{\rm CC}\le x^{C_H}\,b$.
Assumed here, from the Lean development. (The source states it as
$s_u+\mu_{\max}(\mathcal R)\le x^{C_H}b$, with $\mu_{\min}(\mathcal R^\vee)=-\mu_{\max}(\mathcal R)$.) It
is a crude bound, but it is not elementary. In the source, the statement that yields it has among its
hypotheses the bound for denominators, that is, the conclusion \eqref{gc:div} of Proposition~G2 of
\S III.9, and the statement behind (H2). We read these statements, not their proofs (\S9, Remark~(6)).

\medskip
\noindent\textbf{Hypothesis (H10).} $0\le\delta_{\rm p}\le\frac12$. This follows from Statement~(I) of
\S III.8, \eqref{tr:I}, for $x\ge\max(x_0,1)$, where $x_0$ is the threshold of that statement: then
$x^{\varepsilon_z/7}\ge1>\log2$. Statement (I) is quoted there; it holds for $x\ge x_0$ only.

\medskip
\noindent\textbf{Hypothesis (H11).} \emph{For every word, $d_{\mathsf w}\ge1$.} \emph{Status.} It follows
from statements that \S III.8 already quotes. In the notation of \S III.8 our $d_{\mathsf w}$ is
$d^{\rm ram}_{\mathsf w}=\dim\mathcal M_{\mathsf w}$, since $U_{\mathsf w}$ is spanned over $L_0$ by a
basis of $\mathcal M_{\mathsf w}$ (O25(e)). The lower bound in \eqref{tr:proviso} gives
$\dim\mathcal C_{\mathsf w}>0$, and the second condition of \eqref{tr:side}, with $\alpha_{\rm ram}<1$,
gives $d^{\rm ram}_{\mathsf w}\ge(1-\alpha_{\rm ram})\dim\mathcal C_{\mathsf w}>0$. So (H11) follows from
\eqref{tr:proviso} and \eqref{tr:side}, which the status of (H8) uses in any case. (That
$\mu(\Lambda^{\rm gl}_{\mathsf w})\le\mu_{\max}(\Lambda^{\rm gl}_{\mathsf w})$ is part of Definition~M2;
in the source it is a separate clause.)

\medskip
\noindent\emph{Remarks.} (1) The Hypotheses speak of real numbers attached to the data: slopes, heights,
logarithms of primes. The deduction of (M) from them is arithmetic of real numbers. The lattices
themselves enter only in the status of (H6)(iii) and of (H8). (2) The gain $G_{\rm top}$ occurs in (H6) and (H7) and
nowhere else; $F_0$ occurs in (H1) and (H6); $\mu_{\rm CC}$ in (H1) and (H8); $\log\Omega_E$ in (H1) only.
(3) Proposition~G2, the sets $T_2$ and $T$, Theorem~J and the Bessel identity (Statement~(P) of
\S III.8) do not occur in the list. They are used inside the Hypotheses: Proposition~G2 in (H1) and in
(H9); the sets $T_2$ and $T$ in (H0) and (H7); Theorem~J in (H6); Statement~(P) in (H8). (4) The proof
of Theorem~M uses (H0)--(H11) with one exception, (H6)(i). That statement and (8b) serve only to derive
(H8) from Proposition~M4.

\subsection*{4a.\quad The largest slope of the glued lattice}

This subsection proves the statement on which the status of (H6)(iii) rests. It is not used in the
proofs of \S5 and \S6, which take (H6)(iii) as a Hypothesis.

\medskip
\noindent\textbf{Lemma M10 (the largest slope of the glued lattice; proved here from three statements of
the source and the inclusion $\Lambda^{\rm div}_{\mathsf w}\subset\Lambda^{\rm src}_{\mathsf w}$ of O25(f)).} \emph{Fix a word $\mathsf w$; it is a member $\alpha$ in the language of \S III.6, whose
notation \textup{(\S1--2 there)} we use. Assume:}
\begin{enumerate}[label=\textup{(\alph*)},leftmargin=2.2em,itemsep=1pt,topsep=2pt]
\item \emph{the inclusion in \textup{(Lat)}: for every place $w$ of $L_0$ with
$p_w\ne\mathrm{Lv}(\mathsf w)$ the number $\tau(w)\ge0$ is an integer and
$(\Lambda^{\rm div}_{\mathsf w})_w\subset\mathfrak p_w^{e_w\tau(w)}(\Lambda^{\rm src}_{\mathsf w})_w$;}
\item \emph{$\tau(w)=0$ at the places $w$ with $p_w=\mathrm{Lv}(\mathsf w)$ \textup{(}statement
\textup{(b)} of \S III.6, \S2\textup{)};}
\item $\mu_{\max}(\Lambda^{\rm src}_{\mathsf w})\le\mu_{\max}(\Lambda^0_{\mathsf w})$.
\end{enumerate}
\emph{Then}
\begin{equation}\label{as:M10}
 \mu_{\max}(\Lambda^{\rm gl}_{\mathsf w})\ \le\ \mu_{\max}(\Lambda^0_{\mathsf w})+G_{\rm top}+F_0-E_{\rm err}(\mathsf w).
\end{equation}
\emph{In particular \textup{(H6)(iii)} holds, for every $c_6\ge0$.}

\noindent Here $\Lambda^{\rm div}_{\mathsf w}\subset\Lambda^{\rm src}_{\mathsf w}$ are the lattices of
O25(f) in $U_{\mathsf w}$, with the norms of O25(f), and
$\Lambda^{\rm gl}_{\mathsf w}=I_{\mathsf w}\Lambda^{\rm div}_{\mathsf w}$ as in \eqref{lt:N}. The numbers
$\mathrm{sg}_{\mathsf w}(w)$ of \eqref{lt:sg} and $\tau(w)$ vanish outside a finite set of places.
Statement (c) compares the source lattice with the first lattice of the chain, of which it is the
preimage. It is a statement of the source (\S9), a companion of the equality for the slopes that
\S III.6, \S2(d) quotes. All three statements are assumed. The proof uses a fourth statement of the
source: the inclusion $\Lambda^{\rm div}_{\mathsf w}\subset\Lambda^{\rm src}_{\mathsf w}$, at the places
above the level. In this paper it is part of the data (O25(f); \S III.6, \S2), so it is already assumed.
In the source it is a separate clause (\S9). The lemma is the deduction. (It is numbered M10 because it
was written after Corollary~M9.)

\emph{Proof.} Two remarks on a lattice with norms $\Lambda$ in an $L_0$-vector space $V$ come first.

\emph{\textup{(1)} A sublattice.} Let $\Lambda'\subset\Lambda$ be a nonzero submodule. Every submodule
of $\Lambda'$ is a submodule of $\Lambda$, with the same norms. So
$\mu_{\max}(\Lambda')\le\mu_{\max}(\Lambda)$, by Definition~M2. Next, for a nonzero subspace
$V_1\subset V$ the module $\Lambda\cap V_1$ is a
lattice in $V_1$, and every submodule $\Lambda_1\subset\Lambda$ with span $V_1$ lies in it. Then
$\mathrm{covol}(\Lambda_1)=[\Lambda\cap V_1:\Lambda_1]\cdot\mathrm{covol}(\Lambda\cap V_1)\ge\mathrm{covol}(\Lambda\cap V_1)$,
so $\mu(\Lambda_1)\le\mu(\Lambda\cap V_1)$. Hence $\mu_{\max}(\Lambda)$ is the maximum of
$\mu(\Lambda\cap V_1)$ over the nonzero subspaces $V_1$.

\emph{\textup{(2)} Multiplication by an ideal.} Let $\nu_w$ be integers, zero for almost all places $w$,
and $\mathfrak a=\prod_w\mathfrak p_w^{-\nu_w}$. Then
\begin{equation}\label{D:410}
 \mu_{\max}(\mathfrak a\Lambda)=\mu_{\max}(\Lambda)+2\sum_w\lambda_w\nu_w .
\end{equation}
Indeed $\mathfrak a\Lambda\cap V_1=\mathfrak a(\Lambda\cap V_1)$ for every subspace $V_1$: the right side
is contained in the left; and $\mathfrak a^{-1}(\mathfrak a\Lambda\cap V_1)\subset\Lambda\cap V_1$, which
gives the other inclusion after multiplication by $\mathfrak a$. By Lemma~L2(ii) of \S III.6, applied in
$V_1$, $\mu(\mathfrak a(\Lambda\cap V_1))=\mu(\Lambda\cap V_1)+2\sum_w\lambda_w\nu_w$. Now take the
maximum over $V_1$ and use the second half of (1) for $\Lambda$ and for $\mathfrak a\Lambda$.

\emph{The inclusion.} Put $\mathfrak b=\prod_w\mathfrak p_w^{e_w\tau(w)}$, the product over all places; by
(a) and (b) it is an ideal of the ring of integers, and by (b) the
places above the level do not occur in it. We claim that
$\Lambda^{\rm div}_{\mathsf w}\subset\mathfrak b\Lambda^{\rm src}_{\mathsf w}$. By Lemma~L2(i),
$(\mathfrak b\Lambda^{\rm src}_{\mathsf w})_w=\mathfrak p_w^{e_w\tau(w)}(\Lambda^{\rm src}_{\mathsf w})_w$ for every
$w$. At a place with $p_w\ne\mathrm{Lv}(\mathsf w)$ this contains $(\Lambda^{\rm div}_{\mathsf w})_w$ by
(a). At a place above the level the exponent is $0$, and
$(\Lambda^{\rm div}_{\mathsf w})_w\subset(\Lambda^{\rm src}_{\mathsf w})_w$ because
$\Lambda^{\rm div}_{\mathsf w}\subset\Lambda^{\rm src}_{\mathsf w}$. An element that lies in all
localisations of a lattice lies in the lattice (proof of Lemma~L2(i)). This proves the claim.

\emph{Conclusion.} By \eqref{lt:N} and the claim,
$\Lambda^{\rm gl}_{\mathsf w}=I_{\mathsf w}\Lambda^{\rm div}_{\mathsf w}\subset I_{\mathsf w}\mathfrak b\,\Lambda^{\rm src}_{\mathsf w}$,
and $I_{\mathsf w}\mathfrak b=\prod_w\mathfrak p_w^{-\nu_w}$ with
$\nu_w=\lfloor e_w\,\mathrm{sg}_{\mathsf w}(w)\rfloor-e_w\tau(w)$. Use the first half of (1), then (2), then
$\lfloor c\rfloor\le c$ for real $c$, $\lambda_w>0$, and the definition \eqref{lt:lam} of $\mathrm{sz}_w$:
\begin{equation}\label{D:411}
 \mu_{\max}(\Lambda^{\rm gl}_{\mathsf w})\ \le\ \mu_{\max}(\Lambda^{\rm src}_{\mathsf w})+2\sum_w\lambda_w\nu_w
 \ \le\ \mu_{\max}(\Lambda^{\rm src}_{\mathsf w})+\sum_w\mathrm{sz}_w\big(\mathrm{sg}_{\mathsf w}(w)-\tau(w)\big).
\end{equation}
By \eqref{lt:sgsum} the last sum is $G_{\rm top}+F_0-E_{\rm err}(\mathsf w)$. Now use (c). \qed

\noindent The lemma is the upper counterpart of Theorem~L4(iii). There the lengths \eqref{lt:L} are used
and the inclusion is not; here it is the other way round. No term for the rounding and none for the rank
occur, because $\lfloor c\rfloor\le c$ costs nothing and no index is needed. The source proves its
statement in the same way (\S9). \emph{Example} over $\Q$, where $\lambda_p=\log p$, $e_p=1$ and the
slope is $-\frac2{\mathrm{rk}}\log\mathrm{covol}$: let $\Lambda^{\rm src}=\Z^2$, $\tau(2)=1$,
$\Lambda^{\rm div}=2\Z\oplus4\Z\subset2\Z^2$ and $\mathrm{sg}(2)=3$. Then
$\Lambda^{\rm gl}=\frac18\Lambda^{\rm div}=\frac14\Z\oplus\frac12\Z$. Its largest slope is that of
$\frac14\Z$, namely $2\log4$ (the lattice itself has slope $\log8$). The bound is
$\mu_{\max}(\Z^2)+2\log2\cdot(3-1)=4\log2$: equality.

\subsection*{5.\quad The inequality for one translate}

\textbf{Lemma M5 (the cancellation; proved here).} \emph{Let $a_t\ge0$ $(t\in T)$, $A=\sum_ta_t$, let the
weights be positive, and let $\ell_g$ and $\varrho>0$ be real numbers. Put
$\log\Omega=\gamma+\log(\pi/\varrho)-c_{\chi_d}$,
$S=\sum_tm(a_t\log a_t-a_t-a_t\log\Omega)$ and
$\Lambda_{\rm c}=\sum_{j=1}^m\big(\langle\log w_{u,j}\rangle+\log(\varrho/\pi)\big)+m\,\ell_g$. Then}
\begin{equation}\label{as:cancel}
 S-A\,\Lambda_{\rm c}+m\,A\,\ell_g\ =\ \sum_{t\in T}\sum_{j=1}^mM_{u,t,j}=M_u .
\end{equation}
\emph{Here $a_t\log a_t$ means $0$ for $a_t=0$.}

\emph{Proof.} $S=m\sum_ta_t(\log a_t-1-\gamma+c_{\chi_d})-mA\log(\pi/\varrho)$, and
$-A\Lambda_{\rm c}=-A\sum_j\langle\log w_{u,j}\rangle-mA\log(\varrho/\pi)-mA\ell_g$. In the sum the terms
$\mp mA\log(\pi/\varrho)$ cancel, and so do $\mp mA\ell_g$. What remains is
$\sum_j\sum_ta_t(\log a_t-1-\gamma+c_{\chi_d}-\langle\log w_{u,j}\rangle)$. \qed

\noindent The lemma is used with $\varrho=\mathrm{rd}_E/2$. It is \eqref{II:t6.6} of Part~II, summed over $t$ and
$j$. Three things cancel in the assembly. (a) \emph{The period against the count.} The number
$\log\Omega_E$ comes from the size bound (H1), with the sign $-m\,a_t\log\Omega_E$. Its part
$\log(2\pi/\mathrm{rd}_E)$ meets $+\log(\mathrm{rd}_E/2\pi)$ from the lower bound (H4) for the level, and
the two cancel. What is left of $\log\Omega_E$ is $\gamma-c_{\chi_d}$, and it gives the terms
$-\gamma+c_{\chi_d}$ of $M_{u,t,j}$. So in the assembly $\log\Omega_E$ enters once, and one part of it
cancels against the count. (Inside the proof of (H1) the periods occur with both signs, in the height and
in the slope, and cancel up to an error: see the table after Theorem~\ref{sa:thm}.) (b) \emph{The level
primes $g_0,g_1$.} The term $m\,A\,\ell_g$ of the count (H4) cancels against the main part of the level
correction $E_{\rm err}$ (H2). (c) \emph{The weights.} The term $-A\sum_j\langle\log w_{u,j}\rangle$ comes
from (H4) alone. The main term $a_t\log a_t-a_t$ comes from (H1) alone and does not cancel against
anything here: it is matched by the prime sum, through (H6) and (H7).

\textbf{Lemma M6 (proved here).} \emph{Let $0\le\mathrm{sp}\le U_{\rm sp}$, $b>0$, $X>0$ and
$0<\mathsf R\le Xb$. Then $\mathrm{sp}\log(1+\mathsf R/\mathrm{sp})\le\log(1+X)\,(U_{\rm sp}+b)$, the left side being
$0$ for $\mathrm{sp}=0$.}

\emph{Proof.} Let $\mathrm{sp}>0$. For every real $c>0$ we have $\log c\le c-1$. Apply this to
$c=(1+\mathsf R/\mathrm{sp})/(1+X)$:
\begin{equation}\label{D:413}
 \mathrm{sp}\log\Big(1+\frac{\mathsf R}{\mathrm{sp}}\Big)\le\mathrm{sp}\log(1+X)+\frac{\mathrm{sp}+\mathsf R}{1+X}-\mathrm{sp}
 \le\mathrm{sp}\log(1+X)+\frac{Xb}{1+X}.
\end{equation}
And $X/(1+X)=1-1/(1+X)\le\log(1+X)$, by the same inequality for $c=1/(1+X)$. \qed

\noindent The lemma converts the term $4\,\mathrm{sp}\log(1+\mathsf R/\mathrm{sp})$ of (H8) into
$\log(1+x^{C_H})$ times a bound for the spread. Without a bound for the spread one has only
$\mathrm{sp}\log(1+\mathsf R/\mathrm{sp})\le\mathsf R=x^{C_H}b$, which is too large.

\textbf{Proposition M7 (one translate; proved here from (H0)--(H11)).} \emph{Assume
\textup{(H0)}, except its part \textup{(h)}, and \textup{(H1)--(H11)} for the translate $u$. Then}
\begin{equation}\label{as:Err}
 P_u-M_u\ \le\ \mathrm{Err}_u:=K_0+\log|\mathsf W_u|+\mathrm{LoP}+\mathrm{MLo}
 +8\log\big(1+x^{C_H}\big)\,(U_{\rm sp}+b)+12\,b ,
\end{equation}
\emph{where $\mathrm{LoP}$ is as in \eqref{as:H7} and}
\begin{align}\label{D:415}
 K_0&=c_4A\Big(L_a+\frac{m^2L_a}k\Big)+c_1|T|aL_a+c_1m(B_0+N_c)x^{C_{\rm abs}\varepsilon}+c_2mB_0+c_5|T|aL_a\Big(1+\frac mk\Big),\\
 \mathrm{MLo}&=c_6\,\mathrm{LoP}+c_6c_4\,|T|\,m^2\Big(\frac ak+b'\Big)x^{\varepsilon}+c_6X''+1,\\
 V_{\rm var}&=2c_2mB_0+c_3\big(|T|aL_a+mB_0\big)+c_4A\log k,\\
 U_{\rm sp}&=c_5|T|aL_a\Big(2+\frac mk\Big)+V_{\rm var}+c_6mB_0+\mathrm{MLo}.
\end{align}

\emph{Proof.} Write $G=G_{\rm top}$, $\mu_{\mathsf w}=\mu(\Lambda^{\rm gl}_{\mathsf w})$,
$\mu^{\max}_{\mathsf w}=\mu_{\max}(\Lambda^{\rm gl}_{\mathsf w})$, $E_{\mathsf w}=E_{\rm err}(\mathsf w)$,
$\ell_{\mathsf w}=\log\mathrm{Lv}(\mathsf w)$, and let $\overline{\mathrm{Lo}}$ be the mean of the numbers
$\mathrm{Lo}_{\mathsf w}$ with the weights $d_{\mathsf w}$. The set $\mathsf W_u$ is not empty by (H0)(g), and the weights
$d_{\mathsf w}$ are positive by (H11); so the means over the words that occur below are defined. And
$A\ge0$ by (H0)(c).

\emph{Step 1: a lower bound for $s_u$.} Fix a word. By (H6)(ii), and by (H1) solved for $-h_{\mathsf w}$,
\begin{equation}\label{D:416}
 \mu_{\mathsf w}-h_{\mathsf w}\ \ge\ \mu(\Lambda^0_{\mathsf w})+G+F_0-E_{\mathsf w}-\mathrm{Lo}_{\mathsf w}
 +\mu_{\rm CC}-F_0-S+A\ell_{\mathsf w}-c_1|T|aL_a-c_1m(B_0+N_c)x^{C_{\rm abs}\varepsilon}.
\end{equation}
The two terms $F_0$ cancel. Now use the lower bound of (H5) for $\mu(\Lambda^0_{\mathsf w})$, the upper bound
$E_{\mathsf w}\le mA\ell_g+c_2mB_0$ of (H2), and the lower bound of (H4), which reads
$\ell_{\mathsf w}\ge\Lambda_{\rm c}-c_4(L_a+m^2L_a/k)$ with $\Lambda_{\rm c}$ as in Lemma~M5 for
$\varrho=\mathrm{rd}_E/2$. Since $A\ge0$,
\begin{equation}\label{D:417}
 \mu_{\mathsf w}-h_{\mathsf w}\ \ge\ G+\mu_{\rm CC}-\big(S-A\Lambda_{\rm c}+mA\ell_g\big)-K_0-\mathrm{Lo}_{\mathsf w}
 =G+\mu_{\rm CC}-M_u-K_0-\mathrm{Lo}_{\mathsf w},
\end{equation}
by Lemma~M5. Taking the mean over the words with the weights $d_{\mathsf w}$,
\begin{equation}\label{as:step1}
 s_u\ \ge\ G+\mu_{\rm CC}-M_u-K_0-\overline{\mathrm{Lo}} .
\end{equation}

\emph{Step 2: the spread.} Let $\mathsf w,\mathsf w'$ be words. By (H6)(iii) for $\mathsf w$ and (H6)(ii) for
$\mathsf w'$, the terms $G+F_0$ cancel and
\begin{equation}\label{D:419}
 (\mu^{\max}_{\mathsf w}-h_{\mathsf w})-(\mu_{\mathsf w'}-h_{\mathsf w'})\le
 \big(\mu_{\max}(\Lambda^0_{\mathsf w})-\mu(\Lambda^0_{\mathsf w'})\big)+c_6mB_0+\mathrm{Lo}_{\mathsf w'}
 +\big(E_{\mathsf w'}-E_{\mathsf w}\big)+\big(h_{\mathsf w'}-h_{\mathsf w}\big).
\end{equation}
By (H5) the first bracket is at most $c_5|T|aL_a(2+m/k)$. By (H2), $E_{\mathsf w'}-E_{\mathsf w}\le2c_2mB_0$.
By (H3) and (H4),
$h_{\mathsf w'}-h_{\mathsf w}\le c_3(|T|aL_a+mB_0)+A(\ell_{\mathsf w}-\ell_{\mathsf w'})\le c_3(|T|aL_a+mB_0)+c_4A\log k$.
So the left side is at most $c_5|T|aL_a(2+m/k)+V_{\rm var}+c_6mB_0+\mathrm{Lo}_{\mathsf w'}$. Take the mean
over $\mathsf w'$ with the weights $d_{\mathsf w'}$, and then the maximum over $\mathsf w$. By
\eqref{as:s},
\begin{equation}\label{as:step2}
 0\le\mathrm{sp}_u\ \le\ c_5|T|aL_a\Big(2+\frac mk\Big)+V_{\rm var}+c_6mB_0+\overline{\mathrm{Lo}} .
\end{equation}

\emph{Step 3: the losses.} By the upper bound of (H4) and the second inequality of (H7),
$\mathrm{Lo}_{\mathsf w}\le c_6\mathrm{LoP}+c_6|T|m(a/k+b')\,c_4mx^{\varepsilon}+c_6X''+1=\mathrm{MLo}$ for
every word. Hence $\overline{\mathrm{Lo}}\le\mathrm{MLo}$, and $\mathrm{sp}_u\le U_{\rm sp}$ by \eqref{as:step2}.

\emph{Step 4: the mean-slope inequality.} Put $\mathsf R=x^{C_H}b>0$. By (H9), $s_u-\mu^{\min}_{\rm CC}\le\mathsf R$,
so \eqref{as:H8} holds for this $\mathsf R$. By (H10), $(1-\delta_{\rm p})^{-1}\le2$. By (H8),
$(\eta_u+4\sqrt{\mathsf e_u})\mathsf R\le6b$. By Lemma~M6 with $X=x^{C_H}$,
$\mathrm{sp}_u\log(1+\mathsf R/\mathrm{sp}_u)\le\log(1+x^{C_H})(U_{\rm sp}+b)$. So
\begin{equation}\label{D:421}
 s_u\ \le\ \mu_{\rm CC}+\log|\mathsf W_u|+8\log\big(1+x^{C_H}\big)(U_{\rm sp}+b)+12\,b .
\end{equation}
With \eqref{as:step1} the terms $\mu_{\rm CC}$ cancel:
\begin{equation}\label{D:422}
 G\ \le\ M_u+K_0+\overline{\mathrm{Lo}}+\log|\mathsf W_u|+8\log\big(1+x^{C_H}\big)(U_{\rm sp}+b)+12\,b .
\end{equation}

\emph{Step 5: the prime sum.} By the first inequality of (H7), $P_u\le G+\mathrm{LoP}$. With
$\overline{\mathrm{Lo}}\le\mathrm{MLo}$ this is \eqref{as:Err}. \qed

\noindent\emph{Which hypothesis was used where.} Step~1: (H1), (H2), (H4) lower bound, (H5) lower bound,
(H6)(ii). Step~2: (H2), (H3), (H4) third part, (H5), (H6)(ii) and (iii). Step~3: (H4) upper bound, (H7)
second part. Step~4: (H8), (H9), (H10). Step~5: (H7) first part. The means over the words exist by
(H0)(g) and (H11), and (H0)(c) gives the sign of $A$. So every part of (H1)--(H11) is used \emph{except}
(H6)(i): the proof never uses that $\beta_{\mathsf w}$ maps $\Lambda^{\rm gl}_{\mathsf w}$ into
$\mathcal R^\vee$. That statement, and the statement (8b) that $h_{\mathsf w}$ bounds the height of
$\beta_{\mathsf w}$, enter only through (H8): they are two of the three inputs from which Proposition~M4
gives (H8). Three smaller points. (H2) is used only in its summed form
$|E_{\rm err}(\mathsf w)-mA\ell_g|\le c_2mB_0$. That $\mathrm{sp}_u\ge0$ is part of Definition~M2
($\mu\le\mu_{\max}$), not a Hypothesis. And of (H0), besides (c) and (g), only $x\ge1$ and $b\ge1$ are
used in this proof (in Step~4). Steps~1, 4 and~5 give the inequality; Steps~2 and~3 only make its error
small.

\subsection*{6.\quad The size of the error; Theorem~M}

\textbf{Lemma M8 (the size of the error; proved here).} \emph{Assume \textup{(H0)}, except its part \textup{(h)}. Let $K\ge0$ be
a real number with $c_i\le K$ for $1\le i\le7$, and suppose}
\begin{equation}\label{as:big}
 x^{\varepsilon}\ \ge\ 79\,(1+K)^2\,(1+C_H)\,(1+e_K).
\end{equation}
\emph{Then, with $\mathfrak E$ as in \eqref{as:E},}
\begin{equation}\label{as:size}
 \mathrm{Err}_u\ \le\ x^{(C_{\rm abs}+e_L+3)\varepsilon}\ \mathfrak E .
\end{equation}

\emph{Proof.} Write $z=x^{\varepsilon}$, $z_D=x^{C_{\rm abs}\varepsilon}$, $z_L=x^{e_L\varepsilon}$; all are
$\ge1$. Write $\kappa_1=(1+K)^2$, $\kappa_2=1+C_H$, $\kappa_3=1+e_K$. We collect the facts that follow
from (H0). (1) $m\le x$, because $m\le x^{1/2+\varepsilon/2}\le x$; so $mb\le a$, and $m/k\le1$. (2)
$\mathfrak E\ge|T|a(1+m^2/k)\ge|T|a\ge a\ge b\ge1$, because $|T|\ge1$; and $mB_0\le\mathfrak E$,
$m(B_0+N_c)\le\mathfrak E$. (3) $A\le2|T|a$. (4) $\log k\le e_K\log x\le\kappa_3\,z$. (5)
$\log(1+x^{C_H})\le\log2+C_H\log x\le\kappa_2\,z$. (6) $\log|\mathsf W_u|\le k\le a\le\mathfrak E$. (7)
$|T|m^2(a/k+b')\le2|T|m^2a/k\le2\mathfrak E$, because $b'\le a/k$. (8) $X''\le a\le\mathfrak E$.
Now bound the terms of \eqref{as:Err}.
\begin{align}\label{D:425}
 K_0&\le4Kz\,|T|a(1+m^2/k)+2Kz\,|T|a+Kz_D\,m(B_0+N_c)+K\,mB_0+4Kz\,|T|a\\
 &\le(10K\,z\,z_D+K)\,\mathfrak E,\\
 \mathrm{LoP}&\le K\,z_L\,|T|\,a\le Kz_L\,\mathfrak E,\\
 \mathrm{MLo}&\le\big(K^2z_L+2K^2z+K+1\big)\mathfrak E,\\
 U_{\rm sp}-\mathrm{MLo}&\le6Kz\,\mathfrak E+\big(2K+K(2z+1)+2K\kappa_3\,z\big)\mathfrak E+K\mathfrak E\\
 &=\big(8Kz+4K+2K\kappa_3z\big)\mathfrak E .
\end{align}
Hence $K_0+\log|\mathsf W_u|+\mathrm{LoP}+\mathrm{MLo}\le(14K+3K^2+2)\,z\,z_D\,z_L\,\mathfrak E\le7\kappa_1\,z\,z_D\,z_L\,\mathfrak E$,
and $U_{\rm sp}+b\le(15K+3K^2+2)\,\kappa_3\,z\,z_L\,\mathfrak E\le7.5\,\kappa_1\,\kappa_3\,z\,z_L\,\mathfrak E$. So
\begin{equation}\label{D:426}
 \mathrm{Err}_u\le\big(7\kappa_1\,z\,z_Dz_L+8\kappa_2z\cdot7.5\,\kappa_1\,\kappa_3\,z\,z_L+12\big)\mathfrak E
 \le79\,\kappa_1\,\kappa_2\,\kappa_3\cdot z^2z_Dz_L\,\mathfrak E\le z^3z_Dz_L\,\mathfrak E,
\end{equation}
by \eqref{as:big}. And $z^3z_Dz_L=x^{(C_{\rm abs}+e_L+3)\varepsilon}$. \qed

\medskip
\noindent\textbf{Theorem M.} \emph{Fix $\varepsilon$, $x$ and a translate $u\in H$. Assume the Hypotheses
\textup{(H0)--(H11)} of \S4 for $u$, with constants $c_1,\dots,c_7\le K$ and $C_{\rm abs},e_L,C_H,e_K\ge0$, and
assume \eqref{as:big}. Then}
\begin{equation}\label{as:thmM}
 \sum_{t\in T}\sum_{j=1}^m\big(P_{u,t,j}-M_{u,t,j}\big)\ \le\ x^{(C_{\rm abs}+e_L+3)\varepsilon}
 \Big(|T|\,a\Big(1+\frac{m^2}k\Big)+m\,(B_0+N_c)\Big),
\end{equation}
\emph{and, with $B=|T|\,m\,b\,\delta/2$ and $\delta=|T_2|/n$,}
\begin{equation}\label{as:four}
 \frac1B\sum_{t\in T}\sum_{j=1}^m\big(P_{u,t,j}-M_{u,t,j}\big)\ \le\
 8\,x^{(C_{\rm abs}+e_L+3+\frac{17}8)\varepsilon}\Big(\frac xm+\frac{xm}k+m+\frac k{n^2}\Big).
\end{equation}
\emph{So the main inequality \textup{(M)}, in the form \eqref{as:M} for the translate $u$, follows
from \textup{(H0)--(H11)}.}

\emph{Proof.} \eqref{as:thmM} is Proposition~M7 with Lemma~M8. For \eqref{as:four}, by (H0)(f),
\begin{equation}\label{as:div}
 \frac{\mathfrak E}B=\frac2\delta\Big(\frac xm+\frac{xm}k\Big)+\frac{4n}{|T|\,\delta}\,m\,x^{\varepsilon}
 +\frac{2n}{|T|}\,x^{2\varepsilon}\,\frac k{n^2}:
\end{equation}
the first term is $|T|\,xb\,(1+m^2/k)$ divided by $|T|mb\delta/2$; the second is $m\cdot2nmbx^{\varepsilon}$
divided by the same; and for the third, $N_c=\delta\,k\,b\,x^{2\varepsilon}/n$. By (H0)(h),
$1/\delta\le x^{\varepsilon}$ and $n/|T|\le2x^{\varepsilon/8}$. So the three terms are at most
$2x^{\varepsilon}(x/m+xm/k)$, $8x^{17\varepsilon/8}m$ and $4x^{17\varepsilon/8}k/n^2$. \qed

\noindent The theorem will be used as the statement (M) of Parts~I and~II. It is an implication: it
reduces (M) to (H0)--(H11), and it does not prove (M). Of its hypotheses, the paper proves none
outright; the status of each is in the table of \S1. \emph{A limitation.} The theorem is a statement for
one scale $x$ and one set of data. Parts~I and~II need (M) for every $x$ in intervals $[Z,Z^{3/2}]$
with $Z$ arbitrarily large, each time with auxiliary primes $q$, $d$ for which Estimates~B and~C hold at
that $x$. That data satisfying the Hypotheses exist for these $x$, with such primes and with constants
that do not depend on $x$, is not proved here. It is Hypothesis (H12) of \S7.

\textbf{Corollary M9 (the mean over the translates).} \emph{Assume the hypotheses of Theorem~M for every
$u\in H$, with the same constants $K$, $C_{\rm abs}$, $e_L$, $C_H$, $e_K$ for all $u$. Assume also that
$\varepsilon$, $x$, $b$, $n$, $m$, $k$ and the sets $T$, $T_2$ are the same for all $u$, as the last
sentence of Hypothesis~\textup{O23} says. Then}
\begin{gather}\label{D:430}
 \frac1n\sum_{u\in H}\sum_{t\in T}\sum_{j=1}^m\big(P_{u,t,j}-M_{u,t,j}\big)\le\mathcal E,\qquad
 \mathcal E=x^{(C_{\rm abs}+e_L+3)\varepsilon}\,\mathfrak E,\\
 \frac{\mathcal E}B\le8\,x^{(C_{\rm abs}+e_L+3+\frac{17}8)\varepsilon}\Big(\frac xm+\frac{xm}k+m+\frac k{n^2}\Big).
\end{gather}
\emph{This is the display \textup{(M)} of Part~II, \S4.3, with $C=C_{\rm abs}+e_L+6$, because
$8\le x^{7\varepsilon/8}$ by \eqref{as:big}.}

\emph{Proof.} The right side of \eqref{as:thmM} is formed from $|T|$, $a=xb$, $m$, $k$, $B_0$, $N_c$ and
the constants, which are common to all translates; so it does not depend on $u$. Take the mean of
\eqref{as:thmM} over the $n$ translates. \qed

\noindent On the route of the final theorem $e_L=0$, so that $C=C_{\rm abs}+6$ (\S2, item (3)).

\noindent\emph{On the two means.} The step from one translate to the mean uses no further hypothesis. The
mean over the \emph{words}, weighted by the ranks $d_{\mathsf w}$, is inside the proof for one translate:
it is the number $s_u$, and it is forced by Proposition~M4, where the degrees of the sources are added.
After Step~4 of Proposition~M7 no word is left. The error is the same for every translate, except for the
term $\log|\mathsf W_u|$, which is bounded by $k$ in Lemma~M8. If one wants only the mean over $u$, it is
enough to assume (H7) on average over $u$; the source contains this variant, and the final theorem does
not use it (\S9).

\textbf{The four errors.} The table says which term of $\mathrm{Err}_u$ becomes which error of (M),
which Hypothesis it comes from, and its size for the parameters of \S\ref{ob:sec}
($n=x^{2+8\varepsilon}$, $m=x^{1/2-\varepsilon/16}$, $k=x^{9/2+10\varepsilon}$, $|T|=nx^{-\varepsilon/8}$,
$\delta=x^{-\varepsilon}$, up to bounded factors). The last column is the term divided by $B$, as a power
of $x$; constants, and the factors $x^{C_{\rm abs}\varepsilon}$, $x^{e_L\varepsilon}$, $\log(1+x^{C_H})$,
are left out.

{\small\setlength{\LTleft}{0pt}\setlength{\LTright}{0pt}
\begin{longtable}{@{}p{.25\textwidth}p{.34\textwidth}p{.12\textwidth}p{.20\textwidth}@{}}
\toprule
Term of $\mathrm{Err}_u$ & Hypothesis & Error of (M) & Size divided by $B$\\
\midrule\endhead
\bottomrule\endfoot
$|T|\,a\,L_a$ & (H1) first error; (H5); (H3), through the spread; (H4), $c_4AL_a$ & $x/m$ & $x^{1/2+33\varepsilon/16}$\\
$A\log k$ & (H4) third part, through the spread & $x/m$ & $x^{1/2+17\varepsilon/16}$\\
$|T|\,a\,L_a\,m^2/k$ & (H4): the defect $c_4m^2L_a/k$ of the count & $xm/k$ & $x^{-3-129\varepsilon/16}$\\
$|T|\,m^2(a/k+b')\,x^{\varepsilon}$ & (H6)(ii): the places above the level, with (H4) upper bound & $xm/k$ & $x^{-3-129\varepsilon/16}$\\
$m\,B_0$ & (H1) second error; (H2); (H3). (The term $c_6mB_0$ of (H6)(iii) also has this form, but by Lemma~M10 it can be taken to be $0$; it is kept only to match the source) & $m$ & $x^{1/2+33\varepsilon/16}$\\
$m\,N_c$ & (H1) second error only & $k/n^2$ & $x^{1/2-31\varepsilon/8}$\\
$\mathrm{LoP}=c_7|T|mbx^{e_L\varepsilon}$ & (H7) & (absorbed in $x/m$) & $x^{\varepsilon}$\\
$b$; $X''$; $\log|\mathsf W_u|$ & (H8) with (H9); (H6)(ii), rounding; (H8) & (absorbed in $x/m$) & $\le x^{-5/2}$\\
\end{longtable}}

\noindent For comparison, the four errors themselves are $x/m=x^{1/2+\varepsilon/16}$,
$xm/k=x^{-3-161\varepsilon/16}$, $m=x^{1/2-\varepsilon/16}$ and $k/n^2=x^{1/2-6\varepsilon}$. So three of
them are of size $x^{1/2+O(\varepsilon)}$ and the fourth is negligible, as Remark~(1) of \S\ref{ob:sec}, \S3
says. \emph{Which choice of parameters produces which error.} (i) $x/m$ is a cost of order $|T|\,a$
incurred once for the whole chain; it is divided by $B$, which grows with the number $m$ of lifts. (ii)
$m$ is a cost of order $b\,x^{\varepsilon}$ for each of the $2n$ places and each of the $m$ lifts, that
is $B_0$, incurred for each lift; so it is $mB_0$, and after division it grows with $m$. The choice
$m\approx\sqrt x$ balances (i) and (ii). (iii) $xm/k$ comes from the terms with a factor $m^2/k$, in (H4) and (H6); it is made negligible
by $k\ge xm$. Why the factor has this form is not visible in this section. (iv) $k/n^2$ is the term $N_c$ of the size bound, which is proportional to $k\,|T_2|/n^2$;
it forces $k\le n^2x^{1/2}$. Since the number of words requires $k\ge8C_{\mathsf G}m^5f\,x^{\varepsilon}$
(\S III.8, after Corollary~T2) and $f=2n$, this gives $n\ge16C_{\mathsf G}m^5x^{\varepsilon-1/2}$, about
$x^2$, and then $k$ about $x^{9/2}$. This explains the sizes of $n$ and $k$; the remark is ours.

\subsection*{7.\quad The scales and the primes: Hypothesis (H12) and the Riemann hypothesis}

Theorem~M is a statement about one scale $x$ and one set of data. Parts~I and~II need more. This
subsection says what they need, states it as one further Hypothesis, (H12), and proves that (H12)
implies the Riemann hypothesis (Corollary~M11). Nothing in it is used in the proofs above.

\emph{What Parts~I and~II need.} Recall from \S I.1 the function
$G(x)=\sum_{r\le x}\Lambda(r)(x/r-1)-x\log x+(1+\gamma)x$, where $\Lambda$ is von Mangoldt's function. The
criterion of Part~II, \S1 needs a bound $G(x)\le x^\theta$ for every $x$ of intervals $[Z,Z^{3/2}]$
with $Z$ arbitrarily large. An arbitrary unbounded set of $x$ would not do. For one $x$, Part~II,
\S3--5 obtain the bound for $G(x)$ from three things at that $x$: the inequality (M); Estimate~B, for
the three characters $\chi_q$, $\chi_d$, $\chi_q\chi_d$; and Estimate~C, for the characters of the field
$E$ of the data, which is formed with $q$, $d$ and $n$. The three must be for the same primes $q$, $d$
and the same $n$. But the argument is for each $x$ separately, and nothing ties the primes used at
different $x$. So the need is this: \emph{for each $x$ of the interval, some data and some pair $q$, $d$
for which \textup{(M)} and Estimates~B and~C hold at that $x$; and constants that do not depend on $x$,
so that one exponent serves on the whole interval.} Part~II, \S3 takes one pair for the whole interval.
That is more than its argument uses. Only small $\varepsilon$ matter: the conclusion lets $\varepsilon$
tend to $0$.

\medskip
\noindent\textbf{Hypothesis (H12) (supply on windows; assumed here, from the Lean development).}
\emph{There are real numbers $C_{\rm abs}\ge0$, $e_L\ge0$ and $\varepsilon_1$ with
$0<\varepsilon_1\le\frac1{52}$, and for every $\varepsilon$ with $0<\varepsilon<\varepsilon_1$ there are
real numbers $K\ge0$, $C_H\ge0$, $e_K\ge0$ and $J\ge1$, such that the following holds. For every real
number $y$ there is a real number $Z\ge y$ such that for every real $x$ with $Z\le x\le Z^{3/2}$ there
exist}
\begin{enumerate}[label=\textup{(\alph*)},leftmargin=2.2em,itemsep=1pt,topsep=2pt]
\item \emph{a real number $Y$ with $x^{1/J}\le Y\le x$;}
\item \emph{primes $q$ and $d$ that have the properties \textup{(a), (b), (c)} and \textup{(e)} of
Theorem~\ref{thm:pair} for $M=1$, for this $\varepsilon$ and for this $Y$, and the property
\textup{(d)} of that theorem for $M=1$ at all $y'$ with $Y\le y'\le x$: for $\chi=\chi_q$, $\chi_d$ and
$\chi_q\chi_d$,}
\begin{equation}\label{D:431}
 |G_\chi(y')|\le y'^{\,1/2+\varepsilon/2}\qquad\text{and}\qquad\Big|\frac{L'}{L}(1,\chi)\Big|\le4\log y'
 \qquad(Y\le y'\le x);
\end{equation}
\item \emph{data of the construction of \S\ref{ob:sec} for this $\varepsilon$, this $x$ and these $q$,
$d$: an integer $n=2^e$ with $e\ge2$ and $x^{2+8\varepsilon}\le n\le2x^{2+8\varepsilon}$; the field $E$
and the group $H$ of order $n$ formed with $q$, $d$ and $n$ \textup{(\S\ref{ob:sec}, \S1)}; and for
every translate $u\in H$ the objects listed at the start of \S4. The numbers $b$, $m$, $k$, the sets $T$
and $T_2$, and the orders $a_t$ are the same for all $u$; the set $T_2$ is regarded as a set of
columns, that is, as a subset of $H$;}
\item \emph{such that \textup{(H0)--(H11)} hold for every $u\in H$, with these $C_{\rm abs}$, $e_L$,
$C_H$, $e_K$ and with constants $c_1,\dots,c_7\le K$;}
\item \emph{and such that $b=\exp(x^{\varepsilon})$; $a_t>0$ and $|a_t/b-x|\le1$ for $t\in T$;
$m\ge x^{1/2-\varepsilon}$ and $k\le n^2x^{1/2-\varepsilon}$; and, for every $u\in H$ and every lift $j$,
$b\le w_{u,j}(1,s)\le b\,(1+1/(nx^2))$ for $s\in uT_2$, and
$a_t/w_{u,j}(v,s)\ge\frac13\exp(x^{\varepsilon/3})$ for every other cell $(v,s)$ and every $t\in T$.}
\end{enumerate}

\noindent\emph{How to read the quantifiers.} The numbers $C_{\rm abs}$, $e_L$, $\varepsilon_1$ come first
and depend on nothing. The numbers $K$, $C_H$, $e_K$, $J$ may depend on $\varepsilon$. None of these
depends on $x$, on the translate or on the word. The interval $[Z,Z^{3/2}]$ depends on $\varepsilon$ and
$y$. Everything in (a)--(e) may depend on $x$: the scale $Y$, the primes, $n$, the field and all the
data. The Hypothesis does not say that two points $x$ of one interval have the same primes. It does not
say that $Y$ is the left end $Z$ of the interval.

\noindent\emph{What ``data'' means.} In (c), ``data'' means the objects listed at the start of \S4, with
the types stated there, together with the sets $T,T_2\subset H$, the numbers $b$, $m$, $k$, $a_t$
($t\in T$) and the weights $w_{u,j}$ on the $2n$ cells. Corollary~M11 uses no property of them beyond
(H0)--(H11) and (e). In the source they are the objects of \S\ref{ob:sec}, whose existence is Hypotheses
O12, O19--O21, O23--O25 there; so (H12) contains those hypotheses for the scales of its intervals, to the
extent that the dictionary of \S\ref{ob:sec}, \S7 is correct.

\noindent\emph{What each clause is for.} Clauses (c) and (d) are the hypotheses of Corollary~M9 at the
scale $x$; they give (M) at $x$. Clause (d) is for every $u\in H$; so the translates over which (M)
takes its mean form the whole group $H$, which contains $T_2$. Clause (b) gives Estimate~B at $x$ (take
$y'=x$ in (d) of Theorem~\ref{thm:pair}), and it gives the hypotheses of Theorem~\ref{thm:C}, hence
Estimate~C at $x$. Clause (a) ties the sizes of $q$ and $d$ to $x$. The proof of Corollary~M11
uses only $1\le Y\le x$ from it; the lower bound $x^{1/J}\le Y$, and with it $J$, is what makes clause (b)
available from Theorem~\ref{thm:pair} (Status, (1)), and is not used in the proof. Likewise the lower bounds
in (a), (b) of Theorem~\ref{thm:pair}, the small primes in (c) and the height of the region in (e) are not
used there. Clause (e) lists the conditions
under which Part~II, \S4--5 pass from (M) to Estimate~A: the shape of the weights (Part~II, \S4.2), and
the two inequalities that make the four errors of (M) at most $x^{1/2+\varepsilon}$.

\noindent\emph{Status.} Assumed here, from the Lean development; \S9, Remark~(7) gives the statements of
the source to which it corresponds. Three remarks say how much of it the paper supplies.

(1) \emph{The pair.} Theorem~\ref{thm:pair}, proved in \S III.1 from four cited theorems, gives for
$M=1$, $A=J$ and every large $Y$ a pair $q$, $d$ with the properties (a)--(e) there. Such a pair has the
properties of clause (b) for every $x$ with $Y\le x\le Y^J$, because then $[Y,x]\subset[Y,Y^J]$. And
$Y\le x\le Y^J$ is clause (a). So the paper proves that pairs as in (b) exist, for every large scale $Y$
and every $x$ in $[Y,Y^J]$. \S III.2 proves that Estimate~C holds for such a pair
(Theorem~\ref{thm:C}; this is used in the proof of Corollary~M11 and is not part of the Hypothesis).

(2) \emph{What is not supplied.} First, that the construction can be carried out at $x$ with such a
pair: that the primes $q$, $d$ of the data at $x$ are a pair as in (b). Hypothesis~O23 asserts data at
its admissible scales and says of their primes what (a)--(c) of Theorem~\ref{thm:pair} say, at a scale
$Y$ with $x^{1/J}\le Y\le x$; it does not say (d) or (e). The construction also asks more of the primes
than Theorem~\ref{thm:pair} gives, for instance $q,d>4t_0r+1$ (table of O23). Second, that these data
satisfy (H0)--(H11); the table of \S1 says what is known for each. Third, the uniformity: that $K$,
$C_H$, $e_K$, $C_{\rm abs}$, $e_L$ do not depend on $x$, and that the scales $x$ at which all this holds
fill intervals $[Z,Z^{3/2}]$ with $Z$ arbitrarily large. O23 speaks of an unbounded set of admissible
scales and does not say that it contains intervals.

(3) \emph{Clause \textup{(e)} is not a new assumption on data as in Hypothesis~O23.} It follows from the
conditions of the table of O23 for $x$ large in terms of $\varepsilon$. Indeed $b=\exp(x^\varepsilon)$
is in the table. From $a\le a_t\le a(1+1/(kx^{13}))$ and $a=xb$ we get $0\le a_t/b-x\le1$. From
$m\ge\sqrt{x\delta_1}=x^{1/2-\varepsilon/16}$ we get $m\ge x^{1/2-\varepsilon}$. From
$k\le2x^{9/2+10\varepsilon}$ and $n\ge x^{2+8\varepsilon}$ we get $k\le n^2x^{1/2-\varepsilon}$ as soon as
$x^{5\varepsilon}\ge2$. A thick cell has $b\le w\le b+\eta+1$ with $6f(x+1)(\eta+1)\le b/x$ and $f=2n$, so
$\eta+1\le b/(nx^2)$. A top cell has $w\le b_t+b'/x+\eta\le3b'$ and a thin cell has
$w\le w''+\eta\le b'$, for large $x$; and $a_t/(3b')\ge a/(3b')=\frac13X'=\frac13\exp(x^{\varepsilon/3})$.

\medskip
\noindent\textbf{Corollary M11 (the Riemann hypothesis from the Hypotheses; proved here from (H12),
Theorem~M, Part~II and Theorem~\ref{thm:C}).} \emph{Assume \textup{(H12)}; it contains
\textup{(H0)--(H11)} for the data at each scale. Put $C_*=C_{\rm abs}+e_L+8$. Then for every
$\varepsilon$ with $0<\varepsilon<\varepsilon_1$ and every real number $y$ there is $Z\ge y$ such that}
\begin{equation}\label{as:Gbound}
 G(x)\ \le\ x^{1/2+C_*\varepsilon}\qquad\text{for all }x\text{ with }Z\le x\le Z^{3/2}.
\end{equation}
\emph{Consequently $\zeta(s)\ne0$ for $\re s>\frac12$: the Riemann hypothesis holds.}

\emph{Proof.} Fix $\varepsilon$ with $0<\varepsilon<\varepsilon_1$, and let $K$, $C_H$, $e_K$, $J$ be as
in (H12). Below, ``for large $x$'' means: for $x\ge x_1(\varepsilon)$, where $x_1(\varepsilon)$ depends
on $\varepsilon$, $K$, $C_H$, $e_K$ and on nothing else. There are finitely many such conditions. Given
$y$, apply (H12) with $\max(y,x_1(\varepsilon))$ in place of $y$. This gives $Z$. Fix $x$ with
$Z\le x\le Z^{3/2}$, and let $Y$, $q$, $d$, $n$ and the data be as in (a)--(e).

\emph{Step 1: \textup{(M)} at $x$.} By (c) and (d) the hypotheses of Corollary~M9 hold, and
\eqref{as:big} holds for large $x$. So, with $B=|T|mb\delta/2$,
\begin{equation}\label{D:433}
 \frac1n\sum_{u\in H}\sum_{t\in T}\sum_{j=1}^m\big(P_{u,t,j}-M_{u,t,j}\big)\le\mathcal E,\qquad
 \frac{\mathcal E}B\le x^{(C_{\rm abs}+e_L+6)\varepsilon}\Big(\frac xm+\frac{xm}k+m+\frac k{n^2}\Big).
\end{equation}
By (e), $x/m\le x^{1/2+\varepsilon}$ and $k/n^2\le x^{1/2-\varepsilon}$. By (H0)(b),
$m\le x^{1/2+\varepsilon/2}$ and $xm/k\le1$. So
\begin{equation}\label{as:EB}
 \mathcal E/B\ \le\ 4\,x^{1/2+(C_{\rm abs}+e_L+7)\varepsilon}.
\end{equation}

\emph{Step 2: Estimate~B at $x$.} The primes $q\equiv3$ and $d\equiv7\bmod8$ are odd and different. So
$\chi_q$, $\chi_d$ and $\eta=\chi_q\chi_d$ are primitive characters, different from the trivial one. By
(b) with $y'=x$,
\begin{equation}\label{as:estB}
 |G_\chi(x)|\le x^{1/2+\varepsilon/2}\qquad\text{and}\qquad|c_\chi|\le4\log x\qquad(\chi=\chi_q,\ \chi_d,\ \eta),
\end{equation}
where $c_\chi=\frac{L'}L(1,\chi)$. This is Estimate~B (statement 2 of Part~II, \S3) at $x$.

\emph{Step 3: Estimate~C at $x$.} Apply Theorem~\ref{thm:C} with $Q=4n$, with the bound $1$ for the odd
part of $Q$, and with $c'=1$. Its hypotheses hold. The number $Q=2^{e+2}$ is divisible by $8$ and prime
to $qd$, and its odd part is $1$. By (a), (b) of Theorem~\ref{thm:pair} and $1\le Y\le x$,
$\log q\le x^{\varepsilon/5}$ and $\log d\le x^{\varepsilon/5}(\log x)^2$; and
$\log(4n)\le\log8+(2+8\varepsilon)\log x$. So $\log(Qqd)\le2x^{\varepsilon/5}(\log x)^2$ for large $x$.
Let $\chi\ne1$ be a real primitive character whose conductor divides $Qqd$ and is divisible by $q$ or
by $d$. Then $\chi=\rho\,\chi_q^i\chi_d^j$ with $(i,j)\ne(0,0)$ and $\rho$ a real primitive character
whose conductor is a power of $2$, hence $1$, $4$ or $8$. By (e) of Theorem~\ref{thm:pair} with $M=1$,
$L(\beta,\chi)\ne0$ for real $\beta>1-\frac1{400}$. So Theorem~\ref{thm:C} applies, for large $x$: for
every primitive character $\chi$ whose conductor divides $4nqd$, the trivial one included with $G_1=G$,
\begin{equation}\label{as:estC}
 |G_\chi(y')|\le y'\exp\big(-x^{\varepsilon/8}\big)\qquad\text{for all }y'\ge\tfrac13\exp\big(x^{\varepsilon/3}\big).
\end{equation}
This is \eqref{II:t5.6} of Part~II, with $\kappa=1$. The characters of the cell expansion of Part~II, \S5.2 and
\S5.4 are the primitive characters that induce $\xi\chi_q^i\chi_d^j$ with $\xi$ a character of $H$; their
conductors divide $4nqd$.

\emph{Step 4: Part~II, \S5 at $x$.} We go through \S5 of Part~II and say where each of its hypotheses
comes from. The sums $P_{u,t,j}$ and $M_{u,t,j}$ of Definition~M1 are \eqref{II:t4.1} and \eqref{II:t4.2} of Part~II.
\begin{itemize}[leftmargin=1.6em,itemsep=1pt,topsep=2pt]
\item \emph{\textup{\eqref{II:t5.1}--\eqref{II:t5.3}}} hold for any positive weights and any $a_t>0$. They use the prime
number theorem for progressions to the fixed modulus $4nqd$, which is cited in Part~II. The weights are
positive by (H0)(e), and $a_t>0$ by (e). The result is
$P_{u,t,j}-M_{u,t,j}=\sum_iw_iR_i(a_t/w_i)$, the sum over the $2n$ cells.
\item \emph{\textup{\eqref{II:t5.4}, \eqref{II:t5.5}}} need the mean over all $u\in H$, and that the distinguished cells of
the translate $u$ are the cells $(1,s)$ with $s\in uT_2$. The first is in (c) and (d). The second is
the way (e) names the cells.
\item \emph{The other cells} (Part~II, \S5.4, first part) need \eqref{as:estC} at $y'=a_t/w_i$ and
$a_t\le(x+1)b$. By (e) the argument is at least $\frac13\exp(x^{\varepsilon/3})$, and $a_t\le(x+1)b$.
Each term is then at most $2a_t\exp(-x^{\varepsilon/8})$. There are $n\,|T|\,m$ triples $(u,t,j)$ and at
most $2n$ cells. After division by $nB$ the total is at most
$8n(x+1)\delta^{-1}\exp(-x^{\varepsilon/8})$. By (c) and (H0)(h), $n\le2x^{2+8\varepsilon}$ and
$\delta^{-1}\le x^\varepsilon$. So the total is at most $1$ for large $x$.
\item \emph{The distinguished cells} (\S5.4, second part) need $w_i=b+\Delta_i$ with
$0\le\Delta_i\le b/(nx^2)$. This is in (e). Part~II uses Chebyshev's bound, and the total after division
by $nB$ is at most a constant times $1/x$.
\item \emph{The argument} (\S5.4, third part) needs $|a_t/b-x|\le1$ and $|c_\chi|\le4\log x$ for the
three characters. These are in (e) and in \eqref{as:estB}. The change is at most a constant times
$1+\log x$.
\end{itemize}
The constants in the last two items are absolute. So \eqref{II:t5.7} of Part~II holds at $x$ in the form
\begin{equation}\label{D:437}
 \frac1n\sum_{u\in H}\sum_{t\in T}\sum_{j=1}^m\big(P_{u,t,j}-M_{u,t,j}\big)\ \ge\ B\,R_{q,d}(x)-C_\varepsilon\,B\,(1+\log x),
\end{equation}
where $R_{q,d}$ is the function of Part~II, \S3 and $C_\varepsilon$ depends on $\varepsilon$ only. With
Step~1 and \eqref{as:EB},
\begin{equation}\label{D:438}
 R_{q,d}(x)\ \le\ 4\,x^{1/2+(C_{\rm abs}+e_L+7)\varepsilon}+C_\varepsilon(1+\log x).
\end{equation}
This is Estimate~A at $x$.

\emph{Step 5: the bound for $G(x)$.} The identity of Part~II, \S3 is
$G(x)=R_{q,d}(x)-G_{\chi_q}(x)+G_{\chi_d}(x)+G_\eta(x)$. By Step~4 and \eqref{as:estB},
\begin{equation}\label{D:439}
 G(x)\ \le\ 4\,x^{1/2+(C_{\rm abs}+e_L+7)\varepsilon}+C_\varepsilon(1+\log x)+3\,x^{1/2+\varepsilon/2}
 \ \le\ x^{1/2+(C_{\rm abs}+e_L+8)\varepsilon}
\end{equation}
for large $x$, because $x^\varepsilon\ge8+C_\varepsilon(1+\log x)$ for large $x$.

\emph{Step 6: one exponent on the whole interval.} Every condition ``for large $x$'' above holds for
$x\ge x_1(\varepsilon)$, and $x_1(\varepsilon)$ does not depend on $x$, $Y$, $q$, $d$, $n$ or the data:
this is where the uniformity of the constants in (H12) is used. Since $Z\ge x_1(\varepsilon)$, Step~5
holds for every $x$ of $[Z,Z^{3/2}]$. This is \eqref{as:Gbound}.

\emph{Step 7: the Riemann hypothesis.} Let $\frac12<\theta<1$. Choose $\varepsilon$ with
$0<\varepsilon<\varepsilon_1$ and $\frac12+C_*\varepsilon\le\theta$. For $x\ge1$, \eqref{as:Gbound} gives
$G(x)\le x^\theta$ on intervals $[Z,Z^{3/2}]$ with $Z$ arbitrarily large. By the Lemma of Part~II, \S1
(the one-sided criterion on intervals), $\zeta$ has no zero in $\re s>\theta$. Since $\theta>\frac12$ is
arbitrary, $\zeta$ has no zero in $\re s>\frac12$. By the functional equation the zeros in
$0<\re s<1$ are symmetric about the line $\re s=\frac12$. So all of them lie on it. \qed

\noindent\emph{What the proof uses.} (i) Hypothesis (H12). (ii) Theorem~M and Corollary~M9, proved in
this section; they do not use (C6) or Proposition~M4, which serve only the status of (H8). (iii)
Part~II: the Lemma of \S1, the identity of \S3, and \S5; they are proved there, from the prime number
theorem for progressions to a fixed modulus, Chebyshev's bound, the continuation of $\zeta$ with its
expansion at $s=1$, and quadratic reciprocity (\S4.1). (iv) Theorem~\ref{thm:C}, proved in
\S III.2 from the cited theorems (Z1)--(Z3), the functional equation of $L(s,\chi)$ and Dirichlet's theorem
$L(1,\chi)\ne0$. (v) The functional equation of $\zeta$, in Step~7. Theorem~\ref{thm:pair} itself is not used: its conclusions
for the primes of the data are clause (b) of (H12). So in Corollary~M11 Estimate~B is an assumption about
the primes of the data. What \S III.1 adds is that primes with these properties exist. We found no step
of Part~II, \S3--5 that needs more than (H12) and the cited theorems. The conditions of Part~II, \S4.2
on the data that are not among (H0)--(H11) are exactly clause (e) and the bounds for $n$ in clause (c).

\noindent\emph{The exponent.} The source states the bound with $C_{\rm abs}+e_L+7$ in the place of
$C_*$. The $8$ here is not optimised.

\emph{What the statements of the source say} (the names are in \S9, Remark~(7)). The source has a
statement of the shape of (H12), and its quantifiers are as in (H12). Two exponents, our $C_{\rm abs}$
and $e_L$, and a bound $\varepsilon_1$ come first. For every $\varepsilon<\varepsilon_1$ there is a
family of data with a set of admissible scales, and there are $K$, $C_H$, $e_K$, a fourth constant $C_L$
that (H12) does not need, and $x_0$, such that for
every admissible $x\ge x_0$ the hypotheses of the source's form of Theorem~M hold for all translates.
The family is required to have \emph{windows with exponent $\frac32$}: for every $y$ there is $Z\ge y$
such that every $x$ with $Z\le x\le Z^{3/2}$ is admissible. A scale $x$ is admissible when there are a
scale $Y$ with $x^{1/J}\le Y\le x$, primes $q$, $d$ and an exponent $e$ with four properties: the
properties (a)--(c) of Theorem~\ref{thm:pair} at the scale $Y$; the property (d) for all $y'$ with
$Y\le y'\le x$; the property (e); and $x^{2+8\varepsilon}\le2^e\le2x^{2+8\varepsilon}$, together with the
bound \eqref{II:shabound} for Galois cohomology for the field formed with $q$, $d$, $e$. The source
proves that such windows exist from two inputs: its form of Theorem~\ref{thm:pair} with $A=J$, and the
bound for Galois cohomology on windows inside $[X^{1/J},X]$ for large $X$, which is where $J$ comes from.
In that proof one pair is chosen for a whole window, at the scale $Y=X^{1/J}$, with $X$ above the window. So $Y$ is
not the left end of the window, and need not be. The source then builds the field of the data at $x$
from the primes $q$, $d$ attached to $x$, and derives the bounds \eqref{as:estB} that its passage from
(M) to $G(x)$ uses from the property (d) at $y'=x$. In the statement that passes from (M) to $G(x)$,
the bound of (M) and the character estimates are taken at the same $x$ for the same data, hence for the
same $q$, $d$. All primes at one $x$ are the same for all translates; nothing ties the primes of
different $x$. The conclusion of the source is $G(x)\le x^{1/2+(C_{\rm abs}+e_L+7)\varepsilon}$ for the
admissible $x\ge x_1$, and then the criterion on windows. We read these statements, and the proofs named
in \S9, Remark~(7). We did not read how the source obtains its form of Estimate~C from property (e), nor
the proofs behind the source's form of Part~II, \S5. That the data of the source are the objects of
\S\ref{ob:sec} is the dictionary of \S\ref{ob:sec}, \S7, with the reservations made there.

\emph{Summary.} Theorem~M reduces (M) at one sufficiently large scale to (H0)--(H11). Hypothesis (H12) says that data
satisfying (H0)--(H11) exist at every $x$ of intervals $[Z,Z^{3/2}]$ with $Z$ arbitrarily large, with
constants that do not depend on $x$ and with primes for which Estimates~B and~C hold at $x$.
Corollary~M11 deduces the Riemann hypothesis from (H12). So the Riemann hypothesis is reduced to the
Hypotheses (H0)--(H12) and the cited theorems. None of the Hypotheses is proved outright in this paper,
and (H12) is assumed; of (H12) the paper proves that primes as in its clause (b) exist.

\subsection*{8.\quad Numerical checks}

The scripts are in \texttt{outline/examples/assembly/}. They test the arithmetic of this section on
made-up numbers. They do not test the Hypotheses for the construction.

\emph{Proposition M7} (\texttt{m1\_bookkeeping.py}). Random data were drawn so that (H1)--(H11) hold: the
free numbers (sizes, constants, weights, ranks, levels) at random, and then $h_{\mathsf w}$,
$\mu(\Lambda^0_{\mathsf w})$, $E_{\rm err}(\mathsf w)$, $\mu(\Lambda^{\rm gl}_{\mathsf w})$,
$\mu_{\max}(\Lambda^{\rm gl}_{\mathsf w})$, $G_{\rm top}$ and $P_u$ from the Hypotheses. In $20000$ cases
all the inequalities of (H1), (H2), (H5)--(H8) were equalities (``the extremes''), in $20000$ they had
random slack. The conclusion \eqref{as:Err} holds in all $40000$; the smallest value of
$(M_u+\mathrm{Err}_u-P_u)/\mathrm{Err}_u$ is $0.36$. At the extremes the inequality obtained before
Lemma~M6 and Step~3 are applied, $P_u\le M_u+K_0+\log|\mathsf W_u|+\mathrm{LoP}+\overline{\mathrm{Lo}}+{}$the
error of \eqref{as:H8}, is an equality up to $10^{-15}\,\mathrm{Err}_u$. So Steps~1, 4, 5 lose nothing,
and all the loss is in the bound for the spread. \emph{Controls.} One Hypothesis was violated by a
multiple of $\mathrm{Err}_u$ and the others kept ($2000$ cases each). With a violation of
$3\,\mathrm{Err}_u$ in (H1), (H2), (H4) (lower bound), (H5), (H6), (H7) or (H8) the conclusion fails in
all cases; so it does for (H9), when $\mathsf R$ is increased by $3\,\mathrm{Err}_u/(\eta_u+4\sqrt{\mathsf e_u})$,
and for (H10) with $\delta_{\rm p}$ near $1$. For (H3) it fails in about one case in five. This is as it
should be: (H3) enters only through the spread, and by the remark after Lemma~M6 a large spread costs at
most $8x^{C_H}b$.

\emph{Lemma M8 and the division} (\texttt{m2\_error.py}). On $400000$ random data satisfying (H0), half
of them with every bound of (H0) attained and with equality in \eqref{as:big}, the ratio
$\mathrm{Err}_u/(x^{(C_{\rm abs}+e_L+3)\varepsilon}\mathfrak E)$ is at most $0.063$. With the exponent
$C_{\rm abs}+e_L+2$ in place of $C_{\rm abs}+e_L+3$ the ratio reaches $1.4\cdot10^5$. This happens on the
data in which $\log k$ and $\log(1+x^{C_H})$ have the largest values that (H0) allows,
$(1+e_K)x^{\varepsilon}$ and $(1+C_H)x^{\varepsilon}$; that is, when $\log x$ is as large as
$x^{\varepsilon}$. There $\mathrm{Err}_u$ contains the product $8\log(1+x^{C_H})\cdot c_4A\log k$, of size
$K(1+C_H)(1+e_K)\,x^{2\varepsilon}\,|T|\,a$. So the third power of $x^{\varepsilon}$ pays for these two
logarithms when $\log x$ is as large as (H0)(a) allows, and for the constants. The scripts do not test
separately the case in which $\log x$ is small compared with $x^{\varepsilon}$; there a smaller exponent
may suffice. The identity \eqref{as:div} and the bound
\eqref{as:four} were checked on $10^5$ random data; the largest ratio of the two sides of the last step found is $1.000$.

\emph{Proposition M4} (\texttt{m3\_slopes.py}). Lattices over $\Z$ in place of $\mathfrak o_K$: a target
$Q$ of rank $2$ in $\mathbb R^2$, between one and six sources of rank $1$ or $2$ with integral maps into
$Q$, heights equal to the logarithms of the squared operator norms; the filtration of $Q$ computed from
its shortest vector. In $20000$ configurations, with $5389$ levels $\vartheta>0$ at which the filtration is not
trivial: \eqref{as:count} holds in all; the images of the spaces $V^\circ_{\mathsf w}$ lie in the step of
the filtration in all; and the inequality of Step~4 holds in all. \emph{Controls.} With the heights
lowered by $2$, \eqref{as:count} fails in $4472$ configurations. With $\mu$ in place of $\mu_{\max}$ in
the spread it fails in $42$.

\emph{Lemma M10} (\texttt{m10\_check.py}). A model over $\Z$, with the primes $2,3,5,7$: a lattice
$\Lambda^{\rm src}$ of rank $2$ in $\mathbb R^2$, integers $\tau(p)\in\{0,1,2\}$, a sublattice
$\Lambda^{\rm div}$ of $\prod_pp^{\tau(p)}\Lambda^{\rm src}$, real numbers $\mathrm{sg}(p)\in[-2,4]$, and
$\Lambda^{\rm gl}=\prod_pp^{-\lfloor\mathrm{sg}(p)\rfloor}\Lambda^{\rm div}$; the largest slope is computed
from the shortest vector. In $20000$ cases the inequality
$\mu_{\max}(\Lambda^{\rm gl})\le\mu_{\max}(\Lambda^{\rm src})+\sum_p2\log p\,(\mathrm{sg}(p)-\tau(p))$ of
the proof holds in all, and it is an equality in all $5000$ cases with
$\Lambda^{\rm div}=\prod_pp^{\tau(p)}\Lambda^{\rm src}$ and integral $\mathrm{sg}(p)$. \emph{Control.}
With $2\tau(p)$ in place of $\tau(p)$ it fails in $13536$ cases.

\emph{Lemmas M5 and M6, and the sizes} (\texttt{m4\_identity.py}). The identity \eqref{as:cancel} was
verified symbolically for $m=3$, $f=4$, $|T|=2$. Over $10^6$ random data the ratio of the two sides of
Lemma~M6 is at most $0.955$. The last column of the table of \S6 is computed by this script.

\subsection*{9.\quad Correspondence with the Lean development}

{\footnotesize\sloppy Files, under \path{lean/src/}: \texttt{MB} is
\path{Formalism/RankGapBounds/MeanBoundSecondForm.lean}; \texttt{QS} is
\path{Formalism/SizeWindowBrackets/QuotientSlopeBoundsWith.lean}; \texttt{LC} is
\path{Formalism/SizeWindowBrackets/ChainSlopeWindows/LevelOfCongruence.lean}; \texttt{WM} is
\path{Formalism/RankGapBounds/RankGapProofs/WeightedMeanData.lean}; \texttt{EF} is
\path{Formalism/SizeWindowBrackets/ErrorFormsWeakWindowsFalse.lean}; \texttt{PC} is
\path{Formalism/SizeWindowBrackets/ChainSlopeWindows/PChain.lean}; \texttt{SC} is
\path{Formalism/SizeWindowBrackets/SChain.lean}; \texttt{FPNS} is
\path{Formalism/SizeWindowBrackets/FirstPrimeNeSecondPrime.lean}; \texttt{ChTop}, \texttt{ChCore} are in
\path{Development/AutomorphicArithmetic/DedekindLFunctions/LandauTheorems/}; \texttt{SI} is
\path{Development/AutomorphicArithmetic/DedekindLFunctions/SplitInertia.lean}; \texttt{CoreGm},
\texttt{CoreAtH}, \texttt{LGm} (\path{LevelGmLeOnFamilyScales.lean}) are in
\path{Development/AutomorphicArithmetic/PrimitiveCosets/}; \texttt{QLDC}, \texttt{TieGlobal}, \texttt{LSS}
(\path{LossesSmallShares.lean}) are in \path{Formalism/SharpRestriction/GlobalSchurBounds/}; \texttt{GTI} is
\path{Formalism/SharpRestriction/GlobalTheoremInputs.lean}; \texttt{DP2} is
\path{Formalism/LocalRingCardinality/DataParams₂.lean}; \texttt{HN} is the folder
\path{Formalism/HarderNarasimhan/}.\par}

The last column: \emph{yes}, the declaration is in the lists of the dependency computation as one the
final theorem \texttt{riemannHypothesis} depends on; \emph{ref}, it is not in the lists but occurs in the
statement or in the proof term of a declaration that is: (a) \texttt{ch\_top}, (b)
\texttt{each\_unit\_bound\_statement}, (c) \texttt{linkCore}, (d) \texttt{linkGLOBAL7\_items}, (e)
\texttt{linkDIV\_at\_words}; \emph{no}, it is in the lists as not used; \emph{?}, not determined. To be
in the lists as used is necessary for a statement to be consulted by the proof of the final theorem; it
is not sufficient. The lists are by the last component of a name; where several declarations share it,
the one meant is the one in the file given.

{\footnotesize\setlength{\LTleft}{0pt}\setlength{\LTright}{0pt}
\begin{longtable}{@{}p{.30\textwidth}p{.535\textwidth}p{.10\textwidth}@{}}
\toprule
This section & Declarations in the source (file:line) & Dep.\\
\midrule\endhead
\bottomrule\endfoot
$P_{u,t,j}$, $M_{u,t,j}$ & \texttt{characterSumData.ASK}, \texttt{.MAIN} (FPNS:84, 89); the sum \texttt{ASK} is cut at a bound $N\ge a_t/w$ & yes\\
$P_u-M_u$ & \texttt{gapA} (SI:50) & yes\\
$\mathfrak E$ & \texttt{ErrText} (EF:16); \texttt{errorBracketSizeWindows} (EF:133) & ref (a); yes\\
the abstract numbers: slopes, heights, levels of one translate & \texttt{Data}, \texttt{Consts} (MB:37, 15); the fields \texttt{d}, \texttt{T}, \texttt{muN}, \texttt{muMaxN}, \texttt{h}, \texttt{muM}, \texttt{muMaxM}, \texttt{EeJ}, \texttt{muCC}, \texttt{muMinCC}, \texttt{Fo}, \texttt{Ga}, \texttt{ASK}, \texttt{IO}, \texttt{HH}, \texttt{del} are $d_{\mathsf w}$, $\mathrm{Lv}$, $\mu(\Lambda^{\rm gl})$, $\mu_{\max}(\Lambda^{\rm gl})$, $h_{\mathsf w}$, $\mu(\Lambda^0)$, $\mu_{\max}(\Lambda^0)$, $E_{{\rm err},j}$, $\mu_{\rm CC}$, $\mu^{\min}_{\rm CC}$, $F_0$, $G_{\rm top}$, $P_u$, $I_u$, $\mathsf R$, $\delta_{\rm p}$ & ref (b)\\
what the abstract numbers are for the records & \texttt{TieLocal}, \texttt{TieGlobal}, \texttt{TieArith} in \texttt{ChainTied} (TieGlobal:50--97); \texttt{Mod.data} (QLDC:464) & yes: \texttt{TieLocal}, \texttt{ChainTied}; others ref (c)\\
$s_u$, $\mathrm{sp}_u$ & \texttt{Data.s}, \texttt{Data.uF} (MB:92, 94); \texttt{sFrame}, \texttt{uFrame} (\path{DataArith.lean}:131, 134) & ref (b); yes\\
(H0) & \texttt{ErrWindows0}, \texttt{ErrWindows} (EF:28, 43); \texttt{BNBoundsU} (SI:74); \texttt{RestSize} (SC:53); \texttt{Win\_at} (QLDC:231); \texttt{HidaSetting} (MB:111) & yes; the first ref (a)\\
(H0)(g), $1\le|\mathsf W_u|$ & the instance hypothesis \texttt{Nonempty (W u0)}, for all \texttt{u0}, of \texttt{ch\_top} (ChTop:37) & ref (a)\\
(H1), (H2), (H3) & \texttt{Data.DIV} (MB:118), \texttt{EComp5} (QS:24), \texttt{Data.HComp} (MB:126); clauses 7, 8, 9 of \texttt{LinkDIV\_at} (GTI:77), supplied by \texttt{linkDIV\_at\_words} (\path{RecBridge.lean}:327); \texttt{SizeData}$_2$\texttt{.MAIN} (DP2:232) is $S-A\log\mathrm{Lv}$ & yes; \texttt{Data.DIV} ref (a)\\
the number $h_{\mathsf w}$ & \texttt{hExpl6} (\path{KernelPlaceMeanBound.lean}:25), by the hypothesis \texttt{hh} of \texttt{linkDIV\_at\_words}; (H3) is \texttt{Compl3} (\path{RecBridge.lean}:315) & ref (e)\\
(H4), (H5) & \texttt{Count5} (QS:27), \texttt{Data.Slopes} (MB:139); \texttt{Count\_at}, \texttt{SlopesM\_at} (QLDC:167, 173) & yes\\
(H6)(ii),(iii) & \texttt{QuotientSlopeBoundsWith}, \texttt{Lo5} (QS:19, 15); \texttt{SlopesN\_at}, \texttt{LoU} (QLDC:181, 177) & yes; \texttt{Lo5}, \texttt{LoU} ref (a), (c)\\
Lemma M10 & the first clause of \texttt{SlopesN\_at}, proved by \texttt{slopesN\_at\_of} (\path{TopDensityNGeSixteen.lean}:1529) from \texttt{Places\_at} (:1516) with \texttt{claimHigh\_valid} (\path{WeightedSlopeExcessDecomposition.lean}:191); its hypothesis \texttt{HypHigh} (\path{ClaimGaFree.lean}:94) is filled in \texttt{places\_at\_of} (\path{PlacesP.lean}:297) by \texttt{muMax\_le} (:171). Statement (c) is the third clause of \texttt{LocalQuadraticParts.adm26} (\path{LocalQuadraticParts.lean}:626). The inclusion $\Lambda^{\rm div}\subset\Lambda^{\rm src}$ is the fourth clause of \texttt{adm26} at the level places and a clause of \texttt{NuDef} elsewhere. In the lists as used: \texttt{places\_at\_of} and \texttt{adm26}. Ref (d): \texttt{slopesN\_at\_of}. In no list: \texttt{muMax\_le}, \texttt{claimHigh\_valid}, \texttt{Places\_at}, \texttt{HypHigh} & yes, ref or ?: see left\\
(H6)(i); (8b) & clauses 1 and 6 of \texttt{LinkDIV\_at}; (8b) and (H1) are consecutive parts of \texttt{WordFacts} (\path{RecBridge.lean}:128) & yes\\
(H7) & \texttt{LossBounds} (QS:34), \texttt{Data.LoP} (MB:105); \texttt{Loss\_at} (QLDC:203); proved by \texttt{stLoss7} & yes; \texttt{Data.LoP} ref (a)\\
(H8) & \texttt{Data.SlopeIneq} (MB:151); \texttt{Frame\_at}, \texttt{FrameF} (QLDC:188, 258); proved by \texttt{frame\_at\_of} (\path{TopDensityNGeSixteen.lean}:538) from \texttt{TheoremF} (HN, \path{StatementFrameF.lean}:26) & yes\\
(H9), (H10), (H11) & \texttt{Data.HHBound}, \texttt{Data.FParams} (MB:158, 161); clause 2 of \texttt{LinkDIV\_at}, from \texttt{FrameNumberBound} (\path{LocalQuadraticParts.lean}:363); \texttt{Lat\_at} (QLDC:185); \texttt{LinkIRR} (GTI:55) & yes; the third ref (e)\\
(H1)--(H11) together & \texttt{SizeData}$_2$ (SC:17), the conjunction of eleven clauses & yes\\
Lemma M3; (C6) & \texttt{SlopeIneq}, \texttt{MuChar}, \texttt{HNStep}, \texttt{HNSlopes} (HN, \path{Statement.lean}:91, 54, 64, 74); \texttt{htMap}, \texttt{muDef} (HN, \path{StatementPl.lean}:60, 64) & yes\\
Proposition M4 & \texttt{TheoremFSharp}, \texttt{TheoremF} (HN, \path{StatementFrameF.lean}:64, 26), with \texttt{FrameLA1} (:11) for the rank bound and \texttt{LayerCake} for Step~4 & yes\\
Lemma M5 & \texttt{key\_identity} (WM:170), \texttt{key\_identity5} (LC:105) & yes\\
Lemma M6 & \texttt{phi\_bound} (WM:81) & yes\\
Proposition M7 & \texttt{each\_unit\_bound\_statement} (LC:175), with \texttt{s\_lower5}, \texttt{Vvar\_bound5}, \texttt{u\_function\_le\_mean}, \texttt{Ga\_bound5}, \texttt{weighted\_mean\_lower\_bound\_le} (LC:50--173); the error is \texttt{ErrX5at} (QS:45) & yes\\
Lemma M8 & \texttt{errX5at\_le} (PC:613), \texttt{core\_ineq} (\path{SubThreeMulLogTwoLe.lean}:13); \texttt{CErr}, \texttt{KBig}, \texttt{LargeX} (EF:19, 21, 52) & yes; the last three ref (a)\\
Theorem M, \eqref{as:thmM} & \texttt{Deduction.ch\_top} (ChTop:37) & yes\\
\eqref{as:four} & the division is in the proof of \texttt{proposition\_size\_bound}; not read & yes\\
the use of (M): the sum over $u$ & \texttt{ch\_core} (ChCore:343), \texttt{sum\_Th} (ChCore:56); \texttt{core\_gm} (CoreGm:12); \texttt{CoreAtH} (CoreAtH:29) & yes\\
the Hypotheses hold for the records & \texttt{linkCore} (LSS:38): from fourteen statements on the family it gives \texttt{ChainTied}; then \texttt{tiedCore\_of\_chainTied} (LGm:13). The fourteen statements: \texttt{linkGLOBAL7\_items} (\path{LinkMin.lean}:693). The whole: \texttt{FamilyClassConnected} (\path{DivinNOfHypV.lean}:13) & yes\\
(H12), the constants and the windows & \texttt{FamilyClassConnected} (\path{DivinNOfHypV.lean}:13); \texttt{FamCore.Win} (\path{FamCore.lean}:78, a field of the structure at line 71); \texttt{HasWindows} (\path{Window.lean}:7); \texttt{ChebyshevInequalityAtWindowScale} (\path{TopWin.lean}:30) & yes; the second and third ref\\
(H12)(a)--(c), the scale $Y$, the primes and $n$ at one $x$ & \texttt{WinDatum} (\path{TwistForm.lean}:19) with \texttt{PairGood} (\path{PairGood.lean}:18); the windows: \texttt{winPairs} (\path{WinPairs.lean}:5); the admissible set and the field at $x$: \texttt{supplyOfAdmOf} (\path{supplyOfAdmOf.lean}:13, 56), with \texttt{fieldChi} (\path{FieldChi.lean}:20) and \texttt{boundedTw} (\path{BoundedTw.lean}:45), which occur in the proof of \texttt{supplyOfAdmOf} & yes; the last two ref\\
Theorem~\ref{thm:pair} with $A=J$ & \texttt{PairEx} (\path{PairGood.lean}:28), \texttt{QH.pairEx} & yes\\
Corollary M11, \eqref{as:Gbound} & \texttt{ChebyshevInequalityOnWindows} (\path{TopWin.lean}:26), proved by \texttt{windowInequality\_of\_artinInputs} (\path{WindowInequalityOfArtinInputs.lean}:17) from \texttt{core\_gm}; then \texttt{level\_of\_tiedWin} (\path{TopWin.lean}:68) and \texttt{zeroFreeAll\_of\_topG} (\path{ZeroFreeAllOfTopG.lean}:26) & yes\\
the variant for the mean & \texttt{weighted\_mean\_bound\_statement}, \texttt{uniform\_bound\_statement} (LC:199, 237); \texttt{MeanBoundSecondForm} (MB:187) & no\\
\end{longtable}}

\emph{Remarks.} (1) \emph{The declarations named \texttt{ASK} and \texttt{MAIN}.} There are several. On the
route of the final theorem: \texttt{characterSumData.ASK} and \texttt{.MAIN} (FPNS:84, 89) are $P_{u,t,j}$
and $M_{u,t,j}$; the end of the proof (\texttt{ch\_balance}, \texttt{ch\_core}, \texttt{ch\_top}) is stated
with them. \texttt{Data.ASK} is a field of the abstract record; the hypothesis \texttt{IdentU} (SI:67) says
that it equals $\sum_{t,j}$\texttt{characterSumData.ASK}, that \texttt{Data.gam} is Euler's constant and
that \texttt{Data.LpL} is the real part of $\frac{L'}L(1,\chi_d)$. \texttt{LocalQuadraticParts.ASK}
(\path{LocalQuadraticParts.lean}:470) is the same sum written with the numbers of one record; it is tied to
\texttt{characterSumData.ASK} by the statement \texttt{askIsF}
(\path{FrameCodeOfAdditive.lean}:546), whose proof was not read. \texttt{DeficitSplitLemmas.ASK}
(\path{SecondKindBucket.lean}:13) and \texttt{LD.ASK} (\path{WindowSettingSatisfiable.lean}:275) are
abstract forms used to split $P_u-G_{\rm top}$ by primes, in the proof of (H7). All these are in the
dependency lists as used. \texttt{SizeData}$_2$\texttt{.MAIN} (DP2:232) is not $M_u$: it is the main term
$S-A\log\mathrm{Lv}(\mathsf w)$ of (H1). Not used by the final theorem: \texttt{Deduction.BFrame.ASK},
\texttt{.MAIN}; \texttt{Arith.ASK} (\path{LocalToGlobal/GlobalTheory/IsResidueClassZero.lean}:517);
\texttt{FamRec.FieldEq.ASK}; and the identity \texttt{AskMainBalance} (FPNS:98) in that form (the proved
form is \texttt{ch\_balance}).

(2) \emph{Constants.} In the source $c_1=c_2=c_3$ (\texttt{mkC}, QLDC:839: they are the constant of
\texttt{LinkDIV}), $C_{\rm abs}$ is \texttt{eDiv}, $e_L$ is \texttt{eLoss}, $C_H$ is \texttt{cHH}. The
condition \eqref{as:big} is, in the source, $200(1+K)^2(1+C_H)(1+e_K)\le x^{\varepsilon}$ with a $K$ that is
also at least $6$ (\texttt{KBig}, \texttt{LargeX}). The source keeps $(1-\bar\delta)^{-1}$ and
$(\bar\eta+4\sqrt{\bar e})x^{C_H}b$ in the error, where \eqref{as:Err} has $2$ and $6b$; and it has a
number $L\le2x^{\varepsilon}$ where $\mathrm{MLo}$ has $x^{\varepsilon}$. With these changes $\mathrm{Err}_u$
is \texttt{ErrX5at}. Lemma~M8 was proved here directly; the proof of \texttt{core\_ineq} was not read
beyond its statement. \emph{The values of $e_L$ and $C_{\rm abs}$.} The theorem
\texttt{linkGLOBAL7\_items} (\path{LinkMin.lean}:693) ends by applying \texttt{linkCore} with the value
$0$ for \texttt{eLoss}, and its conclusion has $e_D=C_{\rm abs}$ and $e_L=0$. So $e_L=0$ on the route of
the final theorem. In that theorem $C_{\rm abs}$ is a parameter, with the hypotheses $0\le C_{\rm abs}$,
$0<C_H$ and $3\le C_0$; it is still a parameter in \texttt{linkGLOBAL7\_min} (\path{LinkMin.lean}:723),
which uses it. The statement \texttt{FamilyClassConnected} asserts the existence of $e_D$ and $e_L$. Where
the value of $C_{\rm abs}$ is fixed was not determined.

(3) \emph{The route.} \texttt{riemannHypothesis} $\leftarrow$ \dots $\leftarrow$
\texttt{zeroFreeAll\_of\_topG} $\leftarrow$ \texttt{level\_gm\_le\_on\_family\_scales} $\leftarrow$
\texttt{windowInequality\_of\_artinInputs} $\leftarrow$ \texttt{core\_gm} $\leftarrow$ \texttt{ch\_core}
and \texttt{ch\_top} $\leftarrow$ \texttt{each\_unit\_bound\_statement} and \texttt{errX5at\_le}. Each of
the names before \texttt{errX5at\_le} is in the dependency lists as used. The hypotheses of
\texttt{core\_gm} are supplied by \texttt{tiedCore\_of\_chainTied} from \texttt{ChainTied}, which
\texttt{linkCore} proves from: \texttt{CountF}, \texttt{SlopesMF}, \texttt{SlopesNF}, \texttt{LatF},
\texttt{NodesF}, \texttt{FrameF}, \texttt{LossFS}, \texttt{CharF}, \texttt{WinF}, \texttt{ChosenTT},
\texttt{ChosenQD}, \texttt{ChosenThin}, \texttt{LinkIRR} and \texttt{LinkDIV}. The first nine are the
statements \texttt{*\_at} of the table for all records of the family; they are proved in the source from
\texttt{LinkG}, \texttt{LinkS}, \texttt{LinkIRR}, \texttt{LinkDIV}, \texttt{LinkMu} and conditions on the
family (\texttt{linkGLOBAL7\_items}). We read the statement of \texttt{linkCore} and the part of its proof
that fills the eleven clauses; we did not read how the nine statements are proved. \texttt{LinkV\_at}
is in the lists as not used; the bound for Galois cohomology enters as \texttt{HypV}
(\path{DivinNOfHypV.lean}:10).

(4) \emph{Proposition M4 and the source.} The source proves \texttt{TheoremFSharp} from \texttt{FrameLA1}
and the statements on slopes of the folder \texttt{HN}; we read the statements, not the proof. Our proof
of Proposition~M4 is not a translation of it. \texttt{TheoremF} is \texttt{TheoremFSharp} weakened by the
term $\log|\mathsf W|$ (\texttt{theoremF\_of\_sharp}). In \texttt{TheoremF} the rank bound is not a
hypothesis: its hypotheses are the bounds for the diagonal blocks and for $\operatorname{tr}(Q^2)$ of
\S III.8, at one complex place. The statement \texttt{FrameF} (QLDC:258) carries the bound
$\mathsf e_u\le e_B\le\frac14$, and the theorem behind it the conditions $\alpha_{\rm ram}\le\frac14$,
$x^{-C_H}\le\frac1{20}$; these belong to what is quoted when (H8) is taken from the source.

(5) \emph{The translates.} \texttt{ch\_top} gives the bound for every element of the group: in
\texttt{tiedCore\_of\_chainTied} (LGm:13) the translates are indexed by the whole group. \texttt{ch\_core}
uses the bound for $u$ in the subgroup $A_0$ generated by one element
(\path{ZsmulGeneratorEqZeroOfProjection.lean}:63), sums over $A_0$ and divides by $|A_0|$ (\texttt{sum\_Th}),
and requires $T_2\subset A_0$ (\texttt{CoreAtH}). For the frame \texttt{cyclicFrame}
(\path{FrameCyclicGenerator.lean}:150) the group is $\Z/2^e$, the generator corresponds to the class of
$5$ (\texttt{cycIso\_one}, :80) and $A_0$ has index $1$ (\texttt{cyclicFrame\_index}, :158);
\texttt{cyclicFrame} is in the dependency lists as used. So for this frame $A_0$ is the whole group and
the mean is $\frac1n\sum_{u\in H}$. For the trivial character, that is, for $\zeta$, the index is $1$
for every frame that the statements admit, whether or not it is literally \texttt{cyclicFrame}. Indeed
\texttt{Level.cyc} $\psi$ is defined as \texttt{Level.ofSub (Subgroup.zpowers} $\psi.\chi$\texttt{)}
(\path{ExplicitFormulaErrorOfHostConditions.lean}:121), with \texttt{o} the order of that group of
characters; for the trivial character the order is $1$. The structure \texttt{EAtH} has the field
\texttt{o\_eq}: $\varphi.o=\Lambda.o$ (:145), and \texttt{Frame.o} is the index of $A_0$. So every frame
whose data satisfy \texttt{EAtH (Level.cyc 1)} has $A_0$ equal to the whole group. (Statements only.) The source
also has frames whose group has order $c\cdot2^e$ and whose $A_0$ has index $o>1$ (\texttt{fieldChi};
\texttt{core\_gm} then has a weight $(3o+1)^{-1}$). They serve the Dirichlet characters.

(6) \emph{The bound for denominators.} \eqref{gc:div} is \texttt{DIVinN\_at} (GTI:59); its left side,
\texttt{etaSum}, is that of \texttt{LinkN\_at} (GTI:40). The theorem \texttt{inEta\_num}
(\path{SumLeEtaSumNum.lean}:60) turns \texttt{DIVinN\_at} into a row of the package from which the size
bound is proved. The clause that \S III.7 calls (S8) is a field of the structure \texttt{ArithAt}
(\path{ConstantGeTwentyEightAlt.lean}:15); its left side is a sum over $j$ of terms that contain
\texttt{etaSumBound} (\path{LinkDef.lean}:93), which is \texttt{bd}$_j$ plus two logarithms. That
\texttt{bd} is, on the route, the bound of \texttt{DIVinN\_at} was not followed. \texttt{divinN\_of\_hypV}
is in the lists as used. For (H9): clause 2 of \texttt{LinkDIV\_at} is supplied from
\texttt{FrameNumberBound}, and the statement that yields \texttt{FrameNumberBoundFamily} has
\texttt{DIVinN} and \texttt{ComplementTwoFamily}, the source of (H2), among its hypotheses
(\path{FinalAll.lean}:34). Statements read; proofs not read.

(7) \emph{The scales and the primes: the statements behind \textup{(H12)}.} A family
(\texttt{FamCore}) has a set \texttt{Win} of scales. \texttt{FamilyClassConnected} reads: there are
$e_D$, $e_L$, $\varepsilon_1>0$ such that for every $0<\varepsilon<\varepsilon_1$ there is a family and
there are $K$, $c_{HH}$, $e_K$, $C_L$, $x_0$ with \texttt{ChainTied}, that is, the Hypotheses, at every
$x\ge x_0$ in \texttt{Win}. So $e_D=C_{\rm abs}$, $e_L$ and $\varepsilon_1$ are chosen before
$\varepsilon$, and $K$, $C_H$, $e_K$ after it; this is the order of the quantifiers in (H12).
\texttt{ChebyshevInequalityAtWindowScale} requires in addition \texttt{HasWindows Win la}: for every $y$
there is $Z\ge y$ with $[Z,Z^{la}]\subset\texttt{Win}$; \texttt{zeroFreeAll\_of\_topG} is applied with
$la=\frac32$ (\path{EveryCharacter.lean}:10). On the route the set \texttt{Win} is
$\{x:\exists\,Y,q,d,e,\ \texttt{WinDatum}\ \psi\ \varepsilon\ J\ c\ x\ Y\ q\ d\ e\}$
(\path{supplyOfAdmOf.lean}:56). \texttt{WinDatum} has the fields $x^{1/J}\le Y$, $Y\le x$,
\texttt{good : PairGood} $\psi.N$ $\varepsilon$ $Y$ $x$ $q$ $d$, the two bounds for $c\cdot2^e$, and the
bound for Galois cohomology for the field. \texttt{PairGood} $M$ $\varepsilon$ $Y$ $X$ $q$ $d$ has four
fields: \texttt{pair} (parts (a)--(c) of Theorem~\ref{thm:pair} at $Y$), \texttt{coprime},
\texttt{bounds} (part (d) for all $y$ with $Y\le y\le X$) and \texttt{boxes} (part (e)). So the condition
on $(q,d)$ at $x$ contains Estimate~B on $[Y,x]$, in particular at $x$. \texttt{winPairs} states: there
is $J\ge1$ such that for every $y$ there is $Z\ge y$ such that for every $x\in[Z,Z^{3/2}]$ there are $Y$,
$q$, $d$, $e$ with \texttt{WinDatum}; $J$ is found after $\varepsilon$. We read its proof: the pair comes
from \texttt{PairEx} with $A=J$ at the scale $X^{1/J}$, once for the window, and $Y=X^{1/J}$ for every $x$
of the window. We read the proof of \texttt{supplyOfAdmOf} from line 56 on: the frame and the field at
$x$ are produced by \texttt{fieldChi} from the $q$, $d$ of the \texttt{WinDatum} of $x$ (its conclusions
\texttt{H.q = q}, \texttt{H.d = d}), and the character bounds required of the data (the fields
\texttt{bdG}, \texttt{bdL} of \texttt{EAtH}) are obtained by \texttt{boundedTw} from \texttt{PairGood},
by taking \texttt{bounds} at $y=x$. So the primes of the data at $x$ are the primes of the
\texttt{WinDatum} of $x$. In \texttt{windowInequality\_of\_artinInputs} (proof read) the
character estimates and the bound of (M) (\texttt{CoreAtH}) are taken for the same $x$ and the same
record, hence the same $q$, $d$. \texttt{FamCore} has fields \texttt{q\_eq}, \texttt{d\_eq}: all records
at one $x$ have the same primes; nothing ties the primes of different $x$. The conclusion is
$G(x)\le x^{1/2+(e_D+e_L+7)\varepsilon}$ for $x\ge x_1$ in \texttt{Win}
(\texttt{ChebyshevInequalityOnWindows}). \emph{Not read:} how the source's form of Estimate~C
(\texttt{ThinH}) is supplied from \texttt{boxes}; the proof of \texttt{core\_gm}, which is the source's
form of Part~II, \S5; the other parts of \texttt{EAtH}; the first 55 lines of the proof of
\texttt{supplyOfAdmOf}. \emph{Differences from \textup{(H12)}.} (H12) is stated for the objects of
\S\ref{ob:sec}; the source speaks of its records. The source has $\psi.N$ in the place of $M=1$ and
$c\cdot2^e$ in the place of $2^e$; for $\zeta$, $\psi$ is trivial and $\psi.N=1$. For $\zeta$, $c=1$: the number
$c$ comes from \texttt{ordAny} (\path{OrdAny.lean}:86), which gives $1\le c\le|\langle\psi\rangle|$, and
$|\langle\psi\rangle|=1$ for the trivial character. (The statement of \texttt{fieldChi} keeps only
$1\le c$; the bound is in its proof.) Clause (e) of (H12) is not one statement of the source; it was checked against the
table of O23 (\S7, Status, (3)).

\textbf{What remains.} (i) Proofs of the Hypotheses. The table of \S1 gives the status of each. For
(H2), (H3), (H5), (H9) and the third part of (H4) no section of this paper gives a proof or a reduction,
and the lower bound of (H4) is quoted in Part~II in shape only, the size of its allowance being
assumed. In the source (H2), (H3) and (H9) are clauses of the
statement that contains (H1). For (H12) see \S7: it is assumed, and the paper supplies only the
existence of the primes. (ii) That the abstract numbers of
\S4 are the slopes, heights and levels of the objects of \S\ref{ob:sec}: this was read from the
definitions \texttt{TieLocal} and \texttt{Mod.data} of the source, not verified further. Proposition~M4
requires, for (H8), that the slopes of $\Lambda^{\rm gl}_{\mathsf w}$ and $\mathcal R^\vee$ and the height
in (8b) are taken for the same norms. In the source this is so by definition: a lattice is a record with
a basis, a matrix at each complex place and a module, and both its slopes and the height in (8b) are
defined from these matrices (Lean: \texttt{HLat.mu}, \texttt{HLat.muDual}, \texttt{htMap} in
\path{DataArith.lean} and \texttt{HN}). We did not check that the Lean definition of the slope agrees
with \eqref{as:mu} beyond its shape, twice a slope for the Euclidean norm plus the constant
$\mathrm{cst}_K$. (iii) The value of $C_{\rm abs}$ on the route of the final theorem; and the one link
of Remark~(6). (iv) The step from \eqref{as:M} to Estimate~A of Part~I is in Part~II, \S5 and \S3. Its hypotheses
were compared one by one with (H12) in the proof of Corollary~M11; the proofs of Part~II were not
re-examined. The question whether the mean for $\zeta$ is over all translates is settled at the level of
statements (Remark~(5)). (v) (C6) is cited from memory. (vi) The identification of the pairings (Comment~(e) to O25; \S III.8, \S2)
was not checked in the source; see the status of (H8).


\section{The simplest case in full: the lemniscate}\label{el:sec}

\begin{illbox}
The whole of this section is an illustration. The Lean development does not treat the lemniscate; the
divisibility statement that the proof needs, for the theta series of the construction, is hypothesis (B),
taken from the Lean development (\S III.5). Nothing here is used in the deduction of the Riemann
hypothesis (see ``What the proof uses'' after Corollary~M11 in \S\ref{as:sec}).
\end{illbox}

Part~I used one example twice. In \S I.3 the derivatives $A_n$ of the lemniscate sine at the origin showed
what \emph{divisibility at inert primes} means, and in \S I.4.1 the theta series of the Gaussian integers
showed what \emph{size through a period} means. Several statements about this example were given there as
observations of a computation, and in \S III.5 a hypothesis \textup{(B)} was verified for an example built
from the same curve only up to a finite degree. This section proves these statements. It uses the classical
theory of elliptic functions and congruences for binomial coefficients; what is taken from books is listed
below under ``What is cited''. The reader will then have seen one complete instance of the two mechanisms of
the proof, with no hypothesis left open.

Nothing in this section is taken from the Lean development. The primes $p\equiv3\bmod4$ are called \emph{inert} and the primes
$p\equiv1\bmod4$ \emph{split}, as in \S I.3.

\emph{What was observed, and what is proved here.} The statements of Part~I and of \S III.5 that rested on
computation are these. The list serves as a table of contents. It uses some words that are defined only
later: $v_p$ is the exponent of the prime $p$ (\S5), and formal groups, their logarithms and their height
are defined in \S4.
\begin{enumerate}[label=\textup{(O\arabic*)},leftmargin=3.2em,itemsep=1pt,topsep=2pt]
\item \S I.3: the numbers $A_n$ are integers, and $A_n=0$ unless $n\equiv1\bmod4$ (cited there from
Hurwitz \cite{Hur99}). \emph{Proved here}: Lemma~Le4.
\item \S I.3: the exponent of $p$ in $A_n$ is at least $n\,p/(p^2-1)-1$ (in the first version of the paper
this was an observation, for the inert primes $3,7,11,19,23,31$ and $n\le201$). \emph{Proved here for every inert prime and every $n$}, in the form
$v_p(A_n)\ge\lfloor n\,p/(p^2-1)\rfloor$: Theorem~Le18. For the numbers $A_n$ this is the same inequality as
the one observed (see the comparison after Theorem~Le18); what is added is the range, all inert $p$ and
all $n$.
\item \S I.3: in that range the exponent exceeds $n\,p/(p^2-1)$ by at most $1.65$, and this ceiling does
not persist (the excess is $4.625$ at $n=1945$ for $p=3$). Both values are confirmed. That the excess is
unbounded is \emph{not proved}.
\item \S I.3: the split primes $5,13,17,29$ divide none of the $A_n$ with $n\le201$, $n\equiv1\bmod4$
(the other $A_n$ are zero), and $157$ divides
$A_{101}$. Confirmed, and extended to $n\le2001$. \emph{Proved here}: a split prime $p$ does not divide $A_p$
(Proposition~Le20); and $5$, $13$, $17$, $29$ divide no $A_n$ with $n\equiv1\bmod4$ at all
(Proposition~Le20a\,(iii)).
\item \S I.3: multiplication by an inert prime $p$ on the formal group is $X\mapsto X^{p^2}$ modulo $p$,
and the logarithm behaves like the model series $X+X^{p^2}/p+X^{p^4}/p^2+\cdots$ (cited there from
\cite{Deu41,Hon70,LT65}). \emph{Proved here} for this curve, with the signs determined: Theorems~Le12
and~Le13 and Corollary~Le16. The signs alternate: the model series is
$X-X^{p^2}/p+X^{p^4}/p^2-\cdots$, and multiplication by $p$ is $X\mapsto-X^{p^2}$ modulo $p$, with no
further terms (Corollary~Le16\,(ii)).
\item \S I.3, defect (D3): $|A_n|$ is about $n!/R^n$ with $R=1.854\ldots$ \emph{Proved here}, with
$R=\varpi/\sqrt2$: Theorem~Le9.
\item \S I.4.1: $T_0=\sqrt2\,\varpi/\pi$ and $T_4/T_0=\frac14(\varpi/\pi)^4$. \emph{Proved here}, from the
two-squares theorem, which is cited: Proposition~Le7.
\item \S I.4.1: the numbers $r_k=T_{4k}/(T_0(\varpi/\pi)^{4k})$ are rational with denominator $4^k$
(identified numerically for $k\le30$). \emph{Proved here}: they are rational, with denominator dividing
$4^k$ (Corollary~Le8). That the denominator is exactly $4^k$ is observed for $k\le120$ and \emph{not
proved}.
\item \S I.4.1: the identity $\Theta(z)/T_0=\sigma(u)^2\wp(u)$ with $u=\varpi z/2\pi$ (observed to $400$
digits). \emph{Proved here}: Theorem~Le6. The normalisation printed in \S I.4.1 is correct.
\item \S I.4.1: the exponent of $3$ in $r_{30}$ is $14$. Confirmed; by Corollary~Le8 the rational number
$r_{30}$ is now known exactly and not only to $250$ digits.
\item \S III.5, ``An example'': hypothesis \textup{(B4)} was checked up to total degree $84$, $72$, $144$
for $\ell=3,7,11$, and \textup{(B)} in all degrees was not proved. \emph{Proved here in all degrees}:
Proposition~Le22.
\item \S III.5, same paragraph: equality in \eqref{dv:beta} at the exceptional monomials up to degree
$800$, and the numerical checks of Theorem~J on this example. \emph{Not treated here.}
\end{enumerate}
The section has one more result, which was not stated in Part~I: a Chebyshev inequality for the inert
primes that follows from Theorems~Le9 and~Le18 (Proposition~Le19). It makes the three defects of \S I.3
precise.

\emph{What is cited.} The proofs of this section use the following published statements and no others.
\begin{enumerate}[label=\textup{(C\arabic*)},leftmargin=3.2em,itemsep=1pt,topsep=2pt]
\item The first chapter of Apostol's book on modular functions \cite{Apo90}, for elliptic functions and
the Weierstrass function $\wp$: Theorems 1.4, 1.7, 1.8 (used only through its corollary that $\wp$ takes
every value twice) and 1.10 to 1.14, with the definitions of $\wp$, $g_2$, $g_3$ that stand beside them, and
Lemma~1 of its \S1.5 on the
convergence of $\sum|\omega|^{-\alpha}$ (all read in a local copy; the numbers are given where they are
used). That chapter does not treat the Weierstrass functions $\sigma$ and $\zeta_{\rm w}$; what we need of
them is proved in \S1a.
\item The theorem that the number of representations of $n\ge1$ as a sum of two squares is
$4\sum_{d\mid n}\chi_4(d)$: Hardy and Wright \cite{HW}, Theorem~278 (read in a local copy). For $n$ odd it is
also the case $d=-4$, $w=4$ of formula~(12) in Chapter~6 of Davenport \cite{Dav00}, which is stated there,
without proof, for $n$ prime to $d$. (This was read in a local copy of the second edition, 1980, and the
page, p.~47, is that of the second edition. The entry \cite{Dav00} of the bibliography is the third edition,
2000, which was not consulted; for this reason the chapter and the formula number are given.) The theorem is
used once, in Proposition~Le7.
\item Lucas's congruence $\binom{ap+b}{cp+d}\equiv\binom ac\binom bd\pmod p$ for $0\le b,d<p$ (Lucas
\cite{Luc78}; cited from memory). \S III.5 also uses it. It is used in Proposition~Le20, and through it in
Proposition~Le20a\,(iii).
\item Four facts of a first course, in the numbering of Apostol's \emph{Introduction} \cite{Apo76} (read in a
local copy): Legendre's formula for the exponent of $p$ in $n!$ (Theorem~3.14; it is \eqref{leg}), the
theorem of Euler and Fermat (Theorem~5.17), Chebyshev's bound $\pi(n)=O(n/\log n)$ (Theorem~4.6), used in
the form $\sum_{p\le n}\log p=O(n)$ in Proposition~Le19, and, in a remark after Proposition~Le19 only,
$\sum_{p\le n,\ p\equiv3\,(4)}(\log p)/p=\frac12\log n+O(1)$ (Theorem~7.3).
\end{enumerate}
Four more works are named for comparison, and nothing is taken from them. The functional equation lemma of
Hazewinkel \cite{Haz78}, I.2.2 (read in a local copy): the case needed is proved here as Lemma~Le14. The
congruences of Atkin and Swinnerton-Dyer in the form of \cite{Haz78}, (33.2.3) (read in a local copy):
Theorem~Le12 is proved here directly. Whittaker and Watson \cite{WW27} for the addition theorem (from
memory): it is proved here as Lemma~Le5. Deuring's theorem \cite{Deu41} on the height at split primes (from
memory): it is mentioned after Proposition~Le20 and not used. Lemma~1 of \S III.5 is not used either:
Proposition~Le15 contains it in one variable.

\emph{Letters.} They are local to this section. $\varpi$ is the lemniscate constant; $\mathcal P=\varpi\Z[i]$
the period lattice; $\wp,\zeta_{\rm w},\sigma$ the Weierstrass functions of $\mathcal P$ (the index
${\rm w}$ distinguishes $\zeta_{\rm w}$ from the Riemann zeta function); $u_0=(1+i)\varpi/2$;
$\mathrm{sl}$, $\mathrm{arcsl}$, $A_n$ as in \S I.3; $\mathrm{sl}^\flat(x)=(1+i)\,\mathrm{sl}(x/(1+i))$;
$J(u)=\int_0^u\mathrm{sl}^2$ and $M(u)=\int_0^uJ$; $\Phi=\sigma^2\wp$; $\Theta,T_n,r_k$ as in \S I.4.1;
$\theta$ the theta series in one variable, and $G$ an auxiliary entire function in the proof of
Corollary~Le6a; $S(u)$ a series of hyperbolic secants; $C_0,C_1$ the constants and $\delta_n$ the relative
error of Theorem~Le9, with the poles $u_1,u_2$, the radius $\varrho$ and the angle $\vartheta=\arctan\frac12$ of its proof; $F$ the formal group law
of the lemniscate; $\kappa_k=\binom{2k}{k}4^{-k}$; $p$ a prime, $q=p^2$, $v_p$ the $p$-adic valuation,
$\Z_{(p)}$ the ring of rational numbers with denominator prime to $p$; $N!'$ the product of the integers
$j\le N$ prime to $p$; $L_+$ and $L_-$ the two model logarithms \eqref{el:Lpm}, $g_\pm$ their inverse series,
$H_\pm$ the rescaled inverses, $\mathfrak r=p^{p/(q-1)}$; $U_p=\mathrm{sl}\circ L_-$; $G_-$ the model group
law; $\mathcal K$ the cube function; and in \S9 the letters $\ell,c,\sigma_c,t,m,\tilde\beta,\beta,\zeta,K,K_0$ of
the example of \S III.5, with $l_c$, $h_c$, $\tilde z=c^{1/4}z$, the valuation $v_\ell$ of $\Q(i)$, the
auxiliary series $\mathcal K_m$, $\mathcal K^\ast$, $\mathcal E$, $\mathcal N$ of Steps~4 to~6, the
cyclotomic polynomial $\Phi_{q-1}$, its degree $D$, the auxiliary ring $\mathcal A$ and the class $\psi$ of
$x$ in it ($\psi$ is not the root of unity $\zeta$ of \S III.5). In \S1a, $\eta(\omega)$ and $C_\omega$
are the constants in the behaviour of $\zeta_{\rm w}$ and $\sigma$ under a period $\omega$. In \S\S1--8 the
inert prime is called $p$, as in Part~I; it is the prime $\ell$ of \S III.5. The series $L_+$ is the
logarithm \eqref{dv:log} in one variable. The letter $v$ without index occurs only for the valuation of the field $K$ of \S III.5, in \S9 and
\S10; the $p$-adic valuation is always written $v_p$, and the second variable of the addition theorem is
$u'$. Some letters have another
meaning elsewhere in the paper: $q$ is $p^2$ here and an auxiliary prime in Parts~I and~II; $S(u)$ is not
the sum $S(x)$ of Part~I; $T$ is the variable of $L_\pm$, and $T_n$ are the derivatives of $\Theta$; $G$
(Corollary~Le6a) and the model group law $G_-$ are not the $G$ of Part~I.

\subsection*{1.\quad The lemniscate functions}

We start from the Weierstrass function of a square lattice and obtain the lemniscate sine from it. This is
shorter than the historical route, and it gives at once that $\mathrm{sl}$ is meromorphic in the whole
plane. For a lattice $\mathcal L\subset\C$ recall the definitions of $\wp$ and of the invariants $g_2$,
$g_3$, and the differential equation (\cite{Apo90}, Chapter~1: the series are definitions there, and the
properties are Theorems 1.10 to 1.14)
\begin{gather}\label{D:440}
 \wp(u)=\frac1{u^2}+\sum_{\omega\in\mathcal L,\ \omega\ne0}\Big(\frac1{(u-\omega)^2}-\frac1{\omega^2}\Big),\qquad
 \wp'^2=4\wp^3-g_2\wp-g_3,\\
 g_2=60\sum_{\omega\in\mathcal L,\ \omega\ne0}\omega^{-4},\qquad g_3=140\sum_{\omega\in\mathcal L,\ \omega\ne0}\omega^{-6}.
\end{gather}
The function $\wp$ is even, has the periods $\mathcal L$, has a double pole at each point of $\mathcal L$
and no other pole, and takes every value twice in a period parallelogram, counted with multiplicity. The
last statement is a corollary of \cite{Apo90}, Theorem 1.8, which says that an elliptic function has as many
zeros as poles in a period parallelogram: apply it to $\wp-a$ for a constant $a$, which has one double pole
there. The derivative $\wp'$ vanishes exactly at
the points $\omega/2\notin\mathcal L$ with $\omega\in\mathcal L$ (this is shown in the proof of
Theorem~1.14 of \cite{Apo90}); the values of $\wp$ at the three classes
of these points are the roots $e_1,e_2,e_3$ of $4x^3-g_2x-g_3$, and they are distinct. We also use two
facts about an elliptic function $f$ with period lattice $\mathcal L$: if $f$ has no poles it is constant,
and the sum of its residues in a period parallelogram is zero (\cite{Apo90}, Theorems 1.4 and 1.7).

\textbf{Lemma Le1 (the square lattice).} \emph{There is exactly one real number $\varpi'>0$ such that the
lattice $\mathcal P=\varpi'\Z[i]$ has $g_2=4$ and $g_3=0$. For this lattice $\wp(iu)=-\wp(u)$, the function
$\wp$ is real on the real axis, and}
\begin{equation}\label{el:evalues}
 \wp(\tfrac12\varpi')=1,\qquad \wp(\tfrac12i\varpi')=-1,\qquad \wp(u_0)=0\quad\text{where }u_0=\tfrac12(1+i)\varpi' .
\end{equation}
\emph{On the interval $0<u\le\frac12\varpi'$ the function $\wp$ decreases from $+\infty$ to $1$.}

\emph{Proof.} Let $\mathcal L=\lambda\Z[i]$ with $\lambda>0$. Multiplication by $i$ permutes $\mathcal L$
and multiplies $\omega^{-6}$ by $-1$; so $g_3=0$. In the series for $\wp(iu)$ put $\omega=i\omega'$: each term
is multiplied by $i^{-2}=-1$, so $\wp(iu)=-\wp(u)$. Complex conjugation permutes $\mathcal L$, so
$\wp(\bar u)=\overline{\wp(u)}$ and $\wp$ is real on the real axis. Put $e_1=\wp(\lambda/2)$, a real number.
Then $\wp(i\lambda/2)=-e_1$. The sum of the roots of $4x^3-g_2x$ is $0$, so the third root is $0$; the roots
are distinct, so $e_1\ne0$; and $4x^3-g_2x=4x(x-e_1)(x+e_1)$ gives $g_2=4e_1^2>0$. Since
$g_2(\lambda\Z[i])=\lambda^{-4}g_2(\Z[i])$, exactly one $\lambda=\varpi'$ gives $g_2=4$. Then $e_1=\pm1$,
and the third half-period value is $\wp(u_0)=0$. On $0<u<\varpi'/2$ the function $\wp$ is real and finite
and $\wp'\ne0$. Since $\wp'^2=4\wp(\wp^2-1)$, the function $\wp$ does not take the values $0,1,-1$ there. It
tends to $+\infty$ as $u\to0$. So $\wp>1$ on this interval, $\wp'<0$, and $e_1=\wp(\varpi'/2)\ge1$. Hence
$e_1=1$. \qed

\textbf{Definition (the lemniscate sine).} By Lemma~Le1 the function $1/\wp$ has a double zero at each
point of $\mathcal P$, a double pole at each point of $u_0+\mathcal P$ (where $\wp=0=\wp'$), and no other
zero or pole, because $\wp$ takes the value $0$ only twice in a period parallelogram. A meromorphic function
$f$ on $\C$ all of whose zeros and poles have even order has a meromorphic square root: $f'/2f$ has only
simple poles with integer residues, so $\exp\int f'/2f$ is single valued. Near the origin
$1/\wp(u)=u^2(1+O(u^4))$. We define $\mathrm{sl}$ as the meromorphic function on $\C$ with
\begin{equation}\label{el:sldef}
 \mathrm{sl}(u)^2=\frac1{\wp(u)},\qquad \mathrm{sl}(u)=u+O(u^5)\quad(u\to0).
\end{equation}

\textbf{Proposition Le2 (properties of $\mathrm{sl}$).} \emph{Let $\varpi=2\int_0^1dt/\sqrt{1-t^4}=2.62205\ldots$}
\begin{enumerate}[label=\textup{(\alph*)},leftmargin=2.4em,itemsep=1pt,topsep=2pt]
\item $\mathrm{sl}'^2=1-\mathrm{sl}^4$ \emph{and} $\mathrm{sl}''=-2\,\mathrm{sl}^3$.
\item \emph{$\varpi'=\varpi$, and for $0\le u\le\varpi/2$ the number $t=\mathrm{sl}(u)$ is the solution in
$[0,1]$ of $u=\int_0^tds/\sqrt{1-s^4}$. So $\mathrm{sl}$ is the lemniscate sine of \S I.3, and
$\mathrm{sl}(\varpi/2)=1$.}
\item $\mathrm{sl}(-u)=-\mathrm{sl}(u)$, $\mathrm{sl}(iu)=i\,\mathrm{sl}(u)$,
$\mathrm{sl}(u+\varpi)=\mathrm{sl}(u+i\varpi)=-\mathrm{sl}(u)$.
\item \emph{The zeros of $\mathrm{sl}$ are the points of $\mathcal P$ and the poles are the points of
$u_0+\mathcal P$; all are simple.}
\item $\mathrm{sl}(u+u_0)\,\mathrm{sl}(u)=-i$. \emph{The residue of $\mathrm{sl}$ at the pole
$u_0+(a+bi)\varpi$ is $-i(-1)^{a+b}$.}
\end{enumerate}

\emph{Proof.} (a) Differentiate \eqref{el:sldef}: $2\,\mathrm{sl}\,\mathrm{sl}'=-\wp'/\wp^2$. Square and use
$\wp'^2=4\wp^3-4\wp$: $\mathrm{sl}'^2=(\wp^2-1)/\wp^2=1-\mathrm{sl}^4$. Differentiating once more and
dividing by $2\,\mathrm{sl}'$ gives $\mathrm{sl}''=-2\,\mathrm{sl}^3$.

(b) By Lemma~Le1, on $0<u\le\varpi'/2$ the function $1/\wp$ increases from $0$ to $1$. So $\mathrm{sl}$,
which is positive near $0$, increases from $0$ to $1$ there, and $\mathrm{sl}'=+\sqrt{1-\mathrm{sl}^4}$.
Hence $u=\int_0^{\mathrm{sl}(u)}ds/\sqrt{1-s^4}$, and at $u=\varpi'/2$ this reads $\varpi'/2=\varpi/2$.

(c) Each of $-\mathrm{sl}(-u)$, $-i\,\mathrm{sl}(iu)$ has the square $1/\wp(u)$ and the expansion
$u+O(u^5)$; so each equals $\mathrm{sl}(u)$. The quotient $\mathrm{sl}(u+\varpi)/\mathrm{sl}(u)$ is
meromorphic with square $1$, so it is a constant $\pm1$. At $u=-\varpi/2$ it is
$\mathrm{sl}(\varpi/2)/\mathrm{sl}(-\varpi/2)=-1$. Finally
$\mathrm{sl}(u+i\varpi)=i\,\mathrm{sl}(-iu+\varpi)=-i\,\mathrm{sl}(-iu)=-\mathrm{sl}(u)$.

(d) This was shown in the Definition.

(e) The function $\wp(u+u_0)\wp(u)$ is elliptic. At a point of $\mathcal P$ the double pole of $\wp(u)$
meets the double zero of $\wp(u+u_0)$, and at a point of $u_0+\mathcal P$ the roles are exchanged; so it has
no poles and is constant. At $u=\varpi/2$ its value is $\wp(\varpi+\tfrac12i\varpi)\cdot1=-1$ by
\eqref{el:evalues}. Hence $(\mathrm{sl}(u+u_0)\mathrm{sl}(u))^2=-1$, and the product is a constant, $i$ or
$-i$. To find the sign we compute $\wp$ on the diagonal. For real $x$,
$\wp((1+i)x)=(1+i)^{-2}\wp_1(x)=-\tfrac i2\wp_1(x)$, where $\wp_1$ is the Weierstrass function of the lattice
$(1+i)^{-1}\mathcal P$. That lattice is stable under conjugation, so $\wp_1$ is real on the real axis. It is
positive for small $x>0$, and its first positive zero is at $(1+i)x=u_0$, that is at $x=\varpi/2$. So
$\wp(u_0/2)$ is a negative multiple of $i$. Since $\wp(u_0/2)^2=\wp(-u_0/2)\wp(u_0/2)=-1$, we get $\wp(u_0/2)=-i$ and
$\mathrm{sl}(u_0/2)^2=i$. At $u=-u_0/2$ the product is $\mathrm{sl}(u_0/2)\mathrm{sl}(-u_0/2)=-i$. So
$\mathrm{sl}(u_0+x)=-i/\mathrm{sl}(x)=-i/x+O(x^3)$: the residue at $u_0$ is $-i$, and by (c) it is
multiplied by $-1$ under each of the translations $\varpi$, $i\varpi$. \qed

\noindent Numerically\footnote{Script \texttt{el1\_theta\_identity.py}. The scripts named in the footnotes of
this section are in the folder \texttt{outline/examples/elliptic/}.}: $g_2(\varpi\Z[i])=3.99988$ from the lattice sum over
$|a|,|b|\le60$; $\wp-\mathrm{sl}^{-2}$, the product in (e) plus $i$, and the residues agree with the
statements to $60$ digits or to the precision of the difference quotient used.

\textbf{Hurwitz series.} Let $R$ be a subring of a field of characteristic zero. A \emph{Hurwitz series}
over $R$ is a power series $\sum_{n\ge0}a_nx^n/n!$ with all $a_n\in R$. These are the series whose
derivatives of all orders at $0$ lie in $R$. For example $e^x-1=\sum_{n\ge1}x^n/n!$ and
$\log(1+x)=\sum_{n\ge1}(-1)^{n-1}(n-1)!\,x^n/n!$ are Hurwitz series over $\Z$, and so is every power series
with coefficients in $\Z$; but $(e^x-1)/x=\sum_{n\ge0}\frac1{n+1}\,x^n/n!$ is not.

\textbf{Lemma Le3.} \emph{The Hurwitz series over $R$ form a ring, closed under $d/dx$ and under
$\int_0^x$. If $h$ is a Hurwitz series over $R$ with $h(0)=0$, then $h^j/j!$ is a Hurwitz series over $R$ for
every $j\ge0$; hence so is $\exp(h)$.}

\emph{Proof.} The product of $\sum a_nx^n/n!$ and $\sum b_nx^n/n!$ has the coefficients
$\sum_k\binom nka_kb_{n-k}$ in the same form. Differentiation and integration shift the sequence $(a_n)$.
For the last statement use induction on $j$: $h^j/j!$ vanishes at $0$ and has the derivative
$(h^{j-1}/(j-1)!)\,h'$. The series $\exp(h)=\sum_jh^j/j!$ is a sum in which only finitely many terms
contribute to each coefficient. \qed

\textbf{Lemma Le4 (the Taylor coefficients).} \emph{Let $\mathrm{sl}(u)=\sum_{n\ge1}A_nu^n/n!$.}
\begin{enumerate}[label=\textup{(\roman*)},leftmargin=2.4em,itemsep=1pt,topsep=2pt]
\item \emph{$A_n\in\Z$, and $A_n=0$ unless $n\equiv1\bmod4$.}
\item \emph{$(-1)^kA_{4k+1}$ is a positive integer divisible by $4^k$. Equivalently
$\mathrm{sl}^\flat(x)=(1+i)\,\mathrm{sl}(x/(1+i))$ is a Hurwitz series over $\Z$, with positive coefficients
in the degrees $4k+1$.}
\item \emph{The inverse series of $\mathrm{sl}$ is}
$\ \mathrm{arcsl}(t)=\displaystyle\int_0^t\frac{ds}{\sqrt{1-s^4}}=\sum_{k\ge0}\frac{\kappa_k}{4k+1}\,t^{4k+1}$,
\emph{where} $\kappa_k=\dbinom{2k}{k}4^{-k}$.
\end{enumerate}
\emph{Proof.} (i) By Proposition~Le2\,(a) and $A_0=0$, $A_1=1$,
\begin{equation}\label{D:443}
 A_{n+2}=-2\sum_{a+b+c=n}\frac{n!}{a!\,b!\,c!}\,A_aA_bA_c ,
\end{equation}
so all $A_n$ are integers. From $\mathrm{sl}(iu)=i\,\mathrm{sl}(u)$ we get $i^nA_n=iA_n$.
(ii) Since $(1+i)^4=-4$, the function $\mathrm{sl}^\flat$ satisfies $(\mathrm{sl}^\flat)''=\frac12(\mathrm{sl}^\flat)^3$.
Write $\mathrm{sl}^\flat=\sum G_nx^n/n!$; then $G_{4k+1}=(-4)^{-k}A_{4k+1}$, $G_0=0$, $G_1=1$ and
$G_{n+2}=\frac12\sum_{a+b+c=n}\frac{n!}{a!b!c!}G_aG_bG_c$, the sum over $a,b,c\ge1$. Assume $G_a\in\Z$ and $G_a\ge0$ for
$a\le n$. The ordered triples with $a,b,c$ distinct come in sets of six with equal terms. A triple
$(a,a,c)$ with $c\ne a$ comes in three orders, each with the factor $\frac{n!}{a!a!c!}=\binom nc\binom{2a}a$,
and $\binom{2a}{a}=2\binom{2a-1}{a-1}$ is even. For $(a,a,a)$ the factor is
$\binom{3a}{a}\binom{2a}{a}$, again even. So the sum is an even integer $\ge0$, and $G_{n+2}$ is an
integer $\ge0$. The triple $(1,1,4k-3)$ shows $G_{4k+1}>0$ by induction on $k$.
(iii) By Proposition~Le2\,(b), $\mathrm{arcsl}(\mathrm{sl}(u))=u$ for small real $u$, and
$(1-s^4)^{-1/2}=\sum_k\binom{-1/2}{k}(-s^4)^k=\sum_k\kappa_ks^{4k}$. \qed

\noindent So $A_1=1$, $A_5=-12$, $A_9=3024$, $A_{13}=-4390848$, $A_{17}=21224560896$, as in \S I.3, and the
exponent of $2$ in $A_n$ is at least $(n-1)/2$. (For $n\le2001$ the least value of
$v_2(A_n)-(n-1)/2$ is $0$.\footnote{Script \texttt{el4\_divisibility.py}.})

\textbf{Lemma Le5 (Euler's addition theorem).} \emph{As an identity of meromorphic functions of
$(u,u')\in\C^2$,}
\begin{equation}\label{el:add}
 \mathrm{sl}(u+u')=\frac{\mathrm{sl}(u)\,\mathrm{sl}'(u')+\mathrm{sl}(u')\,\mathrm{sl}'(u)}{1+\mathrm{sl}(u)^2\,\mathrm{sl}(u')^2}.
\end{equation}
\emph{Proof.} Here $\mathrm{sl}'$ is the derivative of $\mathrm{sl}$, and $u'$ is a second variable. Write
$s_1=\mathrm{sl}(u)$, $s_2=\mathrm{sl}(u')$, $s_1'=\mathrm{sl}'(u)$, $s_2'=\mathrm{sl}'(u')$, and
let $N=s_1s_2'+s_2s_1'$, $D=1+s_1^2s_2^2$. By Proposition~Le2\,(a), $\partial_uN=s_1's_2'-2s_1^3s_2$ and
$\partial_uD=2s_1s_1's_2^2$. So $D^2\,\partial_u(N/D)$ equals
\begin{equation}\label{D:445}
 (s_1's_2'-2s_1^3s_2)(1+s_1^2s_2^2)-2s_1s_1's_2^2\,(s_1s_2'+s_2s_1')
 =s_1's_2'\,(1-s_1^2s_2^2)-2s_1s_2\,(s_1^2+s_2^2),
\end{equation}
where $s_1'^2=1-s_1^4$ was used once. The result is symmetric in $u$ and $u'$. So
$\partial_u(N/D)=\partial_{u'}(N/D)$, and $N/D$ is a function of $u+u'$ alone. At $u'=0$ it is
$\mathrm{sl}(u)$. \qed

\noindent This is the classical proof; compare Whittaker and Watson \cite{WW27}, \S22.2, for the Jacobi
function $\mathrm{sn}$ (cited from memory; we did not consult a copy). In that notation
$\mathrm{sl}(u)=\mathrm{sn}(u,k)$ with $k^2=-1$.

\subsection*{1a.\quad The functions $\sigma$ and $\zeta_{\rm w}$ of Weierstrass}

The identity of \S I.4.1 and the example of \S III.5 are stated with the Weierstrass function $\sigma$. The
chapter of \cite{Apo90} does not treat $\sigma$, so we define it here and prove every property that is
used later. Let $\mathcal L$ be a lattice in $\C$. One fact is taken from \cite{Apo90} (Lemma~1 of \S1.5):
the series $\sum_{\omega\in\mathcal L,\,\omega\ne0}|\omega|^{-\alpha}$ converges for $\alpha>2$. Put
\begin{equation}\label{D:446}
 \zeta_{\rm w}(u)=\frac1u+\sum_{\omega\ne0}\Big(\frac1{u-\omega}+\frac1\omega+\frac u{\omega^2}\Big),\qquad
 \sigma(u)=u\prod_{\omega\ne0}\Big(1-\frac u\omega\Big)\exp\Big(\frac u\omega+\frac{u^2}{2\omega^2}\Big),
\end{equation}
the sum and the product over the $\omega\ne0$ in $\mathcal L$. When the lattice has to be named we write
$\sigma(u;\mathcal L)$ and $\zeta_{\rm w}(u;\mathcal L)$.

\textbf{Lemma Le5a (the functions $\sigma$ and $\zeta_{\rm w}$).} \emph{Let $\mathcal L$ be a lattice.}
\begin{enumerate}[label=\textup{(\alph*)},leftmargin=2.4em,itemsep=1pt,topsep=2pt]
\item \emph{The product converges for every $u$, and $\sigma$ is an entire function. Its zeros are the
points of $\mathcal L$, and they are simple; $\sigma(u)/u\to1$ as $u\to0$. The series for $\zeta_{\rm w}$
converges for $u\notin\mathcal L$, and $\zeta_{\rm w}$ is meromorphic on $\C$; its poles are the points of
$\mathcal L$, they are simple with residue $1$, and $\zeta_{\rm w}(u)-1/u\to0$ as $u\to0$.}
\item $\sigma'/\sigma=\zeta_{\rm w}$ \emph{and} $\zeta_{\rm w}'=-\wp$. \emph{Hence $(\log\sigma)''=-\wp$.}
\item \emph{$\sigma$ and $\zeta_{\rm w}$ are odd. For every complex $\lambda\ne0$,
$\sigma(\lambda u;\lambda\mathcal L)=\lambda\,\sigma(u;\mathcal L)$ and
$\zeta_{\rm w}(\lambda u;\lambda\mathcal L)=\lambda^{-1}\zeta_{\rm w}(u;\mathcal L)$. In particular, if
$i\mathcal L=\mathcal L$ then $\sigma(iu)=i\,\sigma(u)$ and $\zeta_{\rm w}(iu)=-i\,\zeta_{\rm w}(u)$.}
\item \emph{For every $\omega\in\mathcal L$ there is a constant $\eta(\omega)$ with
$\zeta_{\rm w}(u+\omega)=\zeta_{\rm w}(u)+\eta(\omega)$, and}
\begin{equation}\label{el:sigmaqp}
 \sigma(u+\omega)=C_\omega\,e^{\eta(\omega)u}\,\sigma(u)\qquad\text{for a constant }C_\omega\ne0 .
\end{equation}
\emph{If $\omega/2\notin\mathcal L$ then $C_\omega=-e^{\eta(\omega)\omega/2}$.}
\item \emph{(Legendre's relation.) If $\mathcal L=\Z\omega_1+\Z\omega_2$ with
$\mathrm{Im}(\omega_2/\omega_1)>0$, then $\eta(\omega_1)\,\omega_2-\eta(\omega_2)\,\omega_1=2\pi i$.}
\item \emph{For the square lattice $\mathcal P=\varpi\Z[i]$ of Lemma~Le1: $\eta(\varpi)=\pi/\varpi$,
$\eta(i\varpi)=-i\pi/\varpi$, and}
\begin{equation}\label{el:sigmasq}
 \sigma(u+\varpi)=-e^{\pi/2}\,e^{\pi u/\varpi}\,\sigma(u),\qquad
 \sigma(u+i\varpi)=-e^{\pi/2}\,e^{-i\pi u/\varpi}\,\sigma(u).
\end{equation}
\end{enumerate}

\emph{Proof.} (a) Fix $R>0$ and let $|u|\le R$. Only finitely many $\omega\in\mathcal L$ have $|\omega|<2R$.
For the others $|u-\omega|\ge\frac12|\omega|$, and
\begin{equation}\label{D:449}
 \Big|\frac1{u-\omega}+\frac1\omega+\frac u{\omega^2}\Big|=\Big|\frac{u^2}{\omega^2(u-\omega)}\Big|\le\frac{2R^2}{|\omega|^3}.
\end{equation}
For $|x|\le\frac12$ we have $(1-x)\exp(x+\frac12x^2)=\exp\big(-\sum_{k\ge3}x^k/k\big)$, and
$\big|\sum_{k\ge3}x^k/k\big|\le\frac13|x|^3\sum_{k\ge0}2^{-k}\le|x|^3$. With $x=u/\omega$ this is at most
$R^3/|\omega|^3$. Since $\sum|\omega|^{-3}<\infty$, the sum over the $\omega$ with $|\omega|\ge2R$ of the terms
of $\zeta_{\rm w}$, and the sum of the logarithms of the factors of $\sigma$, converge absolutely and
uniformly in $|u|\le R$. A uniformly convergent series of holomorphic functions has a holomorphic sum. So
in $|u|<R$ the function $\sigma$ is the product of $u$, of the finitely many factors with $|\omega|<2R$, and
of the exponential of a holomorphic function; and $\zeta_{\rm w}$ is the sum of $1/u$, of finitely many
terms with a simple pole of residue $1$ at one point $\omega$, and of a holomorphic function. Since $R$ is
arbitrary, this gives the statements on convergence, on the zeros and on the poles. For the two limits:
$\sigma(u)/u$ is the product without the factor $u$, it is holomorphic by what was said, hence continuous at
$u=0$, and every factor of it is $1$ at $u=0$; so $\sigma(u)/u\to1$. In the same way
$\zeta_{\rm w}(u)-1/u$ is holomorphic near $u=0$, and each term of its series vanishes at $u=0$; so
$\zeta_{\rm w}(u)-1/u\to0$.

(b) A uniformly convergent series of holomorphic functions may be differentiated term by term. The
logarithmic derivative of the factor $(1-u/\omega)\exp(u/\omega+u^2/2\omega^2)$ is
$\frac1{u-\omega}+\frac1\omega+\frac u{\omega^2}$, and the derivative of this is
$-\frac1{(u-\omega)^2}+\frac1{\omega^2}$. Compare with the series for $\zeta_{\rm w}$ and for $\wp$.

(c) The terms may be rearranged, because the convergence is absolute. Replacing $\omega$ by $-\omega$ in
the product for $\sigma(-u)$ gives $-\sigma(u)$; in the same way $\zeta_{\rm w}(-u)=-\zeta_{\rm w}(u)$. The
points of $\lambda\mathcal L$ are the $\lambda\omega$, and $\lambda u/\lambda\omega=u/\omega$; this gives
the two rules for $\lambda$. For $\lambda=i$ and $i\mathcal L=\mathcal L$ they are the last statement.

(d) The function $\zeta_{\rm w}(u+\omega)-\zeta_{\rm w}(u)$ has the derivative $-\wp(u+\omega)+\wp(u)=0$ on
the connected set $\C\setminus\mathcal L$. So it is a constant $\eta(\omega)$. By (a) the quotient
$\sigma(u+\omega)/\sigma(u)$ is an entire function without zeros, and by (b) its logarithmic derivative is
$\eta(\omega)$. So $\sigma(u+\omega)\,e^{-\eta(\omega)u}/\sigma(u)$ has derivative $0$; it is a constant
$C_\omega\ne0$. Let $\omega/2\notin\mathcal L$ and put $u=-\omega/2$ in \eqref{el:sigmaqp}. Since $\sigma$ is
odd, $\sigma(\omega/2)=-C_\omega e^{-\eta(\omega)\omega/2}\sigma(\omega/2)$, and $\sigma(\omega/2)\ne0$ by (a).

(e) Choose $\alpha\in\C$ such that the parallelogram with the corners $\alpha$, $\alpha+\omega_1$,
$\alpha+\omega_1+\omega_2$, $\alpha+\omega_2$ has no point of $\mathcal L$ on its boundary. (Such $\alpha$
exist: a point $\omega\in\mathcal L$ lies on a side only if $\alpha$ lies on one of four lines through
$\omega$, $\omega-\omega_1$ or $\omega-\omega_2$ in the directions $\omega_1$, $\omega_2$; these are countably many lines
in all, and they do not cover the plane.) It contains
exactly one point of $\mathcal L$ in its interior (every complex number is congruent modulo $\mathcal L$
to exactly one point $\alpha+s\omega_1+t\omega_2$ with $0\le s,t<1$), and the corners in this order go round it in the
positive sense, because $\mathrm{Im}(\omega_2/\omega_1)>0$. By (a) and the residue theorem the integral of
$\zeta_{\rm w}$ over the boundary is $2\pi i$. The first and the third side together give
$\int_\alpha^{\alpha+\omega_1}\big(\zeta_{\rm w}(u)-\zeta_{\rm w}(u+\omega_2)\big)du=-\eta(\omega_2)\,\omega_1$,
and the second and the fourth give
$\int_\alpha^{\alpha+\omega_2}\big(\zeta_{\rm w}(u+\omega_1)-\zeta_{\rm w}(u)\big)du=\eta(\omega_1)\,\omega_2$.

(f) By (c), $\zeta_{\rm w}(iu+i\varpi)=-i\,\zeta_{\rm w}(u+\varpi)=\zeta_{\rm w}(iu)-i\,\eta(\varpi)$. So
$\eta(i\varpi)=-i\,\eta(\varpi)$. Legendre's relation with $\omega_1=\varpi$, $\omega_2=i\varpi$ reads
$2i\varpi\,\eta(\varpi)=2\pi i$. So $\eta(\varpi)=\pi/\varpi$ and $\eta(i\varpi)=-i\pi/\varpi$. Now (d) gives
$C_\varpi=-e^{\pi/2}$ and $C_{i\varpi}=-e^{(-i\pi/\varpi)(i\varpi/2)}=-e^{\pi/2}$. \qed

\noindent \emph{Example and checks.} For $\mathcal P$ one finds $\eta(\varpi)=\pi/\varpi=1.19814\ldots$ The
truncated sum for $\zeta_{\rm w}$ over the $\omega=\varpi(a+bi)$ with $|a|,|b|\le120$ gives
$\zeta_{\rm w}(u+\varpi)-\zeta_{\rm w}(u)=1.19815$ at $u=0.37+0.21i$, and the truncated product for $\sigma$
satisfies \eqref{el:sigmasq} within $3\cdot10^{-5}$. An independent formula for $\sigma$ by a rapidly
convergent series satisfies every statement of the lemma to $38$ digits; with the sign $+$ in
\eqref{el:sigmasq}, or without the factor $e^{\pi/2}$, the test fails.\footnote{Script
\texttt{el9\_sigma\_zeta.py}.} The lemma is used as follows: (a), (b), (c) in Theorem~Le6; (a), (b), (c), (d)
in Proposition~Le21; the rule for $\lambda$ in (c) in Proposition~Le22. Parts (e) and (f) make the
factors in \eqref{el:sigmaqp} explicit; they give a check of Theorem~Le6 below.

\subsection*{2.\quad The theta series of the Gaussian integers}

Recall from \S I.4.1 the series $\Theta(z)=\sum_{\lambda\in\Z[i]}\exp(-\pi|\lambda|^2+\lambda z)$ and its
derivatives $T_n=\Theta^{(n)}(0)$. We now prove the identity stated there. From now on $\sigma$ and
$\zeta_{\rm w}$ are the functions of \S1a for the lattice $\mathcal P$.
Put $J(u)=\int_0^u\mathrm{sl}(s)^2\,ds$ and $M(u)=\int_0^uJ(s)\,ds$, power series in $u$ that begin with
$u^3/3$ and $u^4/12$.

\textbf{Theorem Le6 (the identity of \S I.4.1).} \emph{Let $\Phi=\sigma^2\wp$. It is an entire function,
because the double zeros of $\sigma^2$ are at the double poles of $\wp$ (Lemma~Le5a\,(a)). Then}
\begin{equation}\label{el:theta}
 \Phi(u)=\Big(\frac{\sigma(u)}{\mathrm{sl}(u)}\Big)^2=\exp\big(2M(u)\big)\quad(\text{$u$ near }0),\qquad
 \Theta(z)=T_0\,\sigma(u)^2\,\wp(u)\quad\text{with }u=\frac{\varpi z}{2\pi}\quad(z\in\C).
\end{equation}

\emph{Proof.} \emph{The function $\Phi$.} By \eqref{el:sldef}, $\Phi=(\sigma/\mathrm{sl})^2$. It is even, and
$\Phi(0)=1$, by Lemma~Le5a\,(a) and~(c). Since $(\log\sigma)''=-\wp=-\mathrm{sl}^{-2}$ by Lemma~Le5a\,(b)
and, by Proposition~Le2\,(a),
\begin{equation}\label{el:logsl}
 (\log\mathrm{sl})''=\frac{\mathrm{sl}''\,\mathrm{sl}-\mathrm{sl}'^2}{\mathrm{sl}^2}
 =\frac{-2\,\mathrm{sl}^4-(1-\mathrm{sl}^4)}{\mathrm{sl}^2}=-\mathrm{sl}^2-\mathrm{sl}^{-2},
\end{equation}
we get $(\log\Phi)''=2\,\mathrm{sl}^2$. With $\log\Phi(0)=0=(\log\Phi)'(0)$ this is $\log\Phi=2M$; in
particular $\Phi(u)=1+O(u^4)$.

\emph{The theta series as a product.} The double series converges absolutely, and with $\lambda=a+bi$ it
factors:
\begin{equation}\label{D:452}
 \Theta(z)=\theta(z)\,\theta(iz),\qquad \theta(z)=\sum_{a\in\Z}e^{-\pi a^2+az}.
\end{equation}
Replacing $a$ by $a+1$ in the sum gives $\theta(z+2\pi)=e^{\pi+z}\theta(z)$; clearly
$\theta(z+2\pi i)=\theta(z)$. Hence
\begin{equation}\label{el:thetaqp}
 \Theta(z+2\pi)=e^{\pi+z}\,\Theta(z),\qquad \Theta(z+2\pi i)=e^{\pi-iz}\,\Theta(z).
\end{equation}
\emph{The zeros of $\theta$.} Let $Q$ be a square with corners $\alpha,\alpha+2\pi,\alpha+2\pi+2\pi i,\alpha+2\pi i$
and no zero of $\theta$ on its boundary. The function $\theta'/\theta$ has the period $2\pi i$, and
$(\theta'/\theta)(z+2\pi)=(\theta'/\theta)(z)+1$. In $\frac1{2\pi i}\oint_{\partial Q}\theta'/\theta$ the
horizontal sides cancel, and the two vertical sides together give $\frac1{2\pi i}\int_\alpha^{\alpha+2\pi i}1\,dz=1$.
So $\theta$ has exactly one zero in $Q$, and it is simple. The point $z_0=\pi(-1+i)$ is a zero: in
$\theta(z_0)=\sum_a(-1)^ae^{-\pi a^2-\pi a}$ the substitution $a\mapsto-1-a$ changes the sign of each term.
So the zeros of $\theta$ are the points of $z_0+2\pi\Z[i]$. The zeros of $\theta(iz)$ are the points
$-iz_0+2\pi\Z[i]$, and $-iz_0=\pi(1+i)=z_0+2\pi$. Therefore $\Theta$ has a double zero at each point of
$\pi(1+i)+2\pi\Z[i]$ and no other zero. In the variable $u=\varpi z/2\pi$ these are the points of
$u_0+\mathcal P$: the poles of $\mathrm{sl}$.

\emph{Comparison.} Let $\Psi(u)=\Theta(2\pi u/\varpi)$. By \eqref{el:thetaqp}, $\log\Psi$ changes by a
polynomial of degree one in $u$ under $u\mapsto u+\varpi$ and under $u\mapsto u+i\varpi$. So $(\log\Psi)''$
is an elliptic function for $\mathcal P$. Its poles are at the double zeros of $\Psi$, with principal part
$-2/(u-u_0)^2$ and no term in $1/(u-u_0)$. By Proposition~Le2\,(e), $\mathrm{sl}(u_0+x)=-i/\mathrm{sl}(x)$, so
$2\,\mathrm{sl}(u_0+x)^2=-2x^{-2}+O(x^2)$: the same principal part. Also $\mathrm{sl}^2$ is elliptic for
$\mathcal P$ by Proposition~Le2\,(c). So $(\log\Psi)''-2\,\mathrm{sl}^2$ is an elliptic function without poles: a
constant. Now $\Theta$ is even, and $T_2=\sum\lambda^2e^{-\pi|\lambda|^2}=0$ because $\lambda\mapsto i\lambda$
changes its sign. So $(\log\Psi)''(0)=0=2\,\mathrm{sl}(0)^2$ and the constant is $0$. Hence
$(\log\Psi)''=(\log\Phi)''$. Both $\Psi$ and $\Phi$ are even, so $\log\Psi-\log\Phi$ is a constant near $0$;
at $u=0$ it is $\log T_0$. Thus $\Psi=T_0\Phi$ near $0$, and so everywhere, both sides being entire. \qed

\noindent The identity \eqref{el:theta} was checked at three complex points, to $60$, $52$ and $46$ digits
(the last two are limited by the truncated series used for $\sigma$ and $\wp$), with the wrong scaling
$u=\varpi z/\pi$ as a control.\footnote{Script \texttt{el1\_theta\_identity.py}.} There is also a check by
Lemma~Le5a\,(f): by \eqref{el:sigmasq} and the periodicity of $\wp$,
$\Phi(u+\varpi)=e^{\pi+2\pi u/\varpi}\Phi(u)$ and $\Phi(u+i\varpi)=e^{\pi-2\pi iu/\varpi}\Phi(u)$, and with
$z=2\pi u/\varpi$ these are the two factors of \eqref{el:thetaqp}.

The theorem gives $(\sigma/\mathrm{sl})^2$ as a product $\theta(z)\theta(iz)$. To take the square root one
needs to know how $\theta(iz)$ and $\theta(z)$ are related.

\textbf{Corollary Le6a (the functional equation of $\theta$; a square root of \eqref{el:theta}).}
\emph{For all $z\in\C$, and with $u=\varpi z/2\pi$,}
\begin{equation}\label{el:thetafe}
 \theta(iz)=e^{-z^2/4\pi}\,\theta(z),\qquad\qquad
 \frac{\sigma(u)}{\mathrm{sl}(u)}=e^{-z^2/8\pi}\,\frac{\theta(z)}{\theta(0)} .
\end{equation}

\emph{Proof.} In the proof of Theorem~Le6 we saw that $\theta(z)$ and $\theta(iz)$ have the same zeros, the
points of $\pi(1+i)+2\pi\Z[i]$, and that these are simple for both. So
$G(z)=e^{z^2/4\pi}\,\theta(iz)/\theta(z)$ is an entire function. We also saw that
$\theta(z+2\pi)=e^{\pi+z}\theta(z)$ and $\theta(z+2\pi i)=\theta(z)$. The first of these, applied at the
point $iz-2\pi$, gives $\theta(iz-2\pi)=e^{\pi-iz}\theta(iz)$. Since
$e^{(z+2\pi)^2/4\pi}=e^{z^2/4\pi}e^{z+\pi}$ and $e^{(z+2\pi i)^2/4\pi}=e^{z^2/4\pi}e^{iz-\pi}$, we get
\begin{equation}\label{D:455}
 G(z+2\pi)=\frac{e^{z^2/4\pi}e^{z+\pi}\,\theta(iz)}{e^{\pi+z}\,\theta(z)}=G(z),\qquad
 G(z+2\pi i)=\frac{e^{z^2/4\pi}e^{iz-\pi}\,e^{\pi-iz}\,\theta(iz)}{\theta(z)}=G(z).
\end{equation}
An entire function with the two periods $2\pi$ and $2\pi i$ is bounded, so it is constant by Liouville's
theorem (or by \cite{Apo90}, Theorem~1.4). Since $G(0)=1$, this is the first equation. For the second, by
\eqref{el:theta} with $\Theta(z)=\theta(z)\theta(iz)$ and $T_0=\theta(0)^2$,
\begin{equation}\label{D:456}
 \Big(\frac{\sigma(u)}{\mathrm{sl}(u)}\Big)^2=\frac{\theta(z)\,\theta(iz)}{\theta(0)^2}
 =\Big(e^{-z^2/8\pi}\,\frac{\theta(z)}{\theta(0)}\Big)^2 .
\end{equation}
The quotient of the two meromorphic functions $\sigma(u)/\mathrm{sl}(u)$ and
$e^{-z^2/8\pi}\theta(z)/\theta(0)$ has the square $1$, so it is a constant, $1$ or $-1$. Both functions
tend to $1$ as $z\to0$. \qed

\noindent The first equation of \eqref{el:thetafe} is a case of Jacobi's transformation formula for the
theta series. It is usually proved by Poisson summation; the proof above uses only the zeros of $\theta$.
Both equations were checked at complex points with $|z|\le6.1$, to $40$ and to $30$ digits; with $e^{+z^2/4\pi}$
in the first, or $e^{-z^2/4\pi}$ in the second, the test fails.\footnote{Script
\texttt{el9\_sigma\_zeta.py}.} The corollary is not used in the rest of the section.

\textbf{Proposition Le7 (the two values of \S I.4.1).}
\begin{equation}\label{el:T0}
 T_0=\sum_{n\in\Z}\frac1{\cosh\pi n}=\frac{\sqrt2\,\varpi}{\pi}=1.18034\ldots,\qquad
 \frac{T_4}{T_0}=\frac14\Big(\frac\varpi\pi\Big)^4=0.12131\ldots
\end{equation}

\emph{Proof.} \emph{The second value.} By \eqref{el:theta}, $\Phi(u)=\exp(2M(u))=1+u^4/6+O(u^8)$, so
$T_4=4!\,T_0\,(\varpi/2\pi)^4/6$.

\emph{The first equality.} $T_0=\sum_{N\ge0}r_2(N)e^{-\pi N}$, where $r_2(N)$ is the number of
$\lambda\in\Z[i]$ with $|\lambda|^2=N$. By the two-squares theorem (C2),
$r_2(N)=4\sum_{d\mid N}\chi_4(d)$ for $N\ge1$, where $\chi_4$ is the nonprincipal character modulo $4$. With
$x=e^{-\pi}$ and $\sum_{d\ge1}\chi_4(d)y^d=y/(1+y^2)$,
\begin{equation}\label{D:458}
 T_0=1+4\sum_{d,m\ge1}\chi_4(d)x^{dm}=1+4\sum_{m\ge1}\frac{x^m}{1+x^{2m}}=1+2\sum_{m\ge1}\frac1{\cosh\pi m}.
\end{equation}
\emph{The second equality.} Let $S(u)=\sum_{n\in\Z}1/\cosh\big(\pi(u+n\varpi)/\varpi\big)$. The series
converges uniformly on compact sets away from its poles. It has the period $\varpi$, and
$S(u+i\varpi)=-S(u)$ because $\cosh(y+i\pi)=-\cosh y$. Its poles are the points of
$\frac12i\varpi+\mathcal P$; at $u=\frac12i\varpi+x$ the term $n=0$ is $1/(i\sinh(\pi x/\varpi))$, so the
principal part is $-i\varpi/(\pi x)$. Hence $S^2$ is elliptic for $\mathcal P$ with a double pole of
leading coefficient $-(\varpi/\pi)^2$ at $\frac12i\varpi$. An elliptic function with one pole in a period
parallelogram has residue $0$ there (\cite{Apo90}, Theorem~1.7). So $S(u)^2+(\varpi/\pi)^2\wp(u-\frac12i\varpi)$ has no poles and is a
constant $C$. Next, $S(\frac12i\varpi+x)=\sum_n1/(i\sinh(\pi(x+n\varpi)/\varpi))$ is odd in $x$ and has the
period $\varpi$ in $x$; so it vanishes at $x=\varpi/2$. At that point
$\wp(u-\frac12i\varpi)=\wp(\varpi/2)=1$, so $C=(\varpi/\pi)^2$. At $u=0$, by \eqref{el:evalues},
\begin{equation}\label{D:459}
 S(0)^2=C-(\varpi/\pi)^2\,\wp(-\tfrac12i\varpi)=2\,(\varpi/\pi)^2,
\end{equation}
and $S(0)>0$. \qed

\noindent Apart from the statements of \cite{Apo90} quoted in \S1 and \S1a, the two-squares theorem (C2) is
the only cited input of \S\S1--3. Numerically, $T_0$,
$\sum_n1/\cosh\pi n$ and $\sqrt2\,\varpi/\pi$ agree to $60$ digits, as do the two sides of
$S^2=(\varpi/\pi)^2(1-\wp(u-\frac12i\varpi))$ at a complex point.\footnote{Script
\texttt{el1\_theta\_identity.py}; $\varpi$ is computed there as $\pi$ divided by the arithmetic-geometric
mean of $1$ and $\sqrt2$.}

\textbf{Corollary Le8 (the rational numbers $r_k$).} \emph{Write $\Phi(u)=\sum_{k\ge0}\phi_ku^{4k}$. Then
$(4k)!\,\phi_k$ is an integer divisible by $4^k$, and}
\begin{equation}\label{el:rk}
 r_k=\frac{T_{4k}}{T_0\,(\varpi/\pi)^{4k}}=\frac{(4k)!\,\phi_k}{16^k}\in\frac1{4^k}\,\Z .
\end{equation}
\emph{Proof.} The equality is \eqref{el:theta}: $T_{4k}=(4k)!\,T_0\,\phi_k\,(\varpi/2\pi)^{4k}$. Let
$\mu=1+i$, so $\mu^4=-4$ and $\mathrm{sl}(u)=\mu^{-1}\mathrm{sl}^\flat(\mu u)$. Then
$2M(u)=E_1(\mu u)$ with $E_1(x)=-\frac12\int_0^x\!\!\int_0^y\mathrm{sl}^\flat(s)^2\,ds\,dy$. By
Lemma~Le4\,(ii) and Lemma~Le3, $\mathrm{sl}^\flat\cdot(\mathrm{sl}^\flat)'$ is a Hurwitz series over $\Z$,
so $(\mathrm{sl}^\flat)^2=2\int_0^x\mathrm{sl}^\flat(\mathrm{sl}^\flat)'$ is twice a Hurwitz series over $\Z$,
and $E_1$ is a Hurwitz series over $\Z$ with $E_1(0)=0$. By Lemma~Le3, $\exp(E_1)=\sum_ne_nx^n/n!$ with
$e_n\in\Z$. So $\Phi(u)=\exp(E_1(\mu u))$ gives $(4k)!\,\phi_k=e_{4k}\mu^{4k}=(-4)^ke_{4k}$. \qed

\noindent So $r_k=(-1)^ke_{4k}/4^k$, and the proof is a recipe for computing it. We computed $\Phi$
exactly in two ways, from $\mathrm{sl}$ as here and from the Laurent coefficients of
$\wp$ (\cite{Apo90}, Theorem~1.13, with $g_2=4$, $g_3=0$); the results agree up to
$u^{480}$.\footnote{Scripts \texttt{el2\_rk.py} and, for $r_{11}$ and $r_{12}$, \texttt{el10\_examples.py}.}
The values are
\begin{equation}\label{el:rktable}
\begin{aligned}
 &r_1=\tfrac14,\quad r_2=\tfrac{17}{16},\quad r_3=-\tfrac{111}{64},\quad r_4=\tfrac{12513}{256},\quad
 r_5=-\tfrac{146079}{1024},\quad r_6=\tfrac{81727857}{4096},\quad r_7=\tfrac{22284691569}{16384},\\
 &r_8=-\tfrac{751645880127}{65536},\quad r_9=-\tfrac{305713085239359}{262144},\quad
 r_{10}=\tfrac{274920417491193297}{1048576},\\
 &r_{11}=-\tfrac{12320223978927877551}{4^{11}},\quad r_{12}=\tfrac{12397456849784472393633}{4^{12}}.
\end{aligned}
\end{equation}
The first five are the values of \S I.4.1. For $k\le10$ they agree with the lattice sums to $59$ digits.
For $k\le120$ the denominator of $r_k$ is exactly $4^k$; that it is never smaller is observed, not
proved. The exponent of $3$ in $r_{30}$ is $14$.

\subsection*{3.\quad The size of $A_n$}

\textbf{Theorem Le9 (size through a period).} \emph{Let $n=4k+1$.}
\begin{enumerate}[label=\textup{(\roman*)},leftmargin=2.4em,itemsep=1pt,topsep=2pt]
\item \emph{For every constant $1<C_0<\sqrt5$,}
\begin{equation}\label{el:size}
 \frac{A_n}{n!}=(-1)^k\,4\,\Big(\frac{\sqrt2}{\varpi}\Big)^{n+1}\big(1+O(C_0^{-n})\big),\qquad
 \frac{\varpi}{\sqrt2}=1.8540746\ldots
\end{equation}
\emph{Consequently $\log|A_n|=\log n!-(n+1)\log(\varpi/\sqrt2)+\log4+O(C_0^{-n})$.}
\item \emph{Let $\delta_n$ be the relative error in \eqref{el:size}, that is
$A_n/n!=(-1)^k\,4\,(\sqrt2/\varpi)^{n+1}(1+\delta_n)$. For every constant $1<C_1<\sqrt{13}$,}
\begin{equation}\label{el:size2}
 \delta_n=-2\cdot5^{-(n+1)/2}\cos\big((n+1)\arctan\tfrac12\big)+3^{-(n+1)}+O(C_1^{-n}) .
\end{equation}
\emph{Consequently $|\delta_n|\le2\cdot5^{-(n+1)/2}+O(3^{-n})$. The cosine in \eqref{el:size2} is never $1$
or $-1$, so the first term of \eqref{el:size2} is smaller than $2\cdot5^{-(n+1)/2}$ in absolute value for
every $n$: the constant $2$ is not attained. It cannot be replaced by a smaller constant.}
\end{enumerate}

\emph{Proof.} By Proposition~Le2\,(d) and~(e) the poles of $\mathrm{sl}$ are the points
$P=u_0+(a+bi)\varpi=\frac12\varpi\,((2a+1)+(2b+1)i)$ with integers $a,b$; they are simple, and the residue at
$P$ is $-i\,s(P)$ with $s(P)=(-1)^{a+b}$. The distance of $P$ from $0$ is
$\frac12\varpi\sqrt{(2a+1)^2+(2b+1)^2}$. The sums of two odd squares are, in order, $2,10,18,26,\ldots$, and
the poles on the first four circles are these: the four points $i^ju_0$ ($j=0,1,2,3$) for the sum $2$; the
eight points $i^ju_1$ and $i^ju_2$ for the sum $10$, where $u_1=\frac12\varpi(1+3i)$ and
$u_2=\frac12\varpi(3+i)$; the four points $3\,i^ju_0$ for the sum $18$; and the eight points
$\frac12\varpi(\pm1\pm5i)$, $\frac12\varpi(\pm5\pm i)$ for the sum $26$. So the numbers of poles on these
circles are $4,8,4,8$, and their radii are $\varpi/\sqrt2$ times $1,\sqrt5,3,\sqrt{13}$. If $P$ is a pole
then so is $iP$, and $s(iP)=-s(P)$: for $P$ as above, $iP=u_0+((-b-1)+ai)\varpi$.

Let $\varrho$ be one of the radii $\frac12\varpi\sqrt{10}$, $\frac12\varpi\sqrt{26}$. The function
\begin{equation}\label{D:464}
 \mathrm{sl}(u)+\sum_{|P|<\varrho}\frac{i\,s(P)}{u-P},
\end{equation}
the sum over the poles $P$ with $|P|<\varrho$, is holomorphic in $|u|<\varrho$, and by Cauchy's inequality its
coefficient of $u^n$ is $O(r^{-n})$ for every $r<\varrho$. The coefficient of $u^n$ in $-i\,s(P)/(u-P)$ is
$i\,s(P)\,P^{-n-1}$. For $n=4k+1$ one has $i^{-(n+1)}=-1$, so $s(iP)(iP)^{-n-1}=s(P)P^{-n-1}$: the four poles
$i^jP$ contribute the same, together $4i\,s(P)P^{-n-1}$.

(i) Take $\varrho=\frac12\varpi\sqrt{10}$. The only poles with $|P|<\varrho$ are the $i^ju_0$, and $s(u_0)=1$. So
$A_n/n!=4i\,u_0^{-n-1}+O(r^{-n})$ for every $r<\frac12\varpi\sqrt{10}$. Now
$u_0=(\varpi/\sqrt2)e^{i\pi/4}$ and $e^{-i\pi(4k+2)/4}=-i(-1)^k$, so
$4i\,u_0^{-n-1}=(-1)^k4(\sqrt2/\varpi)^{n+1}$, and $r\sqrt2/\varpi$ is any number below $\sqrt5$.

(ii) Take $\varrho=\frac12\varpi\sqrt{26}$. The poles with $|P|<\varrho$ are the sixteen points $i^ju_0$,
$i^ju_1$, $i^ju_2$, $3\,i^ju_0$, with $s(u_0)=1$, $s(u_1)=s(u_2)=-1$ and $s(3u_0)=1$ (here $(a,b)$ is
$(0,0)$, $(0,1)$, $(1,0)$, $(1,1)$). So
\begin{equation}\label{D:465}
 \frac{A_n}{n!}=4i\,u_0^{-n-1}\Big(1-\Big(\frac{u_0}{u_1}\Big)^{n+1}-\Big(\frac{u_0}{u_2}\Big)^{n+1}+3^{-(n+1)}\Big)+O(r^{-n})
\end{equation}
for every $r<\frac12\varpi\sqrt{26}$. Here $u_0/u_1=(1+i)/(1+3i)=(2-i)/5$ and $u_0/u_2=(1+i)/(3+i)=(2+i)/5$,
and $(2\pm i)/5=5^{-1/2}e^{\pm i\vartheta}$ with $\vartheta=\arctan\frac12$. So the two middle terms add up
to $-2\cdot5^{-(n+1)/2}\cos((n+1)\vartheta)$. This is \eqref{el:size2}, and the bound for $|\delta_n|$
follows.

Suppose that $m\vartheta$ is a multiple of $\pi$ for an integer $m\ge1$. Then
$\big((2+i)/(2-i)\big)^m=e^{2im\vartheta}=1$, that is $(3+4i)^m=5^m$. But $(3+4i)^2=(3+4i)+5(-2+4i)$, so by
induction $(3+4i)^m\equiv3+4i$ modulo $5\Z[i]$ for all $m\ge1$, and $3+4i\notin5\Z[i]$. So $\vartheta/\pi$ is
irrational. In particular $\cos((n+1)\vartheta)\ne\pm1$. And the numbers $(n+1)\vartheta=2\vartheta+k\cdot4\vartheta$ come
arbitrarily close to multiples of $\pi$, because the multiples of an irrational number come arbitrarily
close, modulo $1$, to every real number (by the pigeonhole principle). For such $n$ the first term of
\eqref{el:size2} is close to $\pm2\cdot5^{-(n+1)/2}$, and the other terms are $O(3^{-n})$. So no constant
smaller than $2$ will do. \qed

\noindent Part (i) is the statement of \S I.3, defect (D3), with $R=\varpi/\sqrt2$. Part (ii) shows that
the error is not of the exact order $5^{-(n+1)/2}$: the cosine can be small. For
$n=5,9,13,17,41,101$ the relative error $\delta_n$ is $1.6\cdot10^{-2}$, $7.1\cdot10^{-5}$, $-2.5\cdot10^{-5}$,
$4.9\cdot10^{-7}$, $-3.4\cdot10^{-15}$, $4.4\cdot10^{-36}$; multiplied by $5^{(n+1)/2}$
these are, in absolute value, $1.94$, $0.22$, $1.94$, $0.95$, $1.62$, $1.97$, and for all $n\equiv1\bmod4$ with $5\le n\le401$ this
product lies between $0.0070$ (at $n=165$) and $1.9998$ (at $n=121$); it is less than $2$ throughout this
range.\footnote{Scripts \texttt{el6\_inverse\_log.py} and
\texttt{el10\_examples.py}. The second also checks that the error differs from the contribution of the eight
poles by less than $3^{-n}$ in this range. A separate computation made for this revision, from the exact
$A_n$, gave the two extreme values and $|\delta_n-(\text{the first two terms of \eqref{el:size2}})|\le2.2\cdot13^{-(n+1)/2}$
for $5\le n\le401$.} Each derivative costs one factor $\sqrt2/\varpi$: the reciprocal
of the distance to the nearest pole, which is a period divided by $\sqrt2$.

\subsection*{4.\quad Formal groups}

The addition theorem \eqref{el:add} expresses $\mathrm{sl}(u+u')$ through $\mathrm{sl}(u)$ and
$\mathrm{sl}(u')$ by power series. The notion that records this is the following.

\textbf{Definition.} Let $R$ be a subring of a field $K$ of characteristic zero. A \emph{formal group law}
over $R$ is a power series $F(X,Y)\in R[[X,Y]]$ with
\begin{equation}\label{D:466}
 F(X,0)=X,\qquad F(X,Y)=F(Y,X),\qquad F(X,F(Y,Z))=F(F(X,Y),Z).
\end{equation}
A \emph{logarithm} of $F$ is a series $l(X)=X+\cdots\in K[[X]]$ with $l(F(X,Y))=l(X)+l(Y)$; its inverse
series $e$ is the \emph{exponential}, and $F(X,Y)=e(l(X)+l(Y))$. Conversely, for any $l(X)=X+\cdots\in K[[X]]$
with inverse $e$, the series $e(l(X)+l(Y))$ satisfies the three identities; it is a formal group law over
$K$, and the question is whether its coefficients lie in $R$. For a formal group law $F$ with a logarithm
$l$ and an integer $m$ put $[m](X)=e(m\,l(X))$:
the multiplication by $m$. If $F$ is a formal group law over $R$, then $[m](X)$ has coefficients in $R$ for
every $m\ge0$, by induction from $[0](X)=0$ and $[m+1](X)=F([m](X),X)$. The series $[m]$, and the height
defined below, are defined here through a logarithm, so they presuppose one. Every formal group law of this
section is given together with its logarithm. (A formal group law over a field of characteristic zero always
has a logarithm; this is not proved here and not needed.)
Two formal group laws $F,F'$ over $R$ with logarithms $l,l'$ are \emph{strictly isomorphic over $R$} if the
series $\iota=l'^{-1}\circ l$, which satisfies $\iota(F(X,Y))=F'(\iota(X),\iota(Y))$ and $\iota(X)=X+\cdots$,
has coefficients in $R$; then $\iota^{-1}$ has coefficients in $R$ as well. Let $p$ be a prime and let $F$
be a formal group law over $R=\Z_{(p)}$ that has a logarithm, so that $[p](X)\in\Z_{(p)}[[X]]$. If
$[p](X)\equiv aX^{p^h}+(\text{higher terms})$ modulo $p$ with $a\not\equiv0$, then $h$ is the \emph{height}
of $F$ at $p$.

\textbf{Example (the multiplicative group).} $F(X,Y)=X+Y+XY=(1+X)(1+Y)-1$ is a formal group law over $\Z$,
with $l(X)=\log(1+X)=\sum_{n\ge1}(-1)^{n-1}X^n/n$ and $e(z)=e^z-1=\sum_{n\ge1}z^n/n!$. Here
$[p](X)=(1+X)^p-1\equiv X^p$ modulo $p$: the height is $1$ at every prime. The logarithm has the
denominator $p^j$ at the degree $p^j$, and the exponential has the denominators $n!$, with
$v_p(n!)=\sum_{h\ge1}\lfloor n/p^h\rfloor$ by Legendre's formula: the factorial of \S I.2. Multiplying the
coefficient of $z^n$ by $n!$, as an $n$th derivative does, gives $1$. There is no divisibility left.

\textbf{Example (strict isomorphism).} The additive law $F'(X,Y)=X+Y$ has the logarithm $l'(X)=X$. For the
multiplicative law $F$ of the last example, $\iota=l'^{-1}\circ l$ is the series
$\log(1+X)=X-\frac12X^2+\frac13X^3-\cdots$. So $F$ and $F'$ are strictly isomorphic over $\Q$, but not over
$\Z$ and not over $\Z_{(p)}$ for any prime $p$, because of the coefficient $\pm1/p$.

\textbf{Proposition Le10 (the formal group of the lemniscate).} \emph{The series}
\begin{equation}\label{el:F}
 F(X,Y)=\frac{X\sqrt{1-Y^4}+Y\sqrt{1-X^4}}{1+X^2Y^2},\qquad \sqrt{1-X^4}=\sum_{k\ge0}\binom{1/2}{k}(-X^4)^k,
\end{equation}
\emph{is a formal group law over $\Z[\frac12]$. Its logarithm is $\mathrm{arcsl}$ and its exponential is
$\mathrm{sl}$: $\ \mathrm{sl}(u+u')=F(\mathrm{sl}\,u,\mathrm{sl}\,u')$. Moreover $F(iX,iY)=iF(X,Y)$ and
$F(X,Y)=(X+Y)\,V(X,Y)$ with $V\in\Z[\frac12][[X,Y]]$, $V(0,0)=1$.}

\emph{Proof.} $\binom{1/2}{k}=(-1)^{k-1}\binom{2k}{k}/(4^k(2k-1))$ lies in $\Z[\frac12]$, because
$\binom{2k}{k}/(2k-1)=2\binom{2k-1}{k}/(2k-1)=2\binom{2k-2}{k-1}/k$ is twice a Catalan number. As power series in $u$,
$\mathrm{sl}'=\sqrt{1-\mathrm{sl}^4}$ with the series \eqref{el:F}, since both sides have square
$1-\mathrm{sl}^4$ and constant term $1$. So \eqref{el:add} reads
$\mathrm{sl}(u+u')=F(\mathrm{sl}\,u,\mathrm{sl}\,u')$; put $u=\mathrm{arcsl}\,X$, $u'=\mathrm{arcsl}\,Y$. The rule
$F(iX,iY)=iF(X,Y)$ is visible in \eqref{el:F}. Since $F(X,-X)=0$, the series $F(X,Y'-X)$ has no term free
of $Y'$; dividing it by $Y'$ and substituting back $Y'=X+Y$ gives $V$. \qed

\noindent The rule $F(iX,iY)=iF(X,Y)$ says that $X\mapsto iX$ is compatible with the law $F$. This is the
symmetry $x\mapsto ix$ of the curve $y^2=1-x^4$, which \S I.3 calls its complex multiplication by $\Z[i]$. Since $2$ is the only prime in the denominators, $F$ is a formal group
law over $\Z_{(p)}$ for every odd $p$. The next two paragraphs determine its height.

\subsection*{5.\quad The functional equation of the logarithm at an inert prime}

Throughout this paragraph $p$ is a prime with $p\equiv3\bmod4$, except where another condition is stated,
and $q=p^2$. For a rational number $x\ne0$ the \emph{$p$-adic valuation} $v_p(x)$ is the exponent of $p$ in
$x$: $x=p^{v_p(x)}a/b$ with integers $a,b$ prime to $p$; and $v_p(0)=+\infty$. For example $v_3(18)=2$ and
$v_3(\frac5{12})=-1$. So $\Z_{(p)}$ is the set of the $x$ with $v_p(x)\ge0$. A congruence
$x\equiv y\pmod{p^r}$ between elements of $\Z_{(p)}$ means $v_p(x-y)\ge r$. For an
integer $N\ge0$ let $N!'$ be the product of the integers $j$ with $1\le j\le N$ and $p\nmid j$. Sorting
the factors of $N!$ by divisibility by $p$ gives
\begin{equation}\label{el:fact}
 N!=p^{\lfloor N/p\rfloor}\cdot\lfloor N/p\rfloor!\cdot N!' .
\end{equation}

\textbf{Lemma Le11 (Wilson's theorem for prime powers).} \emph{Let $p$ be an odd prime, $b\ge1$ and
$P=p^b$.}
\begin{enumerate}[label=\textup{(\roman*)},leftmargin=2.4em,itemsep=1pt,topsep=2pt]
\item \emph{$(N+P)!'\equiv-N!'\pmod P$ for every $N\ge0$.}
\item \emph{$N!'\cdot(P-1-N)!'\equiv(-1)^{N+1-\lfloor N/p\rfloor}\pmod P$ for $0\le N<P$.}
\end{enumerate}
\emph{Proof.} (i) The integers prime to $p$ in an interval of length $P$ represent each invertible class
modulo $P$ once. In the product of all invertible classes, each class cancels against its inverse, except
the classes with $x^2=1$. If $P\mid(x-1)(x+1)$ then $p$ divides only one factor, since $p$ is odd; so
$x=\pm1$, and the product is $-1$. (ii) By (i) with $N=0$, $(P-1)!'\equiv-1$. The factors $j$ of
$(P-1)!'$ with $j\ge P-N$ are $j=P-j'$ with $1\le j'\le N$, $p\nmid j'$; there are $N-\lfloor N/p\rfloor$ of
them, and their product is $\equiv(-1)^{N-\lfloor N/p\rfloor}N!'$. \qed

\noindent \emph{Example.} $p=3$, $P=9$. The product of the integers below $9$ prime to $3$ is
$8!'=1\cdot2\cdot4\cdot5\cdot7\cdot8=2240=9\cdot249-1$. For (i) with $N=4$: $13!'=3203200\equiv1$ and
$-4!'=-8\equiv1\pmod9$. For (ii) with $N=2$: $2!'\cdot6!'=2\cdot40=80\equiv-1\pmod 9$, and the exponent
$N+1-\lfloor N/p\rfloor=3$ is odd.\footnote{Script \texttt{el10\_examples.py}, which also tests the lemma for
$p\le11$, $b\le3$ and all $N<3P$.}

\textbf{Theorem Le12 (congruences for $\kappa_k=\binom{2k}{k}4^{-k}$).} \emph{Let $p\equiv3\bmod4$, let
$k\ge0$, and let $a=v_p(4k+1)$.}
\begin{enumerate}[label=\textup{(\roman*)},leftmargin=2.4em,itemsep=1pt,topsep=2pt]
\item \emph{If $a\ge1$ then $v_p(\kappa_k)\ge1$.}
\item \emph{If $a\ge2$, write $4k+1=p^2(4k'+1)$. Then}
\begin{equation}\label{el:asd}
 \kappa_k\equiv-p\,\kappa_{k'}\pmod{p^{a}} .
\end{equation}
\end{enumerate}

\emph{Proof.} (i) Let $4k+1=pm_1$. Then $m_1\equiv3\bmod4$; write $m_1=4j+3$. One checks
\begin{equation}\label{D:470}
 2k=p(2j+1)+\tfrac{p-1}2,\qquad k=pj+\tfrac{3p-1}4,\qquad 0\le\tfrac{p-1}2,\ \tfrac{3p-1}4<p,
\end{equation}
so $\lfloor2k/p\rfloor=2j+1$ and $\lfloor k/p\rfloor=j$. By \eqref{el:fact} applied to $(2k)!$ and to $k!$,
\begin{equation}\label{el:step1}
 \binom{2k}{k}=p\,(2j+1)\binom{2j}{j}\,\frac{(2k)!'}{(k!')^2}.
\end{equation}
The last factor is a quotient of integers prime to $p$. This proves (i).

(ii) Now $p\mid m_1$; write $m_1=pm_2$ with $m_2=4k'+1$. Then $2j=p\cdot2k'+\frac{p-3}2$ and
$j=pk'+\frac{p-3}4$, so $\lfloor2j/p\rfloor=2k'$, $\lfloor j/p\rfloor=k'$, and \eqref{el:fact} gives
$\binom{2j}{j}=\binom{2k'}{k'}(2j)!'/(j!')^2$. With \eqref{el:step1},
\begin{equation}\label{el:step2}
 \kappa_k=p\,\kappa_{k'}\cdot(2j+1)\cdot4^{k'-k}\cdot\frac{(2k)!'\,(2j)!'}{(k!'\,j!')^2}.
\end{equation}
Let $b=a-1\ge1$ and $P=p^b$. Since $\kappa_{k'}\in\Z_{(p)}$, it suffices to show that the product of the last
three factors of \eqref{el:step2} is $\equiv-1$ modulo $P$. Note that $P$ divides $m_1=4j+3$; write $m_1=Ps$
with $s$ odd.

\emph{The factor $2j+1$.} $2(2j+1)=m_1-1\equiv-1$, so $2j+1\equiv-\frac12\pmod P$.

\emph{The power of $4$.} $4(k-k')=(q-1)m_2$, so $4^{k-k'}=2^{(q-1)m_2/2}$. Here $p^{b-1}$ divides $m_2$,
so the exponent is a multiple of $(p-1)p^{b-1}$, and $4^{k-k'}\equiv1\pmod P$ by Euler's theorem.

\emph{The numerator.} $2k=\frac12(pPs-1)=P\cdot\frac{ps-1}2+\frac{P-1}2$ and
$2j+1=\frac12(Ps-1)=P\cdot\frac{s-1}2+\frac{P-1}2$. By Lemma~Le11\,(i),
\begin{equation}\label{D:473}
 (2k)!'\,(2j+1)!'\equiv(-1)^{\frac{ps-1}2+\frac{s-1}2}\Big(\big(\tfrac{P-1}2\big)!'\Big)^2 \pmod P.
\end{equation}
The exponent is $\frac{(p+1)s}2-1$, which is odd because $4\mid p+1$. By Lemma~Le11\,(ii) with
$N=\frac{P-1}2=P-1-N$, the square is $\equiv(-1)^{N+1-\lfloor N/p\rfloor}$, and
$N+1-\lfloor N/p\rfloor=\frac{p^b+1}2-\frac{p^{b-1}-1}2=p^{b-1}\cdot\frac{p-1}2+1$ is even. So
$(2k)!'(2j+1)!'\equiv-1$, and $(2k)!'(2j)!'\equiv-1/(2j+1)\equiv2\pmod P$.

\emph{The denominator.} $k+j=\frac14(pPs+Ps)-1=Pw-1$ with $w=\frac{(p+1)s}4\in\Z$. Write $k=Pk_1+k_0$ with
$0\le k_0<P$; then $j=P(w-1-k_1)+(P-1-k_0)$. By Lemma~Le11, $k!'\,j!'\equiv\pm k_0!'\,(P-1-k_0)!'\equiv\pm1$,
so $(k!'j!')^2\equiv1\pmod P$.

The product of the three factors is therefore $\equiv(-\frac12)\cdot1\cdot2=-1\pmod P$. \qed

\noindent \emph{Examples.} $p=3$, $k=2$, $k'=0$: $\kappa_2=\frac38$ and $\frac38+3=\frac{27}8$. $p=7$,
$k=12$, $k'=0$: $\kappa_{12}+7=\frac{2704156}{4^{12}}+7=7^3\cdot\frac{350276}{4^{12}}$. Both (i) and (ii)
were tested for all $4k+1\le4\cdot10^6$ and the $24$ inert primes $p\le199$, in exact
$p$-adic arithmetic;\footnote{Script \texttt{el3\_congruence.py}.} in every case \eqref{el:asd} holds even
modulo $p^{a+1}$, which we do not prove. With the sign $+p$ in place of $-p$ it fails at $k=(q-1)/4$.

\textbf{Theorem Le13 (the functional equation).} \emph{Let $p\equiv3\bmod4$ and $q=p^2$. Then}
\begin{equation}\label{el:fe}
 \mathrm{arcsl}(t)+\frac1p\,\mathrm{arcsl}\big(t^{q}\big)\ \in\ \Z_{(p)}[[t]] .
\end{equation}
\emph{The series $\mathrm{arcsl}(t)-\frac1p\,\mathrm{arcsl}(t^{q})$ does not have coefficients in
$\Z_{(p)}$.}

\emph{Proof.} By Lemma~Le4\,(iii) the coefficient of $t^m$ in \eqref{el:fe} is $0$ unless $m=4k+1$. Let
$a=v_p(m)$. If $a=0$ the coefficient is $\kappa_k/m\in\Z_{(p)}$. If $a=1$ it is $\kappa_k/m$, and
$v_p(\kappa_k)\ge1$ by Theorem~Le12\,(i). If $a\ge2$ and $m=q(4k'+1)$, it is
\begin{equation}\label{D:475}
 \frac{\kappa_k}{m}+\frac1p\cdot\frac{\kappa_{k'}}{m/q}=\frac{\kappa_k+p\,\kappa_{k'}}{m},
\end{equation}
which lies in $\Z_{(p)}$ by \eqref{el:asd}. With the other sign the coefficient of $t^q$ is
$(\kappa_{(q-1)/4}-p)/q$, and $\kappa_{(q-1)/4}-p\equiv-2p\pmod{p^2}$. \qed

\noindent Congruences of this kind for the coefficients of a differential on an elliptic curve go back to
Atkin and Swinnerton-Dyer \cite{ASD71} (cited from memory). In the form of Hazewinkel \cite{Haz78},
(33.2.3), ``Theorem (Atkin--Swinnerton-Dyer conjectures)'' (read in a local copy), with the parameters
$a_p=0$, $b_p=1$ there, the statement is \eqref{el:asd} with the same modulus $p^a$; according to the notes
(E.5.2) on that section in \cite{Haz78}, p.~476 (read in a local copy), proofs were announced and sketched by
Cartier, and another proof was given by Ditters. We have not checked
that the lemniscate with these parameters satisfies the hypotheses of that theorem; Theorem~Le12 is proved
above without it. The sign in \eqref{el:fe} belongs to the curve: in the theory of curves over finite
fields, which we do not use, it expresses that the map $(x,y)\mapsto(x^p,y^p)$ of the curve modulo $p$, called
the Frobenius map, applied twice, is the multiplication by $-p$. In \eqref{el:fe} the exponent is $p^2$,
not $p$: this is what height two means for a logarithm.
Now we compare $\mathrm{arcsl}$ with the simplest series that satisfy the same equation. Put
\begin{equation}\label{el:Lpm}
 L_-(T)=\sum_{e\ge0}(-1)^e\,\frac{T^{q^e}}{p^e},\qquad L_+(T)=\sum_{e\ge0}\frac{T^{q^e}}{p^e};\qquad
 L_\mp(T)\pm\frac1p\,L_\mp(T^{q})=T .
\end{equation}
$L_+$ is the model series of \S I.3 and the logarithm \eqref{dv:log} of \S III.5 in one variable. We call
$L_+$ and $L_-$ the two \emph{model logarithms of height two}.

\textbf{Lemma Le14 (the functional equation lemma, a special case).} \emph{Let $\mathcal O$ be $\Z_{(p)}$,
or the ring $\Z[i]_{(p)}$ of the numbers $x/y$ with $x\in\Z[i]$, $y\in\Z$, $p\nmid y$, where
$p\equiv3\bmod4$. Let $\varepsilon\in\{1,-1\}$, and let $l(t)=t+\cdots$ and $\tilde l(t)=t+\cdots$ be power
series with coefficients in the fraction field of $\mathcal O$ such that}
\begin{equation}\label{D:477}
 l(t)-\frac\varepsilon p\,l(t^q)\in\mathcal O[[t]]\qquad\text{and}\qquad
 \tilde l(t)-\frac\varepsilon p\,\tilde l(t^q)\in\mathcal O[[t]].
\end{equation}
\begin{enumerate}[label=\textup{(\alph*)},leftmargin=2.4em,itemsep=1pt,topsep=2pt]
\item \emph{If $A,B$ are power series without constant term, in any number of variables, with coefficients
in $\mathcal O$ and $A-B\in p\,\mathcal O[[\ldots]]$, then $l(A)-l(B)\in p\,\mathcal O[[\ldots]]$.}
\item \emph{The series $l^{-1}\circ\tilde l$ has coefficients in $\mathcal O$.}
\end{enumerate}

\emph{Proof.} In both rings every element $x$ satisfies $x^q\equiv x$ modulo $p\mathcal O$: for $\Z_{(p)}$
by Fermat's theorem, and for $\Z[i]_{(p)}$ because $\Z[i]/p\Z[i]$ is a field with $q$ elements when
$p\equiv3\bmod4$, so that its $q-1$ nonzero elements form a group and satisfy $x^{q-1}=1$. (This ring has the
$p^2$ elements $a+bi$ with $0\le a,b<p$. The number $-1$ is not a square modulo $p$: from $y^2\equiv-1$ one
would get $y^{p-1}\equiv(-1)^{(p-1)/2}=-1$, against Fermat's theorem. So $a^2+b^2\not\equiv0$ modulo $p$
unless $p$ divides $a$ and $b$, and a nonzero class $a+bi$ has the inverse $(a-bi)/(a^2+b^2)$.) So for a polynomial $V$ over $\mathcal O$ we have $V(T)^q\equiv V(T^q)$ modulo
$p\,\mathcal O[T]$. Write $l=\sum l_mt^m$ and $\varphi(t)=l(t)-\frac\varepsilon pl(t^q)$.

(a) Comparing coefficients, $l_m=\varphi_m+\frac\varepsilon pl_{m/q}$, where $l_{m/q}=0$ if $q\nmid m$. By
induction, $p^el_m\in\mathcal O$ if $q^{e+1}\nmid m$. If $A=B+pC$ then $A^p=B^p+p^2(\ldots)$, and by
induction $A^m-B^m\in p^{1+e'}\mathcal O[[\ldots]]$ where $p^{e'}$ is the exact power of $p$ in $m$. If
$q^e$ divides $m$ then $e'\ge2e\ge e$. So every term $l_m(A^m-B^m)$ lies in $p\,\mathcal O[[\ldots]]$.

(b) Let $U=l^{-1}\circ\tilde l=T+\sum_{n\ge2}u_nT^n$. Let $n\ge2$, suppose $u_j\in\mathcal O$ for $j<n$, and
let $V$ be the polynomial $T+\cdots+u_{n-1}T^{n-1}$. We compute modulo $T^{n+1}$. Since $l(U)=\tilde l$,
\begin{equation}\label{D:478}
 \tilde l(T)=\varphi(U(T))+\frac\varepsilon p\,l\big(U(T)^q\big)\equiv\varphi(V(T))+u_nT^n+\frac\varepsilon p\,l\big(V(T)^q\big),
\end{equation}
because $U^q-V^q$ is divisible by $T^{n+q-1}$. By (a), $l(V(T)^q)=l(V(T^q))+p\,W$ with $W$ over
$\mathcal O$. And $V(T^q)\equiv U(T^q)$ modulo $T^{qn}$, so $l(V(T^q))\equiv l(U(T^q))=\tilde l(T^q)$. Hence
\begin{equation}\label{D:479}
 u_nT^n\equiv\Big(\tilde l(T)-\frac\varepsilon p\,\tilde l(T^q)\Big)-\varphi(V(T))-\varepsilon W(T)\pmod{T^{n+1}}.
\end{equation}
The right side has coefficients in $\mathcal O$. So $u_n\in\mathcal O$. \qed

\noindent This is a special case of the functional equation lemma of Hazewinkel \cite{Haz78}, I.2.2 (read
in a local copy): part (b) is its part (ii), and part (a) is the analogue of its part (iv), which is stated
there for series in one variable; both for the equation
$l(t)-\frac\varepsilon pl(t^q)\in\mathcal O[[t]]$.

\emph{Example.} Let $p=3$, so $q=9$. Take $l=\mathrm{arcsl}$ and $\tilde l=L_-$, with $\varepsilon=-1$; the
hypotheses are \eqref{el:fe} and \eqref{el:Lpm}. Then $l^{-1}\circ\tilde l=\mathrm{sl}\circ L_-$, and by
Lemma~Le4, with $A_5/5!=-\frac1{10}$ and $A_9/9!=\frac1{120}$,
\begin{equation}\label{D:480}
 \mathrm{sl}\Big(T-\frac{T^9}3+\cdots\Big)=T-\frac{T^5}{10}+\Big(\frac1{120}-\frac13\Big)T^9+\cdots
 =T-\frac{T^5}{10}-\frac{13}{40}\,T^9+\frac{863}{5200}\,T^{13}-\cdots .
\end{equation}
No denominator is divisible by $3$, as (b) asserts. With $L_+$ in place of $L_-$ the coefficient of $T^9$
is $\frac1{120}+\frac13=\frac{41}{120}$, which is not in $\Z_{(3)}$: the sign matters. The pattern is the same
as in a better known case of height one, which we do not use: the series $\log(1+X)$ and
$X+X^p/p+X^{p^2}/p^2+\cdots$ both satisfy $l(X)-\frac1p\,l(X^p)\in\Z_{(p)}[[X]]$, and
$\exp(X+X^p/p+\cdots)$ has coefficients in $\Z_{(p)}$, although $\exp(X)$ has not (a classical fact, not
proved here; this series is called the Artin--Hasse exponential). For $p=2$ it is
$1+X+X^2+\frac23X^3+\frac23X^4+\frac7{15}X^5+\cdots$.\footnote{Script \texttt{el10\_examples.py}.}

The next proposition holds for every prime $p$. It describes the inverse series of the two model
logarithms, and it is the source of the divisibility in \S6.

\textbf{Proposition Le15 (the inverse of the height-two logarithm).} \emph{Let $p$ be any prime, $q=p^2$,
and let $g=\sum_{n\ge1}g_nz^n$ be the inverse series of $L_+$ or of $L_-$. Then $g_n=0$ unless
$n\equiv1\bmod(q-1)$, and for $n=1+j(q-1)$ the number}
\begin{equation}\label{el:Hur}
 \frac{n!\,g_n}{p^{\,pj}}\qquad\text{is an integer.\quad Thus } v_p(n!\,g_n)\ge\frac{p\,(n-1)}{p^2-1}.
\end{equation}

\emph{Proof.} Let $\varepsilon=\pm1$ be the sign. The equation $L_\varepsilon(g(z))=z$ is
$g=z-\sum_{e\ge1}\varepsilon^ep^{-e}g^{q^e}$. It determines $g_n$ from the $g_j$ with $j<n$, and by induction $g$
has terms only in the degrees $\equiv1\bmod(q-1)$, since $q^e\equiv1$. Let $\mathfrak r=p^{p/(q-1)}$, a
positive real number, and put $H(x)=\mathfrak r\,g(x/\mathfrak r)=\sum_ng_n\mathfrak r^{1-n}x^n$. Its
coefficients are the rational numbers $g_np^{-pj}$, $n=1+j(q-1)$. Substituting $z=x/\mathfrak r$ and
multiplying by $\mathfrak r$,
\begin{equation}\label{D:482}
 H=x-\sum_{e\ge1}\varepsilon^e\,\frac{H^{q^e}}{p^{d_e}},\qquad d_e=e+\frac{p\,(q^e-1)}{q-1} .
\end{equation}
By Legendre's formula, $v_p(q^e!)=\frac{q^e-1}{p-1}=(p+1)\frac{q^e-1}{q-1}=d_e+\big(1+q+\cdots+q^{e-1}\big)-e\ge d_e$.
So $w_e=q^e!/p^{d_e}$ is an integer, and
\begin{equation}\label{el:Heq}
 H=x-\sum_{e\ge1}\varepsilon^e\,w_e\,\frac{H^{q^e}}{q^e!} .
\end{equation}
We show by induction on $n$ that the coefficient of $x^n$ in $H$ is an integer divided by $n!$. Let $H_{<n}$
be the sum of the terms of $H$ of degree $<n$. By induction it is a Hurwitz series over $\Z$. The
coefficient of $x^n$ on the right of \eqref{el:Heq} is the same for $H$ and for $H_{<n}$, because $q^e\ge2$
and $H$ has no constant term. By Lemma~Le3 each $H_{<n}^{\,q^e}/q^e!$ is a Hurwitz series over $\Z$. \qed

\noindent So $H$ is a Hurwitz series over $\Z$: the inverse of the logarithm, after the change of scale
$\mathfrak r$, has integral derivatives. This is the exact analogue of $e^z-1=\sum z^n/n!$ for the
multiplicative group, and the scale $\mathfrak r$ carries the gain: $v_p(\mathfrak r)=p/(p^2-1)$ per
derivative. For $p=3$: $g_-(z)=z+z^9/3+\cdots$, and $9!/3=3^3\cdot4480$. In a computation of
$g_\pm$ up to $z^{1500}$ for $p=3,5,7,11$ the bound \eqref{el:Hur} holds with equality at \emph{every}
$n\equiv1\bmod(q-1)$ in the range.\footnote{Scripts \texttt{el6\_inverse\_log.py} and
\texttt{el10\_examples.py}.} Since $v_p(n!)\le(n-1)/(p-1)$ by Legendre's formula, and
$\frac p{q-1}-\frac1{p-1}=-\frac1{q-1}$, the bound \eqref{el:Hur} implies
\begin{equation}\label{el:ginv}
 v_p(g_n)\ \ge\ -\frac{n-1}{q-1}\qquad(n\ge1).
\end{equation}
This is Lemma~1 of \S III.5 in one variable, that is \eqref{dv:ginv}. The proposition is used twice: the
bound \eqref{el:ginv} in the proof of Corollary~Le16, and the full statement \eqref{el:Hur}, with the
factorial, in Theorem~Le18.

Divisibility statements of this kind, for the Taylor coefficients of functions on a formal group of height
two at the primes concerned (for an elliptic curve these are called the supersingular primes; here they
are the inert primes), are in the literature: see N.~Katz \cite{Kat81}, and for the numbers $A_n$ themselves
Hurwitz \cite{Hur99}. Both are cited from memory. We have not compared Proposition~Le15 or Theorem~Le18
below with the statements proved there, and we do not claim that either is new.

\textbf{Corollary Le16 (the lemniscate at an inert prime has height two).} \emph{Let $p\equiv3\bmod4$.}
\begin{enumerate}[label=\textup{(\roman*)},leftmargin=2.4em,itemsep=1pt,topsep=2pt]
\item \emph{The series $U_p=\mathrm{sl}\circ L_-$ and its inverse $L_-^{-1}\circ\mathrm{arcsl}$ have
coefficients in $\Z_{(p)}$, and $U_p(T)=T+\cdots$ is odd. So the formal group law $F$ of \eqref{el:F} is
strictly isomorphic over $\Z_{(p)}$ to the formal group law $G_-(X,Y)=g_-(L_-(X)+L_-(Y))$ with the
logarithm $L_-$, where $g_-$ is the inverse series of $L_-$.}
\item \emph{$[p]_F(X)\equiv-X^{p^2}$ modulo $p$, that is, $[p]_F(X)+X^{p^2}\in p\,\Z_{(p)}[[X]]$: there are
no further terms. In particular the height of $F$ at $p$ is $2$.}
\end{enumerate}
\emph{Proof.} (i) Apply Lemma~Le14\,(b) with $\mathcal O=\Z_{(p)}$, $\varepsilon=-1$, to the pair
$(\mathrm{arcsl},L_-)$ and to the pair $(L_-,\mathrm{arcsl})$; the hypotheses are \eqref{el:fe} and
\eqref{el:Lpm}. The series $U_p$ is odd because $\mathrm{sl}$ and $L_-$ are odd. (ii) $G_-(X,Y)=U_p^{-1}(F(U_p(X),U_p(Y)))$ has coefficients in $\Z_{(p)}$. By
\eqref{el:Lpm}, $-pL_-(T)=L_-(T^q)-pT$. So $[-p]_{G_-}(T)=g_-(L_-(T^q)+L_-(g_-(-pT)))=G_-(T^q,g_-(-pT))$.
By \eqref{el:ginv} the coefficient $g_n$ of $z^n$ in $g_-$ has $v_p(g_n)\ge-(n-1)/(q-1)$, so
$v_p(g_n(-p)^n)\ge n-(n-1)/(q-1)\ge1$ for $n\ge1$. Hence $g_-(-pT)\in p\,\Z_{(p)}[[T]]$ and
$[-p]_{G_-}(T)\equiv G_-(T^q,0)=T^q$ modulo $p$. Since $L_-$ is odd, $[p]_{G_-}(T)=-[-p]_{G_-}(T)\equiv-T^q$.
Finally $[p]_F=U_p\circ[p]_{G_-}\circ U_p^{-1}$, because $\mathrm{arcsl}=L_-\circ U_p^{-1}$. Write
$V=U_p^{-1}$; by (i) the series $U_p$ and $V$ have coefficients in $\Z_{(p)}$ and no constant term. A
congruence modulo $p$ between power series over $\Z_{(p)}$ remains true when $V(X)$ is substituted for the
variable, and when both sides are substituted in $U_p$. So $[p]_{G_-}(V(X))\equiv-V(X)^q$. Next
$V(X)^q\equiv V(X^q)$ modulo $p$: this was shown for polynomials at the beginning of the proof of
Lemma~Le14, and the coefficient of $X^n$ on either side depends only on the terms of $V$ of degree $\le n$.
Since $U_p$ is odd,
\begin{equation}\label{D:485}
 [p]_F(X)\equiv U_p\big(-V(X)^q\big)\equiv U_p\big(-V(X^q)\big)=-U_p\big(V(X^q)\big)=-X^q\pmod p .
\end{equation}
\qed

\noindent For example, a computation of $\mathrm{sl}(p\,\mathrm{arcsl}(X))$ in exact arithmetic, made for
this revision, shows that modulo $p$ it is the single term $-X^{9}$ up to degree $40$ for $p=3$, and the
single term $-X^{49}$ up to degree $60$ for $p=7$.

\noindent \emph{Remarks.} (1) Part (i) is the analogue, for one function of one variable, of hypothesis
\textup{(B4)} of \S III.5: the function $\mathrm{sl}$, written in the coordinate $T$ of the height-two
logarithm, has $p$-integral coefficients. (Hypothesis (B4) itself concerns the series
$\exp(\Delta_i\beta)\circ L$; see \S9.) A computation confirms part (i) up to $T^{800}$ for
$p=3,7,11,19,23$; for $L_+$ in place of $L_-$ it fails at $T^{q}$, and for $p=5,13$ it fails at $T^p$ with
either sign.\footnote{Script \texttt{el6\_inverse\_log.py}.}

(2) \emph{The sign.} The model series of \S I.3 is $L_+$, and the lemniscate has $L_-$. The two differ by a
change of variable: if $\xi^{q-1}=-1$ then $L_+(T)=\xi^{-1}L_-(\xi T)$, because
$\xi^{q^e-1}=(-1)^{1+q+\cdots+q^{e-1}}=(-1)^e$. Such a $\xi$ is a root of unity of order $2(q-1)$, prime to
$p$. So over a ring that contains $\xi$ the two formal group laws are isomorphic, and every statement
about valuations of coefficients is the same for both. Over $\Z_{(p)}$ itself they are not strictly
isomorphic: otherwise $\mathrm{arcsl}$ would satisfy the equation of Lemma~Le14 with $\varepsilon=+1$, by
Lemma~Le14\,(a), against Theorem~Le13. In \S9 the sign is changed by replacing the curve
$Y^2=X^3-X$ by a curve $Y^2=X^3-cX$ with a suitable number $c$.

\subsection*{6.\quad Divisibility at inert primes}

By Corollary~Le16, $\mathrm{sl}(z)=U_p(g_-(z))$ with $U_p$ integral at $p$. So the denominators of the
Taylor coefficients of $\mathrm{sl}$ are those of the inverse $g_-$ of $L_-$. The bound \eqref{el:ginv} for
the coefficients of $g_-$ would give the theorem of this paragraph up to an error of order $\log n$ in the
exponent. To have the exponent exactly we use Proposition~Le15, in which the factorial is taken into
account, and one more lemma on factorials.

\textbf{Lemma Le17.} \emph{Let $p$ be a prime and $j\ge1$ an integer, $j\ne p^2-1$. Then
$v_p(j!)\ge\lfloor j\,p/(p^2-1)\rfloor$. For $j=p^2-1$ one has $v_p(j!)=p-1=j\,p/(p^2-1)-1$.}

\emph{Proof.} Let $j=\sum_{i\ge0}d_ip^i$ with digits $0\le d_i<p$. By Legendre's formula
$v_p(j!)=\sum_id_i\frac{p^i-1}{p-1}$, so
\begin{equation}\label{D:486}
 (p^2-1)\Big(v_p(j!)-\frac{j\,p}{p^2-1}\Big)=\sum_id_i\big((p+1)(p^i-1)-p^{i+1}\big)
 =-p\,d_0-d_1+\sum_{i\ge2}d_i\,(p^i-p-1).
\end{equation}
Now $p\,d_0+d_1\le p^2-1$, with equality only if $d_0=d_1=p-1$, and $p^i-p-1>0$ for $i\ge2$. So the right
side is $>-(p^2-1)$ unless $j=p^2-1$. Hence $v_p(j!)>j\,p/(p^2-1)-1$, and $v_p(j!)$ is an integer. \qed

\noindent \emph{Example.} $p=3$, so $j\,p/(p^2-1)=3j/8$. For $j=5$: $v_3(5!)=1=\lfloor15/8\rfloor$. For
$j=9$: $v_3(9!)=4\ge3=\lfloor27/8\rfloor$. For the exceptional $j=8$: $v_3(8!)=2$, and $3\cdot8/8=3$. For
$p=7$ the exception is $j=48$: $v_7(48!)=6$ and $7\cdot48/48=7$.\footnote{Script \texttt{el10\_examples.py}:
for $p<60$ and $j<20000$ the only exception is $j=p^2-1$.}

\textbf{Theorem Le18 (divisibility at inert primes).} \emph{Let $p\equiv3\bmod4$, and let
$f(t)\in\Z_{(p)}[[t]]$. For $n\ge0$ let $D_n$ be the $n$th derivative of $f(\mathrm{sl}(z))$ at $z=0$. Then}
\begin{equation}\label{el:div}
 v_p(D_n)\ \ge\ \Big\lfloor\frac{n\,p}{p^2-1}\Big\rfloor\quad\text{if }(p^2-1)\nmid n,\qquad
 v_p(D_n)\ \ge\ \frac{n\,p}{p^2-1}-1\quad\text{if }(p^2-1)\mid n .
\end{equation}
\emph{If $f$ is odd, the first bound holds for all $n$. In particular, for every $n\ge1$,}
\begin{equation}\label{el:divA}
 v_p(A_n)\ \ge\ \Big\lfloor\frac{n\,p}{p^2-1}\Big\rfloor\ >\ n\Big(\frac1p+\frac1{p^3}+\frac1{p^5}+\cdots\Big)-1 .
\end{equation}

\emph{Proof.} By Corollary~Le16, $f(\mathrm{sl}(z))=(f\circ U_p)(g_-(z))$ and
$f\circ U_p=\sum_{j\ge0}b_jT^j$ with $b_j\in\Z_{(p)}$. With $H$ and $\mathfrak r$ as in the proof of
Proposition~Le15, $g_-(z)^j=\mathfrak r^{-j}H(\mathfrak rz)^j$. By Lemma~Le3, $H^j=j!\sum_nh_{j,n}x^n/n!$ with
integers $h_{j,n}$, and $h_{j,n}=0$ unless $n\equiv j\bmod(q-1)$ and $n\ge j$. So
\begin{equation}\label{D:489}
 D_n=\sum_j b_j\;j!\;p^{\,p(n-j)/(q-1)}\,h_{j,n},
\end{equation}
the sum over the $j$ with $0\le j\le n$ and $j\equiv n\bmod(q-1)$. For $n=0$ the claim \eqref{el:div} is
trivial, because $D_0=f(0)\in\Z_{(p)}$. Let $n\ge1$. The term $j=0$ vanishes, because $H^0=1$ and so
$h_{0,n}=0$. A term with $j\ge1$ has valuation at least
$v_p(j!)+p(n-j)/(q-1)$, where $p(n-j)/(q-1)$ is an integer. If $j\ne q-1$, Lemma~Le17 gives
$v_p(j!)+p(n-j)/(q-1)\ge\lfloor jp/(q-1)\rfloor+p(n-j)/(q-1)=\lfloor np/(q-1)\rfloor$. The term $j=q-1$ occurs
only if $(q-1)\mid n$, and then its valuation is at least $np/(q-1)-1$. If $f$ is odd then $f\circ U_p$ is
odd, and $j=q-1$, an even number, does not occur. \qed

\noindent \emph{Comparison with \S I.3.} The number $x=n\,p/(p^2-1)$ is the prediction (the display after
the table of \S I.3), and \S I.3 observed that the exponent is ``at least the prediction minus $1$''. For
an integer $e$ the two statements $e\ge\lfloor x\rfloor$ and $e\ge x-1$ compare as follows. If $x$ is not
an integer they are equivalent, because $\lfloor x\rfloor$ is then the least integer that is $\ge x-1$. If
$x$ is an integer, the first says $e\ge x$ and the second $e\ge x-1$: the first is stronger by exactly
$1$. Now $x$ is an integer only if $(p^2-1)\mid n$, because $p$ is prime to $p^2-1$; then $n$ is even.
So for odd $n$ the bound \eqref{el:divA} is the same statement as the one observed in \S I.3, and for even
$n$ both are empty, because $A_n=0$. What Theorem~Le18 adds to the observation is the range: every inert
prime and every $n$. For a general $f$ the two cases are those of \eqref{el:div}, and there the loss of $1$
at the multiples of $p^2-1$ does occur. The bound \eqref{el:divA} cannot be improved in this form:

\begin{center}\small
\begin{tabular}{r|rr|rr|rr}
\toprule
$n$ & $v_3(A_n)$ & $\lfloor3n/8\rfloor$ & $v_7(A_n)$ & $\lfloor7n/48\rfloor$ & $v_{11}(A_n)$ & $\lfloor11n/120\rfloor$\\
\midrule
17 & 6 & 6 & 2 & 2 & 1 & 1\\
49 & 19 & 18 & 7 & 7 & 4 & 4\\
101 & 37 & 37 & 14 & 14 & 9 & 9\\
201 & 75 & 75 & 29 & 29 & 18 & 18\\
\bottomrule
\end{tabular}
\end{center}

\noindent We computed $A_n$ exactly for all $n\le2001$ and compared with
\eqref{el:divA} for the $24$ inert primes $p<200$:\footnote{Script \texttt{el4\_divisibility.py}.} there is
no violation, and equality holds for $418$ of
the $501$ values of $n$ when $p=3$, for $490$ when $p=7$, and for at least $465$ for each of the others.
The largest excess over $n\,p/(p^2-1)$ is $4.625$, at $p=3$, $n=1945$, as stated in \S I.3; we do not prove
that the excess is unbounded. For general $f$ the loss of $1$ in \eqref{el:div} occurs (computed with
random $f$).\footnote{Script \texttt{el6\_inverse\_log.py}.} As said after Proposition~Le15, divisibilities
of this kind at supersingular primes are in the literature (Katz \cite{Kat81}; Hurwitz \cite{Hur99} for the
numbers $A_n$; both from memory), and we have not compared the statements. Two weaker forms of \eqref{el:divA} are worth recording. First,
$\lfloor np/(p^2-1)\rfloor\ge\sum_{h\ \mathrm{odd}}\lfloor n/p^h\rfloor$: this is Legendre's formula
\eqref{leg} with the odd powers of $p$ only. Second, $v_p(A_n)\ge\lfloor n/p\rfloor$: every inert prime
$p\le n$ divides $A_n$.

\emph{The mechanism in one sentence.} In the coordinate $T$ of the height-two logarithm the function is
integral; the exponential of that logarithm is a Hurwitz series in $\mathfrak rz$ with
$v_p(\mathfrak r)=p/(p^2-1)$; so an $n$th derivative in $z$ carries the factor $\mathfrak r^n$ up to
rounding. This is the second of the three reasons for differentiating given in \S I.4.1, and it is the
one-variable case of \S III.5, \S\S2--3: there the loss is $\rho=1/(\ell^2-1)$ per degree, the factorial
gives $1/(\ell-1)$, and the difference is $\gamma=\ell/(\ell^2-1)$.

\subsection*{7.\quad The Chebyshev inequality that results, and its three defects}

\textbf{Proposition Le19.} \emph{For every $n\equiv1\bmod4$, and every constant $1<C_0<\sqrt5$,}
\begin{equation}\label{el:cheb}
 \sum_{p\equiv3\,(4)}\Big\lfloor\frac{n\,p}{p^2-1}\Big\rfloor\log p\ +\ \frac{n-1}2\,\log2\ \le\ \log|A_n|
 \ =\ \log n!-(n+1)\log\frac{\varpi}{\sqrt2}+\log4+O(C_0^{-n}).
\end{equation}
\emph{Consequently $\displaystyle\sum_{p\le n,\ p\equiv3\,(4)}\frac{\log p}{p}\le\log n+O(1)$.}

\emph{Proof.} By Lemma~Le4\,(ii), $A_n$ is a nonzero integer divisible by $2^{(n-1)/2}$, and by
\eqref{el:divA} by $p^{\lfloor np/(p^2-1)\rfloor}$ for each inert $p$. This is an application of the principle \eqref{det}
of \S I.3: a nonzero integer is at least as large as the product of the prime powers that divide it. The equality is
Theorem~Le9. For the last statement use $\lfloor np/(p^2-1)\rfloor\ge n/p-1$, the terms with $p>n$ being
$\ge0$; then $\sum_{p\le n}\log p\le\pi(n)\log n=O(n)$ by Chebyshev's bound (C4), and
$\log n!\le n\log n$. \qed

\noindent \emph{What it gives.} The inequality \eqref{el:cheb} is true and gives nothing. By the theorem
on primes in a progression quoted in (C4) (\cite{Apo76}, Theorem~7.3, with modulus $4$) the sum
$\sum_{p\le n,\,p\equiv3(4)}(\log p)/p$ is
$\frac12\log n+O(1)$. So \eqref{el:cheb} has the constant $1$ where the truth is $\frac12$: it is off by
the factor $2$ in the leading term. Numerically,\footnote{Script \texttt{el6\_inverse\_log.py}.} with the
sums over the inert primes $p\le n$:

\begin{center}\small
\begin{tabular}{r|rrrr}
\toprule
$n$ & $\log|A_n|$ & left side of \eqref{el:cheb}, inert part & the same with the true exponents & $\frac12n\log n$\\
\midrule
101 & 306.8 & 163.6 & 167.7 & 233.1\\
201 & 745.2 & 393.4 & 393.4 & 533.0\\
1001 & 5301.8 & 2749.7 & 2753.8 & 3457.8\\
2001 & 11979.5 & 6188.3 & 6188.3 & 7605.2\\
\bottomrule
\end{tabular}
\end{center}

\noindent The values $167.7$ and $306.8$ at $n=101$ are those of \S I.3. The three defects named in
\S I.3 and \S I.4 are now visible in \eqref{el:cheb}.
\begin{enumerate}[label=\textup{(D\arabic*)},leftmargin=3em,itemsep=1pt,topsep=2pt]
\item \emph{Rounding.} The exponent is $\lfloor np/(p^2-1)\rfloor$, and by the table after Theorem~Le18
the floor is attained for most $n$. Against the smooth exponent $np/(p^2-1)$ this loses up to $\log p$ at
each inert prime $p\le n$, in total up to $\sum_{p\le n,\,p\equiv3(4)}\log p\approx n/2$. That is of the
order of the linear term $n\,(1+\log(\varpi/\sqrt2))$ of the right side.
\item \emph{Only half the primes.} A split prime $p$ does not divide $A_p$ (Proposition~Le20 below), and in
the range computed most split primes divide no $A_n$ with $n\equiv1\bmod4$ (for $5,13,17,29$ this
is proved for all such $n$: Proposition~Le20a\,(iii)).
\item \emph{Size.} The right side is $n\log n+O(n)$, because the nearest pole of $\mathrm{sl}$ is at a
finite distance; the left side is $\frac12n\log n+O(n)$. An entire function of Gaussian growth, such as
$\Theta$, has derivatives of size $\sqrt{n!}$ (\S I.4.1); but for $\Theta$ the divisibility is lost: the
exponent of $3$ in $r_{30}$ is $14$, and $\lfloor3\cdot120/8\rfloor=45$.
\end{enumerate}
In the proof both properties are obtained together, for other theta series, and this is the statement that
\S III.5 reduces to hypothesis \textup{(B)}.

\subsection*{8.\quad The split primes}

\textbf{Proposition Le20.} \emph{Let $p\equiv1\bmod4$.}
\begin{enumerate}[label=\textup{(\roman*)},leftmargin=2.4em,itemsep=1pt,topsep=2pt]
\item \emph{If $4k+1=p\,(4k'+1)$ then $\kappa_k\equiv\alpha_p\,\kappa_{k'}\pmod p$, where
$\alpha_p=\kappa_{(p-1)/4}$ is prime to $p$. So the coefficient of $t^{p^j}$ in $\mathrm{arcsl}(t)$ has the
valuation exactly $-j$, for every $j\ge0$.}
\item \emph{$p$ does not divide $A_p$. In contrast, an inert prime $p$ divides $A_n$ for every $n\ge p$.}
\end{enumerate}
\emph{Proof.} (i) Here $k=pk'+\frac{p-1}4$ and $2k=p\cdot2k'+\frac{p-1}2$, with $\frac{p-1}4$ and
$\frac{p-1}2$ less than $p$. By Lucas's congruence (C3)
$\binom{2k}{k}\equiv\binom{2k'}{k'}\binom{(p-1)/2}{(p-1)/4}\pmod p$, and
$4^k=(4^p)^{k'}4^{(p-1)/4}\equiv4^{k'}4^{(p-1)/4}$. The binomial coefficient $\binom{(p-1)/2}{(p-1)/4}$ is
prime to $p$. By induction $\kappa_{(p^j-1)/4}\equiv\alpha_p^{\,j}$, and the coefficient is
$\kappa_{(p^j-1)/4}/p^j$. (ii) Let $\mathrm{arcsl}(t)=\sum l_mt^m$ and $\mathrm{sl}(z)=\sum a_nz^n$. From
$\mathrm{arcsl}(\mathrm{sl}(z))=z$, for $n\ge2$,
\begin{equation}\label{D:491}
 a_n=-\sum_{m=2}^{n}l_m\cdot\big(\text{the coefficient of $z^n$ in }\mathrm{sl}(z)^m\big).
\end{equation}
For $m<p$ the number $l_m$ lies in $\Z_{(p)}$. So $a_n\in\Z_{(p)}$ for $n<p$, by induction, and for $n=p$
all terms with $m<p$ lie in $\Z_{(p)}$, while the term $m=p$ is $l_p$, of valuation $-1$ by (i). So
$v_p(a_p)=-1$ and $v_p(A_p)=v_p(p!)-1=0$. The statement on inert primes is \eqref{el:divA}. \qed

\noindent So at a split prime the logarithm has the denominator $p^j$ already at the degree $p^j$, as for
the multiplicative group, and the first factor $p$ of $n!$ is cancelled at $n=p$. This is the behaviour
of height one described in \S I.3; that the height is one is Deuring's theorem \cite{Deu41} (cited from
memory), which we do not prove and do not use. A computation confirms $p\nmid A_p$ for the $147$ split
primes $p<2000$.\footnote{Script \texttt{el6\_inverse\_log.py}.} For
$n\le2001$, $n\equiv1\bmod4$, the only split primes $p<200$ that divide some $A_n$ are $157$ (thirteen values of $n$, the
first $n=101$) and $181$ (first at $n=105$), each to the first power.\footnote{Script
\texttt{el4\_divisibility.py}.}

\textbf{Proposition Le20a (the $p$-th power of the derivation).} \emph{Let $R=\Z[X,Y]/(Y^2-1+X^4)$, with
the derivation $D$ given by $DX=Y$, $DY=-2X^3$. For every odd prime $p$ and every $f\in R$,
\begin{equation}\label{el:Dp}
 D^pf\equiv A_p\,Df\pmod{pR}.
\end{equation}
Consequently:}
\begin{enumerate}[label=\textup{(\roman*)},leftmargin=2.4em,itemsep=1pt,topsep=2pt]
\item \emph{$A_{n+p-1}\equiv A_p\,A_n\pmod p$ for every $n\ge1$.}
\item \emph{If $p\equiv3\bmod4$ then $D^pR\subset pR$, and $p^{\lfloor n/p\rfloor}$ divides $A_n$ for every
$n\ge1$.}
\item \emph{If $p\equiv1\bmod4$ then $p$ divides some $A_n$ with $n\equiv1\bmod4$ if and only if it divides
one of $A_1,A_5,\ldots,A_{p-4}$. In particular $A_{4k+1}\equiv3^k\pmod5$ for every $k\ge0$, and none of
$5$, $13$, $17$, $29$ divides an $A_n$ with $n\equiv1\bmod4$.}
\end{enumerate}

\emph{Proof.} The rule is well defined on $R$: it is a derivation of $\Z[X,Y]$, and
$D(Y^2-1+X^4)=2Y\cdot(-2X^3)+4X^3Y=0$. Since $\mathrm{sl}'^2=1-\mathrm{sl}^4$ and
$\mathrm{sl}''=-2\,\mathrm{sl}^3$, the map $X\mapsto\mathrm{sl}$, $Y\mapsto\mathrm{sl}'$ carries $D$ to
$d/du$. So $\mathrm{sl}^{(n)}=(D^nX)(\mathrm{sl},\mathrm{sl}')$, and $A_n$ is the value of $D^nX$ at $X=0$,
$Y=1$; this evaluation is a ring homomorphism $R\to\Z$.

Let $\bar R=R/pR=\mathbb F_p[X,Y]/(Y^2-1+X^4)$. Every element of $\bar R$ is $P+P_1Y$ with unique
polynomials $P,P_1\in\mathbb F_p[X]$. The ring $\bar R$ has no zero divisors, because $1-X^4$ is not a
square in $\mathbb F_p[X]$: it has no repeated root, its derivative $-4X^3$ vanishing only at $0$. We prove
\eqref{el:Dp} in four steps.

\emph{Step 1: $D^p$ is a derivation of $\bar R$.} By Leibniz's rule
$D^p(fg)=\sum_{i=0}^p\binom piD^if\cdot D^{p-i}g$, and $p$ divides $\binom pi$ for $0<i<p$.

\emph{Step 2: $D^pX=QY$ with a polynomial $Q\in\mathbb F_p[X]$ of degree at most $p-1$.} For a polynomial
$P$ in $X$ we have $DP=P'Y$ and
\begin{equation}\label{D:493}
 D^2P=P''Y^2+P'\cdot(-2X^3)=(1-X^4)P''-2X^3P'.
\end{equation}
So $D^2$ maps polynomials in $X$ of degree at most $k$ to polynomials in $X$ of degree at most $k+2$. Since
$p-1$ is even, $P=D^{p-1}X$ is a polynomial in $X$ of degree at most $p$, and $D^pX=P'Y$. Put $Q=P'$.

\emph{Step 3: $D^p=Q\,D$ on $\bar R$.} Both sides are derivations, and they agree on $X$. For any
derivation $\delta$ of $\bar R$, applying $\delta$ to $Y^2=1-X^4$ gives $2Y\,\delta(Y)=-4X^3\,\delta(X)$. So
$2Y\,(D^pY-Q\,DY)=-4X^3\,(D^pX-Q\,DX)=0$, and since $2Y\ne0$ and $\bar R$ has no zero divisors, the two
derivations agree on $Y$ also. A derivation is determined by its values on $X$ and $Y$.

\emph{Step 4: $Q$ is the constant $A_p$.} The maps $D$ and $D^p=Q\,D$ commute. Now
$D(Q\,DX)=D(QY)=Q'Y^2-2X^3Q$ and $Q\,D(DX)=Q\,DY=-2X^3Q$. So $Q'\,(1-X^4)=0$, hence $Q'=0$. A polynomial
over $\mathbb F_p$ with zero derivative is a polynomial in $X^p$; as $Q$ has degree less than $p$, it is a
constant. The value of $D^pX=QY$ at $X=0$, $Y=1$ is $A_p$ modulo $p$. This proves \eqref{el:Dp}.

(i) Apply \eqref{el:Dp} to $f=D^{n-1}X$ and evaluate at $(0,1)$. (ii) By Lemma~Le4, $A_p=0$ when
$p\not\equiv1\bmod4$, so $D^pR\subset pR$. Since $D$ is additive, $D^{kp}R\subset p^kR$ for every $k\ge0$,
and with $k=\lfloor n/p\rfloor$ the element $D^nX=D^{n-kp}(D^{kp}X)$ lies in $p^kR$; evaluate. (iii) By
Proposition~Le20\,(ii), $A_p$ is prime to $p$. Every $n\equiv1\bmod4$ is $n_0+j(p-1)$ with
$n_0\in\{1,5,\ldots,p-4\}$ and $j\ge0$, because $p-1$ is divisible by $4$; and
$A_n\equiv A_p^{\,j}A_{n_0}$ by (i). For $p=5$: $n_0=1$, $A_1=1$ and $A_5=-12\equiv3$. The residues of
$A_1,A_5,\ldots,A_{p-4}$ are $(1,1,8)$ modulo $13$, $(1,5,15,14)$ modulo $17$ and
$(1,17,8,13,6,20,28)$ modulo $29$; none is zero.\footnote{Script \texttt{el11\_split\_recurrence.py},
which also computes $D^pX$ in $\bar R$ directly, as a check of \eqref{el:Dp}, for the odd primes $p<200$. The
statement \eqref{el:Dp} and its proof are due to a reviewer.} \qed

\noindent For example, for $p=5$,
\begin{equation}\label{D:494}
\begin{aligned}
 D^2X&=-2X^3,\qquad D^3X=-6X^2Y,\qquad D^4X=-12X(1-X^4)+12X^5=-12X+24X^5,\\
 D^5X&=(-12+120X^4)\,Y\equiv3Y\pmod5,
\end{aligned}
\end{equation}
and $A_5=-12\equiv3$. For $p=157$ one finds $A_{157}\equiv22$ and $A_{101}\equiv0\pmod{157}$, in agreement
with the observation above.

\noindent\emph{Remarks.} (1) Statement (ii) is the form in which the announcement states the divisibility
at an inert prime: the $p$-th power of the invariant derivation vanishes modulo $p$ (\S I.9, Step~2). Here
it is proved for the lemniscate curve. (2) The exponent $\lfloor n/p\rfloor$ of (ii) is weaker than the
exponent $\lfloor np/(p^2-1)\rfloor$ of Theorem~Le18: the two agree for $n<p^2-1$; among the $n\equiv1\bmod4$, the only ones with $A_n\ne0$, they first
differ at $n=p^2+3p-1$ (for example $n=17$ for $p=3$). So (ii) is a second, shorter proof of the
divisibility for $n<p^2+3p-1$, and Theorem~Le18 gives more beyond; (ii) does not replace it.

\subsection*{9.\quad The example of \S III.5: hypothesis \textup{(B)} in all degrees}

In this paragraph the inert prime is called $\ell$, as in \S III.5, and $q=\ell^2$. We recall the example.
Let $\ell\equiv3\bmod4$ and let $c\in\Z[i]$ with $c^{(q-1)/4}\equiv-1$ modulo $\ell$. Let $\wp_c$,
$\sigma_c$ be the Weierstrass functions with $g_2=4c$, $g_3=0$, let $t(z)=-2\wp_c(z)/\wp_c'(z)$, which is
the function $-X/Y$ on the curve $Y^2=X^3-cX$ when $X=\wp_c(z)$, $Y=\frac12\wp_c'(z)$, and
\begin{equation}\label{D:495}
 m(z)=\log\frac{\sigma_c(z)}{t(z)},\qquad \tilde\beta(z,w)=m(z+w)-m(z)-m(w).
\end{equation}
The curve $Y^2=X^3-cX$ is called a \emph{twist} of the curve $Y^2=X^3-X$, which is the case $c=1$ and
belongs to the lattice $\mathcal P$: the substitution $X=c^{1/2}X'$, $Y=c^{3/4}Y'$ turns one equation into
the other, but it uses a fourth root of $c$, which does not lie in $\Q(i)$ in general. For example, for
$\ell=3$ one may take $c=-i$: then $c^{(q-1)/4}=(-i)^2=-1$.
The series $\beta(X_0,X_1)$ is the sum of the terms $b_{a,b}X_0^aX_1^b$ of $\tilde\beta$ with
$a-\ell b\equiv0\bmod(q-1)$. Here $d=2$. Hypothesis \textup{(B4)} asks that the
two series $\exp(\Delta_i\beta)\circ L$ have coefficients of valuation $\ge0$, where $\Delta_1,\Delta_2$ are
the differences \eqref{dv:de} and $L$ replaces each variable $X$ by $L_+(X)$.

\textbf{The field.} Hypothesis \textup{(B)} is stated in \S III.5 for a field $K$ with a valuation $v$,
$v(\ell)=1$, a subfield $K_0$ with value group $\Z$, and a primitive root of unity $\zeta\in K$ of order
$\ell^d-1=q-1$. For the example we take $K=K_0=\Q_\ell(i)$. This field is defined as follows. For
$x\in\Z[i]$, $x\ne0$, let $v_\ell(x)$ be the largest $e$ with $x\in\ell^e\Z[i]$. Since $\ell\equiv3\bmod4$,
the ring $\Z[i]/\ell\Z[i]$ is a field with $q$ elements (see the proof of Lemma~Le14). So if $\ell$ divides a product in $\Z[i]$ it
divides a factor, and $v_\ell(xy)=v_\ell(x)+v_\ell(y)$. Put $v_\ell(x/y)=v_\ell(x)-v_\ell(y)$. This is the
\emph{$\ell$-adic valuation} of $\Q(i)$. On $\Q$ it is the valuation $v_p$ of \S5 for $p=\ell$, and the elements of
$\Q(i)$ with $v_\ell\ge0$ form the ring $\Z[i]_{(\ell)}$ of Lemma~Le14. For example $v_3(3+6i)=1$,
$v_3(2+i)=0$ and $v_3((1+i)/9)=-2$. (For a split prime the definition fails: $(2+i)(2-i)=5$.) The field
$\Q_\ell(i)$ is the completion of $\Q(i)$ for the absolute value $\ell^{-v_\ell}$, formed from Cauchy
sequences as $\mathbb R$ is formed from $\Q$. It contains $\Q(i)$, and $v_\ell$ extends to it with the same
value group $\Z$: if a Cauchy sequence $(x_n)$ in $\Q(i)$ does not tend to $0$, then $v_\ell(x_n)$ is constant from
some $n$ on, because $v_\ell(x_n-x_m)>v_\ell(x_n)$ implies $v_\ell(x_m)=v_\ell(x_n)$; this constant is the
value at the limit. It contains a primitive root of unity of order $q-1$: if the class of $x\in\Z[i]$
generates the cyclic group of the nonzero elements of the field $\Z[i]/\ell\Z[i]$, then
$x^{q^{k+1}}\equiv x^{q^k}$ modulo $\ell^{k+1}$ for all $k\ge0$ (by induction, since for $k\ge1$ a congruence $a\equiv b$ modulo
$\ell^k$ implies $a^\ell\equiv b^\ell$ modulo $\ell^{k+1}$), so the sequence $x^{q^k}$ converges, and its limit
$\zeta$ satisfies $\zeta^q=\zeta$ and $\zeta\equiv x$ modulo $\ell$; so $\zeta^{q-1}=1$, and
$\zeta^j\equiv x^j\not\equiv1$ modulo $\ell$ for $0<j<q-1$. For $\ell=3$ and $x=1+i$, an element of
order $8$ modulo $3$, one finds $\zeta\equiv7+7i$ modulo $9$, and indeed $(7+7i)^8=7^8\cdot16\equiv1$ modulo
$9$.\footnote{Script \texttt{el10\_examples.py}.} The completion and the root $\zeta$ enter only through the
wording of \textup{(B)}. Since $\mathrm{lw}(0)=1$ and $\mathrm{lw}(1)=-\ell$, condition \textup{(B2)} says
that $\beta(\zeta X_0,\zeta^{-\ell}X_1)=\beta(X_0,X_1)$, and this holds if and only if every monomial
$X_0^aX_1^b$ of $\beta$ has $a-\ell b\equiv0\bmod(q-1)$. All coefficients that occur below lie in $\Q(i)$,
and for these ``valuation $\ge0$'' means ``in $\Z[i]_{(\ell)}$''. So the proof works in $\Q(i)$ and never
uses the completion.

The proof has two parts. The first is a statement about elliptic functions: the third difference of
$\log(\sigma/\mathrm{sl})$ is the logarithm of a power series with integral coefficients. The second is
Lemma~Le14 for the curve $Y^2=X^3-cX$.

\textbf{Proposition Le21 (the cube function).} \emph{For $M(u)=\int_0^u\!\int_0^y\mathrm{sl}^2=\log(\sigma(u)/\mathrm{sl}(u))$ put}
\begin{equation}\label{D:496}
 \mathcal K(x,y,w)=\exp\Big(M(x+y+w)-M(x+y)-M(x+w)-M(y+w)+M(x)+M(y)+M(w)\Big).
\end{equation}
\emph{Then $\mathcal K$, written as a power series in $t_x=\mathrm{sl}(x)$, $t_y=\mathrm{sl}(y)$,
$t_w=\mathrm{sl}(w)$, has coefficients in $\Z[\frac12]$ and the constant term $1$. Explicitly, with
$t_{x+y}=F(t_x,t_y)$ and so on, $s=x+y+w$, and $t'=\sqrt{1-t^4}$,}
\begin{equation}\label{el:cube}
 \frac1{\mathcal K}=\frac{t_s\,t_xt_yt_w}{t_{x+y}\,t_{x+w}\,t_{y+w}}
 \Big(\frac{t_x'}{t_x}+\frac{t_y'}{t_y}+\frac{t_w'}{t_w}-\frac{t_s'}{t_s}-t_xt_y\,t_{x+y}-t_{x+y}\,t_w\,t_s\Big).
\end{equation}

\emph{Proof.} The functions $\sigma$, $\zeta_{\rm w}$ are those of the lattice $\mathcal P$, and
$M=\log(\sigma/\mathrm{sl})$ near $0$ by \eqref{el:theta}.

\emph{Step 1: $\zeta_{\rm w}=\mathrm{sl}'/\mathrm{sl}+J$.} By \eqref{el:logsl} the derivative
of $\mathrm{sl}'/\mathrm{sl}+J$ is $-\mathrm{sl}^{-2}=-\wp$, which is the derivative of $\zeta_{\rm w}$
(Lemma~Le5a\,(b)); and both
$\zeta_{\rm w}(u)-1/u$ and $\mathrm{sl}'(u)/\mathrm{sl}(u)+J(u)-1/u$ tend to $0$ as $u\to0$
(Lemma~Le5a\,(a), and $\mathrm{sl}(u)=u+O(u^5)$).

\emph{Step 2: $J(x+y)=J(x)+J(y)+\mathrm{sl}(x)\,\mathrm{sl}(y)\,\mathrm{sl}(x+y)$.} Let $\delta(x,y)$ be the
difference of the two sides, and $s_1=\mathrm{sl}(x)$, $s_2=\mathrm{sl}(y)$ as in the proof of Lemma~Le5.
By \eqref{el:add},
\begin{equation}\label{D:498}
 \mathrm{sl}(x+y)\,(s_1's_2-s_1s_2')=\frac{s_2^2s_1'^2-s_1^2s_2'^2}{1+s_1^2s_2^2}
 =\frac{s_2^2(1-s_1^4)-s_1^2(1-s_2^4)}{1+s_1^2s_2^2}=s_2^2-s_1^2 .
\end{equation}
Hence $(\partial_x-\partial_y)\,\delta=-s_1^2+s_2^2-(s_1's_2-s_1s_2')\,\mathrm{sl}(x+y)=0$. So $\delta$
depends on $x+y$ only, and $\delta(x,0)=0$.

\emph{Step 3: a classical identity.} For $x,y,w\in\C$ in general position,
\begin{equation}\label{el:zetasigma}
 \zeta_{\rm w}(x)+\zeta_{\rm w}(y)+\zeta_{\rm w}(w)-\zeta_{\rm w}(x+y+w)
 =\frac{\sigma(x+y)\,\sigma(x+w)\,\sigma(y+w)}{\sigma(x)\,\sigma(y)\,\sigma(w)\,\sigma(x+y+w)} .
\end{equation}
Fix $x,y$ and regard both sides as functions of $w$. By Lemma~Le5a\,(d) the left side is elliptic for
$\mathcal P$: under $w\mapsto w+\omega$ the two terms that contain $w$ change by $\eta(\omega)$ and
$-\eta(\omega)$. By Lemma~Le5a\,(a) it has
simple poles at $w\in\mathcal P$ (residue $1$) and at $w\in-x-y+\mathcal P$ (residue $-1$); it vanishes at
$w=-x$ and at $w=-y$, since $\zeta_{\rm w}$ is odd (Lemma~Le5a\,(c)). By \eqref{el:sigmaqp} the right side is elliptic in $w$: under
$w\mapsto w+\omega$ numerator and denominator are multiplied by the same factor
$C_\omega^2e^{\eta(\omega)(x+y+2w)}$. By Lemma~Le5a\,(a) it has simple poles at the same points, simple
zeros at $-x+\mathcal P$, $-y+\mathcal P$, and no other zero. So the quotient of the left side by the right
side is an elliptic function of $w$ without poles; it is constant in $w$ (\cite{Apo90}, Theorem~1.4). As
$w\to0$ both sides are $1/w+O(1)$, because $\sigma(w)/w\to1$. So the constant is $1$.

\emph{Step 4.} Since $M=\log(\sigma/\mathrm{sl})$ by \eqref{el:theta}, the function $\mathcal K$ is the
reciprocal of the right side of \eqref{el:zetasigma} times $t_{x+y}t_{x+w}t_{y+w}/(t_st_xt_yt_w)$. Insert
Step~1 in the left side of \eqref{el:zetasigma}, and use Step~2 twice:
$J(s)-J(x)-J(y)-J(w)=t_{x+y}t_wt_s+t_xt_yt_{x+y}$. This is \eqref{el:cube}.

\emph{Step 5: integrality.} By Proposition~Le10 the series $t_{x+y}$, $t_s$, $t_x'$, $t_s'=\sqrt{1-t_s^4}$
have coefficients in $\Z[\frac12]$, and $t_{x+y}=(t_x+t_y)V(t_x,t_y)$ with $V$ invertible. Multiply
\eqref{el:cube} by $t_{x+y}t_{x+w}t_{y+w}$: the result is $t_st_xt_yt_w$ times the bracket, a power series
with coefficients in $\Z[\frac12]$. So $(t_x+t_y)(t_x+t_w)(t_y+t_w)/\mathcal K$ has coefficients in
$\Z[\frac12]$. Now $1/\mathcal K$ is a power series with rational coefficients and constant term $1$. If
$(t_x+t_y)B$ has coefficients in $\Z[\frac12]$ for a power series $B$, so has $B$: substitute
$t_y=Y_1-t_x$ with a new variable $Y_1$, which is an invertible change of variables over $\Z$, and divide by $Y_1$. Three such
steps show that $1/\mathcal K$, and then $\mathcal K$, has coefficients in $\Z[\frac12]$. \qed

\noindent A computation verifies \eqref{el:cube}, Step~2, and the integrality of $\mathcal K$ up
to total degree $22$:\footnote{Script \texttt{el7\_cube.gp}.} the denominators are powers of $2$, while the exponent
$\log\mathcal K$ has the primes $3$ and $5$ in its denominators already in this range.
$\mathcal K=1+t_xt_yt_w(t_x+t_y+t_w)+\cdots$.

\textbf{Proposition Le22 (hypothesis (B) for the example of \S III.5).} \emph{Let $\ell\equiv3\bmod4$ and
$c\in\Z[i]$ with $c^{(\ell^2-1)/4}\equiv-1\pmod\ell$. Then the series $\beta$ recalled above satisfies
\textup{(B1)}, \textup{(B2)}, \textup{(B3)} and \textup{(B4)} of \S III.5, in all degrees.}

\emph{Proof.} Here $K=K_0=\Q_\ell(i)$ with the valuation $v_\ell$, as explained above. Let
$\mathcal O=\Z[i]_{(\ell)}$ as in Lemma~Le14, and fix a fourth root $c^{1/4}$ in $\C$.

\emph{Step 1: reduction to $\mathrm{sl}$.} The lattice with $g_2=4c$, $g_3=0$ is $c^{-1/4}\mathcal P$,
because $g_2(\lambda\mathcal L)=\lambda^{-4}g_2(\mathcal L)$. Put $\tilde z=c^{1/4}z$. By the series for
$\wp$ and by Lemma~Le5a\,(c) with $\lambda=c^{-1/4}$,
$\wp_c(z)=c^{1/2}\wp(\tilde z)$ and $\sigma_c(z)=c^{-1/4}\sigma(\tilde z)$. Since $\wp=\mathrm{sl}^{-2}$,
we have $-2\wp/\wp'=\mathrm{sl}/\mathrm{sl}'$. Hence
\begin{equation}\label{el:twist}
 t(z)=c^{-1/4}\,\frac{\mathrm{sl}(\tilde z)}{\mathrm{sl}'(\tilde z)},\qquad m(z)=M(\tilde z)+\log\mathrm{sl}'(\tilde z),\qquad \tilde z=c^{1/4}z .
\end{equation}
\emph{Step 2: the logarithm of the curve $Y^2=X^3-cX$.} Let
$\tau=\mathrm{sl}(\tilde z)/\mathrm{sl}'(\tilde z)=c^{1/4}t$. By
Proposition~Le2\,(a), $d\tau/d\tilde z=(1+\mathrm{sl}^4)/(1-\mathrm{sl}^4)$ and
$1+4\tau^4=(1+\mathrm{sl}^4)^2/(1-\mathrm{sl}^4)^2$. So $d\tilde z/d\tau=(1+4\tau^4)^{-1/2}$, and
\begin{equation}\label{el:lc}
 z=l_c(t),\qquad l_c(t)=\sum_{k\ge0}\frac{\kappa_k\,(-4c)^k}{4k+1}\,t^{4k+1}=\sum_{k\ge0}\binom{2k}{k}\frac{(-c)^k}{4k+1}\,t^{4k+1}.
\end{equation}
We claim $l_c(t)-\frac1\ell\,l_c(t^q)\in\mathcal O[[t]]$. Let $m_0=4k+1$ and $a=v_\ell(m_0)$. For $a\le1$ the
coefficient of $t^{m_0}$ is in $\mathcal O$ by Theorem~Le12\,(i). Let $a\ge2$ and $m_0=q(4k'+1)$. The
coefficient is $(-4c)^{k'}\big(\kappa_k(-4c)^{k-k'}-\ell\,\kappa_{k'}\big)/m_0$. By \eqref{el:asd},
$\kappa_k\equiv-\ell\kappa_{k'}$ modulo $\ell^a$, so it suffices that $(-4c)^{k-k'}\equiv-1$ modulo
$\ell^{a-1}\mathcal O$. Now $k-k'=\frac{q-1}4(4k'+1)$. The number $\omega=(-4c)^{(q-1)/4}$ is $\equiv-1$
modulo $\ell$: indeed $(-4)^{(q-1)/4}=2^{(q-1)/2}=(2^{\ell-1})^{(\ell+1)/2}\equiv1$, because $(q-1)/4$ is even, and
$c^{(q-1)/4}\equiv-1$ by hypothesis. And $4k'+1$ is odd and divisible by $\ell^{a-2}$. So
$\omega^{4k'+1}=-(-\omega)^{4k'+1}$ with $-\omega\equiv1$ modulo $\ell$, and $(-\omega)^{\ell^{a-2}}\equiv1$
modulo $\ell^{a-1}$. This proves the claim. It is here, and only here, that the condition on $c$ is used.

\emph{Step 3.} By Lemma~Le14\,(b) with $\varepsilon=+1$, applied to $l_c$ and $L_+$, the series
$h_c=l_c^{-1}\circ L_+$ has coefficients in $\mathcal O$. By \eqref{el:lc} it is the function $t$ in the
coordinate of the logarithm: $t(L_+(X))=h_c(X)$.

\emph{Step 4: the cube function in the coordinate $t$.} Let $\mathcal K_m(x,y,w)$ be defined as
$\mathcal K$ in Proposition~Le21, with $m$ in place of $M$ and with $x,y,w$ values of $z$. By \eqref{el:twist} it is the product of
$\mathcal K$ and of the same expression formed with $\mathrm{sl}'$, both taken at $c^{1/4}x,c^{1/4}y,c^{1/4}w$. In the
coordinates $T=\mathrm{sl}(\tilde z)$ the first is integral over $\Z[\frac12]$ by Proposition~Le21. So is the
second: $\mathrm{sl}'(\tilde z_1+\tilde z_2)=\sqrt{1-F(T_1,T_2)^4}$ has coefficients in $\Z[\frac12]$ and constant term
$1$. Next, $\tau=T(1-T^4)^{-1/2}$ has an inverse series $T=\tau+\cdots$ over $\Z[\frac12]$. So
$\mathcal K_m=\mathcal K^\ast(\tau_x,\tau_y,\tau_w)$ with a series $\mathcal K^\ast$ over $\Z[\frac12]$,
where $\tau_x=c^{1/4}t_x$ and so on. The substitution $\tilde z\mapsto i\tilde z$ in all three variables multiplies each
$\tau$ by $i$ and leaves $M$ and $\mathrm{sl}'$ unchanged, by Proposition~Le2\,(c). So $\mathcal K^\ast$
has terms only in total degrees divisible by $4$, and a term of degree $4n$ in $\tau$ is $c^n$ times a
monomial in $t_x,t_y,t_w$. Therefore $\mathcal K_m$, as a power series in $t_x=t(x)$, $t_y=t(y)$,
$t_w=t(w)$, has coefficients in $\Z[\frac12,c]\subset\mathcal O$.

\emph{Step 5: \textup{(B4)} for $\tilde\beta$.} From the definitions,
$\Delta_1\tilde\beta(x,y;w)=\tilde\beta(x+y,w)-\tilde\beta(x,w)-\tilde\beta(y,w)=\log\mathcal K_m(x,y,w)$.
With Steps 3 and 4,
\begin{equation}\label{D:502}
 \mathcal E(X,Y,W):=\exp(\Delta_1\tilde\beta)\big(L_+(X),L_+(Y),L_+(W)\big)
\end{equation}
is the series $\mathcal K_m$ of Step~4 with $h_c(X),h_c(Y),h_c(W)$ substituted for $t_x,t_y,t_w$. It is
a power series over $\mathcal O$ with constant term $1$.

\emph{Step 6: the invariant part.} The series $\beta$ consists of the monomials of $\tilde\beta$ whose
exponents satisfy a congruence modulo $q-1$. Such a part of a series is cut out by averaging over roots of
unity of order $q-1$. We need them in a ring in which integrality can still be read off, and we take the
following one. Let $\Phi_{q-1}(x)\in\Z[x]$ be the cyclotomic polynomial of order $q-1$. Recall that it is
monic, and that $x^N-1=\prod_{\nu\mid N}\Phi_\nu(x)$ for every $N\ge1$. Let
\begin{equation}\label{D:503}
 \mathcal A=\mathcal O[x]/(\Phi_{q-1}(x)),\qquad \mathcal A_{\Q}=\Q(i)[x]/(\Phi_{q-1}(x)),\qquad
 \psi=\text{the class of }x .
\end{equation}
We write $\psi$ and not $\zeta$: the letter $\zeta$ stands for the root of unity in $\Q_\ell(i)$ of the
paragraph ``The field'', and $\psi$ is a different object. Let $D$ be the degree of $\Phi_{q-1}$. Division with remainder by the monic polynomial $\Phi_{q-1}$ shows
that every element of $\mathcal A_\Q$ is $c_0+c_1\psi+\cdots+c_{D-1}\psi^{D-1}$ with unique
$c_0,\ldots,c_{D-1}\in\Q(i)$, and that it lies in $\mathcal A$ exactly when all $c_j$ lie in
$\mathcal O$. In particular
\begin{equation}\label{el:AO}
 \mathcal A\cap\Q(i)=\mathcal O .
\end{equation}
For example, for $\ell=3$: $\Phi_8=x^4+1$, the elements of $\mathcal A$ are
$c_0+c_1\psi+c_2\psi^2+c_3\psi^3$ with $c_j\in\Z[i]_{(3)}$, and $\psi^4=-1$. The ring $\mathcal A_\Q$ need not be a field, and $\psi$ need not
be a root of unity in a field: for $\ell=3$ we have $x^4+1=(x^2-i)(x^2+i)$ over $\Q(i)$, so
$(\psi^2-i)(\psi^2+i)=0$ in $\mathcal A_\Q$ with both factors different from $0$. No field is needed: what
follows uses only the basis $1,\psi,\ldots,\psi^{D-1}$, the two relations $\psi^{q-1}=1$ and \eqref{el:orth}, and rules for
power series over a commutative ring. In $\mathcal A$ we have
$\psi^{q-1}=1$, because $\Phi_{q-1}$ divides $x^{q-1}-1$, and for every integer $k$
\begin{equation}\label{el:orth}
 \sum_{j=0}^{q-2}\psi^{jk}=q-1\quad\text{if }(q-1)\mid k,\qquad \sum_{j=0}^{q-2}\psi^{jk}=0\quad\text{otherwise.}
\end{equation}
The first case is clear. For the second let $k_\ast$ be the greatest common divisor of $k$ and $q-1$, so $k_\ast<q-1$.
Modulo $q-1$ the numbers $jk$, $0\le j\le q-2$, run $k_\ast$ times through the multiples of $k_\ast$. So
$\sum_jx^{jk}$ is congruent modulo $x^{q-1}-1$ to $k_\ast\,(x^{q-1}-1)/(x^{k_\ast}-1)$, and
$(x^{q-1}-1)/(x^{k_\ast}-1)$ is the product of the $\Phi_\nu$ with $\nu\mid q-1$, $\nu\nmid k_\ast$; one of them is
$\Phi_{q-1}$. (For $\ell=3$, $k=2$: $1+\psi^2+\psi^4+\psi^6+1+\psi^2+\psi^4+\psi^6=2(1+\psi^2)(1+\psi^4)=0$.)

By \eqref{el:orth}, as an identity of power series with coefficients in $\mathcal A_\Q$,
\begin{equation}\label{D:506}
 \beta(X_0,X_1)=\frac1{q-1}\sum_{j=0}^{q-2}\tilde\beta\big(\psi^jX_0,\psi^{-\ell j}X_1\big) ,
\end{equation}
because the monomial $X_0^aX_1^b$ of $\tilde\beta$ is multiplied by $\frac1{q-1}\sum_j\psi^{j(a-\ell b)}$.
The same holds for $\Delta_1$ of both sides. Since $\psi^{q}=\psi$, we have
$L_+(\psi^jX)=\sum_e(\psi^j)^{q^e}X^{q^e}/\ell^e=\psi^jL_+(X)$, and in the same way
$\psi^{-\ell j}L_+(W)=L_+(\psi^{-\ell j}W)$. So by Step~5
\begin{equation}\label{D:507}
 (\Delta_1\beta)\big(L_+(X),L_+(Y),L_+(W)\big)=\frac1{q-1}\sum_{j=0}^{q-2}\log\mathcal E\big(\psi^jX,\psi^jY,\psi^{-\ell j}W\big).
\end{equation}
The ring $\mathcal A_\Q$ is commutative and contains $\Q$, and the rules $\exp(a+b)=\exp(a)\exp(b)$ and
$\exp(\log(1+a))=1+a$ hold for power series without constant term over any such ring. So
\begin{equation}\label{D:508}
 \exp(\Delta_1\beta)\big(L_+(X),L_+(Y),L_+(W)\big)=\exp\Big(\frac1{q-1}\log\mathcal N\Big),\qquad
 \mathcal N=\prod_{j=0}^{q-2}\mathcal E\big(\psi^jX,\psi^jY,\psi^{-\ell j}W\big).
\end{equation}
By Step~5 each factor of $\mathcal N$ has coefficients in $\mathcal A$ and the constant term $1$; so has
$\mathcal N$. Now $\exp\big(\frac1{q-1}\log(1+Y)\big)=1+\sum_{n\ge1}c_nY^n$ is the series whose $(q-1)$th
power is $1+Y$, and its coefficients lie in $\Z[\frac1{q-1}]\subset\Z_{(\ell)}$: comparing the coefficients
of $Y^n$ in the $(q-1)$th power gives $(q-1)c_n$ as a polynomial with integer coefficients in
$c_1,\ldots,c_{n-1}$ (plus $1$ if $n=1$). Substituting $Y=\mathcal N-1$ we see that
$\exp(\Delta_1\beta)\circ L$ has coefficients in $\mathcal A$. They also lie in $\Q(i)$, because $\beta$ has
coefficients in $\Q(i)$. By \eqref{el:AO} they lie in $\mathcal O$. This is \textup{(B4)} for $\Delta_1$.
For $\Delta_2$ the argument is the same, because $\tilde\beta(z,w)=\tilde\beta(w,z)$ and
$\Delta_2\tilde\beta(x;w,w')=\log\mathcal K_m(x,w,w')$.

\emph{Step 7: \textup{(B1)--(B3)}.} The coefficients of $\beta$ lie in $\Q(i)\subset K_0=\Q_\ell(i)$, whose
value group is $\Z$: this is (B1). (B2) holds by the definition of $\beta$, with $\mathrm{lw}(0)=1$ and
$\mathrm{lw}(1)=-\ell$, as explained before the proposition. For (B3): $m$ is an even power series with $m(z)=-\frac{5c}{12}z^4+\cdots$, so
every monomial of $\tilde\beta$ has positive degree in each of $z$ and $w$ and total degree at least $4$.
\qed

\noindent \emph{Remarks.} (1) \emph{The role of $c$.} For $c=1$ the curve is the lemniscate itself in the
parameter $\mathrm{sl}/\mathrm{sl}'$, and Step~2 gives the sign $\varepsilon=-1$, as in Theorem~Le13. The
condition $c^{(\ell^2-1)/4}\equiv-1$ changes the sign, so that the logarithm \eqref{dv:log}, which has the
signs $+$, applies. With the alternating logarithm $L_-$ in \eqref{dv:log} the curve with $c=1$ would do as
well.

(2) \emph{Tests.}\footnote{Scripts \texttt{el5\_example\_B4.gp}, \texttt{el7\_cube.gp},
\texttt{el8\_twist\_log.gp}.} The coefficient of $X^NYW$ in
$\exp(\Delta_1\beta)\circ L$ equals the coefficient of $X^N$ in the series obtained from $m''$ by keeping
the terms of degree $\equiv\ell-1\bmod(q-1)$ and composing with $L_+$. For $\ell=3$, $c=-i$ it is $\ell$-integral for all $N\le170$, and for $\ell=7$, $c=-1-i$ for all
$N\le150$; with $L_-$ in place of $L_+$, or with $c=1$, it fails at $N=10$ for $\ell=3$ and at $N=54$ for
$\ell=7$. The series $\mathcal K_m$ in the coordinates $t$ has no denominators at all up to total degree
$22$, for $c=-i$, $-1-i$ and $1$. The series $h_c$ is integral up to degree $400$ for $(\ell,c)=(3,-i)$
and up to degree $260$ for $(7,-1-i)$ and $(11,-i)$, and the functional equation of $l_c$ of Step~2 holds
up to degree $20000$ in these three cases; with the sign $\varepsilon=-1$ both fail at the degree $\ell^2$, and
for $c=1$ the two signs change roles.

(3) \emph{What is not covered.} The statement of \S III.5 that equality holds in \eqref{dv:beta} at the
exceptional monomials of this $\beta$, and the tests of Theorem~J on it, are not treated. With
Proposition~Le22 the conclusions of Lemmas~3 and~4 of \S III.5 hold for this $\beta$ in all degrees. So
the term $\mathrm{ex}(n)$ of \eqref{dv:beta} cannot be omitted: equality holds at every exceptional monomial
of degree $\le800$, by the computation reported there. That equality holds at all exceptional monomials is
not proved. No example with $d\ge4$ is given here.

\subsection*{10.\quad Correspondence with the Lean development}

Nothing in this section is taken from the Lean development, and no statement of the Lean development is
proved here. The table records which objects of this section also occur there. The files are in the
directory \texttt{Formalism/GaussJetBounds} (the second file of the middle row in its subdirectory
\texttt{GaussJetProofs}). Only the statements were read. The last column says whether the name occurs in
the lists of declarations for which the dependency of the final theorem was computed. All of them do except \texttt{de1} and \texttt{de2},
marked as used (\texttt{genLog}, \texttt{GenLogWeight}, \texttt{generalized\_log\_weight}, \texttt{BHyp},
\texttt{letterLog}); as explained in Part~III, this is necessary and not sufficient for a use in the
argument. (By \S III.5 the final theorem depends on
\texttt{weightQuantization}, \texttt{lemmaS\_of} and \texttt{thmAOp}; how these use the three names of
the table was not followed.)

\begin{center}\small
\begin{tabular}{p{0.27\textwidth}p{0.47\textwidth}p{0.16\textwidth}}
\toprule
this section & Lean name, file:line & in the dependency lists\\
\midrule
$L_+$ of \eqref{el:Lpm} & \texttt{genLog} with $\pi=\mathrm{id}$; \texttt{GenLogInverse.lean}:12 & yes\\[2pt]
the bound \eqref{el:ginv}, that is $v(g_n)\ge-(n-1)\rho$ with $\rho=1/(q-1)$, implied by Proposition~Le15 & \texttt{GenLogWeight}, \texttt{GenLogInverse.lean}:24; its proof
\texttt{generalized\_log\_weight}, in \texttt{LogarithmWeightComparison.lean}:252 & yes\\[2pt]
hypothesis (B), verified in Proposition~Le22 for one $\beta$ & \texttt{BHyp}, \texttt{Statement.lean}:85; with \texttt{letterLog} (:74), \texttt{de1}, \texttt{de2} & yes (\texttt{de1}, \texttt{de2}: not in the lists)\\
\bottomrule
\end{tabular}
\end{center}

\noindent The Lean development contains no statement about the lemniscate, as far as a search for the
names above and their uses shows; Proposition~Le22 concerns the example of \S III.5, which is ours and
not the source's. For the bound \eqref{el:Hur} with the factorial, and for Theorem~Le18, we found no
counterpart in the Lean development.

\textbf{What this section proves, and what it cites.} \emph{Proved here:} all numbered statements Le1 to
Le22, with Lemma~Le5a, Corollary~Le6a and Proposition~Le20a. \emph{Cited:} the four items (C1) to (C4) of the list at the
beginning of the section and nothing else: the theory of elliptic functions and of $\wp$ from \cite{Apo90},
Chapter~1; the two-squares theorem, used only in the first equality of \eqref{el:T0}; Lucas's congruence,
used only in Proposition~Le20 (and through it in Proposition~Le20a\,(iii)); and Legendre's formula, the theorem of Euler and Fermat and Chebyshev's
bound from a first course. The properties of $\sigma$ and $\zeta_{\rm w}$ are proved in \S1a. The definitions
of \S III.5 (the hypothesis \textup{(B)} and the example) are used in \S9, but no result of \S III.5 is used.
None of the statements is claimed to be new; see the remark after Proposition~Le15. \emph{Observed and
not proved:} that the denominator of $r_k$ is exactly $4^k$; that \eqref{el:asd} holds modulo $p^{a+1}$;
that equality holds in \eqref{el:Hur} for every $n$; that the excess of $v_p(A_n)$ over $n\,p/(p^2-1)$ is
unbounded. (That $5,13,17,29$ divide no $A_n$ with $n\equiv1\bmod4$ is proved: Proposition~Le20a\,(iii).)


\section{What the proof takes from outside analytic number theory: a census}\label{cs:sec}

This section proves nothing. It is a list. The proof of the main inequality (M) in the Lean development
is given many theorems from parts of mathematics that a first course in analytic number theory does not
touch. At one place the development collects them: 195 statements, of which 177 carry in their names
the words ``of literature'' (176 end in them). We opened every one and wrote down, in classical
language, what it says, which published theorem it points to, and to which part of the proof of (M) it is
passed. \emph{How it was obtained.} Each statement is a named proposition defined elsewhere in the
development. We read each definition and, where it is a conjunction of further named
propositions, those too, about ninety in all. Notions below that level (for instance ``is a cusp form'')
were \emph{read by name}: we took the name to mean what it says. An independent checker then re-read
106 of the 195 rows against the development, one to three levels further down. In its sample of
102 rows it judged four wrong and 42 incomplete; among the 45 rows drawn at random
from the 128 then graded ``medium'', one wrong and 19 incomplete. Every row in which
it found a fault has been corrected here. The 79 other rows of the first eleven groups are unchanged, and at that
rate many of them must be incomplete in the same ways: a hypothesis left out, or a special case not said
to be one. \emph{What it can and cannot tell.} The tables of subsection~2 show, on a few pages, how much is
taken from outside and from where. (1)~\emph{The statement.} Each row is
our reading of one statement of the development. (2)~\emph{The source.} The publication named in a row is
the one we believe the statement corresponds to. Every reference in this section is from memory. None was
compared with the printed text, and theorem numbers in particular may be wrong. Where the confidence is
not ``high'' we do not claim that the formal statement equals the published one, and in several rows it is
known to be weaker (subsection~4). (3)~\emph{The proof.} In the development these statements are not
assumptions: each is a theorem with a proof there. We read none of these proofs. So this section does not
say that the published theorems have been formally verified, and it does not say that they are proved in
this paper. They are not. It says what the proof of (M) is given, at the place where the development
gathers its inputs. It is not needed for the logic of Parts~I--III; it makes precise the list (b) of
\S I.11.

\emph{Dependence.} A computation on the compiled development shows that the final theorem depends on each of
the 195 statements. Here ``depends on'' means membership in a dependency closure: the statement is among
the declarations reached from the final theorem by following references. That is necessary for a statement
to be consulted in the proof. It is not sufficient: the statements are packed into records, the records
are passed to one theorem, and so the final theorem depends on all of them whether or not their content is
read. Two examples. First, row~122 is the proposition ``True''; the
development itself names it an unused slot. The final theorem depends on it, and the entry of the record
that holds it is passed on like any other. Second, the ten rows 9, 11, 54, 55, 59, 97, 106, 140, 147, 148. Eight of them are entries of the
first large record (row~192), and the other two
(97, 147) occur only in the
one-line proofs of these. That record travels through five theorems of the development, and in the source
text of those five theorems these eight entries are never taken out of it. So the ten rows are never read in
the theorems they are passed through, as far as this was followed; yet the final theorem depends on each
of them. (The search was of the source text; a use that does not name the entry would escape it.) For the other rows, to learn which are consulted, and for what, one would have to read the
receiving proofs. We did not. What was traced is where each row is \emph{passed}.

The full statements, the reasons for the grades, the checker's verdicts and the names of the statements in
the development are in the companion file ``zeta-paper-census.pdf''. This section uses no identifier of the
development and no notation beyond the codes of subsection~1.

\subsection*{1.\quad Where the rows are passed}\label{cs:codes}

The development reaches its statement for all Dirichlet characters by one implication
(row~195). It applies a single theorem to seven arguments: two large
records of statements (rows~192, 193); the two halves of
row~121; a bound for covolumes, itself obtained
from a third record (row~194) and fourteen further statements, of which
twelve are rows and two are theorems of the development that are not rows; the local correspondence of
row~66; and the bound for Galois cohomology, Theorem~G of
\S III.9. We read the proof of that theorem and of the theorems it calls, two levels in all. Each takes
named entries of the records and produces a named intermediate statement. Comparing these with
\S III.7--III.10 gives the codes below. For S3, S5, S6 and S7 the checker verified that the intermediate
statements are exactly the hypotheses of the development that \S III.7 lists under these inputs.

\begin{itemize}[leftmargin=3.2em,itemsep=0pt,topsep=2pt]
\item[S3] The size bound, input (S3) of \S III.7: the lower bound for the eigenvalues of a lift (the
 inner-product formula and its local factors). The result is passed on towards (S7) as well.
\item[S5] Input (S5): the local index bounds in the slope bound at the lower level. In the development this
 statement has as a hypothesis the bound for denominators that \S III.9 derives from Theorem~G.
\item[S6] Input (S6): the covolume bound that is the first hypothesis of the theorem quoted under (S5).
\item[S7] Input (S7): the bound for the maximal slope in the passage between levels.
\item[R] The trace-formula statements of \S III.8: (R1) for row~114
 (compare (R2)), and (R4) for row~29.
\item[G] The bound for Galois cohomology, but not Theorem~G: the inequality \eqref{gc:link} of \S III.9 for
 the sum over twists.
\item[O] The construction of the objects of \S III.10: the statements that the spaces, lattices, forms and
 weights of the construction exist and have their properties. They were not separated further, and
 not matched with the single Hypotheses O12, O19--O21, O23--O25.
\item[T] Statements on theta lifts that are passed together, as one record, to theorems on test vectors,
 on the separation of linear forms, and on the fit of Fourier coefficients; their conclusions go to the
 construction (O) and to the input (S7). We could not tell which row goes to which.
\item[LL] The local correspondence of row~66, a parameter of
 the whole construction. Rows marked LL are passed to its construction in the development.
\item[n.r.] Never read: one of the ten rows of the second example above.
\end{itemize}

Several codes mean that the row is passed to several places. A star, as in ``(G, O)$^*$'', means that the
row is not itself passed to any of the theorems we read: it occurs in the one-line proof of another row
and inherits the codes of that row. The codes say where a statement is \emph{passed}. Whether the
receiving proof uses it, and how, was not checked. Theorem~G itself receives none of the rows. This holds
through the file that gathers the rows only: there Theorem~G is deduced from two of the three open
statements of subsection~4. Its proof was not read, and theorems of the development that stand behind
rows may be cited inside it under their own names. No row is assigned to the divisibility statements of
\S III.5--III.6: the theorems we read do not show the connection.

\subsection*{2.\quad The census}\label{cs:tab}

The rows are numbered from 1 to 195 and grouped by subject; the grouping is ours. Each group has a
paragraph for a reader who has not met the subject, with the idea and one small computed example. Each
table has one line for each row: its number; the theorem in a few words; the author and year of the
published source, from memory in every row; the codes of subsection~1; and our confidence that the
statement of the development, as far as it was read, is the published theorem. The grades are ``high''
(the statement is explicit and we recognise it), ``med.''\ (it has the shape of the published theorem, but
notions in it were read by name, or it is a special case) and ``low'' (it is a bundle read in part, or a
special or quantitative form that we could not match with a printed statement, or it is much weaker than
the theorem cited). ``Not matched'' means that we know no printed statement of this form.

\smallskip\noindent\textbf{Group A. Analytic number theory and $L$-functions} (rows 1--6).\label{cs:grp:A}
\emph{What it is.} The reader's own subject, carried from $\mathbb Q$ to a number field. A \emph{Hecke
character} is to a number field what a Dirichlet character is to $\mathbb Q$: a multiplicative function on
the ideals prime to a modulus whose value on a principal ideal $(\alpha)$ with $\alpha\equiv1$ is given by a
fixed formula in the complex embeddings of $\alpha$. For $\mathbb Q(i)$ the function
$(\alpha)\mapsto(\alpha/|\alpha|)^{4k}$ is one; it is well defined because the units $\pm1,\pm i$ have fourth
power $1$. Its $L$-function has an Euler product, a continuation and a functional equation.
\emph{A typical theorem} is row~2, Linnik's theorem: the least prime in a class
modulo $q$ is at most $Cq^L$. (For $q\le300$ it is always below $q^2$; the largest ratio is $7/9$, for
$7\equiv1\pmod3$.) \emph{Where it goes.} The theta series of the construction belong to lattices with
complex multiplication, and the $L$-functions attached to such series are those of Hecke characters. All
six rows are passed with the statements on theta lifts (T), three of them also to G and~O; what the
receiving proofs do with them was not traced.

{\scriptsize\renewcommand{\arraystretch}{0.9}\setlength{\tabcolsep}{4pt}\setlength{\LTpre}{2pt}\setlength{\LTpost}{2pt}
\begin{longtable}{@{}r@{\hspace{6pt}}p{.43\textwidth}p{.285\textwidth}p{.15\textwidth}l@{}}
\toprule
no. & theorem, in a few words & author and year (from memory) & passed to & conf.\\
\midrule
\endhead
1 & \raggedright $L'/L(1,\chi)$ and the zeros, uniformly in the modulus & \raggedright after Davenport, Ch.~12 & \raggedright T & med.\tabularnewline
2 & \raggedright Least prime in a progression is $\le Cq^L$ & \raggedright Linnik 1944 & \raggedright T & high\tabularnewline
3 & \raggedright Functional equation of Hecke $L$-functions (CM fields) & \raggedright Hecke 1918--20; Tate 1950 & \raggedright G, O, T & med.\tabularnewline
4 & \raggedright Functional equation of partial zeta functions & \raggedright Hecke 1917 & \raggedright (G, O, T)$^{*}$ & med.\tabularnewline
5 & \raggedright Hecke $L$-functions are entire away from unitary norm powers; $L(1,\chi)\ne0$ for $\chi\ne1$ & \raggedright Hecke 1920 & \raggedright T & med.\tabularnewline
6 & \raggedright The same with Euler factors removed; $\re s\ge1$ & \raggedright Hecke 1920; Landau 1903 & \raggedright G, O, T & med.\tabularnewline
\bottomrule
\end{longtable}}

\smallskip\noindent\textbf{Group B. Algebraic number theory, class field theory, complex multiplication} (rows 7--27).\label{cs:grp:B}
\emph{What it is.} A number field $K$ has a ring of integers in which ideals, not numbers, factor uniquely
into primes. The ideal classes form a finite group; its order is the class number, $1$ for $\mathbb Q(i)$
and $3$ for $\mathbb Q(\sqrt{-23})$. For a Galois extension $L/K$ and a prime $\mathfrak p$ of $K$
unramified in $L$ there is a \emph{Frobenius element} in the Galois group, which acts as
$x\mapsto x^{N\mathfrak p}$ modulo a prime above $\mathfrak p$. For $\mathbb Q(i)/\mathbb Q$ it is trivial
exactly when $p\equiv1\pmod4$, that is, when $p=a^2+b^2$ (eleven of the odd primes up to $100$, against
thirteen with $p\equiv3$). \emph{Typical theorems.} Chebotarev's theorem: every element of the Galois group
is a Frobenius element for infinitely many primes. \emph{Class field theory} describes the abelian
extensions of $K$ by congruence conditions on ideals, as the cyclotomic fields do for $\mathbb Q$.
\emph{Complex multiplication} generates the abelian extensions of an imaginary quadratic field $K$ by
special values: the value of the $j$-function at the ring of integers generates the Hilbert class field over
$K$ (for $\mathbb Q(\sqrt{-23})$ the number $j(\frac{1+\sqrt{-23}}2)$ has degree $3$), and adjoining the
values of Weber's function at division points gives the ray class fields. The Hasse
principle says that a Hermitian form over a number field is determined by its behaviour over the real
numbers and over the $p$-adic numbers (see group~D). \emph{Where it goes.} The lattices of the
construction are Hermitian lattices over the integers of the field $E$ of \S III.10. Hermitian spaces are
classified place by place
(rows~15--17); a lattice is determined
place by place only up to its genus.

{\scriptsize\renewcommand{\arraystretch}{0.9}\setlength{\tabcolsep}{4pt}\setlength{\LTpre}{2pt}\setlength{\LTpost}{2pt}
\begin{longtable}{@{}r@{\hspace{6pt}}p{.43\textwidth}p{.285\textwidth}p{.15\textwidth}l@{}}
\midrule
\endhead
7 & \raggedright Artin reciprocity (existence of the Artin map) & \raggedright Artin 1927; Chevalley 1940 & \raggedright S5, G, O, LL & med.\tabularnewline
8 & \raggedright Existence of ray class fields & \raggedright Takagi 1920; Artin 1927 & \raggedright S6, S7, G, O & med.\tabularnewline
9 & \raggedright Characters of ray class groups as Galois characters & \raggedright Hasse's \emph{Bericht} & \raggedright n.r. & med.\tabularnewline
10 & \raggedright Chebotarev: existence of primes of degree one & \raggedright Chebotarev 1926 & \raggedright S5 & med.\tabularnewline
11 & \raggedright The same for cosets of open subgroups & \raggedright Chebotarev 1926 & \raggedright n.r. & med.\tabularnewline
12 & \raggedright Units $\equiv1$ modulo $a$ are $m$th powers & \raggedright Chevalley 1951 & \raggedright S5, G, O, LL & med.\tabularnewline
13 & \raggedright Norms from a quadratic extension: local to global & \raggedright Hasse 1923--24 & \raggedright S6 & high\tabularnewline
14 & \raggedright Class field theory of quadratic extensions (five parts) & \raggedright Hasse; Cassels--Fr\"ohlich & \raggedright T & low\tabularnewline
15 & \raggedright Hasse principle for Hermitian forms & \raggedright Landherr 1936 & \raggedright T & high\tabularnewline
16 & \raggedright Hermitian forms: local classification; CM case & \raggedright Landherr 1936; Jacobowitz 1962 & \raggedright T & high\tabularnewline
17 & \raggedright Hermitian forms over CM fields: the invariants & \raggedright Landherr 1936 & \raggedright S6 & high\tabularnewline
18 & \raggedright Strong approximation for special unitary groups & \raggedright Kneser 1965 & \raggedright S6, O, T & med.\tabularnewline
19 & \raggedright $\ell$-integrality of $L$-values of totally real fields & \raggedright Deligne--Ribet 1980 & \raggedright (G, O, T)$^{*}$ & med.\tabularnewline
20 & \raggedright The same with rationality of the values & \raggedright Klingen 1962; Siegel 1969 & \raggedright (G, O, T)$^{*}$ & med.\tabularnewline
21 & \raggedright Main theorem of complex multiplication & \raggedright Shimura 1971 & \raggedright (G, O, T)$^{*}$ & med.\tabularnewline
22 & \raggedright $K(j(\mathcal O_K))$ is the Hilbert class field & \raggedright Weber; Hasse 1927 & \raggedright (G, O, T)$^{*}$ & high\tabularnewline
23 & \raggedright Artin symbol acts by $j(\mathfrak a)\mapsto j(\mathfrak b^{-1}\mathfrak a)$ & \raggedright Hasse 1927 & \raggedright T & high\tabularnewline
24 & \raggedright The $j(\mathfrak a)$ form one orbit of conjugates & \raggedright Hasse 1927 & \raggedright T & high\tabularnewline
25 & \raggedright Ray class fields from $j$ and Weber's function & \raggedright Weber; Hasse 1927 & \raggedright O & med.\tabularnewline
26 & \raggedright CM curve over $\mathbb Q(j)$ with discriminant $-p^3$ & \raggedright Gross 1980 & \raggedright O & med.\tabularnewline
27 & \raggedright Special values: denominators, algebraicity (five parts) & \raggedright not matched; cf.\ Shimura 1975 & \raggedright G, O, T & low\tabularnewline
\bottomrule
\end{longtable}}

\smallskip\noindent\textbf{Group C. Finite groups, algebra and invariant theory} (rows 28--36).\label{cs:grp:C}
\emph{What it is.} A representation of a finite group $G$ is a homomorphism to invertible complex matrices;
its \emph{character} is the trace. The squares of the degrees of the irreducible characters add up to $|G|$.
Green (1955) found them all for $G=GL_n(\mathbb F_q)$. The \emph{cuspidal} ones correspond to the
monic irreducible polynomials of degree $n$ over $\mathbb F_q$ with nonzero constant term. \emph{Example.}
$GL_2(\mathbb F_3)$ has order $48$. Its irreducible characters are two of degree $1$, two of degree $3$, one
of degree $4$ and three cuspidal ones of degree $2$, one for each irreducible quadratic over $\mathbb F_3$;
and $2+18+16+12=48$. The subject also contains the representations of $GL_n(\mathbb C)$ and $U(n)$, indexed
by highest weights $\lambda_1\ge\dots\ge\lambda_n$, and the rule for restricting them: the representation of
$GL_3$ with weight $(2,1,0)$ has dimension $8$ and restricts to $GL_2$ as the sum of those with weights
$(1,0),(1,1),(2,0),(2,1)$, of dimensions $2+1+3+2$ (row~33). \emph{Where it
goes.} At the two level primes of \S III.8 each form of the construction carries a label, a polynomial over
a finite field; that these labels are distributed like random polynomials is statement (R4) there. Green's
parametrisation (row~29) is passed to the theorems that give (R4).
The branching rule is passed with the theta lifts (T); where it is used we do not know.

{\scriptsize\renewcommand{\arraystretch}{0.9}\setlength{\tabcolsep}{4pt}\setlength{\LTpre}{2pt}\setlength{\LTpost}{2pt}
\begin{longtable}{@{}r@{\hspace{6pt}}p{.43\textwidth}p{.285\textwidth}p{.15\textwidth}l@{}}
\midrule
\endhead
28 & \raggedright Brauer induction in degree zero & \raggedright Brauer 1947; Deligne 1973 & \raggedright S5, G, O, LL & high\tabularnewline
29 & \raggedright Cuspidal characters of $GL_d(\mathbb F_g)$: values & \raggedright Green 1955 & \raggedright R, G, O & med.\tabularnewline
30 & \raggedright Cuspidal characters of $GL_d(\mathbb F_g)$: labels, degree & \raggedright Green 1955 & \raggedright S5, O, T & med.\tabularnewline
31 & \raggedright Generic characters of $GL_N(\mathbb F_g)$ of given degree & \raggedright Green 1955; Deligne--Lusztig 1976 & \raggedright S6, T & med.\tabularnewline
32 & \raggedright Cuspidal representations restricted to the mirabolic & \raggedright S. I. Gelfand 1970 & \raggedright T & med.\tabularnewline
33 & \raggedright Branching from $GL_{n+1}$ to $GL_n$ & \raggedright Weyl 1931 & \raggedright T & med.\tabularnewline
34 & \raggedright $(GL_n,GL_{n'})$ duality on polynomials & \raggedright Howe 1989 & \raggedright T & med.\tabularnewline
35 & \raggedright Shuffle algebra is polynomial on Lyndon words & \raggedright Radford 1979 & \raggedright T & high\tabularnewline
36 & \raggedright Simple induced modules in the lowest alcove & \raggedright Andersen 1980; Jantzen & \raggedright S6 & med.\tabularnewline
\bottomrule
\end{longtable}}

\smallskip\noindent\textbf{Group D. Representations of $p$-adic groups} (rows 37--66).\label{cs:grp:D}
\emph{What it is.} The field $\mathbb Q_p$ of $p$-adic numbers is the completion of $\mathbb Q$ for the
absolute value $|p^ka/b|=p^{-k}$.
The group $GL_n(\mathbb Q_p)$ has representations on infinite-dimensional spaces; those in which every
vector is fixed by a subgroup of matrices $\equiv1$ modulo a power of $p$ are called \emph{smooth}. A modular
form that is an eigenfunction of the Hecke operator $T_p$, with eigenvalue $a_p$, generates such a
representation of $GL_2(\mathbb Q_p)$, and the two roots of $X^2-a_pX+p^{k-1}$ are its \emph{Satake
parameters}. \emph{Typical theorems.} Bernstein and Zelevinsky classified the irreducible smooth
representations of $GL_n$ over a $p$-adic field: each is built from finitely many \emph{cuspidal} ones by
induction, and the building blocks are indexed by `segments'
(rows~43--51). The \emph{local
Langlands correspondence} matches them with $n$-dimensional representations of a group built from the Galois
group of $\mathbb Q_p$ (the Weil--Deligne group); for $n=1$ it is local class field theory.
Row~65 states a weak form of it. \emph{Where it goes.} The norm of
a theta lift is a product over all primes of local factors (input (S3)), and the index of one lattice in
another is a product of local indices (input (S5)); each factor belongs to the representation theory of a
$p$-adic group. Most rows of this group carry the code~G. That code comes from the passing of one whole
record to the theorems that give \eqref{gc:link}; which of its rows those proofs consult was not traced.

{\scriptsize\renewcommand{\arraystretch}{0.9}\setlength{\tabcolsep}{4pt}\setlength{\LTpre}{2pt}\setlength{\LTpost}{2pt}
\begin{longtable}{@{}r@{\hspace{6pt}}p{.43\textwidth}p{.285\textwidth}p{.15\textwidth}l@{}}
\midrule
\endhead
37 & \raggedright Irreducible smooth representations are admissible & \raggedright Jacquet 1975 & \raggedright (G, O)$^{*}$ & high\tabularnewline
38 & \raggedright Induction and Jacquet functor preserve finite length & \raggedright Bernstein--Zelevinsky 1977 & \raggedright (G, O)$^{*}$ & high\tabularnewline
39 & \raggedright The two previous rows; contragredient of an induced & \raggedright Bernstein--Zelevinsky 1976--77 & \raggedright S3, S5, G, O, T, LL & med.\tabularnewline
40 & \raggedright Product of unitary irreducibles is irreducible & \raggedright Bernstein 1984 & \raggedright G, O, T & high\tabularnewline
41 & \raggedright Geometric lemma: filtration of a Jacquet module & \raggedright Bernstein--Zelevinsky 1977 & \raggedright O & med.\tabularnewline
42 & \raggedright The same with the multiplicities & \raggedright Bernstein--Zelevinsky 1977; Zelevinsky 1980 & \raggedright S5, G, O, LL & med.\tabularnewline
43 & \raggedright Representation attached to a segment & \raggedright Zelevinsky 1980 & \raggedright (S5, G, O, T, LL)$^{*}$ & med.\tabularnewline
44 & \raggedright Grothendieck ring is polynomial in the segments & \raggedright Zelevinsky 1980 & \raggedright (S5, G, O, T, LL)$^{*}$ & med.\tabularnewline
45 & \raggedright The same with existence included & \raggedright Zelevinsky 1980 & \raggedright S5, G, O, T, LL & med.\tabularnewline
46 & \raggedright Essentially square-integrable $=$ segment & \raggedright Bernstein; Zelevinsky 1980 & \raggedright G, O & med.\tabularnewline
47 & \raggedright Square-integrable implies segment & \raggedright Zelevinsky 1980 & \raggedright G, O & med.\tabularnewline
48 & \raggedright Generic representations: unlinked segments & \raggedright Zelevinsky 1980 & \raggedright G, O & med.\tabularnewline
49 & \raggedright Every irreducible is a Langlands quotient & \raggedright Rodier 1982; Zelevinsky 1980 & \raggedright S5, G, O, LL & med.\tabularnewline
50 & \raggedright Uniqueness of the multiset of segments & \raggedright Zelevinsky 1980 & \raggedright S5, G, O, LL & med.\tabularnewline
51 & \raggedright Standard module has a unique irreducible quotient & \raggedright Silberger 1978 & \raggedright S5, G, O, LL & med.\tabularnewline
52 & \raggedright Tempered $=$ induced from square-integrable & \raggedright Jacquet 1977; Zelevinsky 1980 & \raggedright G, O & med.\tabularnewline
53 & \raggedright Irreducibility of a product, supports far apart & \raggedright after Zelevinsky 1980 & \raggedright G, O & med.\tabularnewline
54 & \raggedright Irreducibility of a product, supports unrelated & \raggedright after Zelevinsky 1980 & \raggedright n.r. & med.\tabularnewline
55 & \raggedright Generic unitary dual of $GL_N$ & \raggedright Tadi\'c 1986; Bernstein 1984 & \raggedright n.r. & med.\tabularnewline
56 & \raggedright Depth zero passes to the cuspidal support & \raggedright Moy--Prasad 1996 & \raggedright G & med.\tabularnewline
57 & \raggedright Depth zero passes from the cuspidal support & \raggedright Moy--Prasad 1996 & \raggedright S5, G, O, LL & med.\tabularnewline
58 & \raggedright Fixed vectors and the Jacquet module (level one) & \raggedright Moy--Prasad 1996 & \raggedright G, O & med.\tabularnewline
59 & \raggedright The injectivity half of the previous row & \raggedright Moy--Prasad 1996 & \raggedright n.r. & med.\tabularnewline
60 & \raggedright Iwahori-fixed vectors and the Jacquet module & \raggedright Borel 1976; Casselman 1980 & \raggedright T & med.\tabularnewline
61 & \raggedright Local converse theorem for cuspidals & \raggedright Henniart 1993 & \raggedright S5, G, O, LL & med.\tabularnewline
62 & \raggedright Unramified doubling zeta integral & \raggedright Piatetski-Shapiro--Rallis 1987 & \raggedright S3, G, O, T & med.\tabularnewline
63 & \raggedright Stone--von Neumann theorem over a $p$-adic field & \raggedright M\oe glin et al.\ 1987 & \raggedright T & high\tabularnewline
64 & \raggedright Weil representation of $U(\nu,\nu)$ without a cover & \raggedright Kudla 1994 & \raggedright T & med.\tabularnewline
65 & \raggedright Weak form of the local Langlands correspondence & \raggedright Harris--Taylor 2001; Henniart 2000 & \raggedright S5, G, O, LL & low\tabularnewline
66 & \raggedright The same statement & \raggedright as for 65 & \raggedright S5, G, O, LL & low\tabularnewline
\bottomrule
\end{longtable}}

\smallskip\noindent\textbf{Group E. Representations of real groups} (rows 67--84).\label{cs:grp:E}
\emph{What it is.} The group $U(p,q)$ consists of the complex matrices that preserve the Hermitian form
$|z_1|^2+\dots+|z_p|^2-|z_{p+1}|^2-\dots-|z_{p+q}|^2$. For $pq>0$ its irreducible representations are
infinite-dimensional, and one studies them through their restriction to the compact subgroup
$K=U(p)\times U(q)$. The \emph{discrete series} are the
irreducible representations that occur as summands of $L^2$ of the group. A cusp form of weight $k$ on the
upper half-plane is a vector in a discrete series representation of $SL_2(\mathbb R)$. Harish-Chandra
classified them by a strictly decreasing sequence $\lambda$. The \emph{formal degree} replaces
the dimension, and it is given by Weyl's formula: $\prod_{a<b}|\lambda_a-\lambda_b|/\prod_{a<b}(b-a)$ is $1$,
$3$, $6$ for $\lambda=(1,0,-1)$, $(2,0,-1)$, $(3,0,-1)$, the dimensions of the trivial, the standard and the
symmetric square representation of $U(3)$ (row~82). \emph{Where it
goes.} The forms of the construction have prescribed weights at the real places: they are vectors of a
prescribed discrete series. The dimension and hook-product terms in the main term of \S I.7 have the shape
of formal degrees; \S III.4 obtains the hook product as an exact Gaussian integral, and the identification
with formal degrees is not made anywhere in this paper.

{\scriptsize\renewcommand{\arraystretch}{0.9}\setlength{\tabcolsep}{4pt}\setlength{\LTpre}{2pt}\setlength{\LTpost}{2pt}
\begin{longtable}{@{}r@{\hspace{6pt}}p{.43\textwidth}p{.285\textwidth}p{.15\textwidth}l@{}}
\midrule
\endhead
67 & \raggedright Schur orthogonality; formal degree & \raggedright Godement 1947 & \raggedright O, T & high\tabularnewline
68 & \raggedright Smooth functions are sums of convolutions & \raggedright Dixmier--Malliavin 1978 & \raggedright G, O, T & med.\tabularnewline
69 & \raggedright $K$-finite vectors of unitary representations of $U(p,q)$ & \raggedright Harish-Chandra 1953 & \raggedright S6, S7, O, T & med.\tabularnewline
70 & \raggedright Admissibility, analytic vectors, irreducibility & \raggedright Harish-Chandra 1953 & \raggedright T & med.\tabularnewline
71 & \raggedright Closures of $(\mathfrak g,K)$-submodules; equivalence & \raggedright Harish-Chandra 1953 & \raggedright T & med.\tabularnewline
72 & \raggedright Matrix coefficients are real analytic & \raggedright Harish-Chandra 1953 & \raggedright G, O & med.\tabularnewline
73 & \raggedright Automatic continuity of $(\mathfrak g,K)$-maps & \raggedright Casselman 1989; Wallach 1992 & \raggedright S5, S6, S7, O, T & med.\tabularnewline
74 & \raggedright Existence of discrete series of $U(p,q)$ & \raggedright Harish-Chandra 1966 & \raggedright T & med.\tabularnewline
75 & \raggedright The same statement & \raggedright as for 74 & \raggedright G, O, T & med.\tabularnewline
76 & \raggedright Classification of discrete series of $U(p,q)$ & \raggedright Harish-Chandra 1966 & \raggedright G, O & med.\tabularnewline
77 & \raggedright A family covering all discrete series & \raggedright Harish-Chandra 1966 & \raggedright S5, S6, S7, O, T & med.\tabularnewline
78 & \raggedright Existence for every admissible parameter & \raggedright Harish-Chandra 1966 & \raggedright O, T & med.\tabularnewline
79 & \raggedright Discrete series determined by lowest $K$-type & \raggedright Hecht--Schmid 1975 & \raggedright T & med.\tabularnewline
80 & \raggedright Lowest $K$-type occurs; uniqueness & \raggedright Hecht--Schmid 1975 & \raggedright T & med.\tabularnewline
81 & \raggedright Holomorphic discrete series are lowest weight & \raggedright Harish-Chandra 1955--56 & \raggedright T & med.\tabularnewline
82 & \raggedright Formal degree $=$ Weyl's dimension formula & \raggedright Harish-Chandra 1966 & \raggedright S3, O, T & med.\tabularnewline
83 & \raggedright Metaplectic representation of $Sp_{2n}(\mathbb R)$ & \raggedright Shale 1962; Weil 1964 & \raggedright T & med.\tabularnewline
84 & \raggedright The same: Gaussian and its matrix coefficient & \raggedright Shale 1962; Weil 1964 & \raggedright T & med.\tabularnewline
\bottomrule
\end{longtable}}

\smallskip\noindent\textbf{Group F. Automorphic forms} (rows 85--106).\label{cs:grp:F}
\emph{What it is.} An automorphic form is a function on a group that is invariant under an arithmetic
subgroup, as a modular form is under $SL_2(\mathbb Z)$. For the definite unitary groups of \S I.7 it is a
function on the finite set of classes of lattices in a genus. For the indefinite ones the space need not be compact, and
\emph{cusp forms} are those that decay at infinity (a cusp form on the upper half-plane is
$O(e^{-2\pi y})$ as $y\to\infty$). \emph{Typical theorems.} The space of
forms of a given weight and level has finite dimension (row~86). A cusp
form on $GL_n$ is determined by its Hecke eigenvalues at almost all primes
(rows~95, 97). \emph{Base change} (row~102, in a weak form)
attaches to a form on a unitary group forms on general linear groups over $E$ with matching Hecke
eigenvalues. \emph{Where it goes.} Galois
representations are attached to forms on $GL_n$ over CM fields. One expects that the forms of the
construction are first moved there by base change, on the way to the Galois cohomology of \S III.9.
Multiplicity one
(row~97) would make the result unambiguous, but that row is among
the ten that are never read.

{\scriptsize\renewcommand{\arraystretch}{0.9}\setlength{\tabcolsep}{4pt}\setlength{\LTpre}{2pt}\setlength{\LTpost}{2pt}
\begin{longtable}{@{}r@{\hspace{6pt}}p{.43\textwidth}p{.285\textwidth}p{.15\textwidth}l@{}}
\midrule
\endhead
85 & \raggedright Reduction theory: Siegel sets for unitary groups & \raggedright Borel 1963; Borel--Harish-Chandra 1962 & \raggedright T & med.\tabularnewline
86 & \raggedright Finiteness of spaces of automorphic forms & \raggedright Harish-Chandra 1968 & \raggedright G, O, T & med.\tabularnewline
87 & \raggedright The same for $GL_a$ & \raggedright Harish-Chandra 1968 & \raggedright G, O & med.\tabularnewline
88 & \raggedright Cusp forms are rapidly decreasing & \raggedright Gelfand--Piatetski-Shapiro 1963 & \raggedright S3, S5, S6, S7, G, O, T & med.\tabularnewline
89 & \raggedright Cusp forms and their derivatives are bounded & \raggedright Gelfand--Piatetski-Shapiro 1963 & \raggedright T & med.\tabularnewline
90 & \raggedright Decay of cusp forms, uniform in the form & \raggedright not matched; method of 88 & \raggedright T & low\tabularnewline
91 & \raggedright The same with one norm of the derivatives & \raggedright not matched; method of 88 & \raggedright T & low\tabularnewline
92 & \raggedright Convolution on cusp forms is of trace class & \raggedright Gelfand--Piatetski-Shapiro 1963 & \raggedright G, O, T & med.\tabularnewline
93 & \raggedright Cusp forms on $GL_a$ are bounded & \raggedright Gelfand--Piatetski-Shapiro 1963 & \raggedright (S5, G, O)$^{*}$ & med.\tabularnewline
94 & \raggedright Cuspidal spectrum of $GL_a$ is discrete & \raggedright Gelfand--Piatetski-Shapiro 1963 & \raggedright G, O & med.\tabularnewline
95 & \raggedright Strong multiplicity one for isobaric sums & \raggedright Jacquet--Shalika 1981 & \raggedright (S5, G, O)$^{*}$ & med.\tabularnewline
96 & \raggedright The same; bound for Satake parameters & \raggedright Jacquet--Shalika 1981 & \raggedright G, O & med.\tabularnewline
97 & \raggedright Multiplicity one for $GL_a$ & \raggedright Shalika 1974 & \raggedright n.r. & med.\tabularnewline
98 & \raggedright Contragredient realised by $g\mapsto{}^tg^{-1}$ & \raggedright Gelfand--Kazhdan 1975 & \raggedright S5, G, O & med.\tabularnewline
99 & \raggedright Tensor product decomposition & \raggedright Flath 1979 & \raggedright (G, O, T)$^{*}$ & med.\tabularnewline
100 & \raggedright Unitary irreducibles of $p$-adic groups are admissible & \raggedright Bernstein 1974 & \raggedright (G, O, T)$^{*}$ & med.\tabularnewline
101 & \raggedright Factorisation with spherical vectors & \raggedright Flath 1979 & \raggedright G, O, T & med.\tabularnewline
102 & \raggedright Weak base change from $U(n)$ to $GL_n$ & \raggedright Labesse 2011 & \raggedright S5, G, O & low\tabularnewline
103 & \raggedright The same for the full induced representation & \raggedright Labesse 2011 & \raggedright S5 & low\tabularnewline
104 & \raggedright Regular algebraic implies cohomological & \raggedright Clozel 1990 & \raggedright S5, G, O & med.\tabularnewline
105 & \raggedright Purity lemma & \raggedright Clozel 1990 & \raggedright G, O & med.\tabularnewline
106 & \raggedright The two previous rows, for CM fields & \raggedright Clozel 1990 & \raggedright n.r. & med.\tabularnewline
\bottomrule
\end{longtable}}

\smallskip\noindent\textbf{Group G. Theta correspondence, Eisenstein series, the Siegel--Weil formula} (rows 107--122).\label{cs:grp:G}
\emph{What it is.} The subject of \S I.7 in its general form. A theta series in two sets of variables is the
kernel of an integral operator between the automorphic forms on two groups, the theta lift. The \emph{Weil
representation} is the action of the two groups on the functions from which the theta series are formed
(rows~107, 108).
\emph{Typical theorems.} The \emph{Siegel--Weil formula} says that the average of the theta series over the
classes in a genus is an Eisenstein series. Its ancestor is Siegel's theorem that the number of
representations of $n$ by the forms of a genus, averaged with weights $1/|{\rm Aut}|$, is a product of local
densities; Jacobi's formula $r_4(n)=8\sum_{4\nmid d\mid n}d$ is the case of one class. The \emph{inner product
formula} of Rallis (row~118) expresses the norm of a theta
lift as a value of an $L$-function times local factors; in the literature it is proved by the
\emph{doubling method} (row~117). \emph{Where it goes.}
Input (S3) of \S III.7 is a lower bound for the eigenvalues of a lift composed with its adjoint; the
four rows with code S3 in this group are passed to the theorem that gives it
(row~115 only through another row). The Siegel--Weil formula
in the form of row~114 is passed to the theorems that give
(R1) of \S III.8; compare (R2) there.

{\scriptsize\renewcommand{\arraystretch}{0.9}\setlength{\tabcolsep}{4pt}\setlength{\LTpre}{2pt}\setlength{\LTpost}{2pt}
\begin{longtable}{@{}r@{\hspace{6pt}}p{.43\textwidth}p{.285\textwidth}p{.15\textwidth}l@{}}
\midrule
\endhead
107 & \raggedright Weil representation of a unitary pair exists & \raggedright Weil 1964; Kudla 1994 & \raggedright T & med.\tabularnewline
108 & \raggedright The same for the spaces of the construction & \raggedright Weil 1964; Kudla 1994 & \raggedright G, O, T & med.\tabularnewline
109 & \raggedright Weil's criterion of convergence & \raggedright Weil 1965 & \raggedright (T)$^{*}$ & med.\tabularnewline
110 & \raggedright Siegel--Weil formula, range of convergence & \raggedright Weil 1965 & \raggedright S3, G, O, T & med.\tabularnewline
111 & \raggedright Siegel--Weil beyond convergence (Gaussian data) & \raggedright Ichino 2004, 2007 & \raggedright S3, O, T & med.\tabularnewline
112 & \raggedright The same in the Schr\"odinger model & \raggedright Ichino 2004, 2007 & \raggedright O & med.\tabularnewline
113 & \raggedright Siegel--Weil, one space definite & \raggedright Ichino 2004, 2007; Yamana 2011 & \raggedright (T)$^{*}$ & med.\tabularnewline
114 & \raggedright Finite volume; Siegel's formula for Hermitian forms & \raggedright Weil 1965; Braun 1941 & \raggedright R, O & med.\tabularnewline
115 & \raggedright Continuation of Siegel Eisenstein series & \raggedright Langlands 1976 & \raggedright S3, G, O, T & med.\tabularnewline
116 & \raggedright The same in general, with moderate growth & \raggedright Bernstein--Lapid 2024 (a guess) & \raggedright (T)$^{*}$ & med.\tabularnewline
117 & \raggedright Doubling method: the basic identity & \raggedright Piatetski-Shapiro--Rallis 1987 & \raggedright S3, G, O, T & med.\tabularnewline
118 & \raggedright Rallis inner product formula (convergent range) & \raggedright Rallis 1987; Li 1992 & \raggedright T & med.\tabularnewline
119 & \raggedright Howe duality at the real places, trivial representation & \raggedright Howe 1989 & \raggedright T & low\tabularnewline
120 & \raggedright Theta functions under rational points & \raggedright not matched; cf.\ Weil 1964 & \raggedright T & low\tabularnewline
121 & \raggedright Two theta integrals are proportional (three parts) & \raggedright not matched; cf.\ Kudla--Rallis 1994 & \raggedright T & low\tabularnewline
122 & \raggedright The proposition `True' & \raggedright none & \raggedright O & high\tabularnewline
\bottomrule
\end{longtable}}

\smallskip\noindent\textbf{Group H. Galois cohomology} (rows 123--134).\label{cs:grp:H}
\emph{What it is.} Let a group $G$ act on an abelian group $M$. A \emph{crossed homomorphism} is a map
$c:G\to M$ with $c(gh)=c(g)+g\,c(h)$; those of the form $g\mapsto gm-m$ are called principal, and $H^1(G,M)$
is the quotient. $H^0$ is the group of fixed elements, and there are higher groups $H^2,\dots$ \emph{Example.}
For the Galois group of $\mathbb Q_p$, $p$ odd, and $M=\mu_p$, the $p$th roots of unity: $H^0$ is trivial,
$H^1$ is $\mathbb Q_p^\times$ modulo $p$th powers, of order $p^2$, and $H^2$ has order $p$. This agrees with
the formula of row~125: $p^2=1\cdot p\cdot p$. \emph{Typical theorems.}
Tate's local duality pairs $H^i$ of $M$ with $H^{2-i}$ of a dual module. Globally one imposes local
conditions and obtains \emph{Selmer groups}; the ideal class group is the simplest one. The formula of
row~133 compares the order of a Selmer group with that of its dual.
\emph{Where it goes.} The groups whose orders \S III.9 bounds are of this kind. All twelve rows are passed,
in one record, to the theorems that give the inequality \eqref{gc:link} and to the construction.

{\scriptsize\renewcommand{\arraystretch}{0.9}\setlength{\tabcolsep}{4pt}\setlength{\LTpre}{2pt}\setlength{\LTpost}{2pt}
\begin{longtable}{@{}r@{\hspace{6pt}}p{.43\textwidth}p{.285\textwidth}p{.15\textwidth}l@{}}
\midrule
\endhead
123 & \raggedright Local Galois cohomology is finite & \raggedright Tate 1962 & \raggedright G, O & med.\tabularnewline
124 & \raggedright Local Tate duality & \raggedright Tate 1962 & \raggedright G, O & med.\tabularnewline
125 & \raggedright Local Euler characteristic formula & \raggedright Tate 1962 & \raggedright G, O & med.\tabularnewline
126 & \raggedright Finiteness at a place of a number field & \raggedright Tate 1962 & \raggedright G, O & med.\tabularnewline
127 & \raggedright Duality at a place of a number field & \raggedright Tate 1962 & \raggedright G, O & med.\tabularnewline
128 & \raggedright Euler characteristic at a place of a number field & \raggedright Tate 1962 & \raggedright G, O & med.\tabularnewline
129 & \raggedright Unramified classes are exact annihilators & \raggedright Tate 1962; NSW 7.2.15 & \raggedright G, O & med.\tabularnewline
130 & \raggedright Finiteness of $H^i(G_S,M)$ & \raggedright Tate 1962 & \raggedright G, O & high\tabularnewline
131 & \raggedright Global Euler characteristic formula & \raggedright Tate 1965 & \raggedright G, O & med.\tabularnewline
132 & \raggedright Poitou--Tate exact sequence & \raggedright Tate 1962; Poitou 1967 & \raggedright G, O & med.\tabularnewline
133 & \raggedright Formula for the quotient of two Selmer groups & \raggedright Wiles 1995 & \raggedright G, O & med.\tabularnewline
134 & \raggedright Local conditions $H^1_f$: dimension and duality & \raggedright Bloch--Kato 1990 & \raggedright G, O & med.\tabularnewline
\bottomrule
\end{longtable}}

\smallskip\noindent\textbf{Group I. Galois representations and $p$-adic Hodge theory} (rows 135--151).\label{cs:grp:I}
\emph{What it is.} To a Dirichlet character one attaches a character of the Galois group of a cyclotomic
field. In the same way, to a Hecke eigenform one attaches a two-dimensional representation of the Galois
group of $\overline{\mathbb Q}$ over $\mathbb Q$ with coefficients in the $\ell$-adic numbers, in which the
Frobenius element at $p$ has trace $a_p$. Rows~135--140
are the case of Hecke characters. Rows~145--148 concern forms
on $GL_a$ over a CM field; they state weak forms of the theorem. \emph{$p$-adic Hodge theory} studies such
a representation at the prime $\ell$ itself, where it is not determined by a Frobenius element. Fontaine's
rings (row~142) single out the \emph{crystalline} representations, the analogue
at $\ell$ of `unramified'; the theorem of Bhatt and Scholze (row~143)
describes the lattices in them. \emph{Where it goes.} One expects that at a prime $\ell$ dividing the index of two lattices integral
information at $\ell$ is needed, and that statements on lattices in crystalline representations provide it.
Ten of the rows are passed to the local index bounds (S5) and nine to the inequality for Galois
cohomology (G); three are never read.

{\scriptsize\renewcommand{\arraystretch}{0.9}\setlength{\tabcolsep}{4pt}\setlength{\LTpre}{2pt}\setlength{\LTpost}{2pt}
\begin{longtable}{@{}r@{\hspace{6pt}}p{.43\textwidth}p{.285\textwidth}p{.15\textwidth}l@{}}
\midrule
\endhead
135 & \raggedright Galois character of an algebraic Hecke character & \raggedright Weil 1956; Serre 1968 & \raggedright S5 & med.\tabularnewline
136 & \raggedright It is de Rham with weights $-n_\tau$ & \raggedright Serre 1968; Fontaine & \raggedright (S5, G, O)$^{*}$ & med.\tabularnewline
137 & \raggedright It is crystalline where $\chi$ is unramified & \raggedright Serre 1968; Fontaine & \raggedright (S5, G, O)$^{*}$ & med.\tabularnewline
138 & \raggedright The two previous rows; potential crystallinity & \raggedright Serre 1968; Fontaine & \raggedright S5, G, O & med.\tabularnewline
139 & \raggedright Its Weil--Deligne representation above $\ell$ & \raggedright local class field theory; Fontaine & \raggedright S5, G, O & med.\tabularnewline
140 & \raggedright The same through the correspondence of row 65 & \raggedright local class field theory; Fontaine & \raggedright n.r. & med.\tabularnewline
141 & \raggedright De Rham implies potentially semistable & \raggedright Berger 2002 & \raggedright S5, G, O & med.\tabularnewline
142 & \raggedright Seven properties of $B_{\rm cris}$ & \raggedright Fontaine 1982, 1994 & \raggedright G & low\tabularnewline
143 & \raggedright Lattices in crystalline representations & \raggedright Bhatt--Scholze 2023 & \raggedright S5, G & med.\tabularnewline
144 & \raggedright The same over a finite extension & \raggedright Bhatt--Scholze 2023 & \raggedright S5 & med.\tabularnewline
145 & \raggedright Weak form: Galois representations for $GL_a$ & \raggedright Shin 2011; Chenevier--Harris 2013 & \raggedright S5, G, O & low\tabularnewline
146 & \raggedright The same, polarisation in a second form & \raggedright as for 145 & \raggedright S5 & low\tabularnewline
147 & \raggedright Weak form of compatibility at $\ell$ & \raggedright Caraiani 2014 & \raggedright n.r. & low\tabularnewline
148 & \raggedright The same statement & \raggedright as for 147 & \raggedright n.r. & low\tabularnewline
149 & \raggedright A de Rham representation has $a$ Hodge--Tate weights & \raggedright Fontaine 1994 & \raggedright G, O & med.\tabularnewline
150 & \raggedright Kernel of the Hodge--Tate map and duality & \raggedright Tate 1967; Fargues 2008 & \raggedright S6 & med.\tabularnewline
151 & \raggedright Dieudonn\'e functor for $p$-divisible groups over $\mathcal O_C$ & \raggedright Scholze--Weinstein 2013 & \raggedright S6 & med.\tabularnewline
\bottomrule
\end{longtable}}

\smallskip\noindent\textbf{Group J. Abelian varieties and their moduli} (rows 152--165).\label{cs:grp:J}
\emph{What it is.} An elliptic curve over $\mathbb C$ is a torus $\mathbb C/\Lambda$. An abelian variety of
dimension $g$ is a torus $\mathbb C^g/\Lambda$ that can be embedded in projective space; the condition is
the existence of a \emph{polarisation}, a positive Hermitian form whose imaginary part is integral on
$\Lambda$. It has $n^{2g}$ points of order dividing $n$ (row~152): on the curve
$y^2=x^3-x$ the points of order dividing $2$ are $(0,0)$, $(\pm1,0)$ and the point at infinity. Abelian
varieties with a fixed kind of polarisation are themselves the points of a variety, a \emph{moduli space};
for elliptic curves it is the $j$-line. The unitary groups $U(p,q)$ of the construction
have such spaces, and modular forms are sections of powers of a line bundle $\omega$ on them, the
\emph{Hodge bundle}. \emph{Typical theorems.} These spaces are not compact; a \emph{toroidal
compactification} adds a boundary over which the abelian varieties degenerate in a controlled way
(rows~163--165). \emph{Where it goes.} Every row of this group is passed to one place: the theorem that gives the
bound for the covolume of the lattice of integral forms, input (S6) of \S III.7. That lattice is the set of
integral sections of a power of $\omega$. One expects that counting them needs a compact space, a bundle
defined over the integers, and control of the metric near the boundary; we did not read the proof.

{\scriptsize\renewcommand{\arraystretch}{0.9}\setlength{\tabcolsep}{4pt}\setlength{\LTpre}{2pt}\setlength{\LTpost}{2pt}
\begin{longtable}{@{}r@{\hspace{6pt}}p{.43\textwidth}p{.285\textwidth}p{.15\textwidth}l@{}}
\midrule
\endhead
152 & \raggedright $n^{2g}$ points of order dividing $n$ & \raggedright Mumford 1970 & \raggedright S6 & high\tabularnewline
153 & \raggedright Principal polarisation: one section & \raggedright Mumford 1970 & \raggedright S6 & med.\tabularnewline
154 & \raggedright Dual abelian scheme & \raggedright Grothendieck; Mumford 1965 & \raggedright (S6)$^{*}$ & med.\tabularnewline
155 & \raggedright Zarhin's trick over a valuation ring & \raggedright Zarhin 1974 & \raggedright S6 & med.\tabularnewline
156 & \raggedright De Rham cohomology of abelian schemes (five parts) & \raggedright Berthelot--Breen--Messing 1982 & \raggedright S6 & low\tabularnewline
157 & \raggedright Lifting abelian schemes with the Hodge filtration & \raggedright Grothendieck 1970; Messing 1972 & \raggedright (S6)$^{*}$ & med.\tabularnewline
158 & \raggedright The resulting equivalence of categories & \raggedright Grothendieck 1970; Messing 1972 & \raggedright S6 & med.\tabularnewline
159 & \raggedright The same when $p$ is nilpotent & \raggedright Grothendieck 1970; Messing 1972 & \raggedright S6 & med.\tabularnewline
160 & \raggedright Semistable reduction of abelian varieties & \raggedright Grothendieck 1972 (SGA 7) & \raggedright S6 & med.\tabularnewline
161 & \raggedright Extension of homomorphisms of semi-abelian schemes & \raggedright Faltings--Chai 1990 & \raggedright S6 & high\tabularnewline
162 & \raggedright Seminegativity of the Lie algebra bundle & \raggedright Griffiths 1970; Fujita 1978 & \raggedright S6 & med.\tabularnewline
163 & \raggedright Toroidal compactification, unitary case & \raggedright Ash et al.\ 1975; Lan 2013 & \raggedright S6 & low\tabularnewline
164 & \raggedright The same for a polarisation function; ampleness & \raggedright as for 163 & \raggedright S6 & low\tabularnewline
165 & \raggedright Toroidal compactification, Siegel case, integral & \raggedright Faltings--Chai 1990 & \raggedright S6 & med.\tabularnewline
\bottomrule
\end{longtable}}

\smallskip\noindent\textbf{Group K. Algebraic and arithmetic geometry} (rows 166--185).\label{cs:grp:K}
\emph{What it is.} On a projective variety of dimension $D$ the sections of the $m$th power of a line bundle
$L$ form a space whose dimension grows like $L^D\,m^D/D!$, where $L^D$ is an intersection number. The
simplest case is the count of monomials: there are $\binom{m+2}2$ monomials of degree $m$ in three
variables, $5151$ for $m=100$, against $m^2/2=5000$. \emph{Vanishing theorems} (Kawamata--Viehweg, Nadel)
say when this count is exact. \emph{Arakelov
geometry} adds the arithmetic: the sections with integer coefficients form a lattice, a metric on $L$ makes
it a lattice in a normed space, and the \emph{arithmetic Hilbert--Samuel theorem}
(row~183) says that the number of lattice points of norm at most $1$ in
the $m$th power grows like $\exp\bigl(a\,m^{D+1}/(D+1)!\bigr)$, with $a$ an arithmetic intersection number.
By Minkowski's theorem, a count of lattice points in a ball is a statement about a covolume. \emph{Where it
goes.} The theorem of the development that gives the covolume bound has these statements as hypotheses. We
did not read its proof. One expects that the covolume is turned into an arithmetic intersection number of
the Hodge bundle, and that the other rows (resolution of singularities, positivity and vanishing theorems,
an $L^2$ extension theorem) serve this.
All twenty rows are passed to the covolume bound and to nothing else. In \S III.7, $E_3=m\,C_3\,B_0\log\log x$ by definition; what receives the rows is the
covolume bound, the first hypothesis of the theorem quoted under (S5), and it is derived in the
development from a statement that takes all the rows of this group and of the previous one as hypotheses.

{\scriptsize\renewcommand{\arraystretch}{0.9}\setlength{\tabcolsep}{4pt}\setlength{\LTpre}{2pt}\setlength{\LTpost}{2pt}
\begin{longtable}{@{}r@{\hspace{6pt}}p{.43\textwidth}p{.285\textwidth}p{.15\textwidth}l@{}}
\midrule
\endhead
166 & \raggedright Resolution of singularities & \raggedright Hironaka 1964 & \raggedright S6 & med.\tabularnewline
167 & \raggedright Principalisation of an ideal & \raggedright Hironaka 1964 & \raggedright S6 & med.\tabularnewline
168 & \raggedright The same, compatible with a divisor & \raggedright Hironaka 1964; Koll\'ar 2007 & \raggedright S6 & med.\tabularnewline
169 & \raggedright Existence of the Chow form & \raggedright Chow--van der Waerden 1937 & \raggedright S6 & med.\tabularnewline
170 & \raggedright Algebraisation of formal morphisms and schemes & \raggedright Grothendieck 1961 (EGA III) & \raggedright S6 & high\tabularnewline
171 & \raggedright Holomorphic sections are algebraic & \raggedright Serre 1956 & \raggedright S6 & med.\tabularnewline
172 & \raggedright Nef bundles: positivity, bigness, growth of sections & \raggedright Kleiman 1966 & \raggedright S6 & high\tabularnewline
173 & \raggedright The growth of sections, in coordinates over $\mathbb Q$ & \raggedright not matched; cf.\ Lazarsfeld & \raggedright S6 & low\tabularnewline
174 & \raggedright Kawamata--Viehweg vanishing & \raggedright Kawamata 1982; Viehweg 1982 & \raggedright S6 & high\tabularnewline
175 & \raggedright Nadel vanishing; a divisor not klt at one point & \raggedright Koll\'ar 1997 & \raggedright S6 & med.\tabularnewline
176 & \raggedright Sections separating two points & \raggedright Koll\'ar 1997 & \raggedright S6 & med.\tabularnewline
177 & \raggedright $L^2$ extension from points, constant $34/\delta$ & \raggedright Demailly 2016 & \raggedright S6 & low\tabularnewline
178 & \raggedright Chern numbers from good metrics & \raggedright Mumford 1977 & \raggedright S6 & low\tabularnewline
179 & \raggedright Poincar\'e--Lelong formula, in coordinates & \raggedright Lelong 1957; Bost et al.\ 1994 & \raggedright S6 & low\tabularnewline
180 & \raggedright Arithmetic intersection numbers (seven parts) & \raggedright Gillet--Soul\'e 1990 & \raggedright S6 & med.\tabularnewline
181 & \raggedright Heights with respect to Hermitian line bundles & \raggedright Bost--Gillet--Soul\'e 1994 & \raggedright S6 & med.\tabularnewline
182 & \raggedright Birational invariance; degree as an integral & \raggedright Gillet--Soul\'e 1990 & \raggedright S6 & med.\tabularnewline
183 & \raggedright Arithmetic Hilbert--Samuel theorem & \raggedright Gillet--Soul\'e 1992; Moriwaki 2009 & \raggedright S6 & med.\tabularnewline
184 & \raggedright Arithmetic volume of a nef bundle & \raggedright Moriwaki 2009 & \raggedright S6 & med.\tabularnewline
185 & \raggedright Lower bound for the arithmetic volume & \raggedright Moriwaki 2009; Yuan 2008 & \raggedright S6 & med.\tabularnewline
\bottomrule
\end{longtable}}

\smallskip\noindent\textbf{Group L. Records and the final step} (rows 186--195).\label{cs:grp:L}
These ten rows are not theorems of mathematics in their own right. Nine are records that pack other rows
into one statement, and the tenth is the implication that takes them as arguments. No destination is shown for a
record: a record goes where its entries go.

{\scriptsize\renewcommand{\arraystretch}{0.9}\setlength{\tabcolsep}{4pt}\setlength{\LTpre}{2pt}\setlength{\LTpost}{2pt}
\begin{longtable}{@{}r@{\hspace{6pt}}p{.43\textwidth}p{.285\textwidth}p{.15\textwidth}l@{}}
\midrule
\endhead
186 & \raggedright Record: $p$-adic and finite groups & \raggedright --- & \raggedright --- & ---\tabularnewline
187 & \raggedright Record: automorphic forms and Galois representations & \raggedright --- & \raggedright --- & ---\tabularnewline
188 & \raggedright Record: Galois cohomology & \raggedright --- & \raggedright --- & ---\tabularnewline
189 & \raggedright Record: the three previous records and two rows & \raggedright --- & \raggedright --- & ---\tabularnewline
190 & \raggedright Record: theta correspondence, real groups, $L$-functions & \raggedright --- & \raggedright --- & ---\tabularnewline
191 & \raggedright Record: $p$-adic Hodge theory, local cohomology & \raggedright --- & \raggedright --- & ---\tabularnewline
192 & \raggedright The first large record & \raggedright --- & \raggedright --- & ---\tabularnewline
193 & \raggedright The second large record: theta lifts & \raggedright --- & \raggedright --- & ---\tabularnewline
194 & \raggedright Record: abelian varieties and geometry & \raggedright --- & \raggedright --- & ---\tabularnewline
195 & \raggedright The implication that takes the records & \raggedright --- & \raggedright --- & ---\tabularnewline
\bottomrule
\end{longtable}}


\subsection*{3.\quad Summary}\label{cs:sum}

The table counts the rows of each group, the rows that carry each code, and the rows of each grade. A row
with several codes is counted under each, so the columns of codes do not add up to the number of rows.

{\footnotesize\setlength{\tabcolsep}{3pt}
\begin{center}
\begin{tabular}{@{}lr|rrrrrrrrrr|rrr@{}}
\toprule
Group & rows & S3 & S5 & S6 & S7 & R & G & O & T & LL & n.r. & high & med. & low\\
\midrule
A. Analytic number theory & 6 & $\cdot$ & $\cdot$ & $\cdot$ & $\cdot$ & $\cdot$ & 3 & 3 & 6 & $\cdot$ & $\cdot$ & 1 & 5 & $\cdot$\\
B. Algebraic number theory, CM & 21 & $\cdot$ & 3 & 4 & 1 & $\cdot$ & 8 & 11 & 11 & 2 & 2 & 7 & 12 & 2\\
C. Finite groups, algebra & 9 & $\cdot$ & 2 & 2 & $\cdot$ & 1 & 2 & 3 & 6 & 1 & $\cdot$ & 2 & 7 & $\cdot$\\
D. $p$-adic groups & 30 & 2 & 12 & $\cdot$ & $\cdot$ & $\cdot$ & 23 & 23 & 9 & 12 & 3 & 4 & 24 & 2\\
E. Real groups & 18 & 1 & 2 & 3 & 3 & $\cdot$ & 4 & 10 & 16 & $\cdot$ & $\cdot$ & 1 & 17 & $\cdot$\\
F. Automorphic forms & 22 & 1 & 7 & 1 & 1 & $\cdot$ & 15 & 15 & 10 & $\cdot$ & 2 & $\cdot$ & 18 & 4\\
G. Theta correspondence & 16 & 4 & $\cdot$ & $\cdot$ & $\cdot$ & 1 & 4 & 8 & 13 & $\cdot$ & $\cdot$ & 1 & 12 & 3\\
H. Galois cohomology & 12 & $\cdot$ & $\cdot$ & $\cdot$ & $\cdot$ & $\cdot$ & 12 & 12 & $\cdot$ & $\cdot$ & $\cdot$ & 1 & 11 & $\cdot$\\
I. Galois representations & 17 & $\cdot$ & 10 & 2 & $\cdot$ & $\cdot$ & 9 & 7 & $\cdot$ & $\cdot$ & 3 & $\cdot$ & 12 & 5\\
J. Abelian varieties & 14 & $\cdot$ & $\cdot$ & 14 & $\cdot$ & $\cdot$ & $\cdot$ & $\cdot$ & $\cdot$ & $\cdot$ & $\cdot$ & 2 & 9 & 3\\
K. Algebraic and arithmetic geometry & 20 & $\cdot$ & $\cdot$ & 20 & $\cdot$ & $\cdot$ & $\cdot$ & $\cdot$ & $\cdot$ & $\cdot$ & $\cdot$ & 3 & 13 & 4\\
L. Records and the final step & 10 & \multicolumn{13}{c}{---}\\
\midrule
Total & 195 & 8 & 36 & 46 & 5 & 2 & 80 & 92 & 71 & 15 & 10 & 22 & 140 & 23\\
\bottomrule
\end{tabular}
\end{center}}

\noindent By confidence: 22 rows ``high'', 140 ``medium'', 23 ``low'', and ten rows of
the last group, which are lists and have no grade. In the development,
152 rows are a theorem cited by its name, 31 a theorem applied to other rows, three involve an
open statement directly, and nine are records. Of the rows of the first eleven groups, 20 carry
their codes with a star, and ten are never read.

Three things stand out. First, the covolume input (S6) receives 46 rows, among them all fourteen rows on
abelian varieties and all twenty rows of algebraic and arithmetic geometry, and no row of these two groups
is passed anywhere else. So the whole of this geometry is passed to one inequality, the covolume bound. Second, the inequality \eqref{gc:link} for Galois cohomology receives
80 rows, from nine of the eleven subjects; all twelve rows of Galois cohomology are passed there, and
otherwise only to the construction. Third, the inner-product formula (S3) receives only eight rows, and
the trace-formula statements only two; as far as we read, the other statements behind \S III.8 are obtained
in the development from theorems that are not among these rows.

\subsection*{4.\quad Discussion}\label{cs:disc}

\emph{What kind of object a row is.} The names say ``literature'', and the top theorem takes an argument
called the open literature statements. The rows are nevertheless not assumptions. The argument is a record
with three entries. Two are the same proposition: that a polarised regular algebraic cuspidal
representation of $GL_a$ over a CM field has a Galois representation with certain local properties. The
third is the existence of the local correspondence (row~65). At
the top of the development all three are supplied by theorems of the development. Apart from the records,
only eight rows (65, 66, 137, 138, 145--148) reach these entries, directly or through other rows; two of them
(147, 148) are never read. Every other row is
proved by naming a theorem of the development, in 152 cases with no reference to another row. A
search of the source text, made twice, found no axiom and no marker of an unfinished proof. We did not run
the proof checker or read the proofs. So the accurate description is: \emph{the development states its
versions of these theorems in its own language and contains proofs of them; this census records the
statements.} The versions differ from the published theorems in many rows, in both directions.

\emph{Rows that look like routine citations.} The 22 rows graded ``high'' (2, 13, 15--17, 22--24, 28, 35, 37, 38, 40, 63, 67, 122, 130, 152, 161, 170, 172, 174) are explicit
statements of standard theorems: for example Linnik's theorem (2),
Brauer's theorem (28) and the vanishing theorem of Kawamata and Viehweg
(174). Row~122 is the proposition ``True'';
its grade means only that the reading is certain. Most of the 140 rows graded ``medium'' are of the
same kind; their grade is lower because some notion in them was read by name, or because they are special
cases. In particular eleven rows (10--12, 29, 100, 123--128) are explicit and were first graded ``high'', but are
special cases: Chebotarev's theorem as a statement of existence, Chevalley's theorem for the group of
units, Green's theorem for prime fields, admissibility for unitary groups that come from a number field,
and Tate's local theorems for modules of $\ell$-power order.

\emph{Rows in a special form.} Many rows are the published theorem for a special class of objects. The
theorems on real groups (group~E) are stated for the groups $U(p,q)$ only, those on automorphic forms on
general linear groups for totally complex fields, and Green's theorems
(29--31) for prime fields. A common
restriction, above all in the group on theta lifts, is to the data of the construction:
rows~92, 109,
110, 117,
118 (and also 120,
149, and in a similar way 163,
164) are stated for one
particular kind of CM field (Galois over $\mathbb Q$, containing two given imaginary quadratic fields) and
for spaces that carry a unimodular lattice, not for Hermitian spaces over CM fields in general.
Rows~107, 116 and
121 assume that $E$ is Galois over $\mathbb Q$ and
unramified over its real subfield. Howe duality at the real places (119) is
stated for the trivial representation only. These are still citations; but a reader who checks one must
first confirm that the special case stated is the one printed.

\emph{Rows in a quantitative form.} Three rows of group~K carry explicit constants, and all three are in
print in the form stated: the divisor that fails to be klt at one point, with the condition
$\sum e/c_e\le1$ (row~175: Koll\'ar, \emph{Singularities
of pairs}, Theorem~6.4); the separation of two points, with the condition $\sum2^{1/e}e/c_e\le1$
(row~176: the same paper, Theorem~5.9); and the $L^2$ extension theorem
with the constant $34/\delta$ (row~177: Demailly's extension
theorem of 2015/16, for finitely many points). The papers, the constants and the theorem numbers are all
from memory, the checker's as well as ours; the theorem number that the checker gives for the last, 2.8,
is from memory too. The formula for $L'/L(1,\chi)$ with an explicit dependence on the modulus
(1) is a short consequence of a formula in Davenport's book. The bound for Sobolev norms in
the theorem of Casselman and Wallach (73) is not quantitative: it is what
the open mapping theorem gives.

\emph{Rows that could not be matched to print.} For these we know no printed statement of the same form, and
a reader should expect to have to extract the statement from published proofs: the bounds for cusp forms
uniform in the form (90, 91); the
bundle on special values with its bounds for denominators (27), which the
development derives from five other rows; the behaviour of theta functions under rational points
(120) and the proportionality of two theta integrals
(121), both specific to the construction, the second
derived in the development from rows~116 and~113; and the
statements in coordinates of group~K (173, 179, and
the numerical form of Mumford's theorem in 178). For
row~116 the attribution to Bernstein and Lapid is a guess from a name.

\emph{Weak forms of deep theorems.} Six rows cite some of the deepest theorems in the list and state much
less. Rows~65 and~66
state a weak form of the local Langlands correspondence: the dimension, the case $n=1$, compatibility with
segments and with unramified twists, and further laws for a restricted class of representations only.
Bijectivity is not stated, nor are the $L$- and $\varepsilon$-factors of pairs that characterise the
correspondence. Rows~145--148 state a weak form of the
theorems on Galois representations of regular algebraic cuspidal representations over CM fields and on
their behaviour at $\ell$: full compatibility with the local correspondence at the unramified places only;
elsewhere away from $\ell$ equal dimensions and, under a tameness condition, equal characteristic
polynomials on inertia; and at $\ell$ no isomorphism of Weil--Deligne representations. What is cited is
the full theorem; what is stated is a list of consequences. To a lesser degree the same holds for base
change (102, 103: nothing at ramified
inert places), for the statement on the Hodge--Tate map (150: no bound for
torsion), and for the criterion for irreducibility of a product
(53: a sufficient condition only).

\emph{Repetitions, and what is not here.} Three pairs of rows are one proposition under two names
(74, 75;
65, 66;
147, 148), and several rows are deduced from other
rows by a one-line proof. So the number of independent inputs is smaller than 195; we did not count it.
Theorems from the literature may also be used under other names inside the other proofs of the
development, and inside Mathlib, the library on which it is built; these are not in the list. And the
census does not say what any receiving proof does with a row.

{\setlength{\parindent}{0pt}


\section*{Computations}\addcontentsline{toc}{section}{Computations}

The numerical statements of Parts~I and~II were computed with the following scripts (Python, with
\texttt{mpmath} and \texttt{sympy}); each prints the numbers quoted. The sections of Part~III name their
own scripts where the numbers occur (Python and PARI/GP); those of \S\ref{ob:sec}, \S\ref{as:sec} and
\S\ref{el:sec} are in the folders \texttt{outline/examples/objects/}, \texttt{outline/examples/assembly/}
and \texttt{outline/examples/elliptic/}.

\begin{center}\small
\begin{tabular}{p{.34\textwidth}p{.58\textwidth}}
\texttt{e1\_running\_example.py} & $G$, $G_\chi$, $R_{q,d}$, the table for $q=11$, $d=1319$ (the first version used $q=3$, $d=23$; scripts \texttt{\_q3\_d23}), $P$ and $M$ (Part I, \S I.1, I.5, I.6; Part II, \S4.1--4.3, \S5.2--5.3, \S6.7)\\
\texttt{e1b\_check\_cells.py} & the eight cell constants, by a second method (Part II, \S5.2)\\
\texttt{e2\_lemniscate.py}, \texttt{e2b\_size.py} & the numbers $A_n$: divisibility and size (Part I, \S I.3; Part II, \S6.3)\\
\texttt{e4\_sigma.py} & the Weierstrass $\sigma$-function: Gaussian size without the divisibility (Part II, \S6.5)\\
\texttt{e5\_kronecker.py} & the two-variable function $\sigma(z+w)/(\sigma(z)\sigma(w))$ of the lemniscatic lattice: divisibility at inert primes and size of its mixed derivatives (an experiment; its numbers are not quoted in the text)\\
\texttt{e6\_hecke\_theta.py} & $T_n$, the period, the rational numbers $r_k$ (Part I, \S I.4.1)\\
\end{tabular}
\end{center}
}
\section*{Acknowledgements}\addcontentsline{toc}{section}{Acknowledgements}

We thank Manvir Jaswal for comments on an earlier version of this paper.


\phantomsection\addcontentsline{toc}{section}{References}
\begin{thebibliography}{GGP12}\small
\bibitem[Alp26]{Alp26} L.~Alp\"oge, \emph{[Remark on a proof of the Riemann hypothesis]}, with a formal proof in Lean 4, 20 November 2026. \url{https://alpo.ge/remark.pdf}.
\bibitem[Apo76]{Apo76} T.~M.~Apostol, \emph{Introduction to Analytic Number Theory}, Springer, 1976.
\bibitem[Apo90]{Apo90} T.~M.~Apostol, \emph{Modular Functions and Dirichlet Series in Number Theory}, 2nd ed., Graduate Texts in Math. 41, Springer, 1990.
\bibitem[ASD71]{ASD71} A.~O.~L.~Atkin and H.~P.~F.~Swinnerton-Dyer, Modular forms on noncongruence subgroups, in \emph{Combinatorics}, Proc. Sympos. Pure Math. 19, Amer. Math. Soc., 1971, 1--25.
\bibitem[Bar61]{Bar61} V.~Bargmann, On a Hilbert space of analytic functions and an associated integral transform, \emph{Comm. Pure Appl. Math.} 14 (1961), 187--214.
\bibitem[BC04]{BC04} J.~Bella\"iche and G.~Chenevier, Formes non temp\'er\'ees pour $U(3)$ et conjectures de Bloch--Kato, \emph{Ann. Sci. \'Ecole Norm. Sup.} 37 (2004), 611--662.
\bibitem[Beh46]{Beh46} F.~A.~Behrend, On sets of integers which contain no three terms in arithmetical progression, \emph{Proc. Nat. Acad. Sci. U.S.A.} 32 (1946), 331--332.
\bibitem[BK10]{BK10} K.~Bannai and S.~Kobayashi, Algebraic theta functions and the $p$-adic interpolation of Eisenstein--Kronecker numbers, \emph{Duke Math. J.} 153 (2010), 229--295.
\bibitem[BK16]{BK16} K.~Bannai and S.~Kobayashi, Integral structures on $p$-adic Fourier theory, \emph{Ann. Inst. Fourier} 66 (2016), 521--550.
\bibitem[Bom65]{Bom65} E.~Bombieri, On the large sieve, \emph{Mathematika} 12 (1965), 201--225.
\bibitem[Bom74]{Bom74} E.~Bombieri, Counting points on curves over finite fields (d'apr\`es S.~A.~Stepanov), S\'eminaire Bourbaki, exp.~430 (1973), Lecture Notes in Math. 383, Springer, 1974.
\bibitem[Bom87]{Bom87} E.~Bombieri, \emph{Le grand crible dans la th\'eorie analytique des nombres}, 2nd ed., Ast\'erisque 18, Soc. Math. France, 1987.
\bibitem[Bos96]{Bos96} J.-B.~Bost, P\'eriodes et isog\'enies des vari\'et\'es ab\'eliennes sur les corps de nombres (d'apr\`es D.~Masser et G.~W\"ustholz), S\'eminaire Bourbaki, exp.~795, \emph{Ast\'erisque} 237 (1996), 115--161.
\bibitem[BS23]{BS23} B.~Bhatt and P.~Scholze, Prismatic $F$-crystals and crystalline Galois representations, \emph{Camb. J. Math.} 11 (2023), 507--562.
\bibitem[Car12]{Car12} A.~Caraiani, Local-global compatibility and the action of monodromy on nearby cycles, \emph{Duke Math. J.} 161 (2012), 2311--2413.
\bibitem[Car14]{Car14} A.~Caraiani, Monodromy and local-global compatibility for $l=p$, \emph{Algebra Number Theory} 8 (2014), 1597--1646.
\bibitem[CH13]{CH13} G.~Chenevier and M.~Harris, Construction of automorphic Galois representations, II, \emph{Camb. J. Math.} 1 (2013), 53--73.
\bibitem[Che52]{Che52} P.~L.~Tchebichef, M\'emoire sur les nombres premiers, \emph{J. Math. Pures Appl.} 17 (1852), 366--390.
\bibitem[CS67]{CS67} S.~Chowla and A.~Selberg, On Epstein's zeta-function, \emph{J. reine angew. Math.} 227 (1967), 86--110.
\bibitem[Dav00]{Dav00} H.~Davenport, \emph{Multiplicative Number Theory}, 3rd ed., revised by H.~L.~Montgomery, Springer, 2000.
\bibitem[Dam70]{Dam70} R.~M.~Damerell, $L$-functions of elliptic curves with complex multiplication. I, \emph{Acta Arith.} 17 (1970), 287--301.
\bibitem[Deu41]{Deu41} M.~Deuring, Die Typen der Multiplikatorenringe elliptischer Funktionenk\"orper, \emph{Abh. Math. Sem. Hansischen Univ.} 14 (1941), 197--272.
\bibitem[DH66]{DH66} H.~Davenport and H.~Halberstam, The values of a trigonometrical polynomial at well spaced points, \emph{Mathematika} 13 (1966), 91--96.
\bibitem[Erd32]{Erd32} P.~Erd\H{o}s, Beweis eines Satzes von Tschebyschef, \emph{Acta Litt. Sci. Szeged} 5 (1932), 194--198.
\bibitem[FW79]{FW79} B.~Ferrero and L.~C.~Washington, The Iwasawa invariant $\mu_p$ vanishes for abelian number fields, \emph{Ann. of Math.} 109 (1979), 377--395.
\bibitem[Gal70]{Gal70} P.~X.~Gallagher, A large sieve density estimate near $\sigma=1$, \emph{Invent. Math.} 11 (1970), 329--339.
\bibitem[GGP12]{GGP12} W.~T.~Gan, B.~H.~Gross and D.~Prasad, Symplectic local root numbers, central critical $L$-values, and restriction problems in the representation theory of classical groups, \emph{Ast\'erisque} 346 (2012), 1--109.
\bibitem[GHY01]{GHY01} W.~T.~Gan, J.~P.~Hanke and J.-K.~Yu, On an exact mass formula of Shimura, \emph{Duke Math. J.} 107 (2001), 103--133.
\bibitem[Gra84]{Gra84} D.~R.~Grayson, Reduction theory using semistability, \emph{Comment. Math. Helv.} 59 (1984), 600--634.
\bibitem[Haz78]{Haz78} M.~Hazewinkel, \emph{Formal Groups and Applications}, Academic Press, 1978.
\bibitem[HB95]{HB95} D.~R.~Heath-Brown, A mean value estimate for real character sums, \emph{Acta Arith.} 72 (1995), 235--275.
\bibitem[Hon70]{Hon70} T.~Honda, On the theory of commutative formal groups, \emph{J. Math. Soc. Japan} 22 (1970), 213--246.
\bibitem[How79]{How79} R.~Howe, $\theta$-series and invariant theory, in \emph{Automorphic Forms, Representations and $L$-functions}, Proc. Sympos. Pure Math. 33, part 1, Amer. Math. Soc., 1979, 275--285.
\bibitem[Hur99]{Hur99} A.~Hurwitz, \"Uber die Entwicklungskoeffizienten der lemniskatischen Funktionen, \emph{Math. Ann.} 51 (1899), 196--226.
\bibitem[HW]{HW} G.~H.~Hardy and E.~M.~Wright, \emph{An Introduction to the Theory of Numbers}, 4th ed., Oxford University Press, 1960 (printing of 1975).
\bibitem[Jac86]{Jac86} H.~Jacquet, Sur un r\'esultat de Waldspurger, \emph{Ann. Sci. \'Ecole Norm. Sup.} 19 (1986), 185--229.
\bibitem[KP86]{KP86} J.~Kaczorowski and J.~Pintz, Oscillatory properties of arithmetical functions.~I, \emph{Acta Math. Hungar.} 48 (1986), 173--185.

\bibitem[Kat04]{Kat04} K.~Kato, $p$-adic Hodge theory and values of zeta functions of modular forms, \emph{Ast\'erisque} 295 (2004), 117--290.
\bibitem[Kat81]{Kat81} N.~M.~Katz, Divisibilities, congruences, and Cartier duality, \emph{J. Fac. Sci. Univ. Tokyo Sect. IA Math.} 28 (1981), 667--678.
\bibitem[Kud84]{Kud84} S.~S.~Kudla, Seesaw dual reductive pairs, in \emph{Automorphic Forms of Several Variables} (Katata, 1983), Progr. Math. 46, Birkh\"auser, 1984.
\bibitem[Kur92]{Kur92} M.~Kurihara, Some remarks on conjectures about cyclotomic fields and $K$-groups of $\mathbb Z$, \emph{Compositio Math.} 81 (1992), 223--236.
\bibitem[Lan05]{Lan05} E.~Landau, \"Uber einen Satz von Tschebyschef, \emph{Math. Ann.} 61 (1905), 527--550.
\bibitem[Leg08]{Leg08} A.-M.~Legendre, \emph{Essai sur la th\'eorie des nombres}, 2nd ed., Paris, 1808.
\bibitem[LT65]{LT65} J.~Lubin and J.~Tate, Formal complex multiplication in local fields, \emph{Ann. of Math.} 81 (1965), 380--387.
\bibitem[Luc78]{Luc78} \'E.~Lucas, Th\'eorie des fonctions num\'eriques simplement p\'eriodiques, \emph{Amer. J. Math.} 1 (1878), 184--240 and 289--321.
\bibitem[Mon71]{Mon71} H.~L.~Montgomery, \emph{Topics in Multiplicative Number Theory}, Lecture Notes in Math. 227, Springer, 1971.
\bibitem[MV07]{MV07} H.~L.~Montgomery and R.~C.~Vaughan, \emph{Multiplicative Number Theory I. Classical Theory}, Cambridge Univ. Press, 2007.
\bibitem[MW84]{MW84} B.~Mazur and A.~Wiles, Class fields of abelian extensions of $\mathbb Q$, \emph{Invent. Math.} 76 (1984), 179--330.
\bibitem[Neu99]{Neu99} J.~Neukirch, \emph{Algebraic Number Theory}, Grundlehren Math. Wiss. 322, Springer, 1999.
\bibitem[NSW00]{NSW00} J.~Neukirch, A.~Schmidt and K.~Wingberg, \emph{Cohomology of Number Fields}, 1st ed., Grundlehren Math. Wiss. 323, Springer, 2000.
\bibitem[Pag35]{Pag35} A.~Page, On the number of primes in an arithmetic progression, \emph{Proc. London Math. Soc.} 39 (1935), 116--141.
\bibitem[Ral87]{Ral87} S.~Rallis, \emph{$L$-functions and the Oscillator Representation}, Lecture Notes in Math. 1245, Springer, 1987.
\bibitem[Rub00]{Rub00} K.~Rubin, \emph{Euler Systems}, Ann. of Math. Studies 147, Princeton Univ. Press, 2000.
\bibitem[Sie35]{Sie35} C.~L.~Siegel, \"Uber die analytische Theorie der quadratischen Formen, \emph{Ann. of Math.} 36 (1935), 527--606.
\bibitem[Sou03]{Sou03} C.~Soul\'e, A bound for the torsion in the $K$-theory of algebraic integers, \emph{Doc. Math.} Extra Vol. Kato (2003), 761--788.
\bibitem[Sou79]{Sou79} C.~Soul\'e, $K$-th\'eorie des anneaux d'entiers de corps de nombres et cohomologie \'etale, \emph{Invent. Math.} 55 (1979), 251--295.
\bibitem[Ste69]{Ste69} S.~A.~Stepanov, On the number of points of a hyperelliptic curve over a finite prime field, \emph{Izv. Akad. Nauk SSSR Ser. Mat.} 33 (1969), 1171--1181.
\bibitem[Stu76]{Stu76} U.~Stuhler, Eine Bemerkung zur Reduktionstheorie quadratischer Formen, \emph{Arch. Math.} 27 (1976), 604--610.
\bibitem[Tit86]{Tit86} E.~C.~Titchmarsh, \emph{The Theory of the Riemann Zeta-function}, 2nd ed., revised by D.~R.~Heath-Brown, Oxford University Press, 1986.
\bibitem[TZ21]{TZ21} J.~Thorner and A.~Zaman, Refinements to the prime number theorem for arithmetic progressions, \emph{Math. Z.} 306 (2024), article 54; arXiv:2108.10878.
\bibitem[Was82]{Was82} L.~C.~Washington, \emph{Introduction to Cyclotomic Fields}, Graduate Texts in Math. 83, Springer, 1982.
\bibitem[Wei52]{Wei52} A.~Weil, Sur les ``formules explicites'' de la th\'eorie des nombres premiers, \emph{Comm. S\'em. Math. Univ. Lund}, tome suppl\'ementaire (1952), 252--265.
\bibitem[Wei64]{Wei64} A.~Weil, Sur certains groupes d'op\'erateurs unitaires, \emph{Acta Math.} 111 (1964), 143--211.
\bibitem[Wei65]{Wei65} A.~Weil, Sur la formule de Siegel dans la th\'eorie des groupes classiques, \emph{Acta Math.} 113 (1965), 1--87.

\bibitem[Wei76]{Wei76} A.~Weil, \emph{Elliptic Functions according to Eisenstein and Kronecker}, Springer, 1976.
\bibitem[Wey25]{Wey25} H.~Weyl, Theorie der Darstellung kontinuierlicher halb-einfacher Gruppen durch lineare Transformationen. I, \emph{Math. Z.} 23 (1925), 271--309.
\bibitem[WW27]{WW27} E.~T.~Whittaker and G.~N.~Watson, \emph{A Course of Modern Analysis}, 4th ed., Cambridge University Press, 1927.
\end{thebibliography}

\end{document}
